\documentclass[11pt]{article}
\usepackage[a4paper,margin=1.45cm]{geometry}
\usepackage[T1]{fontenc}
\usepackage[utf8]{inputenc}
\usepackage{lmodern}
\usepackage[spanish,es-noquoting,es-noshorthands]{babel}
\usepackage{amsmath,amssymb,mathtools,array,booktabs,pdflscape,adjustbox,diagbox}
\usepackage[protrusion=true,expansion=false]{microtype}
\setlength{\parindent}{1.25em}
\setlength{\parskip}{0pt}
\newcommand{\widebar}[1]{\overline{#1}}
\newcommand{\A}{\Delta}
\newcommand{\D}{\Delta}
\allowdisplaybreaks
\setlength{\emergencystretch}{12em}
\sloppy
\hbadness=10000
\vbadness=10000
\hfuzz=2pt
\vfuzz=10pt
\raggedbottom
\newcommand{\R}{\varrho}
\newcommand{\hata}{\widehat{a}}
\newcommand{\hatv}{\widehat{v}}
\newcommand{\modu}[1]{\;\operatorname{mod} #1}
\newcommand{\mcell}[1]{\ensuremath{\begin{gathered}#1\end{gathered}}}
\newcommand{\pair}[2]{(#1\,|\,#2)}
\newcommand{\eps}{\varepsilon}

\providecommand{\Hgrp}{\mathfrak{H}}
\providecommand{\GaloisRes}{\Phi}
\providecommand{\pol}[2]{\left(#1\frac{\partial}{\partial #2}\right)}
\providecommand{\Deltaop}{\Delta\!\left(\frac{\partial}{\partial x},\frac{\partial}{\partial y},\frac{\partial}{\partial z}\right)}
\providecommand{\Omegaop}{\Omega}

\makeatletter
\renewcommand\footnotesize{\@setfontsize\footnotesize{7.5}{8.7}\selectfont}
\makeatother
\newcommand{\bdelta}{\bar{\delta}}
\providecommand{\dd}{\mathrm d}
\providecommand{\Div}{\operatorname{Div}}
\providecommand{\G}{\mathfrak G}
\providecommand{\bd}{\bar{\delta}}
\providecommand{\p}{\partial}
\providecommand{\ideal}[1]{\mathfrak{#1}}
\providecommand{\eqzero}[1]{\equiv 0\,(#1)}
\providecommand{\gtr}{>}
% Paper 24 cumulative macro support
\providecommand{\widebar}[1]{\overline{#1}}
\providecommand{\A}{\Delta}
\providecommand{\D}{\Delta}
\providecommand{\R}{\varrho}
\providecommand{\hata}{\widehat{a}}
\providecommand{\hatv}{\widehat{v}}
\providecommand{\modu}[1]{\;\operatorname{mod} #1}
\providecommand{\mcell}[1]{\ensuremath{\begin{gathered}#1\end{gathered}}}
\providecommand{\pair}[2]{(#1\,|\,#2)}
\providecommand{\eps}{\varepsilon}
\providecommand{\rhoidx}{\varrho}
\providecommand{\dd}{\mathrm d}
\providecommand{\Div}{\operatorname{Div}}
\providecommand{\G}{\mathfrak G}
\providecommand{\bd}{\bar{\delta}}
\providecommand{\p}{\partial}
\providecommand{\frS}{\mathfrak S}
\providecommand{\frp}{\mathfrak p}
\providecommand{\fra}{\mathfrak a}
\providecommand{\frb}{\mathfrak b}
\providecommand{\frc}{\mathfrak c}
\providecommand{\frf}{\mathfrak f}
\providecommand{\frr}{\mathfrak r}
\providecommand{\frm}{\mathfrak m}
\providecommand{\frD}{\mathfrak D}
\providecommand{\calA}{\mathcal A}
\providecommand{\calB}{\mathcal B}
\providecommand{\calN}{\mathcal N}
\providecommand{\mM}{\mathfrak{M}}
\providecommand{\mN}{\mathfrak{N}}
\providecommand{\mL}{\mathfrak{L}}
\providecommand{\mG}{\mathfrak{G}}
\providecommand{\mA}{\mathfrak{A}}
\providecommand{\mB}{\mathfrak{B}}
\providecommand{\mX}{\mathfrak{X}}
\providecommand{\mY}{\mathfrak{Y}}
\providecommand{\ma}{\mathfrak{a}}
\providecommand{\mb}{\mathfrak{b}}
\providecommand{\mx}{\mathfrak{x}}
\providecommand{\my}{\mathfrak{y}}
\providecommand{\mz}{\mathfrak{z}}
\providecommand{\me}{\mathfrak{e}}
\providecommand{\mbarA}{\overline{\mathfrak{A}}}
\providecommand{\mbarB}{\overline{\mathfrak{B}}}
\providecommand{\bara}{\overline{\mathfrak{a}}}
\providecommand{\mU}{\mathfrak{U}}
\providecommand{\mP}{\mathfrak{P}}
\providecommand{\mQ}{\mathfrak{Q}}
\providecommand{\mS}{\mathfrak{S}}
\providecommand{\mD}{\mathfrak{D}}
\providecommand{\mR}{\mathfrak{R}}
\providecommand{\mV}{\mathfrak{V}}
\providecommand{\ideal}[1]{\mathfrak{#1}}
\providecommand{\Amod}{\mathfrak A}
\providecommand{\Bmod}{\mathfrak B}
\providecommand{\Gmod}{\mathfrak G}
\providecommand{\Rmod}{\mathfrak R}
\providecommand{\Smod}{\mathfrak S}
\providecommand{\Mmod}{\mathfrak M}
\providecommand{\Nmod}{\mathfrak N}
\providecommand{\mideal}{\mathfrak m}
\providecommand{\nideal}{\mathfrak n}
\providecommand{\gideal}{\mathfrak g}
\providecommand{\hideal}{\mathfrak h}
\providecommand{\jideal}{\mathfrak j}
\providecommand{\tideal}{\mathfrak t}
\providecommand{\aideal}{\mathfrak a}
\providecommand{\bideal}{\mathfrak b}
\providecommand{\bb}{\mathfrak b}
\providecommand{\cc}{\mathfrak c}
\providecommand{\zero}[1]{\equiv 0\pmod{#1}}
\providecommand{\Norm}{N}
\providecommand{\frakS}{\mathfrak S}
\providecommand{\frakM}{\mathfrak M}
\providecommand{\frakG}{\mathfrak G}
\providecommand{\frakm}{\mathfrak m}
\providecommand{\frakn}{\mathfrak n}
\providecommand{\frakp}{\mathfrak p}
\providecommand{\frakq}{\mathfrak q}
\providecommand{\frakr}{\mathfrak r}
\providecommand{\fraka}{\mathfrak a}
\providecommand{\frakb}{\mathfrak b}
\providecommand{\frakc}{\mathfrak c}
\providecommand{\frakd}{\mathfrak d}
\providecommand{\frake}{\mathfrak e}
\providecommand{\frakg}{\mathfrak g}
\providecommand{\fraks}{\mathfrak s}
\providecommand{\frako}{\mathfrak o}
\providecommand{\frK}{\mathfrak K}
\providecommand{\barfrakm}{\overline{\mathfrak m}}
\providecommand{\frakh}{\mathfrak h}
\providecommand{\frakt}{\mathfrak t}
\providecommand{\barfraka}{\overline{\mathfrak a}}
\providecommand{\barfrakb}{\overline{\mathfrak b}}
\providecommand{\barfrakc}{\overline{\mathfrak c}}
\providecommand{\calP}{\mathcal P}
\providecommand{\calR}{\mathcal R}
\providecommand{\calG}{\mathcal G}
\providecommand{\calH}{\mathcal H}
\providecommand{\barr}{\overline}
\providecommand{\tagg}[1]{\tag{#1}}



% Paper 25--29 cumulative macro support
\providecommand{\frakR}{\mathfrak R}
\providecommand{\frakT}{\mathfrak T}
\providecommand{\frakK}{\mathfrak K}
\providecommand{\frakL}{\mathfrak L}
\providecommand{\frakH}{\mathfrak H}
\providecommand{\frakP}{\mathfrak P}
\providecommand{\frakQ}{\mathfrak Q}
\providecommand{\frakZ}{\mathfrak Z}
\providecommand{\frakN}{\mathfrak N}
\providecommand{\frakB}{\mathfrak B}
\providecommand{\frakF}{\mathfrak F}
\providecommand{\frakW}{\mathfrak W}
\providecommand{\Trdeg}{\operatorname{trdeg}}
\providecommand{\Gal}{\operatorname{Gal}}
\providecommand{\Norm}{\operatorname{N}}


% Paper 30 macro support
\providecommand{\widebar}[1]{\overline{#1}}
\providecommand{\mR}{\mathfrak{R}}
\providecommand{\mK}{\mathfrak{K}}
\providecommand{\mL}{\mathfrak{L}}
\providecommand{\mS}{\mathfrak{S}}
\providecommand{\mT}{\mathfrak{T}}
\providecommand{\mM}{\mathfrak{M}}
\providecommand{\mN}{\mathfrak{N}}
\providecommand{\mA}{\mathfrak{A}}
\providecommand{\mB}{\mathfrak{B}}
\providecommand{\mC}{\mathfrak{C}}
\providecommand{\mD}{\mathfrak{D}}
\providecommand{\mE}{\mathfrak{E}}
\providecommand{\mF}{\mathfrak{F}}
\providecommand{\mG}{\mathfrak{G}}
\providecommand{\mH}{\mathfrak{H}}
\providecommand{\mU}{\mathfrak{U}}
\providecommand{\mV}{\mathfrak{V}}
\providecommand{\mbar}[1]{\overline{#1}}
\providecommand{\ma}{\mathfrak{a}}
\providecommand{\mb}{\mathfrak{b}}
\providecommand{\mc}{\mathfrak{c}}
\providecommand{\md}{\mathfrak{d}}
\providecommand{\me}{\mathfrak{e}}
\providecommand{\mf}{\mathfrak{f}}
\providecommand{\mm}{\mathfrak{m}}
\providecommand{\mn}{\mathfrak{n}}
\providecommand{\mo}{\mathfrak{o}}
\providecommand{\mpideal}{\mathfrak{p}}
\providecommand{\mq}{\mathfrak{q}}
\providecommand{\mt}{\mathfrak{t}}
\providecommand{\mv}{\mathfrak{v}}
\providecommand{\frakp}{\mathfrak p}
\providecommand{\frakq}{\mathfrak q}
\providecommand{\modu}[1]{\;\operatorname{mod} #1}
\begin{document}

\begin{center}
{\Large\bfseries 1. Sobre la formación del sistema de formas de la forma\\ ternaria bicuadrática}\par
\vspace{1.5em}
\emph{Sitzungsberichte der Physikalisch-medizinischen Sozietät in Erlangen} 39 (1907), pp. 176--179
\end{center}

\vspace{3em}
\begin{center}
(Extracto de la disertación de la autora.)
\end{center}

Los trabajos de Gordan, Maisano y Pascal\footnote{P. Gordan, ``Über das volle Formensystem der ternären biquadratischen Form'' \(f=x_1^3x_2+x_2^3x_3+x_3^3x_1\) (Math. Annalen, vol. XVII, pp. 217--233, 1880); G. Maisano, 1. ``Sistemi completi dei primi cinque gradi della forma ternaria biquadratica e degl' invarianti, covarianti e contravarianti di sesto grado'' (Giorn. di Battaglini XIX). 2. ``Sui covarianti indipendenti di \(6^\circ\) grado nei coefficienti della forma biquadratica ternaria'' (Rend. Circ. Mat. di Palermo I, 1887); E. Pascal, ``Contributo alla teoria della forma ternaria biquadratica e delle sue varie decomposizioni in fattori'' (memoria premiada por la R. Accademia delle scienze fisiche e matematiche di Napoli, 1905).} se ocupan del sistema de formas de la forma ternaria bicuadrática. Gordan establece, para la forma automorfa especial
\[
f=x_1^3x_2+x_2^3x_3+x_3^3x_1,
\]
el sistema completo de formas, compuesto de 54 formaciones, sobre la base de principios análogos a los que había dado para los sistemas de formas en el dominio binario.

En el trabajo de Maisano, para la forma bicuadrática general, se establecen las formas hasta el quinto orden inclusive\footnote{Por ``orden'' debe entenderse la dimensión en los coeficientes, y por ``grado'' la dimensión en las variables.}, junto con algunos invariantes, covariantes y contravariantes de orden superior, según el método empleado por Gordan en el volumen I de los \emph{Mathematische Annalen} para la forma ternaria cúbica. Pascal, utilizando los resultados de Maisano, se ocupa principalmente de la cuestión de la descomposición de la forma bicuadrática en factores\footnote{En su segunda nota breve Maisano intenta demostrar una dependencia lineal entre los tres covariantes de orden 6 y grado 6. Frente a esta demostración no del todo completa, Pascal cree haber probado la independencia lineal de los tres covariantes de orden 6, así como la de los tres invariantes de orden 9. Pero que en efecto existe una relación lineal entre los tres covariantes, por una parte, y los tres invariantes, por otra, resultó de mis cálculos; en la notación de Pascal, las fórmulas explícitas son
\[
0=20\Omega_1+6\Omega_2-3\Omega_3+4fC_1-12fC_2-4A\cdot\Delta,
\]
\[
0=4A^3-15AB+30C-90D+9E.
\]}.

\textbf{El objetivo de mis investigaciones es establecer el sistema de formas de la forma ternaria bicuadrática general; por lo pronto se dan sólo los fundamentos principales y se establece un llamado ``sistema relativamente completo''\footnote{Para la terminología, véase Gordan--Kerschensteiner, \emph{Vorlesungen über Invariantentheorie}, vol. II, p. 227.}.}

La idea básica para la formación de sistemas de formas ternarias es la misma que en el dominio binario. Partiendo de un primer sistema relativamente completo, a saber, el sistema de formas tomado del caso binario, se pasa, según una ley determinada, a sistemas con módulo cada vez más alto, hasta que o bien el sistema de un módulo, tomado éste como forma fundamental, se vuelve finito y conocido, o bien un módulo puede reducirse a formas que tienen invariantes como factores. Por transvección del sistema relativamente completo con el sistema del módulo, en el primer caso se obtiene el sistema absolutamente completo; en el segundo, el sistema relativamente completo y el absolutamente completo coinciden. Por la demostración general de Hilbert de la finitud de los sistemas de formas, este procedimiento debe necesariamente llegar a término.

En nuestro caso, el módulo \((abc)\) del primer sistema relativamente completo puede reducirse a los módulos
\[
\Delta=(abc)^2a_x^2b_x^2c_x^2\quad\text{y}\quad \nu=(abu)^4;
\]
de aquí se obtiene fácilmente el sistema relativamente completo módulo \(\nu\). Como sucesión de módulos elegimos ahora las formas
\[
\nu;\qquad \nu^{(\nu)}=(\nu\nu_1x)^4=s_x^4,
\]
\[
\nu^{(s)}=(ss'u)^4\ldots,
\]
y añadimos, como módulos adicionales, dos formas cuadráticas que aparecen en la formación del sistema relativamente completo módulo \(s\), a saber, \(u_\sigma^2\) y \(t_x^2\). Entonces se puede mostrar que el módulo que sigue a \(s\), es decir \((ss'u)^4\), es reducible al módulo \((\varrho,t)\). Pero como el sistema simultáneo de dos formas cuadráticas es finito y conocido, se alcanza así el final descrito arriba. Como sistema relativamente completo módulo \((\varrho,t)\) se obtienen 331 formaciones.

Añado algunos detalles sobre el método empleado.

El proceso fundamental para producir formas, el proceso de plegamiento, puede definirse en el dominio ternario de la siguiente manera.

Si en un producto simbólico
\[
s_x^m t_x^n u_\sigma^\mu u_\tau^\nu
\]
se reemplazan los pares de factores
\[
\begin{array}{c|c|c|c|c}
 & s_xt_x & u_\sigma u_\tau & s_xu_\sigma\ \text{o}\ s_xu_\tau & t_xu_\tau\ \text{o}\ t_xu_\sigma\\
\text{respectivamente por} & (stu) & (\sigma\tau x) & s_\sigma\ \text{o}\ s_\tau & t_\tau\ \text{o}\ t_\sigma\\
\text{plegamiento} & \mathrm{I} & \mathrm{II} & \mathrm{III} & \mathrm{IV},
\end{array}
\]
entonces las formas así obtenidas han surgido de la forma original por plegamiento\footnote{Gordan, Math. Annalen, vol. XVII, p. 219.}.

Aquí todavía se puede poner siempre \(s_\sigma=0;\ t_\tau=0\). Para la relación entre los plegamientos individuales vale el siguiente teorema:

\textbf{Los plegamientos I y II son plegamientos fundamentales, a partir de los cuales, independientemente del orden de composición, pueden componerse los plegamientos III y IV. En otras palabras: para formar todas las formas que surgen de una expresión dada por plegamiento, sólo hay que aplicar los plegamientos I y II}\footnote{El teorema tiene una excepción en el caso en que \(n\) y \(\mu\), respectivamente \(m\) y \(\nu\), se anulan simultáneamente; la ``sucesión de formas'' se reduce entonces al término diagonal.}.

Por una ``sucesión de formas'' entendemos una forma inicial junto con todas las formas que surgen de ella por plegamiento en sí misma. Por el teorema, una sucesión de formas
\[
s_x^m t_x^n u_\sigma^\mu u_\tau^\nu
\]
puede disponerse en un esquema rectangular. Al avanzar en él se obtiene:
\begin{enumerate}
\item al pasar una columna hacia la derecha, todas las formas que surgen por plegamiento I de las formas adyacentes;
\item al pasar una fila hacia abajo, todas las formas que surgen por plegamiento II de las formas situadas encima, y por tanto todas las formas que surgen por plegamiento.
\end{enumerate}

Bajo estas hipótesis, los teoremas sobre reductores pueden formularse en su forma más general, con la definición
\begin{center}
``Un reductor es una sucesión de formas reducible,''
\end{center}
del siguiente modo:

\begin{center}
\textbf{Si la forma inicial de una sucesión de formas es reducible porque uno de sus términos ha surgido por plegamiento con un reductor, y si la forma final de la sucesión de formas ha surgido de ese mismo término por plegamiento, entonces toda la sucesión de formas es reducible.}
\end{center}

Como métodos adicionales de reducción hay que mencionar:
\begin{enumerate}
\item la llamada ``doble reducción'', es decir, la reducción de una forma de dos maneras para obtener una relación entre las formas superiores;
\item el ``plegamiento con formas descompuestas'', que en parte tiene el carácter de reducción por reductores y en parte el de ``doble reducción''.
\end{enumerate}

\clearpage

\begin{center}
{\Large\bfseries 2. Sobre la formación del sistema de formas de la forma\\ bicuadrática ternaria}\par
\vspace{1.5em}
Journal f. d. reine u. angew. Math. 134 (1908), pp. 23--90 y dos tablas
\end{center}

\vspace{3em}

\begin{center}
\textbf{Introducción.}
\end{center}

Del sistema de formas de la forma bicuadrática ternaria tratan trabajos de Gordan, Maisano y Pascal\footnote{Un breve extracto de la introducción y del capítulo I apareció en los \emph{Sitzungsberichte der phys.-med. Sozietät Erlangen 1907}, pp. 176--179.}\footnote{P. Gordan, sobre el sistema completo de formas de la forma bicuadrática ternaria
\[
f=x_1^3x_2+x_2^3x_3+x_3^3x_1.
\]
(\emph{Math. Annalen}, vol. XVII (1880), pp. 217--233.)
G. Maisano, 1) \emph{Sistemi completi dei primi cinque gradi della forma ternaria biquadratica e degl' invarianti, covarianti e contravarianti di sesto grado}. (Giorn. di Battaglini XIX.)
2) \emph{Sui covarianti indipendenti di $6^\circ$ grado nei coefficienti della forma biquadratica ternaria}. (Rend. Circ. Mat. di Palermo I, 1887.)
E. Pascal, \emph{Contributo alla teoria della forma ternaria biquadratica e delle sue varie decomposizioni in fattori}. (Memoria premiada por la R. Accademia delle scienze fisiche e matematiche di Napoli. 1905.)}. El señor Gordan establece el sistema completo de formas, compuesto de 54 formaciones, de la forma automorfa especial:
\[
f=x_1^3x_2+x_2^3x_3+x_3^3x_1
\]
sobre la base de principios análogos a los que había dado para los sistemas de formas en el dominio binario.

En el señor Maisano se establecen, para la forma bicuadrática general, las formas hasta el quinto orden\footnote{Por ``orden'' debe entenderse la dimensión en los coeficientes, y por ``grado'' la dimensión en las variables.} inclusive, así como algunos invariantes, covariantes y contravariantes de orden superior, según el método aplicado por el señor Gordan en el volumen I de los \emph{Math. Annalen} a la forma cúbica ternaria. El señor Pascal, usando los resultados de Maisano, se ocupa principalmente de la cuestión de la descomposición de la forma bicuadrática en factores\footnote{En la segunda nota breve, el señor Maisano intenta demostrar una dependencia lineal de los tres covariantes de sexto orden y sexto grado. Frente a esta demostración no del todo completa, el señor Pascal cree haber demostrado la independencia lineal de los tres covariantes de sexto orden, así como la de los tres invariantes de noveno orden. Pero que en efecto existe una relación lineal entre los tres covariantes, por una parte, y los tres invariantes, por otra, lo muestran las fórmulas (13), § 11, y (22), nota del § 17, del presente trabajo, donde esas relaciones se dan explícitamente. Según una comunicación del señor Pascal, la contradicción se explica por un error numérico en la expresión de su contravariante $p$ (aquí llamada $\varrho$), p. 46 de su memoria.}.

\emph{El objetivo de las investigaciones que siguen es establecer el sistema de formas de la forma bicuadrática ternaria general; más precisamente, en este trabajo sólo se dan los fundamentos principales y se establece un llamado ``sistema relativamente completo''\footnote{Para la denominación, véase § 3a.}.}

El trabajo se vincula estrechamente con el trabajo de Gordan; no obstante, los principios que allí sólo se daban en forma muy general debían primero desarrollarse en detalle. Con ayuda de un teorema sobre la conexión de los plegamientos y mediante la introducción de la ``sucesión de formas'' (§ 1), los teoremas de reducción se formulan de manera precisa y se completan (§ 3), mientras que el establecimiento recurrente de desarrollos en serie especiales (§ 2) da el medio de cálculo para realizar efectivamente las reducciones.

La idea fundamental de la formación de sistemas de formas ternarias es la misma que en el dominio binario. Partiendo de un primer sistema relativamente completo ---el sistema de las formas tomadas del caso binario--- se llega, según una ley determinada, a sistemas con módulo cada vez más alto, hasta que el sistema de un módulo ---tomado el módulo como forma fundamental--- se vuelva finito y conocido, o bien hasta que un módulo pueda reducirse a formas que tengan invariantes como factores. Por transvección del sistema relativamente completo con el sistema del módulo surge, en el primer caso, el sistema absolutamente completo, mientras que en el segundo caso el sistema relativamente completo y el sistema absolutamente completo se vuelven idénticos. A consecuencia de la finitud de los sistemas de formas, demostrada en general por el señor Hilbert, este procedimiento debe necesariamente terminar.

En nuestro caso, el módulo $(abc)$ del primer sistema relativamente completo (§ 4) se reduce a los módulos
\[
\Delta=(abc)^2a_x^2b_x^2c_x^2\quad\text{y}\quad \nu=(abu)^4
\]
(§ 5), y luego se encuentra el sistema relativamente completo módulo $\nu$ (§ 6). Estos resultados ya se dan en el trabajo de Gordan para la forma bicuadrática general; sin embargo, nuestro modo de derivación se transporta más fácilmente a formas de grado superior.

Como sucesión de módulos elegimos ahora las formas (§ 7):
\[
\nu,\qquad \nu(\nu)=(\nu\nu_1x)^4=s_x^4,\qquad \nu(s)=(ss'u)^4,\ldots .
\]

En la formación del sistema relativamente completo módulo $s$, se adjuntan como módulos dos formas cuadráticas que aparecen en el sistema, $u_\varrho^2$ y $t_x^2$ (§§ 9 y 10), y por consiguiente se forma el sistema relativamente completo módulo $(s,\varrho,t)$; después se muestra (§ 17) que el módulo $(ss'u)^4$ es reducible a los módulos $\varrho$ y $t$. De ese modo, el sistema inmediatamente superior pasa a un sistema relativamente completo módulo $(\varrho,t)$. Pero como el sistema simultáneo de dos formas cuadráticas es finito y conocido, y en el caso general consta de 20 formaciones\footnote{Cf. Clebsch--Lindemann, \emph{Vorlesungen über Geometrie}. (vol. I, p. 288 ss.) Sistema de dos formas cogredientes. Gordan, \emph{Über Büschel von Kegelschnitten}. (\emph{Math. Ann.} vol. XIX, p. 530.) Sistema de dos formas contragredientes.}, con el establecimiento del sistema relativamente completo módulo $(\varrho,t)$ se alcanza la terminación caracterizada arriba. Queda reservada la transvección de este sistema con el sistema de $(\varrho,t)$ para formar el sistema absolutamente completo.

Como sistema relativamente completo módulo $(\varrho,t)$ se encuentran 331 formaciones, ordenadas en la tabla adjunta según su grado en las variables $x$ y $u$.

Finalmente damos todavía una visión general de los símbolos introducidos:
\begin{align*}
f&=a_x^4,\\
\theta&=\theta_x^4u_\vartheta^2=(abu)^2a_x^2b_x^2,\\
K&=K_x^6u_k^3=(a\theta u)u_\vartheta^2a_x^3\theta_x^3,\\
j&=u_j^6=(a\theta u)^4u_\vartheta^2,\\
\Delta&=\Delta_x^6=a_\vartheta^2a_x^2\theta_x^4=(abc)^2a_x^2b_x^2c_x^2,\\
N&=N_x^8u_\eta=(a\Delta u)a_x^3\Delta_x^5,\\
\nu&=u_\nu^4=(abu)^4,\\
H&=H_x^2u_\eta^4=(\nu\nu_1x)^2u_\nu^2u_{\nu_1}^2,\\
L&=L_x^3u_l^6=(\nu\eta x)H_x^2u_\nu^3u_\eta^3,\\
g&=g_x^6=(\nu\eta x)^4H_x^2,\\
\sigma&=u_\sigma^6=H_\nu^2u_\nu^2u_\eta^4,\\
s&=s_x^4=(\nu\nu_1x)^4,\\
Z&=Z_x^4u_\zeta^2=(ss'u)^2s_x^2s_x'^2,\\
\varrho&=u_\varrho^2=\theta_\nu^4u_\vartheta^2,\\
t&=t_x^2=a_\eta^4H_x^2,\\
i&=a_\nu^4,\\
J&=s_\nu^4.
\end{align*}

\begin{center}
{\Large\bfseries Capítulo I. \quad Teoremas generales sobre formas ternarias.}
\end{center}

\begin{center}
\textbf{§ 1. \quad Proceso de plegamiento. \quad Sucesiones de formas.}
\end{center}

El proceso fundamental para producir formas es el proceso de plegamiento, que en el dominio ternario se define de la siguiente manera. Sea dado un producto simbólico:
\[
s_x^m t_x^n u_\sigma^\mu u_\tau^\nu .
\]
Si se reemplazan los pares de factores
\[
s_xt_x\qquad u_\sigma u_\tau\qquad s_xu_\sigma\ \text{o}\ s_xu_\tau
\qquad t_xu_\tau\ \text{o}\ t_xu_\sigma
\]
respectivamente por
\[
(stu)\qquad (\sigma\tau x)\qquad s_\sigma\ \text{o}\ s_\tau
\qquad t_\tau\ \text{o}\ t_\sigma,
\]
plegamiento
\[
\text{I}\qquad\text{II}\qquad\text{III}\qquad\text{IV},
\]
entonces las formas resultantes han surgido de la forma originaria por plegamiento\footnote{Gordan, \emph{Math. Annalen} vol. XVII, p. 219.}.

La expresión $s_x^m t_x^n u_\sigma^\mu u_\tau^\nu$ puede concebirse ya sea como el producto efectivo de dos formas $S\cdot T$, ya sea como una sola forma con varias series de símbolos:
\[
s_x^m t_x^n u_\sigma^\mu u_\tau^\nu=A_x^{m+n}u_x^{\mu+\nu}.
\]
En el primer caso hablamos del plegamiento de la forma $S$ con $T$, en el segundo del ``plegamiento de la forma $A$ en sí misma''. Para abreviar, en lo que sigue supondremos siempre (lo cual de hecho ocurre para todas las formas consideradas más adelante) que
\[
\tag{a}
s_\sigma^\lambda=0,\qquad t_\tau^\lambda=0
\]
para todo $\lambda$ no nulo y cualquier $x$, y junto con cualesquiera otros plegamientos. Esto dice que la forma $s_x^m t_y^n u_\sigma^\mu u_\tau^\nu$ es una llamada ``forma normal'', es decir, que se anula al aplicar el proceso $\Omega$, tanto con respecto a las variables $x$ y $u$ como con respecto a las variables $y$ y $v$.

Por tanto, la condición (a) no representa ninguna restricción, sino que siempre puede alcanzarse\footnote{Cf. Gordan, \emph{Math. Annalen} vol. V, p. 104.}.

Para la conexión entre los distintos plegamientos vale el siguiente teorema:

\emph{Teorema I. Los plegamientos I y II son plegamientos fundamentales a partir de los cuales, independientemente del orden de composición, pueden componerse los plegamientos III y IV. En otras palabras: para formar todas las formas que surgen por plegamiento de una expresión dada, basta aplicar los plegamientos I y II.}

Demostración: por el teorema de identidad, respectivamente por el teorema del producto para matrices, teniendo en cuenta (a), se obtiene
\[
(\widehat{st}\sigma x)=-t_\sigma,\qquad
(\widehat{\sigma\tau}su)=-s_\tau,
\]
\[
(\widehat{st}\tau x)=s_\tau,\qquad
(\widehat{\sigma\tau}tu)=t_\sigma,
\]
\[
(stu)(\sigma\tau x)=s_\tau+t_\sigma-u_x\,s_\tau t_\sigma .
\]
Así, por composición de los plegamientos I y II surgen los plegamientos III y IV, tanto cuando se aplica el plegamiento II al plegamiento I, esto es, a los factores $(stu)u_\sigma u_\tau$, como cuando se aplica el plegamiento I al plegamiento II, a los factores $(\sigma\tau x)s_xt_x$.

Llamamos \emph{sucesión de formas}\footnote{Cf. Clebsch, \emph{Über eine Fundamentalaufgabe der Invariantentheorie} (Abh. der Gött. Ges. d. Wiss. vol. XVII): § 17 y final del § 18. La ``sucesión de formas'' difiere del ``sistema equivalente propiamente reducido'' introducido allí por el principio de ordenación según formas superiores.} a una forma inicial junto con todas las formas que han surgido de ella por plegamiento en sí misma, y designamos la sucesión de formas por la forma inicial. Aquí se entenderá por ``forma superior'' toda forma con más plegamiento en sí misma; entre los plegamientos igualmente legítimos $s_\tau$ y $t_\sigma$, uno ha de normalizarse como superior. De la ley de formación de la sucesión de formas se sigue:

1) Cada forma de la sucesión de formas es lineal en los símbolos de la forma inicial.

2) La sucesión de formas está determinada unívocamente para formas iniciales que tienen dos series de símbolos contragredientes cada una; para formas iniciales con más series de símbolos debe normalizarse unívocamente distinguiendo plegamientos especiales entre aquellos que están conectados por el teorema de identidad.

Por el teorema I, una sucesión de formas $s_x^m t_x^n u_\sigma^\mu u_\tau^\nu$ ($m>n;\ \mu>\nu$) puede disponerse según el siguiente esquema rectangular:
\begingroup\small
\[
\begin{array}{c|c|c|c}
s_x^m t_x^n u_\sigma^\mu u_\tau^\nu
& (stu)
& (stu)^2
& \cdots\\[0.3em]
(\sigma\tau x)
& t_\sigma,\ s_\tau
& t_\sigma(stu),\ s_\tau(stu)
& \cdots\\[0.3em]
(\sigma\tau x)^2
& t_\sigma(\sigma\tau x),\ s_\tau(\sigma\tau x)
& t_\sigma^2,\ t_\sigma s_\tau,\ s_\tau^2
& \cdots\\[0.3em]
\vdots&\vdots&\vdots&\ddots\\[0.3em]
(\sigma\tau x)^\nu
& t_\sigma(\sigma\tau x)^{\nu-1},\ s_\tau(\sigma\tau x)^{\nu-1}
& t_\sigma^2(\sigma\tau x)^{\nu-2},\ t_\sigma s_\tau(\sigma\tau x)^{\nu-2},\ s_\tau^2(\sigma\tau x)^{\nu-2}
& \cdots\\[0.3em]
\vdots
& t_\sigma(\sigma\tau x)^\nu
& t_\sigma^2(\sigma\tau x)^{\nu-1},\ t_\sigma s_\tau(\sigma\tau x)^{\nu-1}
& \cdots
\end{array}
\]
\endgroup
y la parte inferior correspondiente, continuada como sigue\footnote{Aquí, y análogamente en lo que sigue, la parte inferior marcada con un asterisco debe colocarse en el lugar igualmente marcado, arriba a la derecha del esquema.}
\begingroup\small
\[
\begin{array}{c|c|c}
(stu)^n & \cdots & \cdots\\[0.3em]
t_\sigma(stu)^{n-1},\ s_\tau(stu)^{n-1}
& s_\tau(stu)^n & \cdots\\[0.3em]
t_\sigma^2(stu)^{n-2},\ t_\sigma s_\tau(stu)^{n-2},\ s_\tau^2(stu)^{n-2}
& t_\sigma s_\tau(stu)^{n-1},\ s_\tau^2(stu)^{n-1}
& s_\tau^2(stu)^n .
\end{array}
\]
\endgroup

Aquí se obtiene, cuando se avanza

1) una columna hacia la derecha, todas las formas que surgen por plegamiento I de las formas adyacentes,

2) una fila hacia abajo, todas las formas que surgen por plegamiento II de las formas situadas encima,

y por tanto todas las formas que surgen por plegamiento.

Una forma cualquiera de la sucesión total, $t_\sigma^\kappa s_\tau^\lambda(stu)^\mu$ o $t_\sigma^\kappa s_\tau^\lambda(\sigma\tau x)^\nu$, constituye el comienzo de una nueva sucesión de formas, compuesta por todas aquellas formas del esquema rectangular que tienen $t_\sigma^\kappa s_\tau^\lambda(stu)^\mu$ o $t_\sigma^\kappa s_\tau^\lambda(\sigma\tau x)^\nu$ como esquina superior izquierda y poseen el factor $t_\sigma^\kappa s_\tau^\lambda$.

Observación final. El teorema I tiene una excepción en el caso en que los números $n$ y $\mu$ (respectivamente $m$ y $\nu$) se vuelven nulos, de modo que los plegamientos I, II y IV ya no existen, mientras que el plegamiento III (respectivamente IV) todavía existe. En este caso designamos como sucesión de formas el miembro diagonal del esquema:
\[
\begin{array}{ccccc}
s_x^m u_\tau^\nu &&&& t_x^n u_\sigma^\mu\\
& s_\tau && t_\sigma &\\
& s_\tau^2 && t_\sigma^2 &\\
&&\ddots&&\\
&&s_\tau^\nu && t_\sigma^\mu .
\end{array}
\]
De acuerdo con la ausencia de los plegamientos I y II, el desarrollo en serie según las polares de la sucesión de formas se reduce entonces siempre al miembro inicial; sin embargo, los teoremas sobre reductores siguen siendo válidos.

\begin{center}
\textbf{§ 2. \quad Desarrollos en serie según polares de la sucesión de formas.}
\end{center}

Para formas con $n$ variables se han establecido explícitamente dos clases de desarrollos en serie\footnote{Cf. Gordan, \emph{Über Kombinanten}, §§ 2 y 5 (\emph{Math. Ann.} vol. V).}:

1) para formas con una serie cada una de variables contragredientes, $s_x^m u_\tau^\nu$, una serie que progresa según potencias de $u_x$, cuyos coeficientes son ``formas normales'';

2) para formas con dos series de variables cogredientes, $s_x^m t_x^n$, un desarrollo según las polares de los covariantes elementales (polares de la sucesión de formas $s_x^m t_x^n$), que en el caso ternario tiene la forma siguiente\footnote{Cf. Study, \emph{Methoden zur Theorie der ternären Formen} (Teubner 1889) II. § 7. A diferencia de la derivación dada allí, la nuestra proporciona el método para determinar rápidamente por cálculo los coeficientes numéricos $C_{pr\varrho}$. El desarrollo de Clebsch--Study, bajo la especialización a), a'), sería
\[
(stu)^\mu u_\tau^{\nu-\kappa}v_\tau^\kappa s_x^m t_x^{\,n-\lambda}(tuv)^\lambda
=\sum_{p,r,\varrho} C_{pr\varrho}
\bigl[s_\tau^\varrho(stu)^{\mu+r-\varrho}\bigr]_{\widehat{uv}^{\lambda-r+\varrho}v^{\nu-\varrho},\,u^p,\,(vw=\widehat{xw})},
\]
y por tanto no permite la simplificación que en nuestro desarrollo aparece ya durante el cálculo, cuando se sabe que han de omitirse todos los términos que contienen el factor $u_x$.}:
\[
\tag{I}
s_x^m t_x^{n-\lambda}t_y^\lambda
=\sum_{\varrho}
\frac{\binom{m}{\varrho}\binom{\lambda}{\varrho}}
{\binom{m+n-\varrho+1}{\varrho}}
\bigl[(st\widehat{xy})^\varrho s_x^{m-\varrho}t_x^{n-\varrho}\bigr]_{y^{\lambda-\varrho}}.
\]

Combinamos los dos desarrollos estableciendo, para formas con dos series cada una de variables contragredientes,
\[
s_x^m t_x^{n-\lambda}t_y^\lambda u_\sigma^\mu u_\tau^{\nu-\kappa}v_\tau^\kappa
\]
desarrollos según las potencias de $u_x$, cuyos coeficientes son polares compuestas de la sucesión de formas $s_x^m t_x^n u_\sigma^\mu u_\tau^\nu$, tomadas según las variables $x$ y $u$.

Basta establecer la serie para dos series contragredientes de símbolos (respectivamente de variables), pues, reuniendo por pares las series de símbolos (de variables) en una nueva serie única, las formas con más series de símbolos (de variables) se reducen a este caso.

Para que todas las formas de la sucesión de formas contengan sólo dos series de símbolos, respectivamente de variables, especializamos:
\[
\begin{array}{ll}
\text{a) } \sigma=\widehat{st}, & \text{a') } y=\widehat{uv},\\[0.3em]
\text{b) } s=\widehat{\sigma\tau}, & \text{b') } v=\widehat{xy},
\end{array}
\]
y efectuamos el desarrollo para la especialización a), a').

Por transferencia directa del desarrollo I se obtienen polares compuestas para las formas $(stu)^\mu s_x^m t_x^{n-\lambda}t_y^\lambda$, mientras que para las formas $(stu)^\mu u_\tau^{\nu-\kappa}v_\tau^\kappa s_x^m t_x^{n-\lambda}t_y^\lambda$, en lugar de polares compuestas, aparecen términos cuya evaluación hace necesaria la introducción de variables auxiliares siempre nuevas, que sólo han de identificarse al final, lo cual complica innecesariamente el cálculo.

Por eso seguimos un camino indirecto: representamos las polares compuestas de todas las formas de la sucesión de formas, es decir, las expresiones
\[
\bigl[s_\tau^\varrho(stu)^{\mu-\varrho+r}\bigr]_{y^\lambda}^{v^\kappa}
\]
por los términos iniciales de esas polares, y obtenemos después el desarrollo buscado según polares mediante inversión del sistema recurrente de ecuaciones que resulta.

Según el principio de transferencia, para la $\lambda$-ésima polar de una forma $s_x^m t_x^n:[s_x^m t_x^n]_{y^\lambda}$ ($\lambda\le n$) vale el desarrollo siguiente (el caso $\lambda>n$ se reduce al primero por el desarrollo de $[s_y^m t_y^n]_{x^{m+n-\lambda}}$)\footnote{Gordan, \emph{Die Resultante binärer Formen} (cap. I, § 5). Rend. Circ. Mat. di Palermo XXII. 1906.}:
\[
\bigl[s_x^m t_x^n\bigr]_{y^\lambda}
=\sum_r (-1)^r\,
\frac{\binom{m}{r}\binom{\lambda}{r}}{\binom{m+n}{r}}\,
(st\widehat{xy})^r s_x^{m-r}t_x^{n-\lambda}t_y^{\lambda-r},
\]
y del mismo modo para la $\kappa$-ésima polar de $u_\sigma^\mu u_\tau^\nu:[u_\sigma^\mu u_\tau^\nu]_{v^\kappa}$,
\[
\bigl[u_\sigma^\mu u_\tau^\nu\bigr]_{v^\kappa}
=\sum_\varrho (-1)^\varrho\,
\frac{\binom{\mu}{\varrho}\binom{\kappa}{\varrho}}{\binom{\mu+\nu}{\varrho}}\,
(\sigma\tau\widehat{uv})^\varrho u_\sigma^{\mu-\varrho}u_\tau^{\nu-\kappa}v_\tau^{\kappa-\varrho}.
\]

Por multiplicación se sigue de ello, introduciendo las especializaciones a) y a'),
\[
\begin{aligned}
\bigl[s_x^m t_x^n(stu)^\mu u_\tau^\nu\bigr]_{\widehat{uv}^{\lambda}}^{v^\kappa}
={}&
\sum_{r,\varrho}c_{r\varrho}\,
(st\widehat{uv})^r(st\widehat{uv})^\varrho
s_x^{m-r}t_x^{n-\lambda}(tuv)^{\lambda-r}(stu)^{\mu-\varrho}
u_\tau^{\nu-\kappa}v_\tau^{\kappa-\varrho},
\end{aligned}
\]
y de ahí, desarrollando según potencias de $u_x$, distinguiendo el primer término ($\sum'_{r,\varrho}$ significa que se omite el término $r=0$, $\varrho=0$) y teniendo en cuenta (a) p. 26,
\[
\tag{II}
\begin{aligned}
&s_x^m t_x^{n-\lambda}(tuv)^\lambda(stu)^\mu u_\tau^{\nu-\kappa}v_\tau^\kappa\\
&\quad =
\bigl[s_x^m t_x^n(stu)^\mu u_\tau^\nu\bigr]_{\widehat{uv}^{\lambda}}^{v^\kappa}\\
&\qquad
-\sum_{r,\varrho}' c_{r\varrho}\,s_\tau^\varrho(stu)^{\mu-\varrho+r}
u_\tau^{\nu-\kappa}v_\tau^{\kappa-\varrho}
s_x^{m-r}t_x^{n-\lambda}(tuv)^{\lambda-r+\varrho}v_x^r\\
&\qquad
+u_x\sum_{\varrho}\sum_{r=1}^{\lambda} c_{r\varrho}\,
s_\tau^\varrho(stu)^{\mu-\varrho+r-1}(stv)
u_\tau^{\nu-\kappa}v_\tau^{\kappa-\varrho}
s_x^{m-r}t_x^{n-\lambda}(tuv)^{\lambda-r+\varrho}v_x^{r-1}\\
&\qquad
\vdots\\
&\qquad
-(-1)^\lambda u_x^\lambda\sum_\varrho c_{\lambda\varrho}\,
s_\tau^\varrho(stu)^{\mu-\varrho}(stv)^\lambda
u_\tau^{\nu-\kappa}v_\tau^{\kappa-\varrho}
s_x^{m-\lambda}t_x^{n-\lambda}(tuv)^\varrho,
\end{aligned}
\]
donde
\[
c_{r\varrho}=(-1)^{r+\varrho}\cdot
\frac{\binom{m}{r}\binom{\lambda}{r}}{\binom{m+n}{r}}\cdot
\frac{\binom{\mu}{\varrho}\binom{\kappa}{\varrho}}{\binom{\mu+\nu}{\varrho}}
\qquad
\begin{array}{l}
\text{para }\lambda\le n,\\
\kappa\le\nu.
\end{array}
\]
y de manera análoga para las demás combinaciones de los números
\[\lambda,n;\kappa,\nu.\]
La expresión
\[s_x^m t_x^{n-\lambda}(tuv)^\lambda(stu)^\mu u_\tau^{\nu-\kappa}v_\tau^\kappa\]
queda, por tanto, reducida a una polar compuesta y, aplicando la identidad
\[
(stv)u_\tau=(stu)v_\tau-s_\tau(tuv),
\]
a una suma de expresiones formadas de manera análoga, que corresponden a formas superiores de la sucesión de formas.

Por iteración se llega al desarrollo:
\[
\tag{III}
s_x^m t_x^{n-\lambda}(tuv)^\lambda(stu)^\mu u_\tau^{\nu-\kappa}v_\tau^\kappa
=
\sum_{p,r,\varrho}u_x^p\,C_{pr\varrho}\cdot
\bigl[s_\tau^\varrho(stu)^{\mu-\varrho+r}\bigr]_{\widehat{uv}^{\lambda-r+\varrho}v^{\kappa-\varrho+p},\,v_x^{r-p}}.
\]
Los coeficientes numéricos $C_{pr\varrho}$ se calculan de manera unívoca y recursiva a partir de los coeficientes $c_{r\varrho}$, $c'_{r\varrho}$ del sistema de ecuaciones que resulta de iterar (II.) (véase el ejemplo, p. 42). Desarrollos análogos se obtienen para las demás especializaciones.

\begin{center}
\textbf{§ 3. \quad Teoremas de reducción.}
\end{center}

Los teoremas conocidos sobre la reducción de formas y sistemas de formas han de ponerse ahora en relación con los §§ 1 y 2; para ello son necesarias algunas definiciones tomadas de la teoría de las formas binarias.

a) Designamos como \emph{sistema relativamente completo}, módulo una sucesión prefijada de formas, a un sistema de formas con la propiedad de que todas las formaciones que surgen por plegamiento de un producto arbitrario de esas formas pueden expresarse como funciones racionales enteras de las formas del sistema, salvo un término aditivo formado por expresiones que han surgido por plegamiento de las formas del sistema con el sistema del módulo\footnote{Gordan--Kerschensteiner, \emph{Vorlesungen über Invariantentheorie}. vol. II, p. 227.}.

b) Llamamos \emph{reducible} a una forma cuando puede expresarse por formas que tienen invariantes como factores, o por ``formas superiores''; esto es, por formas superiores de la sucesión total de formas a la que pertenece la forma, o por formas que contienen los símbolos del módulo en orden superior.

Los teoremas sobre reductores\footnote{Gordan, \emph{Math. Annalen} vol. XVII, p. 222.} pueden entonces enunciarse en su forma más general de la siguiente manera:

Definición: \emph{Un reductor es una sucesión de formas reducible.}

\emph{Teorema II. Si la forma inicial de una sucesión de formas es reducible porque uno de sus miembros ha surgido por plegamiento con un reductor (tiene un reductor como factor), y si la forma final de la sucesión de formas ha surgido de ese mismo miembro por plegamiento, entonces toda la sucesión de formas es reducible}\footnote{La simplificación que introduce el teorema II puede verse en el siguiente ejemplo: para la forma especial $f=x_1^3x_2+x_2^3x_3+x_3^3x_1$, $a_y^2a_x^2u_y^2$ es un reductor. De ello se sigue, por el teorema II, la reducción de la sucesión de formas
\[
\theta_\nu(\vartheta\nu x)\theta_x^3u_\vartheta u_\nu^2
=
a_\nu^2(abu)-a_\nu b_\nu(aba),
\]
pues $\theta_\nu^4=a_\nu^2b_\nu^2(abu)^2$ ha surgido del miembro $a_\nu^2(abu)$ por plegamiento. En Gordan, loc. cit. p. 231, en cambio, se efectúa separadamente la reducción de las formas
\[
\theta_\nu(\vartheta\nu x),\quad \theta_\nu(\vartheta\nu x)^2,\quad
\theta_\nu^2,\quad \theta_\nu^2(\vartheta\nu x),\quad
\theta_\nu^2(\vartheta\nu x)^2
\]
.}.

Demostración: la sucesión de formas surge del miembro inicial por plegamiento en sí mismo, por tanto, por un lado, mediante plegamiento superior con el reductor y, por otro, mediante plegamiento del reductor en sí mismo. En ambos casos, según § 2, todas las formas de la sucesión de formas pueden representarse como polares (transvecciones) sobre la sucesión de formas del reductor.

La conclusión desde la forma inicial a toda la sucesión de formas ya no se da en los demás métodos de reducción. Éstos son:

1) la llamada ``reducción doble''.

Entendemos por ella la reducción, de dos maneras, de una expresión que tiene dos reductores independientes entre sí como factores, con lo cual surgen relaciones entre las formas superiores a las que se redujo.

Mediante la aplicación sistemática de la ``reducción doble'' se deben obtener, por una parte, todas las relaciones\footnote{Así se encontraron, por ejemplo, mediante ``reducción doble'': la reducción de la forma $a_\eta^4$ (§ 10), la única forma que Maisano incluyó superfluamente en el sistema; luego, las relaciones mencionadas en la nota a la introducción entre los tres covariantes y los tres invariantes.}; por otra parte, cuando la ``reducción doble'', aplicada a distintas expresiones, no da nuevas relaciones, se tiene a la vez un criterio para la independencia de las formas superiores y un control del cálculo.

2) \emph{Plegamiento con formas descomponibles.}

Por ``formas descomponibles'' entenderemos ---con una ligera generalización del concepto usual--- formas que pueden expresarse por productos de formas de grado inferior y por formas ``superiores''. Los productos de formas pasan, bajo un plegamiento simple, a productos; asimismo, los productos de formas no lineales bajo un plegamiento doble en las especializaciones a') y b') (§ 2), según el teorema del producto:
\[
a_yb_y a_xb_x=\frac12a_y^2b_x^2+\frac12b_y^2a_x^2-\frac12(ab\widehat{yx})^2.
\]
Las primeras polares (transvecciones) y ciertas segundas polares de formas descomponibles son, por tanto, de nuevo formas descomponibles. Se sigue:

Un miembro de una transvección simple (respectivamente doble) sobre una forma descomponible, surgido por plegamiento simple (respectivamente doble) con una forma descomponible, se vuelve él mismo una forma descomponible:

a) si las formas inmediatamente superiores en la sucesión de formas de la forma descomponible se descomponen o se vuelven reducibles, o también si, en un plegamiento simple, se presentan las especializaciones $y=\widehat{uv}$ o $v=\widehat{xy}$;

b) si los demás plegamientos simples conducen a formas superiores (véase § 12, B. $H_k$ y $H_k(k\eta x)$).

Un tal miembro de una transvección puede también descomponerse:

c) si los demás plegamientos simples conducen a formas inferiores que se descomponen independientemente del plegamiento con la forma descomponible, es decir, que pueden expresarse por productos y por la forma que se ha de reducir, pero verificando en cada caso por cálculo que el coeficiente numérico del miembro correspondiente no se anula idénticamente. Esto no es otra cosa que la ``reducción doble'' de la forma inferior (véase § 23, C. $H_q(\theta H u)(\theta s u)$).

Formas descomponibles surgen por todos los plegamientos simples con determinantes funcionales\footnote{cf. Gordan, loc. cit. § 3.}, y además por ciertos plegamientos superiores. Se tiene
\[
(st\widehat{yx})s_y=\frac12\,t\,s_y^2-\frac12\,s\,t_y^2+\frac12(st\widehat{yx})^2,
\]
\[
(st\widehat{yx})^2s_y^2=\frac13\,t\,s_y^3-\frac13\,s\,t_y^3+s_y(st\widehat{yx})^2-\frac13(st\widehat{yx})^3.
\]
Fórmulas análogas valen para las formas dualísticas
\[
(\sigma\tau\widehat{\nu u})v_\sigma
\quad\text{y}\quad
(\sigma\tau\widehat{\nu u})v_\sigma^2.
\]

De los teoremas de este parágrafo resulta la siguiente regla para establecer todas las formas irreducibles que surgen del plegamiento de $S$ con $T$:

Se procede, comenzando por los plegamientos más bajos, por cualquier camino fijado de antemano (por ejemplo, fila por fila o simétricamente respecto del término diagonal), en la formación de la sucesión de formas $ST$, hasta llegar a una forma que tiene un reductor como factor. La sucesión de formas definida por esta forma debe omitirse tan pronto como la forma final satisfaga la condición establecida en el teorema II. Después de excluir todas las sucesiones de formas reducibles, se han de eliminar, mediante la reducción doble, todas las demás formas reducibles, así como todas las formas descomponibles. Los lugares de la sucesión total de formas que quedan doblemente cubiertos por sucesiones de formas reducibles independientes entre sí dan lugar a la reducción de formas superiores (véase § 11, D).

\clearpage
\noindent\textbf{Teorema III.}
\emph{Las formas tomadas del dominio binario, es decir, aquellas formas que surgen de las formas del sistema de la forma binaria correspondiente, según el principio de transferencia de Clebsch, por bordado, forman un sistema relativamente completo módulo $(abc)$.}

\noindent\textbf{Demostración:}
A una relación binaria que afirma que las formas $Q'_1,\ldots,Q'_n$ forman un sistema completo de una forma fundamental binaria,
\[
Q'-F(Q')=0=\sum_\lambda\bigl\{(ab)c_x+(bc)a_x+(ca)b_x\bigr\}\varphi_\lambda^{*}
\]
corresponde por bordado la relación ternaria:
\[
Q-F(Q_i)=\sum_\lambda u_x\cdot(abc)\varphi_\lambda .
\]
Pero ésta es precisamente la ecuación definitoria de un sistema relativamente completo módulo $(abc)$. Para la forma bicuadrática, el sistema de las formas tomadas consta de las siguientes formas:
\[
\begin{gathered}
f=a_x^4,\qquad
\theta=\theta_x^4u_\vartheta^2=(abu)^2a_x^2b_x^2,\qquad
K=K_x^6u_k^3=(a\theta u)u_\vartheta^2a_x^3\theta_x^3,\\
j=u_j^6=(a\theta u)^4u_\vartheta^2,\qquad
\nu=u_\nu^4=(abu)^4,
\end{gathered}
\]
entre las cuales las cuatro primeras forman un sistema relativamente completo módulo $((abc),\nu)$.

\begin{center}
\textbf{§ 5. \quad Reducción del módulo $(abc)$ a los módulos $\A$ y $\nu$.}
\end{center}

El módulo $(abc)$, dado sólo por un factor de paréntesis, debe, de acuerdo con la definición (§ 3, a), reducirse a formas del sistema. En otras palabras: la sucesión de formas $(abc)$ ha de normalizarse unívocamente (cf. § 1).

\[
a_s^2a_xu_x=\frac13 a_s^3u_x=\frac13 S\cdot u_x. \tag{2}
\]
\[
\left.
\begin{aligned}
(a\A u)^2&=a_s-\frac{S}{6}u_x^2,\\
(a\A u)^3&=0.
\end{aligned}
\right\} \tag{3}
\]
\[
\left.
\begin{aligned}
\theta(s)=(ss,x)^2&=2a_t+\frac13 S\cdot\theta-\frac13 Tu_x^2,\\
\theta(t)=(tt,x)^2&=-\frac13 S\cdot a_t+\frac23 T\cdot a_s+\frac1{18}S^2\cdot\theta-\frac1{18}u_x^2\cdot S\cdot T.
\end{aligned}
\right\} \tag{4}
\]

Sucesión de los módulos: $(abc)$; $\A$; $s$; $t$ (el sistema del módulo $t$ consta de la única forma $t$).\footnote{Gordan--Kerschensteiner, loc. cit. p. 134.}

Mostraremos que las formas
\[
\A=(abc)^2a_x^2b_x^2c_x^2,\qquad \nu=(abu)^4
\]
bastan como módulos, es decir, que las formas
\[
\begin{gathered}
\A=(abc)^2a_x^2b_x^2c_x^2,\qquad
a_\nu=(abc)(bcu)^3a_x^3,\\
a_\nu^2=(abc)^2(bcu)^2u_x^2,\qquad
i=a_\nu^4=(abc)^4
\end{gathered}
\]
definen la sucesión de formas $(abc)$. En la sucesión de formas $a_\nu$ falta la forma $a_\nu^3$, según la identidad
\[
a_\nu^3=(abc)^3a_x(bcu)=\frac13 u_x\cdot(abc)^4=\frac13 i\cdot u_x.
\]

Para reducir un módulo a un módulo superior debemos, por definición, reducir los plegamientos más generales, o bien todos los plegamientos especiales que entren en consideración con el módulo, a plegamientos con el módulo superior.

Los plegamientos más generales surgen, en nuestro caso, de las expresiones
\[
(abc)a_x^3b_y^3c_z^3,\qquad
(abc)a_x^2b_y^2c_z^2a_zb_xc_x
\]
(tomadas de las formas cúbicas),
\[
(abc)a_xb_yc_za_x^2b_x^2c_x^2
\]
(tomada de las formas cuadráticas), y por polarización de estas expresiones.

Para calcular estas expresiones por las polares de $\A$ y de la sucesión de formas $a_\nu$, respectivamente por sus términos iniciales, seguimos, como en § 2, el camino indirecto. Desarrollamos los términos polares
\[
(abc)^2a_x^2b_y^2c_z^2,\qquad
a_\nu(\nu xy)(\nu yz)(\nu zx)a_xa_ya_z
=(abc)(bc\widehat{x}u)(bc\widehat{y}z)(bc\widehat{z}x)a_xa_ya_z\quad\text{etc.}
\]
como agregados lineales de expresiones que tienen el factor simbólico $(abc)$, y por inversión del sistema de ecuaciones obtenemos las relaciones buscadas, así como una relación entre las polares tomadas según diferentes combinaciones de las variables. Para la evaluación sirven el teorema de identidad y el teorema del producto para determinantes.

El sistema de ecuaciones para $(abc)a_x^3b_y^3c_z^3$ es:
\begingroup\small
\[
\begin{aligned}
1)\quad
\sum_3 a_\nu(\nu yz)^3a_x^3
&=6(abc)a_x^3b_y^3c_z^3
-6\sum_3(abc)a_x^3b_y^2b_zc_z^2c_y^{*},\\[0.4em]
2)\quad
3(abc)^2a_x^2b_y^2c_z^2\cdot(xyz)
-\frac12\sum_3 a_\nu^2(\nu yz)^2\nu_x^2(xyz)
&=2\sum_3(abc)a_x^3b_y^2b_zc_z^2c_y\\
&\quad-4\sum_3(abc)a_xa_ya_zb_y^2b_zc_z^2c_x,\\[0.4em]
3)\quad
a_\nu(\nu xy)(\nu yz)(\nu zx)a_xa_ya_z
&=2\sum_3(abc)a_xa_ya_zb_y^2b_zc_z^2c_x,\\[0.4em]
4)\quad
\sum_3 a_\nu(\nu yz)^3u_x^3
-\frac32\sum_3 a_\nu^2(\nu yz)^2\nu_x^2(xyz)
+\frac16 i\cdot(xyz)^3
&=6\sum_3(abc)a_xa_ya_zb_y^2b_zc_z^2c_x.
\end{aligned}
\]
\endgroup

De ello resulta para
\[
(abc)a_x^3b_y^3c_z^3,
\]
y, por un cálculo análogo, para
\[
(abc)a_x^2b_y^2c_z^2a_zb_xc_x,\qquad
(abc)a_xb_yc_za_z^2b_x^2c_x^2:
\]
\[
\begin{aligned}
(abc)a_x^3b_y^3c_z^3
&=\frac32(abc)^2a_x^2b_y^2c_z^2(xyz)
+\frac32 a_\nu(\nu xy)(\nu yz)(\nu zx)a_xa_ya_z\\
&\quad-\frac1{12}i\cdot(xyz)^3,\\[0.4em]
(\text{IV.})\quad
(abc)a_x^2b_y^2c_z^2a_zb_xc_x
&=\frac23(abc)^2a_x^2b_y^2c_z^2c_x\cdot(xyz)\\
&\quad+\frac13a_\nu(\nu xy)(\nu yz)(\nu zx)a_x^3
-\frac16 a_\nu^2(\nu xy)(\nu zx)a_x^2\cdot(xyz),\\[0.4em]
(abc)a_xb_yc_za_z^2b_x^2c_x^2
&=\frac16(abc)^2a_x^2b_z^2c_z^2\cdot(xyz).
\end{aligned}
\]

Hemos obtenido así un sistema relativamente completo módulo $(\A,\nu)$, compuesto de las cuatro formas $f,\theta,K,j$. De ello resulta que todas las transvecciones de $f$ sobre $\theta$ que no se han tomado del dominio binario, así como la transvección $(a\theta u)^2$, reducible en el binario a $(ab)^4$, pueden expresarse mediante los símbolos $\A$ y $\nu$.\footnote{$\sum_3$ se refiere a la suma, de tres términos, obtenida a partir del término inicial por permutación cíclica de las variables. Las ecuaciones valen también separadamente para cada término de la suma.}

Obtenemos las fórmulas de reducción fundamentales para lo que sigue por especialización del desarrollo IV o, más brevemente, por cálculo directo (permutación de los símbolos), para $a_\vartheta(a\theta u)^2$, $a_\vartheta^2((a\theta u)^2$, $(a\theta u)^2$ según el siguiente planteamiento:
\[
\begin{aligned}
a_\vartheta(a\theta u)^2
&=(abc)(bcu)(abu)^2-\frac13(abc)(bcu)\{(bcu)a_x-u_x(abc)\}^2,\\
a_\vartheta^2(a\theta u)^2
&=(abc)^2(abu)^2-\frac13(abc)^2\{(bcu)a_x-u_x(abc)\}^2,
\end{aligned}
\]
para $((a\theta u)^2)^{*}$:
\[
\begin{aligned}
(abu)a_y^3b_x^3&=\frac32u_\vartheta(\vartheta yx)\theta_y^2+\frac14(\nu yx)^3,\\
((abu)(acu))b_x^3&=\frac32u_\vartheta(\vartheta\widehat{a}ux)(\theta au)^2+\frac14(\nu\widehat{a}ux)^3.
\end{aligned}
\]
\[
\left.
\begin{aligned}
(a\theta u)^2&=\frac16\nu\cdot f-\frac23u_x\cdot a_\nu
+\frac23u_x^2\cdot a_\nu^2-\frac{i}{18}u_x^4,\\
a_\vartheta&=\frac13u_x\cdot\A,\\
a_\vartheta(a\theta u)&=0,\\
a_\vartheta^2&=\A,\\
(a\theta u)^3&=0,\\
a_\vartheta(a\theta u)^2&=-\frac56a_\nu+\frac76u_x\cdot a_\nu^2-\frac{i}{9}u_x^3,\\
a_\vartheta^2(a\theta u)&=0,\\
(a\theta u)^4&=j,\\
a_\vartheta(a\theta u)^3&=0,\\
a_\vartheta^2(a\theta u)^2&=\frac23a_\nu^2-\frac{i}{9}u_x^2.
\end{aligned}
\right\} \tag{1}
\]

Añadimos la fórmula, procedente igualmente de la reducción del módulo $(abc)$,
\[
a_\nu^3a_xu_\nu=\frac13 i\cdot u_x. \tag{2}
\]

\emph{De las fórmulas (1.) y (2.) y de las fórmulas construidas de manera análoga para la forma fundamental $\nu$, todas las fórmulas de reducción posteriores se deducen por polarización}, con excepción de las relaciones para las formas descomponibles. De las fórmulas (1.) vemos:

La sucesión de formas:
\[
\begin{array}{rcl}
a_\vartheta(a\theta u)^2 &\text{conduce a}& \text{símbolos }\nu\text{ solamente},\\[0.2em]
a_\vartheta &\text{conduce a}& \text{símbolos }\A\text{ y }\nu,\\[0.2em]
(a\theta u)^2 &\text{conduce a}& \text{símbolos }j,\A\text{ y }\nu,\\[0.2em]
(a\theta u)^2\text{ bajo plegamiento I} &\text{conduce a}& \text{símbolos }j\text{ y }\nu,\\[0.2em]
(a\theta u)^2\text{ bajo plegamiento II} &\text{conduce a}& \text{símbolos }\A\text{ y }\nu.
\end{array}
\]

Por tanto podemos reemplazar:
\begin{enumerate}
\item módulo $(\A,\nu)$: los plegamientos con $j$ por los plegamientos correspondientes con $(a\theta u)^2$ o $(a\theta u)^3$;
\item módulo $(\nu)$: los plegamientos con $\A$ por los plegamientos correspondientes con $a_\vartheta$ o $a_\vartheta(a\theta u)$, o también por plegamientos especiales con $(a\theta u)^2$ (que sólo conducen a un plegamiento I simple).
\end{enumerate}
En particular se tiene:
\[
\begin{aligned}
(a\theta u)^2\theta_y^2\theta_x^2&=\frac13(jyx)^2\pmod{\nu},&
(a\theta u)^2u_y\theta_y&=-\frac16(jyx)^2\pmod{\nu},\\
(a\theta u)^2\nu_y^2&=\frac13(\A u\nu)^2\pmod{\nu},&
(a\theta u)(a\theta\nu)u_\vartheta\nu_\vartheta&=-\frac16(\A u\nu)^2\pmod{\nu},\\
(a\theta u)^2u_\nu^2\vartheta_\nu^2&=\frac13(\A u\nu)^2\A_\nu\pmod{\nu}.
\end{aligned}
\]

\begin{center}
\textbf{§ 6. \quad Reducción del módulo $(\A,\nu)$ al módulo $(\nu)$.}
\end{center}

Las fórmulas de reducción del parágrafo anterior nos dan el medio para establecer directamente el sistema relativamente completo módulo $\nu$; veremos que al sistema de las formas tomadas sólo tenemos que adjuntar las formas $\A$ y $(a\A u)=N$ (de manera análoga al caso de la forma cúbica, y válida más generalmente).

Para reducir el módulo $\A$, consideramos todos los plegamientos especiales que entran en cuenta, es decir, los plegamientos del sistema relativamente completo módulo $(\A,\nu)$ con el sistema de $\A$, y partimos de las formas de menor orden en los coeficientes, los plegamientos de $f$ con $\A$.

Para reducir las formas $(a\A u)^2$, $(a\A u)^3$, $(a\A u)^4$, reemplazamos, según § 5, los símbolos $\A$ por formas inferiores de la sucesión de formas, es decir, desarrollamos las expresiones
\[
a_\vartheta b_\vartheta(b\theta u)^2,\qquad
a_\vartheta(a\theta u)b_\vartheta(b\theta u)^2,\qquad
a_\vartheta(a\theta u)b_\vartheta(b\theta u)^3
\]
1) según las polares de la sucesión de formas $a_\vartheta$ (símbolos $\A$ y $\nu$),

2) según las polares de la sucesión de formas $b_\vartheta(b\theta u)^2$ (símbolos $\nu$),

y obtenemos las fórmulas de reducción (cálculo, véase abajo):
\begin{align}
\text{a)}\quad (a\A u)^2
&=\frac12(\vartheta\nu x)^2-\frac25 f\cdot a_\nu^2
-\frac45u_x\cdot\theta_\nu(\vartheta\nu x)^2 \notag\\
&\quad+u_x^2\left\{\frac16 i\cdot f-\frac1{40}S\right\},\notag\\
\text{b)}\quad (a\A u)^3
&=-\frac3{10}\theta_\nu(\vartheta\nu x),\tag{3}\\
\text{c)}\quad (a\A u)^4
&=\frac{21}{10}\theta_\nu^2-\frac3{10}\Pi
-\frac65u_x\cdot\theta_\nu^3+\frac1{10}u_x^2\cdot\theta .\notag
\end{align}

(Los mismos resultados se obtienen también por reducción doble de las expresiones $a_\vartheta^2(b\theta u)^2$, $a_\vartheta^2(b\theta u)^3$, $a_\vartheta^2((a\theta u)^2(b\theta u)^2$ y $a_\vartheta^2((a\theta u)(b\theta u))^3$, después de eliminar $a_\vartheta^2$.)

De las fórmulas (3.) resulta que $(a\A u)^2$ es un reductor relativamente al módulo $\nu$. De ello se deduce:

1) La reducción del sistema de $\A$: la sucesión de formas $(\A\A' u)^2=a_\vartheta^2((a\A u)^2$ tiene como factor el reductor $(a\A u)^2$.

2) La reducción del sistema nacido por plegamiento con el módulo $\A$:

La sucesión de formas $(\theta\A u)^2$ tiene como factor el reductor $(a\A u)^2$.

Para las formas $(\theta\A u)$ y $\A_\vartheta$, se obtiene:
\[
(a\A u)^2(abu)=(\theta\A u)-u_x\cdot\A_\vartheta(\theta\A u),
\]
\[
(a\A u)(a\A b)=-\frac12\A_\vartheta+\frac12u_x\cdot\A_\vartheta^2.
\]

Pero del reductor $(\theta\A u)$ resulta la reducción del sistema nacido por plegamiento con el módulo $\A$.

\bigskip

El sistema relativamente completo módulo $\nu$ se compone, por tanto, de las seis formas:
\[
f,\ \theta,\ K,\ j,\ \A,\ N.
\]

Como ejemplo del desarrollo en serie del § 2, damos en detalle la deducción de la fórmula (3.)a; en los cálculos posteriores se escribirá directamente la fórmula final III. (Las polares de las formas nulas ya se han omitido durante el cálculo; la primera línea da los valores insertados según el desarrollo II, la segunda línea los valores finales encontrados recursivamente empezando por la última ecuación del sistema.) Del desarrollo II se sigue:

\begin{align*}
\text{1)}\quad&m=3,\quad n=4,\quad \lambda=2,\quad \mu=0,\quad \nu=\chi=1,\\
a_\vartheta b_\vartheta(b\theta u)^2
&=[a_\vartheta]_{b^2}b+\frac67\{a_\vartheta b_\vartheta(a\theta u)(b\theta u)-u_xa_\vartheta b_\vartheta(a\theta b)(b\theta u)\}\\
&\quad-\frac17\{a_\vartheta(a\theta u)^2b_\vartheta-2u_xa_\vartheta(a\theta u)(a\theta b)b_\vartheta+u_x^2a_\vartheta(a\theta b)^2b_\vartheta\}\\
&=\frac23(a\Delta u)^2+\frac15[a_\vartheta(a\theta u)^2]b+\frac1{30}f[a_\vartheta^2(a\theta u)^2]-\frac25u_x[a_\vartheta(a\theta u)^2]b^2\\
&\quad-\frac1{15}u_x[a_\vartheta^2(a\theta u)^2]b+\frac15u_x^2[a_\vartheta(a\theta u)^2]b^3+\frac1{30}u_x^2[a_\vartheta^2(a\theta u)^2]b^2;\\[.6em]
\text{2)}\quad&m=2,\quad n=3,\quad \lambda=1,\quad \mu=1,\quad \nu=\chi=1,\\
a_\vartheta(a\theta u)b_\vartheta(b\theta u)
&=\frac25\{a_\vartheta(a\theta u)^2b_\vartheta-u_xa_\vartheta(a\theta u)(a\theta b)b_\vartheta\}+\frac12a_\vartheta^2(b\theta u)^2\\
&\quad-\frac15\{a_\vartheta^2(b\theta u)(a\theta u)-u_xa_\vartheta^2(a\theta b)(b\theta u)\}\\
&=\frac12(a\Delta u)^2+\frac25[a_\vartheta(a\theta u)^2]b+\frac1{15}[a_\vartheta^2(a\theta u)^2]f\\
&\quad-\frac25u_x[a_\vartheta(a\theta u)^2]b^2-\frac1{10}u_x[a_\vartheta^2(a\theta u)^2]b+\frac1{30}u_x^2[a_\vartheta^2(a\theta u)^2]b^2;\\[.6em]
\text{3)}\quad&m=2,\quad n=3,\quad \lambda=1,\quad \mu=1,\quad \nu=1,\quad \chi=2,\\
a_\vartheta(a\theta b)b_\vartheta(b\theta u)
&=\frac25\{a_\vartheta(a\theta b)(a\theta u)b_\vartheta-u_xa_\vartheta(a\theta b)^2b_\vartheta\}\\
&=\frac25[a_\vartheta(a\theta u)^2]b^2+\frac1{30}[a_\vartheta^2(a\theta u)^2]b-\frac25u_x[a_\vartheta(a\theta u)^2]b^3-\frac1{30}u_x[a_\vartheta^2(a\theta u)^2]b^2;\\[.6em]
\text{4)}\quad&m=1,\quad n=2,\quad \lambda=0,\quad \mu=2,\quad \nu=\chi=1,\\
a_\vartheta(a\theta u)^2b_\vartheta
&=[a_\vartheta(a\theta u)^2]b+\frac23a_\vartheta^2(a\theta u)(b\theta u)\\
&=[a_\vartheta(a\theta u)^2]b+\frac16[a_\vartheta^2(a\theta u)^2]f-\frac16u_x[a_\vartheta^2(a\theta u)^2]b;\\[.6em]
\text{5)}\quad&m=1,\quad n=2,\quad \lambda=0,\quad \mu=2,\quad \nu=1,\quad \chi=2,\\
a_\vartheta(a\theta u)(a\theta b)b_\vartheta
&=[a_\vartheta(a\theta u)^2]b^2+\frac13a_\vartheta^2(a\theta b)(b\theta u)\\
&=[a_\vartheta(a\theta u)^2]b^2+\frac1{12}[a_\vartheta^2(a\theta u)^2]b-\frac1{12}u_x[a_\vartheta^2(a\theta u)^2]b^2;\\[.6em]
\text{6)}\quad&m=1,\quad n=2,\quad \lambda=0,\quad \mu=2,\quad \nu=1,\quad \chi=3,\\
a_\vartheta(a\theta b)^2b_\vartheta&=[a_\vartheta(a\theta u)^2]b^3;\\[.6em]
\text{7)}\quad&m=2,\quad n=4,\quad \lambda=2,\quad \mu=\nu=\chi=0,\quad (\text{según el desarrollo I}),\\
a_\vartheta^2(b\theta u)^2&=[a_\vartheta^2]_{ub^2}+\frac1{10}\{[a_\vartheta^2(a\theta u)^2]f-2u_x[a_\vartheta^2(a\theta u)^2]b+u_x^2[a_\vartheta^2(a\theta u)^2]b^2\};\\[.6em]
\text{8)}\quad&m=1,\quad n=3,\quad \lambda=1,\quad \mu=1,\quad \nu=\chi=0,\quad (\text{desarrollo I}),\\
a_\vartheta^2(a\theta u)(b\theta u)&=\frac14[a_\vartheta^2(a\theta u)^2]f-\frac14u_x[a_\vartheta^2(a\theta u)^2]b;\\[.6em]
\text{9)}\quad&m=1,\quad n=3,\quad \lambda=1,\quad \mu=\chi=1,\quad \nu=0,\quad (\text{desarrollo I}),\\
a_\vartheta^2(a\theta b)(b\theta u)&=\frac14[a_\vartheta^2(a\theta u)^2]b-\frac14u_x[a_\vartheta^2(a\theta u)^2]b^2.
\end{align*}

De 1) y 4) resulta:
\[
\begin{aligned}
(a\Delta u)^2
&=\frac65[a_\vartheta(a\theta u)^2]b+\frac15f[a_\vartheta^2(a\theta u)^2]+\frac35u_x[a_\vartheta(a\theta u)^2]b^2-\frac3{20}u_x[a_\vartheta^2(a\theta u)^2]b\\
&\quad-\frac3{10}u_x^2[a_\vartheta(a\theta u)^2]b^3-\frac1{20}u_x^2[a_\vartheta^2(a\theta u)^2]b^2,\\
(a\Delta u)^2
&=-a_\nu b_\nu+\frac35fa_\nu^2+\frac45u_xa_\nu^2b_\nu-\frac3{20}u_x^2a_\nu^2b_\nu^2-\frac1{12}u_x^2if.
\end{aligned}
\]

(Para la conversión en símbolos $\theta$, véanse § 8 y § 9.) La fórmula (3.)a podría calcularse aún más brevemente por transferencia directa del desarrollo I, introduciendo la variable auxiliar $w=x\widehat{u}b$, mientras que, ya desde las fórmulas (3.)b y (3.)c, el camino que hemos seguido se vuelve más simple; para (3.)b se sabe de antemano, por § 9, que todos los términos que contienen el factor $u_x$ deben omitirse. Para calcular (3.)b y (3.)c se obtiene:
\[
\begin{aligned}
(3.)\mathrm{b)}\quad 1)&\quad a_\vartheta(a\theta u)b_\vartheta(b\theta u)^2=\frac12(a\Delta u)^3+\frac45[a_\vartheta(a\theta u)^2]_{ub}b+\frac7{360}[a_\vartheta^2(a\theta u)^2]_{ub^2},\\
2)&\quad a_\vartheta(a\theta u)b_\vartheta(b\theta u)^2=[a_\vartheta(a\theta u)^2]_{ub}b+\frac{13}{36}[a_\vartheta^2(a\theta u)^2]_{ub^2};\\[.5em]
(3.)\mathrm{c)}\quad 1)&\quad a_\vartheta(a\theta u)b_\vartheta(b\theta u)^3=\frac12(a\Delta u)^4+\frac65[a_\vartheta(a\theta u)^2]_{ub^2}b+\frac{53}{120}[a_\vartheta^2(a\theta u)^2]_{ub^3}\\
&\qquad-\frac65u_x[a_\vartheta(a\theta u)^2]_{ub^3}b^2-\frac35u_x[a_\vartheta^2(a\theta u)^2]_{ub^2}b+\frac{19}{120}u_x^2[a_\vartheta^2(a\theta u)^2]_{ub^3}b^2,\\
2)&\quad a_\vartheta(a\theta u)b_\vartheta(b\theta u)^3=\frac34[a_\vartheta^2(a\theta u)^2]_{ub^2}.
\end{aligned}
\]

Observación final: las transvecciones de $f$ sobre $j$ que son reducibles según § 5 se calculan de manera análoga. Se tiene
\[
\tag{4}
\begin{aligned}
g_\nu&=-\theta_\nu+u_x\left(\frac92\theta_\nu^2-\frac12H\right)-3u_x^2\theta_\nu^3+\frac12u_x^3\theta,\\
g_\nu^2&=\frac{21}{10}\theta_\nu^2-\frac3{10}H-2u_x\theta_\nu^3+\frac12u_x^2\theta,\\
g_\nu^3&=-\frac7{10}\theta_\nu^3+\frac12u_x^2\theta,\qquad g_\nu^4=\frac35\theta.
\end{aligned}
\]

\[
\tag{4}
g_\nu^3=-\frac7{10}\theta_\nu^3+\frac12u_x^2\theta,
\qquad
g_\nu^4=\frac35\theta.
\]

De (3.) y (4.) resulta, por transferencia desde el dominio binario,
\[
(\theta\theta'u)^2=\frac13\,j\cdot f\pmod{\nu}.\footnotemark
\]
Las demás formas de la sucesión de formas $(\theta\theta'u)^2$ y la forma $\theta'_\nu$ son reducibles a los símbolos $\nu$. Por tanto podemos reemplazar:
\[
\bmod\ \nu:\quad \text{los plegamientos con } f\cdot j
\text{ por los plegamientos correspondientes con }(\theta\theta'u)^2.
\]
Además se tiene:
\[
(\vartheta\vartheta'x)^2=\frac43 f\cdot\A,
\qquad
\theta_{g\nu}(\vartheta\vartheta'x)=-N.
\]
\footnotetext{Gordan--Kerschensteiner, loc. cit. p. 181.}



\begin{center}
{\Large\bfseries Capítulo III.\quad El sistema relativamente completo módulo $(s,\varrho,t)$.}
\end{center}

\begin{center}
\textbf{\S\ 7.\quad Visión general de la formación del sistema.}
\end{center}

Para pasar del sistema relativamente completo módulo $\nu$ al sistema relativamente completo con el módulo inmediatamente superior, según la definición de \S\ 3 hay que formar el sistema de $\nu$ con respecto a este módulo superior y plegarlo con las formas del sistema relativamente completo módulo $\nu$. Pero $\nu=u_\nu^4$, como contravariante de cuarto grado de la forma $f$, se halla en relación dualística con ella. Si, por tanto, consideramos $\nu$ como forma fundamental, obtenemos un sistema relativamente completo de seis formas módulo $\nu(\nu)=(\nu\nu_1x)^4=s_x^4$.

Así, para formar el sistema relativamente completo módulo $s$ (sistema III), debemos transveccionar
\[
\begin{array}{ll}
\text{Sistema I:} & f,\ \theta,\ K,\ j,\ \A,\ N\quad \text{con}\vphantom{\dfrac11}\\[0.3em]
\text{Sistema II:} & \nu,\ H=(\nu\nu_1x)^2u_\nu^2u_{\nu_1}^2,\ L=(\nu\eta x)H_x^2u_\nu^3u_\eta^3,\\
& g=(\nu\eta x)^4H_x^2,\ \sigma=H_\eta^2u_\nu^2u_\eta^4,\ (\nu\sigma x).
\end{array}
\]
Se verá (fórmula (13)) que la forma $g$ es reducible, de modo que el sistema II consta sólo de cinco formas.

En el curso del cálculo se adjuntan como módulos las formas cuadráticas
\[
\varrho=u_\varrho^2=\theta_\nu^4u_\vartheta^2,
\qquad
 t=t_x^2=a_\eta^4H_x^2,
\]
de manera que se obtiene un sistema relativamente completo módulo $(s,\varrho,t)$; y el módulo $(\varrho,t)$ debe considerarse como módulo superior al módulo $s$.

La ordenación de las formas individuales debe seguir el orden en los coeficientes. Normalizamos el plegamiento
\[
\theta_\nu(K_\nu,\theta_\nu,\ldots)
\]
como superior al plegamiento
\[
H_y(H_y,L_y,\ldots).
\]
Por esto, según \S\ 1 y \S\ 3b, queda determinada de modo unívoco la sucesión de las formas individuales del mismo orden.

La formación de las formas de igual orden se llevará a cabo en forma tabular:

I. Indicación de qué formas de los sistemas I y II han de plegarse entre sí para alcanzar el orden determinado en los coeficientes.

II. Disposición de las formas irreducibles según el esquema de la sucesión de formas.

III. Reducción de las sucesiones de formas o formas reducibles o descomponibles, es decir, de los lugares en los que se interrumpe la sucesión total de formas; con ello se demuestra que todas las formaciones irreducibles quedan agotadas por las indicadas.

IV. Consecuencias: 1) indicación de los reductores recién surgidos; 2) indicación de aquellas formas que no entran, en productos con formas del mismo sistema, en plegamiento con formas del otro sistema.

Se verá que sólo aparecen:

1) plegamientos de una forma del sistema I con una forma del sistema II;

2) plegamientos de potencias de $\theta$, respectivamente de $H$, o de dos formas de uno de los sistemas con una forma del segundo sistema.

Obtendremos el siguiente esquema de plegamiento:
\[
\begin{array}{r@{\ }c@{\ }l@{\qquad}c@{\qquad}r@{\ }c@{\ }l}
\multicolumn{3}{c}{\text{Sistema I}} & \text{plegado con} & \multicolumn{3}{c}{\text{Sistema II}}\\[0.25em]
1. & \text{orden} & f & & 2. & \text{orden} & \nu\\
2. & ,, & \theta & & 4. & ,, & H\\
3. & ,, & K,j,\A & & 6. & ,, & L,\sigma\\
4. & ,, & N,\theta^2,f\cdot j & & 8. & ,, & (\nu\sigma x),H^2\\
5. & ,, & \theta\cdot K & & 10. & ,, & H\cdot L\\
6. & ,, & \theta^3,\A\cdot j & & 12. & ,, & H^3.
\end{array}
\]

Además de la ``doble reducción'', sirven como control del cálculo las dos formas especiales:

\[
\begin{array}{ll}
\text{I.}&\begin{array}{rlrlrl}
f&=\displaystyle\sum_3x_1^3x_2,& u_\nu^2&=\frac{i}{6}u_x^2,& s&=\frac{i}{3}f,\\
a_\eta^2&=-\frac19 i\theta,& H_\nu^2&=\frac{i}{6}\theta,& \sigma&=-\frac{i}{6}j,\\
g&=-\frac23i\Delta,& \varrho&=0,& t&=0,
\end{array}\\[1.2em]
\text{II.}&\begin{array}{rlrlrl}
f&=\displaystyle\sum_3x_1^4,& s&=\frac43if,& a_\eta^2&=\frac29i\theta,\\
\Delta_\nu^2&=\frac1{15}i\theta,& \sigma&=\frac43ij,& g&=\frac43i\Delta,\\
\varrho&=0,& t&=0.
\end{array}
\end{array}
\]

Observación final: a las fórmulas (1), (2), (3), (4) les corresponden, para la forma fundamental $\nu$, las siguientes:
\[
\tag{5}
\begin{array}{c|c|c}
(\nu\eta x)^2=A & H_\nu(\nu\eta x)=0 & H_\nu^2=\sigma\\[0.4em]
(\nu\eta x)^3=0 & H_\nu(\nu\eta x)^2=B & H_\nu^2(\nu\eta x)=0\\[0.4em]
(\nu\eta x)^4=g & H_\nu(\nu\eta x)^3=0 & H_\nu^2(\nu\eta x)^2=C,
\end{array}
\]
\[
A=\frac16s\nu-\frac23u_xs_\nu+\frac23u_x^2s_\nu^2-\frac{J}{18}u_x^4,
\quad B=-\frac56s_\nu+\frac76u_xs_\nu^2-\frac{J}{9}u_x^3,
\quad C=\frac23s_\nu^2-\frac{J}{9}u_x^2.
\]
\[
\tag{6}
s_\nu^3u_xs_\nu=\frac13u_xs_\nu^4=\frac13Ju_x.
\]

\[
\tag{7}
\begin{aligned}
\text{a)}\quad (\nu\sigma x)^2
&=\frac12(Hsu)^2-\frac25\nu\cdot s_\nu^2-
\frac45u_x\,s_\eta(Hsu)^2
+u_x^2\left\{\frac16J\cdot\nu-\frac1{40}(ss'u)^4\right\},\\
\text{b)}\quad (\nu\sigma x)^3&=-\frac3{10}s_\eta(Hsu),\\
\text{c)}\quad (\nu\sigma x)^4&=\frac{21}{10}s_\eta^2-\frac3{10}Z-
\frac65u_xs_\eta^3+\frac1{10}u_x^2s_\eta^4.
\end{aligned}
\]
\[
\tag{8}
\begin{aligned}
\text{a)}\quad g_\nu&=-s_\eta+u_x\left(\frac92s_\eta^2-\frac12Z\right)-3u_x^2s_\eta^3+\frac12u_x^3s_\eta^4,\\
\text{b)}\quad g_\nu^2&=\frac{21}{10}s_\eta^2-\frac3{10}Z-2u_xs_\eta^3+\frac12u_x^2s_\eta^4,\\
\text{c)}\quad g_\nu^3&=-\frac7{10}s_\eta^3+\frac12u_xs_\eta^4,\\
\text{d)}\quad g_\nu^4&=\frac35s_\eta^4.
\end{aligned}
\]

\begin{center}
\textbf{\S\ 8.\quad Formas de tercer orden (sistema III).}
\end{center}

\noindent\textbf{Plegamiento de $f$ con $\nu$.}

\noindent\emph{Formas irreducibles:}
\[
\begin{array}{cccc}
a_\nu&\cdot&\cdot&\cdot\\
\cdot&a_\nu^2&\cdot&\cdot\\
\cdot&\cdot&\cdot&\cdot\\
\cdot&\cdot&\cdot&a_\nu^4=i.
\end{array}
\]

\noindent\emph{Reducción:}
\[
a_\nu^3=\frac13 i u_x\qquad\text{(fórmula (2)).}
\]
\noindent\emph{Consecuencias:} 1) $a_\nu^3$ es reductor. 2) $f$ sólo entra en plegamiento con $\nu$ en productos con formas contragredientes $(j)$. Lo mismo vale para $\nu$ con respecto a $f$.

\begin{center}
\textbf{\S\ 9.\quad Formas de cuarto orden (sistema III).}
\end{center}

\noindent\textbf{Plegamiento de $\theta$ con $\nu$.}

\noindent\emph{Formas irreducibles:}
\[
\begin{array}{cccc}
(\theta\nu x)&\theta_\nu&\cdot&\cdot\\
(\theta\nu x)^2&\theta_\nu(\theta\nu x)&\theta_\nu^2&\cdot\\
\cdot&\theta_\nu(\theta\nu x)^2&\cdot&\theta_\nu^3.
\end{array}
\]

\noindent\emph{Reducciones:} Del reductor $a_\nu^3$, por doble reducción de las expresiones $a_\nu^3(abu)$, $a_\nu^3b_\nu$, $a_\nu^3b_\nu(abu)$ según el planteamiento:

\[
\begin{aligned}
\theta_\nu^2(\vartheta\nu x)
&=a_\nu^2(\widehat{a}b\nu x)(abu)-\frac13(\nu\nu_1x)^3=a_\nu b_\nu(\widehat{a}b\nu x)(abu)+\frac16(\nu\nu_1x)^3\\
&=a_\nu^3(abu)-a_\nu^2b_\nu(abu)=2a_\nu^2b_\nu(abu)=\frac23a_\nu^3(abu),\\[.4em]
\theta_\nu^2(\vartheta\nu x)^2
&=a_\nu^2(\widehat{a}b\nu x)^2-\frac13(\nu\nu_1x)^4=a_\nu b_\nu(\widehat{a}b\nu x)^2+\frac16(\nu\nu_1x)^4\\
&=if-2a_\nu^3b_\nu+a_\nu^2b_\nu^2-\frac13s=2a_\nu^3b_\nu-2a_\nu^2b_\nu^2+\frac16s=\frac23if-\frac23a_\nu^3b_\nu-\frac16s,\\[.4em]
\theta_\nu^3(\vartheta\nu x)&=a_\nu^2b_\nu(\widehat{a}b\nu x)(abu)=a_\nu^3b_\nu(abu).
\end{aligned}
\]

Se obtienen las fórmulas:
\[
\tag{9}
\begin{array}{c|c|c}
\theta_\nu^2(\theta\nu x)=0
& \theta_\nu^3(\theta\nu x)=0
& \theta_\nu^4u_x^2=u_\varrho^2=\varrho\quad(\text{módulo adjunto, \S\ 7}),\\[0.6em]
\theta_\nu^2(\theta\nu x)^2=\dfrac49 i f-\dfrac16s
& &
\end{array}
\]

\noindent\emph{Consecuencias:} 1) $\theta_\nu^2(\theta\nu x)^2$ es reductor. 2) $\nu$ sólo entra en plegamiento con $\theta$ en productos con formas contragredientes $(g)$.

\begin{center}
\textbf{\S\ 10.\quad Formas de quinto orden (sistema III).}
\end{center}

\[
\text{Plegamiento de }\left\{\begin{array}{l}
K,j,\A \text{ con }\nu,\\
f\text{ con }H.
\end{array}\right.
\]

\noindent\emph{Formas irreducibles:}
\[
\begin{array}{c@{\qquad}c@{\qquad}c@{\qquad}c}
\text{A)} & \begin{array}{cccc}
\cdot&K_\nu&\cdot&\cdot\\
(k\nu x)^2&\cdot&K_\nu^2&\cdot\\
(k\nu x)^3&K_\nu(k\nu x)^2&\cdot&\cdot\\
\cdot&K_\nu(k\nu x)^3&\cdot&\cdot
\end{array}
&
\text{B)}\ \begin{array}{c}(j\nu x)\\(j\nu x)^2\\(j\nu x)^3\\ \cdot\end{array}
&
\text{C)}\ \begin{array}{c}\A_\nu\\\cdot\\\cdot\\\cdot\end{array}
\end{array}
\]
\[
\text{D)}\qquad
\begin{array}{cccc}
(aHu)&(aHu)^2&\cdot&\cdot\\
a_\eta&a_\eta(aHu)&a_\eta(aHu)^2&\cdot\\
\cdot&a_\eta^2&\cdot&\cdot
\end{array}
\]

\noindent\emph{Reducciones:}

\noindent A) Sucesión de formas $K_\nu^3$: reductor $a_\nu^3$.

Forma $K_\nu^2(k\nu x)^2$: reductor $\theta_\nu^2(\vartheta\nu x)^2$.
\[
(k\nu x),\quad K_\nu(k\nu x),\quad K_\nu^2(k\nu x)
\]
son formas descomponibles según \S\ 3.

\noindent B) Forma
\[
\begin{aligned}
(j\nu x)^4&=(\hata\theta\nu x)^4
=a_\nu^4\theta+\theta_\nu^4 f-4a_\nu^3\theta_\nu
-4a_\nu\{a_\nu\theta_x-(\hata\theta\nu x)\}^3\\
&\quad+6a_\nu^2\{a_\nu\theta_x-(\hata\theta\nu x)\}^2
=\R f+\frac53 i\theta-6a_\nu^2(\hata\theta\nu x)^2
+4a_\nu(\hata\theta\nu x)^3,
\end{aligned}
\]
y por tanto
\[
(j\nu x)^4=\R f+\frac53 i\theta\modu{(\A,\nu)},\qquad (\text{cf. final de \S\ 5}).
\]

\noindent C) Forma $\A_\nu^4$: reductor $\theta_\nu^4$.

Formas $\A_\nu^2$ y $\A_\nu^3$: reductor $a_\nu^3$ o $\theta_\nu^4$ por sustitución de la forma $\A$ por formas inferiores de la sucesión de formas (final de \S\ 5), es decir, por doble reducción de las expresiones
\[
a_\vartheta a_\nu^3,
\qquad a_\vartheta a_\nu^3\theta_\nu,
\qquad a_\vartheta\theta_\nu^4.
\]
El doble cálculo de $\A_\nu^3$ da ocasión a la reducción de la forma superior $a_\eta^3$.

\noindent\emph{Fórmulas de reducción:}
\[
\tag{10}
\begin{aligned}
\text{a)}\quad \A_\nu^2
&=\frac35a_\eta^2-\frac3{10}(asu)^2+\frac13 i\theta
+\frac15u_x^2(2a_\R^2-t),\\
\text{b)}\quad \A_\nu^3
&=-\frac3{10}a_\R+\frac35u_x\left(\frac{13}{10}a_\R^2-\frac25t\right),\\
\text{c)}\quad \A_\nu^4&=a_\R^2-\frac25t.
\end{aligned}
\]

\noindent\emph{Derivación de las fórmulas:}
\[
\begin{aligned}
\text{a)}\quad
 a_\vartheta a_\nu^3
&=\frac13 i\theta
=[a_\vartheta]_{\nu^3}+\frac67[a_\vartheta^2]_{\nu^2}
+\frac65[a_\vartheta(a\theta u)^2]_{\nu^2}\,
u x^2
+\frac{11}{30}[a_\vartheta^2(a\theta u)^2]_\nu\,\nu x^2\\
&\quad-\frac67u_x[a_\vartheta^2]_{\nu^3}
-\frac16u_x[a_\vartheta^2(a\theta u)^2]_{\nu^2}\,\nu x^2,\\
\frac13 i\theta
&=\A_\nu^2-\frac23u_x\A_\nu^3-a_\nu a_{\nu_1}(\nu\nu_1x)^2
+\frac25a_\nu^2(\nu\nu_1x)^2
+\frac15u_x^2a_\nu^2a_{\nu_1}(\nu\nu_1x)^2;
\end{aligned}
\]
\[
\begin{aligned}
\text{b)}\quad
 a_\vartheta a_\nu^3\theta_\nu
&=0=[a_\vartheta]_{\nu^4}+\frac9{14}[a_\vartheta^2]_{\nu^3}
+\frac35[a_\vartheta(a\theta u)^2]_{\nu^2}\,\nu x^2
+\frac{17}{120}[a_\vartheta^2(a\theta u)^2]_{\nu}\,\nu x^2\\
&\quad-\frac9{14}u_x[a_\vartheta^2]_{\nu^4}
-\frac1{24}u_x[a_\vartheta^2(a\theta u)^2]_{\nu^2}\,\nu x^2,\\
0&=\frac56\A_\nu^3-\frac12u_x\A_\nu^4-\frac14a_\eta^3+\frac1{20}u_xa_\eta^4;
\end{aligned}
\]
\[
\begin{aligned}
\text{b')}\quad
 a_\vartheta\theta_\nu^4
&=a_\R=a_\vartheta\theta_\nu\{a_\nu-(a\theta\nu x)\}^3\\
&=-\frac32[a_\vartheta]_{\nu^3}+\frac35[a_\vartheta(a\theta u)^2]_{\nu^2}\,\nu x^2
-\frac{37}{120}[a_\vartheta^2(a\theta u)^2]_{\nu}\,\nu x^2\\
&\quad+\frac32u_x[a_\vartheta^2]_{\nu^4}
+\frac{49}{120}u_x[a_\vartheta^2(a\theta u)^2]_{\nu^2}\,\nu x^2,\\
a_\R&=-\frac32\A_\nu^3+\frac32u_x\A_\nu^4-\frac{11}{20}a_\eta^3+\frac7{20}u_xa_\eta^4;
\end{aligned}
\]
\[
\begin{aligned}
\text{c)}\quad
 a_\vartheta^2\theta_\nu^4
&=a_\R^2=[a_\vartheta^2]_{\nu^4}+\frac35[a_\vartheta^2(a\theta u)^2]_{\nu^2}\,\nu x^2,\\
a_\R^2&=\A_\nu^4+\frac25a_\eta^4.
\end{aligned}
\]

\noindent D) Reducción de $a_\eta^3$ por doble reducción de $\A_\nu^3$ (C).

Las demás reducciones se obtienen por un cálculo dualísticamente opuesto al de \S\ 9.

\noindent\emph{Fórmulas de reducción:}
\[
\tag{11}
\begin{array}{rl|rl}
a_\eta^2(aHu)&=-\dfrac16(asu)^3,&
 a_\eta^2(aHu)^2&=\dfrac49 i\nu-\dfrac16(asu)^4,\\[0.7em]
a_\eta^3&=-a_\R+\dfrac35u_xa_\R^2+\dfrac15u_xt,
& a_\eta^3(aHu)&=0,\\[0.7em]
&&a_\eta^4&=t\quad(\text{adjuntado como módulo}).
\end{array}
\]

\noindent\emph{Consecuencias:}

1) $(j\nu x)^4$, $\A_\nu^2$, $a_\eta^2(aHu)$ son reductores.

2) $\nu$ entra sólo en productos con formas contragredientes en plegamiento con $K,j,\A$; lo mismo vale para $f$ con respecto a $H$ y para $j$ con respecto a $\nu$, pues $\theta_\nu(\theta\cdot j,\nu,x^3)$ se vuelve reducible según \S\ 11 B.

\begin{center}
\textbf{\S\ 11.\quad Formas de sexto orden (sistema III).}
\end{center}
\[
\text{Plegamiento de }\left\{\begin{array}{l}
\theta^2,
 f\cdot j,
 N \text{ con }\nu,\\
\theta\text{ con }H.
\end{array}\right.
\]

\noindent\emph{Formas irreducibles:}
\[
\begin{array}{c@{\quad}c@{\quad}c}
\text{A)}\ (\vartheta^2\nu x)^3,\quad \theta_\nu(\vartheta^2\nu x)^3
&
\text{B)}\ a_\nu(j\nu x),\quad a_\nu^2(j\nu x)
&
\text{C)}\ N_\nu
\end{array}
\]
\[
\text{D)}\quad
\begin{array}{c|c|c|c}
 & (\theta Hu)&(\theta Hu)^2&\cdot\\
(\vartheta\eta x)&\theta_\eta&H_\vartheta(\theta Hu),\ \theta_\eta(\theta Hu)&\theta_\eta(\theta Hu)^2\\
(\vartheta\eta x)^2&H_\vartheta(\vartheta\eta x),\ \theta_\eta(\vartheta\eta x)&\theta_\eta^2&\theta_\eta^2(\theta Hu)\\
\cdot&\theta_\eta(\vartheta\eta x)^2&\cdot&\cdot
\end{array}
\]

\noindent\emph{Reducciones:}

\noindent A) $(\vartheta^2\nu x)^4$ se descompone según la fórmula (3)a:
\[
(a\A\vartheta x)^4=\frac12(\vartheta^2\nu x)^4-\frac35 f\cdot a_\nu^2(\nu\vartheta x)^2
=\frac13\A^2+f\A_\vartheta^2.
\]

\noindent B) $a_\nu(j\nu x)^2$ se descompone según la fórmula (9)a, sustituyendo, de acuerdo con el final de \S\ 6, $f\cdot j$ por $(\theta\theta'u)^2$ y teniendo en cuenta la reducibilidad de $\theta_\vartheta'$:
\[
\frac13a_\nu(j\nu x)^2=(\theta\theta'\nu x)^2\theta_\nu
=\theta_\nu^3\theta-\theta_\nu^2\theta_\nu'
=\theta_\nu^3\theta+\theta_\nu^2(\widehat{j}\nu\theta'u)
=\theta_\nu^3\theta-\frac23\theta_\nu^3\theta.
\]
Formas $a_\nu(j\nu x)^3$ y $a_\nu^2(j\nu x)^2$: reductores $a_j$ y $(j\nu x)^4$:
\[
a_\nu(j\nu x)^3=a_j(j\nu x)^3-(j\nu x)^3(j\nu\widehat{a}u),
\]
\[
a_\nu^2(j\nu x)^2=a_j^2(j\nu x)^2-2a_j(j\nu x)^2(j\nu\widehat{a}u)+(j\nu x)^2(j\nu\widehat{a}u)^2.
\]

\noindent C) Sucesión de formas $N_\nu^2$; reductor $\A_\nu^2$. $(n\nu x)$ y $N_\nu(n\nu x)$ se descomponen como transvecciones sobre determinantes funcionales según \S\ 3.

\noindent D) \emph{Vista general de las reducciones:}

a) Sucesión de formas $\theta_\eta^2H_\vartheta$: reductor $a_\eta^2(aHu)$. (Conduce a formas con símbolos superiores.)

b) Sucesión de formas $\theta_\eta H_\vartheta$: reductor $a_\eta^2(aHu)$ por doble reducción. (Conduce a formas con símbolos superiores y a formas superiores de la sucesión total de formas, es decir, a la sucesión de formas $\theta_\eta^2$.) En lugar de $\theta_\eta H_\vartheta$ consideramos la sucesión de formas
\[
\theta_\eta(\vartheta\eta x)(\theta Hu)=\theta_\eta H_\vartheta+\theta_\eta^2,
\]
y en ella reemplazamos los símbolos $\theta$ por la forma inferior $(abu)$, llegando así a la doble reducción de la sucesión de formas
\[
a_\eta^2(aHu)(abu)=\theta_\eta H_\vartheta+\frac32\theta_\eta^2\modu{s}
\]
hasta la forma final
\[
a_\eta^3b_\eta(bHu)(abu)=\theta_\eta^3H_\vartheta+\frac12\theta_\eta^4.
\]

c) Sucesión de formas $\theta_\eta^3$: por doble reducción de aquellas formas que pertenecen simultáneamente a las sucesiones $\theta_\eta H_\vartheta$ y $\theta_\eta^2H_\vartheta$, es decir, de las formas de la sucesión $\theta_\eta^2H_\vartheta$, hacia formas con símbolos superiores por una parte, y hacia formas con símbolos superiores y hacia la sucesión de formas superior $\theta_\eta^3$ por otra.

d) Formas $\theta_\eta^2(\vartheta\eta x)$, $\theta_\eta^2(\vartheta\eta x)^2$, $\theta_\eta^3(\vartheta\eta x)$: por doble reducción mediante el reductor $a$, de manera análoga al cálculo de \S\ 9.

e) Formas $\theta_\eta^2(\theta Hu)^2$ y $\theta_\eta^2(\vartheta\eta x)^2$: por doble reducción de las expresiones $\A_\nu^2(a\A u)^4$ y $(a\A\nu x)^4$ mediante los reductores $\A_\nu^2$ y $(a\A u)^4$.

f) Formas $H_\vartheta$ y $H_\vartheta^2$: formas descomponibles. Para $H_\vartheta^2$ esto resulta de la doble reducción de la expresión $\A_\nu^2(a\A u)^2$, o también del cálculo directo de los productos $(a_\nu^2)^2$, $a_\nu^2b_{\nu_1}$ por formas del sistema. (El cálculo análogo de $(a_\nu)^2$ da una relación entre formas descomponibles.)

g) Reducción de formas superiores porque formas de la sucesión $\theta\cdot H$ admiten dos posibilidades de reducción (pertenecen a dos sucesiones de formas reducibles), a saber:
\begin{quote}\small
$g$: por doble reducción de $\theta_\eta^2(\vartheta\eta x)^2$; admite las reducciones d) y e).\par
$s_\vartheta^2(\theta su)$: por doble reducción de $\theta_\eta^3(\vartheta\eta x)$; admite las reducciones c) y d).\par
$s_\vartheta^2(\theta su)^2$: por doble reducción de $\theta_\eta^2H_\vartheta^2$; admite la reducción a) y tiene como factor el reductor $\theta_\nu^2(\vartheta\nu x)^2$.\par
$s_\nu^2$: por doble reducción de $\theta_\eta^3H_\vartheta$; admite las reducciones a) y c).
\end{quote}
La demostración de la independencia de las formas restantes resulta por doble reducción de la sucesión de formas $\A_\nu^2(a\A u)^2=\theta_\eta^2$.

\noindent\emph{Ejecución de las reducciones:}

La existencia de las reducciones a), b), c), d) es clara a priori, pues según la ley de formación indicada las relaciones que surgen en la doble reducción son independientes entre sí. Para las reducciones e), f), g) debe verificarse por cálculo que no surjan identidades.

Damos, simétricamente respecto al término diagonal en la sucesión de formas $\theta\cdot H$ y avanzando hasta la forma $\theta_\eta$, en parte con indicación del cálculo, también las fórmulas obtenidas mediante las reducciones a), b), c), d), pues se emplean en parte en las reducciones e), f), g), y en parte más adelante.

\noindent 1) $H_\vartheta$ y $H_\vartheta^2$ según f). Con el planteamiento\footnote{El signo $\sim$ en lugar del signo de igualdad significa que se omiten productos de $u_x$ con formas superiores a las que surgen por los plegamientos III y IV.}
\[
\begin{aligned}
a_\nu b_{\nu_1}^2
&=\nu a_\nu b_\nu^2-(aHu)(bHu)b_\eta
\sim \nu a_\nu b_\nu^2-f a_\eta(aHu)^2-(abu)(aHu)b_\eta+(abu)^2b_\eta,\\
a_\nu^2b_{\nu_1}^2
&=b_\nu^2\{a_\nu u_\nu-(a\nu\nu_1u)\}^2
\sim \nu\{a_\nu^2b_\nu^2-2a_\nu b_\nu(aHu)(bHu)
-\frac16(asu)^2(bsu)^2\}\\
&\sim \nu a_\nu^2b_\nu^2-2a^2(aHu)(abu)-\frac16(asu)^2(bsu)^2+(\vartheta\eta x)^2(\theta Hu)^2
\end{aligned}
\]
surgen, al sustituir el valor de $\theta_\eta H_\vartheta$, las fórmulas:
\[
\tag{12}
\begin{aligned}
\text{a)}\quad H_\vartheta&\sim 2a_\nu a_\R^2+2\nu b_\nu(\vartheta\eta x)^2+2f a_\eta(aHu)^2-2\theta_\eta,\\
\text{b)}\quad H_\vartheta^2&\sim (a^2)^2+2\theta_\eta^2+\frac23(\theta su)^2
+\frac1{12}s\nu-\frac19if\nu-\frac16f(asu)^4.
\end{aligned}
\]

\noindent 2) $\theta_\eta H_\vartheta$ según b):
\[
\begin{aligned}
a_\eta^2(aHu)(abu)&\sim [a_\eta^2(aHu)]_{bu}+\frac12 f[a_\eta^2(aHu)^2]
=a_\eta b_\eta(aHu)(abu)\\
&\quad+a_\eta(a\widehat{b}\eta x)(abu)(aHu)
\sim \theta_\eta(\vartheta\eta x)(\theta Hu)+\frac12\theta_\eta^2+\frac14(\nu\eta x)^2.
\end{aligned}
\]
De las fórmulas (5) y (11) se sigue:
\[
\theta_\eta H_\vartheta\sim -\frac32\theta_\eta^2-\frac14(\theta su)^2-\frac1{12}s\nu+\frac29if\nu.
\]

\noindent 3) $\theta_\eta H_\vartheta(\theta Hu)$ se calcula según b) de modo análogo a $\theta_\eta H_\vartheta$:
\[
\theta_\eta H_\vartheta(\theta Hu)\sim-\theta_\eta^2(\theta Hu)-\frac16(\theta su)^3.
\]

\noindent 4) $\theta_\eta H_\vartheta(\vartheta\eta x)$ según b); $\theta_\eta^2(\vartheta\eta x)$ según d) (cf. \S\ 9):
\[
\theta_\eta H_\vartheta(\vartheta\eta x)\sim-\theta_\eta^2(\vartheta\eta x)-\frac16s_\vartheta(\theta su)
\sim-\frac13(\vartheta\R x)-\frac16s_\vartheta(\theta su),
\]
\[
\theta_\eta^2(\vartheta\eta x)\sim\frac23a_\eta^3(abu)\sim-\frac13(\vartheta\R x).
\]

\noindent 5)
\[
\begin{array}{ccc}
\theta_\eta H_\vartheta^2\text{ según b)},&
\theta_\eta^2H_\vartheta\text{ según a)},&
\theta_\eta^2H_\vartheta\text{ según b) y con ello }\theta_\eta^3\text{ según c)},\\
\text{de }a_\eta^2(aHb)(bHu);&
\text{de }a_\eta^3(Hab)(abu);&
\text{de }a_\eta^3(bHu)(abu).
\end{array}
\]
\[
\theta_\eta H_\vartheta^2\sim-\frac32\theta_\eta^2H_\vartheta-\frac14s_\vartheta(\theta su)^2+\frac16s_\nu-\frac49ia_\nu
\sim-\theta_\R-\frac16s_\nu-\frac19ia_\nu,
\]
\[
\theta_\eta^2H_\vartheta\sim\frac23\theta_\R-\frac16s_\vartheta(\theta su)^2+\frac29s_\nu-\frac29ia_\nu,
\]
\[
\theta_\eta^3H_\vartheta\sim\frac23\theta_\R-\frac23\theta_\eta^3+\frac5{36}s_\nu,
\qquad
\theta_\eta^3\sim\frac14s_\vartheta(\theta su)^2-\frac18s_\nu+\frac13ia_\nu.
\]

\noindent 6) $\theta_\eta^2(\theta Hu)^2$ según e):
\[
\A_\nu^2(a\A u)^2=[(a\A u)^4]_{\nu^2}
=-\frac{12}{5}\theta_\eta^2(\theta Hu)^2-\frac3{10}\sigma-\frac1{20}(\theta su)^4+\nu\R
\quad(\text{según la fórmula (3)c}),
\]
\[
\A_\nu^2(a\A u)^4=[\A_\nu^2]_{aa}^4=-\frac35\theta_\eta^2(\theta Hu)^2+\frac15\sigma-\frac3{10}(\theta su)^4+\frac13ij
\quad(\text{según la fórmula (10)a}),
\]
\[
\theta_\eta^2(\theta Hu)^2=-\frac16\sigma+\frac1{12}(\theta su)^4-\frac19ij+\frac13\nu\R.
\]

\noindent 7) $\theta_\eta^2(\vartheta\eta x)^2$ según d) y e). De aquí resulta la reducción de la forma superior $g$.

Por un cálculo análogo al de \S\ 9 se sigue:
\[
\begin{aligned}
\theta_\eta^2(\vartheta\eta x)^2
&=\frac12 i f-\frac12a_\eta^2b_\eta-\frac1{12}g
=\frac23tf-\frac23a_\eta^3b_\eta-\frac16g\\
&=-\frac16g-\frac13(\vartheta\R x)^2+\frac4{15}fa_\R^2+\frac8{15}ft
\quad(\text{según la fórmula (11)}).
\end{aligned}
\]
Por el cálculo correspondiente como arriba se sigue:
\[
((a\A\nu x)^4)=[(a\A u)^4]_{\nu x^4}
=\frac{21}{10}\theta_\eta^2(\vartheta\eta x)^2+\frac7{10}s_\vartheta^2-\frac3{10}g
\quad(\text{según la fórmula (3)c}),
\]
\[
\begin{aligned}
(a\A\nu x)^4&=-\frac13i\A+6[\A_\nu^2]a^2-4[\A_\nu^3]a+\A_\nu^4f\\
&= -\frac{36}{5}\theta_\eta^2(\vartheta\eta x)^2-\frac95s_\vartheta^2-\frac35g-\frac35(\vartheta\R x)^2\\
&\quad+\frac53i\A+\frac{37}{25}fa_\R^2+\frac{74}{25}ft\quad(\text{según la fórmula (10)}).
\end{aligned}
\]
Por tanto
\[
\begin{aligned}
\frac{93}{10}\theta_\eta^2(\vartheta\eta x)^2
&=-\frac3{10}g-\frac52s_\vartheta^2-\frac35(\vartheta\R x)^2\\
&\quad+\frac53i\A+\frac{37}{25}fa_\R^2+\frac{74}{25}ft,
\end{aligned}
\]
y al comparar los dos valores de $\theta_\eta^2(\vartheta\eta x)^2$ según g):
\[
\tag{13}
0=g-2s_\vartheta^2+2(\vartheta\R x)^2+\frac43i\A-\frac45fa_\R^2-\frac85ft.
\]

\noindent 8) $\theta_\eta^2H_\vartheta(\theta Hu)$ según a) y b); de aquí, $\theta_\eta^3(\theta Hu)$ según c) (las formas con el factor $u_x$ desaparecen):
\[
\theta_\eta^2H_\vartheta(\theta Hu)=-\frac16s_\vartheta(\theta su)^3-\frac13(\nu\R x),
\]
\[
\theta_\eta^2H_\vartheta(\theta Hu)=-\frac13\theta_\eta^3(\theta Hu)-\frac13(\nu\R x),
\]
\[
\theta_\eta^3(\theta Hu)=\frac12s_\vartheta(\theta su)^3.
\]

\noindent 9) $\theta_\eta^2H_\vartheta(\vartheta\eta x)$ según a) y b); de aquí, $\theta_\eta^3(\vartheta\eta x)$ según c). $\theta_\eta^3(\vartheta\eta x)$ según d); de aquí, la reducción de la forma superior $s_\vartheta^2(\theta su)$ según g) (las formas con el factor $u_x$ desaparecen):
\[
\theta_\eta^2H_\vartheta(\vartheta\eta x)=\frac13(atu),
\]
\[
\theta_\eta^2H_\vartheta(\vartheta\eta x)=\frac13(atu)+\frac13\theta_\eta^3(\vartheta\eta x)-\frac16s_\vartheta^2(\theta su),
\]
\[
\theta_\eta^3(\vartheta\eta x)=\frac12s_\vartheta^2(\theta su),
\]
\[
\theta_\eta^3(\vartheta\eta x)=\frac25\theta_\R(\vartheta\R x)+\frac15(atu),
\]
\[
\tag{14}
s_\vartheta^2(\theta su)=\frac45\theta_\R(\vartheta\R x)+\frac25(atu).
\]

\noindent 10) $\theta_\eta^2H_\vartheta^2$ según a) y mediante el reductor $\theta_\nu^2(\vartheta\nu x)^2$; $\theta_\eta^3H_\vartheta$ según a) y c); $\theta_\eta^4$ según c); de aquí, la reducción de las formas superiores $s_\vartheta^2(\theta su)^2$ y $s_\nu^2$ según g):
\[
\begin{aligned}
\text{Según a)}\quad
\theta_\eta^2H_\vartheta^2&=[a_\eta^2(aHu)^2]b^2+\frac13[a_\eta^4]_{bu^2}-\frac13H_\nu^2(\nu\eta x)^2\\
&=\frac49ia_\nu^2-\frac16s_\vartheta^2(\theta su)^2-\frac5{18}s_\nu^2+\frac13(atu)^2+\frac1{54}Ju_x^2.
\end{aligned}
\]
\[
\begin{aligned}
\text{Mediante }\theta_\nu^2(\vartheta\nu x)^2:
\quad\theta_\eta^2H_\vartheta^2
&=[\theta_\nu^2(\vartheta\nu x)^2]_{\nu_1}+\frac13[\theta_\nu^4]_{\nu_1}\hat\nu_1x^2-\frac13s_\vartheta^2(\theta su)^2\\
&=\frac49ia_\nu^2-\frac16s_\nu^2-\frac13s_\vartheta^2(\theta su)^2+\frac13(\nu\R x)^2.
\end{aligned}
\]
\[
\begin{aligned}
\text{Según a)}\quad
\theta_\eta^3H_\vartheta&=-[a_\eta^2(aHu)]_{bu}b^2-\frac12[a_\eta^2(aHu)^2]b^2-\frac13[a_\eta^3]b+\frac1{12}[a_\eta^4]_{bu^2}\\
&\quad+\frac12u_x[a_\eta^2(aHu)^2]b^3\\
&=-\frac29ia_\nu^2+\frac1{12}s_\nu^2+\frac4{15}\theta_\R^2-\frac7{90}(\nu\R x)^2+\frac1{30}(atu)^2+\frac2{27}i^2u_x^2-\frac1{36}Ju_x^2.
\end{aligned}
\]
\[
\begin{aligned}
\text{Según c)}\quad
\theta_\eta^3H_\vartheta:
 a_\eta^3(Hab)(bHu)&=\frac12(atu)^2
=\theta_\eta^2H_\vartheta(\theta Hu)(\vartheta\eta x)+\frac14H_\nu^2(\nu\eta x)^2\\
&\quad+\frac12\theta_\eta^2H_\vartheta^2
=\frac32\theta_\eta^2H_\vartheta^2+\theta_\eta^3H_\vartheta-u_x[a_\eta^2(aHu)^2]b^3\\
&\quad+\frac14H_\nu^2(\nu\eta x)^2,
\end{aligned}
\]
y por tanto
\[
\begin{aligned}
\theta_\eta^3H_\vartheta&=-\frac23ia_\nu^2+\frac5{12}s_\nu^2+\frac12(\nu\R x)^2-\frac12(atu)^2\\
&\quad+\frac4{27}i^2u_x^2-\frac1{12}Ju_x^2.
\end{aligned}
\]
Según c),
\[
\begin{aligned}
\theta_\eta^4&=\frac49ia_\nu^2-\frac16s_\nu^2+\frac{16}{15}\theta_\R^2\\
&\quad-\frac{14}{45}(\nu\R x)^2+\frac2{15}(atu)^2.
\end{aligned}
\]
Según g), a partir de los dos valores para $\theta_\eta^2H_\vartheta^2$:
\[
\tag{15}
\begin{aligned}
s_\vartheta^2(\theta su)^2
&=\frac23s_\nu^2-2(atu)^2+2(\nu\R x)^2-\frac19Ju_x^2\\
&=\frac89ia_\nu^2+\frac8{15}\theta_\R^2-\frac{14}{15}(atu)^2+\frac{38}{45}(\nu\R x)^2-\frac4{27}i^2u_x^2.
\end{aligned}
\]
Según g), a partir de los dos valores para $\theta_\eta^3H_\vartheta$:
\[
\tag{16}
s_\nu^2=\frac43ia_\nu^2+\frac45\theta_\R^2+\frac85(atu)^2-\frac{26}{15}(\nu\R x)^2+\frac16Ju_x^2-\frac29i^2u_x^2.
\]
La fórmula (16) proporciona la base para la reducción del módulo $(ss'u)^4$, y la fórmula (13) para la reducción del sistema de $s$ (\S\ 17).

\noindent\emph{Consecuencias:}

1) $a_\nu(j\nu x)^3$, $\theta_\eta H_\vartheta$, $\theta_\eta^2(\vartheta\eta x)$, $\theta_\eta^2(\theta Hu)^2$ son reductores; $a_\nu(j\nu x)^2$, $H_\vartheta$ y $H_\vartheta^2$ son formas descomponibles.

2) La forma del sistema II, $j(\nu)=g=(\nu\eta x)^4H_x^2$, es reducible.

3) Como la única forma contragrediente $g$ se vuelve reducible, $\nu$ no entra en absoluto, ni en potencias ni en productos con formas del mismo sistema, en plegamiento con el sistema I.

$\theta$ entra en plegamiento sólo considerada como contravariante en productos o potencias, y correspondientemente $H$ sólo considerada como covariante (cf. \S\ 13 B).




\begin{center}
\textbf{\S\ 12.\quad Formas de séptimo orden.}
\end{center}
\[
\text{Plegamiento de }\left\{\begin{array}{l}
\theta\cdot K\text{ con }\nu,\\
K,j,\A\text{ con }H,\\
f\text{ con }L,\sigma.
\end{array}\right.
\]

\noindent\emph{Formas irreducibles:}
\[
\text{B)}\quad
\begin{array}{c|c|c|c|c}
\cdot&\cdot&(KHu)^2&\cdot&\cdot\\
\cdot&K_\eta&\cdot&K_\eta(KHu)^2&\cdot\\
(k\eta x)^2&\cdot&K_\eta^2&\cdot&\cdot\\
(k\eta x)^3&H_k(k\eta x)^2,\ K_\eta(k\eta x)^2&\cdot&\cdot&\cdot\\
\cdot&K_\eta(k\eta x)^3&\cdot&\cdot&\cdot
\end{array}
\]
\[
\text{C)}\quad
\begin{array}{cc}
(j\eta x)&H_j\\
\cdot&H_j(j\eta x)\\
\cdot&H_j^2
\end{array}
\qquad
\text{D)}\quad
\begin{array}{c}
(\A Hu)\\\A_\eta
\end{array}
\]
\[
\text{E)}\quad
\begin{array}{cccc}
a_i&(aLu)^2&(aLu)^3&\cdot\\
\cdot&\cdot&a_i(aLu)^2&a_i(aLu)^3\\
\cdot&a_i^2&\cdot&\cdot
\end{array}
\]

\noindent\emph{Reducciones:}

\noindent A) $(\vartheta\cdot k,\nu,x)^4$ se descompone como una transvección simple sobre la forma descomponible $(\vartheta^2\nu x)^4$.

\noindent B) Sucesión de formas $H_kK_\eta$: reductor $H_\vartheta\theta_\eta$.

Sucesión de formas $K_\eta^2(KHu)$: reductor $a_\eta^2(aHu)$.

Sucesión de formas $K_\eta^2(k\eta x)$: reductor $\theta_\eta^2(\vartheta\eta x)$.

La doble reducción de la sucesión de formas $K_\eta^3$ que resulta de ahí da ocasión a la reducción de la sucesión superior $(j\eta x)^3$.

$(KHu)$, $(k\eta x)$, $K_\eta(k\eta x)$ son formas descomponibles según \S\ 3.

$H_k(KHu)$, $K_\eta(KHu)$ se descomponen como surgidas por plegamiento simple con la forma descomponible $K_\nu(k\nu x)$ y el determinante funcional
\[
K_\nu^2=\theta_\nu^2(a\theta u)\equiv a_\nu^2(a\theta u)\modu{\nu^2}.
\]
A la doble representación de $K_\nu^2$ corresponde una relación entre formas descomponibles. Se obtiene:
\[
H_k(KHu)=2K_\nu(\nu\nu_1k)=2\theta_{\nu_1}(\nu\nu_1a\theta)
=2\theta_\nu^2a_\nu-a_\nu^2\theta_\nu+f\theta_\eta(\theta Hu)^2-f\nu\theta_\vartheta^3
\modu{\nu^3},
\]
\[
\begin{aligned}
K_\eta(KHu)&=-K_\nu^2(\nu\nu_1x)
=-a_\nu^2(a\theta u)(\nu\nu_1x)\\
&=a_\eta(aHu)^2\theta+a_\nu^2\theta_\nu
=-\theta_\nu(\theta Hu)^2f-\theta_\nu^2a_\nu.
\end{aligned}
\]
Por tanto
\[
\tag{17}
0=a_\nu^2\theta_\nu+\theta_\nu^2a_{\nu_1}+f\theta_\eta(\theta Hu)^2+\theta a_\eta(aHu)^2\modu{\nu^3}.
\]
Las formas $H_k$, $H_k(k\eta x)$, $H_k^2$, $H_k^2(k\eta x)$ se descomponen como surgidas por plegamiento simple con las formas descomponibles $H_\vartheta$ y $H_\vartheta^2$ (cf. \S\ 3. 2), a) y b)).

Se tiene
\[
H_k=H_\vartheta(a\theta u)\sim [H_\vartheta]_{ua}-fH_\vartheta(\theta Hu)
\]
(especialización $y=uv$ de 2)a).
\[
H_k(k\eta x)=H_\vartheta a_\eta-H_\vartheta\theta_\eta f,
\qquad
H_\vartheta a_\eta=\frac54[H_\vartheta]a-\frac14H_\vartheta a_\vartheta
\quad(\S\ 3. 2)b),
\]
\[
H_\vartheta^2(a\theta u)=[H_\vartheta^2]_{ua},
\qquad
H_\vartheta^2a_\eta=[H_\vartheta^2]a=H_k^2(k\eta x)+H_\vartheta^2\theta_\eta f.
\]

\noindent C) Sucesión de formas $(j\eta x)^3$: reducible por doble reducción de la sucesión de formas $K_\eta^3$ según B), de acuerdo con el planteamiento:
\[
0=a_\eta^3(a\theta u)=\theta_\eta^3(a\theta u)
=\theta_\eta+(a\theta\eta x)^3(a\theta u)-\theta_\eta^3(a\theta u)
=\frac12(j\eta x)^3\modu{(\nu^3,s,\R,t,i,J)}.
\]
El planteamiento análogo vale hasta la forma final de la sucesión:
\[
0=H_\vartheta^2(a\theta\eta x)a_\eta^3
=H_\vartheta^2(a\theta\eta x)\theta_\eta^3
-\frac12H_\vartheta^2(j\eta x)^4\modu{(\nu^3,s,\R,t,i,J)}.
\]
Forma $(j\eta x)^2$: se descompone 1) por doble polarización de la relación para la forma descomponible $H_\vartheta^2$ ((12.) b);

2) por cálculo directo del producto $a_\nu\theta_\nu^3$; de ahí la descomposición de la forma superior $(\A Hu)^2$. Se obtiene:
\[
\tag{18}
\left.\begin{array}{ll}
\text{a)}&\dfrac13(\A Hu)^2+\dfrac13(j\eta x)^2=\theta_\nu^2a_{\nu_1}^2,\\[0.6em]
\text{b)}&(j\eta x)^2=\dfrac43\theta_\nu^3a_{\nu_1},
\end{array}\right\}\modu{(\nu^3,s,\R,t,i,J)}.
\]
según el planteamiento:
\[
H_\vartheta^2(a\theta u)^2-\frac13(\A Hu)^2
=2\theta_\eta^2(a\theta u)(aHu)+\theta_\nu^2a_{\nu_1}^2
=(a\theta u)(aHu)\{a_\eta-(a\theta\eta x)^2\}+\theta_\nu^2a_{\nu_1}^2,
\]
\[
\begin{aligned}
a_{\nu_1}\theta_\nu^3
&=-(a\widehat{\nu}\nu_1u)\theta_\nu^2(\theta\widehat{\nu}\nu_1u)
-\frac12(u\nu\nu_1u)\theta_\nu\theta_{\nu_1}(\vartheta\nu\nu_1u)\\
&=-\frac32\theta_\eta^2(\theta Hu)(aHu)
=-\frac32\theta_\eta^2(\theta Hu)(a\theta u)\\
&=-\frac32\{a_\eta-(a\theta\eta x)^2\}\{(aHu)-(a\theta u)\}(a\theta u).
\end{aligned}
\]

\noindent D) Sucesión de formas $\A_\eta(\A Hu)$: reductor $\A_\nu^2$.

$(\A Hu)^2$ es forma descomponible según C).

\noindent E) Sucesión de formas $a_i^2(aLu)$: reductor $a_\eta^2(aHu)$. Las formas $(aLu)$ y $a_i(aLu)$ se descomponen como transvecciones sobre determinantes funcionales según \S\ 3.

\noindent F) Sucesión de formas $a_\sigma^3$: reductor $a_\R^3$.

Formas $a_\sigma^2$ y $a_\sigma^3$: reductores $a_\nu^3$ y $a_\eta^3$ por doble reducción de las expresiones correspondientes
\[
H_\nu a_\nu^3\quad\text{y}\quad H_\nu a_\eta^3,
\qquad
H_\nu a_\nu^3a_\nu\quad\text{y}\quad H_\nu a_\eta^4
\]
(cf. \S\ 10. C), reducción de $\A_\nu^2$ y $\A_\nu^3$).

De ahí la reducción para las formas superiores $a_\nu s_\nu(asu)^2$, $a_\nu^2s_\nu(asu)^2$ y para formas en los símbolos $\R,t(a_\nu,a_\eta^2)$, teniendo en cuenta (16.):
\[
\tag{19}
\left.\begin{array}{rcl}
a_\nu s_\nu(asu)^2&=&\dfrac13 i\theta_\nu^2-\dfrac1{18}iH,\\[0.5em]
a_\nu^2s_\nu(asu)^2&=&0,
\end{array}\right\}\modu{(\R,t)}.
\]
La forma $a_\sigma$ se descompone por polarización de (12.)b o, más brevemente, por desarrollo directo del producto $a_\nu^2\theta_\nu^3$ según el planteamiento:
\[
a_\nu^2\theta_{\nu_1}^3=(a_\nu^2\theta_{\nu_1})a_{\nu_1}-(a\theta\nu_1x)^2
=-2a_\eta(aHu)(a\theta H)(a\theta\eta x)+(a\theta\nu_1x)^2a_\nu^2\theta_{\nu_1}
\]
y teniendo en cuenta la reducción de $H_j(j\eta x)^2$ (C):
\[
\tag{20}a_\sigma=2a_\nu^2\theta_{\nu_1}^3\modu{(s,\R,t,i)}.
\]

\noindent\emph{Consecuencias:}

1) $(j\eta x)^3$, $\A_\eta(\A Hu)$, $a_\sigma^2$ son reductores; $(j\eta x)^2$, $(\A Hu)^2$, $a_\sigma$ son formas descomponibles.

2) $f$ entra en plegamiento con $L$ sólo en productos con formas contragredientes; $f$, y por tanto también $K$ y $N$, no entran en absoluto en plegamiento con $\sigma$ y, por consiguiente, tampoco con $(\nu\sigma x)$. $j$ entra sólo en productos con formas contragredientes; $\A$ no entra en absoluto en productos en plegamiento con $H$ y $L$ (esto se sigue de $\A_\eta(\A Hu)$ y $(\A Hu)^2$). Del mismo modo $\sigma$, y por tanto $(\nu\sigma x)$, no entra en productos en plegamiento con ninguna forma del sistema I. En los productos del sistema II, teniendo en cuenta las consecuencias de \S\ 11, sólo quedan potencias de $H$ y productos de $H$ con $L$.

\begin{center}
\textbf{\S\ 13.\quad Formas de octavo orden (sistema III).}
\end{center}
\[
\text{Plegamiento de }\left\{\begin{array}{l}
\A\cdot j\text{ con }\nu,\\
\theta^2,\\
 f\cdot j,\\
 N\text{ con }H,\\
\theta\text{ con }L,\sigma.
\end{array}\right.
\]

\noindent\emph{Formas irreducibles:}
\[
\begin{array}{c@{\quad}l@{\qquad}c@{\quad}l}
\text{A)}&\A_\nu(j\nu x)&
\text{B)}&\begin{array}{c}(\vartheta^2\eta x)^3\\(\vartheta^2\eta x)^4\end{array}
\quad H_\vartheta(\vartheta^2\eta x)^3,
\ \theta_\eta(\vartheta^2\eta x)^3,\\[1em]
\text{C)}&(aHu)H_j,
\quad a_\eta(aHu)H_j&
\text{D)}&H_n,
\ N_\eta,
\end{array}
\]
\[
\text{E)}\quad
\begin{array}{c|c|c|c}
\cdot&(\theta Lu)^2&(\theta Lu)^3&\cdot\\
\theta_i&\cdot&L_\vartheta(\theta Lu)^2,
\ \theta_i(\theta Lu)^2&\theta_i(\theta Lu)^3\\
(\vartheta lx)^2&\theta_i^2&\cdot&\cdot\\
\cdot&\theta_i(\vartheta lx)^2&\cdot&\cdot
\end{array}
\]

\noindent\emph{Reducciones:}

\noindent A) $\A_\nu(j\nu x)^3$: reductor $a_\nu(j\nu x)^3$.

$\A_\nu(j\nu x)^2$ es reducible por la forma descomponible $a_\nu(j\nu x)^2$ según \S\ 3. 2)b, teniendo en cuenta el reductor $\theta_{\vartheta j}$. En efecto, según \S\ 11 B) se obtiene:
\[
\frac3{10}\A_\nu(j\nu x)^2+\frac35a_\nu(j\nu x)a_\vartheta(j\nu x)
+\frac1{10}a_\nu(j\nu x)^2
=\frac35\theta_\nu^3\theta_{\vartheta j}+\frac25\theta_\nu^3\theta_{\vartheta j}\theta_{\vartheta j},
\]
(reductores $a_j$ y $\theta_{\vartheta j}$).

\noindent B) Las formas $H_\vartheta(\vartheta'\eta x)^2$, $H_\vartheta^2(\vartheta'\eta x)$, $H_\vartheta^2(\vartheta'\eta x)^2$ se descomponen como surgidas por plegamiento simple sucesivo con las formas descomponibles $H_\vartheta$ y $H_\vartheta^2$, respectivamente $H_\vartheta a_\eta$ y $H_\vartheta^2a_\eta$, según \S\ 3. 2)b) y según la relación:
\[
H_\vartheta(\vartheta'\eta x)^2=2H_\vartheta a_\eta^2f-2H_\vartheta a_\eta b_\eta.
\]

\noindent C) Sucesión de formas $a_\eta(j\eta x)^2$: reductores $a_j$ y $(j\eta x)^3$.

La forma $a_\eta(j\eta x)$ se descompone: reductor $a_j$ y plegamiento simple con la forma descomponible $(j\eta x)^2$ según \S\ 3, 2)a (especialización $y=uv$).

La forma $a_\eta H_j$ se descompone: 1) por plegamiento simple con la forma descomponible $a_\nu(j\nu x)^2$;

2) por plegamiento doble especial con la forma descomponible $(j\eta x)^2$; de ahí una relación entre formas descomponibles.

A partir de $a_\nu(j\nu x)^2$:
\[
0=\theta_\nu'\theta_\nu+2a_\eta H_j\modu{(s,\R,t,i,J)}.
\]
A partir de $(j\eta x)^2$:
\[
0=\theta_\nu'\theta_\nu-\frac32a_\eta H_j\modu{(s,\R,t,i,J)}
\]
(se omiten los productos con contravariantes), según el planteamiento:
\[
0=\frac12a_\nu(abu)^2\theta_\nu^3
+\frac12u_x(ubu)\theta_{\nu_1}^3(\theta bu)
-\frac34(\widehat{j}\eta bu)^2+\frac34H_j(\widehat{j}\eta bu)
-\frac18fH_j^2
\]
y permutando los símbolos en $a_\nu(ubu)(\theta bu)\theta_{\nu_1}^3$.

\noindent D) Sucesión de formas $N_\eta(NHu)$: reductor $\A_\nu^2$.

$(NHu)$, $(n\eta x)$, $N_\eta(n\eta x)$, $H_\eta(NHu)$ son formas descomponibles (cf. \S\ 12 B); $(NHu)^2$ se descompone por plegamiento simple con la forma $(\A Hu)^2$.

\noindent E) Sucesiones de formas $\theta_iL_\vartheta$, $\theta_i^2(\theta Lu)^2$, $\theta_i^2(\vartheta lx)$:

Reductores: $\theta_\eta H_\vartheta$, $\theta_\eta^2(\theta Hu)^2$, $\theta_\eta^2(\vartheta\eta x)$.

Las formas $(\theta Lu)$, $(\vartheta lx)$, $L_\vartheta$, $L_\vartheta(\theta Lu)$, $\theta_i(\theta Lu)$, $L_\vartheta(\vartheta lx)$, $\theta_i(\vartheta lx)$, $L_\vartheta^2$, $L_\vartheta^2(\theta Lu)$, $\theta_i^2(\theta Lu)$ se descomponen. (Dualísticamente opuestas a las formas descomponibles de \S\ 12 B).)

\noindent F) Sucesión de formas $\theta_\sigma(\vartheta\sigma x)$: reductor $a_\sigma^2$.

Las formas $(\vartheta\sigma x)$ y $(\vartheta\sigma x)^2$ se descomponen, como surgidas por plegamiento simple con la forma descomponible $a_\sigma$; $\theta_\sigma$, surgida por plegamiento doble, se descompone, pues la sucesión de formas
\[
a_{\nu_2}^2(a\theta u)\theta_{\nu_1}^3\equiv H_j^2(j\eta x)\equiv0\modu{u_x(s,\R,t,i,J)}:
\]
\[
\theta_\sigma=a_\sigma(abu)^2=a_\nu^2(abu)^2\theta_\nu^3.
\]

\noindent\emph{Consecuencias:}

$\theta^3$ o $\theta^2K$ no entra en plegamiento con $H$; $\theta$, considerada sólo como contravariante, entra en potencias o productos en plegamiento con $L$, pero no entra en absoluto en plegamiento con $\sigma$ y $(\nu\sigma x)$.

$N$ no entra en productos en plegamiento con ninguna forma del sistema II, ni tampoco con potencias o productos de formas del sistema II. En los productos del sistema I, teniendo en cuenta las consecuencias de los últimos párrafos, sólo quedan los productos $f\cdot j$, $\A\cdot j$, las potencias de $\theta$ y los productos $\theta\cdot K$.

\begin{center}
\textbf{\S\ 14.\quad Formas de noveno orden (sistema III).}
\end{center}
\[
\text{Plegamiento de }\left\{\begin{array}{l}
\theta\cdot K\text{ con }H,\\
K,j,\A\text{ con }L,\\
j,\A\text{ con }\sigma,\\
f\text{ con }H^2.
\end{array}\right.
\]

\noindent\emph{Formas irreducibles:}
\[
\begin{array}{c@{\quad}l@{\qquad}c@{\quad}l}
\text{A)}&(\vartheta\cdot k,\eta,x)^4,
\ H_k(\vartheta\cdot k,\eta,x)^4&
\text{B)}&(KLu)^3,
\ K_i(KLu)^3,
\ K_i^2,
\ (klx)^3,
\ K_i(Klx)^3,\\[0.7em]
\text{C)}&L_j,
\ L_j^2&
\text{D)}&\A_i,\\[0.7em]
\text{E)}&(j\sigma x)&
\text{G)}&(aH^2u)^3,
\ (aH^2u)^4,
\ a_\eta(aH^2u)^3.
\end{array}
\]

\noindent\emph{Reducciones:}

A) Las formas $H_\vartheta(k\eta x)^3$, $H_\vartheta^2(k\eta x)^2$, $H_\vartheta^2(k\eta x)^3$ se descomponen como surgidas por plegamiento simple sucesivo con formas descomponibles.

B) Sucesiones de formas $K_iL_k$, $K_i^2(KLu)$, $K_i^2(klx)$:

Reductores: $\theta_\eta H_\vartheta$, $a_\eta^2(aHu)$, $\theta_\eta^2(\vartheta\eta x)$.

Las formas $L_i$, $K_i(KLu)$, $K_i(klx)$, $(KLu)^2$, $(klx)^2$ se descomponen como transvecciones sobre determinantes funcionales (\S\ 3).

Las formas $L_k$, $L_k(KLu)$, $L_k(KLu)^2$, $L_k(klx)$, $L_k(klx)^2$, $L_k^2$, $L_k^2(KLu)$, $L_k^2(klx)$, $L_k^3$, $K_i(KLu)^2$, $K_i(klx)^2$ se descomponen como surgidas por plegamiento simple con las formas descomponibles:
\[
L_\vartheta,
\ H_k(KHu),
\ L_\vartheta(\vartheta lx),
\ H_k^2,
\ L_\vartheta^2,
\ L_\vartheta^2(\theta Lu),
\ K_\eta(KHu),
\ \theta_i(\vartheta lx)
\]
según \S\ 3. 2), a) y b).

C) Sucesión de formas $(jlx)^3$: reductor $(j\eta x)^3$.

Las formas $(jlx)^2$ y $L_j(jlx)$ surgen por plegamiento simple con la forma descomponible $(j\eta x)^2$.

D) Sucesión de formas $\A_i(\A Lu)$: reductor $\A_\nu^2$.

Formas $(\A Lu)^2$, $(\A Lu)^3$: plegamiento simple con la forma descomponible $(\A Hu)^2$.

E) Sucesión de formas $(j\sigma x)^2$: reductor $a_\sigma^2$ por sustitución de $j$ por formas inferiores de la sucesión de formas (\S\ 5):
\[
a_\sigma^2(a\theta u)^2=\frac13(j\sigma x)^2,
\qquad
a_\sigma^2(a\theta\sigma x)^2=\frac13(j\sigma x)^4.
\]
F) Sucesión de formas $\A_\sigma^2$: reductor $a_\sigma^2$.

Forma $\A_\sigma$: reductor $a_\sigma^2$ por sustitución de $\A$ por formas inferiores de la sucesión:
\[
a_\sigma a_\R^2=\frac23\A_\sigma\modu{\nu}.
\]

\noindent\emph{Consecuencias:} Sólo la forma $j$ entra en plegamiento con $\sigma$ y $(\nu\sigma x)$.

$N$ no entra en plegamiento con $L$, pues la forma $N_i$ correspondiente a $\A_i$ se descompone.

\begin{center}
\textbf{\S\ 15.\quad Formas de décimo orden (sistema III).}
\end{center}
\[
\text{Plegamiento de }\left\{\begin{array}{l}
\theta^2,\\
\ f\cdot j\text{ con }L,\\
\theta\text{ con }H^2.
\end{array}\right.
\]

\noindent\emph{Formas irreducibles:}
\[
\begin{array}{c@{\quad}l@{\qquad}c@{\quad}l}
\text{A)}&(\vartheta^2lx)^4,
\ \theta_i(\vartheta^2lx)^4&
\text{B)}&(aLu)^2L_j,
\ a_i(aLu)^2L_j,\\[0.8em]
\text{C)}&(\theta H^2u)^3,
\ (\theta H^2u)^4,
\ H_\vartheta(\theta H^2u),
\ \theta_\eta(\theta H^2u)^3.
\end{array}
\]

\noindent\emph{Reducciones:}

A) y B): descomposición de las restantes formas por plegamiento simple con formas descomponibles.

C) Descomposición de las restantes formas según \S\ 13 B), por la consideración dualísticamente opuesta.

\noindent\emph{Consecuencias:}

$\theta^2K$ no entra en plegamiento con $L$; correspondientemente, $H^2L$ no entra en plegamiento con $K$. $H^3$ y $H^2L$ no entran en plegamiento con $\theta$. Con ello queda demostrado el esquema de plegamiento de \S\ 7.

\begin{center}
\textbf{\S\ 16.\quad Formas de orden 11, 12, 13 y 15 (sistema III).}
\end{center}

\noindent\emph{Orden 11.}
\[
\text{Plegamiento de }\left\{\begin{array}{l}
\theta\cdot K\text{ con }L,\\
K,j\text{ con }H^2,\\
j\text{ con }(\nu\sigma x),\\
f\text{ con }H\cdot L.
\end{array}\right.
\]
\noindent\emph{Formas irreducibles:}
\[
\begin{array}{c@{\quad}l@{\qquad}c@{\quad}l}
\text{A)}&(\vartheta\cdot k,l,x)^5,
\ \theta_i(\vartheta\cdot k,l,x)^5&
\text{B)}&(KH^2u)^4,
\ K_\eta(KH^2u)^4,\\[0.8em]
\text{C)}&(\nu\sigma j),&
\text{D)}&(a_\eta H\cdot L,u)^4.
\end{array}
\]

\noindent\emph{Orden 12.}
\[
\text{Plegamiento de }\left\{\begin{array}{l}
\theta^3\text{ con }L,\\
\theta\text{ con }H\cdot L.
\end{array}\right.
\]
\noindent\emph{Formas irreducibles:}
\[
\text{E)}\ (\vartheta^3lx)^6,
\qquad
\text{F)}\ (\theta,H\cdot L,u)^4,
\qquad L_\vartheta(\theta,H\cdot L,u^4).
\]

\noindent\emph{Orden 13.}
\[
\text{Plegamiento de }K\text{ con }H\cdot L.
\]
\noindent\emph{Formas irreducibles:}
\[
\text{G)}\ (K,H\cdot L,u)^5,
\qquad K_\eta(K,H\cdot L,u)^5.
\]

\noindent\emph{Orden 15.}
\[
\text{Plegamiento de }K\text{ con }H^3.
\]
\noindent\emph{Forma irreducible:}
\[
\text{H)}\ (KH^3u)^6.
\]

\noindent\emph{Reducciones:}

Las formas $H_j^2$, $H_j^2$ tienen el reductor $L_j^3$, a partir de la relación:
\[
(\nu lx)L_x^3=-\frac12H^2\modu{(\sigma,s)}.
\]
La descomposición o reducibilidad de las restantes formas sigue de las consideraciones de los párrafos anteriores.

\noindent\emph{Consecuencias:}

Según el esquema de plegamiento de \S\ 7 y las consecuencias de \S\S\ 8--15, quedan con esto agotadas las formas irreducibles del sistema III (117 formas).

\clearpage

\begin{center}
{\Large\bfseries Capítulo IV.\quad El sistema relativamente completo módulo $(\varrho,t)$.}
\end{center}

\begin{center}
\textbf{\S\ 17.\quad Reducción del módulo $(ss'u)^4$ y del sistema de $s$.}
\end{center}

Para formar el sistema relativamente completo inmediatamente superior, por analogía con \S\ 7 hay que formar el sistema de $s$ con respecto al módulo siguiente
\[
\nu(s)=(ss'u)^4
\]
y plegarlo con las formas del sistema relativamente completo módulo $s$. Se mostrará que, al introducir el módulo $(\varrho,t)$ como módulo superior,
\[
\begin{array}{ll}
\text{A)}&\text{el módulo }(ss'u)^4\text{ se vuelve reducible,}\vphantom{\dfrac11}\\
\text{B)}&\text{el sistema de }s\text{ consta de la única forma }s.
\end{array}
\]

A) La reducción del módulo $(ss'u)^4$ se obtiene inmediatamente del reductor $s_\nu^2$ encontrado mediante la fórmula (16.), \S\ 11. De
\[
 s_\nu^2=\frac43 i a_\nu^2+\frac45\theta_\varrho^2+\frac85(atu)^2-
 \frac{26}{15}(\nu\varrho x)^2+\frac16J\cdot u_x^2-
 \frac29i^2\cdot u_x^2
\]
se sigue, por polarización de $x$ hacia $\nu_1$:
\[
\tag{21}
\begin{aligned}
(ss'u)^4&=-\frac23J\cdot\nu+6s_\nu^2s_{\nu_1}^2\\
&=\frac{24}{5}\theta_\nu^2\theta_\varrho^2+
\frac{48}{5}a_\nu^2(atu)^2-
\frac{52}{5}H_\varrho^2+
\frac43i\cdot(asu)^4+\frac13J\cdot\nu-
\frac49i^2\cdot\nu,
\end{aligned}
\]
y con ello la reducción del módulo $(ss'u)^4$.\footnote{De la fórmula (21.) se sigue la relación, mencionada en la nota a la introducción, entre los tres invariantes de orden 9, aplicando la fórmula (10.)c.}

B) La reducción de
\[
Z=(ss'u)^2=\theta(s)
\]
y de ahí la reducción del sistema de $s$ se obtiene sustituyendo la forma $Z$ por la forma inferior $g_\nu^2$ según la fórmula (8.)b, \S\ 7, teniendo en cuenta la relación
\[
 s_\eta^2=s_\nu^2(\nu\nu_1x)^2-\frac13Z.
\]
La forma $g_\nu^2$ tiene, según la fórmula (13.), \S\ 11, el reductor $s_\nu^2$ como factor. Se obtiene, por las fórmulas (8.) y (16.):
\[
 g_\nu^2=-Z+\frac{14}{5}ia_\eta^2+\frac{14}{15}i(asu)^2
 \modu{(\varrho,t)},
\]
y, por las fórmulas (13.), (14.), (15.), (16.):
\[
\begin{aligned}
g_\nu^2&=2[s_\vartheta^2]_{\nu^2}-\frac43 i\A_\nu^2
      =2s_\vartheta^2s_\nu^2-\frac65[s_\vartheta^2(\theta su)^2]\widehat{\nu x}^{\,2}
        -\frac43i\A_\nu^2\\
&=\frac45i\cdot a_\eta^2-\frac25i\cdot(asu)^2+\frac13J\cdot\theta
\modu{(\varrho,t)}.
\end{aligned}
\]
De aquí se sigue
\[
\tag{23}
 Z=2i\cdot a_\eta^2+\frac43i\cdot(asu)^2-\frac13J\cdot\theta
 \modu{(\varrho,t)}.
\]

El sistema relativamente completo módulo $((ss'u)^4,\varrho,t)$ pasa por tanto a un sistema relativamente completo módulo $(\varrho,t)$; y de éste surge, por transvección sobre el sistema conocido de dos formas cuadráticas, el sistema absolutamente completo. Para formar el sistema relativamente completo módulo $(\varrho,t)$, puesto que según la fórmula (23.) el sistema de $s$ se reduce a la forma $s$, debemos transveccionar el sistema relativamente completo módulo $(s,\varrho,t)$ sobre $s$ y sobre potencias de $s$. Se verá que las potencias de $s$ sólo entran en plegamiento con el sistema I.

La relación anunciada es:
\[
\tag{22}
\begin{aligned}
(ass')^4&=\frac{24}{5}\theta_\nu^2\theta_\varrho^2a_\nu^2a_\varrho^2+
\frac{48}{5}a_\nu^2b_\nu^2(tab)^2-
\frac{52}{5}H_\varrho^2a_\nu^4+
\frac53J\cdot i-\frac49i^3,\\
(ass')^4&=\frac{24}{5}a_\varrho^2a_{\varrho'}^2-
\frac{12}{5}t_\varrho^2+\frac53J\cdot i-\frac49i^3.
\end{aligned}
\]

Como ``formas superiores'' definimos, entre las formas del mismo orden y del mismo grado, aquellas que han surgido de formas superiores del sistema relativamente completo módulo $(s,\varrho,t)$ por plegamiento con $s$, considerando el plegamiento con $s$ sólo como polarización de las formas originarias. Con ello queda determinada de modo unívoco la ordenación de las formas individuales (cf. \S\ 1, sucesiones de formas 2). Por ejemplo:
\[
(\theta_\eta(\theta Hu),s,u)^2\text{ es superior a }s_\eta((\theta Hu)^2,s,u),
\]
\[
(\theta_\eta(\theta Hu)^2,s,u)\text{ es superior a }(\theta_\eta(\theta Hu),s,u)^2.
\]

Dividimos el sistema resultante en los cuatro subsistemas:
\[
\begin{array}{ll}
\text{Sistema }s_I:&\text{surgido por plegamiento de }s\text{ y potencias de }s\text{ con el sistema I;}\vphantom{\dfrac11}\\
\text{Sistema }s_{II}:&\text{surgido por plegamiento de }s\text{ con el sistema II;}\vphantom{\dfrac11}\\
\text{Sistema }s_{III}:&\text{surgido por plegamiento de }s\text{ con las formas individuales}\\
&\text{(sin productos) del sistema III y con productos de una}\\
&\text{forma de cada uno de los sistemas I y II;}\vphantom{\dfrac11}\\
\text{Sistema }s_{IV}:&\text{surgido por plegamiento de }s\text{ con productos de una forma de cada}\\
&\text{uno de los sistemas I y III, II y III, III y III.}
\end{array}
\]
Observación: de ahora en adelante, todas las fórmulas deben entenderse módulo $(\varrho,t)$; desde la fórmula (27.) en adelante, módulo $(\varrho,t,i,J,\text{ formas superiores})$, sin que esto se indique expresamente cada vez.

Para la reducción reunimos las fórmulas obtenidas anteriormente:
\[
\tag{24}
\begin{aligned}
s_\nu^2&=\frac43ia_\nu^2+\frac16J\cdot u_x^2-\frac29i^2\cdot u_x^2,
& s_\nu^3&=\frac13J\cdot u_x,
& s_\nu^4&=J,\\[0.3em]
g&=2s_\vartheta^2-\frac43i\A,
& s_\vartheta^2(\theta su)&=0,\\[0.3em]
s_\vartheta^2(\theta su)^2&=\frac89i\cdot a_\nu^3-\frac4{27}i^2\cdot u_x^2,\\[0.3em]
Z&=2i\cdot a_\eta^2+\frac43i(asu)^2-\frac13J\cdot\theta,\\[0.3em]
g_\nu&=-s_\eta+u_x\left\{2ia_\eta^2-\frac23i(asu)^2+\frac23J\cdot\theta\right\},\\[0.3em]
g_\nu^2&=\frac45ia_\eta^2-\frac25i(asu)^2+\frac13J\cdot\theta,
&g_\nu^3&=0,
&g_\nu^4&=0.
\end{aligned}
\]

\noindent\emph{Consecuencias:}
\[
 s_\nu^3,
\quad s_\vartheta^2(\theta su),
\quad s_\vartheta^2s_\nu,
\quad s_\vartheta^2\theta_\nu
\]
son reductores.



\clearpage
\begin{center}
\textbf{\S\ 18.\quad Sistema $s_I$.}
\end{center}

\[
\text{Plegamiento de}\quad
\left\{\begin{array}{l}
s\text{ con } f;\ \theta;\ K,j,\A;\ f\cdot j,N;\ \A\cdot j,\\
s^3\text{ con }\theta,K.
\end{array}\right.
\]

\noindent\emph{Formas irreducibles.}

\noindent Orden 5:
\[
(asu),\qquad (asu)^2,\qquad (asu)^3,\qquad (asu)^4.
\]

\noindent Orden 6:
\[
\begin{array}{cccc}
(\theta su)&(\theta su)^2&(\theta su)^3&(\theta su)^4\\
s_\vartheta&s_\vartheta(\theta su)&s_\vartheta(\theta su)^2&s_\vartheta(\theta su)^3\\
\cdot&s_\vartheta^2&\cdot&\cdot
\end{array}
\]

\noindent Orden 7:
\[
\begin{array}{c@{\qquad}c@{\qquad}c@{\qquad}c}
\cdot&(Ksu)^2&(Ksu)^3&(Ksu)^4\\
\text{A)}\ s_k&s_k(Ksu)&s_k(Ksu)^2&s_k(Ksu)^3\\
\cdot&s_k^2&\cdot&\cdot
\end{array}
\qquad
\begin{array}{l}
\text{B)}\ s_j,\\[0.4em]
\text{C)}\ (\A su),\ (\A su)^2.
\end{array}
\]

\noindent Orden 8:
\[
\begin{array}{ll}
\text{A)}&s_j(asu),\ s_j(asu)^2,\ s_j(asu)^3,\\[0.4em]
\text{B)}&(Nsu)^2,\quad s_n,\quad s_n(Nsu).
\end{array}
\]

\noindent Orden 10:
\[
\text{A)}\ s_j(\A su),\qquad \text{B)}\ s_\vartheta(\theta s'u)^4.
\]

\noindent Orden 11:
\[
(Ks^2u)^6,\qquad s_k(Ks^2u)^5,
\qquad s_k(Ks^2u)^6.
\]

\noindent\emph{Reducciones:}

Las sucesiones de formas
\[
 s_\vartheta^2(\theta su),\quad s_k^2(Ksu),\quad s_j^2,\quad (\A su)^3,
 \quad (Nsu)^3
\]
tienen como factor el reductor $s_\vartheta^2(\theta su)$; $s_j^2$ y $(\A su)^3$ se obtienen retrocediendo de $j$ y de $\A$ a formas inferiores de la sucesión:
\[
\begin{aligned}
s_\vartheta^2(\theta su)(a\theta u)^3&=-\frac12s_j^2,\\
s_\vartheta^2(\theta su)(a\theta u)^2&=\frac13(\A su)^3,\\
s_\vartheta^2(\theta su)^2(a\theta u)^2&=\frac13(\A su)^4+\frac13s_j^2.
\end{aligned}
\]
Formas $s_\vartheta^2s_{\vartheta'}$ y $s_\vartheta^2s_{\vartheta'}^2$: reductores $s_\vartheta^2(\theta su)$ y $\theta_{\vartheta'}$:
\[
\begin{aligned}
s_\vartheta^2s_{\vartheta'}&=s_\vartheta^2\theta_{\vartheta'}-s_\vartheta^2(\theta s\widehat{\vartheta'}x),\\
s_\vartheta^2s_{\vartheta'}^2&=s_\vartheta^2s_{\vartheta'}\theta_{\vartheta'}-s_\vartheta^2s_{\vartheta'}(\theta s\widehat{\vartheta'}x).
\end{aligned}
\]
Sucesión de formas $s_j(\A su)^2$: reductores $\A_j$ y $(\A su)^3$.

\noindent\emph{Consecuencias:}

$s$ no entra en potencias en plegamiento con $f,j,\A,N$. Las formas $f,j,\A,N$ sólo entran en plegamiento en productos con formas contragredientes; sin embargo, en los productos con $f$ dichas formas pueden seguir siendo cuadráticas en $x$, y en los productos con $\A$ pueden seguir siendo lineales en $x$. Por tanto hay que considerar los siguientes productos de formas:
\[
 f\cdot(0,\lambda),\quad f\cdot(1,\lambda),\quad f\cdot(2,\lambda),\quad
 j\cdot(\mu,0),\quad N,\quad
 \A\cdot(0,\lambda),\quad \A\cdot(1,\lambda)\quad(\lambda>0),
\]
donde $(\mu,\lambda)$ designa una forma de grado $\mu$ en $x$ y de grado $\lambda$ en $u$.

Ni $\theta$ ni $K$ entran en potencias ni en productos con formas arbitrarias en plegamiento con $s$ o con $s^2$. En efecto, para productos con el determinante funcional $K$, $K\cdot\varphi_x$ ($\varphi\ne f$ o $\theta$), vale la relación entre formas descomponibles
\[
 K\cdot\varphi_x=(a\theta u)\varphi_x=f\cdot(\varphi\theta u)-\theta\cdot(\varphi au).
\]
Con ello queda demostrado el esquema anterior para el plegamiento del sistema $s_I$.

\begin{center}
\textbf{\S\ 19.\quad Sistema $s_{II}$.}
\end{center}

Plegamiento de $s$ con $\nu;\ H;\ L;\ H^2$.

\noindent\emph{Formas irreducibles:}
\[
\begin{array}{c@{\qquad}c@{\qquad}c@{\qquad}c}
\text{Orden 6} & \text{Orden 8} & \text{Orden 10} & \text{Orden 12}\\[0.3em]
s_\nu & (Hsu),\ (Hsu)^2 & (Lsu)^2,\ (Lsu)^3 & (H^2su)^3\\
\cdot & s_\eta & s_l & \cdot\\
\cdot & s_\nu^4=J & \cdot & \cdot
\end{array}
\]

\noindent\emph{Reducciones:}

Sucesiones de formas $s_\eta(Hsu)$ y $s_l(Lsu)$: reductor $s_\nu^2$.

Sucesión de formas $s_\sigma$: reductor $s_\nu^2$ procedente de $H$, $s_\nu^2=\dfrac23s_\sigma$.

$(H^2su)^4$ se descompone por la fórmula (7.)a (cf. \S\ 11 A). De ahí se sigue también la descomposición de $(H\cdot L,s,u)^4$.

\noindent\emph{Consecuencias:}

$s^2$ no entra en plegamiento con el sistema II. La forma $\nu$ entra sólo en productos con formas contragredientes. La forma $H$, considerada como covariante, entra en plegamiento en productos con formas $(1,\lambda)$, $(2,\lambda)$, $(3,\lambda)$, con $\lambda$ arbitrario. Las formas $\sigma$ y $(\nu\sigma x)$ no entran en absoluto, y $L$, como determinante funcional, no entra en productos en el plegamiento; las transvecciones completas sobre covariantes del sistema III se vuelven reducibles. De los productos de los sistemas I y II quedan:
\[
 f\cdot\nu,\qquad \A\cdot\nu.
\]

\begin{center}
\textbf{\S\ 20.\quad Formas de séptimo orden (sistema $s_{III}$).}
\end{center}

Plegamiento de $s$ con $f\cdot\nu$, $a_\nu$, $a_\nu^2$.

\noindent\emph{Formas irreducibles:}
\[
\begin{gathered}
s_\nu(asu),\quad a_\nu(asu),\quad s_\nu(asu)^2,\quad a_\nu(asu)^2,
\quad s_\nu(asu)^3,\quad a_\nu(asu)^3,\\
a_\nu s_\nu,
\quad a_\nu s_\nu(asu),
\quad a_\nu^2(asu),
\quad a_\nu^2(asu)^2.
\end{gathered}
\]

\noindent\emph{Reducciones:}

No mencionaremos las reducciones evidentes que surgen por plegamiento directo con los reductores $s_\nu^2$ y $s_\eta(Hsu)$, ni tampoco la descomposición de transvecciones con determinantes funcionales.

Para las formas $a_\nu s_\nu(asu)^2$ y $a_\nu^2s_\nu(asu)^2$, véase la fórmula (19.), \S\ 12; además,
\[
 a_\nu s_\nu(asu)^3=a_\nu(a\widehat{\nu}_1\nu_2u)^3(\nu_1\nu_2\nu)=0.
\]
Sucesión de formas $a_\nu^2s_\nu$: reductor $s_\vartheta^2(\theta su)$ obtenido retrocediendo de $a_\nu^2$ a formas inferiores, es decir, por doble reducción de las expresiones
\[
 s_\vartheta^2(a\theta s)(a\theta u),
 \qquad s_\vartheta^2(a\theta s)(a\theta u)^2
\]
según el planteamiento:
\[
\begin{aligned}
s_\vartheta^2(a\theta s)(a\theta u)
&=[(a\theta u)^2]s^3+\frac15[a_\vartheta(a\theta u)^2]s^2+\frac1{60}[a_\vartheta^2(a\theta u)^2]s\\
&=\frac12a_\nu^2s_\nu-\frac23a_\nu^2s_\nu^2,\\
s_\vartheta^2(a\theta s)(a\theta u)^2
&=\frac45[a_\vartheta(a\theta u)^2]_{us}s^2+\frac{23}{180}[a_\vartheta^2(a\theta u)^2]_{us}s\\
&=-\frac12a_\nu^2s_\nu(asu).
\end{aligned}
\]
Aquí vale $a_\nu^2b_\nu(abu)=0$.

Se obtienen las fórmulas:
\[
\tag{25}
\begin{array}{rcl@{\qquad}rcl}
a_\nu^2s_\nu&=&\dfrac49i\,a_\nu^2b_\nu,
& a_\nu s_\nu(asu)^2&=&\dfrac13i\theta_\nu^2-\dfrac1{18}iH,\\[0.7em]
a_\nu s_\nu(asu)^3&=&0,
& a_\nu^2s_\nu(asu)&=&0,\\[0.7em]
a_\nu^2s_\nu(asu)^2&=&0.
\end{array}
\]

\noindent\emph{Consecuencias:}

1) $a_\nu s_\nu(asu)^2$ y $a_\nu^2s_\nu$ son reductores.

2) Se obtienen los valores de las formas $(\A su)^3$, $(\A su)^4$, ya reconocidas como reducibles en \S\ 19; igualmente para la correspondiente sucesión de formas $(agu)^2$, que al plegarse con $\nu$ da lugar a una doble reducción mediante el reductor $g_\nu$:
\[
\tag{26}
\begin{aligned}
\text{a)}\quad (\A su)^3&=-\frac65a_\nu^2(asu)+3a_\nu s_\nu(asu)+2i\theta_\nu(\theta\nu x),\\
\text{b)}\quad (\A su)^4&=-\frac25a_\nu^2(asu)^2+\frac43i\theta_\nu^2+\frac19iH.
\end{aligned}
\]
De las fórmulas (24.) y (26.), módulo $(i,J)$ (cf. p. 71), se sigue
\[
\tag{27}
\begin{aligned}
\text{a)}\quad (agu)^2&=\frac23(\A su)^2-\frac43a_\nu s_\nu+\frac35s\cdot a_\nu^2,\\
\text{b)}\quad (agu)^3&=2a_\nu s_\nu(asu)-a_\nu^2(asu)^2,\\
\text{c)}\quad (agu)^4&=-a_\nu^2(asu)^2.
\end{aligned}
\]

\begin{center}
\textbf{\S\ 21.\quad Formas de octavo orden (sistema $s_{III}$).}
\end{center}

Plegamiento de $s$ con las formas de cuarto orden (sistema III) (\S\ 9).

\noindent\emph{Formas irreducibles:}
\[
\begin{array}{lll}
(\vartheta\nu x)s_\nu,& ((\vartheta\nu x),s,u)^2,\\[-0.1em]
((\vartheta\nu x)^2,s,u),& ((\vartheta\nu x)^2,s,u)^2,\ \theta s_\nu,\\[-0.1em]
(\vartheta\nu x)^2s_\nu,& ((\vartheta\nu x)^2,s,u)s_\nu,
\end{array}
\]
\[
\begin{array}{lll}
((\vartheta\nu x),s,u)^3,\ (\theta_\nu su)^2,&
((\vartheta\nu x),s,u)^4,\ (\theta_\nu su)^3,\\[-0.1em]
((\vartheta\nu x)^2,s,u)^3,\ (\theta_\nu(\vartheta\nu x),s,u)^2,\ (\theta_\nu^2su),&
((\vartheta\nu x),s,u)^3,\ (\theta_\nu^2su)^2,\ (\theta_\nu(\vartheta\nu x),s,u)^4.
\end{array}
\]

\noindent\emph{Reducciones:}

$((\vartheta\nu x),s,u)s_\nu$ se descompone como transvección sobre determinantes funcionales, por \S\ 3.

Sucesión de formas $((\vartheta\nu x),s,u)^2s_\nu$: reductores $s_\nu^2$ y $s_\vartheta^2\theta_\nu$:
\[
(\widehat\vartheta\nu su)s_\nu(\theta su)=s_\nu s_\vartheta(\theta su)=-\frac12(\widehat\vartheta\nu su)^2(\theta su),
\]
\[
\theta_\nu(\widehat\vartheta\nu su)s_\nu=\theta_\nu s_\nu s_\vartheta=-\frac12\theta_\nu(\widehat\vartheta\nu su)^2.
\]
Sucesión de formas $\theta_\nu s_\nu(\theta su)$: reductor $a_\nu s_\nu(asu)^2$ por doble reducción de las expresiones
\[
a_\nu s_\nu(asu)^2(abu)\quad\text{y}\quad a_\nu s_\nu(asu)^2(abu)(sbu);
\]
$\theta_\nu s_\nu(\theta su)^3$ por doble reducción de $(agu)^3(abu)(bgu)^3$.

Sucesión de formas $\theta_\nu(\vartheta\nu x)s_\nu$: reductor $a_\nu^2s_\nu$ a partir de $\theta_\nu(\vartheta\nu x)s_\nu=a_\nu^2(abu)s_\nu$.

Sucesión de formas $(\theta_\nu(\vartheta\nu x)^2,s,u)$: doble reducción de las formas de la sucesión $\theta_\nu(\widehat\vartheta\nu su)s_\nu$, que al mismo tiempo pertenecen a las sucesiones $((\vartheta\nu x)su)^2s_\nu$ y $\theta_\nu(\vartheta\nu x)s_\nu$.

La forma $(\theta_\nu(\vartheta\nu x),s,u)$ se descompone como transvección simple sobre el determinante funcional $\theta_\nu(\vartheta\nu x)=a_\nu^2(abu)$.

La forma $((\vartheta\nu x)^2,s,u)^4$: doble reducción de la expresión $(a\A u)^2(\A su)^4$, o también de $(\theta gu)^4$, a partir de las fórmulas (27.) y de manera análoga a la reducción de $(j\eta x)^4$ (\S\ 12 C)).

La doble reducción da las fórmulas:
\[
\tag{28}
\begin{aligned}
0&=\theta_\nu s_\nu(\theta su)-\frac16(\widehat\vartheta\nu su)^2(\theta su)-\frac1{12}(Hsu),\\
0&=\theta_\nu s(\theta su)^2-\frac14(\widehat\vartheta\nu su)^2(\theta su)^2-\frac1{24}(Hsu)^2,\\
0&=(\widehat\vartheta\nu su)^2(\theta su)^2+2\theta_\nu(\widehat\vartheta\nu su)(\theta su)^2
  +3\theta_\nu^2(\theta su)^2-\frac13H(su)^2,\\
0&=\theta_\nu s_\nu(\theta su)^3+\frac12\theta_\nu(\widehat\vartheta\nu su)(\theta su)^3,\\
0&=\theta_\nu s_\nu s_\vartheta-\frac14s_\eta,
\quad 0=\theta_\nu s_\nu s_\vartheta(\theta su),
\quad 0=\theta_\nu s_\nu s_\vartheta(\theta su)^2.
\end{aligned}
\]
El último grupo corresponde a $((\theta_\nu(\vartheta\nu x)^2,s,u)^2,\ldots)$.

\noindent\emph{Consecuencias:}

1) $\theta(\vartheta\nu x)s_\nu$, $(\widehat\vartheta\nu su)(\theta su)s_\nu$, $\theta_\nu(\widehat\vartheta\nu su)^2$, $(\widehat\vartheta\nu su)^2(\theta su)^2$ son reductores.

2) Las formas de cuarto orden $(5,\lambda)$, $(6,\lambda)$, etc., no entran en plegamiento con $s^2$.

3) Para la sucesión de formas $(\theta gu)^2$, las fórmulas (27.), teniendo en cuenta las reducciones de este parágrafo, dan las fórmulas siguientes:
\[
\tag{29}
\begin{aligned}
\text{a)}\quad (\theta gu)^2&=\frac43\theta_\nu s_\nu+\frac23s_\nu s_\vartheta-\frac13s\theta_\nu^2-\frac19sH-\frac13f\,a_\nu^2(asu)^2+\frac13a_\nu^2(asu)^2,\\
\text{b)}\quad (\theta gu)^3&=-\frac13(\widehat\vartheta\nu su)^2(\theta su)-\theta_\nu(\widehat\vartheta\nu su)(\theta su)-\theta_\nu^2(\theta su),\\
\text{c)}\quad (\theta gu)^4&=2\theta_\nu^2(\theta su)^2-\frac13(Hsu)^2,\\
g\vartheta(\theta gu)&=(\widehat\vartheta\nu su)(\vartheta\nu x)s_\nu-\theta_\nu(\vartheta\nu x)(\widehat\vartheta\nu su),\\
g\vartheta(\theta gu)^2&=\frac23s\theta_\nu^3,
\quad g\vartheta(\theta gu)^3=\frac13\theta_\nu^3(\theta su),
\quad g\vartheta(\theta gu)^4=0.
\end{aligned}
\]

\begin{center}
\textbf{\S\ 22.\quad Formas de noveno orden (sistema $s_{III}$).}
\end{center}

Plegamiento de $s$ con las formas de quinto orden (sistema III) y con el producto $\A\cdot\nu$ (\S\ 10).

\noindent\emph{Formas irreducibles:}
\[
\begin{array}{ll}
\text{A)}& (k\nu x)^2s_\nu,
((k\nu x)^3,s,u),
(k\nu x)^3s_\nu,
((k\nu x)^2,s,u)^2,
K_\nu s_\nu,\\
& ((k\nu x)^3,s,u)^2,
((k\nu x)^2,s,u)s_\nu,
(K_\nu(k\nu x)^3,s,u),\\
& (K_\nu su)^2,\ (K_\nu su)^3,\ (K_\nu su)^4,
((k\nu x)^2,s,u)^3,
(K_\nu^2su)^4,\\[0.3em]
\text{B)}& (j\nu x)s_\nu,
((j\nu x)^2,s,u),
(j\nu x)^3,s,u,
((j\nu x)^2,s,u)^2,\\[0.3em]
\text{C)}& s_\nu(\A su),\quad (\A_\nu su),\\[0.3em]
\text{D)}& (aHu)s_\eta,
(a_\eta su),
((aHu),s,u)^2,
((aHu)^2,s,u),\\
& (aHu)^2s_\eta,
(a_\eta su)^2,
(a_\eta^2su),
((aHu),s,u)^3,
((aHu)^2,s,u)^2,\\
& ((aHu),s,u)^4,
(a_\eta su)^3,
(a_\eta(aHu),s,u)^2,
(a_\eta(aHu)^2,s,u),
(a_\eta(aHu),s,u)^3.
\end{array}
\]

\noindent\emph{Reducciones:}

A) Las reducciones se siguen directamente de las de \S\ 21 por plegamiento simple. La sucesión de formas $K_\nu s_\nu(Ksu)^2$ tiene como factores los reductores $a_\nu s_\nu(asu)^2$ y $\theta_\nu s_\nu(\theta su)^2$. De ahí que las formas superiores $K_\nu^2(Ksu)^2$, $K_\nu^2(Ksu)^3$ se descompongan según el planteamiento:
\[
0=a_\nu s_\nu(asu)^2(a\theta u)=K_\nu s_\nu(Ksu)^2-\theta_\nu s_\nu(\theta su)^2(a\theta u).
\]

B) Sucesión de formas $(j\nu x)^2s_\nu$: reductor $a_\nu^2s_\nu$ a partir de
\[
(j\nu x)^2s_\nu=3a_\nu^2s_\nu(a\theta u)^2.
\]
$((j\nu x)^3,s,u)^2$: reductor $a_\nu^2s_\nu$ a partir de
\[
((j\nu x)^3,s,u)^2=-6a_\nu^2s_\nu(a\theta u)^2(\theta su).
\]

C) Formas $\A_\nu(\A su)^2$ y $s_\nu(\A su)^2$: reductor $g_\nu$ a partir de la fórmula (27.)a.

Forma $\A_\nu s_\nu$:
1) reductor $a_\nu^2s_\nu$ a partir de $a_g a_\nu^2s_\nu=\frac23\A_\nu s_\nu$;
2) se descompone por la fórmula (27.)a mediante doble reducción de la expresión $(\A s\widehat\nu x)^2$.

De ahí se descompone la forma superior $a_\eta s_\eta$.

D) Sucesión de formas $(aHu)^2s_\eta(asu)$: reductor $a_\nu s_\nu(asu)^2$ por doble reducción de la sucesión $[a_\nu s_\nu(asu)^2]_{\nu_1x^2}$; reducción de $(aHu)^2s_\eta(asu)^2$ a partir de $a_\nu s_\nu(asu)^2$ y $a_\nu s_\nu(asu)^3$. De ahí la reducción de $(a_\eta,s,u)^4$.

Sucesión de formas $a_\eta(aHu)s_\eta$: reductores $a_\nu^2s_\nu$, $a_\eta^2s_\eta$ por doble reducción de la expresión $(\A s\widehat\nu x)^3$, puesto que $a_\eta^2s_\eta$ no surge por plegamiento del término reducible $a_\nu^2s_\nu(\nu\nu_1x)$ de la forma $a_\eta(aHu)s_\eta$.

Sucesión de formas $(a_\eta^2su)^2$: reductor $a_\nu^2s_\nu$ a partir de
\[
a_\nu^2s_\nu s_{\nu_1}=-\frac12(a_\eta^2(Hsu))^2.
\]
La forma $(a_\eta(aHu),s,u)$ se descompone como transvección sobre un determinante funcional.

Las fórmulas de reducción obtenidas son:
\[
\tag{30}
\begin{aligned}
\text{a)}\quad 0&=(aHu)^2(asu)s_\eta-a_\eta(asu)(Hsu)^2-a_\eta(aHu)(asu)(Hsu),\\
\text{b)}\quad 0&=(aHu)^2(asu)^2s_\eta-a_\eta(aHu)(asu)^2(Hsu),\\
\text{c)}\quad 0&=a_\eta(asu)^2(Hsu)^2,\\
0&=a_\eta s_\eta+\frac14a_\nu^2g-\frac14s\cdot a_\eta^2,\\
0&=a_\eta(aHu)s_\eta-\frac12a_\eta^2(Hsu),
\quad 0=a_\eta^2(Hsu)^2,\ldots,
\quad 0=a_\eta^2s_\eta.
\end{aligned}
\]

\noindent\emph{Consecuencias:}

1) $(j\nu x)^2s_\nu$, $\A_\nu s_\nu$, $\A_\nu(\A su)^2$, y por consiguiente $N_\nu(Nsu)^2$, $(aHu)^2(asu)s_\eta$, $a_\eta(aHu)s_\eta$, $a_\eta^2(Hsu)^2$, son reductores; $a_\eta s_\eta$ es una forma descomponible que pasa a formas descomponibles bajo todos los plegamientos simples y dobles.

2) Las formas de quinto orden $(5,\lambda)$, $(6,\lambda)$, etc., no entran en plegamiento con $s^2$.

\begin{center}\textbf{\S\ 23.\quad Formas de décimo orden (sistema $s_{III}$).}\end{center}
Plegamiento de $s$ con las formas de sexto orden (sistema III) (\S\ 11).

\noindent\emph{Formas irreducibles:}
\[
\begin{array}{rl}
\text{A)}& (\vartheta^2\nu x)^3s_\nu,
((\vartheta^2\nu x)^3,s,u),
((\vartheta^2\nu x)^3,s,u)^2,
(\vartheta^2\nu x)^3s_\nu s_\vartheta,
\theta_\nu(\vartheta^2\nu x)^3s_\vartheta,\\[0.3em]
\text{B)}& a_\nu(j\nu x)s_\nu,
(a_\nu(j\nu x),s,u)^2,
(a_\nu(j\nu x),s,u)^3,
(a_\nu(j\nu x),s,u)^4,
(a_\nu^2(j\nu x),s,u)^2,
(a_\nu^2(j\nu x),s,u)^3,\\[0.3em]
\text{C)}& ((\vartheta\eta x)^2,s,u),
((\vartheta\eta x),s,u)^2,
(\theta Hu)s_\eta,
(\theta_\eta su),
((\theta Hu),s,u)^2,
((\theta Hu)^2,s,u),\\
& ((\theta\eta x),s,u)^3,
(\theta_\eta su)^2,
(\theta_\eta^2su),
((\theta Hu),s,u)^3,
((\theta Hu)^2,s,u)^2,
(H_\vartheta(\theta Hu),s,u)^2,\\
& (\theta_\eta(\theta Hu),s,u)^2,
(\theta_\eta(\theta Hu)^2,s,u),
((\theta Hu),s,u)^4,
(\theta_\eta su)^4,
(\theta_\eta(\theta Hu),s,u)^3,
(\theta_\eta^2(\theta Hu),s,u)^2.
\end{array}
\]
\noindent\emph{Reducciones:}
A) Las formas $((\vartheta^2\nu x)^3,s,u)^2$, $((\vartheta^2\nu x)^3,s,u)^3$ se descomponen:
\[
(\vartheta\nu x)(\widehat{\vartheta\nu su})(\widehat{\vartheta\nu su})=(\vartheta\nu x)s_\vartheta s_\nu-(\vartheta\nu x)s_\nu s_\vartheta=-2\theta(\vartheta\nu x)s_\nu s_\vartheta+(\vartheta\nu x)s_\vartheta^2.
\]
Las formas restantes tienen reductores como factores o son plegamientos simples con formas descompuestas.

B) Reductor $a_\nu^2s_\nu$ y $(\widehat{j\nu su})s_\nu$.

C) 1. La sucesión de formas $(\theta Hu)^2s_\eta$: a) reductor $(aHu)^2(asu)s_\eta$; b) doble reducción de las sucesiones $(\theta gu)^3g_\nu$ y $g_\nu(agu)^3(bgu)^3(abu)$. De ahí la reducción de las formas superiores $(\theta,s,u)^3$, $(H_\vartheta(\theta Hu),s,u)^3$ y $(a_\nu b_{\nu_1}^2,s,u)^4$ [esta última forma sustituye a $(H_\vartheta su)^4$].

2. $(\vartheta\eta x,s,u)^4$: reductor $(a_\eta su)^4$.

3. $((\vartheta\eta x)^2,s,u)^3$: reductor $s_\vartheta^2s_\nu$ a partir de $(\widehat{\vartheta\eta su})^2(Hsu)$. Las formas $(\vartheta\eta x)s_\eta$, $(\vartheta\eta x)^2s_\eta$, $((\vartheta\eta x)^2,s,u)^2$, $\theta_\eta s_\eta$ se descomponen por plegamiento con la forma descompuesta $a_\eta s_\eta$.

4. Sucesiones de formas $H_\vartheta(\theta Hu)s_\eta$, $\theta_\eta(\theta Hu)s_\eta$, $H_\vartheta(\vartheta\eta x)s_\eta$, $\theta_\eta(\vartheta\eta x)s_\eta$: reductores $a_\eta(aHu)s_\eta$ y $a_\eta^2s_\eta$.

La sucesión de formas $(\theta_\eta(\vartheta\eta x),s,u)^2$: reductor $a_\eta^2(Hsu)^2$ [excepto $(\theta_\eta^2(\theta Hu),s,u)^2$, que surge del término $a_\eta^2(\theta su)(Hsu)$ por plegamiento].

5. $(H_\vartheta(\vartheta\eta x),s,u)$, $(H_\vartheta(\vartheta\eta x),s,u)^2$: doble reducción de $a_\eta(aHu)s_\eta(abu)$ y $g_\vartheta(\theta gu)g_\nu$.

La sucesión de formas $(H_\vartheta(\vartheta\eta x),s,u)^3$: reductor $((\vartheta\nu x)^2,s,u)^2s_\nu$.

6. $(H_\vartheta(\theta Hu),s,u)$ se descompone por plegamiento simple con la forma descompuesta $(\theta_\nu(\vartheta\nu x),s,u)$ según \S\ 3. 2), c), es decir, por doble reducción de la forma descompuesta inferior $(H_\vartheta su)^2$:\footnote{El signo $\equiv$ significa que se han omitido productos de formas de grado inferior.}
\begin{align*}
0&=\theta_\nu(\vartheta\nu_1)(\theta su)+\frac23\theta_\nu^2(\vartheta\nu_1x)(\theta su)+\theta_\nu(\vartheta\nu x)(\theta su)s_{\nu_1}\\
&\equiv -\frac12H_\vartheta(\theta Hu)(\theta su)+\theta_\eta(\theta su)(Hsu)+\frac12H_\vartheta(\theta su)(Hsu)\\
&\equiv -\frac12H_\vartheta(\theta Hu)(\theta su)+\theta_\eta(\theta su)(Hsu)+\frac12[H_\vartheta]_{sx}^2-\frac12\frac35H_\vartheta(\theta Hu)(\theta su)\\
&\equiv -\frac35H_\vartheta(\theta Hu)(\theta su).
\end{align*}
Por doble reducción, según 1. y 5., se obtienen las fórmulas siguientes:
\begin{equation}
\tag{31}
\begin{aligned}
0&=\theta_\eta(\theta su)(Hsu)^2+\frac14H_\vartheta(\theta Hu)(\theta su)^2\\
&\quad-\frac34\theta_\eta(\theta Hu)(\theta su)(Hsu)-\frac54\theta_\eta(\theta Hu)^2(\theta su)^3\\
&\quad-(asu)^3b_\eta(\theta Hu)^2-a_\nu(asu)^3b_\nu^2,\\
0&=H_\vartheta(\theta Hu)(\theta su)^3-7\theta_\eta(\theta Hu)(\theta su)^2(Hsu),\\
0&=a_\nu(asu)^2b_{\nu_1}^{2}(bsu)^2-\theta_\eta(\theta su)^2(Hsu)^2\\
&\quad+6\theta_\eta(\theta Hu)(\theta su)^2(Hsu),\\
0&\equiv (H_\vartheta(\vartheta\eta x),s,u),
&0&\equiv H_\vartheta(\widehat{\vartheta\eta su})(Hsu)-4\theta_\eta^2(Hsu).
\end{aligned}
\end{equation}
\noindent\emph{Consecuencias:}
\begin{enumerate}
\item $(\theta Hu)^2s_\eta$, $H_\vartheta(\theta Hu)s_\eta$, $\theta_\eta(\theta Hu)s_\eta$, $H_\vartheta(\vartheta\eta x)s_\eta$, $\theta_\eta(\vartheta\eta x)s_\eta$, $(H_\vartheta(\vartheta\eta x),s,u)$, $(\theta_\eta(\vartheta\eta x),s,u)^2$ son reductores.
\item $s^2$ no entra en plegamiento con las formas $(5,\lambda)$, etc., de orden 6.
\end{enumerate}

\begin{center}\textbf{\S\ 24.\quad Formas de undécimo orden (sistema $s_{III}$).}\end{center}
Plegamiento de $s$ con las formas de séptimo orden (sistema III) (\S\ 12).

\noindent\emph{Formas irreducibles:}
\[
\begin{array}{rl}
\text{A)}& ((k\eta x)^3,s,u),\quad (K_\eta(k\eta x)^3,s,u),\quad (K_\eta su)^2,
\quad ((KHu)^2su)^2,\\
& ((KHu)^2,s,u)^3,
\quad ((KHu)^2,s,u)^4,
\quad (K_\eta(KHu)^2,s,u)^3,\\[0.3em]
\text{B)}& ((j\eta x),s,u)^2,
\quad (H_jsu),
\quad ((j\eta x),s,u)^3,\\[0.3em]
\text{C)}& (\D_\eta su),\\[0.3em]
\text{D)}& (aLu)^2s_\eta,
\quad (a_\eta su)^2,
\quad ((aLu)^2,s,u)^2,
\quad ((aLu)^3,s,u)^2,
\quad ((aLu)^3,s,u)^3,
\quad (a_l(aLu)^2,s,u)^2.
\end{array}
\]
\noindent\emph{Reducciones:}

A) Reducción o descomposición por plegamiento simple con las formas correspondientes en los símbolos $\theta-H$.

$(K_\eta(KHu)^2,s,u)$ y $(K_\eta(KHu)^2,s,u)^2$: descomposición según \S\ 3. 2), a), por plegamiento con $K_\nu^2(Ksu)$, $K_\nu^2(Ksu)^2$.

$(K_\eta(KHu),s,u)^4\equiv(a_\nu^2\theta_{\nu_1},s,u)^4$: descomposición según \S\ 3. 2), b), a partir de la forma descompuesta $K_\nu^2(Ksu)^3$.

B) Sucesión de formas $(j\eta x)s_\eta$: reductor $(j\nu x)^2s_\nu$.

Forma $H_j(\widehat{j\eta su})(Hsu)$: reductor $(j\nu x)(\widehat{j\nu su})^2$.

Forma $H_j(j\eta x)(Hsu)$: se descompone reduciendo la forma descompuesta $((j\eta x)^2,s,u)^2$ mediante el reductor $(j\nu x)^2s_\nu$ (fórmula (18.)a):
\[
0=-2(j\nu x)^2s_\nu s_{\nu_1}=[(j\eta x)^2]_{su^2}+H_j(j\eta x)(Hsu).
\]
C) Reductores $\D_\nu s_\nu$ y $\D_\nu(\D su)^2$.

D) Reducción y descomposición por plegamiento simple con las formas correspondientes en los símbolos $f-H$.

E) De (30.)c se obtiene la reducción de la sucesión de formas descompuesta $(a_\sigma su)^2$:
\begin{align*}
0&=a_\eta(Hsu)^2(as\widehat{\nu x})^2
=a_\eta(Hsu)^2(a\widehat{\nu\eta}u)^2+2a_\eta(a\widehat{\nu\eta}u)(\widehat{\eta\sigma u})(Hsu)^2
=\frac13a_\sigma(asu)^2,\\
0&=a_\eta a_\nu(Hsu)^2(as\widehat{\nu x})(asu)
=a_\eta(a\widehat{\nu\eta}u)^2(Hsu)^2(asu)+a_\eta(a\widehat{\nu\eta}u)(\widehat{\eta\sigma u})(Hsu)^2(asu)
=\frac13a_\sigma(asu)^3.
\end{align*}
\noindent\emph{Consecuencia:} $s^2$ no entra en plegamiento con las formas de orden 7.

\begin{center}\textbf{\S\ 25.\quad Formas de los órdenes 12, 13, 14 y 15 (sistema $s_{III}$).}\end{center}
Plegamiento de $s$ con las formas de los órdenes 8, 9, 10, 11 (sistema III) (\S\S\ 13--16).

\noindent\emph{Formas irreducibles:}

\noindent Orden 12.
\[
\begin{array}{rl}
\text{A)}& ((\vartheta^2\eta x)^4,s,u),\quad \theta_\eta(\vartheta^2\eta x)^3s_\vartheta,\\
\text{B)}& ((aHu)H_j,s,u)^2,
\quad ((aH\cdot u)H_j,s,u)^3,\\
\text{C)}& (\theta_\nu su)^2,
\quad ((\theta Lu)^2,s,u)^2,
\quad ((\theta Lu)^3,s,u),
\quad ((\theta Lu)^2,s,u)^3,
\quad (\theta_l(\theta Lu)^2,s,u)^2.
\end{array}
\]
\noindent\emph{Reducciones:} Ya no se mencionarán las reducciones que surgen por plegamiento directo con reductores o con formas descompuestas.

B) Descomposición de $((aHu)H_j,s,u)$ y $(a_\eta(aHu)H_j,s,u)$ como transvecciones sobre determinantes funcionales.

Descomposición de $(a_\eta(aHu)H_j,s,u)^2$: doble reducción de la forma descompuesta $[\theta_\nu\theta_{\nu_1}^3]_{su^3}$ (\S\ 13. C):
\begin{align*}
0&\equiv[\theta_{\nu_1}^3\theta_\nu]_{su^3}
\equiv-\frac34\theta_\nu(\theta su)^2(\theta\theta'u)\theta_{\nu_1}^3\\
&\equiv-\frac98(\theta\theta'\eta su)(\theta\theta'u)(\theta Hu)(\theta'Hu)\theta_\nu'(\theta su)
\equiv-\frac38a_\eta(aHu)H_j(asu)^2.
\end{align*}

C) De las fórmulas (31.) se sigue la reducción de las formas descompuestas
\[
 [L_\vartheta(\theta Lu)]_{su^3}\quad\hbox{y}\quad [\theta_l(\theta Lu)]_{s^2u^3},
\]
es decir, de los productos $[\theta_\nu^2H]_{su^3}$ y $[a_\nu^2(bHu)^2]_{su^3}$:
\begin{equation}
\tag{32}
\begin{aligned}
0&\equiv\theta_\nu^2(\theta su)^2(Hsu),\\
0&\equiv [L_\vartheta(\theta Lu)]_{su^4}-7[\theta_l(\theta Lu)]_{su^4},\\
0&\equiv a_\nu^2(asu)^2(bHu)^2(bsu)\\
&\quad+6\theta_l(\theta Lu)^2(\theta su)(Lsu),\\
0&\equiv\theta_\nu^3(\theta su)(Hsu)^2\\
&\quad\hbox{de }0\equiv\theta_l^2(\theta Lu)(\theta su)(Lsu)^2\\
&\quad(\hbox{reductor }a_\eta^2(Hsu)^2).
\end{aligned}
\end{equation}

\noindent Orden 13.
\[
\text{A)}\ ((KLu)^3,s,u)^2,\ ((KLu)^3,s,u)^3,\qquad
\text{B)}\ (L_jsu)^2,
\qquad
\text{C)}\ ((aH^2u)^3,s,u)^2.
\]
\noindent Orden 14.
\[
\text{A)}\ ((aLu)^2L_j,s,u)^2,
\qquad
\text{B)}\ ((\theta H^2u)^3,s,u)^2.
\]
\noindent Orden 15.
\[
 ((KH^2u)^4,s,u)^2.
\]
\noindent\emph{Reducciones:} Descomposición de $(K_l(KLu)^3,s,u)^2$, $(H_\vartheta(\theta H^2u)^3,s,u)$, $((K,H\cdot L,u)^5,s,u)$ según \S\ 3. 2), c); es decir, considerando simultáneamente las formas descompuestas $(K_l(KLu)^2,s,u)^3$, $(H_\vartheta(\theta H'u)^2,s,u)^2$, $((K,H\cdot L,u)^4,s,u)^2$.

\noindent\emph{Consecuencia:} $s^2$ no entra en plegamiento con formas individuales del sistema III.

\begin{center}\textbf{\S\ 26.\quad Sistema $s_{IV}$.}\end{center}
\noindent 1) \emph{Plegamiento de $s$ con los sistemas I--III.}

\noindent\emph{Reducciones:} Según las reducciones de los últimos parágrafos ($(a_\nu^2s_\nu(aHu)^2(asu)s_\eta$, etc.), las formas $(0,\lambda)$, $(1,\lambda)$, $(2,\lambda)$ no admiten simultáneamente plegamientos I y III; por tanto, según \S\ 18, las formas $f,\D,N$ no entran en productos en el plegamiento que forma el sistema $s_{IV}$.

La forma $j$ no entra en plegamiento, pues, por las mismas reducciones, los plegamientos contragredientes a $j$,
\[
(K_\nu(k\nu x)^3,s,u),\qquad ((\vartheta^2\eta x)^4,s,u)\quad\hbox{etc.},
\]
corresponden a los plegamientos cogredientes a $j$:
\[
K_\nu(k\nu x)^2s_k,
\qquad (\vartheta^2\eta x)^3s_\vartheta\quad\hbox{etc.}
\]

\noindent 2) \emph{Plegamiento de $s$ con los sistemas II--III}, es decir, de $s$ con $H\cdot(1,\lambda)$, $H\cdot(2,\lambda)$, $H\cdot(3,\lambda)$.

\noindent\emph{Formas irreducibles:}
\[
\begin{array}{rl}
\text{A)}&(a_\nu H,s,u)^4,
((aHu)^2H,s,u)^3,
((aHu)^2H,s,u)^4,\\
&((aLu)^2H,s,u)^4,
((aLu)^3H,s,u)^3,
((aH^2u)^3H,s,u)^4,\\[0.3em]
\text{B)}&(a_\nu^2H,s,u)^3,
(a_\eta(aHu)^2H,s,u)^3,
(a_l(aLu)^2H,s,u)^4,\\[0.3em]
\text{C)}&(\theta_\nu H,s,u)^4,
((\theta Hu)^2H,s,u)^3,
((\theta Hu)^2H,s,u)^4,\\
&((\theta Lu)^2H,s,u)^4,
((\theta Lu)^3H,s,u)^3,
((\theta H^2u)^3H,s,u)^4,\\[0.3em]
\text{D)}&(\theta_\eta(\theta Hu)^2H,s,u)^3,
(\theta_l(\theta Lu)^2H,s,u)^4,\\[0.3em]
\text{E)}&((KLu)^3H,s,u)^4,
((KH^2u)^4H,s,u)^4,\\[0.3em]
\text{F)}&((j\nu x)^2H,s,u)^3,
((j\nu x)^2H,s,u)^4,
((j\eta x)H,s,u)^4,\\
&(H_jH,s,u)^3,
(L_jH,s,u)^4.
\end{array}
\]
\noindent\emph{Reducciones:} $\nu$ no entra en plegamiento, de modo análogo a $j$.

B) y D). $(a_\nu^2H,s,u)^4$ y $(\theta_\nu^2H,s,u)^4$ se descomponen por las fórmulas (27.)c) y (29.)c), teniendo en cuenta $(gHu)^2=-\frac13s\sigma$:
\[
(a_\nu^2H,s,u)^4=\frac13(asu)^4\sigma,
\qquad
(\theta_\nu^2H,s,u)^4=-\frac16(\theta su)^4\sigma;
\]
de ahí la descomposición de $(a_\eta(aHu)H,s,u)^4$, $(\theta_\eta(\theta Hu)H,s,u)^4$.

$(\theta_\nu^2H,s,u)^3$, $(\theta_\nu^3H,s,u)^3$ y, por tanto, $(\theta_\eta^2(\theta Hu)H,s,u)^4$, se descomponen según \S\ 25, formas de orden 12, C).

\noindent 3) \emph{Plegamiento de $s$ con el sistema III--III}, es decir, de $s$ con formas del sistema III:
\[
(3,\lambda)(3,\lambda),\quad (3,\lambda)(2,\lambda),\quad (3,\lambda)(1,\lambda),\quad
(2,\lambda)(2,\lambda),\quad (2,\lambda)(1,\lambda),\quad (1,\lambda)(1,\lambda).
\]
\noindent\emph{Formas irreducibles:}
\[
\begin{array}{rl}
\text{A)}&(a_\nu^2a_\nu^2,s,u)^3,
(a_\nu^2a_\nu^2,s,u)^4,
(a^2a_\eta(aHu)^2,s,u)^3,\\[0.3em]
\text{B)}&(a_\nu H_j,s,u)^4,
((aHu)^2H_j,s,u)^3,
((aLu)^2H_j,s,u)^4,
((aLu)^3H_j,s,u)^2,
((aH^2u)^3H_j,s,u)^3.
\end{array}
\]
\noindent\emph{Reducciones:}

Los productos II--III, para el mismo orden en los símbolos $\nu$ y $f$, deben considerarse superiores a los productos III--III.

Las reducciones surgen en parte de relaciones entre productos (sizigias) y en parte de sustituir los productos por las formas descompuestas correspondientes y reducir el plegamiento de estas formas con $s$. Los productos que contienen contravariantes siempre se omiten, puesto que pasan a productos.

A) $f$ cuadrático, símbolos $\nu$ de orden arbitrario. Por \S\ 11, fórmulas (12.), y por cálculo análogo, se tiene:
\begin{align*}
1)\quad 0&=a_\nu b_{\nu_1}+\frac12f(aHu)^2-\frac14\theta H+\frac12u_xH_\vartheta-\frac14u_x^2H_\vartheta^2,\\
2)\quad 0&=a_\nu b_{\nu_1}^2+f a_\eta(aHu)^2-\theta_\eta-\frac12H_\vartheta,\\
3)\quad 0&=a_\nu^2b_{\nu_1}^2-H_\vartheta^2+2\theta_\eta^2.
\end{align*}
De ahí se sigue:
\begin{align*}
4)\quad 0&=a_\nu^2b_{\nu_1}^2(j\nu_1x),\\
5)\quad L_\vartheta(\theta Lu)&=-a_\nu b_\eta(bHu)^2+\frac34a_\nu^2(bHu)^2+\frac34\theta_\nu^2H,\\
6)\quad L_\vartheta(\theta Lu)&=-6a_\nu b_\eta(bHu)^2+2a_\nu^2(bHu)^2-\frac12\theta_\nu^2H,\\
7)\quad 0&=a_\nu(bHu)^2+f(aLu)^3+\theta_\nu H,\\
8)\quad 0&=a_\nu(bLu)^3+\frac34(\theta Hu)^2H,\\
9)\quad 0&=(aHu)^2(bHu)^2+2(\theta Hu)^2H,\\
10)\quad 0&=a_\eta(aHu)^2b_\eta(bHu)^2,\\
11)\quad 0&\equiv [a_\nu^2b_\eta(bHu)]_{\widehat{su^4}}.
\end{align*}
Las reducciones de
\[
(H_\vartheta su)^4,
\qquad (L_\vartheta(\theta Lu),s,u)^3,
\qquad (L_\vartheta(\theta Lu),s,u)^4
\]
(fórmulas (31.) y (32.)) dan, junto con las fórmulas 1) a 11), la reducción de todos los plegamientos de $s$ con productos que son cuadráticos en los símbolos $f$ y arbitrarios en $\nu$.

B) Productos de dos de los símbolos $f,\theta,j$; $\nu$ de orden arbitrario. Por las fórmulas (17.), (18.), (20.) y por cálculo directo valen las relaciones:
\begin{align*}
1)\quad 0&=a_\nu\theta_{\nu_1}+\frac14\theta(aHu)^2+\frac14f(\theta Hu)^2,\\
2)\quad 0&=\theta_\nu\theta_{\nu_1}+\frac12\theta(\theta Hu)^2.
\end{align*}
De ahí se sigue:
\begin{align*}
3)\quad 0&=3a_\nu(\theta Hu)^2+\theta(aLu)^3+2f(\theta Lu)^3,\\
4)\quad 0&=3\theta_\nu(aHu)^2+f(\theta Lu)^3+2\theta(aLu)^3,\\
5)\quad 0&=a_\nu(\theta Lu)^3,
&0&=\theta_\nu(aLu)^3,
&0&=(aHu)^2(\theta Hu)^2,\\
6)\quad 0&=a_\nu\theta_\nu^2+\theta_\nu a_\nu^2+f\theta_\eta(\theta Hu)^2+\theta a_\eta(aHu)^2,\\
7)\quad 0&=\theta_\nu\theta_\nu^2+\theta\theta_\eta(\theta Hu)^2+\frac16fH_j,\\
8)\quad 0&=a_\nu(j\nu x)^2+fH_j,\\
9)\quad 0&=(aHu)^2(j\nu x)^2-2a_\nu H_j,\\
10)\quad 0&=\theta_\nu(j\nu x)^2,
&0&=\theta H_j,\\
11)\quad 0&=a_\nu\theta_\nu^3-\frac34(j\eta x)^2,\\
12)\quad 0&=a_\nu^2\theta_\nu^2-\frac13(j\eta x)^2-\frac13(\D Hu)^2,\\
13)\quad 0&=a_\nu^2\theta_\nu^3-\frac12a_\sigma,\\
14)\quad 0&=\theta_\nu\theta_\nu^3,
&0&=a_\eta H_j.
\end{align*}
Las fórmulas se han llevado hasta el punto en que los productos se vuelven polarizables; esto es, por plegamiento sólo aparece un nuevo producto, porque a) el plegamiento del producto consigo mismo se vuelve reducible, y b) los restantes productos que surgen por plegamiento se vuelven reducibles o conducen a contravariantes (cf. \S\ 3. 2), a) y b)).

Las fórmulas 1) a 5) dan la reducción de todos los plegamientos de $s$ con productos que surgen de $a_\nu\theta_\nu$ por plegamiento. Las fórmulas 6) a 10) dan la reducción de los productos que surgen de $a_\nu\theta_\nu^2$ por plegamiento. Según \S\ 24 A), teniendo en cuenta las fórmulas 6) a 10), se descomponen los plegamientos de $s$ con productos que surgen de $a_\nu^2\theta_\nu$.

Se tiene:
\begin{align*}
0&\equiv a_\nu^2(asu)^2\theta_{\nu_1}(\theta su)^2,
&0&\equiv a_\nu^2(asu)\theta_{\nu_1}(\theta su)^3-K_\eta(KHu)^2(Ksu)^3,\\
0&\equiv a_\nu^2(asu)^2\theta_{\nu_1}(\theta su),
&0&\equiv a_\nu^2(asu)\theta_{\nu_1}(\theta su)^2,\\
0&\equiv a_\nu^2(asu)\theta_{\nu_1}(\theta su).
\end{align*}
Los reductores
\[
 ((j\eta x)^2,s,u)^3\ [\hbox{de }(\widehat{\vartheta\eta su})^2(Hsu),\ \S\ 23.\ C3],
\]
\[
(H_j(j\eta x),s,u)^2\ [\S\ 24,\ B],
\quad ((\D Hu)^2,s,u)^2\ [\S\ 24,\ C],
\quad (a_\sigma su)^2\ [\S\ 24,\ E]
\]
dan, junto con las fórmulas 11) a 14), la reducción de todos los productos que surgen de $a_\nu\theta_{\nu_1}^3$, $a_\nu^2\theta_{\nu_1}^2$, $a_\nu^2\theta_{\nu_1}^3$.

Nuestros cálculos anteriores se resumen en la siguiente \emph{conclusión}:

El sistema relativamente completo módulo $(\R,t)$ consta de las siguientes 331 formas y subsistemas, que reproducimos también a continuación ordenados en tablas.
\clearpage
\begin{landscape}
\thispagestyle{plain}
\begin{center}\bfseries Tabla I (véase p. 90)\end{center}
\vspace{-0.6em}
\begingroup
\setlength{\tabcolsep}{2pt}
\renewcommand{\arraystretch}{1.35}
\tiny
\begin{center}
\begin{adjustbox}{max width=0.98\linewidth,max totalheight=0.78\textheight,center}
\begin{tabular}{|c|c|c|c|c|c|c|c|c|c|c|c|c|c|c|c|c|c|c|c|c|c|c|}
\hline
\diagbox{$u$}{$x$} & 0 & 1 & 2 & 3 & 4 & 5 & 6 & 7 & 8 & 9 & 10 & 11 & 12 & 13 & 14 & 15 & 16 & 17 & 18 & 19 & 20 & 21 \\ \hline
0 & \mcell{i} &  &  &  & \mcell{f} &  & \mcell{\A} &  & \mcell{K_{\nu}(k\nu x)^3} &  & \mcell{K_{\eta}(k\eta x)^3} &  &  &  & \mcell{(\vartheta^3\eta x)^4} & \mcell{H_k(k\vartheta,\eta x)^4} &  & \mcell{\theta_l(\vartheta k,lx)^5} &  &  &  & \mcell{(\vartheta^2lx)^6} \\ \hline
1 &  & \mcell{u_x} &  &  &  & \mcell{\theta_\nu(\vartheta\nu x)^2} &  & \mcell{\theta_\eta(\vartheta\eta x)^2} & \mcell{N} & \mcell{(k\nu x)^3} & \mcell{\theta_\nu(\vartheta^2\nu x)^3} & \mcell{(k\eta x)^3} & \mcell{H_{\vartheta}(\vartheta^2\eta x)^3\\ \theta_\eta(\vartheta^2\eta x)^3} &  & \mcell{\theta_l(\vartheta^2lx)^4} &  & \mcell{(\vartheta k,\eta,x)^4} &  & \mcell{(\vartheta k,l,x)^5} &  &  &  \\ \hline
2 &  &  & \mcell{a_\nu^2} &  & \mcell{\theta\\ a_\eta^2} &  & \mcell{(\vartheta\nu x)^2} & \mcell{K_\nu(k\nu x)^2} & \mcell{(\vartheta\eta x)^2} & \mcell{H_k(k\eta x)^2\\K_\eta(k\eta x)^2} &  & \mcell{(\vartheta^2\nu x)^3\\ K_l(klx)^3} &  & \mcell{(\vartheta^2\eta x)^3} &  & \mcell{(\vartheta^2lx)^4} &  &  &  &  &  &  \\ \hline
3 &  & \mcell{\theta_\nu^3} &  & \mcell{a_\nu} & \mcell{\theta_\nu(\vartheta\nu x)} & \mcell{\A_\nu\\a_\eta} & \mcell{K\\H_{\vartheta}(\vartheta\eta x)\\\theta_\eta(\vartheta\eta x)} & \mcell{\A_\eta} & \mcell{(k\nu x)^2\\\theta_l(\vartheta lx)^2} &  & \mcell{(k\eta x)^2} &  & \mcell{(klx)^3} &  &  &  &  &  &  &  &  &  \\ \hline
4 & \mcell{\nu} &  & \mcell{H\\\theta_\nu^2} & \mcell{(j\nu x)^3\\a_\eta(aHu)} & \mcell{\theta_\eta^2} & \mcell{(\vartheta\nu x)\\a_l^2} &  & \mcell{N_\nu\\(\vartheta\eta x)} &  & \mcell{H_\eta\\N_\eta\\(\vartheta lx)^2} &  &  &  &  &  &  &  &  &  &  &  &  \\ \hline
5 &  & \mcell{a_\eta(aHu)^2} & \mcell{\theta_\eta^2(\theta Hu)} & \mcell{\theta_\nu} & \mcell{\frac{K_\nu^2}{(aHu)}} & \mcell{\theta_\eta} & \mcell{K_\eta^2\\(\A Hu)\\a_l} &  & \mcell{\A_l} &  &  &  &  &  &  &  &  &  &  &  &  &  \\ \hline
6 & \mcell{j\\\sigma} &  & \mcell{(j\nu x)^2\\(aHu)^2} & \mcell{L\\a_\nu^2(j\nu x)\\H_{\vartheta}(\theta Hu)\\\theta_\eta(\theta Hu)} & \mcell{\frac{K_\nu}{\theta_l^2}} &  &  & \mcell{K_\eta} &  &  &  &  &  &  &  &  &  &  &  &  &  &  \\ \hline
7 &  & \mcell{\theta_\eta(\theta Hu)^2} & \mcell{H_j(j\eta x)\\a_l(aLu)^2} &  & \mcell{a_\nu(j\nu x)\\(\theta Hu)} &  & \mcell{\A_\nu(j\nu x)\\\theta_l} & \mcell{K_l^2} &  &  &  &  &  &  &  &  &  &  &  &  &  &  \\ \hline
8 & \mcell{H_j^2\\a_l(aLu)^3} & \mcell{(\nu\sigma x)\\(j\nu x)} & \mcell{(\theta Hu)^2} & \mcell{K_\eta(KHu)^2\\(j\eta x)\\(aLu)^2} &  &  &  &  &  &  &  &  &  &  &  &  &  &  &  &  &  &  \\ \hline
9 &  & \mcell{H_j\\(aLu)^3} & \mcell{a_\eta(aHu)H_j\\L_{\vartheta}(\theta Lu)^2\\\theta_l(\theta Lu)^2} &  & \mcell{(KHu)} &  &  &  &  &  &  &  &  &  &  &  &  &  &  &  &  &  \\ \hline
10 & \mcell{\theta_l(\theta Lu)^3} & \mcell{L_j^2\\(j\sigma x)\\a_\eta(aH^2u)^3} &  & \mcell{H_j(aHu)\\(\theta Lu)^2} &  &  &  &  &  &  &  &  &  &  &  &  &  &  &  &  &  &  \\ \hline
11 &  & \mcell{(\theta Lu)^3} & \mcell{K_l(KLu)^3\\L_j\\(aH^3u)^3} &  &  &  &  &  &  &  &  &  &  &  &  &  &  &  &  &  &  &  \\ \hline
12 & \mcell{(aH^2u)^4} & \mcell{a_\eta(aLu)^2L_j\\H_{\vartheta}(\theta H^2u)^3\\\theta_\eta(\theta H^2u)^3} &  & \mcell{(KLu)^3} &  &  &  &  &  &  &  &  &  &  &  &  &  &  &  &  &  &  \\ \hline
13 & \mcell{(\nu\sigma j)} &  & \mcell{(aLu)^2L_j\\(\theta H^3u)^3} &  &  &  &  &  &  &  &  &  &  &  &  &  &  &  &  &  &  &  \\ \hline
14 & \mcell{(\theta H^2u)^4} & \mcell{K_\eta(KH^2u)^4\\(a_\eta H\!\cdot L,u)^4} &  &  &  &  &  &  &  &  &  &  &  &  &  &  &  &  &  &  &  &  \\ \hline
15 & \mcell{L_9(\theta,H\!\cdot L,u)^4} &  & \mcell{(KH^2u)^4} &  &  &  &  &  &  &  &  &  &  &  &  &  &  &  &  &  &  &  \\ \hline
16 & \mcell{(\theta,H\!\cdot L,u)^5} &  &  &  &  &  &  &  &  &  &  &  &  &  &  &  &  &  &  &  &  &  \\ \hline
17 & \mcell{K_\eta(K,H\!\cdot L,u)^5} &  &  &  &  &  &  &  &  &  &  &  &  &  &  &  &  &  &  &  &  &  \\ \hline
18 &  & \mcell{(K,H\!\cdot L,u)^5} &  &  &  &  &  &  &  &  &  &  &  &  &  &  &  &  &  &  &  &  \\ \hline
19 &  &  &  &  &  &  &  &  &  &  &  &  &  &  &  &  &  &  &  &  &  &  \\ \hline
20 &  &  &  &  &  &  &  &  &  &  &  &  &  &  &  &  &  &  &  &  &  &  \\ \hline
21 & \mcell{(KH^3u)^6} &  &  &  &  &  &  &  &  &  &  &  &  &  &  &  &  &  &  &  &  &  \\ \hline
\end{tabular}
\end{adjustbox}
\end{center}
\vspace{-0.2em}
\begin{center}\footnotesize (Los números de la fila horizontal superior indican el grado en las $x$; los de la primera fila vertical, el grado en las $u$.)\end{center}
\endgroup
\end{landscape}
\clearpage
\begin{landscape}
\thispagestyle{plain}
\begin{center}\bfseries Tabla II (véase p. 90)\end{center}
\vspace{-0.6em}
\begingroup
\setlength{\tabcolsep}{2pt}
\renewcommand{\arraystretch}{1.35}
\tiny
\begin{center}
\begin{adjustbox}{max width=0.98\linewidth,max totalheight=0.78\textheight,center}
\begin{tabular}{|c|c|c|c|c|c|c|c|c|c|c|c|c|c|c|c|c|c|}
\hline
\diagbox{$u$}{$x$} & 0 & 1 & 2 & 3 & 4 & 5 & 6 & 7 & 8 & 9 & 10 & 11 & 12 & 13 & 14 & 15 & 16 \\ \hline
0 & \mcell{J} &  &  &  & \mcell{s} &  & \mcell{s_j^3} &  &  &  &  & \mcell{s_\nu} & \mcell{(k\nu x)^3s_\nu} & \mcell{(\vartheta^3\nu x)^2s_{\vartheta}\\(\vartheta^2\nu x)^2s_{\vartheta}} &  & \mcell{\theta_\eta(\vartheta^2\eta x)^2s_\eta} &  \\ \hline
1 &  &  &  &  &  &  & \mcell{(asu)} & \mcell{s_\eta} & \mcell{s_l^2\\(\A su)} & \mcell{s_\eta(Nsu)\\(\vartheta\nu x)^2s,u} & \mcell{((k\nu x)^3,s,u)_s\\(K_\nu(k\nu x)^3,s,u)} &  & \mcell{K_\eta(k\eta x)^3,s,u} &  & \mcell{(\vartheta^2\nu x)^2s_\nu} &  & \mcell{(\vartheta^2\eta x)^4,s,u} \\ \hline
2 &  &  &  &  & \mcell{(asu)^2} & \mcell{s_\theta(\theta su)} & \mcell{(\A su)^2/a_\nu s_\nu} & \mcell{(\vartheta\nu x)^2,s,u\\(\theta_l(\vartheta\nu x)^2,s,u)} &  & \mcell{\theta_\eta(\vartheta^2\eta x)^2,s,u} &  & \mcell{(k\eta x)^2s_\nu\\((k\nu x)^2,s,u)} &  & \mcell{(k\eta x)^3,s,u} &  &  &  \\ \hline
3 &  &  & \mcell{(asu)^3} & \mcell{s_\vartheta(\theta su)^2\\s_\nu} & \mcell{a_\nu s_\nu(asu)\\a_l^2(asu)} & \mcell{s_\eta} & \mcell{(\theta su)\\a_l^2su} &  & \mcell{(Nsu)^2\\(\vartheta\eta x)s_\nu\\((\vartheta\nu x)^2,s,u)} & \mcell{((k\nu x)^2,s,u)^2} & \mcell{(\vartheta^2\eta x)^2,s,u} &  &  &  &  &  &  \\ \hline
4 & \mcell{(asu)^4} & \mcell{s_\vartheta(\theta su)^3} & \mcell{a_\nu^2(asu)^2} & \mcell{(\theta_\nu^2su)} & \mcell{(\theta su)^2} & \mcell{s_k(Ksu)^3\\a_\nu(asu)\\s_\nu(asu)} & \mcell{((\vartheta\nu x)^2,s,u)^2} & \mcell{s_\nu(\A su)\\(\A su)\\(aHu)s_\eta\\(a_\nu su)} &  & \mcell{(\A_\eta su)} &  &  &  &  &  &  &  \\ \hline
5 &  &  & \mcell{(\theta su)^3} & \mcell{s_k(Ksu)^3,s_j\\s_\vartheta(\theta s^\prime u)^4\\s_\nu(asu)^2\\a_\nu(asu)^2} & \mcell{(Hsu)\\((\vartheta\nu x)^2,s,u)\\(\theta_\nu(\vartheta\nu x)^2,s,u)\\(\theta_\nu^2su)} & \mcell{((j\nu x)^2,s,u)\\(aHu)^2s_\nu\\(a_\eta su)^2} & \mcell{(Ksu)^2\\s_\eta\\(\theta_\eta^2su)} &  & \mcell{((k\nu x)^2,s,u)^2\\K_\nu s_\nu} &  &  &  &  &  &  &  &  \\ \hline
6 & \mcell{(\theta su)^4} & \mcell{s_\nu(asu)^3\\a_\nu(asu)^3} & \mcell{(Hsu)^2\\(\theta_\nu(\vartheta\nu x),s,u)^3\\(\theta_\nu^2su)^2} & \mcell{a_\eta s_\nu^2\\(a_\eta(aHu),s,u)^2\\(a_\eta(aH^2u),s,u)} & \mcell{(Ksu)^3} & \mcell{s_\eta(asu)\\((\vartheta\nu x),s,u)^3\\(\theta_\eta su)} & \mcell{((k\nu x)^2,s,u)^3} & \mcell{s_\nu(\A su)\\a_\nu(j\nu x)s_\nu\\(\theta H u)s_\eta\\(\theta_\eta su)} &  &  &  &  &  &  &  &  &  \\ \hline
7 & \mcell{\theta_\nu(\vartheta\nu x),s,u)^4} & \mcell{(a_\eta(aHu),s,u)^3} & \mcell{(Ksu)^4\\(\theta_\eta^2(\theta Hu),s,u)^2\\(a_\eta^2-a_l^2,s,u)^3} & \mcell{s_j(asu)^3\\s_k(K^2u)^3\\((j\nu x),s,u)^2\\(\theta_\eta su)^2} & \mcell{((j\nu x),s,u)\\((j\eta x),s,u)\\(aHu),s,u)\\(aH^2u),s,u)} & \mcell{((\vartheta\eta x),s,u)^3} & \mcell{(aLu)^3s_l/(asu)^3} &  &  &  &  &  &  &  &  &  &  \\ \hline
8 & \mcell{(a^2-a_l^2,s,u)^4} & \mcell{s_j(asu)^3\\s_k(K^2u)^3\\((j\nu x),s,u)^4\\(\theta_\eta su)^3} & \mcell{((j\nu x)^2,s,u)^2\\(aHu),s,u)^3\\((aHu)^2,s,u)^3} & \mcell{(Lsu)^2\\(a_l^2(j\nu x),s,u)^3\\H_\vartheta(\theta Hu),s,u)^3\\\theta_l(\theta Hu),s,u)^3} & \mcell{(asu)^3} & \mcell{(K,su)^3} &  & \mcell{(K_\vartheta su)^2} &  &  &  &  &  &  &  &  &  \\ \hline
9 & \mcell{K_j^2su^4\\((aHu),s,u)^4} & \mcell{(Lsu)^4\\(a_l^2(j\nu x),s,u)^4\\(\theta_\nu(\theta Hu),s,u)^2} & \mcell{(K_\nu^3u)^6\\(a_\eta(aLu)^2,s,u)^3\\(a_l^2H,s,u)^3} & \mcell{(K,su)^3} & \mcell{(a_\nu(j\nu x),s,u)^3\\(\theta Hu),s,u\\(\theta H u)^2,s,u} &  & \mcell{(\theta_lsu)^3} &  &  &  &  &  &  &  &  &  &  \\ \hline
10 &  &  & \mcell{(K,su)^4\\(a_l^2a_\eta(aHu)^2,s,u)^3} & \mcell{((j\eta x),s,u)^3\\(H_j,su)\\((aLu)^2,s,u)^2\\((aLu)^3,s,u)} &  &  &  &  &  &  &  &  &  &  &  &  &  \\ \hline
11 & \mcell{(a_\nu(j\nu x),s,u)^4\\((\theta Hu),s,u)^4} & \mcell{(K_l(KH u)^2,s,u)^3\\((j\eta x),s,u)^3\\((aLu)^2,s,u)^4\\(a_\eta H,s,u)^3} & \mcell{(H^2su)^3\\\theta_l(\theta Lu)^2,s,u)^3} &  & \mcell{((KHu)^2,s,u)^3} &  &  &  &  &  &  &  &  &  &  &  &  \\ \hline
12 &  & \mcell{(a_\eta(aHu)^2H,s,u)^3} & \mcell{((KH u)^3,s,u)^3} & \mcell{((aHu)H_j,s,u)^2\\((\theta Lu)^2,s,u)^3} &  &  &  &  &  &  &  &  &  &  &  &  &  \\ \hline
13 & \mcell{((KH u)^3,s,u)^4} & \mcell{((aHu)H_j,s,u)^3\\((\theta Lu)^2,s,u)^4\\(\nu\cdot H,s,u)^3} & \mcell{(L_jsu)^2\\((aH^3u)^3,s,u)^3\\((j\eta x)H,s,u)^3\\((aHu)^2H,s,u)^3} &  &  &  &  &  &  &  &  &  &  &  &  &  &  \\ \hline
14 & \mcell{((j\nu x)^2H,s,u)^4\\((aHu)^2H,s,u)^4} & \mcell{(\theta_\eta(\theta Hu)^2H,s,u)^3} &  & \mcell{((KLu)^3,s,u)^2} &  &  &  &  &  &  &  &  &  &  &  &  &  \\ \hline
15 & \mcell{(a_l(aLu)^3H,s,u)^4} & \mcell{((KLu)^3,s,u)^3} & \mcell{((aLu)^2L_j,s,u)^2\\((\theta Hu)^3H,s,u)^3\\((\theta Hu)^2H,s,u)^2} &  &  &  &  &  &  &  &  &  &  &  &  &  &  \\ \hline
16 & \mcell{((\theta Hu)^2H,s,u)^4\\(\sigma\cdot H_j,s,u)^4} & \mcell{((j\eta x)H,s,u)^4\\(H_jH,s,u)^3\\((aLu)^2H,s,u)^4\\((aLu)^3H,s,u)^3} &  &  &  &  &  &  &  &  &  &  &  &  &  &  &  \\ \hline
17 & \mcell{(\theta(\theta Lu)^2H,s,u)^4} &  & \mcell{((KH^3u)^3,s,u)^3} &  &  &  &  &  &  &  &  &  &  &  &  &  &  \\ \hline
18 &  & \mcell{((\theta Lu)^2H,s,u)^4\\((\theta Lu)^3H,s,u)^3\\(aHu)^2H,s,u)^4} &  &  &  &  &  &  &  &  &  &  &  &  &  &  &  \\ \hline
19 & \mcell{((aH^3u)^3H,s,u)^4\\(L_jH,s,u)^5} &  &  &  &  &  &  &  &  &  &  &  &  &  &  &  &  \\ \hline
20 &  & \mcell{((KLu)^3H,s,u)^5} & \mcell{((aLu)^3H,s,u)^5} &  &  &  &  &  &  &  &  &  &  &  &  &  &  \\ \hline
21 & \mcell{((\theta H^3u)^3H,s,u)^4\\((aLu)^2H_j,s,u)^4} &  &  &  &  &  &  &  &  &  &  &  &  &  &  &  &  \\ \hline
22 &  &  &  &  &  &  &  &  &  &  &  &  &  &  &  &  &  \\ \hline
23 & \mcell{((KH^3u)^4H,s,u)^4} & \mcell{((aH^3u)^4H,s,u)^3} &  &  &  &  &  &  &  &  &  &  &  &  &  &  &  \\ \hline
\end{tabular}
\end{adjustbox}
\end{center}
\vspace{-0.2em}
\begin{center}\footnotesize (Los números de la fila horizontal superior indican el grado en las $x$; los de la primera fila vertical, el grado en las $u$.)\end{center}
\endgroup
\end{landscape}




\clearpage

\providecommand{\pair}[2]{(#1\,|\,#2)}
\providecommand{\eps}{\varepsilon}
\providecommand{\rhoidx}{\varrho}

\begin{center}
{\Large\bfseries 3. Sobre la teoría de invariantes de las formas en $n$ variables\footnote{Conferencia pronunciada en la reunión de naturalistas de Salzburgo, 1909.}}\par
\vspace{1.2em}
Jahresbericht der DMV 19 (1910), pp. 101--104
\end{center}

\vspace{1em}
En la teoría proyectiva de invariantes de las formas en $n$ variables está resuelto el problema principal: se ha demostrado la finitud del sistema de formas. Sin embargo, no se ha tratado una serie de cuestiones ulteriores, a saber, las que se refieren a métodos para investigar la conexión entre las formas. Quisiera comunicar brevemente los resultados que he obtenido respecto de tales cuestiones; pero, por claridad, esbozaré primero los planteamientos en el dominio ternario, donde son conocidos.

Pienso ante todo en los dos teoremas fundamentales del método simbólico: el teorema de que todas las formaciones invariantes pueden representarse simbólicamente, y el segundo, de que todas las relaciones existentes entre invariantes se obtienen por aplicación sucesiva de un pequeño número de identidades simbólicas; es decir, que el método simbólico da todas las relaciones sin salir del dominio de los invariantes.\footnote{Study: \emph{Methoden zur Theorie der ternären Formen}. II. \S\ 6. (Leipzig, Teubner, 1889.)} Sobre esta base vienen luego los procesos de generación de las formas, el proceso simbólico de plegamiento y los procesos de diferenciación no simbólicos, la teoría de los desarrollos en serie y análogos. En este contexto quisiera mencionar además un teorema demostrado en mi disertación,\footnote{\emph{Über die Bildung des Formensystems der ternären biquadratischen Form}. \S\ 1. (Crelle's Journal Bd. 134. 1908.)} a saber, que todos los plegamientos pueden generarse por aplicación sucesiva de los dos plegamientos ``cogredientes'' posibles; por ello entiendo, para un producto simbólico $s_x t_x u_\sigma u_\tau$, los dos plegamientos $(stu)$ y $(\sigma\tau x)$. Sobre este teorema puede construirse la teoría de los reductores y de la reducción de los sistemas de formas, análogamente al dominio binario, como he mostrado en mi disertación.

Si pasamos ahora al dominio $n$-ario, ya el primer teorema del método simbólico sólo vale en medida condicionada. Como es sabido, ya en las formas cuaternarias aparecen, además de las coordenadas de punto y de plano $x$ y $u$, coordenadas de recta $p_{ik}$, los subdeterminantes de dos filas de coordenadas de punto. Análogamente, en el dominio $n$-ario tenemos filas de variables $p_\rho$, cuyos elementos individuales son los subdeterminantes de $\rho$ filas $x$, y filas simbólicas $S^\rho$, cuyos elementos se comportan como subdeterminantes de $\rho$ filas $s$ cogredientes con las $u$. La representación simbólica por determinantes, como es sabido, sólo vale cuando las filas simbólicas $S^\rho$ se resuelven en las filas $s$ y éstas se distribuyen entre distintos determinantes; entonces sólo tienen significado real aquellos agregados de determinantes que dependen de las $S^\rho$. Pero no se puede reconocer de conjunto qué agregados de determinantes logran esto; y con ello se pierden todos los demás teoremas mencionados en el dominio ternario.

Quisiera ahora indicar un principio simple para obtener una representación simbólica explícita en las filas $S^\rho$, y con ello recuperar también los restantes teoremas. Se trata de una aplicación del conocido teorema sobre las matrices correspondientes: ``A cada fila de variables $p_\rho$ se le puede asociar biunívocamente otra fila $q_{n-\rho}$, cuyos elementos individuales son subdeterminantes de $n-\rho$ filas $u$ unidas a las $\rho$ filas $x$ por ecuaciones $(u\mid x)=0$.'' Correspondientemente, a cada fila simbólica $S^\rho$ se le asocia otra $S_{n-\rho}$, que se comporta como si estuviera compuesta de $n-\rho$ filas cogredientes con las $x$.

El teorema admite todavía una segunda formulación, apropiada para las aplicaciones: ``A cada producto matricial $(v^{(1)}\cdots v^{(\rho)}\mid y^{(1)}\cdots y^{(\rho)})$,\footnote{Es decir, la suma de los productos de determinantes correspondientes, formados a partir de la matriz de las $\rho$ filas $v$ y de la de las $\rho$ filas $y$.} donde las filas de la segunda matriz son contragredientes con las de la primera, se le asocia un determinante; y, recíprocamente, todo determinante puede transformarse en un producto matricial con un número prefijado $\rho$ de filas.'' Así, determinante y producto matricial aparecen como completamente equivalentes, y el primer teorema del método simbólico toma la forma: ``Todas las formaciones invariantes pueden representarse simbólicamente por productos matriciales.'' A partir de dos filas simbólicas $S^\sigma$, $T^\tau$ se puede ahora, con ayuda de filas de variables, formar muy fácilmente un producto matricial que contenga explícitamente esas filas simbólicas y que, por tanto, represente una formación invariante de la forma $(S^\sigma\mid p_\sigma)(T^\tau\mid p_\tau)$, a saber:
\begin{equation}
 \pair{q_{n-\sigma-\lambda}S^\sigma}{T_{n-\tau}p_{\tau-\lambda}},
 \qquad (\lambda=1,2,\ldots,\tau\ \text{o}\ n-\sigma).
\tag{1}
\end{equation}

Más importante es la recíproca: que todas las formaciones invariantes puedan representarse explícitamente por productos matriciales, de modo que el proceso anterior significa la ampliación del proceso de plegamiento binario y ternario. Para demostrar esta recíproca es necesario trasladar el segundo teorema del método simbólico al dominio $n$-ario, es decir, establecer la totalidad de las identidades independientes e irreducibles en las que todas las filas $S^\rho$, $p_\rho$ se consideran irreducibles. Estas identidades se obtienen de las conocidas, que sólo contienen filas $x$ y $u$,\footnote{E. Pascal en \emph{Memorie d. R. Acc. d. Lincei}. (4) 5 (1888), p. 375.} sustituyendo el teorema del producto para determinantes por un sistema de teoremas del producto para matrices, y contrayendo distintos productos matriciales en uno solo precisamente mediante estos teoremas del producto. Se llega así a dos fórmulas duales, de las cuales sólo se cita una:
\begin{equation}
\begin{aligned}
 \pair{S_1^{\rho_1}S_2^{\rho_2}\cdots S_k^{\rho_k}}{x p_{\rho-1}}
 &=\sum_{i=1}^{k}\eps\,\pair{S_1^{\rho_1}\cdots S_{i-1}^{\rho_{i-1}}S_{i+1}^{\rho_{i+1}}\cdots S_k^{\rho_k}q_{n-\rho_i+1}}{S_{n-\rho_i}x},\\
 \rho&=\sum_{i=1}^{k}\rho_i<n,\qquad \eps=\pm1.
\end{aligned}
\tag{2}
\end{equation}

La única especialización contenida en estas fórmulas, que entre una fila $x$, no puede evitarse; para filas simbólicas y de variables completamente generales se llega a una ``identidad de descomposición'', una identidad en la que una fila $p_\rho$ se descompone en las filas individuales $x$. Clebsch ya utilizó el caso especial más simple de la identidad de descomposición en su ``problema fundamental de la teoría de invariantes'' para derivar los covariantes elementales en los desarrollos en serie. Con ayuda de estas identidades, que median el paso de las expresiones determinantes no explícitas al producto matricial, se puede probar efectivamente que con la representación matricial se obtiene la representación simbólica explícita deseada.

Para el proceso de plegamiento vale, sin embargo, un teorema por completo análogo al del dominio ternario, sobre el cual también pueden construirse análogamente los teoremas de reducción. Si en (1) se designa el número $\lambda$ como defecto del plegamiento, y se llaman cogredientes los plegamientos para los que $\sigma=\tau$, entonces todos los plegamientos pueden generarse por aplicación sucesiva de los plegamientos cogredientes de defecto 1. La prueba de este teorema se basa esencialmente en que las formas más generales pueden sustituirse por ``formas normales'', es decir, por formas que se anulan idénticamente bajo la aplicación de todas las operaciones diferenciales invariantes. En el dominio cuaternario esta última circunstancia fue demostrada por Mertens;\footnote{Über invariante Gebilde quaternärer Formen. Wiener Berichte. Bd. 98. 1889.} nuestro conocimiento de las operaciones diferenciales más generales --que, como es sabido, proceden de los procesos de plegamiento al sustituir las filas simbólicas por símbolos diferenciales-- nos permite, usando la identidad de descomposición, transferir casi literalmente la demostración de Mertens al dominio $n$-ario.

\clearpage
\begin{center}
{\Large\bfseries 4. Sobre la teoría de invariantes de las formas en $n$ variables}\par
\vspace{1.2em}
Journal für die reine und angewandte Mathematik 139 (1911), pp. 118--154
\end{center}

\section*{Introducción.}
En la teoría proyectiva de invariantes de las formas en $n$ variables, las cuestiones que se adjuntan al problema principal --la demostración de la finitud del sistema de formas-- apenas han sido tratadas cuando se sale de los dominios binario y ternario. Se trata de las cuestiones relativas a la dependencia mutua de las formas y a los métodos para dominar esta conexión entre formas. Estos métodos descansan esencialmente --como se ha demostrado por completo en los dominios binario y ternario-- en dos teoremas designados por Study como los ``teoremas fundamentales del método simbólico'', a saber, el teorema sobre la representabilidad simbólica de las formas y el teorema sobre las identidades simbólicas.\footnote{E. Study, \emph{Methoden zur Theorie der ternären Formen} (Leipzig, Teubner, 1889). II. \S\ 6. Sobre la aparición de estos teoremas fundamentales también en la teoría de invariantes de grupos especiales de transformaciones, cf. Study, \emph{Das Apollonische Problem} (Math. Ann. Bd. 49 [1897]) y \emph{Zur Differentialgeometrie der analytischen Kurven} (Transactions of the American Math. Society, Bd. 1 [1909]).} El último teorema afirma que todas las relaciones existentes entre formaciones invariantes enteras se obtienen por aplicación sucesiva de un pequeño número de identidades simbólicas; así, el método simbólico, y sólo él, da conocimiento de la conexión entre formas sin salir del dominio de los invariantes.

Gordan, en su programa de Erlangen,\footnote{P. Gordan, \emph{Über das Formensystem binärer Formen}. S. 46. (Leipzig, Teubner, 1875).} ya señaló la razón de la dificultad del paso a $n$ variables: reside en que el primer teorema del método simbólico ya no es válido en su forma general. En el dominio $n$-ario aparecen, como es sabido, filas de variables de nivel superior $p_\rho$\footnote{Para la notación véase \S\ 1.} y, análogamente, filas simbólicas de nivel superior $S^\rho$; se llega a estas filas tanto cuando, como originalmente en Clebsch,\footnote{A. Clebsch, über eine Fundamentalaufgabe der Invariantentheorie. (Abh. der Ges. d. Wiss. zu Göttingen, Bd. XVII, 1872).} se parte de formas con un número arbitrario de filas $p_\rho$, como cuando se parte, siguiendo a F. Deruyts y A. Capelli,\footnote{Cf. A. Capelli, \emph{Lezioni sulla teoria delle forme algebriche} (Napoli 1902), \S\ 22--24, y la literatura allí citada.} de formas que contienen sólo un número arbitrario de filas cogredientes $x$. No existe una representación simbólica que contenga explícitamente las filas $p_\rho$, $S^\rho$;\footnote{Definición más precisa de ``representación explícita'' en \S\ 2.} es necesario resolver $p_\rho$, $S^\rho$ en las filas individuales $x$ y en las $s$ contragredientes con ellas, y distribuir éstas entre distintos determinantes. Es posible, ciertamente, verificar mediante condiciones diferenciales (cf. fórmula (19.)) que un agregado dado de determinantes tiene la propiedad de depender sólo de las filas $p_\rho$, $S^\rho$; pero de este modo no se obtiene una visión de conjunto de la totalidad de las formaciones invariantes ni de los procesos para generarlas.\footnote{Una ampliación del simbolismo debida a Study y muy valiosa para el cálculo (comunicada por F. Engel, Ber. d. K. Sächs. Ges. d. Wiss., math.-phys. Kl., 1900, p. 126, y para el dominio cuaternario por E. Study, \emph{Geometrie der Dynamen}, p. 135) tampoco es una representación explícita en nuestro sentido. Por ejemplo, difícilmente conduciría a los procesos de diferenciación no simbólicos para generar las formas.}

Mostraremos ahora en \S\S\ 2 y 6 cómo, con ayuda del ``teorema sobre las matrices correspondientes'',\footnote{Aparte de la \emph{Ausdehnungslehre} de Grassmann de 1862, n.º 112, primero en A. Brill, über zwei Berührungsprobleme, \S\ 2 (Math. Ann. Bd. 4. 1871).} se recupera en su generalidad el primer teorema fundamental sustituyendo la representación por determinantes por la representación mediante ``productos matriciales''. En \S\ 2 se muestra que todo producto matricial que contiene explícitamente las filas $S^\rho,p_\rho$ es una formación invariante de la forma fundamental; la recíproca dada en \S\ 6, que toda formación invariante puede representarse explícitamente por productos matriciales, sólo se logra mediante la transferencia del segundo teorema fundamental (\S\S\ 3--5).

Este teorema es conocido para formas que contienen sólo filas de variables $x$.\footnote{E. Pascal, Giorn. di mat. 26 (1888), pp. 33, 102; Rendic. d. Accad. d. Linc. (4) 4 (1888), p. 119; ib. Mem. (4) 5 (1888), p. 375.} El problema general, establecer la totalidad de las identidades irreducibles en las que todas las filas $S^\rho,p_\rho$ se consideran irreducibles, fue formulado por Study;\footnote{``Methoden \dots'', p. 204, nota 17).} al mismo tiempo ve en el número de identidades que aparecen una medida de la dificultad de los métodos en el dominio $n$-ario. Puesto que ahora es posible reunir la totalidad de las identidades en una sola fórmula (36.), esta dificultad queda reducida al menos al mínimo. En las aplicaciones, sin embargo, se presentarán a menudo ciertos casos especiales distinguidos de (36.), como (20.), (21.), caracterizados por admitir representación explícita sin sustitución de nuevas filas; o la fórmula (31.), designada como ``identidad de descomposición'', porque una fila $p_\rho$ parece descomponerse en las filas $y$.

Sobre los dos teoremas fundamentales puede ahora construirse fácilmente la teoría ulterior. El primer teorema proporciona directamente los procesos para generar las formas, el proceso simbólico de plegamiento y los procesos de diferenciación no simbólicos,\footnote{En el trabajo ``Invariante Gebilde von Nullsystemen'' (Sitzungsber. d. Ak. d. Wiss. Wien, Bd. 97, Abt. IIa, 1888), Mertens llegó por otro camino, no directamente generalizable, al proceso de plegamiento cuaternario. El invariante alternante de 6 filas $S^2$ que allí aparece difiere del que surge aquí, por una sola aplicación de la sustitución (11.), en un agregado de invariantes pares. Para un caso especial, el proceso de plegamiento cuaternario también se encuentra en P. Gordan, \emph{Das simultane System von zwei quadratischen quaternären Formen} (Math. Ann. Bd. 56. 1903).} mientras que la identidad de descomposición elimina entre esos procesos todos los equivalentes (\S\ 6). El conocimiento de los procesos de diferenciación proporciona además el traslado del concepto de ``forma normal'' al dominio $n$-ario y, por medio de un procedimiento de demostración dado por Mertens\footnote{F. Mertens, Über invariante Gebilde quaternärer Formen. (Sitzber. d. Ak. d. Wiss. Wien. Bd. 98. Abt. IIa. 1889.)} en el dominio cuaternario, el teorema de que un sistema de formaciones invariantes puede sustituirse por un sistema equivalente de formas normales (\S\ 7). Bajo esta última hipótesis mostré en mi disertación\footnote{Über die Bildung des Formensystems der ternären biquadratischen Form. (Dieses Journal, Bd. 134. 1908.)} que el proceso de plegamiento ternario puede generarse por aplicación sucesiva de dos plegamientos fundamentales. Sobre la base de este teorema y de la teoría de los desarrollos en serie desarrollé luego, introduciendo el concepto de ``fila de formas'', los principales teoremas de reducción para sistemas de formas. De modo completamente análogo, en el dominio $n$-ario el proceso de plegamiento puede generarse a partir de $n-1$ plegamientos fundamentales (\S\ 8), y pueden establecerse los teoremas de reducción (\S\ 9).

\paragraph{Nota final.} Con ocasión de la comunicación de Salzburgo, el Prof. E. Müller, de Viena, llamó mi atención sobre el hecho de que el cálculo con productos matriciales que subyace a este trabajo coincide, en principio, con el cálculo de la teoría de la extensión de Grassmann.\footnote{Hermann Graßmanns gesammelte mathematische und physikalische Werke. Ersten Bandes zweiter Teil: Die Ausdehnungslehre von 1862. In Gemeinschaft mit H. Graßmann d. Jüngeren herausgegeben v. Fr. Engel (Leipzig, Teubner 1896).} En efecto, el ``producto combinatorio'' de Grassmann $[S^{\rho_1}S^{\rho_2}\cdots S^{\rho_i}]$ puede considerarse como una matriz dada por esas filas (cf. \S\ 2): concretamente, el ``producto progresivo'' como la matriz de las propias filas, el ``regresivo'' como la matriz de las filas asociadas $S_{n-\rho_1},S_{n-\rho_2},\ldots,S_{n-\rho_i}$, mientras que la matriz correspondiente al ``producto mixto'' surge cuando varias filas $S$ se reúnen en nuevas filas $P,Q$ mediante la sustitución (9.). Nuestra ``fila asociada'' corresponde al ``complemento'' de Grassmann; nuestro ``producto matricial'', a un ``producto mixto de número de nivel 0''.

Bajo esta correspondencia, el desarrollo dado en la \emph{Ausdehnungslehre} de 1862, n.º 113, se vuelve idéntico a una fórmula dual de nuestra fórmula (32.). En el caso especial $\rho=1$, (32.) pasa a (17.), y la fórmula dual a (16.). A partir de estos casos especiales del desarrollo de Grassmann, Müller\footnote{E. Müller, Beiträge zur Graßmannschen Ausdehnungslehre. 1. Mitteilung: Einige allgemeine Sätze. (Sitzungsber. d. Ak. d. Wiss. Wien; Bd. 118. Abt. IIa. Juli 1909), Satz 16 und 18.} ha deducido recientemente las identidades (20.), (21.) mencionadas arriba.

En contraste con la derivación de Grassmann--Müller, nuestra derivación da al mismo tiempo la prueba de completitud de las identidades, lo cual parece ser el punto esencial para las cuestiones aquí consideradas; y, partiendo de (20.), (21.), avanza además hasta las identidades irreducibles más generales (36.), que no parecen haber sido advertidas hasta ahora.

\section*{\S\ 1. Notaciones y definiciones. Resumen de resultados conocidos.}
Denotamos, como de costumbre, por la fila de variables $x$ la fila de elementos $x_1,x_2,\ldots,x_n$, y análogamente por la fila de variables $p_\rho$ ($\rho=1,2,\ldots,n-1$) la fila de los $\binom n\rho$ elementos $p_{i_1i_2\ldots i_\rho}$, donde $p_{i_1i_2\ldots i_\rho}$ designa los $\binom n\rho$ determinantes formados a partir de la matriz de las $\rho$ filas $x^{(1)},x^{(2)},\ldots,x^{(\rho)}$,
\[
\left|\begin{array}{cccc}
 x_{i_1}^{(1)}&x_{i_2}^{(1)}&\cdots&x_{i_\rho}^{(1)}\\
 x_{i_1}^{(2)}&x_{i_2}^{(2)}&\cdots&x_{i_\rho}^{(2)}\\
 \vdots&\vdots&&\vdots\\
 x_{i_1}^{(\rho)}&x_{i_2}^{(\rho)}&\cdots&x_{i_\rho}^{(\rho)}
\end{array}\right|,
\qquad (i_1<i_2<\cdots<i_\rho=1,2,\ldots,n).
\]
Según el teorema sobre las matrices correspondientes,\footnote{Cf. por ejemplo Gordan--Kerschensteiner, \emph{Vorlesungen über Invariantentheorie}, Bd. I, p. 99.} a cada fila $p_\rho$ se le asocia biunívocamente una segunda fila $q_{n-\rho}$ por la relación, válida para cada elemento de la fila,
\begin{equation}
 p_{i_1i_2\ldots i_\rho}=(-1)^{\frac{\rho(\rho+1)}{2}+i_1+i_2+\cdots+i_\rho}q_{j_1j_2\ldots j_{n-\rho}},
\tag{1}
\end{equation}
donde los $i$ y los $j$ están cada uno en buen orden y juntos forman una permutación completa de los números 1 a $n$. La fila $q_{n-\rho}$ consta de los $\binom n\rho$ determinantes de grado $n-\rho$, $q_{j_1j_2\ldots j_{n-\rho}}$, formados a partir de la matriz de $n-\rho$ filas $u^{(1)},u^{(2)},\ldots,u^{(n-\rho)}$ contragredientes con las $x$. Estas filas satisfacen las $\rho(n-\rho)$ ecuaciones de condición
\begin{equation}
 (u^{(k)}\mid x^{(l)})=\sum_{\alpha=1}^{n}u_\alpha^{(k)}x_\alpha^{(l)}=0;
 \qquad (k=1,2,\ldots,n-\rho;\ l=1,2,\ldots,\rho).
\tag{2}
\end{equation}
Las supondremos normalizadas (con Clebsch\footnote{a. a. O. \S\ 5.}) de modo que el factor de proporcionalidad que en otro caso entraría en (1.) se ponga igual a uno.

Según Clebsch\footnote{a. a. O. \S\ 12.} (y asimismo según las investigaciones posteriores de F. Deruyts y A. Capelli, cf. la Introducción), las formas más generales pueden reducirse a aquellas que contienen a lo sumo una de cada una de las $n-1$ filas $p_\rho$, y por tanto pueden representarse simbólicamente como sigue:
\begin{equation}
 (S^1\mid p_1)^{m_1}(S^2\mid p_2)^{m_2}\cdots(S^{n-1}\mid p_{n-1})^{m_{n-1}},
\tag{3}
\end{equation}
donde, en analogía con (2.), hemos puesto
\[
 (S^\rho\mid p_\rho)=\sum S^{i_1i_2\ldots i_\rho}p_{i_1i_2\ldots i_\rho}.
\]
Como consecuencia de las relaciones idénticas no lineales existentes entre los elementos de una fila $p_\rho$, las filas simbólicas $S^\rho$ pueden normalizarse de modo que satisfagan exactamente las mismas identidades; es decir, de modo que se comporten como los determinantes de una matriz de $\rho$ filas cuyas filas $s^{(1)},s^{(2)},\ldots,s^{(\rho)}$ son contragredientes con las $x$.\footnote{Clebsch, a. a. O. \S\ 4.} Así, a cada fila $S^\rho$ puede asociársele biunívocamente otra $S_{n-\rho}$ mediante la relación
\begin{equation}
 S_{i_1i_2\ldots i_{n-\rho}}=(-1)^{\frac{(n-\rho)(n-\rho+1)}2+i_1+i_2+\cdots+i_{n-\rho}}S^{j_1j_2\ldots j_\rho},
\tag{4}
\end{equation}
donde de nuevo los $i$ y los $j$ están cada uno en buen orden y juntos agotan exactamente los números 1 a $n$. Los elementos $S_{i_1i_2\ldots i_{n-\rho}}$ de la fila $S_{n-\rho}$ se comportan como los determinantes formados de $(n-\rho)$ filas $\sigma^{(1)},\sigma^{(2)},\ldots,\sigma^{(n-\rho)}$ cogredientes con las $x$, y las filas $s$ y $\sigma$ están enlazadas por las $\rho(n-\rho)$ condiciones
\begin{equation}
 (s^{(k)}\mid \sigma^{(l)})=0,\qquad(k=1,2,\ldots,\rho;\ l=1,2,\ldots,n-\rho).
\tag{5}
\end{equation}
En consecuencia de las relaciones (1.) y (4.), cada factor de (3.) admite la representación simbólica cuádruple equivalente
\begin{equation}
 (S^\rho\mid p_\rho)=(S^\rho q_{n-\rho})=(S_{n-\rho}p_\rho)=(-1)^{\rho(n-\rho)}(q_{n-\rho}\mid S_{n-\rho}),
\tag{6}
\end{equation}
donde, como siempre en lo sucesivo, $(S^\rho q_{n-\rho})$ y $(S_{n-\rho}p_\rho)$ son abreviaturas de los determinantes $(s^{(1)}\cdots s^{(\rho)}u^{(1)}\cdots u^{(n-\rho)})$ y $(\sigma^{(1)}\cdots\sigma^{(n-\rho)}x^{(1)}\cdots x^{(\rho)})$; por desarrollo de Laplace éstos resultan depender sólo de las filas $S^\rho,q_{n-\rho}$ y, respectivamente, $S_{n-\rho},p_\rho$, como pretende indicar la notación anterior. Las expresiones $(S^\rho\mid p_\rho)$ y $(q_{n-\rho}\mid S_{n-\rho})$ significan, según la discusión precedente, los productos de las matrices de las filas $s$ y $x$, respectivamente de las filas $u$ y $\sigma$; por producto matricial se entiende la suma de los productos de determinantes correspondientes, es decir, el valor del determinante de la matriz producto.

El número característico $\rho(n-\rho)$ perteneciente a una fila simbólica (respectivamente, a una fila de variables) se llamará, siguiendo a Clebsch, el peso de la fila;\footnote{a. a. O. \S\ 11.} el número $\rho$, que da el número de filas $s$, el índice de peso superior; el número $n-\rho$, que da el número de filas $\sigma$, el índice de peso inferior. De las relaciones (4.) se sigue que una fila simbólica puede aparecer en una misma relación una vez con índice de peso superior y otra con índice de peso inferior; lo mismo vale para la aparición de las filas correspondientes $p_\rho$ y $q_{n-\rho}$. Queda por señalar que, en general, denotamos por $p$ las filas de variables de nivel superior cuya matriz asociada consta de filas $x$, y por $q$ aquellas cuya matriz consta de filas $u$; por letras griegas minúsculas $\sigma,\tau,\alpha,\beta$, las filas simbólicas cogredientes con las $x$; por letras latinas minúsculas $s,t,a,b$, las filas simbólicas cogredientes con las $u$; y todas las demás filas simbólicas por letras latinas mayúsculas con índice de peso: $S^\rho$, $T^\tau$, $A^\rho$ o $S_{n-\sigma}$, $T_{n-\tau}$, $A_{n-\rho}$. De acuerdo con un índice de peso superior o inferior, escribimos también los elementos individuales de una fila con índices superiores o inferiores, como por ejemplo en (4.).
\clearpage
\section*{\S\ 2. Representación mediante productos matriciales.}
En lo que sigue entenderemos siempre por ``producto matricial'' la suma de los productos de determinantes correspondientes de dos matrices con el mismo número de filas, donde las filas de la segunda matriz son contragredientes con las de la primera. Las filas de cada matriz pueden ser: 1. dadas individualmente; 2. reunidas, como en (6.), en una fila $S^\rho$; 3. reunidas en distintas filas $S^\rho,S^\sigma\ldots$. El producto matricial se denotará con el símbolo $(\ \mid\ )$ usado en (6.). La representación (6.) es una aplicación especial de una segunda formulación del teorema sobre las matrices correspondientes: ``A cada producto matricial de $\rho$ filas $(v^{(1)}v^{(2)}\cdots v^{(\rho)}\mid y^{(1)}\cdots y^{(\rho)})$ se le asigna un determinante; y, recíprocamente, a cada determinante se le puede asignar un producto matricial con número prescrito $\rho$ de filas ($\rho=1,2,\ldots,n-1$).''\footnote{Para $\rho>n$ el producto matricial, como se sabe, se anula; para $\rho=n$ se convierte en el producto de dos determinantes.} En efecto, basta reunir las filas $y$ en una fila $p_\rho$ (respectivamente, las filas $v$ en una fila $q_\rho$) y efectuar las sustituciones dadas en (1.) para pasar de un producto matricial dado al determinante, y recíprocamente. Sin embargo, el valor del determinante no está dado por sus elementos individuales, sino que resulta sólo del desarrollo de Laplace. Puesto que determinante y producto matricial se muestran así equivalentes, el teorema sobre la representabilidad simbólica de las formas (el primer teorema fundamental) se expresa también de la manera siguiente:

\emph{Teorema I.} ``Todas las formaciones invariantes pueden representarse simbólicamente por productos matriciales, y recíprocamente todos los productos matriciales son formaciones invariantes.''\footnote{Las formaciones invariantes de formas fundamentales prescritas son, naturalmente, sólo aquellos agregados de productos matriciales que dependen de las filas simbólicas $S^\rho\ldots$.}

La representación mediante productos matriciales permite superar la dificultad esencial de la representación no explícita causada por la representación determinantal.\footnote{Cf. la Introducción.} Por ``representación explícita'' entendemos una representación en la que sólo entran las propias filas $S^\rho,T^\tau,\ldots$, o filas $R^\rho$ derivables de manera única de éstas, pero en la que ya no aparecen las filas componentes individuales $s,t,\ldots$. En \S\ 6 obtendremos, para todas las formaciones invariantes que sólo dependen de las filas simbólicas $S^\rho,T^\tau,\ldots$, una representación explícita en dichas filas simbólicas, y con ello el traslado al dominio $n$-ario de los procesos de generación, del proceso de plegamiento y de los procesos de diferenciación. La posibilidad de la representación explícita es patente a partir del hecho de que sólo a los productos matriciales más simples pueden asignarse precisamente los determinantes que aparecen en la representación determinantal, de modo que todos los demás productos matriciales son combinaciones de distintos determinantes en una sola expresión. Pues si se resuelven las filas simbólicas y de variables en las filas individuales $s,u,x$, entonces, según el primer teorema fundamental en su forma usual, debe surgir necesariamente un agregado de determinantes. La demostración de que, recíprocamente, todos los agregados de determinantes que representan formaciones invariantes de formas fundamentales prescritas pueden reunirse en productos matriciales explícitos se obtiene mediante el segundo teorema fundamental (\S\ 6).

Para obtener, a partir de dos filas simbólicas $S^\sigma,T^\tau$, con ayuda de filas de variables, una formación invariante de la forma fundamental $(S^\sigma\mid p_\sigma)(T^\tau\mid p_\tau)$ que contenga explícitamente todas las filas (el proceso de plegamiento), sólo necesitamos formar un producto matricial de la tercera clase según la definición dada al principio del parágrafo:
\begin{equation}
 \pair{q_{n-\sigma-\lambda}S^\sigma}{T_{n-\tau}p_{\tau-\lambda}},
 \qquad (\lambda=1,2,\ldots,\tau,n-\sigma).
\tag{7}
\end{equation}
Desarrollado, (7.) da:
\begin{equation}
\begin{aligned}
&\pair{v^{(1)}\cdots v^{(n-\sigma-\lambda)}s^{(1)}\cdots s^{(\sigma)}}
       {\alpha^{(1)}\cdots\alpha^{(n-\tau)}y^{(1)}\cdots y^{(\tau-\lambda)}} \\
&\qquad =\sum_{i,k,l}\eps\,q_{k_1\ldots k_{n-\sigma-\lambda}}
 S^{k_{n-\sigma-\lambda+1}\ldots k_{n-\lambda}}
 T_{l_1\ldots l_{n-\tau}}
 p^{l_{n-\tau+1}\ldots l_{n-\lambda}},
\end{aligned}
\tag{8}
\end{equation}
es decir, una expresión que sólo depende de las filas $S,T,q,p$. Aquí $\eps=\pm1$\footnote{Esta notación se conservará en todo lo que sigue.} y la suma debe entenderse de tal modo que los índices de cada fila $S\ldots$ estén en buen orden, y que, para cada combinación $i_1,i_2,\ldots,i_{n-\lambda}$, tanto los $k$ como los $l$ recorran individualmente todas las permutaciones posibles de los números $i$. El número $\lambda$ que entra en (7.) se llamará el defecto del plegamiento en el caso $\sigma>\tau$; la razón de esta última restricción aparecerá en \S\ 6.\footnote{El número $\lambda$ llega siempre hasta el menor de los dos últimos números indicados.}

Los determinantes de las dos matrices que entran en (7.) pueden considerarse como elementos de nuevas filas $Q^{n-\lambda},P_{n-\lambda}$, cuyas filas asociadas son $Q_\lambda,P^\lambda$. Según (8.), estas filas $Q,P$ pueden expresarse por las originales $S$ y $q$, respectivamente $T$ y $p$. Indicamos esto mediante las ecuaciones
\begin{equation}
\begin{aligned}
 q_{n-\sigma-\lambda}S^\sigma&=Q^{n-\lambda}\sim Q_\lambda, &\quad\text{o también}\quad Q_\lambda&=[q_{n-\sigma-\lambda}S^\sigma]_\lambda,\\
 T_{n-\tau}p_{\tau-\lambda}&=P_{n-\lambda}\sim P^\lambda, &\quad\text{o también}\quad P^\lambda&=[T_{n-\tau}p_{\tau-\lambda}]^\lambda.
\end{aligned}
\tag{9}
\end{equation}
Teniendo en cuenta (4.), se tiene entonces, análogamente a (6.), la representación simbólica cuádruple equivalente para el producto matricial (7.):
\begin{equation}
\pair{q_{n-\sigma-\lambda}S^\sigma}{T_{n-\tau}p_{\tau-\lambda}}
=(q_{n-\sigma-\lambda}S^\sigma P^\lambda)=(Q_\lambda T_{n-\tau}p_{\tau-\lambda})=(-1)^{\lambda(n-\lambda)}\pair{P^\lambda}{Q_\lambda}.
\tag{10}
\end{equation}
Sobre (7.) puede efectuarse todavía otra sustitución de nuevas filas poniendo
\begin{equation}
 \pair{q_{n-\sigma-\lambda}S^\sigma}{T_{n-\tau}p_{\tau-\lambda}}
 =\pair{R^{\sigma+\lambda}}{p_{\sigma+\lambda}}(R^{\tau-\lambda}p_{\tau-\lambda}).
\tag{11}
\end{equation}
También aquí (8.) muestra que los productos de los elementos de las filas $R$ pueden expresarse de manera única por las filas $S$ y $T$. Las sustituciones (9.) y (11.) se usarán especialmente cuando con (7.) se efectúe otro plegamiento con otras filas.

La representación matricial nos permite formular invariantemente las relaciones cuadráticas entre los elementos de una fila $p_\rho$ (de las cuales, como se sabe, se siguen todas las demás), y por tanto, según \S\ 1, las relaciones lineales en los coeficientes de la forma fundamental que satisfacen las filas $S^\rho$. Éstas son las identidades
\begin{equation}
 \pair{q_{n-\rho-\lambda}S^\rho}{S_{n-\rho}p_{\rho-\lambda}}=0,
 \qquad (\lambda=1,2,\ldots,\rho,n-\rho),
\tag{12}
\end{equation}
donde las $p$ y las $q$ han de considerarse como cantidades arbitrarias. Veremos en \S\ 5 que las identidades con número de defecto impar $\lambda$ son consecuencias de las de defecto superior. Las relaciones (12.) son consecuencia directa de las relaciones (5.); pues se tiene
\[
 \pair{q_{n-\rho-\lambda}S^\rho}{S_{n-\rho}p_{\rho-\lambda}}
 =(v^{(1)}\cdots v^{(n-\rho-\lambda)}s^{(1)}\cdots s^{(\rho)}\mid \sigma^{(1)}\cdots\sigma^{(n-\rho)}y^{(1)}\cdots y^{(\rho-\lambda)}),
\]
y la aplicación del teorema del producto da, como consecuencia de (5.), un determinante de $(n-\lambda)$ filas que se anula. Las identidades (12.) dicen que, para hallar todos los tipos de plegamientos, basta plegar el producto de formas lineales en cada fila de variables; o también, por \S\ 6, que la forma normalizada $(S^\rho\mid p_\rho)^m$ no posee ninguna formación invariante lineal en los coeficientes.\footnote{Cf. un añadido de Study en las obras de Grassmann, primer tomo, segunda parte, p. 510 (condición de que una cantidad $S^\rho$ sea simple).}

Las condiciones (12.) son también suficientes para caracterizar las filas $p_\rho$. Pues la propiedad de que los elementos individuales de $p_\rho$ sean determinantes de una matriz es una propiedad invariante, y por tanto debe poder expresarse invariantemente; pero, por el Teorema VI, las condiciones (12.) dan la mayor restricción invariante posible para las $p_\rho$, a saber, que la forma lineal formada con la fila $p_\rho$ como coeficientes no posea ninguna formación invariante.

\section*{\S\ 3. Las identidades simbólicas.}
Debemos ahora trasladar el segundo teorema fundamental del método simbólico, es decir, establecer la totalidad de las identidades irreducibles e indecomponibles en las que todas las filas $S^\rho,p_\rho$ se consideran indecomponibles. Estipulamos además lo siguiente. De acuerdo con el significado de las filas simbólicas y de variables, deberán considerarse compuestas aquellas identidades que surgen al especializar la fila de variables $p_\rho$ en $p_{\rho-1}y$ y aplicar una de las identidades del sistema a las expresiones resultantes; por otra parte, las identidades que contienen $k$ filas simbólicas deben contarse como indecomponibles aun cuando surjan de identidades con menos filas simbólicas mediante el proceso anterior. Las filas $S^\rho,p_\rho$ no tienen por qué entrar necesariamente de manera explícita en las identidades; para interpretarlas como cantidades indecomponibles basta mostrar que los productos matriciales que aparecen dependen sólo de estas filas. Sin embargo, mostraremos que en este caso, en parte aplicando la sustitución (11.), siempre puede obtenerse una representación explícita. Las identidades generales se obtienen, mediante el principio de la representación matricial, a partir de las identidades especiales dadas por Pascal, que sólo contienen filas $x$ y $u$ y constituyen una transferencia directa de las dadas por Study para el dominio ternario:\footnote{Para la formulación y la bibliografía del problema, cf. la Introducción.}
\begin{align}
\sum_{i=1}^n(-1)^{i+1}\pair{a^{(i)}}{x}(a^{(1)}\cdots a^{(i-1)}a^{(i+1)}\cdots a^{(n)}u)&=(-1)^{n+1}\pair{u}{x}(a^{(1)}\cdots a^{(n)}),\tag{13}\\
\sum\pm\pair{a^{(1)}}{x^{(1)}}\pair{a^{(2)}}{x^{(2)}}\cdots\pair{a^{(n)}}{x^{(n)}}&=(a^{(1)}\cdots a^{(n)})(x^{(1)}\cdots x^{(n)}),\tag{14}\\
\sum_{i=1}^n(-1)^{i+1}\pair{u}{a^{(i)}}(a^{(1)}\cdots a^{(i-1)}a^{(i+1)}\cdots a^{(n)}x)&=(-1)^{n+1}\pair{u}{x}(a^{(1)}\cdots a^{(n)}).\tag{15}
\end{align}
De (13.) y (15.) surgen otras dos identidades por las sustituciones $(a^{(i)}\mid x)=(a^{(i)}q_{n-1})$ y, respectivamente, $(u\mid a^{(i)})=(p_{n-1}a^{(i)})$; por tanto, según nuestra interpretación, no deben considerarse esencialmente distintas. (14.) representa el teorema del producto para determinantes; mostraremos cómo (13.) y (14.) (y, dualísticamente, (15.) y (14.)) pueden sustituirse por el sistema de teoremas del producto convenientemente formados para matrices. El teorema del producto, formado una vez para matrices de $\rho$ filas y otra para matrices de $(\rho-1)$ filas, dice
\[
\pair{a^{(1)}a^{(2)}\cdots a^{(\rho)}}{xp_{\rho-1}}
=(a^{(1)}a^{(2)}\cdots a^{(\rho)}\mid xy^{(2)}\cdots y^{(\rho)})
=\sum\pm\pair{a^{(1)}}{x}\pair{a^{(2)}}{y^{(2)}}\cdots\pair{a^{(\rho)}}{y^{(\rho)}},
\]
\[
\sum\pm\pair{a^{(2)}}{y^{(2)}}\cdots\pair{a^{(\rho)}}{y^{(\rho)}}
=(a^{(2)}\cdots a^{(\rho)}\mid y^{(2)}\cdots y^{(\rho)})
=(a^{(2)}\cdots a^{(\rho)}p_{\rho-1})=(a^{(2)}\cdots a^{(\rho)}q_{n-\rho+1});
\]
y de aquí se sigue el sistema de teoremas del producto, donde $\rho=2,3,\ldots,n$:
\begin{equation}
\begin{aligned}
\pair{a^{(1)}a^{(2)}\cdots a^{(\rho)}}{xp_{\rho-1}}
 &=\sum_{i=1}^{\rho}(-1)^{i+1}\pair{a^{(i)}}{x}(a^{(1)}\cdots a^{(i-1)}a^{(i+1)}\cdots a^{(\rho)}\mid p_{\rho-1})\\
 &=\sum_{i=1}^{\rho}(-1)^{i+1}\pair{a^{(i)}}{x}(a^{(1)}\cdots a^{(i-1)}a^{(i+1)}\cdots a^{(\rho)}q_{n-\rho+1}).
\end{aligned}
\tag{16}
\end{equation}
Para $\rho=n$, (16.) pasa a (13.); si además especializamos $p_{n-1}=y^{(2)}p_{n-2}$, $p_{n-2}=y^{(3)}p_{n-3}$, etc., y aplicamos las identidades (16.) a las expresiones $(a^{(1)}\cdots a^{(i-1)}a^{(i+1)}\cdots a^{(n)}\mid y^{(2)}p_{n-2})$, $(a^{(1)}\cdots a^{(i-1)}a^{(i+1)}\cdots a^{(h-1)}a^{(h+1)}\cdots a^{(n)}\mid y^{(3)}p_{n-3})$, etc., pasamos de (13.) a (14.). Por tanto, la identidad (14.) es, según nuestra definición, compuesta a partir de las identidades (16.); equivalentemente, el sistema (16.) es equivalente a las identidades (13.) y (14.). El sistema (16.) satisface una de las condiciones requeridas para las identidades generales: contiene la fila $p_{\rho-1}$ como cantidad indecomponible y, al mismo tiempo, en forma explícita. Es el único sistema que satisface esta condición, y sólo ésta. Pues no podemos formar otros determinantes o productos matriciales a partir de las filas $a,x,p$; y entre éstos no pueden existir relaciones independientes de (16.), ya que de lo contrario, al eliminar $(a^{(1)}\cdots a^{(\rho)}\mid xp_{\rho-1})$, surgiría una relación entre $\rho$ ($\rho<n$) expresiones $(a^{(i)}\mid x)(a^{(1)}\cdots a^{(i-1)}a^{(i+1)}a^{(\rho)}q_{n-\rho+1})$, mientras que por (13.) al menos $n+1$ de tales expresiones deben entrar en una relación. De consideraciones análogas se obtiene el sistema equivalente a (15.) y (14.):
\begin{equation}
\begin{aligned}
\pair{u q_{\rho-1}}{a^{(1)}\cdots a^{(\rho)}}
 &=\sum_{i=1}^{\rho}(-1)^{i+1}\pair{u}{a^{(i)}}(q_{\rho-1}\mid a^{(1)}\cdots a^{(i-1)}a^{(i+1)}\cdots a^{(\rho)})\\
 &=\sum_{i=1}^{\rho}(-1)^{i+1}\pair{u}{a^{(i)}}(p_{n-\rho+1}a^{(1)}\cdots a^{(i-1)}a^{(i+1)}\cdots a^{(\rho)}).
\end{aligned}
\tag{17}
\end{equation}
Para pasar de (16.) a identidades entre filas arbitrarias $A^{\rho_1},B^{\rho_2},\ldots,x,p$, observamos que estas identidades deben convertirse en (16.) tan pronto como resolvamos
$A^{\rho_1}=a^{(1)}a^{(2)}\cdots a^{(\rho_1)}$, $B^{\rho_2}=b^{(1)}b^{(2)}\cdots b^{(\rho_2)}$, etc., y que también sólo pueden combinarse en expresiones finales mediante las operaciones determinadas por (16.) tan pronto como las expresiones individuales sean del tipo que entra en (16.). Así obtenemos, si además ponemos
\begin{equation}
 \rho=\rho_1+\rho_2+\cdots+\rho_k\le n,
\tag{18}
\end{equation}
\begin{align*}
&\pair{a^{(1)}\cdots a^{(\rho_1)}b^{(1)}\cdots b^{(\rho_2)}\cdots k^{(1)}\cdots k^{(\rho_k)}}{xp_{\rho-1}}
  =\pair{A^{\rho_1}B^{\rho_2}\cdots K^{\rho_k}}{xp_{\rho-1}} \\
&\quad=\sum_{i=1}^{\rho_1}(-1)^{i+1}\pair{a^{(i)}}{x}(a^{(1)}\cdots a^{(i-1)}a^{(i+1)}\cdots a^{(\rho_1)}B^{\rho_2}\cdots K^{\rho_k}q_{n-\rho+1})\\
&\qquad+(-1)^{\rho_1\rho_2}\sum_{i=1}^{\rho_2}(-1)^{i+1}\pair{b^{(i)}}{x}(b^{(1)}\cdots b^{(i-1)}b^{(i+1)}\cdots b^{(\rho_2)}A^{\rho_1}\cdots K^{\rho_k}q_{n-\rho+1})\\
&\qquad+\cdots\\
&\qquad+(-1)^{(\rho_1+\cdots+\rho_{k-1})\rho_k}\sum_{i=1}^{\rho_k}(-1)^{i+1}\pair{k^{(i)}}{x}(k^{(1)}\cdots k^{(i-1)}k^{(i+1)}\cdots k^{(\rho_k)}A^{\rho_1}\cdots K^{\rho_{k-1}}q_{n-\rho+1}).
\end{align*}
Cada suma individual (pero todavía no una parte de la suma) contiene realmente las filas $A^{\rho_1},B^{\rho_2},\ldots$ como cantidades indecomponibles; pues satisface las condiciones necesarias y suficientes\footnote{Capelli, a. a. O. \S\ 23, da condiciones suficientes: $\bigl(a^{(k)}\mid \frac{\partial}{\partial a^{(i)}}\bigr)=0$, $(k>i,\ i=1,2,\ldots,\rho_1-1)$. Éstas implican del modo más simple las condiciones necesarias y suficientes por aplicación del proceso de corchetes de Poisson según Wellstein, Von den Differentialgleichungen der projektiven Invarianten, Math. Ann. 67 (1909) \S\ 3.}
\begin{equation}
 \left(a^{(i+1)}\middle|\frac{\partial}{\partial a^{(i)}}\right)=0;\ \ldots\ ;
 \left(k^{(j+1)}\middle|\frac{\partial}{\partial k^{(j)}}\right)=0.
 \quad (i=1,2,\ldots,\rho_1-1;\ j=1,2,\ldots,\rho_k-1)
\tag{19}
\end{equation}
Mostramos que también contiene explícitamente todas las filas. En efecto, si ponemos $B^{\rho_2}\cdots K^{\rho_k}q_{n-\rho+1}=q_{n-\rho_1+1}$, entonces por (16.) y (10.) obtenemos
\begin{align*}
\sum(-1)^{i+1}\pair{a^{(i)}}{x}(a^{(1)}\cdots a^{(i-1)}a^{(i+1)}\cdots a^{(\rho_1)}q_{n-\rho+1})
&=\pair{A^{\rho_1}}{xp_{\rho_1-1}}\\
&=(A_{n-\rho_1}xp_{\rho_1-1})=(-1)^{(\rho_1+1)(n+1)}\pair{q_{n-\rho_1+1}}{A_{n-\rho_1}x}\\
&=(-1)^{(\rho_1+1)(n+1)}\pair{B^{\rho_2}\cdots K^{\rho_k}q_{n-\rho+1}}{A_{n-\rho_1}x}.
\end{align*}
Si calculamos las restantes sumas de modo análogo y luego ponemos, puesto que las filas ya están caracterizadas como distintas por los índices de peso, $B^{\rho_2}=A^{\rho_2}$, $C^{\rho_3}=A^{\rho_3},\ldots$, obtenemos las identidades deseadas:
\begin{equation}
\begin{aligned}
\pair{A^{\rho_1}A^{\rho_2}\cdots A^{\rho_k}}{xp_{\rho-1}}
&=\sum_{i=1}^{k}(-1)^{\delta_i}\pair{A^{\rho_1}\cdots A^{\rho_{i-1}}A^{\rho_{i+1}}\cdots A^{\rho_k}q_{n-\rho+1}}{A_{n-\rho_i}x},\\
\delta_i&=(\rho_1+\rho_2+\cdots+\rho_{i-1})\rho_i+(\rho_i+1)(n+1),
\end{aligned}
\tag{20}
\end{equation}
y de (17.) las identidades duales:
\begin{equation}
\begin{aligned}
\pair{u q_{\sigma-1}}{A_{\sigma_1}A_{\sigma_2}\cdots A_{\sigma_k}}
&=\sum_{i=1}^{k}(-1)^{\eps_i}\pair{u A^{n-\sigma_i}}{p_{n-\sigma+1}A_{\sigma_1}\cdots A_{\sigma_{i-1}}A_{\sigma_{i+1}}\cdots A_{\sigma_k}},\\
\eps_i&=(\sigma_1+\sigma_2+\cdots+\sigma_{i-1})\sigma_i+(\sigma_i+1)(n+\sigma).
\end{aligned}
\tag{21}
\end{equation}
\emph{Teorema II.} Con (20.) y (21.) queda agotada la totalidad de las identidades indecomponibles entre las filas $A^{\rho_i},x,p$, respectivamente $A_{\sigma_i},u,q$, y al mismo tiempo la totalidad de aquellas identidades indecomponibles que admiten representación explícita sin efectuar la sustitución (11.).

La primera parte del teorema se sigue de las consideraciones que conducen a (20.) y (21.); la segunda, de la imposibilidad de una identidad entre dos filas simbólicas:
\[
\pair{q_\rho A^\sigma}{B_{\rho+\sigma-\tau}p_\tau}
=\eps_1\pair{q_\rho B^{n-\rho-\sigma+\tau}}{A_{n-\sigma}p_\tau}
+\eps_2\pair{q_\rho q_{n-\tau}}{B_{\rho+\sigma-\tau}A_{n-\sigma}},
\]
donde $\rho\le\tau\le\sigma$ y ninguna de las filas tiene peso $1\cdot(n-1)$. (Otras hipótesis sobre $\rho,\tau,\sigma$ pueden reducirse a ésta permutando las designaciones de las filas.) Esta imposibilidad se sigue, por ejemplo, desarrollando (8.) bajo la especialización $q_\rho=lM^{\rho-1}$, $q_{n-\tau}=lN^{n-\tau-1}$, si se pone $l_1=1$, $M^{2,3,\ldots,\rho}=1$, $A^{\rho+1,\ldots,\rho+\sigma}=1$, todos los demás elementos de las filas $l,M,A$ iguales a cero, mientras que $N$ y $B$ son arbitrarios. Puesto que las identidades entre filas arbitrarias, al resolver una fila $p_\tau$ en $y^{(1)}y^{(2)}\cdots y^{(\tau)}$, pasan a identidades compuestas a partir de (20.), la segunda parte de la afirmación queda probada tan pronto como (20.) pueda generarse a partir de identidades con sólo dos filas simbólicas. De hecho, a partir de
\begin{equation}
\pair{A^{\rho_1}A^{\rho_2}}{xp_{\rho-1}}
=\eps\pair{A^{\rho_1}q_{n-\rho+1}}{A_{n-\rho_2}x}
+\pair{A^{\rho_1}}{xp_{\rho_1-1}},
\quad \text{donde }p_{\rho_1-1}\sim A^{\rho_2}q_{n-\rho+1},
\tag{20a}
\end{equation}
por especialización $A^{\rho_1}=B^{\sigma_1}B^{\sigma_2}$, es decir, por
\[
\pair{B^{\sigma_1}B^{\sigma_2}}{xp_{\rho_1-1}}
=\eps\pair{B^{\sigma_1}q_{n-\rho_1+1}}{B_{n-\sigma_2}x}
+\pair{B^{\sigma_1}}{xp_{\sigma_1-1}},
\quad \text{donde } p_{\sigma_1-1}\sim B^{\sigma_2}q_{n-\rho_1+1},
\]
se obtiene una identidad (20.) con tres filas simbólicas; y por nuevas especializaciones de $B^{\sigma_1}=C^{\tau_1}C^{\tau_2}$, etc., las identidades generales (20.). La inexistencia de una representación explícita de una identidad entre dos filas simbólicas y dos filas de variables arbitrarias implica así la inexistencia en el caso de un número arbitrario de filas simbólicas, siempre que entre las filas de variables no esté contenida ninguna fila $x$ o $u$.

\section*{\S\ 4. Las identidades de descomposición.}
Consideramos ahora el caso más general de identidades entre filas simbólicas y de variables arbitrarias, en cuyas expresiones finales, sin efectuar la sustitución (11.), aparece una fila $p_\tau$ resuelta en sus filas individuales $y^{(1)}\cdots y^{(\tau)}$; llamamos a estas identidades ``identidades de descomposición''. Siguiendo la observación del final del parágrafo precedente, primero las determinamos para el caso de sólo dos filas simbólicas, es decir, desarrollamos la expresión $\pair{q_\rho A^\sigma}{B_{\rho+\sigma-\tau}p_\tau}$. Ponemos:
\begin{equation}
\begin{aligned}
q_{\rho-\alpha}^{(i_1\ldots i_\alpha)}&\sim p_{n-\rho}\,y^{(i_1)}\cdots y^{(i_\alpha)},\qquad
p_{\tau-\alpha}^{(i_1\ldots i_\alpha)}=y^{(i_{\alpha+1})}\cdots y^{(i_\tau)},\qquad
p_\tau=y^{(1)}\cdots y^{(\tau)}.
\end{aligned}
\tag{22}
\end{equation}
Respecto de los números $\rho,\sigma,\tau$ deben distinguirse cuatro casos:
\[
\begin{array}{llll}
1)&\rho+\sigma\le n,& \tau\le\rho,& \tau\le\sigma;\\[2pt]
2)&\rho+\sigma\le n,& \tau\ge\rho,& \tau\le\sigma;\\[2pt]
3)&\rho+\sigma\le n,& \tau\le\rho,& \tau>\sigma;\\[2pt]
4)&\rho+\sigma\le n,& \tau\ge\rho,& \tau\ge\sigma.
\end{array}
\]
Destacaremos el primer caso y luego caracterizaremos brevemente los otros a partir de él. Por (20.), (10.) y (9.) obtenemos, cuando descomponemos la fila $y^{(1)}y^{(2)}\cdots y^{(\tau)}B_{\rho+\sigma-\tau}$ en $y^{(1)}$ y $y^{(2)}\cdots y^{(\tau)}B_{\rho+\sigma-\tau}$:
\begin{equation}
\begin{aligned}
\pair{q_\rho A^\sigma}{y^{(1)}y^{(2)}\cdots y^{(\tau)}B_{\rho+\sigma-\tau}}
&=\pair{q_\rho^{(1)}A^\sigma}{y^{(2)}\cdots y^{(\tau)}B_{\rho+\sigma-\tau}}\\
&\quad+(-1)^\rho\pair{q_\rho[A_{n-\sigma}y^{(1)}]^{\sigma-1}}{y^{(2)}\cdots y^{(\tau)}B_{\rho+\sigma-\tau}}.
\end{aligned}
\tag{23}
\end{equation}
El último término de (23.) tiene ahora exactamente la forma de la expresión original --sólo que $\sigma$ se sustituye por $\sigma-1$, y $\tau$ por $\tau-1$-- y por tanto admite de nuevo la ecuación (23.). Así, puesto que $\tau\le\sigma$, podemos aplicar $\tau$ veces la operación (23.) y, teniendo en cuenta (10.), obtenemos la relación:
\begin{equation}
\begin{aligned}
\pair{q_\rho A^\sigma}{p_\tau B_{\rho+\sigma-\tau}}
&=(-1)^{\rho\sigma+(\sigma+\tau)(n+1)}
  \pair{q_\rho B^{n-\rho-\sigma+\tau}}{A_{n-\sigma}p_\tau}\\
&\quad+\sum_{i_1=1}^{\tau}(-1)^{\rho(i_1-1)}
 \pair{q_{\rho-1}^{(i_1)}[A_{n-\sigma}y^{(1)}\cdots y^{(i_1-1)}]^{\sigma-i_1+1}}
 {y^{(i_1+1)}\cdots y^{(\tau)}B_{\rho+\sigma-\tau}}.
\end{aligned}
\tag{24}
\end{equation}
Cada término bajo el signo de suma vuelve a ser del tipo de la expresión original; $\rho$ se sustituye por $\rho-1$, $\sigma$ por $\sigma-i_1+1$ y $\tau$ por $\tau-i_1$, de modo que la condición 1) queda satisfecha aún más fuertemente. Como $\tau\le\rho$, podemos por tanto aplicar $\tau$ veces la operación (24.) y obtener la identidad de descomposición buscada:
\begin{equation}
\begin{aligned}
\pair{q_\rho A^\sigma}{p_\tau B_{\rho+\sigma-\tau}}
&=\eps_{i_0}\pair{q_\rho B^{n-\rho-\sigma+\tau}}{A_{n-\sigma}p_\tau}\\
&\quad+\sum_{i_1=1}^{\tau}\eps_{i_1}
 \pair{q_{\rho-1}^{(i_1)}B^{n-\rho-\sigma+\tau}}{A_{n-\sigma}p_{\tau-1}^{(i_1)}}\\
&\quad+\sum_{i_2=i_1+1}^{\tau}\eps_{i_2}
 \pair{q_{\rho-2}^{(i_1i_2)}B^{n-\rho-\sigma+\tau}}{A_{n-\sigma}p_{\tau-2}^{(i_1i_2)}}+\cdots\\
&\quad+\eps_{i_\tau}\pair{q_{\rho-\tau}^{(i_1\ldots i_\tau)}B^{n-\rho-\sigma+\tau}}{A_{n-\sigma}}\\
&=\sum_{\alpha=0}^{\tau}\sum_{i_\alpha=i_{\alpha-1}+1}^{\tau}\eps_{i_\alpha}
 \pair{q_{\rho-\alpha}^{(i_1\ldots i_\alpha)}B^{n-\rho-\sigma+\tau}}{A_{n-\sigma}p_{\tau-\alpha}^{(i_1\ldots i_\alpha)}}.
\end{aligned}
\tag{25}
\end{equation}
Para el factor de signo $\eps$ se obtiene de (24.):
\begin{equation}
\begin{aligned}
\eps_{i_\alpha}&=(-1)^{\delta_{i_\alpha}},\qquad
\delta_{i_\alpha}=\sum_{\beta=1}^{\alpha}(\rho-\beta+1)(i_{\beta-1}+i_\beta+1)\\
&\quad +(\rho-\alpha)(\sigma+i_\alpha+\alpha)+(\sigma+\tau+\alpha)(n+1),
\end{aligned}
\tag{26}
\end{equation}
donde debe ponerse $i_0=0$. El signo de suma en (25.) debe entenderse así: a cada combinación $i_1i_2\ldots i_{\alpha-1}$, donde $i_1<i_2<\cdots<i_{\alpha-1}$, se adjuntan todos los valores $i_\alpha$ para los que $i_{\alpha-1}<i_\alpha\le\tau$. Teniendo en cuenta (26.), se verifica fácilmente que las sumas individuales en (25.) satisfacen las ecuaciones diferenciales análogas a (19.):
\begin{equation}
 \left(\frac{\partial}{\partial y^{(i)}}\middle|y^{(i+1)}\right)=0,
 \qquad (i=1,2,\ldots,\tau-1),
\tag{27}
\end{equation}
y por tanto dependen sólo de la fila $p_\tau$. Así, la identidad de descomposición es una identidad indecomponible entre las filas $A,B,q,p$; volveremos a su representación explícita en el parágrafo siguiente.

En el caso 2), la operación (23.) también puede aplicarse $\tau$ veces, pero la operación (24.) se detiene después de la $\rho$-ésima aplicación. Así, en la expresión final de (25.) el primer signo de suma debe reemplazarse por $\sum_{\alpha=0}^{\rho}$. En los casos 3) y 4), la operación (23.) sólo puede realizarse $\sigma$ veces; en lugar de (24.) se tiene entonces la relación:
\begin{equation}
\begin{aligned}
\pair{q_\rho A^\sigma}{p_\tau B_{\rho+\sigma-\tau}}
&=(-1)^{\rho\sigma}
 \pair{q_\rho^{(\sigma+1\ldots\tau)}B^{n-\rho+\tau-\sigma}}{A_{n-\sigma}p_{\tau-\sigma}^{(\sigma+1\ldots\tau)}}\\
&\quad+\sum_{i_1=1}^{\sigma}(-1)^{\rho(i_1-1)}
 \pair{q_{\rho-1}^{(i_1)}[A_{n-\sigma}y^{(1)}\cdots y^{(i_1-1)}]^{\sigma-i_1+1}}
 {y^{(i_1+1)}\cdots y^{(\tau)}B_{\rho+\sigma-\tau}};
\end{aligned}
\tag{28}
\end{equation}
al repetir (28.) $\tau-\sigma$ veces, donde el signo de suma se convierte en $\sum_{i_\beta=i_{\beta-1}+1}^{\sigma+\beta-1}$ --como muestra la aplicación $(\sigma-i_1+1)$ veces de (23.) a los términos bajo el signo de suma en (28.)-- se obtienen todos los términos de la suma individual de (25.) correspondientes al índice $\alpha=\tau-\sigma$, con el signo correspondiente. Entonces los casos 3) y 4) pasan a los casos 1) y 2), pues los números correspondientes a $\sigma,\tau$ se convierten en $\sigma'=\sigma-i_{\tau-\sigma}+\tau-\sigma$, $\tau'=\tau-i_{\tau-\sigma}=\sigma'$; y a medida que continúa el procedimiento, $\tau'<\sigma'$. Así, en la expresión final de (25.) el primer signo de suma debe reemplazarse por $\sum_{\alpha=\tau-\sigma}^{\tau,\rho}$.

Análogamente, de (21.) se obtiene la identidad de descomposición dual, en la que una fila $q_\rho$ se resuelve en las filas individuales $v^{(1)},\ldots,v^{(\rho)}$. Si ponemos
\begin{equation}
 \dot p_{\tau-\beta}^{(i_1\ldots i_\beta)}\sim v^{(i_1)}\cdots v^{(i_\beta)}q_{n-\tau};\qquad
 \dot q_{\rho-\beta}^{(i_1\ldots i_\beta)}=v^{(i_{\beta+1})}\cdots v^{(i_\rho)};\qquad
 q_\rho=v^{(1)}\cdots v^{(\rho)},
\tag{29}
\end{equation}
entonces
\begin{equation}
 \pair{q_\rho A^\sigma}{p_\tau B_{\rho+\sigma-\tau}}
 =\sum_{\beta=\max(0,\tau-\sigma)}^{\min(\tau,\rho)}
   \sum_{i_\beta=i_{\beta-1}+1}^{\rho}\eps
\pair{\dot q_{\rho-\beta}^{(i_1\ldots i_\beta)}B^{n-\rho-\sigma+\tau}}{A_{n-\sigma}\dot p_{\tau-\beta}^{(i_1\ldots i_\beta)}} ,
\tag{30}
\end{equation}
donde el índice de suma $\beta$ recorre desde el mayor de los valores inferiores hasta el menor de los valores superiores. Todos los términos de una suma individual en (25.) y (30.) son polares de una sola forma, producidos por los procesos polares
\[
\left(\frac{\partial}{\partial p_{\tau-\alpha}}\middle|p_{\tau-\alpha}^{(i_1\ldots i_\alpha)}\right),\qquad
\left(q_{\rho-\alpha}^{(i_1\ldots i_\alpha)}\middle|\frac{\partial}{\partial q_{\rho-\alpha}}\right).
\]
Para expresar esto, denotamos por $\Pi$ un agregado de tales operaciones polares, multiplicado por constantes, y obtenemos, en lugar de (25.) y (30.), la relación
\begin{equation}
 \pair{q_\rho A^\sigma}{p_\tau B_{\rho+\sigma-\tau}}
 =\sum_{\alpha=\max(0,\tau-\sigma)}^{\min(\tau,\rho)}
 \Pi\pair{q_{\rho-\alpha}B^{n-\rho-\sigma+\tau}}{A_{n-\sigma}p_{\tau-\alpha}}.
\tag{31}
\end{equation}

\section*{\S\ 5. Las identidades de descomposición en forma explícita.}
Si en (25), en el caso 4), especializamos $\tau$ a $\sigma+\rho$, obtenemos
\begin{equation}
 \pair{q_\rho A^\sigma}{y^{(1)}\cdots y^{(\rho+\sigma)}}
 =\sum_{1\le i_1<\cdots<i_\rho\le \rho+\sigma}
 \eps_{i_\rho}\pair{q_\rho}{y^{(i_1)}\cdots y^{(i_\rho)}}
 \pair{A^\sigma}{y^{(i_{\rho+1})}\cdots y^{(i_{\rho+\sigma})}},
\tag{32}
\end{equation}
donde $\eps_{i_\rho}=(-1)^{\delta_{i_\rho}}$ y $\delta_{i_\rho}=\rho(\rho+1)/2+i_1+\cdots+i_\rho$. Ésta es una forma más general del teorema del producto para matrices que (16.) y (17.), y se sigue por desarrollo de Laplace del determinante de la matriz producto
\[
 \sum\pm(v^{(1)}y^{(1)})\cdots(v^{(\rho)}y^{(\rho)})(a^{(1)}y^{(\rho+1)})\cdots(a^{(\sigma)}y^{(\rho+\sigma)}).
\]
Así, el teorema del producto más general se compone de las identidades (20.) y (21.), respectivamente, y con ello se explica la elección de la forma especial (16.), (17.).

La relación (32.) proporciona ahora el medio para representar explícitamente las identidades de descomposición. Para ello consideramos las formas que aparecen en (31.) como formas fundamentales; es decir, efectuamos la sustitución (11.), que no cambia nada en la representación explícita en las filas $A,B$, puesto que por (8.) las nuevas filas $R$ introducidas pueden expresarse mediante las propias filas $A,B$ y no mediante sus filas componentes $a^{(1)}a^{(2)}\ldots$, $b^{(1)}b^{(2)}\ldots$. Así, sea
\begin{equation}
 \pair{q_{\rho-\alpha}B^{n-\rho-\sigma+\tau}}{A_{n-\sigma}p_{\tau-\alpha}}
 =\pair{R_{\rho-\alpha}p_{n-\rho+\alpha}}{R^{\tau-\alpha}p_{\tau-\alpha}}.
\tag{33}
\end{equation}
Entonces, para la suma individual de (25.) correspondiente al índice $\alpha$, obtenemos
\begin{equation}
\begin{aligned}
&\sum \eps_{i_\alpha}\pair{R_{\rho-\alpha}p_{n-\rho}y^{(i_1)}\cdots y^{(i_\alpha)}}{R^{\tau-\alpha}y^{(i_{\alpha+1})}\cdots y^{(i_\tau)}}\\
&\quad=\sum \eps_{i_\alpha}\pair{[R_{\rho-\alpha}p_{n-\rho}]^\alpha y^{(i_1)}\cdots y^{(i_\alpha)}}{R^{\tau-\alpha}y^{(i_{\alpha+1})}\cdots y^{(i_\tau)}}\\
&\quad=\eps\pair{[R_{\rho-\alpha}p_{n-\rho}]^\alpha R^{\tau-\alpha}}{p_\tau}
 =\eps\pair{R^{\tau-\alpha}q_{n-\tau}}{R_{\rho-\alpha}p_{n-\rho}}.
\end{aligned}
\tag{34}
\end{equation}
Así se obtiene efectivamente una expresión final explícita, y la forma simétrica en que aparecen $q_\rho$ y $p_\tau$ muestra que (30.) conduce a la misma expresión final, que por tanto es independiente de la descomposición inicial. Entre (25.) y (30.) sólo hay diferencia mientras las filas $y$ y $v$ conservan significado independiente, de modo que la identidad de descomposición pasa, según nuestra definición, a ser una identidad descomponible. Para llegar a identidades con más filas simbólicas, basta, según \S\ 3 (final), resolver una fila $q_\rho$ en $q_{\rho_1}C^{\sigma_1}$, o $p_\tau$ en $p_{\tau_1}C_{\tau-\tau_1}$. Las expresiones resultantes
\begin{equation}
 \pair{q_{\rho_1}C^{\sigma_1}}{[R^{\tau-\alpha}q_{n-\tau}]_\lambda R_{\rho+\sigma_1-\alpha-\lambda}},
 \qquad
 \pair{[R_{\rho-\alpha}p_{n-\rho}]^\lambda R^{\tau-\alpha}}{p_{\tau_1}C_{\tau-\tau_1}}
\tag{35}
\end{equation}
son exactamente del tipo que aparece para dos filas simbólicas y, por tanto, admiten nuevamente representación explícita después de una sustitución correspondiente a (33.). La repetición proporciona así representación explícita para un número arbitrario de filas. También puede notarse que la primera y la última suma en (25.) y (30.), respectivamente, dan la representación explícita sin aplicar (33), sólo mediante la fórmula (32), o bien se reducen por completo a un solo término que contiene todas las filas explícitamente. La relación obtenida de (32.) al resolver $q_\rho$ en $q_{\rho_1}C^{\sigma_1}$ y así sucesivamente puede considerarse como el teorema del producto en su forma más general.

En resumen mostraremos:

\emph{Teorema III:} Con las identidades
\begin{equation}
 \pair{q_\rho A^\sigma}{p_\tau B_{\rho+\sigma-\tau}}
 =\sum_{\alpha=\max(0,\tau-\sigma)}^{\min(\tau,\rho)}
 \eps\pair{R^{\tau-\alpha}q_{n-\tau}}{R_{\rho-\alpha}p_{n-\rho}},
\tag{36}
\end{equation}
donde las filas $R$ están determinadas por (33), junto con las que se obtienen de ellas al resolver $q_\rho,p_\tau$ en $q_{\rho_1}C^{\sigma_1}$, etc., queda agotada la totalidad de las identidades indecomponibles.

En efecto, las identidades así obtenidas son las más generales de su especie, puesto que las filas individuales ya no están sometidas a ninguna restricción de peso o número, como aún ocurría en (20.) y (21.). Tampoco existen otras identidades entre estas filas; pues cuando una fila $p_\tau$ se resuelve en $y^{(1)}\cdots y^{(\tau)}$, las identidades más generales se componen de (20.), y por tanto, por \S\ 3 (final), de (20a.); y la composición discutida después de (16.) es precisamente la efectuada en \S\ 4. El paso a un número arbitrario de filas lo proporciona, como muestra la discusión de (20a.), aplicar siempre la misma operación a las expresiones que surgen de resolver $q_\rho$ en $q_{\rho_1}C^{\sigma_1}$, y así sucesivamente. También se sigue que el paso directo desde (20.) en lugar de desde (20a.) no aporta nada nuevo. Esto se verifica fácilmente por cálculo, especializando una vez una fila $q_\rho$ en (25.), y otra vez efectuando el paso desde (20.) por el método de \S\ 4. En ambos casos surgen las mismas sumas, que por aplicación del teorema general del producto (véase arriba) pasan a las identidades obtenidas de (36.). Las identidades (20.), (21.) corresponden a los casos especiales $\tau=1$, $\rho=1$; entonces sólo aparecen dos sumas individuales y por tanto, de acuerdo con una observación anterior, una representación explícita sin la sustitución (33.), coincidente con lo encontrado antes.

Finalmente mencionemos brevemente dos casos especiales de la identidad de descomposición. Póngase en (31): $\tau=\sigma$, $\rho=n-\sigma$, y denótese la fila $p_\tau$ por $p'_\sigma\sim q'_{n-\sigma}$. Entonces
\begin{equation}
\begin{aligned}
\pair{A^\sigma}{p_\sigma}\pair{B^\sigma}{p'_\sigma}
-\pair{A^\sigma}{p'_\sigma}\pair{B^\sigma}{p_\sigma}
=\sum_{\alpha=1}^{\min(\sigma,n-\sigma)}
\Pi\pair{q_{n-\sigma-\alpha}B^\sigma}{A_{n-\sigma}p_{\sigma-\alpha}}.
\end{aligned}
\tag{37}
\end{equation}
Ésta es la fórmula de Clebsch\footnote{Clebsch, loc. cit., pp. 25--32. Clebsch demuestra que los términos individuales de los lados derechos dependen sólo de las filas $A,B$; la representación explícita se obtendría aplicando (32.) a sus expresiones $\Phi_h$.} para desarrollar una forma con dos filas cogredientes $p_\sigma,p'_\sigma$ en covariantes elementales. Los covariantes elementales pertenecientes a una forma $\pair{A^\sigma}{p_\sigma}^m\pair{B^\sigma}{p'_\sigma}^n$ son entonces los siguientes:
\begin{equation}
 \left\{\sum_{\alpha=1}^{\min(\sigma,n-\sigma)}
 \eps\pair{q_{n-\sigma-\alpha}B^\sigma}{A_{n-\sigma}p_{\sigma-\alpha}}
 \right\}^{\rho}
 \pair{A^\sigma}{p_\sigma}^{m-\rho}\pair{B^\sigma}{p'_\sigma}^{n-\rho},
 \qquad (\rho=0,1,\ldots,m,n).
\tag{38}
\end{equation}
Proceden, como diremos brevemente, de la forma fundamental por ``plegamiento cogrediente de defecto $\rho$'' (cf. \S\ 2).

En segundo lugar, póngase $\tau=\sigma-\lambda$, $\rho=n-\sigma-\lambda$, $B^\sigma=A^\sigma$. Entonces
\begin{equation}
\begin{aligned}
\pair{q_{n-\sigma-\lambda}A^\sigma}{A_{n-\sigma}p_{\sigma-\lambda}}
&=(-1)^\lambda\pair{q_{n-\sigma-\lambda}A^\sigma}{A_{n-\sigma}p_{\sigma-\lambda}}\\
&\quad+(-1)^{(n-\sigma)(\sigma-\lambda)}
 \sum_{\alpha=1}^{\min(\sigma-\lambda,n-\sigma-\lambda)}
 \Pi\pair{q_{n-\sigma-\lambda-\alpha}A^\sigma}{A_{n-\sigma}p_{\sigma-\lambda}}.
\end{aligned}
\tag{39}
\end{equation}
Así, como se observó en \S\ 2, las condiciones (12.) que caracterizan las filas simbólicas para defecto impar $\lambda$ son consecuencias de las condiciones de defecto superior. Para una forma lineal $\pair{A^\sigma}{p_\sigma}=\pair{B^\sigma}{p_\sigma}$, (39.) implica que sus formaciones invariantes irreducibles cuadráticas en los coeficientes se reducen a las de defecto par; esto concuerda con el hecho conocido de que una forma lineal en los dominios binario y ternario no tiene formaciones invariantes.

Debe observarse que las identidades generales se derivaron originariamente bajo la hipótesis de que las filas individuales satisfacen las condiciones (12.). Sin embargo, puesto que estas filas individuales entran en las identidades sólo linealmente, éstas valen también para las filas más generales que no satisfacen las condiciones (12.).

\section*{\S\ 6. Representación explícita de las formaciones invariantes.\\Procesos de plegamiento y de diferenciación.}
Los resultados de los parágrafos precedentes nos permiten completar el teorema enunciado en \S\ 2 sobre la representabilidad simbólica (el primer teorema fundamental) añadiendo la representabilidad explícita; a saber, mostramos:

\emph{Teorema IV:} ``Todas las formaciones invariantes de formas fundamentales arbitrarias pueden representarse explícitamente por productos matriciales; es decir, aparte de las filas de variables $p_\rho$, sólo entran en las expresiones finales las filas $S^\sigma,\ldots,T^\tau$ de las formas originales, o las filas $P,Q,R$ derivadas de ellas por sustitución (9.) o (11.).''

Para la prueba, supongamos que una formación invariante depende, además de filas de variables, de las filas $S^\sigma,\ldots,T^\tau$; y supongamos que esas filas se han resuelto suficientemente en filas individuales $S^\sigma=S^{\sigma_1}\cdots S^{\sigma_i},\ldots,T^\tau=T^{\tau_1}\cdots T^{\tau_k}$ para hacer posible la representación por determinantes. Ordenemos las filas en una sucesión $S^\sigma,\ldots,T^\tau$ arbitraria en sí misma pero fijada de una vez por todas. Extraemos un agregado de determinantes que depende de la primera fila total $S^\sigma=S^{\sigma_1}\cdots S^{\sigma_i}$, no sólo de filas componentes individuales. Esta dependencia, sin embargo, es por \S\ 3--\S\ 5 una consecuencia de las identidades generales. Si resolvemos ahora una fila de variables $p_\rho$ en las filas $y,v$, o una de las filas posteriores $T$ en las filas individuales $t$, respectivamente $\tau$, entonces las identidades más generales se componen de (20.) y (21.). Pero estas últimas identidades conducen siempre a una expresión que contiene explícitamente la fila $S^\sigma=S^{\sigma_1}\cdots S^{\sigma_i}$, a saber, el lado izquierdo de (20.) y (21.); pues por (19.) se verifica que sólo la suma completa de la derecha --no una parte de ella correspondiente a un término de (20a.)-- posee la propiedad de depender sólo de $S^\sigma$. A medida que continúa el procedimiento, todas las filas que han entrado una vez explícitamente permanecen explícitas; la aplicación de las identidades puede requerir a lo sumo la reunión de varias filas en una sola, es decir, la sustitución (9.). Si finalmente sólo queda por llevar a forma explícita una única fila $T^\tau$, hay dos casos posibles. Puede aparecer todavía una fila de variables libre, esto es, no conectada con otras filas mediante la sustitución (9.); resolverla en filas componentes $y$ o $v$ completa el procedimiento y da, como en (31.), polares de una o más formas que contienen explícitamente todas las filas. Si no queda ninguna fila de variables libre, entonces por el modo en que surgieron reconocemos, en la totalidad de las expresiones que contienen las filas individuales $t$ o $\tau$ pero dependen de la fila $T^\tau$, los términos de una identidad de descomposición; por \S\ 5, sin embargo, éstas siempre admiten representación explícita en todas las filas después de aplicar la sustitución (11.). Esto prueba el teorema en toda generalidad. También debe señalarse que cuando sólo aparecen dos filas simbólicas, la sustitución (11.) nunca aparece. Pues, dado que $\sigma,n-\sigma,\tau,n-\tau<n$, sólo pueden faltar filas de variables cuando las dos filas simbólicas se unen en un único determinante y por tanto aparecen explícitamente; pero tan pronto como entran filas de variables, (34.) y (33.) muestran que la sustitución (11.) se vuelve innecesaria.

De lo recién dicho se sigue que, para generar todas las formaciones invariantes, basta unir las filas simbólicas, tras adjuntar filas de variables, en todos los productos matriciales posibles. El producto matricial más general tiene ahora la forma
\begin{equation}
 \pair{P^{\mu_1}\cdots P^{\mu_i}}{Q_{\nu_1}\cdots Q_{\nu_k}},\qquad \sum\mu=\sum\nu,
\tag{40}
\end{equation}
donde $P$ y $Q$ pueden ser filas de las formas originales, o pueden haber surgido de las filas originales por una sucesión de sustituciones (9.) o (11.), o, finalmente, pueden ser filas de variables. Pero las expresiones (40.) pueden generarse uniendo sucesivamente dos filas simbólicas cada vez en un producto matricial, con ayuda de filas de variables; pues de
\[
 \pair{q_{n-\mu-\lambda}P^\mu}{Q_{\nu_1}p_{n-\mu-\nu_1-\lambda}}=\pair{R_\lambda Q_{\nu_1}}{p_{n-\mu-\nu_1-\lambda}}
\]
se obtiene, al unir con la fila $Q_{\nu_2}$, una expresión
\[
 \pair{[R_\lambda Q_{\nu_1}]^{n-\lambda-\nu_1}}{Q_{\nu_2}p_{n-\mu-\nu_1-\nu_2-\lambda}}
 =\pair{q_{n-\mu-\lambda}P^\mu}{Q_{\nu_1}Q_{\nu_2}p_{n-\mu-\nu_1-\nu_2-\lambda}},
\]
y por consiguiente, continuando el procedimiento, una expresión (40.). Si $P$ y $Q$ no son filas de las formas originales, entonces (9.) y (11.), junto con lo que acaba de mostrarse, implican que también ellas surgieron por una sucesión de uniones de dos filas simbólicas cada vez. Así:

\emph{Teorema V:} El proceso para generar todas las formaciones invariantes, el proceso de plegamiento, consiste en unir sucesivamente dos filas simbólicas cada vez en un producto matricial, con ayuda de filas de variables.

A partir de dos filas simbólicas $S^\sigma,T^\tau$, donde $\sigma\ge\tau$, pueden formarse los siguientes productos matriciales:
\begin{align}
&\pair{q_{n-\sigma-\lambda}S^\sigma}{T_{n-\tau}p_{\tau-\lambda}}, &&(\lambda=1,2,\ldots,\min(\tau,n-\sigma)),\tag{41}\\
&\pair{q_{n-\tau-\mu}T^\tau}{S_{n-\sigma}p_{\sigma-\mu}}, &&(\mu=\sigma-\tau+\lambda),\tag{42}\\
&\pair{q_{n-\tau-\nu}T^\tau}{S_{n-\sigma}p_{\sigma-\nu}}, &&(\nu=1,2,\ldots,\sigma-\tau).
\tag{43}
\end{align}
De (31.), sin embargo, se obtienen los siguientes valores para (42.) y (43.):
\begin{align}
&\pair{q_{n-\sigma-\lambda}S^\sigma}{T_{n-\tau}p_{\tau-\lambda}}
 +\sum_{\alpha}\Pi\pair{q_{n-\sigma-\lambda-\alpha}S^\sigma}{T_{n-\tau}p_{\tau-\lambda-\alpha}},\tag{42a}\\
&\Pi(q_{n-\sigma}S^\sigma)(T_{n-\tau}p_\tau)
 +\sum_{\beta}\Pi\pair{q_{n-\sigma-\beta}S^\sigma}{T_{n-\tau}p_{\tau-\beta}}.
\tag{43a}
\end{align}
Así, las formas (42.) pueden expresarse linealmente por (41.) y polares de (41.), y las formas (43.) por polares de la forma fundamental y de las formas (41.). Asimismo (31.) da la dependencia lineal de las formas (41.) respecto de (42.) y polares de (42.). Por tanto, según la definición de Clebsch,\footnote{Clebsch, loc. cit., \S\ 7.} las formas (42.) son equivalentes a las formas (41.), mientras que (43.) conduce de vuelta a la forma fundamental. El proceso generador queda entonces dado con igual derecho por (41.) o por (42.); definiremos el proceso (41.), más simple en lo que se refiere al número $\lambda$, como el proceso de plegamiento. El número $\lambda$, el defecto del plegamiento (cf. \S\ 2), da el número de filas que faltan en el producto determinantal, a saber, en aquel producto matricial que, para el plegamiento dado, contiene el número máximo de filas. Como $\lambda$ sólo llega hasta el menor de los dos valores $\tau,n-\sigma$, donde $\sigma\ge\tau$, tenemos
\begin{equation}
 \lambda\le\frac n2,\ n\text{ par};\qquad
 \lambda\le\frac{n-1}{2},\ n\text{ impar}.
\tag{44}
\end{equation}
La reducción de (42.) y (43.) a (41.) sigue siendo válida aun cuando, por plegamientos posteriores, las filas $p,q$ hayan pasado a ser filas simbólicas; pues por (36.) se tiene
\[
 \pair{A^{n-\tau-\varkappa}T^\tau}{S_{n-\sigma}B_{\sigma-\varkappa}}
 =\sum_{\alpha=\max(0,\varkappa-\sigma+\tau)}^{\min(\tau,n-\sigma)}
 \eps\pair{R^{\tau-\alpha}B^{n-\sigma+\varkappa}}{A_{\tau+\varkappa}R_{n-\sigma-\alpha}},
\]
donde las filas $R$ han surgido de $S$ y $T$ según (33.). Por tanto, estas formas que sólo contienen filas simbólicas se obtienen plegando $\pair{q_\tau A^{n-\tau-\varkappa}}{B_{\sigma-\varkappa}p_{n-\sigma}}$ o $\pair{q_{\tau-\alpha}B^{n-\sigma+\varkappa}}{A_{\tau+\varkappa}p_{n-\sigma-\alpha}}$ --según que $\sigma-\tau\gtreqless 2\varkappa$-- con la forma fundamental y las formas (41.).

Del proceso simbólico de plegamiento se obtienen, de la manera usual, los procesos invariantes de diferenciación si simplemente se reemplazan las filas simbólicas $S^\sigma,T_{n-\tau}$ por las filas de símbolos de diferenciación $\frac{\partial}{\partial p_\sigma}$, $\frac{\partial}{\partial q_{n-\tau}}$; como es sabido, el producto $\frac{\partial}{\partial p_{i_1\ldots i_\sigma}}\frac{\partial}{\partial q_{j_1\ldots j_{n-\tau}}}$ tiene entonces el significado real $\frac{\partial^2}{\partial p_{i_1\ldots i_\sigma}\partial q_{j_1\ldots j_{n-\tau}}}$. Como las filas $\frac{\partial}{\partial p_\sigma}$ y $\frac{\partial}{\partial q_{n-\tau}}$ son cogredientes con las filas $S^\sigma,T_{n-\tau}$, satisfacen las mismas identidades que éstas, en particular también las que existen entre (42.) y (42a.), y entre (43.) y (43a.). En resumen obtenemos:

\emph{Teorema VI:} Todas las formaciones invariantes de formas fundamentales arbitrarias se obtienen de éstas por aplicación repetida del proceso simbólico de plegamiento (41.), o también de los procesos invariantes de diferenciación
\begin{equation}
 \left(q_{n-\sigma-\lambda}\frac{\partial}{\partial p_\sigma}\middle|\frac{\partial}{\partial q_{n-\tau}}p_{\tau-\lambda}\right),
 \qquad(\sigma\ge\tau,\ \lambda=1,2,\ldots,\min(\tau,n-\sigma)).
\tag{45}
\end{equation}
Queda por notar que en cada plegamiento (41.) el peso total de la forma, es decir, la suma de los pesos de las filas de variables individuales (cf. \S\ 1), disminuye en el número positivo $2(\lambda^2+\lambda(\sigma-\tau))$. Esto implica, independientemente de la demostración general de finitud, la finitud de la fila de formas, es decir, de la totalidad de las formaciones invariantes de una forma fundamental dada que son lineales en los coeficientes.\footnote{Cf. Clebsch, loc. cit., \S\ 11 y 12: finitud del sistema de covariantes elementales.}

Además debe observarse que el Teorema VI sigue siendo válido para formas que no satisfacen las condiciones (12.). Pues cada factor en (3), $\pair{S^\rho}{p_\rho}^{m_\rho}$, puede reemplazarse por un producto de $m_\rho$ factores lineales $\pair{A^\rho}{p_\rho}\pair{B^\rho}{p_\rho}\cdots\pair{M^\rho}{p_\rho}$; las formaciones invariantes pasan así a las de un sistema de formas lineales, que sólo después deben ponerse iguales entre sí. Para cada forma lineal individual, sin embargo, las condiciones (12.) --que, al introducir símbolos diferenciales, contienen sólo cocientes diferenciales de segundo orden-- se satisfacen siempre.

\section*{\S\ 7. Desarrollo de las formas generales según formas normales.}
En lo que sigue (\S\ 8) derivaremos una propiedad principal del proceso de plegamiento (41.), a saber, su composibilidad a partir de plegamientos fundamentales. Para ello necesitamos trasladar al dominio $n$-ario un importante teorema dado por Mertens\footnote{F. Mertens, \emph{"Uber invariante Gebilde quatern"arer Formen} (véase la Introducción), citado como ``Mertens''.} en el dominio cuaternario. Este teorema dice que todo sistema finito de formas puede reemplazarse por un sistema equivalente de formas normales. Por ``forma normal'' entendemos --generalizando la definición binaria y ternaria\footnote{Study, loc. cit., p. 54.}-- una forma cuyas formaciones invariantes lineales en los coeficientes se anulan idénticamente, con excepción de la propia forma fundamental. Así, por el Teorema VI (\S\ 6), una forma normal $\varphi$ satisface el sistema de ecuaciones diferenciales
\begin{equation}
 \left(q_{n-\sigma-\lambda}\frac{\partial}{\partial p_\sigma}\middle|\frac{\partial}{\partial q_{n-\tau}}p_{\tau-\lambda}\right)\varphi=0,
 \qquad(\sigma\ge\tau,\ \lambda=1,2,\ldots,\min(\tau,n-\sigma)),
\tag{46}
\end{equation}
donde las filas $q_{n-\sigma-\lambda},p_{\tau-\lambda}$ han de considerarse como cantidades arbitrarias. En el dominio cuaternario --teniendo en cuenta (39.) para $\sigma=\tau=2$-- estas ecuaciones diferenciales coinciden por completo con las diez ecuaciones diferenciales dadas por Mertens sin introducir las filas arbitrarias $p,q$ (loc. cit., fórmula 28). Su conocimiento, junto con el Teorema VI, nos permite trasladar su demostración casi palabra por palabra a $n$ variables. Lo único que debe añadirse es la verificación, obtenida mediante la identidad de descomposición (30.), de que las formas construidas son formas normales; en el dominio cuaternario esta verificación se seguía todavía sin transformación ulterior. Por ello basta dar una breve exposición de la prueba, y por simplicidad en el caso relevante más abajo, a saber, el de una sola forma fundamental $f$ y su fila de formas (cf. \S\ 6, final); para formaciones invariantes de grado arbitrario en los coeficientes de un número arbitrario de formas fundamentales, la demostración sigue siendo la misma.

La idea básica de la demostración es reemplazar, mediante la introducción de los coeficientes transformados, una forma con $n-1$ filas $p_\rho$ por formas con $n$ filas $\xi$ cogredientes con las $x$, a saber, las filas de las cantidades de transformación. De estas formas se obtienen directamente las formas normales buscadas mediante procesos invariantes de diferenciación. El desarrollo de las formas generales según formas normales es mediado por un desarrollo en serie para formas con $\rho$ ($\rho<n$) filas cogredientes $\xi$; los procesos invariantes de diferenciación conducen de vuelta a las formas en las filas originales $p_\rho$.

Sean $\xi^{(1)},\ldots,\xi^{(n)}$ las filas de cantidades de transformación, de modo que
\begin{equation}
\begin{aligned}
x_i&=\xi_i^{(1)}x'_1+\xi_i^{(2)}x'_2+\cdots+\xi_i^{(n)}x'_n,\\
p_{i_1\ldots i_\rho}&=\sum_j(\xi_{i_1}^{(j_1)}\cdots \xi_{i_\rho}^{(j_\rho)})p'_{j_1\ldots j_\rho}.
\end{aligned}
\tag{47}
\end{equation}
Los coeficientes transformados de las formas $f,\varphi,\ldots$ se denotarán por $(f),(\varphi),\ldots$, y sea
\[
 f=\pair{S^1}{p_1}^{m_1}\pair{S^2}{p_2}^{m_2}\cdots\pair{S^{n-1}}{p_{n-1}}^{m_{n-1}}.
\]
Entonces
\begin{equation}
\begin{aligned}
(f)&=\pair{S^1}{\xi^{(1)}}^{m_{11}}\cdots\pair{S^1}{\xi^{(n)}}^{m_{1n}}
\pair{S^2}{\xi^{(1)}\xi^{(2)}}^{m_{21}}\cdots
\pair{S^{n-1}}{\xi^{(2)}\cdots\xi^{(n)}}^{m_{n-1,n}}\\
&\quad\cdot\pair{s^{(1)}}{\xi^{(1)}}^{k_1}\pair{s^{(2)}}{\xi^{(2)}}^{k_2}\cdots\pair{s^{(n)}}{\xi^{(n)}}^{k_n},
\end{aligned}
\tag{48}
\end{equation}
donde $m_{11}+\cdots+m_{1n}=m_1$, y así sucesivamente. De cada coeficiente $(f)$ para el cual
\begin{equation}
 k_1\ge k_2\ge k_3\ge\cdots\ge k_n,
\tag{49}
\end{equation}
el proceso de diferenciación que agota exactamente las filas $\xi$,
\begin{equation}
\left(\frac{\partial}{\partial\xi^{(1)}}\middle|p_1\right)^{k_1-k_2}
\left(\frac{\partial}{\partial\xi^{(1)}}\frac{\partial}{\partial\xi^{(2)}}\middle|p_2\right)^{k_2-k_3}
\cdots
\left(\frac{\partial}{\partial\xi^{(1)}}\cdots\frac{\partial}{\partial\xi^{(n-1)}}\middle|p_{n-1}\right)^{k_{n-1}-k_n}
\left(\frac{\partial}{\partial\xi^{(1)}}\cdots\frac{\partial}{\partial\xi^{(n)}}\right)^{k_n}
\tag{50}
\end{equation}
da la forma
\begin{equation}
 \varphi=\bigl(s^{(1)}\mid p_1\bigr)^{k_1-k_2}
 \bigl(s^{(1)}s^{(2)}\mid p_2\bigr)^{k_2-k_3}\cdots
 \bigl(s^{(1)}\cdots s^{(n-1)}\mid p_{n-1}\bigr)^{k_{n-1}-k_n}
 \bigl(s^{(1)}s^{(2)}\cdots s^{(n)}\bigr)^{k_n}.
\tag{51}
\end{equation}
Las formas $\varphi$ son las formas normales buscadas. En efecto, poseen las siguientes propiedades que definen las formas normales.

1. Son formaciones invariantes de la forma fundamental $f$ lineales en los coeficientes. Esto se sigue del hecho de que surgen de $f$ por los procesos invariantes de diferenciación $\left(\frac{\partial}{\partial p_\rho}\middle|\xi^{(i_1)}\xi^{(i_2)}\cdots\xi^{(i_\rho)}\right)$ y (50.). Así, (por \S\ 6, final) pueden expresarse linealmente por polares de la fila finita de formas de $f$. Estos polares se obtienen considerando las variables $p_\rho$ que entran en $\varphi$ como coeficientes de una forma descompuesta en formas lineales con las filas de variables $p_\rho$:
\begin{equation}
 H=(q_{n-1}\mid p_{n-1})^{k_1-k_2}(q_{n-2}\mid p_{n-2})^{k_2-k_3}\cdots(q_1\mid p_1)^{k_{n-1}-k_n}.
\tag{52}
\end{equation}
Los invariantes simultáneos de las filas de formas $f$ y $H$ dan todos los polares pertinentes, puesto que éstos deben tener la misma dimensión en las filas individuales $p_\rho$ que la propia $\varphi$.

2. Satisfacen las ecuaciones diferenciales (46.). Si se aplica a $\varphi$ el proceso diferencial (45.), entonces
\[
 \varphi'=\bigl(q_{n-\sigma-\lambda}s^{(1)}\cdots s^{(\sigma)}\mid [s^{(1)}\cdots s^{(\tau)}]_{n-\tau}p_{\tau-\lambda}\bigr),\qquad (\sigma>\tau).
\]
De la identidad de descomposición (30.) se sigue, si las filas $q_\sigma=s^{(1)}\cdots s^{(\sigma)}$, $p_{n-\tau}=[s^{(1)}\cdots s^{(\tau)}]_{n-\tau}$ formadas a partir de las filas $s$ se identifican con las $q_\rho,p_\tau$ que allí aparecen, que
\[
 \varphi'=\sum_{\beta=\sigma-\tau+\lambda}^{n-\tau,\sigma}\sum_{i_\beta}\eps\,
 \pair{q_{\sigma-\beta}^{(i_1\ldots i_\beta)}q_{n-\tau+\lambda}}{p_{\sigma+\lambda}p_{n-\tau-\beta}^{(i_1\ldots i_\beta)}}.
\]
Aquí
\[
 p_{n-\tau-\beta}^{(i_1\ldots i_\beta)}\sim s^{(1)}\cdots s^{(\tau)}s^{(i_1)}\cdots s^{(i_\beta)},
\]
donde $i_1,\ldots,i_\beta$ designan cualesquiera $\beta$ números distintos elegidos entre los números 1 a $\sigma$. Como $\beta>\sigma-\tau$, se sigue que entre estos números $i_1,\ldots,i_\beta$ debe haber alguno entre 1 y $\tau$; por tanto $p_{n-\tau-\beta}^{(i_1\ldots i_\beta)}$ se anula idénticamente, y con ello se anula cada producto matricial individual y, con él, $\varphi'$.

Para la equivalencia del sistema de formas normales con la fila de formas $f$, sólo queda mostrar que todas las formas de la fila de formas pueden desarrollarse en polares de las formas normales en el sentido fijado bajo 1. Sea $\Pi$ un agregado de operaciones polares
\begin{equation}
 \left(\frac{\partial}{\partial\xi^{(i)}}\middle|\xi^{(k)}\right),
 \qquad (i\ne k;\ i=1,2,\ldots,\rho;\ k=1,2,\ldots,\rho),
\tag{53}
\end{equation}
y sea $\Pi^\sigma$ un agregado de las operaciones especiales (53.) para las que $i=\sigma$ y $k<\sigma$. Para una forma $F$ de las $\rho$ filas $\xi^{(1)},\ldots,\xi^{(\rho)}$ que tiene grado $k_\rho$ en la fila $\xi^{(\rho)}$, vale el desarrollo\footnote{F. Mertens, \emph{"Uber eine Formel der Determinantentheorie}. Sitzungsberichte der Akademie der Wissenschaften, Viena, Math.-Naturw. Kl., vol. 91, 2.ª sección (1885), fórmula 20).}
\begin{equation}
 F=\sum_{\alpha=0}^{k_\rho}\Pi F_\alpha,
 \qquad(\rho\le n),
\tag{54}
\end{equation}
donde $F_\alpha$ tiene grado $\alpha$ en la última fila $\xi^{(\rho)}$ y surge de $F$ por las operaciones diferenciales invariantes
\begin{equation}
 \left(\frac{\partial}{\partial\xi^{(1)}}\cdots\frac{\partial}{\partial\xi^{(\rho)}}\middle|\xi^{(1)}\cdots\xi^{(\rho)}\right)^\alpha
\tag{55}
\end{equation}
y $\Pi^\rho$. Si al construir $F_\alpha$ se reemplaza el proceso (55.) por
\begin{equation}
 \left(\frac{\partial}{\partial\xi^{(1)}}\cdots\frac{\partial}{\partial\xi^{(\rho)}}\middle|p_\rho\right)^\alpha,
\tag{56}
\end{equation}
entonces se obtiene una forma $F_\alpha(p_\rho)$ con sólo $\rho-1$ filas $\xi^{(1)},\ldots,\xi^{(\rho-1)}$, a la que puede aplicarse de nuevo (54.). La continuación del procedimiento conduce a formas $G(p_\rho,p_{\rho-1},\ldots,p_1)$. Para obtener de éstas el desarrollo de $F$ según (54.), se sustituyen las filas individuales $p_\rho$ por $\xi^{(1)},\ldots,\xi^{(\rho)}$ y se aplica a las formas resultantes el proceso polar $\Pi$ (53.); esto da (cf. (48.)) una suma de coeficientes transformados $(G)$. Para $\rho=n$, el proceso (55.) permanece en el primer paso. Si $c$ denota constantes numéricas y ponemos, para abreviar,
\begin{equation}
 (\xi^{(1)}\xi^{(2)}\cdots\xi^{(n)})\sim\Delta;
 \qquad
 \left(\frac{\partial}{\partial\xi^{(1)}}\frac{\partial}{\partial\xi^{(2)}}\cdots\frac{\partial}{\partial\xi^{(n)}}\right)=\nabla,
\tag{57}
\end{equation}
entonces, en el caso $\rho=n$,
\begin{equation}
 F=\sum c\Delta^\alpha(G),
\tag{58}
\end{equation}
mientras que para $\rho<n$ sólo aparece en el lado derecho $\alpha=0$, y por tanto $\sum c(G)$.

Aplicando (58.) a un coeficiente transformado $(f)$ (48.), y usando la conmutatividad de los procesos polares $\Pi^\sigma$ con todas las operaciones (56.) para las que $\rho\ge\sigma$, así como el hecho de que los polares de coeficientes transformados vuelven a dar coeficientes transformados, se obtiene
\[
 G=\left(\frac{\partial}{\partial\xi^{(1)}}\middle|p_1\right)^{\beta_1}
 \left(\frac{\partial}{\partial\xi^{(1)}}\frac{\partial}{\partial\xi^{(2)}}\middle|p_2\right)^{\beta_2}
 \cdots
 \left(\frac{\partial}{\partial\xi^{(1)}}\cdots\frac{\partial}{\partial\xi^{(n-1)}}\middle|p_{n-1}\right)^{\beta_{n-1}}
 \nabla^{\beta_n}\Pi(f)=\sum c\varphi,
\]
y de (58.) se sigue que
\begin{equation}
 (f)=\sum c\Delta^\alpha(\varphi)\qquad\text{(Mertens, fórmula 23)).}
\tag{59}
\end{equation}
Para pasar de (59.) al desarrollo de la propia $f$, consideramos de nuevo las variables $p_\rho$ como coeficientes de una forma $H$ (52.) con los grados $m_1,\ldots,m_{n-1}$ correspondientes a $f$, y ponemos $m=m_1+\cdots+m_{n-1}$. Entonces, por la definición de los invariantes,
\[
 f\cdot\Delta^m=\sum(f)(H)=\sum c\Delta^\alpha(\varphi)(H),
\]
y la aplicación del proceso $\nabla^m$, que agota exactamente las filas $\xi$, da
\begin{equation}
 f=\sum c\Phi,
\tag{60}
\end{equation}
donde las formas $\Phi$ se derivan de $\varphi$ y $H$ por procesos invariantes de diferenciación respecto de las variables, y por tanto son invariantes simultáneos de las filas de formas $\varphi$ y $H$.\footnote{En lugar de la conclusión directa, Mertens, \S\ 7, introduce una serie de procesos diferenciales ulteriores que sirven esencialmente para formar la fila de formas $H$; esto lo hace nuestro Teorema VI.} Pero como $\varphi$ es una forma normal, la fila de formas $\varphi$ se reduce a la propia forma $\varphi$; la fila de formas $H$ contiene todas las combinaciones invariantes posibles de variables. Así, los invariantes simultáneos de $\varphi$ y $H$ son los polares de $\varphi$ en el sentido indicado bajo 1.; en otras palabras, (60.) da el desarrollo de $f$ en polares de las formas normales. Queda verificar el mismo desarrollo para una forma arbitraria $\theta$ de la fila de formas $f$. Sean $\Gamma$ las formas normales derivadas de $(\theta)$ por (50.), de modo que
\begin{equation}
 (\theta)=\sum c\Delta^\beta(\Gamma).
\tag{61}
\end{equation}
Las formas $\theta$ se caracterizan por la relación, válida para todo coeficiente transformado,
\[
 (\theta)\Delta^\sigma=\sum c(f)\qquad\text{(Mertens, fórmula 9)).}
\]
Pero toda forma $F$ de las $n$ filas $\xi$ que contiene el factor $\Delta^\sigma$ contiene también la segunda $\nabla^\sigma\Pi F$ (Mertens, fórmula 20)); por tanto
\[
 (\theta)=\sum c\nabla^\sigma(f),\qquad\text{por consiguiente}\qquad \Gamma=\sum c\varphi,
\]
y de (61.) obtenemos la relación correspondiente a (59.),
\[
 (\theta)=\sum c\Delta^\beta(\varphi),
\]
que implica para $\theta$ el desarrollo (60.) en polares de las formas normales $\varphi$. Así:

\emph{Teorema VII.} Toda fila de formas $f$ puede sustituirse por un sistema equivalente de formas normales.

\section*{\S\ 8. Reducción del proceso de plegamiento a plegamientos fundamentales.}
Después del Teorema VII, llevamos ahora a cabo la reducción del proceso de plegamiento a plegamientos fundamentales mencionada al comienzo de \S\ 7. Llamamos (cf. \S\ 5) a todo plegamiento de dos filas cogredientes $S^\rho,T^\rho$ un ``plegamiento cogrediente'', y en particular al plegamiento cogrediente de defecto 1 de dos filas $S^\rho,T^\rho$ un ``plegamiento fundamental de tipo $\rho$''. Demostramos entonces, completando el Teorema VI:

\emph{Teorema VIII:} El proceso general de plegamiento (41.) puede componerse a partir de los $(n-1)$ plegamientos fundamentales de tipo $\rho$, donde $\rho=1,2,\ldots,n-1$. En otras palabras, para generar todas las formaciones invariantes basta aplicar sucesivamente los $n-1$ plegamientos fundamentales.

En el dominio binario el proceso de plegamiento (41.) coincide directamente con el único plegamiento fundamental posible; en el dominio ternario probé el Teorema VIII en \S\ 1 de mi disertación (cf. la Introducción). Allí, a causa de (44.), los plegamientos cogredientes generales son todavía idénticos a los plegamientos fundamentales, y es notable que el Teorema VIII para $n$ variables dé la reducción a plegamientos fundamentales, no sólo la reducción mucho más débil a plegamientos cogredientes. La prueba se modela en principio sobre la prueba ternaria: los plegamientos más generales se construyen a partir de los plegamientos fundamentales, usando las identidades simbólicas. Generaremos:
\begin{enumerate}
\item plegamientos arbitrarios de defecto 1 mediante plegamientos cogredientes de defecto 1, es decir, mediante plegamientos fundamentales (el análogo del caso ternario),
\item plegamientos arbitrarios de defecto $\lambda$ mediante plegamientos cogredientes de defecto $\lambda$ y plegamientos arbitrarios de defecto 1,
\item plegamientos cogredientes de defecto $\lambda$ mediante plegamientos cogredientes de defecto $\lambda-1$ y plegamientos arbitrarios de defecto 1.
\end{enumerate}
Como es esencial, suponemos por el Teorema VII que las dos formas $S$ y $T$ que han de plegarse son formas normales. Así, por (46.), valen las relaciones:
\begin{equation}
\begin{aligned}
\pair{q_{n-\sigma_1-\lambda}S^{\sigma_1}}{S_{n-\sigma_2}p_{\sigma_2-\lambda}}&=0,
&& (\sigma_1\ge\sigma_2,\ \lambda=1,2,\ldots,\sigma_2,n-\sigma_1),\\
\pair{q_{n-\tau_1-\lambda}T^{\tau_1}}{T_{n-\tau_2}p_{\tau_2-\lambda}}&=0,
&& (\tau_1\ge\tau_2,\ \lambda=1,2,\ldots,\tau_2,n-\tau_1).
\end{aligned}
\tag{62}
\end{equation}
Además, por \S\ 2 (final), basta suponer que las formas $S$ y $T$ son lineales en todas las filas de variables; es decir, las formas construidas pertenecen a la fila de formas
\[
 f=\pair{S^1}{p_1}\pair{S^2}{p_2}\cdots\pair{S^{n-1}}{p_{n-1}}
   \pair{T^1}{p_1}\pair{T^2}{p_2}\cdots\pair{T^{n-1}}{p_{n-1}}.
\]
1) Generación de $\pair{q_{n-\sigma-1}S^\sigma}{T_{n-\tau}p_{\tau-1}}$, donde $\sigma>\tau$, a partir de plegamientos fundamentales de tipo $\tau+\alpha$, $\alpha=0,1,\ldots,\sigma-\tau$. Del plegamiento fundamental de tipo $\tau$
\[
 \pair{q_{n-\tau-1}S^\tau}{T_{n-\tau}p_{\tau-1}}
\]
obtenemos, mediante la sustitución $w\sim T_{n-\tau}p_{\tau-1}$, una forma de la fila de formas $f$ con tres filas de índice de peso superior (cf. \S\ 1) $\tau+1$:
\begin{equation}
 \pair{wS^\tau}{p_{\tau+1}}\pair{S^{\tau+1}}{p_{\tau+1}}\pair{T^{\tau+1}}{p_{\tau+1}}.
\tag{63}
\end{equation}
Un plegamiento fundamental de tipo $\tau+1$ aplicado a (63.) da
\[
 \pair{q_{n-\tau-2}wS^\tau}{S_{n-\tau-1}p_\tau};
\]
por (31.), con la sustitución $q_{n-\tau-2}w=q'_{n-\tau-1}$, éste tiene el valor
\[
 \eps\pair{q_{n-\tau-2}wS^{\tau+1}}{S^\tau p_\tau}
 +\sum_{\alpha=1}^{\tau,n-\tau-1}\Pi\pair{q'_{n-\tau-1-\alpha}S^{\tau+1}}{S_{n-\tau}p_{\tau-\alpha}}.
\]
La expresión bajo el signo de suma se anula idénticamente por (62.); así hemos alcanzado una forma
\[
 \pair{wS^{\tau+1}}{p_{\tau+2}}\pair{S^{\tau+2}}{p_{\tau+2}}\pair{T^{\tau+2}}{p_{\tau+2}},
\]
que se obtiene de (63.) cambiando $\tau$ por $\tau+1$. Repitiendo el proceso $\sigma-\tau$ veces se llega por tanto al plegamiento deseado:
\[
 \pair{wS^\sigma}{p_{\sigma+1}}=\eps\pair{q_{n-\sigma-1}S^\sigma}{T_{n-\tau}p_{\tau-1}}.
\]
Indicamos el proceso que acabamos de efectuar mediante la fórmula
\begin{equation}
 T^\tau S^\tau,
\quad S^{\tau+1},\quad S^{\tau+2},\ldots,S^\sigma;
\tag{64}
\end{equation}
y vemos que sólo se usa la propiedad de forma normal de $S$, no la de $T$. Puede efectuarse un proceso dual a (64.), usando sólo la propiedad de forma normal de $T$, en el orden inverso, a saber:
\begin{equation}
 S^\sigma T^\sigma,
\quad T^{\sigma-1},\quad T^{\sigma-2},\ldots,T^\tau,
\tag{65}
\end{equation}
según las relaciones:
\begin{align*}
\pair{q_{n-\sigma-1}S^\sigma}{T_{1-\sigma}p_{\sigma-1}}&=\eps\pair{T_{n-\sigma}y}{p_{\sigma-1}},\quad y\sim q_{n-\sigma-1}S^\sigma,\\
\pair{q_{n-\sigma}T^{\sigma-1}}{T_{n-\sigma}yp_{\sigma-2}}&=\eps\pair{T_{n-\sigma+1}y}{p_{\sigma-2}},\\
&\vdots\\
\pair{q_{n-\tau-1}T^\tau}{T_{n-\tau-1}yp_{\tau-1}}&=\eps\pair{q_{n-\sigma-1}S^\sigma}{T_{n-\tau}p_{\tau-1}}.
\end{align*}
También debemos señalar la conexión con la disminución del peso total dada en \S\ 6 (final): para plegamientos de defecto 1 esta disminución tiene el valor $2(1+(\sigma-\tau))$, y para plegamientos fundamentales el valor $2\cdot1$. Por tanto, si la composición a partir de plegamientos fundamentales es posible en absoluto, se requieren necesariamente exactamente $\sigma-\tau+1$ tales plegamientos, como se llevó a cabo arriba.

2) Generación de plegamientos generales a partir de plegamientos cogredientes y fundamentales. La disminución de peso para un plegamiento general (41) es $2(\lambda^2+\lambda(\sigma-\tau))=2[\lambda^2+(\sigma-\tau)(1+(\lambda-1))]$. Esto da la posibilidad de descomponerlo en un plegamiento cogrediente de defecto $\lambda$ (disminución de peso $2\lambda^2$) y $\sigma-\tau$ plegamientos de defecto 1 con diferencia de índices de peso $\lambda-1$ (disminución de peso $2\{1+(\lambda-1)\}$). En el caso $\lambda=1$ esta descomposición coincide con la dada bajo 1); mostramos que efectivamente puede realizarse.

Del plegamiento cogrediente de defecto $\lambda$
\[
 \pair{q_{n-\tau-\lambda}S^\tau}{T_{n-\tau}p_{\tau-\lambda}}
\]
obtenemos, mediante la sustitución $R^\lambda\sim T_{n-\tau}p_{\tau-\lambda}$, una forma con dos filas de índices de peso superior $\tau+\lambda$ y $\tau+1$:
\begin{equation}
 \pair{R^\lambda S^\tau}{p_{\tau+\lambda}}\pair{S^{\tau+1}}{p_{\tau+1}}.
\tag{66}
\end{equation}
La fila con el índice de peso inferior $\tau+1$ pertenece a una forma normal; por tanto, (66.) admite un plegamiento de defecto 1 compuesto de $\lambda$ plegamientos fundamentales, en el orden (65):
\[
 (R^\lambda S^\tau)S^{\tau+\lambda},\quad S^{\tau+\lambda-1},\quad S^{\tau+\lambda-2},\ldots,S^{\tau+1}.
\]
Teniendo en cuenta (31.) y (62.), esto da
\[
 \pair{q_{n-\tau-\lambda-1}R^\lambda S^\tau}{S_{n-\tau-1}p_\tau}=\eps\pair{R^\lambda S^{\tau+1}}{p_{\tau+\lambda+1}},
\]
es decir, una forma obtenida de (66.) sustituyendo $\tau$ por $\tau+1$, exactamente como en 1) para (63.). Repitiendo el proceso $\sigma-\tau$ veces se llega al plegamiento deseado:
\[
 \pair{R^\lambda S^\sigma}{p_{\sigma+\lambda}}=\eps\pair{q_{n-\sigma-\lambda}S^\sigma}{T_{n-\tau}p_{\tau-\lambda}}.
\]
En todo el proceso sólo se usó la propiedad de forma normal de $S$. Considerando $T$ como forma normal se obtiene el proceso dual, con inversión completa del orden, según las fórmulas
\begin{align*}
\pair{q_{n-\sigma-\lambda}S^\sigma}{T_{n-\sigma}p_{\sigma-\lambda}}&=\eps\pair{T_{n-\sigma}R_\lambda}{p_{\sigma-\lambda}},\quad R_\lambda\sim q_{n-\sigma-\lambda}S^\sigma,\\
\pair{q_{n-\sigma}T^{\sigma-1}}{T_{n-\sigma}R_\lambda p_{\sigma-\lambda-1}}&=\eps\pair{T_{n-\sigma+1}R_\lambda}{p_{\sigma-\lambda-1}},\\
&\vdots\\
\pair{q_{n-\tau-1}T^\tau}{T_{n-\tau-1}R_\lambda p_{\tau-\lambda}}&=\eps\pair{q_{n-\sigma-\lambda}S^\sigma}{T_{n-\tau}p_{\tau-\lambda}}.
\end{align*}
3) Generación de plegamientos cogredientes de defecto $\lambda$ a partir de plegamientos cogredientes de defecto $\lambda-1$ y plegamientos fundamentales. La disminución de peso para los plegamientos cogredientes de defecto $\lambda$ es $2\lambda^2=2((\lambda-1)^2+2\lambda-1)$, lo que da la posibilidad de una descomposición en un plegamiento cogrediente de defecto $\lambda-1$ y $2\lambda-1$ plegamientos fundamentales. Esto se ve del modo más sencillo poniendo primero $n=2\lambda$. El único plegamiento cogrediente posible de defecto $\lambda$ es $\pair{S^\lambda}{T_\lambda}=\pair{S^\lambda}{T_{n/2}}$; contiene una fila $u$ y una fila $x$ menos que el plegamiento de defecto $\lambda-1$ de las mismas filas, $\pair{uS^\lambda}{T_\lambda x}$. Recíprocamente, componemos el plegamiento de defecto $\lambda$ a partir de plegamientos fundamentales que efectúan la desaparición de una fila $u$ y una fila $x$ (disminución de peso $2(n-1)=2(2\lambda-1)$) y al mismo tiempo generan una nueva fila simbólica $R^\lambda$, junto con un plegamiento cogrediente de defecto $\lambda-1$. El plegamiento más simple que hace esto es el siguiente de defecto 1:
\[
 \pair{sT^\lambda}{\sigma p_\lambda}=\eps\pair{S^{2\lambda-1}}{R_{\lambda-1}p_\lambda}=\eps\pair{R^\lambda}{p_\lambda},
\]
donde
\[
 R^{\lambda+1}=sT^\lambda,
 \quad s=S^1,
 \quad \sigma=S_1,
 \quad R^\lambda\sim\sigma R_{\lambda-1}.
\]
Por (65.) y (64.) se compone de los $2\lambda-1$ plegamientos fundamentales
\begin{equation}
 T^\lambda S^\lambda,S^{\lambda-1},\ldots,S^1;
 \qquad R^{\lambda+1}S^{\lambda+1},S^{\lambda+2},\ldots,S^{2\lambda-1}.
\tag{67}
\end{equation}
Un plegamiento de defecto $\lambda-1$, teniendo en cuenta (21.) y (62.), da
\[
 \pair{uS^\lambda}{\sigma R_{\lambda-1}x}=\eps\pair{R^{\lambda+1}}{S_\lambda x}
 =\eps\pair{sT^\lambda}{S^\lambda x}=\eps\pair{S^\lambda}{T_\lambda}.
\]
El caso general, para dos filas $S^\rho,T^\rho$ y $n$ arbitrario, se obtiene directamente del caso especial. En lugar de (67.) se tienen los siguientes $2\lambda-1$ plegamientos fundamentales:
\begin{equation}
 T^\rho S^\rho,S^{\rho-1},\ldots,S^{\rho-\lambda+1};
 \qquad R^{\rho+1}S^{\rho+1},S^{\rho+2},\ldots,S^{\rho+\lambda-1}.
\tag{68}
\end{equation}
La primera mitad de los plegamientos da la forma
\[
 \pair{q_{n-\rho-1}T^\rho}{S_{n-\rho+\lambda-1}p_{\rho-\lambda}},
\]
de la cual, mediante las sustituciones
\[
 S_{n-\rho+\lambda-1}p_{\rho-\lambda}\sim w,
 \qquad T^\rho w=R^{\rho+1},
\]
y la aplicación de la segunda mitad de los plegamientos, se obtiene
\[
 \pair{q_{n-\rho-\lambda}S^{\rho+\lambda-1}}{R_{n-\rho-1}p_\rho}.
\]
Sustituyendo
\[
 q_{n-\rho-\lambda}S^{\rho+\lambda-1}\sim y,
\]
se obtiene, mediante un plegamiento cogrediente de defecto $\lambda-1$,
\begin{equation}
 \pair{q_{n-\rho-\lambda+1}S^\rho}{yR_{n-\rho-1}p_{\rho-\lambda+1}}.
\tag{69}
\end{equation}
Mostramos que éste es el plegamiento deseado de defecto $\lambda$. Sustitúyase
\[
 R_{n-\rho-1}p_{\rho-\lambda+1}=R_{n-\lambda}
\]
y obsérvese que resolver $y$ da
\[
 \pair{q_{n-\rho-\lambda+1}R^\lambda}{S_{n-\rho}y}
 =\eps\pair{q_{n-\rho-\lambda}S^{\rho+\lambda-1}}{S_{n-\rho}p'_{\rho-1}}=0.
\]
Así, por (20.), el valor de (69.) es
\[
 \eps\pair{S^\rho R^\lambda}{p_{\rho+\lambda-1}y}
 =\eps\pair{q_{n-\rho-\lambda}S^{\rho+\lambda-1}}{[S^\rho R^\lambda]_{n-\rho-\lambda}p_{\rho+\lambda-1}}.
\]
Este plegamiento es de tipo (43.), y por tanto equivalente a
\[
 \pair{q_{n-\rho-\lambda}S^\rho}{R_{n-\rho-1}p_{\rho-\lambda+1}}
 =\eps\pair{wT^\rho}{R_\lambda p_{\rho-\lambda+1}},
 \qquad R_\lambda\sim q_{n-\rho-\lambda}S^\rho.
\]
Aquí $R_\lambda$ ya está en la forma final deseada; la expresión corresponde a la forma $\pair{sT^\lambda}{S_\lambda x}$ del caso especial. Ahora, observando que la resolución de $w$ y $R_\lambda$ da
\[
 \pair{wR^{n-\lambda}}{T_{n-\rho}p_{\rho-\lambda+1}}
 =\eps\pair{q'_{n-1}R^{n-\lambda}}{S_{n-\rho+\lambda-1}p_{\rho-\lambda}}
 =\eps\pair{q''_{n-\sigma-1}S^\rho}{S_{n-\rho+\lambda-1}p_{\rho-\lambda}}=0,
\]
obtenemos por (21.) el valor
\[
 \pair{wq_{n-\rho+\lambda-1}}{T_{n-\rho}R_\lambda}
 =\eps\pair{q_{n-\rho+\lambda-1}[T_{n-\rho}R_\lambda]^{n-\lambda}}{S_{n-\rho+\lambda-1}p_{\rho-\lambda}}
 \quad\text{equivalente a}\quad
 \pair{q_{n-\rho-\lambda}S^\rho}{T_{n-\rho}p_{\rho-\lambda}}.
\]
En todo el proceso se usó sólo la propiedad de forma normal de $S$, no la de $T$. Por tanto, el plegamiento cogrediente de defecto $\lambda-1$ usado para generar (69.) puede sustituirse por $2\lambda-3$ plegamientos fundamentales y un plegamiento cogrediente de defecto $\lambda-2$, que a su vez admite una descomposición en plegamientos fundamentales, y así sucesivamente. Análogamente, la generación se obtiene usando la propiedad de forma normal de $T$ y aplicando el proceso dualístico, esto es, los plegamientos fundamentales
\[
 S^\rho T^\rho,\quad T^{\rho-1},\ldots,T^{\rho-\lambda+1};
 \qquad R^{\rho-1}T^{\rho-1},\quad T^{\rho-2},\ldots,T^{\rho-\lambda+1},
\]
y un plegamiento dual a (69.). Entonces se obtiene el plegamiento
\[
\pair{q_{n-\rho-\lambda}T^\rho}{S_{n-\rho}p_{\rho-\lambda}},
\]
equivalente al anterior.

Añadiendo la descomposición obtenida bajo 2), hemos obtenido en efecto una generación de los plegamientos más generales a partir de plegamientos fundamentales. Resta mostrar que esta generación es única, salvo permutación del orden, es decir, que a cada plegamiento le corresponde uno y sólo un número $\alpha_\rho$ de plegamientos fundamentales de tipo $\rho$, donde $\rho=1,2,\ldots,n-1$.

Sean $m_\rho$ los grados en las filas de variables $p_\rho$ de la forma inicial $f$, y $l_\rho$ los grados de la forma producida por plegamiento. Cada plegamiento fundamental de tipo $\rho$ transforma los grados $m_{\rho-1},m_\rho,m_{\rho+1}$ respectivamente en $m_{\rho-1}+1,m_\rho-2,m_{\rho+1}+1$. Así, los $\alpha$ satisfacen el sistema de ecuaciones
\begin{equation}
 m_\rho-l_\rho=-\alpha_{\rho-1}+2\alpha_\rho-\alpha_{\rho+1},
 \qquad (\rho=1,2,\ldots,n-1),
\tag{70}
\end{equation}
donde debe ponerse $\alpha_0=\alpha_n=0$, y cuyo determinante tiene valor $n$.\footnote{Si $\Delta_n$ denota el determinante de $n$ filas, entonces vale la fórmula de recurrencia: $\Delta_n=2\Delta_{n-1}-\Delta_{n-2}$; es decir, $\Delta_n-\Delta_{n-1}=\Delta_{n-1}-\Delta_{n-2}=\cdots=\Delta_2-\Delta_1=1$, puesto que $\Delta_2=3$, $\Delta_1=2$.} Los $\alpha$ quedan así determinados de manera única, y en efecto, por el Teorema VIII, como enteros positivos; esto da condiciones para las diferencias $m_\rho-l_\rho$.

Los resultados de este parágrafo valen en virtud de las relaciones (62.); sin embargo, si las relaciones (62.) no se cumplen, las expresiones finales no son en modo alguno equivalentes a las expresiones iniciales. La definición de equivalencia dada en \S\ 6 admite también la formulación:\footnote{Cf. la consideración al final del parágrafo siguiente.} ``Dos formas son equivalentes si y sólo si difieren por formas de plegamiento superior, es decir, por formas que exhiben una mayor disminución de peso.'' Esta mayor disminución de peso ocurre siempre cuando las dos filas $q_{n-\sigma-i},p_{\tau-\lambda}$ que entran en el plegamiento son realmente filas de variables, como sucede en (41.), (42.), en los plegamientos que aparecen en este parágrafo como equivalentes, y también en las formas equivalentes del Teorema VII. Pero no tiene por qué ocurrir cuando esas filas se derivan de filas simbólicas por sustitución, como fue el caso de las formas de este parágrafo que pudieron expresarse unas en términos de otras en virtud de (62.). Se ve fácilmente que las formas que se anulan tienen incluso, en parte, un peso total mayor.

\section*{\S\ 9. Filas de formas. Teoremas de reducción.}
Definimos ahora con algo más de precisión la fila de formas introducida en \S\ 6 (final) como ``la totalidad de las formaciones invariantes de una forma inicial dada que son lineales en los coeficientes, es decir, la totalidad de las formas que surgen de la forma inicial por plegamiento en sí misma, ordenadas según formas superiores''.

Para explicar las formas superiores, supongamos por el Teorema VII que la forma inicial está dada como producto de dos formas normales $S\cdot T$. Por forma superior entendemos formas con plegamiento superior en sí mismas, y entre formas con el mismo plegamiento en sí mismas, unas especialmente normalizadas. Más precisamente: supongamos que la forma $A$ ha surgido mediante $\alpha_\rho$ plegamientos fundamentales de tipo $\rho$, y la forma $B$ mediante $\beta_\rho$ tales plegamientos. Entonces $B$ es superior a $A$:
\begin{enumerate}
\item si, en el sistema de desigualdades $\beta_\rho\ge\alpha_\rho$, $\rho=1,2,\ldots,n-1$, el signo de desigualdad es estricto para al menos un valor de $\rho$;\footnote{Si en el sistema de desigualdades el signo $>$ aparece en parte y el signo $<$ en parte, las formas no son comparables, pues una forma que surge de otra por plegamiento tiene siempre un sistema de valores positivo de plegamientos fundamentales, y por tanto en nuestro caso valores positivos o nulos $\beta_\rho-\alpha_\rho$.}
\item si el signo de igualdad vale para todos los valores $\rho$, pero se han combinado menos filas simbólicas en un único producto matricial. Esta normalización resulta necesaria por las reducciones de \S\ 8. Así, por ejemplo, $\pair{S^{n-1}}{T_{n-1}}$ es superior a $\pair{S^{n-1}}{[S^1T^\rho]_{n-\rho-1}p_\rho}$;
\item si el signo de igualdad vale para todos los valores $\rho$ y el número de filas simbólicas combinadas en un producto matricial es el mismo, entonces $B$ se hace superior a $A$ por una normalización arbitraria en sí misma pero fijada de una vez por todas. Esta normalización es siempre posible por el número finito de formas pertenecientes a un sistema de valores $\alpha$; puede realizarse, por ejemplo, privilegiando la única forma $S$ (mayor número de filas simbólicas $S$, mayor suma de los índices de peso superiores de las filas $S$, etc.).
\end{enumerate}
Conviene imaginar la fila de formas como dispuesta en una variedad de puntos reticulares de dimensión $(n-1)$, donde las $n-1$ direcciones corresponden a los $n-1$ plegamientos fundamentales. A cada forma le corresponde entonces uno y sólo un sistema de valores $\alpha$, es decir, un punto reticular entero positivo, mientras que las distintas formas pertenecientes a un sistema de valores $\alpha$ quedan determinadas de manera única en su orden por la normalización anterior. Una forma arbitraria de la fila de formas da la forma inicial de una nueva fila de formas, contenida como parte de la fila total de formas; en efecto, está contenida como la parte que consta de las formas obtenidas a partir de la nueva forma inicial al avanzar en las $n-1$ direcciones positivas.

En virtud del Teorema VIII y de la teoría general de los desarrollos en serie, valen para las filas de formas exactamente los teoremas de reducción de los dominios binario y ternario (cf. disertación, \S\ 3). Derivaremos brevemente el teorema principal. Primero enunciamos la definición:

``En un sistema dado de formas, supóngase que cierto número finito de formas se ha reunido en un módulo. Llamamos reducible a una forma si puede expresarse mediante formas que tengan invariantes como factores, o mediante `formas superiores'; es decir, mediante formas superiores de la fila total de formas a la que pertenece la forma, o mediante formas que contengan los símbolos del módulo en orden superior. Una fila de formas reducible se llama reductor.'' Entonces:

\emph{Teorema IX:} Si la forma inicial de una fila de formas es reducible por haber surgido mediante plegamiento con un reductor, entonces toda la fila de formas que surge de esta forma inicial es reducible.

Para la prueba, obsérvese que una forma $h$ surgida por plegamiento de una forma $f$ con una segunda forma $g$ puede, a causa de (45.), considerarse como una forma $f$ que contiene varias filas de variables cogredientes $p_\rho,p'_\rho$. Tales formas pueden representarse, por la teoría general de los desarrollos en serie, como polares de los covariantes elementales de $f$, aplicando sucesivamente el desarrollo en serie a cada par de filas de variables cogredientes $p_\rho,p'_\rho$. Por (38.), los covariantes elementales que aparecen son las formas generadas por plegamiento cogrediente. Sin embargo, como el orden en que se aplica el desarrollo en serie sigue siendo arbitrario, podemos escoger en particular un orden correspondiente a la generación de los plegamientos generales a partir de plegamientos fundamentales. El sistema resultante de covariantes elementales es entonces idéntico a la fila de formas $f$; las formas que surgen por plegamiento cogrediente de defecto superior se ordenan entre las que surgen por plegamientos fundamentales. Pero también se obtienen polares de la fila de formas $f$ para cualquier orden del desarrollo en serie, al que puede corresponder un segundo sistema de covariantes elementales. Como la fila de formas agota la totalidad de las formaciones invariantes, toda forma del segundo sistema puede expresarse linealmente por polares de la fila de formas, y en efecto por polares de aquellas formas que muestran la misma o una mayor disminución de peso. Esto último se sigue porque (cf. \S\ 7) los polares pueden considerarse como invariantes simultáneos de una fila de formas $H$ (52.), con los números de grado correspondientes a la forma $h$, y de la fila de formas $f$. Cada plegamiento de $H$ en sí misma corresponde a una disminución de peso que, después de efectuar la sustitución (11.), se transfiere a las filas simbólicas y por tanto a las formas de $f$.\footnote{Esta consideración subyace también a la segunda formulación del concepto de equivalencia (\S\ 8, final). Otras explicaciones de los diferentes sistemas de covariantes elementales en el dominio ternario se encuentran en Study, loc. cit., II, \S\ 7, Teoremas 7) y 8). Por nuestro Teorema VII los mismos resultados valen también para $n$ variables.}

Una forma arbitraria de la fila de formas $h$ ha surgido por plegamiento de $h$ en sí misma; es decir, por una parte mediante un plegamiento superior de $g$ con $f$, y por otra mediante plegamiento de $f$ o $g$ en sí mismas. En ambos casos el desarrollo en serie conduce a polares de la fila de formas $f$, tal como se mostró antes para la forma inicial $h$. Si en particular $f$ es un reductor, resulta un desarrollo en polares del reductor; por tanto la fila de formas $h$ se vuelve reducible.

Las aplicaciones del teorema a la reducción de sistemas completos de formas se siguen análogamente a los dominios binario y ternario.

\clearpage
% ---- Continuation appended 2026-06-02: Paper 05 complete and Paper 06 through Section 6, Spanish ----
\providecommand{\Kfield}{\mathfrak{K}}
\providecommand{\Ssys}{\mathfrak{S}}
\providecommand{\Mfield}{\mathfrak{M}}
\providecommand{\tq}{\tau}
\providecommand{\Lsys}{\mathfrak{L}}
\providecommand{\Jdom}{\mathfrak{J}}
\providecommand{\ds}{\displaystyle}

\begin{center}
{\Large\bfseries 5. Cuerpos de funciones racionales}\par
\vspace{1.2em}
\emph{Jahresbericht der DMV} 22 (1913), pp. 316--319
\end{center}

\vspace{1.5em}

Las cuestiones siguientes se remontan originalmente a conversaciones con E. Fischer. Algunas de ellas, por lo demás, ya habían sido planteadas y resueltas por E. Steinitz para el caso especial de una indeterminada, e incluso bajo hipótesis más generales sobre el dominio de coeficientes.\footnote{E. Steinitz: \emph{Algebraische Theorie der Körper}, especialmente \S\ 24. Crelle's Journal 137, 1910.}

Por ``cuerpo de funciones racionales'' entiendo un cuerpo cuyos elementos son funciones racionales de $n$ indeterminadas, con coeficientes en un cuerpo numérico dado, que en particular puede comprender también todos los números complejos. Ejemplos de tales cuerpos son el cuerpo de las funciones simétricas de $n$ cantidades, o, más generalmente, la totalidad de las funciones racionales de $n$ indeterminadas que admiten las permutaciones de cierto grupo (los \emph{Gattungsbereiche} de Lagrange); además, el cuerpo de invariantes.

Entre los dos primeros cuerpos mencionados y el último existe una diferencia característica: los primeros contienen $n$ funciones algebraicamente independientes, las funciones simétricas elementales; mientras que el número de invariantes algebraicamente independientes es siempre menor que el número de indeterminadas. Por eso es importante que, en general, se puedan asociar de manera biunívoca tales cuerpos de la segunda especie con cuerpos de la primera especie, sustituyendo algunas de las indeterminadas por números. Por consiguiente, en lo que sigue me limitaré a cuerpos de la primera especie, es decir, a cuerpos con $n$ funciones algebraicamente independientes; en virtud de esa asociación, los resultados valen entonces en general.

Las cuestiones se agrupan alrededor de tres conceptos de base: base racional, base mínima y base íntegra.

1. Por una ``base racional'' entiendo un número finito de funciones del cuerpo tal que toda función del cuerpo puede representarse como combinación racional de ese número finito, con coeficientes en el dominio de coeficientes prescrito. Mediante consideraciones sencillas, que en el caso de una indeterminada se encuentran también en E. Steinitz, loc. cit., se demuestra la existencia de una base racional para todo cuerpo de funciones racionales. Por ejemplo, según el teorema de Lagrange, las funciones simétricas elementales y una función perteneciente al grupo forman una base racional para los \emph{Gattungsbereiche} de Lagrange antes mencionados.

2. Toda base racional debe contener, para cuerpos de la primera especie, al menos $n$ funciones; si contiene exactamente $n$ funciones, que entonces han de ser algebraicamente independientes, la llamaré ``base mínima''. La cuestión de la existencia de una base mínima puede responderse parcialmente mediante teoremas de Lüroth, Castelnuovo y Enriques.\footnote{J. Lüroth: Beweis eines Satzes über rationale Kurven. \emph{Math. Ann.} 9, 1875. --- G. Castelnuovo: Sulla razionalità delle involuzioni piane. \emph{Math. Ann.} 44, 1893. --- F. Enriques: Sopra una involuzione non razionale dello spazio. \emph{Rend. Acc. Linc.} vol. 21, 21 de enero de 1912.} Según ellos, para cuerpos de una indeterminada existe siempre una base mínima: la función de Lüroth del cuerpo, determinada de manera única salvo transformación lineal fraccionaria (cf. Steinitz \S\ 24). Para cuerpos de dos indeterminadas existe siempre una base mínima, y en general sólo entonces, si se permite una extensión algebraica del cuerpo de coeficientes; mientras que para tres o más indeterminadas, en general, no existe en absoluto una base mínima. Todavía no se ha intentado caracterizar los cuerpos especiales que tienen base mínima.

He seguido estudiando la cuestión de la base mínima para los \emph{Gattungsbereiche} de Lagrange. Aquí adquiere el significado siguiente: ``si el \emph{Gattungsbereich} de Lagrange perteneciente al grupo $G$ posee una base mínima, entonces se puede construir racionalmente, por representación paramétrica, la totalidad de las ecuaciones con grupo prescrito $G$; y esto para todo cuerpo numérico que contenga el dominio de coeficientes de la base mínima, por tanto en particular para cualquier cuerpo numérico arbitrario si dicho dominio de coeficientes es el cuerpo de los números racionales.''

El ejemplo más sencillo lo proporciona el grupo simétrico; aquí las funciones simétricas elementales forman una base mínima con coeficientes numéricos racionales. De ello se obtiene como representación paramétrica simplemente la ecuación con coeficientes indeterminados. El teorema de irreducibilidad de Hilbert muestra cómo se pasa de ésta a arbitrariamente muchas ecuaciones no afectas sobre todo cuerpo numérico.

Ahora bien, de los teoremas de Lüroth y Castelnuovo se deduce directamente la existencia de la base mínima para todos los grupos que aparecen en las ecuaciones de grados 3 y 4; al mismo tiempo se ve además que esta base mínima tiene coeficientes numéricos racionales. Por tanto, sobre cualquier cuerpo numérico arbitrario, se puede construir racionalmente la totalidad de las ecuaciones de grados 3 y 4 con grupo prescrito. Para los grupos cíclicos y diedrales existe una base mínima con el cuerpo de las raíces $n$-ésimas de la unidad como dominio de coeficientes; lo mismo vale para algunos otros grupos metacíclicos. La cuestión de si la base mínima es siquiera característica de ciertos grupos debe quedar indecisa.

3. Para llegar al último concepto de base, considero la totalidad de los polinomios contenidos en el cuerpo. Por una ``base íntegra'' entiendo un número finito de esos polinomios tal que todo polinomio del cuerpo pueda representarse como combinación racional entera de ese número finito, con coeficientes en el dominio de coeficientes prescrito. La cuestión de la existencia de una base íntegra, que por ejemplo contiene como caso especial la finitud del sistema de invariantes, fue planteada por Hilbert en sus ``Problemas matemáticos''\footnote{D. Hilbert: Mathematische Probleme. Conferencia de París, 1900. Problema 14. \emph{Göttinger Nachr.} 1900.} en una forma algo distinta, como problema de las funciones relativamente enteras.

Por ahora sólo puedo indicar una clase de cuerpos para los cuales existe una base íntegra. En efecto, si el cuerpo contiene un sistema de $n$ polinomios tal que la resultante homogénea de los términos de dimensión más alta de esos polinomios es distinta de cero, entonces existe una base íntegra; está dada por los coeficientes de todos los $u$ en la ecuación irreducible, dentro del cuerpo, de la forma lineal:
\[
  u_0=u_1x_1+u_2x_2+\cdots+u_nx_n.\footnotemark
\]
\footnotetext{Esta ecuación irreducible aparece, en el caso de una indeterminada, en la demostración de Steinitz del teorema de Lüroth.}
Si, en particular, el valor de la resultante es igual a una unidad del dominio de coeficientes, queda asegurada también la integralidad de la representación. Un ejemplo de esta clase de cuerpos lo dan los \emph{Gattungsbereiche} de Lagrange, pues la resultante de las funciones simétricas elementales tiene el valor $\pm 1$. Otro ejemplo (sin integralidad) lo dan todos los cuerpos que contienen sólo un polinomio algebraicamente independiente; aquí la base íntegra es idéntica a la función de Lüroth del subcuerpo más pequeño que contiene todos los polinomios. Así, como es sabido, todos los invariantes proyectivos racionales enteros de una forma cuadrática en $n$ variables son funciones enteras de la discriminante.

La condición indicada es sólo suficiente, de ningún modo necesaria, para la existencia de la base íntegra. Se pueden construir fácilmente cuerpos en los que la resultante de cualesquiera $n$ polinomios se anula y que, no obstante, poseen una base íntegra dada por los coeficientes de aquella ecuación irreducible. Surge la pregunta de si quizá también en el caso más general los coeficientes de aquella ecuación suministran la base íntegra.

\clearpage

\begin{center}
{\Large\bfseries 6. Cuerpos y sistemas de funciones racionales}\par
\vspace{1.2em}
\emph{Mathematische Annalen} 76 (1915), pp. 161--196
\end{center}

\vspace{1.5em}

El presente trabajo trata cuestiones de base para sistemas arbitrarios de funciones racionales y racionales enteras. Los métodos de la teoría de cuerpos aquí empleados hacen que la resolución de estas cuestiones para cuerpos de funciones racionales --- cuerpos de funciones racionales\footnote{Cf. una comunicación preliminar con ocasión de la reunión de naturalistas de Viena: \emph{Rationale Funktionenkörper}, \emph{Jahresber. d. D. Math.-Ver.} 22 (1913), donde se da una visión de conjunto de las cuestiones y de los resultados referentes a los cuerpos de funciones.} --- aparezca como lo esencial, mientras que la generalización de los resultados a sistemas arbitrarios se obtiene como consecuencia.

Entre las cuestiones de base para sistemas generales, hasta ahora sólo se conocía la existencia de una base de módulo, garantizada por el teorema de Hilbert (\emph{Math. Ann.} 36). En lo que sigue, la cuestión de la representabilidad racional se responde completamente mediante la existencia de una base racional para todo sistema arbitrario (\S\ 7); la base racional de los cuerpos se obtiene ya en \S\ 4. Esta existencia de la base racional permite partir siempre del cuerpo o sistema definido abstractamente, y evitar así dificultades que son causadas sólo por la elección especial de la base racional, y no por el sistema mismo, tales como la aparición de denominadores especiales o de puntos fundamentales de las funciones de la base.

Para cuerpos se trata además la cuestión de la base mínima, es decir, de una base racional compuesta de funciones algebraicamente independientes (\S\ 6). Los métodos de la teoría de cuerpos conducen además a una base racional distinguida, la base de involución\footnote{El nombre se eligió porque, en su interpretación geométrica, el cuerpo de funciones racionales representa la involución más general en un espacio lineal.} (\S\ 5), que resulta especialmente importante para dominios íntegros compuestos de polinomios (\S\ 8). En ese contexto da una representación con denominador fijo (una potencia de una función determinada por el dominio íntegro), tal como se conoce, por ejemplo, en el caso especial de la representación típica de invariantes.

De la representabilidad racional extraemos ahora, en la medida de lo posible, consecuencias para la representabilidad racional entera, es decir, para la cuestión de finitud en sentido estricto.

Los teoremas de finitud conocidos hasta ahora descansan todos en el hecho de que el teorema de Hilbert sobre bases de módulos garantiza la finitud para todos los sistemas en los que una representación
\[
  F=A_1f_1+A_2f_2+\cdots+A_kf_k
\]
entraña una segunda representación
\[
  F=B_1f_1+B_2f_2+\cdots+B_kf_k,
\]
donde los $B_i$ pertenecen al sistema y son de grado menor que $F$.

En cambio, el teorema sobre la base racional y los métodos de la teoría de cuerpos permiten inferir la finitud bajo hipótesis de una especie completamente distinta.

Así, como respuesta parcial a un problema de Hilbert (Problemas matemáticos 14), se obtiene una clase de funciones relativamente enteras (dominios relativamente enteros de primera especie) que son dominios íntegros finitos y cuya base íntegra está dada por la base de involución antes mencionada (\S\ 10). En analogía con la demostración de Hilbert del teorema de Kronecker sobre el sistema fundamental de enteros algebraicos (\emph{Complete invariant systems}, \S\ 2; \emph{Math. Ann.} 42), se puede mostrar además que todos los sistemas regulares de polinomios, es decir, todos los sistemas que pueden aplicarse en otros sin puntos fundamentales, poseen una base íntegra (\S\ 12), mientras que se pueden especificar fácilmente sistemas no regulares sin base íntegra. Como ejemplo sencillo de este hecho puede mencionarse el teorema: ``en todo sistema arbitrario de infinitos polinomios en una indeterminada existe un número finito de esos polinomios tal que todo polinomio del sistema puede representarse como combinación racional entera de ese número finito.'' Otros ejemplos se encuentran en \S\ 13.

Finalmente, los últimos parágrafos (\S\S\ 14 y 15) dan un refinamiento de los teoremas de base con respecto a la integralidad.

Un instrumento esencial de la investigación es el principio de transferencia de \S\ 3, según el cual basta considerar todas las cuestiones de base aquí planteadas para aquellos sistemas en los que el número de indeterminadas coincide con el número de funciones algebraicamente independientes del sistema.

El impulso para el presente trabajo vino de conversaciones con el señor E. Fischer, en particular de una pregunta suya sobre la base mínima de los \emph{Gattungsbereiche} de Lagrange. Las investigaciones relativas a ello, que condujeron a la construcción de ecuaciones con grupo prescrito (cf. la comunicación citada sobre ``Cuerpos de funciones racionales'', parte 2), quedan reservadas para una publicación ulterior.

La ampliación esencial del presente trabajo respecto de la comunicación original consiste en que los teoremas de base allí enunciados sólo para cuerpos o dominios íntegros especiales se trasladan aquí a sistemas arbitrarios. Esta transferencia se logra de manera sencilla por medio del concepto del menor cuerpo que contiene un sistema arbitrario, como generalización del cuerpo de fracciones de un dominio íntegro (\S\ 7). Me condujo a formar este concepto el hecho de que el señor K. Hentzelt, después de conocer la base racional de los cuerpos, pudo demostrar la existencia de la base racional para sistemas arbitrarios dando un rodeo por la base de módulo de Hilbert. También me llevó al teorema sobre la base íntegra de los sistemas regulares la observación de K. Hentzelt de que los razonamientos que yo había aplicado a dominios íntegros eran válidos para sistemas arbitrarios sujetos a las mismas hipótesis; asimismo debo al señor Hentzelt algunas observaciones aisladas de menor tamaño.

Finalmente, se menciona que en el caso de una indeterminada la base racional y la base mínima de los cuerpos se encuentran en la \emph{Algebraische Theorie der Körper} de E. Steinitz, \S\ 24 (Crelle's Journal 137); allí el dominio de coeficientes se toma como cuerpo en el sentido más general. La cuestión de hasta qué punto los teoremas dados aquí siguen siendo válidos bajo esas hipótesis más generales no se trata en lo que sigue; en cualquier caso, las demostraciones requieren modificación, ya que, por ejemplo, los medios de la teoría de los determinantes funcionales fallan para cuerpos de característica $p$.

\section*{\S\ 1. Cuerpos $\Kfield_{n\rho}$ y sistemas $\Ssys_{n\rho}$.}

En lo que sigue se trata de investigar sistemas de funciones racionales de $n$ indeterminadas, en particular cuerpos de funciones racionales.

Por $f(x_1,\ldots,x_n)$, $g(x_1,\ldots,x_n)$, \ldots, o abreviadamente $f(x)$, $g(x)$, \ldots, se entenderán siempre funciones racionales de $x_1,\ldots,x_n$; se suponen en representación reducida, es decir, con numerador y denominador primos entre sí. Las funciones racionales enteras (polinomios) se designarán con letras mayúsculas, $F(x)$, $G(x)$, \ldots. El dominio de coeficientes $\Omega$ se supone algún cuerpo numérico; en particular puede contener también todos los números complejos. El cuerpo $\Omega$ puede contener también un número finito de parámetros.

Las cantidades de $\Omega$, por tanto todas las funciones de grado cero, se cuentan siempre como pertenecientes al sistema.

Con estas estipulaciones, el cuerpo formado por funciones racionales --- el cuerpo de funciones racionales --- puede definirse abstractamente.

\emph{Definición I: Un sistema de funciones racionales se llama cuerpo si satisface las condiciones:}
\begin{enumerate}
\item \emph{Junto con $f(x)$, y para cada cantidad $c$ de $\Omega$, contiene también $c\cdot f(x)$.}
\item \emph{Junto con $f(x)$ y $g(x)$, contiene siempre también $f(x)+g(x)$, $f(x)\cdot g(x)$ y, para $g(x)\ne0$, el cociente $f(x):g(x)$.}\footnote{Aquí $f(x)$ y $g(x)$ pueden ser también de grado cero, es decir, cantidades de $\Omega$.}
\end{enumerate}
Los cuerpos especiales son los que surgen por adjunción de un número finito de funciones racionales
\[
  f_1(x),\ f_2(x),\ldots, f_k(x)
\]
a $\Omega$; se denotarán por
\[
  \Omega\bigl(f_1(x),\ldots,f_k(x)\bigr)
  \quad\text{o}\quad
  \Omega(f_1,\ldots,f_k).
\]
\footnote{Que estos ``cuerpos especiales'' agotan ya la totalidad de todos los cuerpos se deduce en \S\ 4.}
Entre estos cuerpos especiales mencionamos también el cuerpo que surge por adjunción de $x_1,\ldots,x_n$ a $\Omega$,
\[
  \Omega(x_1,\ldots,x_n),
\]
que es idéntico a la totalidad de las funciones racionales de $x_1,\ldots,x_n$ con coeficientes en $\Omega$.

Introducimos además, siguiendo a E. Steinitz, el concepto de cuerpo intermedio:
\begin{quote}
``un cuerpo $\Mfield$ está entre $\Omega_1$ y $\Omega_2$ si contiene a $\Omega_1$ y está contenido en $\Omega_2$;''
\end{quote}
entonces la Definición I admite también la siguiente formulación:

El cuerpo de funciones racionales de $n$ indeterminadas es un cuerpo intermedio entre $\Omega$ y $\Omega(x_1,\ldots,x_n)$ que contiene efectivamente las $n$ indeterminadas en al menos una función.

Junto al número de indeterminadas aparece para sistemas de funciones racionales otro número característico, el número de funciones algebraicamente independientes, o rango algebraico, mediante la siguiente definición.

\emph{Definición II: Un sistema de funciones racionales tiene rango algebraico $\rho$ si se pueden indicar $\rho$ funciones del sistema que sean algebraicamente independientes, pero cualesquiera $(\rho+1)$ funciones del sistema son algebraicamente dependientes.}\footnote{Se llaman algebraicamente independientes $\rho$ funciones $f_1(x),\ldots,f_\rho(x)$ si de toda relación
\[
F\bigl(f_1(x),\ldots,f_\rho(x)\bigr)=0\quad\text{idénticamente en }x_1,\ldots,x_n
\]
se sigue que
\[
F(\lambda_1,\ldots,\lambda_\rho)=0\quad\text{idénticamente en }\lambda_1,\ldots,\lambda_\rho.
\]
En el caso contrario se llaman algebraicamente dependientes.}

Los sistemas de (finitas o infinitas) funciones racionales en $n$ indeterminadas y de rango algebraico $\rho$ se denotarán mediante el doble índice
\[
  \Ssys_{n\rho}.
\]
En particular, por
\[
  \Kfield_{n\rho}
\]
se entenderá un cuerpo de $n$ indeterminadas y de rango algebraico $\rho$. Se tiene la desigualdad
\[
  1\leq \rho\leq n.
\]

Damos aún algunos ejemplos de cuerpos $\Kfield_{nn}$ y $\Kfield_{n\rho}$:

1. Cuerpos $\Kfield_{nn}$ son los \emph{Gattungsbereiche} de Lagrange, que pueden definirse como
\begin{quote}
``la totalidad de las funciones racionales (y racionales enteras) de $x_1,\ldots,x_n$ que admiten las permutaciones de un grupo de permutaciones $G$ de $x_1,\ldots,x_n$.''
\end{quote}
Siempre contienen el cuerpo de las funciones simétricas, y por tanto también las $n$ funciones simétricas elementales algebraicamente independientes.

2. Cuerpos $\Kfield_{n\rho}$ son los cuerpos de invariantes proyectivos, formados por la totalidad de los invariantes racionales (y racionales enteros) de una o varias formas fundamentales.\footnote{Conviene considerar sólo transformaciones de determinante $1$; entonces toda combinación racional de arbitrariamente muchos invariantes admite de nuevo la transformación, y efectivamente se obtiene un cuerpo. Los invariantes homogéneos e isobáricos tomados por sí solos no forman un cuerpo, pero sí un sistema $\Ssys_{n\rho}$.} Aquí $\rho<n$, pues el número de invariantes algebraicamente independientes es siempre menor que el número de coeficientes.

Lo correspondiente vale para los cuerpos de invariantes respecto de subgrupos del grupo proyectivo.

\section*{\S\ 2. Lema auxiliar sobre funciones algebraicamente dependientes.}

El lema auxiliar de este parágrafo mostrará en \S\ 3 que el rango algebraico de un sistema es su número propiamente característico. En efecto, el lema auxiliar muestra que mediante una especialización numérica sucesiva de algunas indeterminadas, de tal modo que $\rho$ funciones algebraicamente independientes sigan siendo algebraicamente independientes, una función arbitraria algebraicamente dependiente de esas $\rho$ funciones no puede ni anularse idénticamente ni volverse indeterminada.

Sea, pues, el sistema
\begin{equation}
  g_1(x),\ g_2(x),\ldots,g_\rho(x)
\tag{1}
\end{equation}
de rango algebraico $\rho$, donde $\rho<n$, y sea $h(x)$ una función racional algebraicamente dependiente de las funciones (1). Por teoremas conocidos, el rango de la matriz funcional de (1) es igual a $\rho$. Por tanto las indeterminadas pueden llamarse, por ejemplo, de manera que
\begin{equation}
 \frac{\partial(g_1,g_2,\ldots,g_\rho)}{\partial(x_1,x_2,\ldots,x_\rho)}
 =\frac{\Gamma(x)}{G(x)^{\rho+1}}\ne0,
\tag{2}
\end{equation}
donde $G(x)$, como es usual, denota el denominador común de las funciones (1).

\emph{El número $a$ se elige ahora de modo que el producto $G(x)\cdot\Gamma(x)$ no sea divisible por $(x_i-a)$}, entendiendo por $i$ uno de los números $\rho+1,\rho+2,\ldots,n$. Así, para $x_i=a$, el denominador común de las funciones (1) no puede anularse, y las $\rho$ funciones
\begin{equation}
  [g_1(x)]_{x_i=a},\ldots,[g_\rho(x)]_{x_i=a}
\tag{3}
\end{equation}
son igualmente algebraicamente independientes. Como diremos, el rango algebraico del sistema (1) se conserva bajo la sustitución $x_i=a$.

La ecuación irreducible que satisface $h(x)$ respecto de las funciones (1) sea
\begin{equation}
\begin{aligned}
&A_0(g_1,\ldots,g_\rho)\,h(x)^\tau
 +A_1(g_1,\ldots,g_\rho)\,h(x)^{\tau-1}
 +\cdots \\
&\hspace{6em}+A_\tau(g_1,\ldots,g_\rho)=0,
\end{aligned}
\tag{4}
\end{equation}
donde $A_0(g_1,\ldots,g_\rho)$ y $A_\tau(g_1,\ldots,g_\rho)$ no se anulan idénticamente. Para pasar de aquí a una identidad entre polinomios, ponemos
\begin{equation}
  g_i(x)=G_i(x):G(x);\qquad h(x)=H_1(x):H_2(x),
\tag{5}
\end{equation}
y observamos que, por la representación reducida supuesta de $h(x)$, $H_1(x)$ y $H_2(x)$ son primos entre sí. En virtud de (5), (4) se transforma en
\begin{equation}
\begin{aligned}
&B_0(G,G_1,\ldots,G_\rho)\,G(x)^{\sigma_0}H_1(x)^\tau \\
&\quad+B_1(G,G_1,\ldots,G_\rho)\,G(x)^{\sigma_1}H_1(x)^{\tau-1}H_2(x)+\cdots \\
&\quad+B_\tau(G,G_1,\ldots,G_\rho)\,G(x)^{\sigma_\tau}H_2(x)^\tau=0,
\end{aligned}
\tag{6}
\end{equation}
donde los $B_k$ son formas homogéneas en $G,G_i$, y al menos uno de los exponentes no negativos $\sigma_k$ debe ser cero.

Supongamos ahora que $H_1(x)$ fuese divisible por $(x_i-a)$. Entonces, por (6),
\[
  B_\tau(G,G_1,\ldots,G_\rho)\,G(x)^{\sigma_\tau}H_2(x)^\tau
\]
sería también divisible por $(x_i-a)$. Esto está excluido por hipótesis para $G(x)$; por tanto, de la divisibilidad de $B_\tau$ se seguiría también la divisibilidad de
\[
  A_\tau=B_\tau:G^\lambda,
\]
y entre las funciones (3) existiría la relación $A_\tau=0$, contra la independencia algebraica supuesta. Finalmente $H_2(x)$, por ser primo con $H_1(x)$, tampoco puede ser divisible por $(x_i-a)$;\footnote{Esta última conclusión, basada en la representación reducida, ya no sería posible para una sustitución $x_i=a$, $x_k=b$ en la que se conservaran las restantes hipótesis. Entonces $H_1(x)$ y $H_2(x)$ podrían anularse simultáneamente, y el cociente se volvería indeterminado bajo la sustitución, igual a $\frac{0}{0}$; por ejemplo
\[
 h(x)=\frac{x_3(\alpha x_1+\beta x_2)+x_4(\gamma x_1+\delta x_2)}
 {x_3(\alpha' x_1+\beta' x_2)+x_4(\gamma' x_1+\delta' x_2)}
\]
bajo la sustitución $x_3=0$, $x_4=0$.} así, la hipótesis de que $H_1(x)$ sea divisible por $(x_i-a)$ conduce a una contradicción. Del mismo modo se muestra que $H_2(x)$ tampoco puede ser divisible por $(x_i-a)$. Por consiguiente, la función $[h(x)]_{x_i=a}=h'(x)$ no puede ni anularse idénticamente ni volverse indeterminada.

La conclusión precedente puede aplicarse ahora de nuevo al sistema (3) y a la función $h'(x)$, puesta otra vez en representación reducida; continuando el procedimiento se llega finalmente al siguiente lema auxiliar.

\emph{Lema auxiliar: Los dos sistemas}
\[
  g_1(x),\ g_2(x),\ldots,g_\rho(x),
\]
\[
  g_1(x),\ g_2(x),\ldots,g_\rho(x),\ h(x)
\]
\emph{tienen cada uno rango algebraico $\rho$, donde $\rho<n$, y supóngase, por ejemplo, que}
\[
 \frac{\partial(g_1,\ldots,g_\rho)}{\partial(x_1,\ldots,x_\rho)}
 =\frac{\Gamma(x)}{G(x)^{\rho+1}}\ne0.
\]
\emph{Los números}
\[
  a_n,\ a_{n-1},\ldots,a_{\rho+1}
\]
\emph{deben elegirse de tal manera que, bajo la sustitución}
\[
  x_n=a_n,\quad x_{n-1}=a_{n-1},\ldots,\quad x_{\rho+1}=a_{\rho+1},
\]
\emph{el producto $G(x)\cdot\Gamma(x)$ no se anule, es decir, que el rango algebraico del sistema (1) se conserve. En las sustituciones sucesivas}
\[
\begin{aligned}
 [h(x)]_{x_n=a_n}&=h^{(n-1)}(x),\qquad
 [h^{(n-1)}(x)]_{x_{n-1}=a_{n-1}}=h^{(n-2)}(x),\ldots,\\
 [h^{(\rho+1)}(x)]_{x_{\rho+1}=a_{\rho+1}}&=h^*(x_1,x_2,\ldots,x_\rho),
\end{aligned}
\]
\emph{pónganse las funciones $h(x)$, $h^{(n-1)}(x)$, \ldots, $h^{(\rho+1)}(x)$ cada vez en representación reducida. Bajo estas hipótesis, en la función $h^*(x_1,\ldots,x_\rho)$ ni el numerador ni el denominador pueden anularse idénticamente, y la función tampoco puede convertirse en $\frac{0}{0}$.}\footnote{Por ejemplo, en la función $h(x)$ de la nota precedente, la sustitución sucesiva da
\[
 [h(x)]_{x_4=0}=h^{(3)}(x)=\frac{\alpha x_1+\beta x_2}{\alpha' x_1+\beta' x_2},\qquad
 h^*(x_1,x_2)=\frac{\alpha x_1+\beta x_2}{\alpha' x_1+\beta' x_2}.
\]}

\section*{\S\ 3. Aplicación biunívoca de sistemas $\Ssys_{n\rho}$ sobre sistemas $\Ssys^*_{\rho\rho}$.}

Del lema auxiliar obtendremos de inmediato una aplicación de los sistemas $\Ssys_{n\rho}$ sobre sistemas $\Ssys^*_{\rho\rho}$, sustituyendo algunas indeterminadas por números; con ello el rango algebraico $\rho$ se presenta como el número propiamente característico.

Puesto que $\Ssys_{n\rho}$ tiene rango algebraico $\rho$, existen $\rho$ funciones de $\Ssys_{n\rho}$,
\begin{equation}
  g_1(x),\ldots,g_\rho(x),
\tag{1}
\end{equation}
que son algebraicamente independientes. Además, si $f(x)$ denota una función arbitraria de $\Ssys_{n\rho}$, entonces el sistema
\begin{equation}
  g_1(x),\ldots,g_\rho(x),\ f(x)
\tag{2}
\end{equation}
tiene también rango algebraico $\rho$, y los sistemas (1) y (2) satisfacen las condiciones del lema auxiliar.\footnote{Si es necesario, deben renumerarse las indeterminadas.}

Si, por tanto, se determinan los números
\[
  a_n,\ a_{n-1},\ldots,a_{\rho+1}
\]
según el lema auxiliar, de manera que el rango algebraico $\rho$ del sistema (1) se conserve por la sustitución, lo cual es posible de arbitrariamente muchas maneras, entonces en la función definida por
\[
\begin{aligned}
 [f(x)]_{x_n=a_n}&=f^{(n-1)}(x),\qquad
 [f^{(n-1)}(x)]_{x_{n-1}=a_{n-1}}=f^{(n-2)}(x);\ldots,\\
 [f^{(\rho+1)}(x)]_{x_{\rho+1}=a_{\rho+1}}&=f^*(x_1,\ldots,x_\rho),
\end{aligned}
\]
a saber,
\begin{equation}
  f^*(x_1,\ldots,x_\rho),
\tag{3}
\end{equation}
ni el numerador ni el denominador pueden anularse idénticamente.

Hagamos ahora que $f(x)$ recorra todas las funciones, finitas o infinitas, de $\Ssys_{n\rho}$. Consideramos el sistema que consta de la totalidad de todas las funciones (3); tiene las siguientes propiedades:

1. Tiene rango algebraico $\rho$; pues contiene las funciones provenientes de (1),
\[
  g_1^*(x_1,\ldots,x_\rho),\ldots,g_\rho^*(x_1,\ldots,x_\rho),
\]
que, por la elección de $a_n,a_{n-1},\ldots,a_{\rho+1}$, son algebraicamente independientes. Por tanto el sistema debe denotarse por
\[
  \Ssys^*_{\rho\rho}.
\]
\footnote{Correspondientemente, por $f^*(x)$, $g^*(x)$, \ldots, se entenderán en todo momento funciones de aplicación, es decir, funciones de $x_1,\ldots,x_\rho$.}

2. La correspondencia de las funciones $f(x)$ y $f^*(x)$ es sin excepción biunívoca.

En efecto, si $f_1(x)$ y $f_2(x)$ pertenecen a $\Ssys_{n\rho}$, entonces su diferencia
\[
  d(x)=f_1(x)-f_2(x)
\]
es también algebraicamente dependiente del sistema (1); además,
\[
  d^*(x)=f_1^*(x)-f_2^*(x).
\]
Como $d(x)$, junto con el sistema (1), satisface las condiciones del lema auxiliar, de
\[
  d(x)\ne0
\]
se sigue necesariamente
\[
  d^*(x)\ne0;
\]
en otras palabras, a funciones distintas de $\Ssys_{n\rho}$ corresponden funciones distintas de $\Ssys^*_{\rho\rho}$.

3. De toda relación entre funciones de $\Ssys^*_{\rho\rho}$,
\[
  F\bigl(f_1^*(x),f_2^*(x),\ldots,f_k^*(x)\bigr)=0
  \quad\text{idénticamente en }x_1,\ldots,x_\rho,
\]
se sigue, para las funciones unívocamente correspondientes de $\Ssys_{n\rho}$, que
\[
  F\bigl(f_1(x),f_2(x),\ldots,f_k(x)\bigr)=0
  \quad\text{idénticamente en }x_1,\ldots,x_n.
\]
Pues con $f_1(x),f_2(x),\ldots,f_k(x)$ también toda combinación racional de estas funciones es algebraicamente dependiente del sistema (1).

Además, esto se sigue de
\begin{equation}
\begin{aligned}
 h(x)&=F\bigl(f_1(x),\ldots,f_k(x)\bigr):\\
 h^*(x)&=F\bigl(f_1^*(x),\ldots,f_k^*(x)\bigr).
\end{aligned}
\tag{4}
\end{equation}
Por el lema auxiliar, por tanto,
\[
  h(x)\ne0\quad\text{implica}\quad h^*(x)\ne0;
\]
q.e.d.

En resumen obtenemos:

\emph{Teorema I: A todo sistema $\Ssys_{n\rho}$ de funciones racionales (finitas o infinitas) se le puede asignar, de arbitrariamente muchas maneras, un sistema biunívoco $\Ssys^*_{\rho\rho}$, de tal modo que todas las relaciones existentes entre funciones de $\Ssys^*_{\rho\rho}$ se conservan también en $\Ssys_{n\rho}$ y recíprocamente. En particular, por (4), a un cuerpo $\Kfield_{n\rho}$ corresponde de nuevo un cuerpo $\Kfield^*_{\rho\rho}$.}

Del Teorema I extraemos la consecuencia que simplifica todas las demostraciones posteriores: \emph{las propiedades de los sistemas $\Ssys_{n\rho}$ que se caracterizan por finitas o infinitas relaciones entre funciones del sistema valen para los sistemas $\Ssys_{n\rho}$ tan pronto como valen para un sistema de aplicación $\Ssys^*_{\rho\rho}$.}

Llamamos brevemente a esta consecuencia el ``principio de transferencia''.

\providecommand{\Lsys}{\mathfrak{L}}
\providecommand{\Jdom}{\mathfrak{J}}

\clearpage
El ejemplo más sencillo de esta aplicación lo proporciona la teoría de los invariantes algebraicos. En efecto, mediante una transformación lineal se pueden fijar numéricamente algunos coeficientes de la forma fundamental de tal modo que todas las relaciones válidas para la forma especial sigan siendo válidas en general.

\section*{\S\ 4. Base racional de los cuerpos $\Kfield_{n\rho}$.}

En las investigaciones siguientes, agrupadas alrededor de conceptos de base, utilizaremos de manera constante el ``principio de transferencia'' de \S\ 3. Primero demostramos la existencia de la base para cuerpos, pero enunciamos las definiciones en general:

\emph{Definición III: Un número finito de funciones de $\Ssys_{n\rho}$ se llama base racional si toda función de $\Ssys_{n\rho}$ puede representarse como combinación racional de ese número finito, con coeficientes de $\Omega$.}

La base racional debe contener $\rho$ funciones algebraicamente independientes, puesto que el rango algebraico de un sistema derivado de una base racional es igual al rango algebraico de la base racional.

Mostramos primero que la demostración de existencia de la base racional admite el principio de transferencia. En efecto, la Definición III puede formularse como sigue:

Las funciones de $\Ssys_{n\rho}$
\[
  g_1(x),\ g_2(x),\ldots, g_k(x)
\]
forman una base racional si y sólo si se satisfacen las infinitas relaciones
\begin{equation}
  f(x)=\gamma\bigl(g_1(x),\ldots,g_k(x)\bigr),
\tag{1}
\end{equation}
donde $f(x)$ recorre todas las funciones de $\Ssys_{n\rho}$ y $\gamma$ denota una función racional de $k$ argumentos, variable con $f$.

Así, si en el sistema de aplicación $\Ssys^*_{\rho\rho}$ se da una base racional por
\[
  g_1^*(x),\ g_2^*(x),\ldots,g_k^*(x),
\]
lo cual entraña las relaciones
\[
  f^*(x)=\gamma\bigl(g_1^*(x),\ldots,g_k^*(x)\bigr),
\]
entonces, por el Teorema I, las relaciones (1) se cumplen para $\Ssys_{n\rho}$. Así, el \emph{principio de transferencia para la base racional} dice:

\emph{De la existencia de una base racional para un sistema de aplicación $\Ssys^*_{\rho\rho}$ se sigue la base racional para el sistema original $\Ssys_{n\rho}$.}\footnote{El principio de transferencia se formula correspondientemente para toda cuestión de base tratada aquí.}

Para demostrar la existencia de la base racional de todos los cuerpos $\Kfield_{n\rho}$ podemos, por tanto, restringirnos a cuerpos $\Kfield^*_{\rho\rho}$; aquí una generalización directa de un argumento empleado por E. Steinitz conduce al objetivo.\footnote{E. Steinitz, \emph{Algebraische Theorie der Körper}, \S\ 24, Crelle's Journal 137 (1909).}

Sea
\begin{equation}
  g_1^*(x_1,\ldots,x_\rho),\ldots,
  g_\rho^*(x_1,\ldots,x_\rho)
\tag{2}
\end{equation}
un sistema de $\rho$ funciones algebraicamente independientes de $\Kfield^*_{\rho\rho}$. Eliminando $x_1,\ldots,x_\rho$ se obtienen $\rho$ relaciones:
\begin{equation}
\begin{aligned}
  F_1\bigl(g_1^*(x),\ldots,g_\rho^*(x),x_1\bigr)&=0,\ldots,\\
  F_\rho\bigl(g_1^*(x),\ldots,g_\rho^*(x),x_\rho\bigr)&=0,
\end{aligned}
\tag{3}
\end{equation}
donde en cada relación $F_i=0$ aparece efectivamente la indeterminada $x_i$; de lo contrario, contra la hipótesis, habría una dependencia algebraica entre las funciones (2). Las relaciones (3) dicen que
\[
  x_1,x_2,\ldots,x_\rho
\]
son elementos algebraicos respecto de
\[
  \Omega\bigl(g_1^*(x),\ldots,g_\rho^*(x)\bigr).
\]
Por consiguiente, el cuerpo que surge por adjunción de $x_1,\ldots,x_\rho$,
\[
  \Omega\bigl(g_1^*,\ldots,g_\rho^*,x_1,
  \ldots,x_\rho\bigr)=\Omega(x_1,
  \ldots,x_\rho),
\]
es una extensión algebraica finita de $\Omega(g_1^*,\ldots,g_\rho^*)$. Por un teorema conocido,\footnote{Compárese, por ejemplo, Weber, \emph{Lehrbuch d. Algebra} 1, \S\ 144 (1.ª ed.) o \S\ 55 (edición pequeña).} $\Kfield^*_{\rho\rho}$, como cuerpo intermedio entre $\Omega(g_1^*,\ldots,g_\rho^*)$ y $\Omega(x_1,\ldots,x_\rho)$, es él mismo una extensión algebraica de $\Omega(g_1^*,\ldots,g_\rho^*)$, mientras que $\Omega(x_1,\ldots,x_\rho)$ es una extensión algebraica de $\Kfield^*_{\rho\rho}$. Así, $\Kfield^*_{\rho\rho}$ se obtiene de $\Omega(g_1^*,\ldots,g_\rho^*)$ por adjunción de un número finito de elementos de $\Kfield^*_{\rho\rho}$, o también por adjunción de una función convenientemente elegida; es decir, existe una base racional. Por el principio de transferencia se obtiene entonces

\emph{Teorema II: Todo cuerpo $\Kfield_{n\rho}$ posee una base racional.}

\section*{\S\ 5. Forma de involución y base de involución de los cuerpos $\Kfield_{n\rho}$.}

Del argumento precedente obtenemos además una base racional determinada por el cuerpo mismo. La indicamos primero para los cuerpos $\Kfield^*_{\rho\rho}$.

Sea $\Kfield^*_{\rho\rho}$ un cuerpo entre $\Omega$ y $\Omega(x_1,\ldots,x_\rho)$. Como en \S\ 4, se eligen $\rho$ funciones algebraicamente independientes de $\Kfield^*_{\rho\rho}$,
\[
  g_1^*(x),\ldots,g_\rho^*(x),
\]
y se forma la extensión algebraica finita $\Omega(x_1,\ldots,x_\rho)$ de $\Omega(g_1^*,\ldots,g_\rho^*)$. Sea $i$ el grado de $\Omega(x_1,\ldots,x_\rho)$ respecto de $\Kfield^*_{\rho\rho}$.

Introducimos una forma lineal
\[
  u(x)=u_1x_1+\cdots+u_\rho x_\rho
\]
con coeficientes indeterminados $u_1,\ldots,u_\rho$, y sea
\begin{equation}
 \Phi^*(z,u)=z^i+\sum_{\alpha+\alpha_1+\cdots+\alpha_\rho=i}
 g^*_{\alpha\alpha_1\ldots\alpha_\rho}(x)
 z^\alpha u_1^{\alpha_1}\cdots u_\rho^{\alpha_\rho}
\tag{1}
\end{equation}
la ecuación irreducible de $u(x)$ respecto de $\Kfield^*_{\rho\rho}$.\footnote{Noether escribe la suma sobre todos los sistemas de índices con suma total $i$; el término principal $z^i$ se ha separado.} Los coeficientes
\[
  g^*_{\alpha\alpha_1\ldots\alpha_\rho}(x)
\]
pertenecen a $\Kfield^*_{\rho\rho}$ y están determinados por el cuerpo. Llamaremos a $\Phi^*(z,u)$ la forma de involución del cuerpo.

Afirmamos que estos coeficientes forman una base racional. En efecto, la ecuación (1), vista como ecuación de $z$, es irreducible respecto de $\Kfield^*_{\rho\rho}$ y por tanto tanto más irreducible respecto del subcuerpo
\[
  \Omega\bigl(g^*_{\alpha\alpha_1\ldots\alpha_\rho}(x)\bigr).
\]
Por consiguiente $\Phi^*(z,u)$ da la ecuación irreducible de $u(x)$ respecto de ese subcuerpo. De ello se sigue que $\Omega(x_1,\ldots,x_\rho)$ es una extensión algebraica --- de grado $i$ --- de $\Omega(g^*_{\alpha\alpha_1\ldots\alpha_\rho}(x))$. Como $\Kfield^*_{\rho\rho}$ está entre este subcuerpo y $\Omega(x_1,\ldots,x_\rho)$, la igualdad de grados implica, como es sabido, la igualdad de los cuerpos. Por el principio de transferencia, los coeficientes de $\Phi(z,u)$ suministran correspondientemente una base racional de $\Kfield_{n\rho}$. Así tenemos:

\emph{Teorema III: Los coeficientes de la forma de involución constituyen una base racional determinada por el cuerpo, la base de involución; para cuerpos $\Kfield^*_{\rho\rho}$ esta base está determinada de manera única.}\footnote{Para cuerpos $\Kfield_{n\rho}$ la forma de involución puede depender de la aplicación; la unicidad puede obtenerse mediante una aplicación con coeficientes indeterminados.}

Añadimos una interpretación irracional de este hecho, que establece la conexión con la teoría geométrica de las involuciones.

La factorización de $\Phi^*(z,u)$ en factores lineales viene dada en un cuerpo de extensión por
\begin{equation}
\begin{aligned}
 \Phi^*(z,u)&=(z-u_1x_1-\cdots-u_\rho x_\rho)
 (z-u_1x_1^{(1)}-\cdots-u_\rho x_\rho^{(1)})\cdots\\
 &\qquad\cdot
 (z-u_1x_1^{(i-1)}-\cdots-u_\rho x_\rho^{(i-1)}).
\end{aligned}
\tag{4}
\end{equation}
La comparación con (1) muestra que las funciones
\[
  g^*_{\alpha\alpha_1\ldots\alpha_\rho}(x_1,
  \ldots,x_\rho)\quad(\alpha+\alpha_1+\cdots+\alpha_\rho=i)
\]
son idénticas a las funciones simétricas elementales de las filas de cantidades
\begin{equation}
\begin{array}{cccc}
  x_1 & x_2 & \cdots & x_\rho\\
  x_1^{(1)} & x_2^{(1)} & \cdots & x_\rho^{(1)}\\
  \vdots & \vdots & & \vdots\\
  x_1^{(i-1)} & x_2^{(i-1)} & \cdots & x_\rho^{(i-1)}.
\end{array}
\tag{5}
\end{equation}
Así, toda función de $\Kfield^*_{\rho\rho}$ es una combinación racional de estas funciones elementales, simétricas en las filas (5), y recíprocamente toda función simétrica de las filas (5) pertenece a $\Kfield^*_{\rho\rho}$.

Así, cuando
\[
  f^*(x)=\frac{F^*(x)}{H^*(x)}
\]
recorre todas las funciones de $\Kfield^*_{\rho\rho}$, existe un sistema de infinitas relaciones:
\begin{equation}
  \frac{F^*(x)}{H^*(x)}=
  \frac{F^*(x^{(1)})}{H^*(x^{(1)})}=\cdots=
  \frac{F^*(x^{(i-1)})}{H^*(x^{(i-1)})},
\tag{6}
\end{equation}
o, en forma resumida,
\[
  F^*(x^{(k)})H^*(x^{(l)})-
  F^*(x^{(l)})H^*(x^{(k)})=0
  \quad(k,l=0,1,\ldots,i-1),
\]
donde ni $F^*(x^{(k)})$ ni $H^*(x^{(k)})$ se anula idénticamente, pues son conjugadas de $F^*(x)$ y $H^*(x)$ y formalmente simétricas.

Recíprocamente, para toda fila de cantidades $\xi$ que satisfaga las infinitas relaciones
\begin{equation}
  F^*(x)H^*(\xi)-H^*(x)F^*(\xi)=0,
  \qquad H^*(\xi)\ne0,
\tag{7}
\end{equation}
se obtiene la pertenencia a las filas de cantidades (5).

En efecto, si $G^*(x)$ denota el denominador común de los coeficientes de $\Phi^*(z,u)$, entonces por la hipótesis (7) la función
\[
  G^*(\xi)\cdot\Phi^*(z,u;\xi),
\]
que también es entera en $\xi$, no se anula idénticamente, es decir, los coeficientes de $\Phi^*(z,u;\xi)$ no se vuelven indeterminados, y tiene el cero
\[
  z=u_1\xi_1+u_2\xi_2+\cdots+u_\rho\xi_\rho.
\]
Por (7), $\Phi^*(z,u;x)$ tiene por tanto también el cero $z=u(\xi)$; es decir, $\xi$ pertenece a las filas de cantidades (5). La afirmación correspondiente vale, por (6), si $\xi$ satisface un sistema de ecuaciones respecto de una fila conjugada de cantidades:
\[
  F^*(x^{(k)})H^*(\xi)-H^*(x^{(k)})F^*(\xi)=0;
  \qquad H^*(\xi)\ne0.
\]
Así, las filas de cantidades (5) se definen como las soluciones $\xi$ del sistema de ecuaciones
\[
  F^*(x)H^*(\xi)-H^*(x)F^*(\xi)=0;
  \qquad H^*(\xi)\ne0,
\]
donde $F^*(x):H^*(x)$ recorre todas las funciones de $\Kfield^*_{\rho\rho}$; aquí la fila $x$ puede reemplazarse también por cualquier fila conjugada.

Geométricamente esto dice lo siguiente, interpretando cada fila $x$ como un punto: las soluciones de (7) representan, en un espacio lineal de dimensión $\rho$, $\infty^\rho$ grupos de $i$ puntos cada uno, de tal manera que cualquier punto del grupo determina unívocamente el grupo individual. Tal sistema de grupos de puntos se llama una ``\emph{involución en un espacio lineal}''.\footnote{Las soluciones excluidas de (7), para las cuales $F^*(\xi)=0$, $H^*(\xi)=0$, y asimismo las soluciones excluidas del sistema de ecuaciones correspondiente a la base de involución, se llaman puntos fundamentales de la involución; esto incluye también los que surgen por homogenización, los puntos ``situados en el infinito''.}

Correspondientemente, un cuerpo $\Kfield_{n\rho}$ caracteriza una involución en un espacio lineal compuesto de curvas, superficies, etcétera.

Por lo precedente, el grado $i$ de $\Omega(x_1,\ldots,x_\rho)$ respecto de $\Kfield^*_{\rho\rho}$ es igual al número de sistemas de soluciones de (7) que satisfacen la condición lateral $H^*(\xi)\ne0$, a los que podemos llamar los sistemas de soluciones de $\Kfield^*_{\rho\rho}$. Por tanto, el argumento empleado para demostrar la existencia de la base de involución permite también la siguiente formulación irracional:

``Si $\mathfrak{L}^*$ es un divisor de $\Kfield^*_{\rho\rho}$ y todo sistema de soluciones del cuerpo $\mathfrak{L}^*$ es también un sistema de soluciones de $\Kfield^*_{\rho\rho}$, entonces $\mathfrak{L}^*$ y $\Kfield^*_{\rho\rho}$ son idénticos.''

El enunciado correspondiente vale, por el principio de transferencia, para divisores $\mathfrak{L}$ de $\Kfield_{n\rho}$.

\section*{\S\ 6. Base mínima de los cuerpos $\Kfield_{n\rho}$.}

Este parágrafo da una visión de conjunto de teoremas conocidos, pero en una formulación ampliada por el principio de transferencia y por la existencia de la base racional.

\emph{Definición IV: Una base racional de $\Kfield_{n\rho}$ se llama base mínima si contiene sólo el número mínimo $\rho$ de funciones, que entonces deben ser algebraicamente independientes.}

De los teoremas de Lüroth, Castelnuovo y Enriques\footnote{J. Lüroth: Beweis eines Satzes über rationale Kurven, \emph{Math. Ann.} 9 (1875). --- G. Castelnuovo: Sulla razionalità delle involuzioni piane, \emph{Math. Ann.} 44 (1893). --- F. Enriques: Sopra una involuzione non razionale dello spazio, \emph{Rend. Acc. Linc.}, vol. 21, 21 de enero de 1912.} se obtienen los hechos:

I. Todo cuerpo $\Kfield_{n1}$ posee una base mínima, la función de Lüroth del cuerpo, determinada de modo único salvo transformación lineal fraccionaria.

II. Para cuerpos $\Kfield_{n2}$ existe siempre una base mínima, y en general sólo entonces, si se permite una extensión algebraica del cuerpo de coeficientes $\Omega$.

III. Para cuerpos $\Kfield_{n\rho}$ ($\rho\ge3$), en general no existe base mínima. Todavía no se ha intentado una caracterización de los cuerpos especiales $\Kfield_{n\rho}$ que tienen base mínima.

Indicamos brevemente la demostración del teorema de Lüroth para cuerpos $\Kfield^*_{11}$ dada por E. Steinitz:\footnote{Loc. cit., \S\ 24.}

Sea $\Phi^*(z,u)$ la forma de involución de $\Kfield^*_{11}$, y sea $\psi^*(x)$ cualquier coeficiente de $\Phi^*(z,u)$ que contenga efectivamente la indeterminada $x$. Resulta que el grado de $\Omega(x)$ respecto de $\Omega(\psi^*(x))$ es igual al grado de $\Omega(x)$ respecto de $\Kfield^*_{11}$, de donde se sigue la identidad de los cuerpos. Además, si $\Omega(\chi^*(x))$ es igual a $\Omega(\psi^*(x))$, entonces la dependencia racional mutua de $\chi^*(x)$ y $\psi^*(x)$ implica la dependencia lineal fraccionaria de las dos funciones.

Por el principio de transferencia, el teorema de Lüroth admite por tanto también la formulación:

\emph{La base mínima de los cuerpos $\Kfield_{n1}$ consta de cualquier función de la base de involución; y cualesquiera dos funciones de la base de involución de $\Kfield_{n1}$ dependen una de otra por una transformación lineal fraccionaria.}

G. Castelnuovo demuestra el Hecho II, por medio de la teoría geométrica de las involuciones, para cuerpos $\Kfield^*_{22}$ derivados de una base racional de tres funciones. Puesto que el principio de transferencia sigue siendo válido también bajo una extensión algebraica finita del cuerpo de coeficientes $\Omega$, la existencia de la base racional completa la demostración para todos los cuerpos $\Kfield_{n2}$.

Respecto de III, debe observarse que Enriques construye una involución no racional en el espacio tridimensional; es decir, un cuerpo $\Kfield_{33}$ que no tiene base mínima.


\providecommand{\Gdom}{\mathfrak{G}}
\providecommand{\Omegaint}{[\Omega]}
\providecommand{\eps}{\varepsilon}

\clearpage
\section*{\S\ 7. Sistema arbitrario $\Ssys_{n\rho}$, familia lineal $\Lsys_{n\rho}$ y dominio íntegro $\Jdom_{n\rho}$ de funciones racionales. Existencia de la base racional.}

A partir de la existencia de la base racional de los cuerpos $\Kfield_{n\rho}$ obtendremos ahora la base racional para sistemas arbitrarios de funciones racionales y racionales enteras. Ordenamos estos sistemas según el modo en que surgen sucesivamente unos de otros mediante las operaciones elementales de adición, multiplicación y división.

I. Como siempre, $\Ssys_{n\rho}$ designa un sistema arbitrario de funciones racionales (o racionales enteras) de $n$ indeterminadas y de rango algebraico $\rho$. Las cantidades de $\Omega$ se cuentan como pertenecientes al sistema.

II. El sistema $\Lsys_{n\rho}$ se llama una \emph{familia lineal} si satisface las condiciones siguientes:

1. Junto con $f(x)$, $\Lsys_{n\rho}$ contiene siempre también $c\cdot f(x)$, donde $c$ es una cantidad arbitraria de $\Omega$;

2. junto con $f(x)$ y $g(x)$, $\Lsys_{n\rho}$ contiene siempre también $f(x)+g(x)$.

III. La familia lineal $\Jdom_{n\rho}$ se llama un \emph{dominio íntegro} bajo la condición ulterior:

3. junto con $f(x)$ y $g(x)$, $\Jdom_{n\rho}$ contiene siempre también $f(x)\cdot g(x)$.

IV. El dominio íntegro $\Kfield_{n\rho}$ se llama un \emph{cuerpo} bajo la condición ulterior:

4. junto con $f(x)$ y $g(x)$ ---donde $g(x)\ne0$---, $\Kfield_{n\rho}$ contiene siempre también el cociente $f(x):g(x)$.

Las condiciones 1 a 4 son las indicadas en \S\ 1 para el cuerpo.\footnote{$f(x)$ y $g(x)$ pueden ser también de grado cero, es decir, cantidades de $\Omega$.}

Estos dominios deben derivarse ahora sucesivamente unos de otros.

a) Un sistema arbitrario $\Ssys_{n\rho}$ se amplía a una familia lineal $\Lsys_{n\rho}$ mediante la exigencia siguiente:

$\Lsys_{n\rho}$ contiene, además de $\Ssys_{n\rho}$, todas las funciones $h(x)$, y sólo éstas, que pueden representarse linealmente con coeficientes de $\Omega$ por medio de un número finito de funciones de $\Ssys_{n\rho}$:
\[
  h(x)=c_1f_1(x)+c_2f_2(x)+\cdots+c_kf_k(x),
\]
donde $f_1(x),\ldots,f_k(x)$ recorren todas las funciones de $\Ssys_{n\rho}$, y $c_1,
\ldots,c_k$ todas las cantidades de $\Omega$.

En efecto, al añadir combinaciones racionales de funciones del sistema se conserva el rango algebraico de un sistema. $\Lsys_{n\rho}$ satisface las condiciones 1 y 2, y por tanto es una familia lineal. Además, toda familia lineal que contiene $\Ssys_{n\rho}$ contiene también $\Lsys_{n\rho}$ por 1 y 2; por consiguiente $\Lsys_{n\rho}$ es la menor familia lineal que contiene $\Ssys_{n\rho}$.

b) Correspondientemente, de una familia lineal $\Lsys_{n\rho}$ surge el menor dominio íntegro que la contiene, $\Jdom_{n\rho}$, mediante la exigencia siguiente:

$\Jdom_{n\rho}$ contiene, además de $\Lsys_{n\rho}$, todas las funciones $h(x)$, y sólo éstas, que pueden representarse como una \emph{combinación racional entera} ---con coeficientes de $\Omega$--- de un número finito de funciones de $\Lsys_{n\rho}$, $f_1(x),\ldots,f_k(x)$, donde $f_1(x),\ldots,f_k(x)$ recorren todas las funciones de $\Lsys_{n\rho}$.

c) Finalmente, de un dominio íntegro $\Jdom_{n\rho}$ surge el menor cuerpo que lo contiene, $\Kfield_{n\rho}$, mediante la exigencia siguiente: $\Kfield_{n\rho}$ contiene, además de $\Jdom_{n\rho}$, todas las funciones $h(x)$, y sólo éstas, que pueden representarse como cocientes de dos funciones de $\Jdom_{n\rho}$.

Si se reúnen las tres ampliaciones en un solo paso, resulta:

\emph{Todo sistema $\Ssys_{n\rho}$ puede ampliarse al menor cuerpo que lo contiene, $\Kfield_{n\rho}$, mediante la exigencia siguiente:}

\emph{$\Kfield_{n\rho}$ contiene, además de $\Ssys_{n\rho}$, todas las funciones $h(x)$, y sólo éstas, que pueden representarse como combinaciones racionales ---con coeficientes de $\Omega$--- de un número finito de funciones de $\Ssys_{n\rho}$, $f_1(x),\ldots,f_k(x)$, donde $f_1(x),\ldots,f_k(x)$ recorren todas las funciones de $\Ssys_{n\rho}$.}

De este último hecho se sigue inmediatamente la existencia de la base racional para dominios arbitrarios:

Sea $\Kfield_{n\rho}$ el menor cuerpo que contiene $\Ssys_{n\rho}$, y sea una base racional de $\Kfield_{n\rho}$ dada por
\begin{equation}
  h_1(x),\ h_2(x),\ldots,h_\sigma(x).
\tag{1}
\end{equation}
Por la definición de $\Kfield_{n\rho}$ existen relaciones
\begin{equation}
\begin{aligned}
 h_1(x)&=r_1\bigl(g_1(x),\ldots,g_\lambda(x)\bigr);\ldots;\\
 h_\sigma(x)&=r_\sigma\bigl(g_\mu(x),\ldots,g_\nu(x)\bigr),
\end{aligned}
\tag{2}
\end{equation}
donde $r_1,
\ldots,r_\sigma$ son funciones racionales de sus argumentos, y
\begin{equation}
  g_1(x),\ldots,g_\lambda(x),\ldots,
  g_\mu(x),\ldots,g_\nu(x)
\tag{3}
\end{equation}
pertenecen todas a $\Ssys_{n\rho}$. A causa de las relaciones (2), (3) representa, además de (1), también una base racional de $\Kfield_{n\rho}$; y, puesto que las funciones (3) pertenecen a $\Ssys_{n\rho}$, es al mismo tiempo una base racional de $\Ssys_{n\rho}$. Tenemos por tanto:

\emph{Teorema IV: Todo sistema arbitrario $\Ssys_{n\rho}$ de funciones racionales (o racionales enteras) posee una base racional, compuesta por un número finito de funciones del sistema, de rango algebraico $\rho$.}

Sea ahora, en particular, $\Lsys_{n\rho}$ una familia lineal, y sea
\begin{equation}
  g_1(x),\ldots,g_\rho(x),\ g_{\rho+1}(x),\ldots,g_\nu(x)
\tag{4}
\end{equation}
una base racional de $\Lsys_{n\rho}$. Supongamos, por ejemplo, que
\[
  g_1(x),\ldots,g_\rho(x)
\]
son algebraicamente independientes. Si se pone entonces
\[
  g(x)=c_1g_{\rho+1}(x)+c_2g_{\rho+2}(x)+\cdots+c_{\nu-\rho}g_\nu(x),
\]
los $c_i$ pueden elegirse como números de $\Omega$ de tal modo que
\begin{equation}
  g_1(x),\ldots,g_\rho(x),\ g(x)
\tag{5}
\end{equation}
sea una base racional del menor cuerpo que contiene la familia, $\Kfield_{n\rho}$. Pero como $\Lsys_{n\rho}$ es una familia lineal, $g(x)$ pertenece a $\Lsys_{n\rho}$, y (5) representa una base racional de $\Lsys_{n\rho}$; es decir:

\emph{Teorema V: La base racional de toda familia lineal $\Lsys_{n\rho}$ de funciones racionales (o racionales enteras) ---por tanto, en particular, la base racional de todo dominio íntegro $\Jdom_{n\rho}$ y de todo cuerpo $\Kfield_{n\rho}$--- puede elegirse de modo especial para que contenga a lo sumo $\rho+1$ funciones (y al menos $\rho$).}

La cota para el número de funciones en la base racional encontrada en \S\ 4 para los cuerpos descansa, por tanto, en la propiedad del cuerpo de ser una familia lineal, mientras que la restricción a $\rho$ funciones en caso de existir una base mínima utiliza todas las hipótesis del cuerpo.

Queda por señalar que, por \S\ 3 (4), todas las propiedades características de los sistemas y de los sistemas mínimos que los contienen se conservan bajo la aplicación.

\section*{\S\ 8. Base de involución de los dominios íntegros $\Jdom_{n\rho}$ formados por polinomios.}

En los parágrafos siguientes tratamos sistemas $\Ssys_{n\rho}$ cuyos elementos son funciones racionales enteras (polinomios); primero, dominios íntegros $\Jdom_{n\rho}$ formados por polinomios. Para estos dominios $\Jdom_{n\rho}$ deben fijarse ante todo los conceptos de divisibilidad.

Sean $F(x),G(x)$ dos polinomios de $\Jdom_{n\rho}$, y sea $L(x)$ un divisor común de estos polinomios, de modo que se tenga
\[
  F(x)=L(x)\cdot F_1(x),\qquad G(x)=L(x)\cdot G_1(x).
\]
$L(x)$ se llama un \emph{divisor común en $\Jdom_{n\rho}$} si y sólo si los tres polinomios $L(x),F_1(x),G_1(x)$ pertenecen a $\Jdom_{n\rho}$.

Dos polinomios de $\Jdom_{n\rho}$ se llaman \emph{primos entre sí en $\Jdom_{n\rho}$} si no poseen ningún divisor común en $\Jdom_{n\rho}$.

Sea $\Kfield_{n\rho}$ el menor cuerpo que contiene $\Jdom_{n\rho}$. Entonces, por la definición, toda función de $\Kfield_{n\rho}$ posee una representación
\[
  f(x)=F_1(x):F_2(x),
\]
donde $F_1(x),F_2(x)$ pertenecen a $\Jdom_{n\rho}$ y son primos entre sí en $\Jdom_{n\rho}$. Correspondientemente, varias funciones de $\Kfield_{n\rho}$ pueden ponerse sobre un denominador común en $\Jdom_{n\rho}$:
\[
  f_1(x)=\frac{F_1(x)}{F(x)},\ldots,
  f_k(x)=\frac{F_k(x)}{F(x)}.
\]
$F(x)$ se llama \emph{un mínimo denominador común en $\Jdom_{n\rho}$} si y sólo si los polinomios de $\Jdom_{n\rho}$
\[
  F(x),F_1(x),\ldots,F_k(x),
\]
son primos entre sí en $\Jdom_{n\rho}$.

Tal representación puede ser posible de maneras distintas, pues en $\Jdom_{n\rho}$ no valen en general las leyes de la descomposición única en funciones irreducibles.

Investigamos ahora el significado de la base de involución de $\Kfield_{n\rho}$ para $\Jdom_{n\rho}$, transfiriendo los conceptos de forma de involución y base de involución a los dominios íntegros. Sea
\[
  \Phi(z,u)=z^i+\sum{}' g_{\alpha\alpha_1\ldots\alpha_\rho}(x)
  z^\alpha u_1^{\alpha_1}\cdots u_\rho^{\alpha_\rho},
  \qquad
  (\alpha+\alpha_1+\cdots+\alpha_\rho=i)
\]
una forma de involución de $\Kfield_{n\rho}$, y sea $G(x)$ un mínimo denominador común en $\Jdom_{n\rho}$ de los coeficientes de $\Phi(z,u)$. Por multiplicación por $G(x)$, $\Phi(z,u)$ pasa a una forma $\Psi(z,u)$ cuyos coeficientes son polinomios de $\Jdom_{n\rho}$ y son primos entre sí en $\Jdom_{n\rho}$:
\[
  G(x)\cdot\Phi(z,u)=\Psi(z,u)
  =G(x)z^i+
  \sum{}' G_{\alpha\alpha_1\ldots\alpha_\rho}(x)
  z^\alpha u_1^{\alpha_1}\cdots u_\rho^{\alpha_\rho},
  \qquad
  (\alpha+\alpha_1+\cdots+\alpha_\rho=i).
\]

Toda forma $\Psi(z,u)$ que surge de una forma de involución $\Phi(z,u)$ de $\Kfield_{n\rho}$ por multiplicación por un polinomio de $\Jdom_{n\rho}$, de tal manera que los coeficientes de $\Psi(z,u)$ sean polinomios de $\Jdom_{n\rho}$ y primos entre sí en $\Jdom_{n\rho}$, se llamará una \emph{forma de involución} de $\Jdom_{n\rho}$; y los coeficientes de $\Psi(z,u)$ se llamarán una \emph{base de involución} correspondiente.

Ante todo es claro, por el significado de la base de involución de $\Kfield_{n\rho}$, que toda base de involución de $\Jdom_{n\rho}$ es al mismo tiempo una base racional. Para la investigación ulterior consideramos un dominio de aplicación $\Jdom^*_{\rho\rho}$ para cuyo menor cuerpo contenedor $\Kfield^*_{\rho\rho}$ la ecuación irreducible con cero
\[
  z=u_1x_1+\cdots+u_\rho x_\rho
\]
está dada por $\Phi^*(z,u)=0$.

Según \S\ 5, los coeficientes de
\[
  \Phi^*(z,u)=z^i+\sum{}' g^*_{\alpha\alpha_1\ldots\alpha_\rho}(x)
  z^\alpha u_1^{\alpha_1}\cdots u_\rho^{\alpha_\rho},
  \qquad
  (\alpha+\alpha_1+\cdots+\alpha_\rho=i)
\]
representan las funciones simétricas elementales de las filas de cantidades conjugadas con respecto a $\Kfield^*_{\rho\rho}$,
\[
  x,\quad x^{(1)},\ldots,x^{(i-1)}.
\]
Sea $F^*(x)$ un polinomio cualquiera de $\Kfield^*_{\rho\rho}$. Puesto que
\[
  F^*(x)=\frac1i\Bigl(F^*(x)+F^*(x^{(1)})+
  \cdots+F^*(x^{(i-1)})\Bigr),
\]
se sigue que $F^*(x)$ es una función simétrica entera de estas filas de cantidades; por consiguiente puede expresarse como combinación racional entera de las funciones simétricas elementales, es decir, de las funciones
\[
  g^*_{\alpha\alpha_1\ldots\alpha_\rho}(x).
\]
Si ahora
\[
  G^*(x)\cdot\Phi^*(z,u)=\Psi^*(z,u)
  =G^*(x)z^i+
  \sum{}'G^*_{\alpha\alpha_1\ldots\alpha_\rho}(x)
  z^\alpha u_1^{\alpha_1}\cdots u_\rho^{\alpha_\rho},
  \qquad
  (\alpha+\alpha_1+\cdots+\alpha_\rho=i)
\]
es una forma de involución de $\Jdom^*_{\rho\rho}$, entonces todo polinomio de $\Kfield^*_{\rho\rho}$, y en particular todo polinomio de $\Jdom^*_{\rho\rho}$, es una combinación racional entera de los cocientes
\[
  G^*_{\alpha\alpha_1\ldots\alpha_\rho}(x):G^*(x).
\]

Teniendo en cuenta el principio de transferencia, esto da:

\emph{Teorema VI: Para todo dominio íntegro de polinomios $\Jdom_{n\rho}$ existe al menos una base racional distinguida, la base de involución, tal que en la representación racional de todos los polinomios de $\Jdom_{n\rho}$ por medio de las funciones de la base de involución sólo aparece en el denominador una potencia de una función fija (el coeficiente principal de la forma de involución correspondiente).}

\section*{\S\ 9. Dominios relativamente enteros $\Gdom_{n\rho}$ formados por polinomios.}

En lo que sigue consideramos dominios íntegros especiales formados por polinomios, que constituyen una etapa intermedia entre los dominios íntegros y los cuerpos.

Un dominio íntegro de polinomios $\Gdom_{n\rho}$ se llama ``dominio relativamente entero'' bajo la condición:

\emph{4a) $\Gdom_{n\rho}$ contiene, junto con $F(x)$ y $G(x)$ ---donde $G(x)\ne0$---, el cociente $F(x):G(x)$ siempre y sólo cuando este cociente es entero en las indeterminadas, es decir, es un polinomio.}

Primero debe demostrarse la identidad de los dominios relativamente enteros con la totalidad de los polinomios de un cuerpo $\Kfield_{n\sigma}$ ($\sigma\ge\rho$) de funciones racionales.

Sea $\Kfield_{n\rho}$ el menor cuerpo que contiene $\Gdom_{n\rho}$. Para toda función $f(x)$ de $\Kfield_{n\rho}$ hay una representación como cociente de dos polinomios de $\Gdom_{n\rho}$:
\[
  f(x)=F_1(x):F_2(x),
\]
donde se admiten también polinomios de grado cero. Tan pronto como $f(x)$ sea un polinomio, pertenece a $\Gdom_{n\rho}$ por la condición 4a); y, puesto que $\Gdom_{n\rho}$ no puede contener otros polinomios, \emph{todo dominio relativamente entero $\Gdom_{n\rho}$ es idéntico con la totalidad de los polinomios del menor cuerpo que lo contiene.}

Recíprocamente, sea $\Kfield_{n\sigma}$ un cuerpo cualquiera de funciones racionales (por tanto, no necesariamente el menor cuerpo que contiene un sistema dado), y sea $\Jdom$ la totalidad de los polinomios contenidos en $\Kfield_{n\sigma}$. $\Jdom$ satisface las condiciones 1, 2, 3 de un dominio íntegro y la condición 4a), y por tanto es un dominio relativamente entero; es decir, \emph{la totalidad de los polinomios contenidos en un cuerpo $\Kfield_{n\sigma}$ forma un dominio relativamente entero $\Gdom_{n\rho}$ ($\rho\le\sigma$).}\footnote{También puede ocurrir que $\rho=0$; es decir, que la totalidad de los polinomios de $\Kfield_{n\sigma}$ consista en elementos del dominio de coeficientes $\Omega$.}

Mostramos además la identidad de los dominios relativamente enteros con los ``dominios íntegros de funciones relativamente enteras'' de Hilbert.\footnote{D. Hilbert: Mathematische Probleme. Conferencia de París, 1900. Problema 14 (\emph{Göttinger Nachrichten} 1900).}

Sea
\begin{equation}
  G_1(x),\ldots,G_k(x)
\tag{1}
\end{equation}
una base racional de $\Gdom_{n\rho}$. Por las condiciones 1, 2, 3, 4a, toda combinación racional de los polinomios (1) que sea entera en las indeterminadas, es decir, que sea un polinomio, pertenece a $\Gdom_{n\rho}$. Recíprocamente, todo polinomio de $\Gdom_{n\rho}$ puede expresarse racionalmente, por el concepto de base racional, mediante los polinomios (1); de ahí que $\Gdom_{n\rho}$ sea \emph{idéntico con la totalidad de aquellas combinaciones racionales de los polinomios (1) que son enteras en las indeterminadas, es decir, que son polinomios.} Así tenemos la definición de Hilbert a partir de la base racional especial, mientras que nuestra definición usa sólo propiedades abstractas del dominio.

Para dominios relativamente enteros, la definición de divisor común dada en \S\ 8 para dominios íntegros generales se simplifica:

Sean $F(x),G(x)$ dos polinomios de $\Gdom_{n\rho}$, y sea $L(x)$ un divisor común de estos polinomios. Entonces $L(x)$ se llama un \emph{divisor común en $\Gdom_{n\rho}$} si y sólo si $L(x)$ pertenece a $\Gdom_{n\rho}$.

En efecto, la pertenencia de los cocientes $F(x):L(x)$ y $G(x):L(x)$ se sigue entonces de la condición 4a).

Otra propiedad esencial de los dominios relativamente enteros es que el concepto de ser ``algebraicamente entero con respecto a $\Gdom_{n\rho}$'' puede definirse por medio de cualquier ecuación exactamente como se conoce para los cuerpos de números algebraicos o para el cuerpo $\Omega(x_1,\ldots,x_n)$.

\emph{Definición V: Una cantidad $\xi$ se llama algebraicamente entera con respecto a $\Gdom_{n\rho}$ ---o también relativamente algebraicamente entera--- si satisface alguna ecuación cuyo coeficiente principal es la unidad y cuyos restantes polinomios pertenecen a $\Gdom_{n\rho}$.}

Supóngase, a saber, que la cantidad relativamente algebraicamente entera $\xi$ satisface, por ejemplo, la ecuación
\[
  L(z)=z^\lambda+G_1(x)z^{\lambda-1}+\cdots+G_\lambda(x)=0.
\]
$L(z)$ es divisible por la función entera irreducible con respecto al menor cuerpo contenedor $\Kfield_{n\rho}$,
\[
  M(z)=z^\mu+H_1(x)z^{\mu-1}+\cdots+H_\mu(x).
\]
Pero de esta divisibilidad se sigue, por la extensión conocida del teorema de Gauss,\footnote{Compárese, por ejemplo, J. König: \emph{Einleitung in die allgemeine Theorie der algebraischen Größen} (Leipzig, B. G. Teubner, 1903), cap. IX, \S\ 2, o Weber (edición menor), \S\ 20.} que las funciones racionales $H_1(x),\ldots,H_\mu(x)$ son polinomios de $\Kfield_{n\rho}$, y por tanto pertenecen a $\Gdom_{n\rho}$. Así, la definición de la cantidad relativamente algebraicamente entera $\xi$ se obtiene de la ecuación irreducible. En particular, todo polinomio $F(x)$ de $\Gdom_{n\rho}$ es también algebraicamente entero con respecto a $\Gdom_{n\rho}$, pues satisface la ecuación $z-F(x)=0$.

Debe mencionarse además que, para un dominio íntegro $\Jdom_{n\rho}$ de polinomios, análogamente al menor cuerpo contenedor, también se define el menor dominio relativamente entero que lo contiene, $\Gdom_{n\rho}$, mediante la exigencia siguiente:

$\Gdom_{n\rho}$ contiene, junto con $\Jdom_{n\rho}$, todos los polinomios $H(x)$, y sólo éstos, que pueden representarse como cocientes de dos polinomios de $\Jdom_{n\rho}$.

De esta definición se sigue además: de toda forma de involución y base de involución de $\Jdom_{n\rho}$ surge una correspondiente de $\Gdom_{n\rho}$ escindiendo, a lo sumo, un divisor común en $\Gdom_{n\rho}$ de todos los polinomios de la base.

En efecto, $\Jdom_{n\rho}$ y $\Gdom_{n\rho}$ poseen el mismo menor cuerpo contenedor $\Kfield_{n\rho}$; y toda forma de involución de $\Jdom_{n\rho}$ surge de una forma de involución de $\Kfield_{n\rho}$ por multiplicación por un polinomio de $\Jdom_{n\rho}$, y tiene por tanto coeficientes en $\Gdom_{n\rho}$, los cuales sin embargo pueden tener un divisor común en $\Gdom_{n\rho}$.\footnote{Por ejemplo, para el dominio íntegro $\Jdom_{22}$ formado por todas las combinaciones racionales enteras de $x_1^2$ y $x_1x_2$, se obtiene como forma de involución:
\[
\Psi(z,u)=x_1^2z^2-x_1^4u_1^2-2x_1^3x_2u_1u_2-x_1^2x_2^2u_2^2.
\]
Puesto que $x_1^2$ pertenece a $\Jdom_{22}$, y por tanto con mayor razón al menor dominio relativamente entero que lo contiene, $\Gdom_{22}$, y es en consecuencia un divisor común en $\Gdom_{22}$, resulta como forma de involución de $\Gdom_{22}$:
\[
\Psi(z,u)=z^2-x_1^2u_1^2-2x_1x_2u_1u_2-x_2^2u_2^2.
\]}

Finalmente debe señalarse que el \emph{dominio de aplicación} de un dominio relativamente entero $\Gdom_{n\rho}$, que por \S\ 7 vuelve a ser un dominio íntegro $\Jdom^*_{\rho\rho}$ de polinomios, no necesita ser él mismo un dominio relativamente entero. Pues si $F(x)$ y $G(x)$ son dos polinomios de $\Gdom_{n\rho}$ cuyo cociente no es un polinomio y por tanto no pertenece a $\Gdom_{n\rho}$, entonces tampoco pertenece a $\Jdom^*_{\rho\rho}$ el cociente $F^*(x):G^*(x)$. Pero este cociente, obtenido por especialización de algunas indeterminadas, puede ser perfectamente un polinomio, y entonces para $\Jdom^*_{\rho\rho}$ no se cumple la condición 4a).\footnote{Sea $\Gdom_{22}$, por ejemplo, la totalidad de los polinomios que son combinaciones racionales de $F=x_1^2x_2$ y $G=x_1^2(1+x_3)$. La especialización admisible $x_3=0$ da $F^*:G^*=x_2$, mientras que $F:G$ no es un polinomio. Por tanto, $\Jdom^*_{22}$ no es un dominio relativamente entero.}

\section*{\S\ 10. Dominios relativamente enteros de primera especie y su base íntegra.}

El concepto introducido en \S\ 9 (Definición V), ``relativamente algebraicamente entero'', permite dar una clase de dominios relativamente enteros que son dominios íntegros finitos, y responder así positivamente, al menos para esta clase, una cuestión planteada por D. Hilbert (\emph{Problemas matemáticos}, problema 14). Damos la siguiente

\emph{Definición VI: Un número finito de polinomios de $\Gdom_{n\rho}$ se llama una base íntegra si todo polinomio de $\Gdom_{n\rho}$ puede representarse como combinación racional entera de ese número finito, con coeficientes de $\Omega$. Un dominio íntegro de polinomios que posee una base íntegra se llama dominio íntegro finito.}

La clase de dominios relativamente enteros que debe considerarse, para la cual la base de involución resultará ser una base íntegra, la designamos como dominios relativamente enteros de primera especie según la siguiente

\emph{Definición VII: Un dominio relativamente entero $\Gdom_{n\rho}$ se llama de primera especie si, como dominio de aplicación, posee un dominio relativamente entero $\Gdom^*_{\rho\rho}$ tal que todo polinomio arbitrario en $x_1,\ldots,x_\rho$ (con coeficientes de $\Omega$) es algebraicamente entero sobre $\Gdom^*_{\rho\rho}$.}

Según esta definición, las indeterminadas $x_1,\ldots,x_\rho$, y por consiguiente también la forma lineal $u(x)=u_1x_1+\cdots+u_\rho x_\rho$, son algebraicamente enteras sobre $\Gdom^*_{\rho\rho}$. La función entera irreducible con cero $z=u(x)$ ---la forma de involución del menor cuerpo contenedor $\Kfield^*_{\rho\rho}$--- tiene por tanto la unidad como coeficiente principal y, como restantes coeficientes, polinomios de $\Gdom^*_{\rho\rho}$; por consiguiente es al mismo tiempo una forma de involución de $\Gdom^*_{\rho\rho}$, y en $\Gdom^*_{\rho\rho}$ no puede existir ninguna otra forma de involución. Por el principio de transferencia, también $\Gdom_{n\rho}$ posee entonces una forma de involución cuyo coeficiente principal es la unidad; y, puesto que por el Teorema VI en la representación racional de todos los polinomios de $\Gdom_{n\rho}$ por la base de involución correspondiente sólo entra en el denominador ese coeficiente principal, la base de involución se convierte en una base íntegra. Así:

\emph{Teorema VII. Todo dominio relativamente entero de primera especie es un dominio íntegro finito; la base íntegra está dada por la base de involución obtenida del dominio de aplicación que lo caracteriza. En particular, la base íntegra de los dominios relativamente enteros de primera especie $\Gdom_{\rho\rho}$ está dada por la base de involución unívocamente determinada de $\Gdom_{\rho\rho}$.}

Añadimos un criterio para dominios relativamente enteros de primera especie. Sea $\Gdom_{n\rho}$ de primera especie, y sea una base íntegra de $\Gdom_{n\rho}$ dada por
\begin{equation}
  F_1(x),F_2(x),\ldots,F_k(x).
\tag{1}
\end{equation}
Entonces, por un lema auxiliar de Hilbert,\footnote{D. Hilbert: Über die vollen Invariantensysteme, \emph{Math. Ann.} 42 (1893), \S\ 1. El lema auxiliar allí enunciado sólo para polinomios homogéneos vale igualmente para los no homogéneos.} existen $\rho$ polinomios algebraicamente independientes de $\Gdom_{n\rho}$,
\begin{equation}
  G_1(x),G_2(x),\ldots,G_\rho(x),
\tag{2}
\end{equation}
tales que los polinomios (1), y por tanto todo polinomio de $\Gdom_{n\rho}$, son algebraicamente enteros sobre los polinomios (2). Lo mismo vale para todo polinomio del dominio de aplicación $\Gdom^*_{\rho\rho}$ con respecto a los polinomios correspondientes
\begin{equation}
  G_1^*(x),G_2^*(x),\ldots,G_\rho^*(x).
\tag{3}
\end{equation}

Por la Definición VII, todo polinomio arbitrario en $x_1,\ldots,x_\rho$ es entonces también algebraicamente entero sobre los polinomios (3).

Recíprocamente, supóngase que un dominio de aplicación $\Jdom^*_{\rho\rho}$ de un dominio relativamente entero arbitrario $\Gdom_{n\rho}$ contiene $\rho$ polinomios (3) sobre los cuales todo polinomio en $x_1,\ldots,x_\rho$ es algebraicamente entero. Mostramos entonces que $\Jdom^*_{\rho\rho}$ se convierte por ello en un dominio relativamente entero $\Gdom^*_{\rho\rho}$, y además que $\Gdom^*_{\rho\rho}$, y por consiguiente $\Gdom_{n\rho}$, son de primera especie.

En efecto, sea $F^*(x)$ un polinomio del menor cuerpo $\Kfield^*_{\rho\rho}$ que contiene $\Jdom^*_{\rho\rho}$. Por hipótesis, es algebraicamente entero sobre los polinomios (3). Como la función correspondiente de $\Kfield_{n\rho}$ es entonces algebraicamente entera sobre $\rho$ polinomios de $\Gdom_{n\rho}$, es un polinomio $F(x)$ y por tanto pertenece a $\Gdom_{n\rho}$; por consiguiente $F^*(x)$ pertenece a $\Jdom^*_{\rho\rho}$. El dominio de aplicación $\Jdom^*_{\rho\rho}$ contiene por tanto todos los polinomios de $\Kfield^*_{\rho\rho}$ y es un dominio relativamente entero $\Gdom^*_{\rho\rho}$. Para $\Gdom^*_{\rho\rho}$, y por tanto también para $\Gdom_{n\rho}$, se sigue inmediatamente de las Definiciones V y VII y de la hipótesis que son de primera especie; es decir:

\emph{Los dominios relativamente enteros de primera especie $\Gdom_{n\rho}$ se caracterizan por el hecho de que en un dominio de aplicación $\Jdom^*_{\rho\rho}$ existen $\rho$ polinomios sobre los cuales todo polinomio arbitrario en $x_1,\ldots,x_\rho$ es algebraicamente entero. De esta hipótesis resulta que el propio dominio de aplicación es un dominio relativamente entero.}\footnote{Por ejemplo, en el dominio $\Gdom_{22}$ mencionado en la observación final de \S\ 9, la especialización admisible $x_1=1$ da los dos polinomios $F^*=x_2$, $G^*=1+x_3$, sobre los cuales todo polinomio arbitrario en $x_2,x_3$ es algebraicamente entero. Por tanto, $\Gdom_{22}$ es de primera especie, de hecho con $F(x)$ y $G(x)$ como base íntegra.}

\section*{\S\ 11. Lema auxiliar sobre dependencia algebraicamente entera.}

Damos ahora un criterio para que todo polinomio arbitrario en $x_1,\ldots,x_\rho$ sea algebraicamente entero sobre $\rho$ polinomios en $x_1,\ldots,x_\rho$,
\begin{equation}
  G_1^*(x),\ldots,G_\rho^*(x).
\tag{1}
\end{equation}

Si los polinomios (1) son, en particular, formas homogéneas, entonces, para todo sistema especial de valores que haga anularse (1), se anulan simultáneamente todas las formas algebraicamente enteramente dependientes; por tanto, en particular, también $x_1,\ldots,x_\rho$.

Las formas (1), por consiguiente, no pueden anularse para ningún sistema de valores distinto de
\[
  x_1=0,
  \ldots,
  x_\rho=0;
\]
su resultante debe ser distinta de cero.

Recíprocamente, esta condición es también suficiente y admite una extensión a polinomios no homogéneos según el siguiente

\emph{Lema auxiliar: Una condición suficiente para que todo polinomio arbitrario en $x_1,\ldots,x_\rho$ sea algebraicamente entero sobre $\rho$ polinomios en $x_1,\ldots,x_\rho$}
\[
  G_1^*(x),\ldots,G_\rho^*(x)
\]
\emph{es que no se anule la resultante, tomada con respecto a $x_1,\ldots,x_\rho$, de los términos de dimensión más alta}
\[
  H_1^*(x),\ldots,H_\rho^*(x)
\]
\emph{de los polinomios (1):}
\begin{equation}
  R_0=\left[
  \begin{array}{ccc}
    H_1^*&\cdots&H_\rho^*\\
    x_1&\cdots&x_\rho
  \end{array}
  \right]\ne0.
\tag{2}
\end{equation}
\emph{Para formas homogéneas (1), la condición (2) es necesaria.}\footnote{Para la teoría de resultantes, compárese F. Mertens, Zur Theorie der Elimination (I. y II. parte), \emph{Sitzungsber. d. k. Ak. d. Wiss. Wien}, Math.-naturw. Kl. Abt. IIa, Bd. 108 (1899). Reproducido parcialmente en J. König, \emph{Theorie d. algebr. Größen}, cap. VI.}

Para la demostración formamos los polinomios
\begin{equation}
\begin{gathered}
  \bar G_1^*(x)=G_1^*(x)-\lambda_1,
  \ldots,
  \bar G_\rho^*(x)=G_\rho^*(x)-\lambda_\rho,\\
  \bar\Phi^*(x)=\Phi^*(x)-\mu,
\end{gathered}
\tag{3}
\end{equation}
donde $\lambda_1,\ldots,\lambda_\rho,\mu$ designan indeterminadas y $\Phi^*(x)$ un polinomio arbitrario en $x_1,\ldots,x_\rho$. Sea la resultante de los polinomios (3) con respecto a $x_1,\ldots,x_\rho,1$\footnote{Es decir, la resultante homogénea con respecto a $x_1,\ldots,x_\rho,x_0$, si los polinomios (3) se homogeneizan mediante una indeterminada adicional $x_0$ de modo que se conserve el grado en las $x$.} dada por
\begin{equation}
  R(\lambda,\mu)=\left[
  \begin{array}{cccc}
   \bar G_1^*&\cdots&\bar G_\rho^*&\bar\Phi^*\\
   x_1&\cdots&x_\rho&1
  \end{array}
  \right];
\tag{4}
\end{equation}
es una función racional entera de $\lambda_1,\ldots,\lambda_\rho,\mu$. Según F. Mertens, en el desarrollo de (4) en potencias de $\mu$ puede indicarse explícitamente el coeficiente principal;\footnote{A. a. O. \S\ 3 (demostración del Teorema I), y explícitamente \S\ 6 (determinación de la totalidad de todos los términos de la resultante $R$ que contienen como factor un producto de potencias dado $a_{\lambda\lambda}^{p_\lambda}a_{\mu\mu}^{p_\mu}\cdots a_{\sigma\sigma}^{p_\sigma}$). En nuestro caso se toma $p_\lambda=0$, $p_\mu=0,\ldots,p_\rho=\tau$, $a_{\sigma\sigma}=(-\mu)$. Reproducido en König, cap. VI, \S\ 9.} se obtiene
\begin{equation}
  (-1)^\tau R(\lambda,\mu)
  =R_0^\sigma\mu^\tau+A_1(\lambda)\mu^{\tau-1}+\cdots+A_\tau(\lambda),
\tag{5}
\end{equation}
donde $A_1(\lambda),\ldots,A_\tau(\lambda)$ designan funciones enteras con valores enteros de $\lambda_1,\ldots,\lambda_\rho$ y de los coeficientes de los polinomios (3). Este desarrollo muestra primero que bajo la hipótesis (2), $R_0\ne0$, tampoco $R(\lambda,\mu)$ puede anularse idénticamente. La identidad en $x_1,\ldots,x_\rho,\lambda_1,\ldots,\lambda_\rho,\mu$,
\[
  R(\lambda,
\mu)\equiv 0\pmod{\bar G_1^*(x),\ldots,
\bar G_\rho^*(x),\bar\Phi^*(x)},
\]
pasa, por la sustitución
\[
  \lambda_i=G_i^*(x),\qquad \mu=\Phi^*(x),
\]
a
\begin{equation}
\begin{aligned}
  R(G_i^*(x),\Phi^*(x))
  &=R_0^\sigma\Phi^{*\tau}(x)
    +A_1(G_i^*(x))\Phi^{*\tau-1}(x)+\cdots \\
  &\qquad +A_\tau(G_i(x))=0,
\end{aligned}
\tag{6}
\end{equation}
y esto demuestra el lema auxiliar, pues $R_0$ es, por hipótesis, un número de $\Omega$.

\section*{\S\ 12. La base íntegra de sistemas regulares $\Ssys_{n\rho}$ formados por polinomios.}

El lema auxiliar de \S\ 11, junto con la conclusión de \S\ 10, da inmediatamente el hecho siguiente:

\emph{Si en un dominio de aplicación $\Jdom^*_{\rho\rho}$ de $\Gdom_{n\rho}$ existen $\rho$ polinomios tales que la resultante homogénea de sus términos de dimensión más alta no se anula idénticamente ---$R_0\ne0$---, entonces $\Gdom_{n\rho}$ es de primera especie y, por consiguiente, es un dominio íntegro finito, con la base de involución como base íntegra.}

La condición $R_0\ne0$ permite ahora, mediante el lema auxiliar, una segunda demostración de existencia de la base íntegra, debida esencialmente a D. Hilbert;\footnote{Über die vollen Invariantensysteme, \S\ 2. En el caso de Hilbert, bajo la hipótesis de la existencia obtenida de otro modo de la base íntegra del cuerpo de invariantes, se trata sólo de una interpretación de esta base mediante el sistema fundamental de Kronecker de las cantidades algebraicamente enteras de un cuerpo.} ciertamente esta demostración no hace reconocible como tal la base de involución, pero vale uniformemente para todos los sistemas $\Ssys_{n\rho}$ con la propiedad indicada.

En efecto, sean dados los $\rho$ polinomios de $\Jdom^*_{\rho\rho}$ para los cuales $R_0\ne0$, y que en consecuencia son algebraicamente independientes,
\begin{equation}
  G_1^*(x),\ldots,G_\rho^*(x).
\tag{1}
\end{equation}
Por el lema auxiliar, todo polinomio arbitrario en $x_1,\ldots,x_\rho$, y por tanto en particular todo polinomio del menor cuerpo $\Kfield^*_{\rho\rho}$ que contiene $\Jdom^*_{\rho\rho}$, es algebraicamente entero sobre los polinomios (1). Si los polinomios correspondientes en $\Ssys_{n\rho}$ son por tanto
\begin{equation}
  G_1(x),\ldots,G_\rho(x),
\tag{2}
\end{equation}
entonces, por el principio de transferencia, todo polinomio del menor cuerpo $\Kfield_{n\rho}$ que contiene $\Ssys_{n\rho}$, y por tanto en particular todo polinomio de $\Ssys_{n\rho}$, es algebraicamente entero sobre los polinomios (2). Recíprocamente, toda función de $\Kfield_{n\rho}$ que sea algebraicamente entera sobre los polinomios (2) es también un polinomio. Así, si se considera $\Kfield_{n\rho}$ (según \S\ 4) como un cuerpo algebraico sobre $\Omega(G_1,\ldots,G_\rho)$, la totalidad de los polinomios de $\Kfield_{n\rho}$ ---el dominio relativamente entero $\Gdom_{n\rho}$--- es idéntica con las cantidades algebraicamente enteras de este cuerpo. Sea $G(x)$ un polinomio de $\Kfield_{n\rho}$ que completa (2) a una base racional, y sea $D(x)$ el discriminante de la ecuación irreducible de grado $k$ que $G(x)$ satisface con respecto a $\Omega(G_1,\ldots,G_\rho)$. Entonces todo polinomio de $\Kfield_{n\rho}$, y en particular todo polinomio $F(x)$ de $\Ssys_{n\rho}$, como cantidad algebraicamente entera, admite la representación
\begin{equation}
  F(x)=\frac{\Gamma_1(G_i)G(x)^{k-1}+\Gamma_2(G_i)G(x)^{k-2}+\cdots+\Gamma_k(G_i)}{D(x)}.
\tag{3}
\end{equation}

Hagamos recorrer a $F(x)$ todos los polinomios de $\Ssys_{n\rho}$. Entonces la aplicación del Teorema I de Hilbert sobre la base de módulo (\emph{Math. Ann.} 36) al sistema formado por las funciones correspondientes
\[
  v_1\Gamma_1(G_i)+v_2\Gamma_2(G_i)+\cdots+v_k\Gamma_k(G_i)
\]
muestra la existencia de un número finito de polinomios de $\Ssys_{n\rho}$,
\begin{equation}
  F_1(x),F_2(x),\ldots,F_\lambda(x),
\tag{4}
\end{equation}
tal que todo polinomio de $\Ssys_{n\rho}$ admite una representación
\begin{equation}
  F(x)=A_1(G_i)F_1(x)+A_2(G_i)F_2(x)+\cdots+A_\lambda(G_i)F_\lambda(x),
\tag{5}
\end{equation}
donde $A_1,\ldots,A_\lambda$ designan polinomios en los $G_i(x)$. Los polinomios (2) y (4) forman por tanto una base íntegra de $\Ssys_{n\rho}$; es decir:

\emph{Teorema VIII: Si en un sistema de aplicación $\Jdom^*_{\rho\rho}$ del sistema de polinomios $\Ssys_{n\rho}$ existen $\rho$ polinomios tales que la resultante de los términos de dimensión más alta no se anula idénticamente ---$R_0\ne0$---, entonces $\Ssys_{n\rho}$ posee una base íntegra.}

Este teorema admite una ligera generalización. Basta que los $\rho$ polinomios (1) para los cuales $R_0\ne0$ pertenezcan al dominio íntegro $\Jdom^*_{\rho\rho}$ que es el menor dominio contenedor de $\Ssys^*_{\rho\rho}$. En efecto, los argumentos que conducen al Teorema VIII muestran que la representación (5) sigue siendo válida. Pero, por la definición, todo polinomio $L(x)$ del menor dominio íntegro contenedor $\Jdom_{n\rho}$ puede representarse de manera racional entera mediante un número finito ---dependiente de $L(x)$--- de polinomios de $\Ssys_{n\rho}$; por consiguiente existen $\sigma$ polinomios de $\Ssys_{n\rho}$,
\begin{equation}
  K_1(x),K_2(x),\ldots,K_\sigma(x),
\tag{6}
\end{equation}
tales que los polinomios $G_i(x)$ se convierten en combinaciones racionales enteras de (6). Así, (6) y (4) dan una base íntegra. Introduciendo la terminología siguiente:

\emph{Definición VIII: Un sistema $\Ssys_{n\rho}$ de polinomios se llama regular si en un dominio de aplicación $\Jdom^*_{\rho\rho}$ del menor dominio íntegro que lo contiene existen $\rho$ polinomios tales que la resultante de sus términos de dimensión más alta no se anula idénticamente ---$R_0\ne0$---}

tenemos por tanto:

\emph{Teorema IX: Todo sistema regular $\Ssys_{n\rho}$ de polinomios posee una base íntegra.}\footnote{Que hay sistemas no regulares sin base íntegra fue mostrado por Hilbert con un ejemplo (\emph{Gött. Nachr.} 1891, p. 233). El siguiente es aún más sencillo:
\[
  \Ssys_{22}=x_1,
  x_1x_2,
  x_1x_2^2,
  \ldots,
  x_1x_2^\nu,
  \ldots\quad\text{in inf.}
\]
La existencia de la base íntegra aquí equivaldría a la existencia de un entero positivo $\nu$ tal que se verifiquen las identidades
\[
  x_1x_2^{\nu+\sigma}=\sum C\cdot x_1^{i_0}(x_1x_2)^{i_1}(x_1x_2^2)^{i_2}\cdots(x_1x_2^\nu)^{i_\nu}
  \quad(\sigma=1,2,\ldots).
\]
La comparación de las dimensiones en $x_1,x_2$, ya para $\sigma=1$, conduce a las dos ecuaciones condicionales para los exponentes
\[
  1=i_0+i_1+\cdots+i_\nu,
  \qquad
  \nu+1=i_1+2i_2+\cdots+\nu i_\nu,
\]
que claramente no tienen solución en enteros no negativos. Por tanto $\Ssys_{22}$ no posee base íntegra.}

Damos además, análogamente a \S\ 5 para la involución, una interpretación geométrica de los sistemas regulares. Para ello consideramos, como allí en (7), el sistema de ecuaciones
\begin{equation}
  L^*(x)-L^*(\xi)=0,
\tag{7}
\end{equation}
donde $L^*(x)$ recorre todos los polinomios de $\Jdom^*_{\rho\rho}$. Al homogenizar mediante $x_0,\xi_0$, (7) pasa al sistema de ecuaciones entre formas homogéneas
\begin{equation}
  \xi_0^{m}M^*(x)-x_0^{m}M^*(\xi)=0,
\tag{8}
\end{equation}
donde, si $H^*(x)$ designa los términos de dimensión más alta de $L^*(x)$, vale la relación
\begin{equation}
  H^*(x)=M^*(x)\pmod{x_0}.
\tag{9}
\end{equation}

Como \emph{puntos fundamentales} de $\Jdom^*_{\rho\rho}$ (compárese la nota a \S\ 5), designamos las soluciones de (8) que son independientes de $x$, es decir, distintas de las filas conjugadas a $x$. Los puntos fundamentales son, por tanto, los sistemas de valores distintos de $\xi_1=0,\ldots,\xi_\rho=0$ para los cuales se cumple simultáneamente
\[
  \xi_0^m=0,
  \qquad M^*(\xi)=0,
\]
o, por (9),
\[
  \xi_0=0,
  \qquad H^*(\xi)=0.
\]

Si existen ahora $\rho$ formas $H^*(\xi)$ para las cuales $R_0\ne0$, y por tanto sin sistema común de soluciones, entonces $\Jdom^*_{\rho\rho}$ no puede poseer ningún punto fundamental. Si, en particular, $\Ssys^*_{\rho\rho}$ consiste en formas homogéneas, entonces para $R_0\ne0$ tampoco $\Ssys^*_{\rho\rho}$ puede poseer ningún punto fundamental, puesto que tal punto sería al mismo tiempo un punto fundamental de $\Jdom^*_{\rho\rho}$. Pero también vale el recíproco. Para ello consideramos el sistema $\Jdom(H^*)$ formado por la totalidad de las formas $H^*(x)$; éste vuelve a ser un dominio íntegro.

Del Teorema I de Hilbert sobre la base de módulo y del lema auxiliar de Hilbert (compárese la cita en \S\ 10) se sigue la existencia de $\rho$ formas de $\Jdom(H^*)$ cuya anulación tiene como consecuencia la anulación de todas las formas de $\Jdom(H^*)$. Si ahora $\Jdom^*_{\rho\rho}$ no posee ningún punto fundamental, estas $\rho$ formas no pueden anularse simultáneamente; su resultante $R_0\ne0$.

Así:

\emph{Los sistemas regulares $\Ssys_{n\rho}$ de polinomios se caracterizan por el hecho de que un dominio de aplicación $\Jdom^*_{\rho\rho}$ del menor dominio íntegro que los contiene no posee ningún punto fundamental. Para sistemas regulares $\Ssys^*_{\rho\rho}$ de formas homogéneas, el propio sistema no posee ya ningún punto fundamental.}\footnote{El sistema no regular $\Ssys_{22}$ dado en la nota precedente posee el punto fundamental
\[
  \xi_1:\xi_2:\xi_0=0:1:0.
\]}

\section*{\S\ 13. Ejemplos de dominios relativamente enteros de primera especie y de sistemas regulares.}

1. Como primer ejemplo de un dominio relativamente entero de primera especie $\Gdom_{nn}$, consideramos la \emph{totalidad de los polinomios en $x_1,\ldots,x_n$ que admiten un grupo de permutaciones dado $G$}; es decir, la totalidad de los polinomios de un dominio de género de Lagrange. Puesto que $x_1,\ldots,x_n$, en virtud de la identidad
\begin{equation}
  x_k^n-\sigma_1(x)x_k^{n-1}+\sigma_2(x)x_k^{n-2}+\cdots\pm\sigma_n(x)=0
  \qquad(k=1,2,
\ldots,n),
\tag{1}
\end{equation}
son algebraicamente enteros sobre las funciones simétricas elementales pertenecientes a $\Gdom_{nn}$, y por tanto, por la Definición V, algebraicamente enteros sobre $\Gdom_{nn}$, $\Gdom_{nn}$ es de primera especie. Además $\Gdom_{nn}$, como dominio de primera especie, puesto que sus elementos son formas homogéneas y $n=\rho$, forma un sistema regular (por la conclusión de \S\ 10). En efecto, por (1), la resultante $R_0$ de las funciones simétricas elementales no puede anularse idénticamente; por las reglas de cálculo de la teoría de resultantes, se calcula $R_0=\pm1$. La \emph{forma de involución} de $\Gdom_{nn}$ está determinada unívocamente ---pues $\Gdom_{nn}$ es de primera especie y $n=\rho$--- como la forma de involución del correspondiente dominio de género de Lagrange. Como las cantidades conjugadas a $x_1,\ldots,x_n$ con respecto a este dominio surgen por las permutaciones del grupo $G$, esta forma de involución está dada por
\begin{equation}
\begin{aligned}
  \Psi(z,u)&=\prod_i
  (z-u_1x_{i_1}-u_2x_{i_2}-\cdots-u_nx_{i_n})\\
  &=z^i+\sum{}'G_{\alpha\alpha_1\ldots\alpha_n}(x)
  z^\alpha u_1^{\alpha_1}\cdots u_n^{\alpha_n},
\end{aligned}
\tag{2}
\end{equation}
donde el producto debe extenderse sobre todas las permutaciones de $G$. La fórmula (2) representa la ``forma de Galois del grupo $G$'', y por tanto vale el hecho siguiente:

\emph{Todos los polinomios en $x_1,\ldots,x_n$ que admiten un grupo de permutaciones $G$ son combinaciones racionales enteras de los coeficientes $G_{\alpha\alpha_1\ldots\alpha_n}(x)$ de la forma de Galois del grupo.}

2. Como ejemplo de \emph{sistemas regulares}, considérense los sistemas $\Ssys_{n1}$ de polinomios, es decir, sistemas que contienen sólo un polinomio algebraicamente independiente.

En efecto, sea $\Ssys^*_{11}$ un sistema de aplicación de $\Ssys_{n1}$ y sea
\[
  F^*(x)=a_0x^k+a_1x^{k-1}+\cdots+a_k\qquad(a_0\ne0)
\]
un polinomio de $\Ssys^*_{11}$. Entonces
\begin{equation}
  R_0=a_0\ne0.
\tag{3}
\end{equation}
Así $\Ssys_{n1}$ es regular por la Definición VIII, y en consecuencia:

\emph{Todo sistema $\Ssys_{n1}$ de polinomios, en particular todo sistema de infinitos polinomios en una indeterminada, posee una base íntegra.}\footnote{Los sistemas no regulares más sencillos posibles sin base íntegra son por tanto sistemas $\Ssys_{22}$. Que tales sistemas existen efectivamente lo muestran el ejemplo de Hilbert y el ejemplo dado en la nota a \S\ 12.}

3. Especializamos ahora los sistemas $\Ssys_{n1}$ a dominios relativamente enteros $\Gdom_{n1}$, que, como sistemas regulares por \S\ 12, son dominios relativamente enteros de primera especie. Su base íntegra está por tanto dada por una base de involución, que al mismo tiempo es la base de involución del menor cuerpo contenedor $\Kfield_{n1}$. Pero por \S\ 6, cualquier función de la base de involución de $\Kfield_{n1}$ es al mismo tiempo una base mínima ---la llamamos función de Lüroth del cuerpo--- y toda función de la base de involución depende de esta función de Lüroth mediante una transformación lineal fraccionaria. En el caso presente la función de Lüroth, puesto que pertenece a $\Gdom_{n1}$, es un polinomio; todo otro polinomio de la base de involución depende por tanto lineal-fraccionariamente, y por (3) algebraicamente enteramente, luego entera y linealmente, de esta función de Lüroth $L(x)$, que así se convierte en la base íntegra. La forma de involución de $\Gdom^*_{11}$, poniendo también $(u_1=1)$, es
\[
  \Psi^*(z)=L^*(z)-L^*(x),
\]
de donde $L(x)$ aparece de nuevo como base íntegra. Así:

\emph{Los dominios relativamente enteros $\Gdom_{n1}$ son de primera especie; su base íntegra consta de un polinomio, la función de Lüroth del menor cuerpo contenedor.}

4. Como ejemplo de un dominio relativamente entero $\Gdom_{n1}$, mencionemos especialmente los invariantes proyectivos racionales enteros de una forma cuadrática en $s$ variables. De acuerdo con lo precedente, es sabido que pueden representarse como combinaciones racionales enteras del determinante $|a_{ik}|$ de la forma; y este determinante es al mismo tiempo la función de Lüroth del cuerpo de invariantes correspondiente.

\section*{\S\ 14. Sistemas formados por polinomios enteros.}

Los teoremas de base obtenidos deben ahora afinarse con respecto a la integralidad sobre los enteros.

En lo que sigue, el dominio de coeficientes $\Omega$ se supondrá por tanto un cuerpo de números algebraicos (finito), y $[\Omega]$ designará la totalidad de los enteros algebraicos de $\Omega$, que llamaremos brevemente enteros. Por polinomio entero se entiende un polinomio con coeficientes en $[\Omega]$; de ahora en adelante, $F(x),G(x),\ldots$ designarán sólo polinomios enteros.

En analogía con \S\ 7, consideramos los distintos sistemas compuestos de polinomios enteros:

Un sistema arbitrario de polinomios enteros se designará como un sistema entero $\Ssys_{n\rho}$.

La familia lineal entera $\Lsys_{n\rho}$ es un sistema entero que, junto con $F_1(x)$ y $F_2(x)$, contiene siempre $c_1F_1(x)+c_2F_2(x)$, donde $c_1,c_2$ son enteros arbitrarios de $[\Omega]$.

El dominio íntegro entero $\Jdom_{n\rho}$ es una familia lineal entera que, junto con $F(x)$ y $G(x)$, contiene siempre $F(x)\cdot G(x)$.

El dominio relativamente entero entero $\Gdom_{n\rho}$ es un dominio íntegro que, junto con $F(x)$ y $G(x)$ ---donde $G(x)\ne0$---, contiene siempre y sólo contiene el cociente $F(x):G(x)$ cuando este cociente es un polinomio entero.

Correspondientemente a \S\ 7, el menor dominio íntegro entero $\Jdom_{n\rho}$ que contiene un sistema entero $\Ssys_{n\rho}$ consta de todos los polinomios $H(x)$ que son combinaciones racionales enteras, en el sentido entero, de un número finito de polinomios de $\Ssys_{n\rho}$;

el menor dominio relativamente entero entero $\Gdom_{n\rho}$ que lo contiene consta de todos los polinomios enteros $H(x)$ que son combinaciones racionales de un número finito de polinomios de $\Ssys_{n\rho}$.

Pasamos ahora a la transferencia de los teoremas de base. Con respecto a la base racional, debe observarse que los conceptos de cuerpo, menor cuerpo contenedor y base racional, que descansan sólo en la ``racionalidad'', no sufren cambio. Además, el sistema de aplicación $\Ssys^*_{\rho\rho}$ de un sistema entero $\Ssys_{n\rho}$ puede elegirse de nuevo, de maneras arbitrariamente numerosas, como sistema entero. Puesto que también para la familia lineal entera siguen siendo válidos los razonamientos que llevan a la cota del número de generadores, se obtiene:

\emph{Teorema V$'$: Todo sistema entero $\Ssys_{n\rho}$ posee una base racional; la base racional de una familia lineal entera $\Lsys_{n\rho}$ puede elegirse en particular de modo que contenga a lo sumo $\rho+1$ polinomios.}

Investigamos ahora el análogo de la base de involución de los dominios íntegros de polinomios (\S\ 8). Para ello, el teorema sobre las \emph{funciones simétricas de filas de cantidades} debe extenderse con respecto a la \emph{integralidad sobre los enteros} por medio del siguiente

\emph{Lema auxiliar. Las funciones simétricas enteras, en el sentido de los enteros, de $n$ filas de cantidades $(yz\cdots t)$:}
\begin{equation}
\begin{array}{cccc}
  y_1&z_1&\cdots&t_1\\
  y_2&z_2&\cdots&t_2\\
  \vdots&\vdots&&\vdots\\
  y_n&z_n&\cdots&t_n
\end{array}
\tag{1}
\end{equation}
\emph{son funciones enteras, en el sentido de los enteros, de un número finito de entre ellas, a saber, de las funciones simétricas elementales}
\begin{equation}
\begin{gathered}
  \sigma_1(y),\sigma_2(y),\ldots,\sigma_n(y);\quad
  \sigma_1(z),\sigma_2(z),\ldots,\sigma_n(z);\quad\ldots;\\
  \sigma_1(t),\ldots,\sigma_n(t);
\end{gathered}
\tag{2}
\end{equation}
\emph{y de las funciones}
\begin{equation}
  \chi_1(y,z,\ldots,t),\ldots,\chi_k(y,z,\ldots,t),
\tag{3}
\end{equation}
\emph{que son algebraicamente enteras y enteras sobre las funciones (2).}

En efecto, las $n$ cantidades $y_1,\ldots,y_n$ dependen algebraicamente entera e integralmente sobre los enteros de $\sigma_1(y),\ldots,\sigma_n(y)$; la afirmación correspondiente vale para $z_1,\ldots,z_n$ respecto de $\sigma_1(z),\ldots,\sigma_n(z)$, y así sucesivamente. El cuerpo de funciones simétricas de las filas de cantidades (1) forma por tanto un cuerpo algebraico sobre el cuerpo determinado por las funciones (2), de tal manera que las cantidades algebraicamente enteras e integrales sobre los enteros de este cuerpo son idénticas con las funciones simétricas enteras, en el sentido de los enteros, de las filas de cantidades (1). Si se toma entonces (3) como sistema fundamental de Kronecker, el lema auxiliar queda demostrado.

Sea ahora $\Jdom_{n\rho}$ un dominio íntegro entero, $\Jdom^*_{\rho\rho}$ un dominio de aplicación entero, y
\[
  \Psi^*(z,u)=G^*(x)z^i+
  \sum{}'G^*_{\alpha\alpha_1\ldots\alpha_\rho}(x)
  z^\alpha u_1^{\alpha_1}\cdots u_\rho^{\alpha_\rho}
\]
una forma de involución de $\Jdom^*_{\rho\rho}$; aquí $G^*(x)$ puede también tener dimensión cero, es decir, ser un número de $[\Omega]$. Las funciones simétricas elementales correspondientes a las funciones (2) están entonces dadas por ciertas funciones entre los cocientes
\begin{equation}
  G^*_{\alpha\alpha_1\ldots\alpha_\rho}(x):G^*(x).
\tag{4}
\end{equation}
Como por el lema auxiliar las funciones (3) dependen algebraicamente entera e integralmente sobre los enteros de (2), las funciones de $\Kfield^*_{\rho\rho}$ correspondientes a (3) están dadas ---por la extensión del teorema de Gauss a polinomios con coeficientes algebraicamente enteros--- por
\begin{equation}
  K_\nu^*(x):(G^*(x))^\sigma,
\tag{5}
\end{equation}
donde $K_\nu^*(x)$ son polinomios enteros de $\Jdom^*_{\rho\rho}$. Así, $K_\nu^*(x)$ pertenece al dominio si se trata de dominios relativamente enteros; en general son cocientes de polinomios de $\Jdom^*_{\rho\rho}$.

Si ahora $F^*(x)$ es un polinomio entero arbitrario de $\Jdom^*_{\rho\rho}$, entonces se tiene
\[
  F^*(x)=\frac1i\Bigl(F^*(x)+F^*(x^{(1)})+\cdots+F^*(x^{(i-1)})\Bigr),
\]
y por tanto $i\cdot F^*(x)$ se convierte en una función simétrica entera, en el sentido de los enteros, de las filas de cantidades $x,x^{(1)},\ldots,x^{(i-1)}$.

Por el lema auxiliar y el principio de transferencia, por consiguiente:

\emph{Teorema VI$'$: Para dominios íntegros enteros $\Jdom_{n\rho}$, cada forma de involución determina una base racional distinguida de tal manera que todo polinomio $F(x)$ de $\Jdom_{n\rho}$ admite una representación}
\[
  F(x)=\frac{\Gamma(H(x),H_1(x),\ldots,H_s(x))}{i\cdot H(x)^i},
\]
\emph{donde $\Gamma$ designa una función entera en el sentido de los enteros. Para dominios relativamente enteros enteros $\Gdom_{n\rho}$, con dominio de aplicación relativamente entero entero $\Gdom^*_{\rho\rho}$, el denominador $H(x)$ está dado por el coeficiente principal de la forma de involución correspondiente.}

\section*{\S\ 15. Dominios relativamente enteros enteros de primera especie y sistemas regulares.}

Para sistemas enteros, la base íntegra entera ocupa el lugar de la base íntegra; es decir, un número finito de polinomios enteros del sistema tal que todo polinomio entero del sistema puede representarse como combinación entera, en el sentido de los enteros, de ese número finito.

Los teoremas sobre la base íntegra entera corren esencialmente paralelos a los anteriores, y sólo las demostraciones requieren ligeras modificaciones.

Primero, la Definición V (\S\ 9) se transfiere directamente.

\emph{Definición V$'$: Una cantidad $\xi$ se llama algebraicamente entera y entera con respecto al dominio relativamente entero entero $\Gdom_{n\rho}$ si satisface alguna ecuación cuyo coeficiente principal es una unidad de $[\Omega]$ y cuyos restantes polinomios pertenecen a $\Gdom_{n\rho}$.}

La extensión del teorema de Gauss a polinomios con coeficientes algebraicamente enteros muestra, como en \S\ 9, que de esto puede obtenerse la definición por la ecuación irreducible.

De la Definición V$'$ se obtiene además la transferencia de la Definición VII.

\emph{Definición VII$'$: Un dominio relativamente entero entero $\Gdom_{n\rho}$ se llama de primera especie si, como dominio de aplicación, posee un dominio relativamente entero $\Gdom^*_{\rho\rho}$ tal que todo polinomio entero arbitrario en $x_1,\ldots,x_\rho$ es algebraicamente entero y entero sobre $\Gdom^*_{\rho\rho}$.}

Para mostrar la finitud de estos dominios de primera especie, observamos primero que, como en \S\ 10, se sigue de la definición que el coeficiente principal de la forma de involución es una unidad de $[\Omega]$.

Por el Teorema VI$'$ (\S\ 14), todo polinomio $F(x)$ de $\Gdom_{n\rho}$ admite por tanto una representación
\[
  F(x)=\frac{\Gamma(H_1(x),\ldots,H_\sigma(x))}{i}.
\]
Aplicando aquí al sistema formado por los polinomios enteros $\Gamma(H_1,\ldots,H_\sigma)$ el Teorema II de Hilbert sobre la base de módulo entera (\emph{Math. Ann.} 36), enunciado originalmente para coeficientes numéricos racionales enteros, en su extensión a polinomios algebraicamente enteros,\footnote{Si $w_1,\ldots,w_k$ designa un sistema fundamental de los enteros algebraicos de $\Omega$, y se pone
\[
  \Gamma(H_1,
\ldots,H_\sigma)=\Gamma_1(H)w_1+\cdots+\Gamma_k(H)w_k,
\]
entonces $\Gamma_1(H),\ldots,\Gamma_k(H)$ tienen coeficientes numéricos racionales enteros. La aplicación del Teorema II al sistema formado por las formas $v_1\Gamma_1(H)+\cdots+v_k\Gamma_k(H)$ muestra directamente la validez de la extensión.} se sigue la existencia de la base íntegra entera; es decir:

\emph{Teorema VII$'$: Todo dominio relativamente entero entero de primera especie posee una base íntegra entera.}

Transferimos además la definición de sistemas regulares:

\emph{Definición VIII$'$: Un sistema entero $\Ssys_{n\rho}$ se llama entero-regular si en un dominio de aplicación $\Jdom^*_{\rho\rho}$ del menor dominio íntegro entero que lo contiene existen $\rho$ polinomios tales que la resultante de los términos de dimensión más alta es una unidad de $[\Omega]$ ---$R_0=\eps$.}

Para transferir la demostración de existencia de la base íntegra, observamos que el lema auxiliar de \S\ 11 admite la siguiente formulación más precisa, que se sigue de las ecuaciones (5) y (6) de \S\ 11 y del hecho de que $A_1(\lambda),\ldots,A_\tau(\lambda)$ son funciones enteras, en el sentido de los enteros, de $\lambda_1,\ldots,\lambda_\rho$:

\emph{Una condición suficiente para que todo polinomio entero arbitrario en $x_1,\ldots,x_\rho$ dependa algebraicamente entera e integralmente sobre los enteros de $\rho$ polinomios enteros es que la resultante de los términos de dimensión más alta de estos polinomios sea una unidad de $[\Omega]$ ---$R_0=\eps$.}

De este lema auxiliar, exactamente como en \S\ 12, se sigue para todo polinomio $F(x)$ de un sistema entero-regular la representación (3) de \S\ 12, donde ahora $\Gamma_1(G_i),\ldots,\Gamma_k(G_i)$ son polinomios enteros en los $G_i$. Además, como en \S\ 12, la aplicación del Teorema II de Hilbert (\emph{Math. Ann.} 36) en su extensión a polinomios algebraicamente enteros, junto con la definición del menor dominio íntegro entero que lo contiene, da la existencia de la base íntegra entera; es decir:

\emph{Teorema IX$'$: Todo sistema entero-regular posee una base íntegra entera.}

Como ejemplo de dominio relativamente entero entero de primera especie, en analogía con el ejemplo 1 de \S\ 13, nos remitimos a la totalidad $\Gdom_{nn}$ de los polinomios enteros de un dominio de género de Lagrange. Que $\Gdom_{nn}$ es entero de primera especie se sigue de la identidad
\[
  x_k^n-\sigma_1(x)x_k^{n-1}+\cdots\pm\sigma_n(x)=0
  \qquad(k=1,2,
\ldots,n);
\]
puesto que la resultante de las funciones simétricas elementales tiene el valor $\pm1$, $\Gdom_{nn}$ es también un sistema entero regular.

Un ejemplo adicional está dado, por el lema auxiliar de \S\ 14, por la totalidad de las funciones simétricas enteras, en el sentido de los enteros, de $n$ filas de cantidades.

Erlangen, mayo de 1914.

\section*{7. El teorema de finitud para los invariantes de grupos finitos}
\begin{center}
Math. Ann. 77 (1916), pp. 89--92
\end{center}

En lo que sigue se dará una demostración completamente elemental de la finitud de los invariantes de grupos finitos ---una que se apoya sólo en la teoría de las funciones simétricas--- y al mismo tiempo proporciona una especificación efectiva del sistema completo; mientras que la demostración habitual, basada en el teorema de Hilbert sobre bases de módulos (Ann. 36), es sólo una demostración de existencia.\footnote{Cf., por ejemplo, Weber, \emph{Lehrbuch der Algebra} (2.a ed.), vol. 2, \S\ 57.}

Sea el grupo finito $\Hgrp$ compuesto por las $h$ transformaciones lineales (con determinante no nulo) $A_1,\ldots,A_h$, donde $A_k$ denota la transformación lineal
\[
  x_1^{(k)}=\sum_{\nu=1}^{n}a_{1\nu}^{(k)}x_\nu,\ldots,
  x_n^{(k)}=\sum_{\nu=1}^{n}a_{n\nu}^{(k)}x_\nu,
\]
o, abreviadamente, $(x^{(k)})=A_k(x)$. Por tanto, el grupo $\Hgrp$ lleva la fila $(x)$ con elementos $x_1,
\ldots,x_n$ a las filas $(x^{(k)})$ con elementos $x_1^{(k)},
\ldots,x_n^{(k)}$. Puesto que la identidad debe estar contenida entre $A_1,
\ldots,A_h$, la fila $(x)$ también está contenida entre las filas $(x^{(k)})$. Por un invariante racional entero (absoluto) del grupo se entiende una función racional entera de $x_1,
\ldots,x_n$ que permanece idénticamente inalterada bajo la aplicación de $A_1,
\ldots,A_h$; para tal invariante $f(x)$ se tiene, por consiguiente,
\begin{equation}
  f(x)=f(x^{(1)})=\cdots=f(x^{(h)})=
  \frac1h\sum_{k=1}^{h}f(x^{(k)}).
\tag{1}
\end{equation}

\subsection*{1.}
La fórmula (1) dice que $f(x)$ es una función simétrica racional entera de las filas de cantidades $(x^{(1)}),
\ldots,(x^{(h)})$; y, como cada sumando $f(x^{(k)})$ contiene sólo la única fila $(x^{(k)})$, éste es el caso más simple, usualmente llamado el caso de \emph{una forma}. Por el teorema conocido sobre funciones simétricas de filas de cantidades,\footnote{Cf. la nota a 2.} $f(x)$ es, por tanto, representable de manera racional entera mediante las funciones simétricas elementales de esas filas, es decir, mediante los coeficientes $G_{\alpha\alpha_1\ldots\alpha_n}(x)$ de la ``resolvente de Galois''
\[
\begin{aligned}
  \Phi(z,u)
  &=\prod_{k=1}^{h}\bigl(z+u_1x_1^{(k)}+\cdots+u_nx_n^{(k)}\bigr) \\
  &=z^h+
    \sum G_{\alpha\alpha_1\ldots\alpha_n}(x)
    z^{\alpha}u_1^{\alpha_1}\cdots u_n^{\alpha_n},
\end{aligned}
\qquad
\left(\begin{array}{c}
\alpha+\alpha_1+
\cdots+
\alpha_n=h\\
\alpha\ne h
\end{array}\right),
\]
donde los $G_{\alpha\alpha_1\ldots\alpha_n}(x)$ son invariantes de grado $\alpha_1+\cdots+\alpha_n$ en las $x$. Así queda probado:

\emph{Los coeficientes $G_{\alpha\alpha_1\ldots\alpha_n}(x)$ de la resolvente de Galois forman un sistema completo de invariantes del grupo, de tal modo que todo invariante puede representarse racionalmente y enteramente mediante estos invariantes finitamente numerosos.}

\subsection*{2.}
Partiendo de (1), también se puede hacer la siguiente consideración todavía más elemental,\footnote{Esto equivale a una demostración del teorema sobre funciones simétricas de filas de cantidades en el caso de una forma mencionado arriba.} que al mismo tiempo conduce a un segundo sistema completo. Pongamos
\[
  f(x)=a+b x_1^{\mu_1}\cdots x_n^{\mu_n}
       +\cdots+c x_1^{\nu_1}\cdots x_n^{\nu_n},
\]
donde $a,b,\ldots,c$ denotan constantes. Entonces, por (1), se tiene
\[
\begin{aligned}
 h\cdot f(x)
 &=h\cdot a+b\sum_{k=1}^{h}(x_1^{(k)})^{\mu_1}\cdots(x_n^{(k)})^{\mu_n}
     +\cdots
     +c\sum_{k=1}^{h}(x_1^{(k)})^{\nu_1}\cdots(x_n^{(k)})^{\nu_n} \\
 &=h\cdot a+b\cdot J_{\mu_1\ldots\mu_n}+\cdots+c\cdot J_{\nu_1\ldots\nu_n}.
\end{aligned}
\]
Por tanto, todo invariante puede componerse de manera entera y lineal a partir de los invariantes especiales
\[
  J_{\mu_1\ldots\mu_n}
  =\sum_{k=1}^{h}(x_1^{(k)})^{\mu_1}\cdots(x_n^{(k)})^{\mu_n},
\]
y basta probar la afirmación para estos invariantes especiales. Pero, salvo un factor numérico, $J_{\mu_1\ldots\mu_n}$ es el coeficiente de $u_1^{\mu_1}
\cdots u_n^{\mu_n}$ en la expresión
\[
  S_\mu=\sum_{k=1}^{h}
  \bigl(u_1x_1^{(k)}+\cdots+u_nx_n^{(k)}\bigr)^{\mu},
\]
donde $\mu=\mu_1+\cdots+\mu_n$, y esta expresión es la $\mu$-ésima suma de potencias de las $h$ formas lineales
\[
  \xi_1=u_1x_1^{(1)}+\cdots+u_nx_n^{(1)},\ldots,
  \xi_h=u_1x_1^{(h)}+\cdots+u_nx_n^{(h)}.
\]
Se sabe que las infinitamente muchas sumas de potencias $S_\mu$ son combinaciones racionales enteras de
\[
  S_1,S_2,\ldots,S_h,
\]
cuyos coeficientes vienen dados por los invariantes
\[
  J_{\mu_1\ldots\mu_n}\qquad(\mu_1+\cdots+\mu_n\le h).
\]
Así se obtiene un segundo sistema completo:

\emph{Un sistema completo de invariantes del grupo está dado por todos los invariantes $J_{\mu_1\ldots\mu_n}$ para los cuales $\mu_1+\cdots+\mu_n\le h$, donde $h$ denota el orden del grupo.}

La conexión con 1. se establece observando que las sumas de potencias $S_1,
\ldots,S_h$ pueden sustituirse por las funciones simétricas elementales de $\xi_1,
\ldots,
\xi_h$, lo que conduce al sistema completo de los $G_{\alpha\alpha_1\ldots\alpha_n}(x)$ dado allí. Ambos resultados muestran que todos los invariantes pueden expresarse racionalmente y enteramente mediante aquellos cuyo grado en las $x$ no excede el orden $h$ del grupo.

\subsection*{3.}
De lo que se acaba de probar pueden extraerse consecuencias para la representación racional. Todo invariante absoluto racional puede representarse ---como se ve inmediatamente multiplicando numerador y denominador por los conjugados del denominador respecto del grupo--- como cociente de dos invariantes absolutos racionales enteros, no necesariamente primos entre sí. De aquí se sigue que todo invariante absoluto racional puede expresarse racionalmente mediante los coeficientes $G_{\alpha\alpha_1\ldots\alpha_n}(x)$ de la resolvente de Galois $\Phi(z,u)$. Este teorema también puede probarse fácilmente de diversas maneras sin pasar por el teorema de finitud; ya está formulado en Weber II, \S\ 58.\footnote{Sin embargo, la demostración dada allí no es correcta. Sólo se muestra que la función $\Psi(t)$ en la fórmula (7) tiene invariantes como coeficientes, pero no que esos invariantes sean expresables racionalmente mediante los coeficientes de $\Phi(t)$. Esta laguna se evita si, al representar los $x_i^{(k)}$ mediante $\xi_k$, se usa un procedimiento de diferenciación conocido en vez de la fórmula de interpolación de Lagrange. Ésta es la relación obtenida al diferenciar la identidad $\Phi(-\xi_k,u)=0$ con respecto a todos los $u$:
\[
 \left[\frac{\partial\Phi}{\partial u_i}
      -x_i^{(k)}\frac{\partial\Phi}{\partial z}\right]_{z=-\xi_k}=0.
\]
En consecuencia, para cada función racional $\omega(x)$ debe sustituirse a las fórmulas (7) y (8) de Weber la fórmula siguiente:
\[
 \omega(x_1^{(k)},\ldots,x_n^{(k)})=
 \left.
 \omega\!\left(
   \frac{\partial\Phi/\partial u_1}{\partial\Phi/\partial z},\ldots,
   \frac{\partial\Phi/\partial u_n}{\partial\Phi/\partial z}
 \right)\right|_{z=-\xi_k},
\]
la cual contiene sólo coeficientes de $\Phi(z,u)$ y cuya suma sobre todos los $k$, para invariantes $\omega(x)$, da la representación deseada.}

\subsection*{4.}
Finalmente, debe mencionarse que los resultados obtenidos aquí están contenidos implícitamente en mi trabajo ``Cuerpos y sistemas de funciones racionales'';\footnote{Math. Ann. 76, p. 161 (1915).} más precisamente, el sistema completo dado en 1. está contenido en los Teoremas VI y VII, y una demostración de la representación racional independiente de este teorema de finitud está contenida en el Teorema III.\footnote{Para invariantes relativos, los Teoremas VIII y IX contienen una demostración de finitud que igualmente descansa sobre un fundamento distinto del habitual, pero que también es sólo una demostración de existencia; la razón es que los invariantes relativos no forman un cuerpo.} En este caso especial, sin embargo, el curso de la demostración y el resultado se simplifican considerablemente en comparación con la teoría general. Fui llevada a aplicar las investigaciones generales a los invariantes de grupos finitos por conversaciones con E. Fischer, de quien proceden, en sustancia, la demostración dada en 2 ---en el caso general el sistema completo que allí aparece no es definible racionalmente--- y también la nota a 3.

Erlangen, mayo de 1915.

\clearpage
\section*{8. Sobre la representación racional entera de los invariantes de un sistema de arbitrariamente muchas formas fundamentales}
\begin{center}
Math. Ann. 77 (1916), pp. 93--102
\end{center}

Al final de su contribución al volumen de homenaje a Schwarz, ``Sobre los invariantes de un sistema de arbitrariamente muchas formas fundamentales'', D. Hilbert formula la conjetura de que estos invariantes pueden representarse racionalmente y enteramente mediante los finitamente muchos invariantes del sistema $(J,PJ)$, donde $J$ denota el sistema completo de invariantes de las formas fundamentales linealmente independientes, y los invariantes $PJ$ surgen de $J$ mediante procesos polares.

Muestro a continuación, en I, que esta conjetura es de hecho correcta, y ello como una consecuencia simple del teorema de reducción según el cual toda forma en arbitrariamente muchas filas de $n$ variables cada una puede representarse como una suma de polares de formas que contienen sólo $n$ filas fijas y que ellas mismas se derivan de la forma original mediante procesos polares. Este teorema de reducción es un caso especial, formulado por primera vez por F. Mertens, de la generalización, dada por A. Capelli, F. Mertens y J. Deruyts, del desarrollo en serie de Clebsch--Gordan a formas en arbitrariamente muchas filas de $n$ variables cada una; esta generalización se demuestra por transformación formal de los procesos de diferenciación.\footnote{F. Mertens: ``Über eine Formel der Determinantentheorie.'' (Fórmula 5), Sitzb. d. Ak. d. Wiss. Wien, vol. 91, II. Abt. (1885), y más detalladamente (con aplicaciones) en: ``Über ganze Funktionen von $m$ Systemen von je $n$ Unbestimmten''. Monatsh. f. Math. u. Phys. 4.o año (1893). Para algunas de las demostraciones de Capelli (desde 1882), cf., por ejemplo, Math. Ann. 37, o las autografiadas ``Lezioni sulla teoria delle forme algebriche'' (1902). J. Deruyts: \emph{Essai d'une théorie générale des formes algébriques}, Mém. d. 1. Soc. Roy. des Sciences de Liège (2) 17 (1892).}

En II doy además una demostración más conceptual del teorema de reducción, tratando la cuestión como un problema sobre la equivalencia de familias lineales de formas. Este punto de vista, y también una parte esencial de las consideraciones usadas, los tomo de un trabajo ``Sobre sistemas algebraicos de módulos''\footnote{E. Fischer: Über algebraische Modulsysteme und lineare homogene partielle Differentialgleichungen mit konstanten Koeffizienten, \S\ 1--4, J. f. M. 140 (1911).} y de teoremas inéditos de E. Fischer conexos con él, cuyo uso me permitió amablemente.

Finalmente, en III, muestro cómo las consideraciones correspondientes conducen también a los desarrollos en serie más generales.

\subsection*{I.}
Sea $\theta$ una forma homogénea en cada una de las $N>n$ filas
\[
  A_1,A_2,\ldots,A_N,
\]
donde, en general, la fila $A_k$ consta de los $n$ elementos
\[
  A_k^{(1)},A_k^{(2)},\ldots,A_k^{(n)}.
\]
Los coeficientes de $\theta$ pueden ser indeterminadas o cantidades de un dominio de racionalidad dado. Denote además $P$ una función racional entera ---con coeficientes enteros racionales--- de los procesos polares
\[
  P_{hk}=\left(A_h\frac{\partial}{\partial A_k}\right)
  =A_h^{(1)}\frac{\partial}{\partial A_k^{(1)}}+
   A_h^{(2)}\frac{\partial}{\partial A_k^{(2)}}+
   \cdots+
   A_h^{(n)}\frac{\partial}{\partial A_k^{(n)}},
\]
donde el producto $P_{hk}P_{ij}\theta$ debe entenderse como la aplicación de $P_{hk}$ a $P_{ij}\theta$.

Entonces el teorema de reducción mencionado está dado por la identidad
\begin{equation}
  \theta=\sum PZ,
\tag{1}
\end{equation}
donde las formas $Z$ contienen sólo las filas
\[
  A_1,A_2,\ldots,A_n
\]
y se derivan de $\theta$ mediante procesos polares $P$.

Consideremos ahora el sistema de formas fundamentales $(F')$, compuesto por las $N$ formas del mismo orden y con el mismo número de variables
\[
  F_1,F_2,\ldots,F_N,
\]
cuyos coeficientes se toman respectivamente como las $N$ filas
\[
  A_1,A_2,\ldots,A_N.
\]
Sea $(J)$ el sistema completo ---compuesto por finitamente muchos invariantes--- de $F_1,
\ldots,F_n$, de tal modo que todo invariante de $F_1,
\ldots,F_n$ pueda expresarse racionalmente y enteramente mediante los invariantes $(J)$. La conjetura de Hilbert que se ha de probar es entonces que, para los invariantes simultáneos $S$ de $F_1,
\ldots,F_N$, un sistema completo está dado por los finitamente muchos invariantes $(J,PJ)$, donde $PJ$ denota todos los invariantes derivables de $J$ mediante procesos polares $P$.

En efecto, todo invariante simultáneo $S$ con coeficientes enteros racionales de $(F')$ es una forma homogénea en cada una de las $N$ filas $A_1,
\ldots,A_N$ ---con coeficientes enteros racionales--- y por tanto se le puede aplicar la identidad (1). Puesto que se sabe que la aplicación de los procesos polares $P$ a $S$ vuelve a producir invariantes, las formas $Z$ también son invariantes, a saber, invariantes simultáneos de $F_1,
\ldots,F_n$, ya que contienen sólo las filas $A_1,
\ldots,A_n$; de ahí que, por hipótesis, sean funciones racionales enteras de los invariantes $(J)$. Así, la identidad (1) se convierte en
\[
  S=\sum PY,
\]
donde las $Y$ denotan productos de potencias de los invariantes $(J)$. Pero, por la definición de $P$, la ejecución efectiva de los procesos $P$ sobre $Y$ consiste en la ejecución sucesiva, finitamente repetida, de las operaciones polares simples $P_{hk}$, lo que, por la regla de diferenciación de productos, conduce a sumas de productos de invariantes $(J)$ y $(PJ)$. Lo mismo vale para la aplicación de $P_{hk}$ a un producto formado por $(J)$ y $(PJ)$; por tanto, también $PY$ produce sólo combinaciones racionales enteras de $(J,PJ)$. \emph{Con ello queda probada la representabilidad racional entera de todos los invariantes $S$ mediante $(J,PJ)$.}

\subsection*{II.}
Antes de probar el teorema de reducción, recordamos algunos hechos sobre familias lineales de formas. Por una \emph{familia lineal de formas sobre $K$} entendemos un sistema de formas de la misma dimensión, completo en el sentido de que, junto con $\theta_1$ y $\theta_2$, contiene siempre $c_1\theta_1+c_2\theta_2$, donde $c_1,c_2$ son cantidades de un dominio de racionalidad prescrito $K$, y donde $K$ contiene el dominio de coeficientes de las $\theta$. Entre estas formas hay siempre un número finito ---digamos $\rho$--- linealmente independiente sobre $K$, mientras que cualesquiera $\rho+1$ formas satisfacen una relación lineal con coeficientes de $K$; la familia tiene \emph{rango $\rho$} respecto de $K$.\footnote{Se obtienen familias lineales de formas más generales ---cuyo rango no necesita ser finito--- suprimiendo la restricción de igual dimensión, o también suprimiendo la restricción de que $K$ contenga el dominio de coeficientes de las $\theta$.} Usaremos el hecho fácilmente visible de que si $\mathcal L$ y $\mathcal T$ son dos familias lineales sobre $K$ del mismo rango $\rho$, y si $\mathcal T$ es una subfamilia de $\mathcal L$, entonces $\mathcal L$ y $\mathcal T$ son idénticas (equivalentes).

Con esto pasamos al teorema de reducción. La identidad de I,
\[
  \theta=\sum PZ,
\]
pasa a una identidad análoga para $P\theta$ cuando se aplica el mismo proceso polar $P$ a ambos lados; así, el teorema de reducción da al mismo tiempo una representación de todas las polares de $\theta$ por sumas de las polares especiales $PZ$. Como, recíprocamente, estas polares especiales están ciertamente contenidas entre todas las polares, el teorema de reducción afirma la equivalencia de estas dos familias lineales de formas; en particular, también afirma la equivalencia de aquellas subfamilias cuyas formas tienen, en cada fila individual, la misma dimensión que $\theta$. Para probar esta equivalencia introducimos una familia intermedia determinada por las formas $Z$. Por simplicidad damos primero la demostración para el caso de tres filas binarias $(N=3,n=2)$ y luego indicamos brevemente la demostración general, completamente paralela.

Sea entonces $\theta(x,y,z)$ de dimensiones $\alpha,\beta,p$, respectivamente, en las filas
\[
  x_1x_2;\quad y_1y_2;\quad z_1z_2;
\]
y sean primero los coeficientes de $\theta$ indeterminadas $a$ (o números racionales). Consideramos las tres familias lineales de formas siguientes.

1. La familia $\mathcal L$, compuesta por todas las formas $f(x,y,z)$ que surgen de $\theta$ mediante procesos polares $P$ y que, en las filas individuales $x,y,z$, tienen la misma dimensión que $\theta$; en particular, $\mathcal L$ contiene a $\theta$. Si $f(x,y,z)$ se considera como forma en $x,y,z$ y las indeterminadas $a$, entonces el cuerpo de coeficientes es el cuerpo $R$ de los números racionales; también para $K$ debe tomarse $R$, ya que sólo para $\theta_1,\theta_2$ racionales enteras se vuelve $c_1P_1+c_2P_2$ nuevamente un proceso $P$; sea $\mathcal L$ de rango $\rho$ respecto de $R$.

2. La familia $\mathcal S$, compuesta por todas las formas $g(\lambda;x,y)$ definidas por
\[
  g(\lambda;x,y)=f(x,y,\lambda_1x+\lambda_2y),
\]
o, expresadas por procesos polares,
\[
\begin{aligned}
  g(\lambda;x,y)
  &=\frac1{p!}\left[(\lambda_1x+\lambda_2y)\frac{\partial}{\partial z}\right]^p f(x,y,z) \\
  &=g_0(xy)\lambda_1^p+g_1(x,y)\lambda_1^{p-1}\lambda_2+\cdots+g_p(xy)\lambda_2^p,
\end{aligned}
\]
donde
\[
  i!(p-i)!\,g_i(xy)=
  \left(y\frac{\partial}{\partial z}\right)^i
  \left(x\frac{\partial}{\partial z}\right)^{p-i}f(xyz).
\]
\footnote{Como en I,
\[
 \left(t\frac{\partial}{\partial x}\right)=t_1\frac{\partial}{\partial x_1}+t_2\frac{\partial}{\partial x_2},
\]
y por tanto, en particular,
\[
 \left[(\lambda_1x+\lambda_2y)\frac{\partial}{\partial z}\right]
 = (\lambda_1x_1+\lambda_2y_1)\frac{\partial}{\partial z_1}
 + (\lambda_1x_2+\lambda_2y_2)\frac{\partial}{\partial z_2},
\]
\[
 \left[z\left(\frac{\partial}{\partial\lambda_1}\frac{\partial}{\partial x}
 +\frac{\partial}{\partial\lambda_2}\frac{\partial}{\partial y}\right)\right]
 =z_1\left(\frac{\partial}{\partial\lambda_1}\frac{\partial}{\partial x_1}
 +\frac{\partial}{\partial\lambda_2}\frac{\partial}{\partial y_1}\right)
 +z_2\left(\frac{\partial}{\partial\lambda_1}\frac{\partial}{\partial x_2}
 +\frac{\partial}{\partial\lambda_2}\frac{\partial}{\partial y_2}\right).
\]}
Así, los $g_i(xy)$ dan formas $Z$ de la identidad (1). Las $g$ son formas racionales enteras en $x,y,\lambda$ y las indeterminadas $a$; sea $\mathcal S$ de rango $\sigma$ respecto de $R$.

3. La familia $\mathcal T$, obtenida de $\mathcal S$ por el proceso polar inverso del mismo modo que $\mathcal S$ se obtiene de $\mathcal L$; las formas $h$ de $\mathcal T$ están definidas por
\[
\begin{aligned}
 h(x,y,z)
 &=\frac1{p!}\left[z\left(\frac{\partial}{\partial\lambda_1}\frac{\partial}{\partial x}
 +\frac{\partial}{\partial\lambda_2}\frac{\partial}{\partial y}\right)\right]^p g(\lambda;xy)\\
 &=h_0(xyz)+h_1(xyz)+\cdots+h_p(xyz),
\end{aligned}
\]
donde, por 2.,
\[
  h_i(xyz)=
  \left(z\frac{\partial}{\partial x}\right)^{p-i}
  \left(z\frac{\partial}{\partial y}\right)^i g_i(xy).
\]
Las $h_i$ individuales, y por consiguiente también $h$, son por tanto formas $PZ$ y sumas de ellas; sea $\mathcal T$ de rango $\tau$ respecto de $R$.

Como muestra la construcción de las formas $h(x,y,z)$ de $\mathcal T$, estas formas surgen de $\theta$ mediante procesos polares $P$, y en las filas $x,y,z$ tienen individualmente la misma dimensión que $\theta$; de ahí que la familia $\mathcal T$ sea una subfamilia de $\mathcal L$. Por el hecho mencionado arriba, sólo queda probar la igualdad de los rangos. En esta prueba no se usa ya la estructura especial de los $f(xyz)$, a saber, que son polares de una y la misma forma $\theta$.

\emph{a) Para comparar los rangos de $\mathcal L$ y $\mathcal S$ usamos el Lema a): de $f(x,y,\lambda_1x+\lambda_2y)=0$ se sigue necesariamente que $f(x,y,z)=0$.} Esto se sigue inmediatamente de la relación idéntica entre tres filas binarias
\[
  (xy)z+(yz)x+(zx)y=0,
  \qquad (xy)=(x_1y_2-x_2y_1),\ldots,
\]
que implica
\[
  f[x,y,(xy)x+(xz)y]=(xy)^p f(x,y,z).
\]
Así, de $f(x,y,\lambda_1x+\lambda_2y)=0$ (idénticamente en $\lambda$) se sigue, como muestra la sustitución $\lambda_1=(zy)$, $\lambda_2=(xz)$, y puesto que $(xy)\ne0$, que necesariamente $f(x,y,z)=0$.

Por este lema hay una correspondencia biunívoca entre las formas de $\mathcal S$ y las de $\mathcal L$. Pues de
\[
  f_1(x,y,\lambda_1x+\lambda_2y)=g(\lambda;xy),\qquad
  f_2(x,y,\lambda_1x+\lambda_2y)=g(\lambda;xy)
\]
se sigue necesariamente que
\[
  f_1(xyz)=f_2(xyz).
\]
Además, toda relación en $\mathcal S$ se conserva entre las formas correspondientes de manera única en $\mathcal L$ (y recíprocamente). Pues de
\[
  c_1g^{(1)}(\lambda;xy)+c_2g^{(2)}(\lambda;xy)+\cdots+c_rg^{(r)}(\lambda;xy)=0
\]
se sigue por el lema que necesariamente
\[
  c_1f^{(1)}(x,y,z)+c_2f^{(2)}(x,y,z)+\cdots+c_rf^{(r)}(x,y,z)=0.
\]
\emph{Esto prueba la igualdad de los rangos de $\mathcal L$ y $\mathcal S$; $\rho=\sigma$.}

\emph{b) Para comparar los rangos de $\mathcal S$ y $\mathcal T$ usamos el Lema b) correspondiente: de $h(x,y,z)=0$ se sigue necesariamente que $g(\lambda;xy)=0$.} Para la demostración, obsérvese que, por 3., $h(x,y,z)=0$ es equivalente a la anulación de los productos de potencias individuales
\[
 \left(\frac{\partial}{\partial\lambda_1}\frac{\partial}{\partial x_1}
 +\frac{\partial}{\partial\lambda_2}\frac{\partial}{\partial y_1}\right)^{p_1}
 \left(\frac{\partial}{\partial\lambda_1}\frac{\partial}{\partial x_2}
 +\frac{\partial}{\partial\lambda_2}\frac{\partial}{\partial y_2}\right)^{p_2}g(\lambda;x,y),
 \qquad p_1+p_2=p.
\]
Una diferenciación ulterior da entonces también la anulación de los productos de potencias
\[
\begin{aligned}
&\left(\frac{\partial}{\partial x_1}\right)^{\alpha_1}
 \left(\frac{\partial}{\partial x_2}\right)^{\alpha_2}
 \left(\frac{\partial}{\partial y_1}\right)^{\beta_1}
 \left(\frac{\partial}{\partial y_2}\right)^{\beta_2} \\
&\quad\cdot
 \left(\frac{\partial}{\partial\lambda_1}\frac{\partial}{\partial x_1}
 +\frac{\partial}{\partial\lambda_2}\frac{\partial}{\partial y_1}\right)^{p_1}
 \left(\frac{\partial}{\partial\lambda_1}\frac{\partial}{\partial x_2}
 +\frac{\partial}{\partial\lambda_2}\frac{\partial}{\partial y_2}\right)^{p_2}
 g(\lambda;xy).
\end{aligned}
\]
Por consiguiente, toda combinación lineal de estos productos de potencias se anula, y por tanto toda forma
\[
  f\left(\frac{\partial}{\partial x},\frac{\partial}{\partial y},
  \frac{\partial}{\partial\lambda_1}\frac{\partial}{\partial x}
  +\frac{\partial}{\partial\lambda_2}\frac{\partial}{\partial y}\right)g(\lambda,xy).
\]
Escogiendo $f$ en particular de modo que
\[
  f(x,y,\lambda_1x+\lambda_2y)=g(\lambda;xy),
\]
se obtiene también
\[
  g\left(\frac{\partial}{\partial x};
    \frac{\partial}{\partial x},
    \frac{\partial}{\partial y}\right)g(\lambda,x,y)=0;
\]
o bien, si
\[
  g(\lambda,xy)=\sum G_i\lambda_1^{\nu_1}\lambda_2^{\nu_2}
  x_1^{\alpha_1}x_2^{\alpha_2}y_1^{\beta_1}y_2^{\beta_2},
\]
donde los $G_i$ son formas lineales racionales enteras en las indeterminadas $a$, entonces
\[
  \sum \nu_1!\nu_2!\alpha_1!\alpha_2!\beta_1!\beta_2!\,G_i^2=0,
\]
y por tanto $g(\lambda;xy)=0$.

Del Lema b) se sigue, exactamente como en a), que \emph{los rangos de $\mathcal S$ y $\mathcal T$ son iguales; $\sigma=\tau$.}

Así, por a) y b), se tiene $\rho=\tau$, y por consiguiente ---como $\mathcal T$ es una subfamilia de $\mathcal L$--- queda probada la equivalencia de estas dos familias lineales. Por ello toda forma de $\mathcal L$, en particular también $\theta$, puede representarse como una combinación lineal racional-entera de $\rho$ formas linealmente independientes $h(x,y,z)$:
\[
  \theta=\sum c_i h^{(i)}(x,y,z)=\sum PZ.
\]
Esta identidad, obtenida para coeficientes indeterminados $a$ de $\theta$, sigue siendo válida cuando las indeterminadas $a$ se sustituyen por cantidades de cualquier dominio de racionalidad; por tanto, \emph{el teorema de reducción queda probado en general para el caso de tres filas binarias.}

La demostración general, para $N$ filas de $n$ variables cada una, es completamente paralela. Sea $\theta(A)$ de dimensión $\alpha_i$ en la fila $A_i$. Consideramos de nuevo las tres familias lineales de formas:

1. la familia $\mathcal L$, compuesta por todas las formas $f(A)$ que surgen de $\theta(A)$ mediante procesos polares $P$ y, en cada fila individual $A_i$, tienen la misma dimensión que $\theta(A)$;

2. la familia $\mathcal S$, compuesta por todas las formas $g(\lambda;A_1\cdots A_n)$ que surgen de $f(A)$ mediante las sustituciones
\[
\begin{aligned}
 A_{n+1}&=\lambda_{n+1}^{(1)}A_1+\lambda_{n+1}^{(2)}A_2+
        \cdots+\lambda_{n+1}^{(n)}A_n,\\
 &\vdotswithin{=}\\
 A_N&=\lambda_N^{(1)}A_1+\lambda_N^{(2)}A_2+
        \cdots+\lambda_N^{(n)}A_n;
\end{aligned}
\]
en este proceso, los coeficientes de los productos de potencias individuales de las $\lambda$ en $g(\lambda;A_1\cdots A_n)$ dan formas $Z$ de la identidad (1);

3. la familia $\mathcal T$, compuesta por todas las formas $h(A)$ definidas por
\[
\begin{aligned}
 h(A)=&\left[A_{n+1}\left(
 \frac{\partial}{\partial\lambda_{n+1}^{(1)}}\frac{\partial}{\partial A_1}
 +\cdots+
 \frac{\partial}{\partial\lambda_{n+1}^{(n)}}\frac{\partial}{\partial A_n}\right)\right]^{\alpha_{n+1}} \\
&\cdots
\left[A_N\left(
 \frac{\partial}{\partial\lambda_N^{(1)}}\frac{\partial}{\partial A_1}
 +\cdots+
 \frac{\partial}{\partial\lambda_N^{(n)}}\frac{\partial}{\partial A_n}\right)\right]^{\alpha_N}
 g(\lambda;A_1\cdots A_n).
\end{aligned}
\]
Aquí $h(A)$ es una suma de formas $PZ$.

A causa de la relación idéntica entre cualesquiera $n+1$ filas, el Lema a) sigue siendo válido; igualmente lo sigue siendo el Lema b), ya que el número de variables no entró en su demostración. Así $\mathcal L$ y $\mathcal T$ tienen el mismo rango y ---como $\mathcal T$ es una subfamilia de $\mathcal L$--- son idénticas. En consecuencia, la identidad
\[
  \theta=\sum PZ
\]
vale para coeficientes indeterminados y, por ello, para coeficientes que sean cantidades de cualquier dominio de racionalidad; \emph{con ello queda probado en general el teorema de reducción.}

\subsection*{III.}
Indicamos todavía brevemente cómo consideraciones análogas producen también el \emph{desarrollo en serie general} (para $N\le n$). Sea $Z$ una forma separadamente homogénea en las $n$ filas $A_1,
\ldots,A_n$ de $n$ cantidades cada una; y sea $\Delta$ el determinante de estas filas. Probamos primero la identidad correspondiente a (1):
\begin{equation}
  Z=\sum PH \pmod{\Delta},
\tag{2}
\end{equation}
donde las formas $H$ contienen sólo las filas $A_1,
\ldots,A_{n-1}$ y se derivan de $Z$ mediante procesos $P$.

La demostración consiste en convertir las relaciones obtenidas en II en congruencias módulo $\Delta$. Por simplicidad tomemos tres filas ternarias $x,y,z$; se supone que los coeficientes son indeterminados.

Las tres familias lineales $\mathcal L,\mathcal S,\mathcal T$ se definen como en II para el caso de filas binarias; sean sus rangos $\rho,\sigma,\tau$, mientras que $\rho^*$ y $\tau^*$ denotan los rangos de $\mathcal L$ y $\mathcal T$ módulo $\Delta$ (es decir, hay $\rho^*$, pero no $\rho^*+1$, formas de $\mathcal L$ linealmente independientes mod $\Delta$). Entonces, como en II, se tiene

\emph{Lema a): de $f(x,y,\lambda_1x+\lambda_2y)=0$ se sigue necesariamente que $f(xyz)\equiv0\pmod{\Delta}$ (y recíprocamente).} Pues, por la relación entre cuatro filas ternarias,
\[
  \Delta_1x-\Delta_2y+\Delta_3z=\Delta t,
  \qquad \Delta_1=(yzt),\ldots,
\]
siempre existe una relación lineal módulo $\Delta$ entre tres filas ternarias:
\[
  \Delta_1x-\Delta_2y+\Delta_3z\equiv0\pmod{\Delta},
\]
y por tanto, como en II,
\[
  f(x,y,-\Delta_1x+\Delta_2y)=\Delta_3^p f(x,y,z)\pmod{\Delta}.
\]
Se sigue de $f(x,y,\lambda_1x+\lambda_2y)=0$ (idénticamente en $\lambda$), puesto que $\Delta$ es irreducible y $\Delta_3$ no es divisible por $\Delta$, que necesariamente $f(x,y,z)\equiv0\pmod{\Delta}$. La recíproca es clara a partir de $\Delta(x,y,\lambda_1x+\lambda_2y)=0$. De ahí, como en II, que $\rho^*=\sigma$.

El Lema b) sigue siendo válido con su demostración, y por tanto se tiene $\sigma=\tau$.

Para probar $\tau=\tau^*$ usamos

\emph{Lema c): de $h(x,y,z)\equiv0\pmod{\Delta}$ se sigue necesariamente que $h(x,y,z)=0$.} En efecto, como polar, $h(x,y,z)$ se anula bajo la aplicación del conocido proceso $\Omega$; esto se expresa por la ecuación diferencial
\[
  \Delta\!\left(\frac{\partial}{\partial x},\frac{\partial}{\partial y},\frac{\partial}{\partial z}\right)h(x,y,z)=0.
\]
Si ahora $h(x,y,z)=\Delta(x,y,z)\cdot k(x,y,z)$, entonces por nueva diferenciación se tiene también la anulación de
\[
  k\!\left(\frac{\partial}{\partial x},\frac{\partial}{\partial y},\frac{\partial}{\partial z}\right)
  \Delta\!\left(\frac{\partial}{\partial x},\frac{\partial}{\partial y},\frac{\partial}{\partial z}\right)
  =
  h\!\left(\frac{\partial}{\partial x},\frac{\partial}{\partial y},\frac{\partial}{\partial z}\right)h(x,y,z),
\]
y por tanto de $h(x,y,z)$. Así $\rho^*=\sigma=\tau=\tau^*$; como $\mathcal T$ es una subfamilia de $\mathcal L$, queda probada la identidad (2). El mismo argumento da la demostración para $n$ variables.\footnote{Después de 3. se tiene la relación correspondiente, que se anula a causa del factor determinante, para $\Omega(h)$; en notación operatoria es la ecuación que surge de productos de
\[
 \left(\lambda_1\frac{\partial}{\partial \xi}+\lambda_2\frac{\partial}{\partial \eta}\right),
 \qquad
 \left[z\left(\frac{\partial}{\partial\lambda_1}\frac{\partial}{\partial \xi}
 +\frac{\partial}{\partial\lambda_2}\frac{\partial}{\partial \eta}\right)\right],
\]
y se anula por la anulación del factor determinante.}

Aplicando la identidad (2) al multiplicador de $\Delta$, se obtiene un desarrollo
\[
  Z=\varphi_0+\Delta\varphi_1+\Delta^2\varphi_2+\cdots+\Delta^\mu\varphi_\mu,
\]
donde las $\varphi_i$ se anulan cada una bajo la aplicación del proceso $\Omega$. La repetición $i$ veces del proceso da, puesto que en general $\Omega(\Delta\cdot\psi)=c_1\psi+c_2\Delta\cdot\Omega(\psi)$,
\[
  c\cdot \Omega^i(Z)\equiv\varphi_i\pmod{\Delta};
\]
por otra parte, por (2) se tiene
\[
  c\cdot\Omega^i(Z)\equiv\sum PH_i\pmod{\Delta},
\]
donde los $H_i$ se derivan de la forma $\Omega^i(Z)$ del mismo modo que $H$ se deriva de $Z$. Por el Lema c) (extendido a $n$ variables) se obtiene entonces $\varphi_i=\sum PH_i$, y por consiguiente
\begin{equation}
  Z=\sum PH+\Delta\cdot\sum PH_1+\Delta^2\cdot\sum PH_2+
  \cdots+\Delta^\mu\cdot\sum PH_\mu
\tag{3}
\end{equation}
\emph{como el desarrollo en serie más general de una forma en $n$ filas de $n$ variables cada una, según polares de formas en sólo $n-1$ filas.}

El desarrollo en serie para formas en $N$ filas de $n$ variables $(N<n)$ se obtiene entonces, como es sabido, mediante una sustitución simple. Ponemos (para $N=3$, $n>3$)
\[
  u=\xi_1x+\xi_2y+\xi_3z,
  \quad v=\eta_1x+\eta_2y+\eta_3z,
  \quad w=\zeta_1x+\zeta_2y+\zeta_3z,
\]
y desarrollamos $H(u,v,w)=Z(\xi,\eta,\zeta)$ según (3). Entonces, puesto que $\xi,\eta,\zeta$ sólo aparecen en las combinaciones $u,v,w$, se tiene
\[
  \left(\eta\frac{\partial}{\partial \xi}\right)=
  \left(v\frac{\partial}{\partial u}\right)\cdots,
  \qquad
  \Omega=\sum_{ikl}\left(\frac{\partial}{\partial u}
  \frac{\partial}{\partial v}
  \frac{\partial}{\partial w}\right)_{ikl}(xyz)_{ikl}
  =\nabla_{uvw}(xyz).
\]
Bajo la especialización $\xi_1=\eta_2=\zeta_3=1$; $\xi_2=\xi_3=\cdots=\zeta_2=0$, $H(u,v,w)$ pasa a $H(x,y,z)$, y $\nabla_{uvw}(xyz)$ a $\nabla_{xyz}(xyz)=\nabla_{xyz}$ salvo un factor numérico, ya que la aplicación de $(y\partial/\partial x)$, etc., a los determinantes $(xyz)_{ikl}$ los hace anularse. Puesto que la afirmación correspondiente vale para $N$ filas, (3) da el desarrollo
\begin{equation}
  H=\sum P\Phi+\sum P\Phi_1+\cdots+\sum P\Phi_\mu,
\tag{4}
\end{equation}
donde las $\Phi_i$ surgen de $\nabla^i_{A_1\ldots A_N}H$ mediante procesos $P$ y contienen sólo las filas $A_1,
\ldots,A_{N-1}$, aparte de las combinaciones determinantes que aparecen en $\nabla^i$, las cuales, como se mencionó, no entran en consideración bajo la polarización.

Si esos determinantes en $\nabla^i$ se sustituyen entonces por indeterminadas $p$, las formas resultantes pueden desarrollarse de nuevo según (4); de este modo se obtiene finalmente un desarrollo
\begin{equation}
  H=\left[\sum P\psi(p)\right]_{p=\text{determinantes de }A_1\cdots A_N},
\tag{5}
\end{equation}
donde las $\psi$ surgen de la forma original mediante los procesos $P$ y $\nabla_{A_1\ldots A_N}(p)$, $\nabla_{A_1\ldots A_{N-1}}(p)$, etc. Con estos desarrollos y sus combinaciones ---por ejemplo, aplicando (3), y luego (4) y (5), a las formas $Z$ que aparecen en (1)--- se agotan todos los desarrollos conocidos para formas en arbitrariamente muchas filas de $n$ variables cogredientes cada una.

Erlangen, 5 de enero de 1915.

\clearpage
\section*{9. Los dominios más generales de números trascendentes enteros}
\begin{center}
Math. Ann. 77 (1916), pp. 103--128
\end{center}

Mediante el teorema del buen orden, E. Zermelo construyó un dominio de ``números trascendentes enteros'', es decir, un dominio numérico $\mathfrak G$ con las siguientes propiedades abstractas:
\begin{enumerate}
\item[I.] La suma, la diferencia y el producto de dos números de $\mathfrak G$ vuelven a ser números de $\mathfrak G$.
\item[II.] Todo número real o complejo es cociente de dos números de $\mathfrak G$.
\item[III.] Todo número racional entero (respectivamente algebraico entero) es elemento de $\mathfrak G$.
\item[IV.] Ningún número racional no entero (respectivamente algebraico no entero) pertenece a $\mathfrak G$.
\footnote{E. Zermelo, Über ganze transzendente Zahlen, Math. Ann. 75, p. 434 (1914).}
\end{enumerate}

La construcción se apoya en el hecho de que el teorema del buen orden implica la existencia de una \emph{base algebraica} de todos los números: a saber, un \emph{sistema} $H$ \emph{de números} $\eta$ entre los cuales no existen relaciones algebraicas, pero por medio de los cuales todos los demás números pueden expresarse algebraicamente. A causa de su independencia algebraica, estos números de base $\eta$ pueden considerarse como indeterminadas, y por tanto el dominio buscado se convierte en un dominio íntegro de funciones racionales y algebraicas de las indeterminadas $\eta$,
\footnote{Por esta razón, las consideraciones del presente trabajo corren en parte paralelas a las del trabajo anterior de la autora: Körper und Systeme rationaler Funktionen, Math. Ann. 76, p. 161 (1915).}
cuyos coeficientes pertenecen al cuerpo $K$ de todos los números algebraicos. El dominio especial $\mathfrak G_\eta$ construido por Zermelo queda entonces caracterizado por las siguientes propiedades relativas a la base:

$\mathfrak G_\eta$ contiene:
\begin{enumerate}
\item[1)] todos los números de base $\eta$;
\item[2)] todos los polinomios en los $\eta$ con coeficientes enteros;
\footnote{En todo este trabajo, ``entero'' significará que los coeficientes pertenecen al dominio íntegro $[K]$ de todos los enteros algebraicos. Para la definición de $\mathfrak G_\eta$, en 1)--5) bastaría considerar sólo coeficientes enteros racionales; para dominios más generales, sin embargo, las propiedades de base serían entonces menos simples.}
\item[3)] todas las funciones algebraicas enteras de los $\eta$ con coeficientes enteros.
\end{enumerate}
$\mathfrak G_\eta$ excluye:
\begin{enumerate}
\item[4)] todas las funciones racionalmente fraccionarias de los $\eta$ con coeficientes enteros;
\item[5)] todas las funciones algebraicamente fraccionarias de los $\eta$ con coeficientes enteros.
\footnote{Zermelo, loc. cit., § 3.}
\end{enumerate}

Es ahora fácil ver (compárense los ejemplos de §§ 2 y 3) que existen dominios $\mathfrak G$ esencialmente distintos del dominio de Zermelo $\mathfrak G_\eta$: es decir, dominios $\mathfrak G$ para los cuales, bajo ningún buen orden, valen simultáneamente todas las propiedades de base 1)--5), y que por consiguiente no pueden generarse mediante la construcción de Zermelo bajo ningún buen orden. En otras palabras, existen dominios $\mathfrak G$ que no pueden aplicarse biyectiva e isomórficamente sobre $\mathfrak G_\eta$. Así surge la cuestión de los dominios más generales de números trascendentes enteros; será respondida por completo en lo que sigue.

Para ello hay que determinar primero cuáles de las propiedades de base son consecuencia de las propiedades definitorias abstractas y cuáles proceden de la construcción especial. Se muestra (§ 1) que para todo dominio prescrito $\mathfrak G$ existe un buen orden tal que se satisfacen las propiedades de base 1) y 2); por tanto, todo dominio $\mathfrak G$ puede construirse mediante una base algebraica de todos los números. Ninguna de las propiedades de base restantes tiene por qué cumplirse (§§ 2 y 3). En particular se sigue que una propiedad abstracta adicional del dominio de Zermelo no pertenece a todos los dominios $\mathfrak G$: la cerradura algebraicamente entera. Puede formularse así:
\begin{enumerate}
\item[V.] Todo número que depende algebraicamente entero e integralmente de un número finito de elementos de $\mathfrak G$ pertenece a $\mathfrak G$.
\end{enumerate}

Los dominios especiales para los cuales, además de I--IV, vale V se llamarán dominios $\mathfrak H$ de números trascendentes algebraicamente enteros. Para los dominios más generales $\mathfrak H$ se obtienen los resultados simples (§ 4):
\begin{enumerate}
\item[a)] La intersección de todos los dominios $\mathfrak H$ construibles por medio de una misma base $H$ está dada por todas las funciones algebraicas enteras de la base, es decir, por el dominio de Zermelo $\mathfrak G_\eta$, que él mismo es un dominio $\mathfrak H$. (Esto proporciona también una caracterización abstracta del dominio de Zermelo.)
\item[b)] Todo dominio $\mathfrak H$ surge por adjunción algebraicamente entera\footnote{Explicación de esta expresión en § 4.} de un sistema arbitrario $\mathfrak S(\eta)$ a $\mathfrak G_\eta$, donde $\mathfrak S(\eta)$ debe satisfacer sólo la condición única de que por la adjunción no entre en el dominio ningún número algebraicamente fraccionario.
\end{enumerate}

Los resultados para los dominios más generales $\mathfrak G$, cuando se toma como fundamento una base algebraica, son menos simples. Aquí la intersección de todos los dominios $\mathfrak G$ no es ella misma un dominio $\mathfrak G$, y por tanto la construcción de los dominios más generales es más difícil de abarcar (§ 5).

Para obtener resultados análogos a los de los dominios $\mathfrak H$, se sustituye la base algebraica por una \emph{base racional} $\Theta$, es decir, por un \emph{sistema} $\Theta$ \emph{de números} $\vartheta$ que permite expresar racionalmente todos los números, mientras que, bajo al menos un buen orden, ningún número de base puede expresarse racionalmente por un número finito de los precedentes. Esta base racional $\Theta$ debe someterse aún a la condición lateral de que el dominio íntegro $[\Theta]$ formado por todos los polinomios enteros en los $\vartheta$ no contenga ningún número algebraicamente fraccionario. (La construcción de tal base con la condición lateral se da en § 7.) Entonces, para los dominios más generales $\mathfrak G$ ---que también podrían llamarse dominios de números trascendentes racionalmente enteros--- se obtienen de nuevo los resultados simples (§ 6):
\begin{enumerate}
\item[a)] La intersección de todos los dominios $\mathfrak G$ construibles por medio de una misma base racional $\Theta$ está dada por todas las funciones racionales enteras de la base, es decir, por el dominio $[\Theta]$, que él mismo es un dominio $\mathfrak G$.
\item[b)] Todo dominio $\mathfrak G$ surge por adjunción racionalmente entera de un sistema arbitrario $\mathfrak S(\vartheta)$ a $[\Theta]$, donde $\mathfrak S(\vartheta)$ debe satisfacer sólo la condición única de que por la adjunción no entre en el dominio ningún número algebraicamente fraccionario.
\end{enumerate}

Para tener una analogía completa con los dominios $\mathfrak H$, habría que omitir los números algebraicos de las propiedades definitorias III y IV para los dominios $\mathfrak G$. Con esta modificación de III y IV, la construcción de los elementos enteros mediante una base racional puede trasladarse a un cuerpo arbitrario (§ 10), mientras que la construcción de Zermelo presupone la cerradura algebraica del cuerpo.

Cuando $H$ recorre todas las bases algebraicas, respectivamente cuando $\Theta$ recorre todas las bases racionales que satisfacen la condición lateral, los resultados a) y b) dan la totalidad de todos los dominios $\mathfrak H$ y $\mathfrak G$. Si todos los dominios isomorfos se reúnen en una clase, entonces de esta totalidad puede seleccionarse un subconjunto que represente la totalidad de todas las clases ---por § 9, al menos una infinitud numerable (§ 8).

\subsection*{§ 1. Demostración de las propiedades de base 1) y 2) para todos los dominios $\mathfrak G$.}

Sea $\mathfrak G$ un dominio cualquiera prescrito de números trascendentes enteros, de modo que tenga las propiedades abstractas I--IV. Siguiendo el procedimiento aplicado por E. Zermelo al cuerpo de todos los números complejos, escogemos en $\mathfrak G$ una \emph{base algebraica} $H$ \emph{de todos los números de $\mathfrak G$}. Puesto que por II todo número real o complejo puede representarse como cociente de dos números de $\mathfrak G$, $H$ es al mismo tiempo una base algebraica de \emph{todos} los números. Así, para todo dominio prescrito $\mathfrak G$, la base algebraica de todos los números puede escogerse de modo que pertenezca a $\mathfrak G$; por tanto se satisface la propiedad de base 1). Además, por III, $\mathfrak G$ contiene el dominio íntegro $[K]$ de todos los enteros algebraicos, y por tanto, por I, también todos los polinomios en los elementos de base $\eta$ con coeficientes de $[K]$, es decir, todos los polinomios enteros en los $\eta$, que forman el dominio íntegro $[H]$. Por consiguiente, la propiedad de base 2) también se cumple siempre; en otras palabras, todo dominio $\mathfrak G$ puede construirse mediante una base algebraica $H$ de todos los números, de tal modo que $\mathfrak G$ contenga el dominio íntegro $[H]$ como divisor.

\textbf{Observación.} No obstante, partiendo de una base $H$, se puede naturalmente construir un dominio $\mathfrak G$ para el cual las propiedades de base 1) y 2) no valgan respecto de $H$ mismo, sino respecto de alguna otra base $H'$. Por ejemplo, sea $\mathfrak G'$ el dominio que se obtiene del dominio de Zermelo $\mathfrak G_\eta$ reemplazando, en toda la construcción, algún elemento de base $\eta$ por un polinomio $f(\eta)$. Entonces $\mathfrak G'$ no contiene ni $\eta$ ni $[H]$; pero si en la base $H$ se reemplaza el elemento de base $\eta$ por $\xi=f(\eta)$ ---con lo cual $H$ pasa de nuevo a una base algebraica $H'$---, las propiedades de base 1) y 2) quedan satisfechas para $\mathfrak G'$ respecto de $H'$.

\subsection*{§ 2. Exclusión de las propiedades de base 4) y 5).}

Mostramos ahora que los dominios $\mathfrak G$ no necesitan satisfacer ninguna de las restantes propiedades de base. Primero se excluyen 4) y 5) reemplazando, en la construcción de Zermelo, el dominio íntegro $[H]$ por un dominio de ``funcionales'' racionales enteros.

Según Weber,\footnote{H. Weber, Lehrbuch der Algebra. Kleine Ausgabe §§ 96 y 97.} por \emph{funcional racional} se entiende toda función racional con coeficientes numéricos racionales en la representación unívocamente determinada
\[
 A(\eta_1\cdots \eta_t)
   = a\cdot\frac{E_1(\eta_1\cdots \eta_t)}{E_2(\eta_1\cdots \eta_t)},
\]
donde $E_1,E_2$ son polinomios primitivos primos entre sí con coeficientes enteros racionales, y $a$ es un número racional positivo (el valor absoluto de $A$). Si en particular $a$ es un \emph{entero racional}, entonces $A$ se llama un \emph{funcional racional entero}. Como casos especiales, los funcionales racionales enteros incluyen los enteros racionales y los polinomios con coeficientes enteros racionales, mientras que los números racionalmente fraccionarios y los polinomios con coeficientes racionalmente fraccionarios deben contarse entre los funcionales racionales fraccionarios.

Sea ahora el dominio $\mathfrak G$ la totalidad $\mathfrak F$ de todos los funcionales racionales enteros de un número finito de elementos de base $\eta$ cada vez, junto con todas las funciones que dependen algebraicamente enteras de $\mathfrak F$ ---es decir, que satisfacen alguna ecuación cuyo coeficiente principal es la unidad y cuyos coeficientes restantes son funcionales racionales enteros.

La demostración de las propiedades I--IV para $\mathfrak G$ corre exactamente paralela a la demostración correspondiente de Zermelo y sólo se indicará brevemente. Primero, II y III son claros, pues $\mathfrak G$ contiene como divisor al dominio de Zermelo $\mathfrak G_\eta$. Por el teorema de Gauss sobre polinomios con coeficientes enteros racionales, sumas, diferencias y productos de funcionales racionales enteros son de nuevo funcionales racionales enteros; por tanto $\mathfrak F$ forma un dominio íntegro, y por argumentos conocidos\footnote{Compárese, por ejemplo, Weber, § 97, 8.} $\mathfrak G$ también se convierte en un dominio íntegro; así se cumple I. Además, del teorema de Gauss extendido\footnote{Weber, § 20, 5.} se deducen las dos proposiciones auxiliares: ``Si $z$ depende algebraicamente entero de $\mathfrak F$ por medio de \emph{alguna} ecuación, entonces también lo hace por medio de la ecuación \emph{irreducible} que $z$ satisface respecto del cuerpo de todos los funcionales racionales'',\footnote{Compárese Weber, § 97, 4.} y ``la ecuación irreducible sobre el cuerpo de todos los números racionales y satisfecha por un número algebraico es también irreducible sobre el cuerpo de todos los funcionales racionales''. Así $\mathfrak G$ no puede contener ningún número algebraico fraccionario; de este modo se ha verificado IV, y por tanto todas las propiedades abstractas I--IV.

Por otra parte, bajo ningún buen orden puede escogerse una base de $\mathfrak G$ de modo que $\mathfrak G$ excluya todas las funciones racional y algebraicamente fraccionarias de los elementos de base. Sea $z(\eta)$ una función cualquiera del dominio que haya de escogerse como elemento de base; entonces está definida por una ecuación
\[
 z^k+A_1(\eta)z^{k-1}+\cdots+A_k(\eta)=0,
\]
donde
\[
 A_i(\eta)=a_i\cdot\frac{E_1(\eta)}{E_2(\eta)}
\]
son funcionales racionales enteros. Entonces la función $\frac{a_k}{z}$ pertenece a $\mathfrak G$, y lo mismo $\frac{a_k}{\sqrt z}$, $\frac{a_k}{\sqrt[3]z}$, etc.\footnote{De modo completamente análogo, toda función algebraica de los $\eta$ puede transformarse en un elemento del dominio funcional $\mathfrak G$ multiplicándola por un entero racional. Así, este dominio tiene la propiedad de que todo número complejo puede representarse como el cociente de un elemento de $\mathfrak G$ y un entero racional, en analogía con los números algebraicos.} Por tanto, las propiedades de base 4) y 5) quedan excluidas para la totalidad de los dominios $\mathfrak G$.

\textbf{Observación.} Un ejemplo en § 3 mostrará que $\mathfrak G$ puede contener, bajo todo buen orden, funcionales racionalmente fraccionarios sin contener un número racional o algebraicamente fraccionario.

\subsection*{§ 3. Exclusión de la propiedad de base 3).}

Para excluir la propiedad de base 3), utilizamos un \emph{dominio de módulos} generalizado en el cual los argumentos de los polinomios del dominio ya no son las indeterminadas, sino todas las funciones algebraicas enteras de las indeterminadas,\footnote{Por ``funciones algebraicas enteras'' se entienden siempre funciones con coeficientes \emph{enteros}; es decir, funciones que satisfacen alguna ecuación cuyo coeficiente principal es la unidad y cuyos coeficientes restantes son polinomios en los $\eta$ con coeficientes enteros algebraicos: polinomios de $[H]$. En particular, todos los polinomios de $[H]$ pertenecen a las funciones algebraicas enteras.} mientras que las propias indeterminadas asumen el papel de una base de módulo.

Sea $\mathfrak G$ el conjunto de \emph{todos los elementos}
\begin{equation}
 z=\alpha+g_1(\eta)\eta_1+g_2(\eta)\eta_2+\cdots+g_t(\eta)\eta_t,
\tag{1}
\end{equation}
\emph{donde $\alpha$ recorre todos los enteros algebraicos, $\eta_1\cdots\eta_t$ todos los elementos de base, y, para cada combinación fija $\eta_1\cdots\eta_t$, las funciones $g_1(\eta)\cdots g_t(\eta)$ recorren separadamente todas las funciones algebraicas enteras de los $\eta$.}\footnote{La propiedad de módulo pertenece a los elementos $z-\alpha$.}

Es fácil verificar que para este dominio de módulos valen las propiedades I--IV. Primero, I vale porque la suma, la diferencia y el producto de dos elementos $z$ tienen de nuevo la forma (1). El dominio $\mathfrak G$ contiene además todos los elementos de base $\eta$, y también todos los elementos $\eta\cdot g(\eta)$, donde $g(\eta)$ recorre todas las funciones algebraicas enteras de los $\eta$. Por tanto, todas las funciones algebraicas enteras de los $\eta$, y en consecuencia todas las funciones algebraicas de los $\eta$ sin excepción, pueden representarse como cocientes de dos elementos de $\mathfrak G$, lo que prueba II. En la representación (1) están incluidas en particular todas las enteras algebraicas ---para $g_1=0\cdots g_t=0$---; pero $z$ no puede ser igual a un número algebraico fraccionario. Pues si $z$ representa un número algebraico $\beta$, entonces de
\[
 \beta=\alpha+g_1(\eta)\eta_1+\cdots+g_t(\eta)\eta_t
\]
idénticamente en los $\eta$, necesariamente $\beta=\alpha$. Esto se muestra mediante la especialización
\[
 \eta_1=0\cdots \eta_t=0,
\]
bajo la cual los multiplicadores $g_1(\eta)\cdots g_t(\eta)$ permanecen finitos como funciones algebraicas enteras, puesto que pasan a funciones o números algebraicos enteros.\footnote{Esta demostración más detallada de IV se da a causa del ejemplo ampliado al final del párrafo. Aquí bastaría observar que, por construcción, $z$ es una función algebraica entera.} Así se satisfacen las condiciones I--IV para $\mathfrak G$.

Por otra parte, bajo ningún buen orden puede escogerse una base de $\mathfrak G$ de tal modo que todas las funciones algebraicas enteras de los elementos de base pertenezcan a $\mathfrak G$. Sea $z(\eta)$ una función cualquiera del dominio que haya de escogerse como elemento de base ---la cual, por tanto, debe contener efectivamente las indeterminadas $\eta$---. Mostraremos que, desde cierto valor finito de $\nu$ en adelante, dependiente de $z$, las funciones
\[
 \xi=\sqrt[\nu]{z-\alpha}
\]
ya no pueden pertenecer a $\mathfrak G$.

En efecto, supóngase que $\xi$ pertenece a $\mathfrak G$. Entonces, por (1), es de la forma
\[
 \xi=\gamma+h_1(\eta)\eta_{i_1}+h_2(\eta)\eta_{i_2}
      +\cdots+h_\tau(\eta)\eta_{i_\tau},
\]
donde de nuevo $h_1(\eta)\cdots h_\tau(\eta)$ son funciones algebraicas enteras de los $\eta$, y $\eta_{i_1}\cdots\eta_{i_\tau}$ son elementos de base arbitrarios, que en particular pueden coincidir también con $\eta_1\cdots\eta_t$. Primero hay que mostrar que el entero algebraico $\gamma$ es cero. Puesto que $h_1(\eta)\cdots h_\tau(\eta)$ son funciones algebraicas enteras, $\xi$ se vuelve igual a $\gamma$ para
$\eta_{i_1}=0\cdots\eta_{i_\tau}=0$ con valores arbitrarios de los demás $\eta$; en particular, entonces, $\xi$ es igual a $\gamma$ para
\[
 \eta_{i_1}=0\cdots \eta_{i_\tau}=0; \qquad
 \eta_1=0\cdots \eta_t=0.
\]
Pero bajo esta especialización la función $(z-\alpha)$ se anula por (1), y por tanto también se anula $\xi$; de aquí $\gamma=0$, y se tiene
\[
 \xi=h_1(\eta)\eta_{i_1}+h_2(\eta)\eta_{i_2}
      +\cdots+h_\tau(\eta)\eta_{i_\tau}.
\]
Elevando a la $\nu$-ésima potencia resulta
\[
 \xi^\nu=z-\alpha=F_\nu(\eta),
\]
donde $F_\nu(\eta)$ denota una forma homogénea, no idénticamente nula, de dimensión $\nu$-ésima en $\eta_{i_1}\cdots\eta_{i_\tau}$, cuyos coeficientes son funciones algebraicas enteras de los $\eta$. Ahora $z-\alpha$, que por (1) es una función algebraica entera de los $\eta$, satisface una ecuación irreducible respecto de $[H]$:
\[
 (z-\alpha)^\chi+B(\eta)(z-\alpha)^{\chi-1}+\cdots+C(\eta)=0,
\]
donde el último coeficiente $C(\eta)$ ---el producto de $(z-\alpha)$ por sus conjugados--- es un polinomio; sea $\lambda$ su grado. Por otra parte, de $z-\alpha=F_\nu(\eta)$ se obtiene, multiplicando $F_\nu(\eta)$ por los conjugados correspondientes,
\[
 C(\eta)=G_{\nu\chi}(\eta),
\]
donde $G_{\nu\chi}(\eta)$ denota una forma homogénea de dimensión $\nu\chi$ en $\eta_{i_1}\cdots\eta_{i_\tau}$, cuyos coeficientes son polinomios en los $\eta$. Así $G_{\nu\chi}$ tiene al menos grado $\nu\chi$ en $\eta$, o $\lambda\geq\nu\chi$; es decir, las funciones $\sqrt[\nu]{z-\alpha}$ para las cuales $\nu>\lambda/\chi$ no pertenecen a $\mathfrak G$. La propiedad de base 3) queda, por tanto, excluida para la totalidad de los dominios $\mathfrak G$.

\textbf{Observación.} Una ligera generalización del dominio que acaba de construirse conduce a un dominio $\mathfrak G$ que, bajo todo buen orden, contiene también funcionales fraccionarios. Sea $\mathfrak G$ el conjunto de todos los elementos
\[
 z=\alpha+\gamma_1(\eta)\eta_1+\gamma_2(\eta)\eta_2
      +\cdots+\gamma_t(\eta)\eta_t,
\]
donde ahora $\gamma(\eta)$ recorre todas las funciones algebraicamente enteras, pero no necesariamente enteras, de los $\eta$. La demostración de las propiedades I--IV permanece igual que antes. Pero como $\mathfrak G$, además de toda función $z$, contiene también $a\cdot(z-\alpha)$, donde $a$ denota cualquier número racional o algebraicamente fraccionario, bajo todo buen orden aparecen también funcionales fraccionarios.

\subsection*{§ 4. Los dominios más generales de números trascendentes algebraicamente enteros.}

Para formular con más comodidad lo que sigue, introduzcamos la expresión: ``un elemento $z$ depende algebraicamente entero de $\mathfrak T$'', lo cual significa que $z$ satisface alguna ecuación cuyo coeficiente principal es la unidad y cuyos coeficientes restantes son polinomios enteros en los elementos de $\mathfrak T$, donde $\mathfrak T$ denota cualquier dominio prescrito.\footnote{Para un dominio íntegro $\mathfrak T$, el coeficiente principal es la unidad y los coeficientes restantes son elementos de $\mathfrak T$.}

Un dominio $\mathfrak T$ se llama \emph{algebraicamente enteramente cerrado} si satisface la condición (compárese la introducción):
\begin{enumerate}
\item[V.] Todo elemento que depende algebraicamente entero de $\mathfrak T$ pertenece a $\mathfrak T$.
\end{enumerate}

Mostramos primero que la propiedad abstracta V pertenece al dominio de Zermelo $\mathfrak G_\eta$. Sea $z$ un elemento que depende algebraicamente entero de $\mathfrak G_\eta$ por medio de una ecuación $f(z)=0$. Entonces la multiplicación de $f(z)$ por sus conjugados respecto de $[H]$ muestra que $z$ depende también algebraicamente entero de $[H]$, y por tanto pertenece a $\mathfrak G_\eta$. La cerradura algebraicamente entera del dominio funcional considerado en § 2 se muestra de la misma manera.

Que la propiedad V no pertenece a todo dominio $\mathfrak G$ lo muestra el dominio de módulos considerado en § 3, que de hecho no posee la propiedad de base 3). Más precisamente, § 3 mostró que el dominio de módulos no está algebraicamente enteramente cerrado respecto de ningún elemento trascendental del dominio, mientras que, por III y IV, esto naturalmente sí ocurre respecto de todo elemento algebraico del dominio.

Esto nos lleva primero a seleccionar, de la totalidad de los dominios $\mathfrak G$, aquellos que poseen también la propiedad abstracta V; éstos se llamarán ``dominios $\mathfrak H$ de números trascendentes algebraicamente enteros''.

Sea $\mathfrak H$ tal dominio, y sea $H$ una base algebraica de todos los números de $\mathfrak H$; por II, $H$ es al mismo tiempo una base algebraica de todos los números. Por III, I y V, $\mathfrak H$ contiene, además de $H$, todas las funciones algebraicas enteras de los $\eta$, es decir, el dominio de Zermelo $\mathfrak G_\eta$. Así, en analogía con § 1: todo dominio $\mathfrak H$ de números trascendentes algebraicamente enteros puede construirse mediante una base algebraica de todos los números, de tal modo que $\mathfrak H$ contenga el dominio de Zermelo $\mathfrak G_\eta$ como divisor.

Sea ahora $H$ una base algebraica cualquiera de todos los números, y denote $\mathfrak M(H)$ la totalidad de todos los dominios $\mathfrak H$ que contienen a $H$. Todo dominio $\mathfrak H$ en $\mathfrak M(H)$ contiene, como acaba de probarse, al dominio de Zermelo $\mathfrak G_\eta$ como divisor; y como $\mathfrak G_\eta$ mismo es un dominio $\mathfrak H$, $\mathfrak G_\eta$ es el máximo común divisor, o intersección, de todos los dominios $\mathfrak H$ de $\mathfrak M(H)$. Así $\mathfrak G_\eta$ ha sido definido abstractamente. También se sigue directamente que $\mathfrak M(H)$ es cerrado bajo intersección: es decir, la intersección de todos los dominios $\mathfrak H$ de cualquier subsistema de $\mathfrak M(H)$ pertenece de nuevo a $\mathfrak M(H)$. En efecto, como intersección se satisfacen I y V; II y III se siguen de que $\mathfrak G_\eta$ es un divisor; IV se sigue de que la intersección es divisor de $\mathfrak H$.

Como muestra el ejemplo del dominio funcional en § 2, los dominios $\mathfrak H$ contienen en general más funciones que $\mathfrak G_\eta$. Sea $\psi(\eta)$ una función racional o algebraica de los $\eta$ de tal clase. Entonces, por I, $\mathfrak H$ contiene también todas las combinaciones racionales enteras de $\psi$ y de un número finito de funciones de $\mathfrak G_\eta$ cada vez. El dominio íntegro que surge así ---formado por todos los polinomios en $\psi$ con coeficientes de $\mathfrak G_\eta$--- se denotará por $[\mathfrak G_\eta,\psi]$, y el proceso que conduce a él se llamará ``adjunción racionalmente entera de $\psi$ a $\mathfrak G_\eta$''. Por V, $\mathfrak H$ contiene además todas las funciones que dependen algebraicamente enteras de $[\mathfrak G_\eta,\psi]$; el dominio así obtenido se denotará por $\{\mathfrak G_\eta,\psi\}$, y el proceso que conduce a él por ``adjunción algebraicamente entera de $\psi$ a $\mathfrak G_\eta$''. El dominio $\{\mathfrak G_\eta,\psi\}$ consta de todas las funciones algebraicas enteras de $\psi$ con coeficientes de $\mathfrak G_\eta$, y por tanto es un dominio íntegro por argumentos conocidos.\footnote{Compárese, por ejemplo, Weber, § 97, 6 y 7.}

Mostramos que $\{\mathfrak G_\eta,\psi\}$ es también algebraicamente enteramente cerrado, es decir, que satisface la condición V. Sea, en efecto, $z$ un elemento que depende algebraicamente entero de $\{\mathfrak G_\eta,\psi\}$ por medio de una ecuación
\[
 f(z;a\cdots c)=z^\nu+a z^{\nu-1}+\cdots+c=0,
\]
donde $a\cdots c$, como elementos de $\{\mathfrak G_\eta,\psi\}$, satisfacen ecuaciones con coeficientes de $[\mathfrak G_\eta,\psi]$:
\[
\begin{gathered}
 a^\lambda+A_1a^{\lambda-1}+\cdots+A_\lambda=0,\\
 \vdots\\
 c^\mu+C_1c^{\mu-1}+\cdots+C_\mu=0.
\end{gathered}
\]
Sean sus raíces
$a,a'\cdots a^{(\lambda-1)}\cdots c,c'\cdots c^{(\mu-1)}$. Si se forma ahora el producto sobre todas las combinaciones de las raíces,
\[
 g(z)=f(z;a\cdots c)f(z;a'\cdots c')\cdots
      f(z;a^{(\lambda-1)}\cdots c^{(\mu-1)}),
\]
entonces $g(z)$ tiene la unidad como coeficiente principal y elementos de $[\mathfrak G_\eta,\psi]$ como coeficientes restantes. En consecuencia, $z$ pertenece a $\{\mathfrak G_\eta,\psi\}$.\footnote{Así, al definir la integridad algebraica mediante \emph{alguna} ecuación, el concepto de irreducibilidad respecto de un cuerpo base se vuelve innecesario para la prueba de I y V. Compárese también § 2, donde, sólo para la prueba de IV, se necesita la proposición auxiliar que pasa de la definición de integridad algebraica mediante una ecuación arbitraria a la ecuación irreducible. Puesto que esta proposición auxiliar ya no vale en general en el caso presente, IV debe imponerse a la función $\psi$ como condición especial.}

Así, el dominio $\{\mathfrak G_\eta,\psi\}$ tiene las propiedades I y V; y, puesto que $\mathfrak G_\eta$ es divisor del dominio, también tiene II y III. Que se satisfaga IV depende de la elección de $\psi$; por ejemplo, IV no se satisface para $\psi=\frac{1}{n\cdot\eta}$, pues entonces $\frac1n$ pertenece al dominio. En cambio se satisface, por § 2, para $\psi=\frac1\eta$, o también para $\psi=\frac{\eta}{n}$. Pues en este último caso, si una función de $\{\mathfrak G_\eta,\psi\}$ representa un número algebraico, su valor permanece inalterado para $\eta=0$; con ello pasa a una función de $\mathfrak G_\eta$, y por consiguiente a un entero algebraico.

Por tanto, $\{\mathfrak G_\eta,\psi\}$ se convierte en un dominio $\mathfrak H$ tan pronto como se impone a $\psi$ la condición de que por la adjunción algebraicamente entera no entre en el dominio ningún número algebraicamente fraccionario.

De modo completamente análogo se define la adjunción racionalmente entera, respectivamente algebraicamente entera, de un sistema arbitrario $\mathfrak S$ de funciones racionales o algebraicas de los $\eta$ a $\mathfrak G_\eta$. La adjunción racionalmente entera conduce al dominio íntegro $[\mathfrak G_\eta,\mathfrak S]$, formado por todos los polinomios enteros en un número finito de elementos de $\mathfrak G_\eta$ y $\mathfrak S$ cada vez. La adjunción algebraicamente entera da el dominio $\{\mathfrak G_\eta,\mathfrak S\}$, formado por todos los elementos que dependen algebraicamente enteros de $[\mathfrak G_\eta,\mathfrak S]$. Exactamente como antes se muestra que el dominio íntegro $\{\mathfrak G_\eta,\mathfrak S\}$ es algebraicamente enteramente cerrado y también satisface II y III. Por tanto, como antes:

$\{\mathfrak G_\eta,\mathfrak S\}$ se convierte en un dominio $\mathfrak H$ tan pronto como se impone a $\mathfrak S$ la condición de que por la adjunción algebraicamente entera no entre en el dominio ningún número algebraicamente fraccionario; es decir, tan pronto como $\mathfrak S$ se convierte en un sistema admisible.\footnote{Más generalmente, se muestra análogamente: si $\mathfrak T$ es cualquier dominio y $\{\mathfrak T\}$ denota la totalidad de elementos que dependen algebraicamente enteros de $\mathfrak T$, entonces $\{\mathfrak T\}$ es algebraicamente enteramente cerrado.}

Ahora es importante que también valga el recíproco, de modo que los dominios $\{\mathfrak G_\eta,\mathfrak S\}$ agotan efectivamente todos los dominios $\mathfrak H$ de $\mathfrak M(H)$ cuando $\mathfrak S$ recorre todos los sistemas admisibles. En efecto, sea $\mathfrak H$ cualquier dominio de $\mathfrak M(H)$ y denote $\mathfrak L=\mathfrak H-\mathfrak G_\eta$ el sistema formado por todas las funciones de $\mathfrak H$ que no pertenecen a $\mathfrak G_\eta$. Por V, $\mathfrak H$ contiene el dominio $\{\mathfrak G_\eta,\mathfrak L\}$ como divisor; pero $\{\mathfrak G_\eta,\mathfrak L\}$ también es divisible por $\mathfrak H$, porque contiene $\mathfrak G_\eta$ y $\mathfrak L$, y como éstos no tienen elemento común, contiene también a $\mathfrak H$. Por tanto $\mathfrak H=\{\mathfrak G_\eta,\mathfrak L\}$; y puesto que IV vale para $\mathfrak H$, $\mathfrak L$ es naturalmente un sistema admisible.\footnote{La adjunción de un subsistema de $\mathfrak L$ puede ya ser suficiente; por ejemplo, el dominio funcional de § 2 surge por adjunción algebraicamente entera de todos los funcionales racionales enteros ---con excepción de los polinomios--- a $\mathfrak G_\eta$.}

Cuando $H$ recorre todas las bases algebraicas posibles, $\mathfrak M(H)$ recorre, como se probó al comienzo, la totalidad de todos los dominios $\mathfrak H$. Por consiguiente, la cuestión de los dominios más generales $\mathfrak H$ de números trascendentes algebraicamente enteros queda completamente respondida por los resultados de este párrafo:
\begin{enumerate}
\item[a)] La intersección de todos los dominios $\mathfrak H$ de $\mathfrak M(H)$ está dada por el dominio de Zermelo $\mathfrak G_\eta$, que él mismo es un dominio $\mathfrak H$; $\mathfrak M(H)$ es cerrado respecto de la formación de intersecciones.
\item[b)] Todo dominio $\mathfrak H$ de $\mathfrak M(H)$ surge por adjunción algebraicamente entera de un sistema admisible $\mathfrak S$ a $\mathfrak G_\eta$. Cuando $H$ recorre todas las bases algebraicas posibles, $\mathfrak M(H)$ recorre la totalidad de todos los dominios $\mathfrak H$.\footnote{Los sistemas admisibles $\mathfrak S$ pueden encontrarse así. Pónganse todas las funciones algebraicamente fraccionarias de los $\eta$ bajo algún buen orden $\Omega$, y admítanse en $\mathfrak S$ todas las funciones salvo aquellas que, al adjuntarse algebraicamente enteras a $\mathfrak G_\eta$ y a las funciones admitidas anteriormente, generen un número algebraicamente fraccionario. Este procedimiento puede interrumpirse en cualquier lugar arbitrario. Cuando $\Omega$ recorre todos los buenos órdenes, y cuando el proceso se interrumpe en todo lugar arbitrario para cada $\Omega$ fijo, se agotan así todos los sistemas admisibles $\mathfrak S$.}
\end{enumerate}

\section*{§ 5. Los dominios más generales formados con números trascendentes enteros, sobre la base de una base algebraica.}

Sea $H$ una base algebraica cualquiera de todos los números, y denote $\mathfrak M(\mathfrak G)$ la totalidad de todos los dominios $\mathfrak G$ que contienen $H$ y, por tanto, $[H]$. Cuando $H$ recorre todas las bases posibles, $\mathfrak M(\mathfrak G)$ recorre la totalidad de todos los dominios $\mathfrak G$; por ello, como en § 4, podemos restringirnos a considerar los dominios $\mathfrak G$ de $\mathfrak M(\mathfrak G)$.

Los dominios $\mathfrak G$ de $\mathfrak M(\mathfrak G)$ contienen todos $[H]$ como divisor; pero, como $[H]$ mismo no es un dominio $\mathfrak G$, no se puede concluir como en § 4 que $[H]$ sea el máximo común divisor. Que éste es, no obstante, el caso se muestra mediante un conjunto numerable de dominios $\mathfrak G$ de $\mathfrak M(\mathfrak G)$, obtenidos por una ligera generalización del dominio de módulos de § 3, y que de hecho no tienen ningún elemento común fuera de $[H]$.

\emph{Para cada $\nu$ fijo, sean los dominios $\mathfrak G_\nu$ $(\nu=1,2,\ldots\text{ in inf.})$ la totalidad de las cantidades}
\begin{equation}
z=f(\eta)+g_1(\eta)\eta_1^\nu+g_2(\eta)\eta_2^\nu+\cdots+g_t(\eta)\eta_t^\nu,
\tag{1}
\end{equation}
\emph{donde $f(\eta)$ recorre todos los polinomios enteros en los $\eta$; $\eta_1,\ldots,\eta_t$ recorren las cantidades de base, y, para cada combinación fija $\eta_1^\nu,\ldots,\eta_t^\nu$, las funciones $g_1(\eta),\ldots,g_t(\eta)$ recorren independientemente todas las funciones algebraicas enteras de los $\eta$.}

Como en § 3, se ve que los dominios $\mathfrak G_\nu$ tienen las propiedades abstractas I--IV; $\mathfrak G_\nu$ sólo contiene funciones algebraicas enteras y, para $\nu=1$, es idéntico al dominio de § 3. Mostramos ahora que los dominios $\mathfrak G_\nu$ no tienen ningún otro elemento común fuera de los polinomios $f(\eta)$ de $[H]$. Sea, en efecto, $z$ una función algebraica entera ---pero no racional--- de los $\eta$, que satisface la ecuación irreducible respecto de $[H]$,
\begin{equation}
z^\chi+a_1(\eta)z^{\chi-1}+a_2(\eta)z^{\chi-2}+\cdots+a_\chi(\eta)=0,
\tag{2}
\end{equation}
donde $\chi>1$; y sea $\lambda$ el mayor de los grados de los polinomios $a_1(\eta),\ldots,a_\chi(\eta)$. Si $z$ pertenece a $\mathfrak G_\nu$, entonces la cantidad $\xi=z-f(\eta)$ satisface la ecuación igualmente irreducible, de nuevo respecto de $[H]$,
\begin{equation}
\xi^\chi+b_1(\eta)\xi^{\chi-1}+b_2(\eta)\xi^{\chi-2}+\cdots+b_\chi(\eta)=0,
\tag{3}
\end{equation}
donde, por (1), $b_1(\eta),b_2(\eta),\ldots,b_\chi(\eta)$, como funciones simétricas de los $\xi$, contienen como términos de menor dimensión términos de al menos las dimensiones $\nu,2\nu,\ldots,\chi\nu$. Comparar (2) y (3) da
\[
 b_1(\eta)=a_1(\eta)+\chi f(\eta).
\]
Si, por tanto, la cantidad $\xi=z+\frac{a_1(\eta)}{\chi}$ satisface
\begin{equation}
\xi^\chi+c_2(\eta)\xi^{\chi-2}+\cdots+c_\chi(\eta)=0,
\tag{4}
\end{equation}
entonces
\begin{equation}
 \xi-\zeta=\frac{b_1(\eta)}{\chi};
 \qquad
 c_i(\eta)\equiv b_i(\eta)\pmod{\frac{b_1(\eta)}{\chi}},
\tag{5}
\end{equation}
donde $b_1(\eta)$ comienza con términos de al menos dimensión $\nu$. Ahora los $c_i(\eta)$, como muestra la comparación de (2) con (4), son de grado $\chi$ en los coeficientes $a_i(\eta)$ de (2), por tanto de grado a lo sumo $\chi\lambda$ en los $\eta$; por otra parte, todos los $b_i(\eta)$ comienzan con términos de al menos dimensión $\nu$. Así, si se escoge $\nu>\chi\lambda$, (5) da
\[
 c_i(\eta)=0;\qquad
 \xi^\chi=0\quad\text{o}\quad
 \left(z+\frac{a_1(\eta)}{\chi}\right)^\chi=0.
\]
Así $z$, contra la hipótesis, se vuelve racional en los $\eta$. De aquí:

\emph{Si $z$ es una función algebraica entera no racional de los $\eta$, entonces $z$ no puede pertenecer a ningún dominio $\mathfrak G_\nu$ para el cual $\nu>\chi\lambda$, o bien:}

\emph{La intersección de todos los dominios $\mathfrak G$ de $\mathfrak M(\mathfrak G)$ está dada por el dominio íntegro $[H]$, que él mismo no es un dominio $\mathfrak G$; por consiguiente, $\mathfrak M(\mathfrak G)$ no es cerrado respecto de la formación de intersecciones.}

Por consiguiente, la construcción de los dominios más generales $\mathfrak G$ mediante una base algebraica es también más difícil de abarcar que la de los dominios $\mathfrak H$. El dominio $[H]$ tiene las propiedades I, III y IV; si se adjunta racionalmente de manera entera un sistema arbitrario $\mathfrak S$, entonces para el dominio íntegro resultante $[H,\mathfrak S]$ siguen valiendo I y III, mientras que IV forma una condición especial sobre $\mathfrak S$. Para que surja un dominio $\mathfrak G$, debe imponerse también II como condición sobre $\mathfrak S$; o, expresado de otro modo: para todo cuerpo $\mathfrak K$, finito sobre $(H)$\footnote{$(H)$ significa el cuerpo formado por todas las funciones racionales de los $\eta$ con coeficientes numéricos algebraicos.}, debe existir un cuerpo $\mathfrak L$ sobre $\mathfrak K$, finito sobre $\mathfrak K$, tal que todas las funciones de $\mathfrak L$ puedan representarse como cocientes de funciones de $[H,\mathfrak S]$. Esta condición resulta mucho menos manejable que la única condición sobre sistemas admisibles de § 4; por ello pasamos ahora a una base racional.

\section*{§ 6. Los dominios más generales formados con números trascendentes enteros, sobre la base de una base racional.}

Sea $\mathfrak G$ un dominio cualquiera de números trascendentes enteros. Ordenemos bien los elementos de $\mathfrak G$ y elijamos, por inducción transfinita, de los elementos de $\mathfrak G$ aquellos que no puedan expresarse racionalmente por un número finito de los precedentes. El sistema así obtenido se denotará por $\Theta$. Todo elemento de $\mathfrak G$ es racionalmente expresable por $\Theta$; pues si $z_0$ fuese el primer elemento que no lo fuera, entonces $z_0$ tendría que poder expresarse racionalmente por un número finito de elementos precedentes de $\mathfrak G$; sustituyendo éstos, por hipótesis de inducción, por funciones racionales de los $\vartheta$, se obtendría una representación de $z_0$ por los $\vartheta$, contradicción. Además, ningún elemento de $\Theta$ puede expresarse racionalmente por un número finito de los precedentes, por construcción. Así $\Theta$ es una base racional de $\mathfrak G$.

Puesto que por II todo número es cociente de dos elementos de $\mathfrak G$, todo número puede expresarse racionalmente por $\Theta$; es decir, $\Theta$ es al mismo tiempo una base racional de todos los números. Por III e I, $\mathfrak G$ contiene, además de $\Theta$, el dominio íntegro $[\Theta]$ de todos los polinomios enteros en los $\vartheta$ como divisor, mientras que por IV $[\Theta]$ no contiene ningún número algebraicamente fraccionario. Así, la base racional $\Theta$ de todos los números satisface la condición lateral de que el dominio íntegro $[\Theta]$ derivado de $\Theta$ no contiene ningún número algebraicamente fraccionario;\footnote{Que esto es efectivamente una condición se ve, por ejemplo, por el hecho de que las dos cantidades $\eta,\frac1{2\sqrt{\eta}}$ no son admisibles como elementos de base, mientras que $\eta,\frac1{\sqrt{\eta}}$ sí lo son.} $\Theta$ es una \emph{base racional con condición lateral} para todos los números. Así, en analogía con § 2: \emph{Todo dominio $\mathfrak G$ puede construirse mediante una base racional con condición lateral, de tal modo que $\mathfrak G$ contenga el dominio íntegro $[\Theta]$ como divisor.}

Sea $\Theta$ una base racional cualquiera con condición lateral para todos los números ---cuya existencia queda por tanto garantizada por la existencia de los dominios $\mathfrak G$---. Denote $\mathfrak N(\mathfrak G)$ la totalidad de todos los dominios $\mathfrak G$ que contienen $\Theta$ y, por tanto, $[\Theta]$. Cuando $\Theta$ recorre todas las bases posibles con condición lateral, $\mathfrak N(\mathfrak G)$ recorre, por lo anterior, la totalidad de todos los dominios $\mathfrak G$; de ahí que podamos restringirnos otra vez a los dominios $\mathfrak G$ de $\mathfrak N(\mathfrak G)$.

Ahora $[\Theta]$ mismo es un dominio $\mathfrak G$; pues todo número puede expresarse racionalmente por $\Theta$, y por tanto como cociente de dos polinomios de $[\Theta]$; y por su construcción $[\Theta]$ satisface también las restantes condiciones I, III, IV. \emph{Así, $[\Theta]$ es el máximo común divisor, o intersección, de todos los dominios $\mathfrak G$ de $\mathfrak N(\mathfrak G)$; además, $\mathfrak N(\mathfrak G)$ es cerrado respecto de la intersección; es decir, la intersección de todos los dominios $\mathfrak G$ de cualquier subsistema de $\mathfrak N(\mathfrak G)$ pertenece también a $\mathfrak N(\mathfrak G)$.} Pues I vale para la intersección; II y III se siguen de que $[\Theta]$ es un divisor; y IV, de que el dominio de intersección está contenido como divisor en un dominio $\mathfrak G$.

La adjunción racionalmente entera de cualquier sistema $\mathfrak S$ de funciones racionales de los $\vartheta$ a $[\Theta]$ ---toda función algebraica de los $\vartheta$ puede, por la propiedad de la base racional, expresarse racionalmente mediante ciertos otros $\vartheta$--- da un dominio íntegro $[\Theta,\mathfrak S]$ que tiene las propiedades I, II, III, y por tanto se convierte en un dominio $\mathfrak G$ tan pronto como $\mathfrak S$ es un ``sistema admisible'', es decir, tan pronto como satisface la condición de que por la adjunción racionalmente entera no entre en el dominio ningún número algebraicamente fraccionario. Recíprocamente, todo dominio $\mathfrak G$ de $\mathfrak N(\mathfrak G)$ admite la representación $[\Theta,\mathfrak S]$, con tal de que $\mathfrak S$ denote el sistema residual $\mathfrak G-[\Theta]$. Así, para los dominios más generales $\mathfrak G$, en completa analogía con § 4, se tiene:
\begin{enumerate}
\item[a)] \emph{La intersección de todos los dominios $\mathfrak G$ de $\mathfrak N(\mathfrak G)$ está dada por el dominio íntegro $[\Theta]$, que él mismo es un dominio $\mathfrak G$; $\mathfrak N(\mathfrak G)$ es cerrado respecto de la formación de intersecciones.}
\item[b)] \emph{Todo dominio $\mathfrak G$ de $\mathfrak N(\mathfrak G)$ surge por adjunción racionalmente entera de un sistema admisible $\mathfrak S$ a $[\Theta]$. Cuando $\Theta$ recorre todas las bases racionales posibles con condición lateral, $\mathfrak N(\mathfrak G)$ recorre la totalidad de todos los dominios $\mathfrak G$.}\footnote{Para los sistemas admisibles más generales $\mathfrak S$, vale lo dicho en § 4, sustituyendo ``algebraico'' por ``racional''.}
\end{enumerate}

\emph{Observación.} Si de $\Theta$ se extrae una base algebraica $H$ de todas las cantidades de $\Theta$, entonces $H$ es al mismo tiempo una base algebraica de todos los números; recíprocamente, toda base algebraica puede extenderse a una racional poniendo $H$ primero en el buen orden. De acuerdo con esto, la construcción de § 5 de los dominios $\mathfrak G$ mediante una base algebraica puede interpretarse así. Descompóngase el sistema $\mathfrak S$ que se va a adjuntar a $[H]$ en dos subsistemas $\mathfrak S_1$ y $\mathfrak S_2$; la adjunción de $\mathfrak S_1$ equivale a extender $H$ a una base racional con condición lateral, a la que se adjunta todavía racionalmente y de manera entera un sistema admisible $\mathfrak S_2$.

\section*{§ 7. Construcción de la base racional más general con condición lateral.}

Los resultados del párrafo precedente necesitan todavía un suplemento, pues allí se demostró la \emph{existencia} de una base racional con condición lateral para todos los números, pero no un método para construirla a partir de un buen orden prescrito. A causa de la condición lateral, esta construcción toma una forma más complicada cuando, en lugar de un dominio $\mathfrak G$, se tiene el cuerpo de todos los números complejos.\footnote{En § 6 b) bastaría hacer recorrer a $\Theta$ sólo aquellas bases racionales especiales que constan de una base algebraica y de funciones algebraicamente enteras de las cantidades de base; entonces la condición lateral se satisface automáticamente. Para verlo, basta, por § 6, extender una base algebraica de $\mathfrak G$ a una racional y, en la base resultante, multiplicar cada función algebraicamente fraccionaria de los $\eta$ por un polinomio adecuadamente elegido en los $\eta$, con lo cual se conserva la propiedad de base. Esta construcción se traslada sin cambio al cuerpo de todos los números complejos. Sin embargo, la cuestión que surge de § 6 acerca de la base racional más general con condición lateral tiene también interés independiente.}

Sométase el cuerpo de todos los números complejos a un buen orden $\Omega$. Entonces pueden aparecer cantidades de tres clases:
\begin{enumerate}
\item[1.] Cantidades $y$ cuya adjunción racionalmente entera a las cantidades de base precedentes ---definidas recursivamente en 3.--- fuerza la entrada de un número algebraicamente fraccionario en el dominio íntegro resultante; si no precede ninguna cantidad de base, este dominio consta de todos los polinomios enteros en $y$. En particular, todos los números algebraicamente fraccionarios son cantidades $y$, puesto que el dominio íntegro resultante contiene las cantidades $y=\beta$.
\item[2.] Cantidades $z$ que no tienen la propiedad 1, pero que pueden expresarse racionalmente por un número finito de números precedentes distintos de los $y$; así satisfacen una relación $z=\varphi(\xi_1,\ldots,\xi_t)$, donde $\varphi$ es una función racional con coeficientes en $K$ y ninguna de las $\xi$ es un $y$. En particular, todos los números enteros algebraicos $\alpha$ pertenecen aquí, pues la propiedad 1 no vale para ellos y satisfacen la relación $z=\alpha$.
\item[3.] Cantidades $\vartheta$ para las cuales no vale ni 1 ni 2, y que deben llamarse cantidades de base. En particular, el primer número trascendental $\vartheta_0$ en el buen orden es una tal cantidad de base; pues los polinomios enteros en $\vartheta_0$ no pueden representar ningún número algebraicamente fraccionario, y $\vartheta_0$ tampoco puede depender racionalmente de los números algebraicos precedentes.
\end{enumerate}

\emph{Mostraremos ahora que el sistema $\Theta$ de todas las cantidades de base $\vartheta$ es una base racional con condición lateral para todos los números.}

Como ninguna cantidad $y$ aparece entre los $\vartheta$, el dominio íntegro $[\Theta]$ no puede contener ningún número algebraicamente fraccionario; de ahí que se satisfaga la condición lateral. Como, además, ninguna cantidad $z$ aparece entre los $\vartheta$, cada $\vartheta$ es racionalmente independiente de las cantidades de base precedentes. El argumento de inducción muestra, exactamente como en § 6, que toda relación $z=\varphi(\xi)$ ---donde no aparece ningún $y$ entre los $\xi$--- implica una relación $z=\psi(\vartheta)$, de modo que todas las cantidades $z$ son racionalmente expresables por los $\vartheta$.

Sólo queda probar la afirmación más difícil: que los $y$ también son racionalmente expresables por los $\vartheta$. Primero mostramos que los $y$ dependen algebraicamente de los $\vartheta$. Por hipótesis, para cada $y$ existe al menos un número algebraicamente fraccionario $\beta$ que admite la representación
\begin{equation}
\beta=f_0(\vartheta)+f_1(\vartheta)y+\cdots+f_\chi(\vartheta)y^\chi,
\tag{1}
\end{equation}
donde $f_0(\vartheta),\ldots,f_\chi(\vartheta)$ son polinomios de $[\Theta]$ y por tanto no pueden representar números algebraicamente fraccionarios. Si $y$ fuera algebraicamente independiente de los $\vartheta$, entonces (1) valdría idénticamente en $y$ y se tendría $\beta=f_0(\vartheta)$, en contradicción con las propiedades de $[\Theta]$.

Para inferir la dependencia racional a partir de la dependencia algebraica, extraemos de $\Theta$ una base algebraica $H$ de todas las cantidades de $\Theta$. Entonces todo $y$, como función algebraica de los $\vartheta$, es también algebraicamente dependiente de los $\eta$; mediante $y=y(\eta)$ queda por tanto determinado un cuerpo algebraico $(H,y(\eta))$ sobre $(H)$. Como todos los $\eta$ son al mismo tiempo cantidades $\vartheta$, la prueba de la dependencia racional equivale a mostrar que para todo cuerpo de esta clase existen elementos primitivos que son funciones racionales de los $\vartheta$; más precisamente, se mostrará que $y(\eta)$ pierde la propiedad 1 después de multiplicarlo por un polinomio en los $\eta$, y en consecuencia pasa a ser un tal elemento.

Se sabe que entre $(H)$ y $(H,y)$ sólo hay un número finito de cuerpos intermedios. Sean $\Omega_0=(H),\Omega_1,\ldots,\Omega_\sigma$ aquellos entre ellos que poseen un elemento primitivo racionalmente expresable por los $\vartheta$, de modo que todo elemento del cuerpo se convierte en una función racional de los $\vartheta$; esta condición se satisface siempre al menos para $\Omega_0=(H)$. Sea la función primitiva correspondiente a $\Omega_i$ dada por $\frac{g_i(\vartheta)}{h_i(\vartheta)}$, donde $g_i,h_i$ son polinomios de $[\Theta]$. Entonces ---si $\alpha_i$ denota un entero elegido convenientemente--- el polinomio de $[\Theta]$
\[
 \xi_i=g_i(\vartheta)+\alpha_i h_i(\vartheta)
\]
es una función primitiva de $\Omega_i$, o de un cuerpo algebraico sobre $\Omega_i$,\footnote{Cf. el final de § 5.} y toda función $\psi(\eta)$ de $\Omega_i$ admite la representación
\begin{equation}
\psi(\eta)=
\frac{a_1(\eta)\xi_i^{\kappa-1}+a_2(\eta)\xi_i^{\kappa-2}+\cdots+a_\kappa(\eta)}
     {a(\eta)}
=\frac{A(\eta,\xi_i)}{a(\eta)},
\tag{2}
\end{equation}
donde $a(\eta),a_1(\eta),\ldots,a_\kappa(\eta)$ son polinomios enteros en los $\eta$, y por tanto $A(\eta,\xi_i)$ y $a(\eta)$ son cantidades de $[\Theta]$.

La ecuación irreducible que $y$ satisface respecto de $\Omega_i$ es, por (2), de la forma
\[
 f^{(i)}(\eta)y^{\lambda_i}+f_1^{(i)}(\eta,\xi_i)y^{\lambda_i-1}
+\cdots+f_{\lambda_i}^{(i)}(\eta,\xi_i)=0,
\]
y todos los coeficientes de esta ecuación pertenecen a $[\Theta]$. Entonces la función $y_i=f^{(i)}(\eta)\cdot y$ satisface la ecuación igualmente irreducible, respecto de $\Omega_i$,
\[
 y_i^{\lambda_i}+f_1^{(i)}(\eta,\xi_i)y_i^{\lambda_i-1}
+\cdots+\bigl(f^{(i)}(\eta)\bigr)^{\lambda_i-1}
 f_{\lambda_i}^{(i)}(\eta,\xi_i)=0,
\]
en virtud de la cual depende algebraicamente entera de $[H,\xi_i]$; lo mismo vale para el producto de $y_i$ por cualquier polinomio en los $\eta$. En particular, la función
\begin{equation}
Y=f^{(0)}(\eta)f^{(1)}(\eta)\cdots f^{(\sigma)}(\eta)\cdot y
\tag{3}
\end{equation}
depende algebraicamente entera de todos los dominios $[H,\xi_i]$ por medio de la ecuación, irreducible en cada caso respecto de $\Omega_i$,
\begin{equation}
Y^{\lambda_i}+F_1^{(i)}(\eta,\xi_i)Y^{\lambda_i-1}
+\cdots+F_{\lambda_i}^{(i)}(\eta,\xi_i)=0
\qquad (i=0,1,\ldots,\sigma).
\tag{4}
\end{equation}
Mostramos ahora que $Y$ es el elemento primitivo buscado, racionalmente expresable por los $\vartheta$. Supóngase, por ejemplo, que el número algebraico $\gamma$ está representado por $Y$ y $[\Theta]$:
\begin{equation}
\gamma=k_0(\vartheta)+k_1(\vartheta)Y+\cdots+k_\nu(\vartheta)Y^\nu,
\tag{5}
\end{equation}
donde $k_0(\vartheta),\ldots,k_\nu(\vartheta)$ son polinomios de $[\Theta]$. Determinamos ahora la ecuación irreducible que satisface $Y$ respecto del cuerpo
$\mathfrak K=(H;k_0(\vartheta),\ldots,k_\nu(\vartheta))$ que surge del dominio de coeficientes de (5). Es sabido que está dada por la ecuación irreducible de $Y$ respecto del cuerpo de intersección de $\mathfrak K$ y $(H,Y)$, situado entre $(H)$ y $(H,Y)$. Ahora $(H,Y)$ es idéntico a $(H,y)$ por (3); y como todos los elementos de $\mathfrak K$, y por tanto también todos los elementos del cuerpo de intersección, pueden expresarse racionalmente por un número finito de $\vartheta$, este último cuerpo es necesariamente uno de los cuerpos $\Omega_0,\Omega_1,\ldots,\Omega_\sigma$, digamos $\Omega_\tau$. La ecuación irreducible buscada es por tanto, por (4), de grado $\lambda_\tau$ y es
\begin{equation}
Y^{\lambda_\tau}+F_1^{(\tau)}(\eta,\xi_\tau)Y^{\lambda_\tau-1}
+\cdots+F_{\lambda_\tau}^{(\tau)}(\eta,\xi_\tau)=0
\tag{6}
\end{equation}
con polinomios de $[\Theta]$ como coeficientes. Por medio de (6), (5) puede reescribirse en la forma
\begin{equation}
\gamma=K_0(\vartheta)+K_1(\vartheta)Y+\cdots+
K_{\lambda_\tau-1}(\vartheta)Y^{\lambda_\tau-1},
\tag{7}
\end{equation}
donde $K_0(\vartheta),K_1(\vartheta),\ldots$ pertenecen a $\mathfrak K$ y, por otra parte, son polinomios de $[\Theta]$. Pero como $\gamma$ es un número algebraico, la relación (7) no alcanza el grado $\lambda_\tau$ en $Y$ y por consiguiente se satisface \emph{idénticamente} en $Y$. De aquí
\[
 \gamma=K_0(\vartheta),
\]
es decir, $\gamma$ pertenece a $[\Theta]$ y por consiguiente es un entero algebraico.

Cuando $\gamma$ recorre todos los números algebraicos representables por $Y$ y $[\Theta]$, el cuerpo de intersección de $(H,Y)$ y del cuerpo correspondiente $\mathfrak K$ recorre siempre sólo los cuerpos $\Omega_0,\Omega_1,\ldots,\Omega_\sigma$, y todos los argumentos precedentes siguen siendo válidos. Es decir, \emph{por adjunción racionalmente entera de $Y$ a $[\Theta]$ no puede representarse ningún número algebraicamente fraccionario}. Por tanto $Y$ ya no tiene la propiedad 1; según 2) o 3), es una cantidad $z$ o $\vartheta$, y por tanto racionalmente representable por un número finito de $\vartheta$; por (3), lo mismo vale para $y$: \emph{se ha demostrado que $\Theta$ es una base racional de todos los números}. Recíprocamente, como toda base racional con condición lateral puede construirse según 1), 2), 3) poniendo $\Theta$ primero en el buen orden, hemos probado: \emph{La construcción dada por 1), 2), 3) conduce a la base racional más general con condición lateral para todos los números.}

\section*{§ 8. División de los dominios $\mathfrak G$ y $\mathfrak H$ en clases.}

Junto a la cuestión de los dominios más generales $\mathfrak G$ y $\mathfrak H$ surge la cuestión adicional de los dominios \emph{esencialmente diferentes}, que ---en un sentido que se precisará abajo--- no pueden generarse por el mismo principio de construcción. Esto conduce a una división de estos dominios en clases.

\emph{Dos dominios se llamarán pertenecientes a la misma clase, o isomorfos, si sus elementos pueden ponerse en correspondencia uno a uno de tal manera que la suma y el producto de dos elementos cualesquiera de un dominio correspondan siempre a la suma y al producto de los elementos correspondientes del otro.} De esta definición se sigue que bajo toda relación isomorfa el cero y la unidad se corresponden consigo mismos, y que todas las relaciones algebraicas entre un número finito de elementos se conservan para los elementos correspondientes, y recíprocamente. En particular, además de la unidad, todos los números racionales se corresponden consigo mismos; mientras que la totalidad de los números algebraicos ---y asimismo la totalidad de los enteros algebraicos--- pasa a sí misma, aunque un número algebraico individual bien puede corresponder a un conjugado.

Obsérvese primero que todo dominio de la totalidad $\mathfrak M(\mathfrak G)$ de todos los dominios $\mathfrak G$ que contienen una base algebraica fija $H$ es isomorfo a un dominio de la totalidad $\mathfrak M^*(\mathfrak G)$ perteneciente a otra base algebraica $H^*$. Pues el cuerpo de todos los números complejos se obtiene de toda base algebraica por extensión algebraica y por tanto tiene la misma cardinalidad que la base;\footnote{Steinitz, §§ 23 y 24.} por ello $H$ y $H^*$ tienen la misma cardinalidad, y la aplicación isomorfa deseada se obtiene asignando las cantidades de base $\eta_i$ a las cantidades $\eta_i^*$.\footnote{Esta conclusión no vale, naturalmente, para los conjuntos $\mathfrak N(\mathfrak G)$ y $\mathfrak N^*(\mathfrak G)$ pertenecientes a dos bases racionales $\Theta$ y $\Theta^*$; tampoco puede inferirse el isomorfismo de $\mathfrak N(\mathfrak G)$ y $\mathfrak N^*(\mathfrak G)$ de la expresabilidad racional mutua de $\Theta$ y $\Theta^*$. Lo que sí se obtiene, sin embargo, es una relación isomorfa entre los cuerpos $(\Theta)$ y $(\Theta^*)$.} Por otra parte, $\mathfrak M(\mathfrak G)$ contiene dominios $\mathfrak G$ no isomorfos, o esencialmente diferentes, como se sigue de los ejemplos de §§ 2, 3 y 5. Pues, como todas las relaciones algebraicas se conservan bajo la aplicación isomorfa, una base algebraica debe corresponder de nuevo a una base algebraica; pero §§ 2 y 3 mostraron que los dominios funcional y de módulos no poseen ninguna base algebraica para la cual se cumplan todas las propiedades de base del dominio de Zermelo. Como el dominio funcional de § 2 es algebraicamente enteramente cerrado, también el subconjunto $\mathfrak M(\mathfrak H)$ contiene dominios $\mathfrak H$ esencialmente diferentes.

Mostramos ahora, de modo análogo a la construcción de bases, cómo puede extraerse de $\mathfrak M(\mathfrak G)$ un subconjunto $\mathfrak L(\mathfrak G)$ tal que todos los dominios $\mathfrak G$ de $\mathfrak L(\mathfrak G)$ sean esencialmente diferentes, mientras que todo dominio de $\mathfrak M(\mathfrak G)$ ---y por tanto todo dominio $\mathfrak G$ sin excepción--- sea isomorfo a un dominio de $\mathfrak L(\mathfrak G)$. Sométase $\mathfrak M(\mathfrak G)$ a un buen orden $\Omega$; entonces un dominio $\mathfrak G$ de $\mathfrak M(\mathfrak G)$ puede ser o no isomorfo a uno precedente. Los dominios $\mathfrak G$ que no son isomorfos a ninguno precedente forman un conjunto $\mathfrak L(\mathfrak G)$, que por construcción contiene sólo dominios $\mathfrak G$ esencialmente diferentes. El argumento de inducción muestra también que todo dominio de $\mathfrak M(\mathfrak G)$ es isomorfo a un dominio de $\mathfrak L(\mathfrak G)$. Pues si $\mathfrak G_1$ fuese el primer dominio para el cual esto fallara, entonces $\mathfrak G_1$ o bien pertenecería a $\mathfrak L(\mathfrak G)$, o bien sería isomorfo a un dominio precedente $\mathfrak G_0$, que por hipótesis es isomorfo a un dominio de $\mathfrak L(\mathfrak G)$; por consiguiente, lo mismo valdría para $\mathfrak G_1$. Como todo dominio $\mathfrak G$ pertenece a algún conjunto $\mathfrak M^*(\mathfrak G)$ y es por tanto isomorfo a un dominio de $\mathfrak M(\mathfrak G)$, \emph{la totalidad de las clases de dominios $\mathfrak G$ queda representada por $\mathfrak L(\mathfrak G)$.}

Si en un dominio $\mathfrak G$ de $\mathfrak L(\mathfrak G)$ se reemplazan las cantidades de base $\eta$ por las $\eta^*$ de alguna otra base, entonces, como se mostró arriba, surge un dominio isomorfo a $\mathfrak G$. Recíprocamente, sea $\mathfrak G^*$ cualquier dominio y supóngase que es isomorfo a un dominio $\mathfrak G$ de $\mathfrak L(\mathfrak G)$. Si las cantidades de base $\eta$ corresponden a las cantidades $\eta^*$ de $\mathfrak G^*$, entonces éstas forman de nuevo una base algebraica, y $\mathfrak G^*$ contiene las mismas funciones algebraicas de los $\eta^*$ que $\mathfrak G$ contiene de los $\eta$, o sus conjugadas. \emph{Así, todo dominio isomorfo a $\mathfrak G$ surge dejando que la base algebraica $H$ recorra todas las bases algebraicas posibles en $\mathfrak G$ y en sus conjugados respecto de $(H)$}\footnote{Los elementos de estos dominios conjugados de $\mathfrak G$ pueden ciertamente expresarse racionalmente por los elementos de $\mathfrak G$ por II, pero no necesariamente de manera entera y racional; los conjugados pueden por tanto diferir de $\mathfrak G$.} \emph{---si éstos difieren de $\mathfrak G$.} De cada dominio fijo $\mathfrak G_i$ en $\mathfrak L(\mathfrak G)$ se obtiene así la clase $\mathfrak R_i$, que contiene todos y sólo los dominios $\mathfrak G$ isomorfos a $\mathfrak G_i$, aunque posiblemente cada dominio aparezca varias veces o incluso infinitamente a menudo. De $\mathfrak R_i$ se puede extraer otra vez, por el procedimiento anterior, un subconjunto $\mathfrak R_i^*$ que contiene cada dominio isomorfo a $\mathfrak G_i$ una y sólo una vez. Cuando $\mathfrak R_i^*$ recorre todas las clases, se obtiene la totalidad de todos los dominios $\mathfrak G$, cada dominio una y sólo una vez. Para caracterizar la clase $\mathfrak R_i$ puede usarse, por supuesto, en lugar de $\mathfrak G_i$, cualquier otro dominio de la clase a partir del cual todos los dominios de $\mathfrak R_i$ puedan derivarse como antes. En resumen:

\emph{De $\mathfrak M(\mathfrak G)$ se puede extraer un subconjunto $\mathfrak L(\mathfrak G)$ tal que los dominios $\mathfrak G$ de $\mathfrak L(\mathfrak G)$ estén en correspondencia biunívoca con la totalidad de las clases (sin repetición). A partir de $\mathfrak G_i$ surge la clase correspondiente $\mathfrak R_i$ haciendo recorrer a $H$ todas las bases algebraicas posibles en $\mathfrak G_i$ y en sus conjugados respecto de $(H)$. Si $\mathfrak R_i^*$ denota todos los dominios mutuamente distintos de $\mathfrak R_i$, y $\mathfrak R_i^*$ recorre todas las clases, entonces surge la totalidad de todos los dominios $\mathfrak G$, cada dominio una y sólo una vez.}\footnote{Aquí $\mathfrak M(\mathfrak G)$ debe construirse según §§ 6 y 7: por § 7 se extiende una base algebraica $H$ a toda base racional $\Theta$ con condición lateral que surja de ella, y para cada $\Theta$ se forma, por § 6 b), la totalidad $\mathfrak N(\mathfrak G)$ de todos los dominios $\mathfrak G$ que contienen $\Theta$.}

La afirmación análoga vale para los dominios $\mathfrak H$ formados con números trascendentes algebraicamente enteros.

\section*{§ 9. Un ejemplo de infinitas clases numerables de dominios $\mathfrak G$.}

La cardinalidad de las construcciones dadas en § 4 b) y § 6 b) puede leerse, por las observaciones correspondientes, como igual a la cardinalidad de todos los buenos órdenes, es decir, como $2^c$, si $c$ denota la cardinalidad del continuo. Pero de esto no puede extraerse conclusión alguna acerca de la cardinalidad de la totalidad de todos los dominios $\mathfrak G$, contados una y sólo una vez, y menos aún acerca de la cardinalidad de todas las clases. Que el número de clases ciertamente no es finito se mostrará ahora mediante un ejemplo de dominios $\mathfrak G$ infinitos numerables ---análogos a los de § 5--- que resultarán todos esencialmente diferentes y proporcionan por tanto una infinitud numerable de clases.\footnote{Los dominios $\mathfrak G$ de § 6 dan asimismo infinitas clases numerables, pero la demostración es algo más complicada.}

De modo análogo a § 5, (1), sea el dominio $\mathfrak G_\sigma$ la totalidad de las cantidades
\begin{equation}
 z=a_0+a_1(\eta)+\cdots+a_{\sigma-1}(\eta)
 \pmod{\mathfrak M_\sigma(\eta)},
\tag{1}
\end{equation}
\emph{donde $a_i(\eta)$ denota formas homogéneas de dimensión $i$ en los $\eta$; y el módulo $\mathfrak M_\sigma$ contiene como base todos los productos de potencias de dimensión $\sigma$ en los $\eta$\footnote{Naturalmente, en cada $z$ individual sólo aparecen en el módulo finitamente muchos productos de potencias de los $\eta$.} y como multiplicadores todas las funciones algebraicas enteras de los $\eta$.} Mostramos que bajo una relación isomorfa entre dos dominios $\mathfrak G_\sigma$ y $\mathfrak G_\tau$ $(\sigma\ne\tau)$ los módulos $\mathfrak M_\sigma$ y $\mathfrak M_\tau$ deben necesariamente corresponder isomórficamente; la imposibilidad se seguirá entonces fácilmente.

Denótense por $\xi$ las indeterminadas de $\mathfrak G_\tau$; bajo la relación isomorfa, correspondan a estos $\xi$ las cantidades $f(\eta)$ de $\mathfrak G_\sigma$. Como la relación isomorfa se conserva al formar cocientes, el dominio $\mathfrak G_\xi$ de todas las funciones algebraicas enteras de los $\xi$ corresponde a un dominio algebraicamente enteramente cerrado $\mathfrak T$ formado por todas las funciones algebraicamente enteras sobre los $f(\eta)$, y por tanto también algebraicamente enteras en los $\eta$.

Supóngase que (1) corresponde a una cantidad de $\mathfrak M_\tau$. Entonces, por la propiedad de módulo de $\mathfrak M_\tau$, el producto de (1) por una cantidad arbitraria $\psi(\eta)$ de $\mathfrak T$ debe pertenecer también a $\mathfrak G_\sigma$; en fórmulas:
\begin{equation}
(a_0+a_1(\eta)+\cdots+a_{\sigma-1}(\eta))\psi(\eta)
\equiv b_0+b_1(\eta)+\cdots+b_{\sigma-1}(\eta)
\pmod{\mathfrak M_\sigma}.
\tag{2}
\end{equation}
Poniendo todos los $\eta$ iguales a cero se obtiene $a_0\psi(0)=b_0$, de modo que (2) se transforma en
\begin{equation}
(a_0+a_1(\eta)+\cdots+a_{\sigma-1}(\eta))(\psi(\eta)-\psi(0))
\equiv c_1(\eta)+\cdots+c_{\sigma-1}(\eta)
\pmod{\mathfrak M_\sigma}.
\tag{3}
\end{equation}
Además de $\psi(\eta)$, las funciones
\[
 \chi_\nu(\eta)=\sqrt[\nu]{\psi(\eta)-\psi(0)}
\]
pertenecen también a $\mathfrak T$ para todo $\nu$; y como $\chi_\nu(0)=0$, las fórmulas análogas a (3) son:
\begin{equation}
(a_0+a_1(\eta)+\cdots+a_{\sigma-1}(\eta))\chi_\nu(\eta)
\equiv c_1^{(\nu)}(\eta)+\cdots+c_{\sigma-1}^{(\nu)}(\eta)
\pmod{\mathfrak M_\sigma}
\quad(\nu=1,2,\ldots).
\tag{4}
\end{equation}
Sea $A(\eta)$, de grado $\lambda$, el producto de $\chi_1=\psi(\eta)-\psi(0)$ por sus $\kappa-1$ conjugados. Puesto que $\chi_1(0)=0$, $A(\eta)$ debe comenzar con términos de al menos primera dimensión. De (4) se obtiene
\begin{equation}
a_0\chi_\nu(\eta)\equiv0\pmod{\mathfrak M_1};\qquad
 a_0^\nu\chi_1(\eta)\equiv0\pmod{\mathfrak M_\nu};\qquad
 a_0^{\nu\kappa}A(\eta)\equiv0\pmod{\mathfrak M_{\nu\kappa}},
\tag{5}
\end{equation}
y por tanto, para $a_0\ne0$, como en § 3, $\lambda>\nu\kappa$. Pero como $\nu<\frac{\lambda}{\kappa}$ contradice la cerradura algebraicamente entera de $\mathfrak T$, necesariamente $a_0=0$. De (4) se obtiene ahora
\begin{equation}
a_1(\eta)\chi_\nu(\eta)\equiv c_1^{(\nu)}(\eta)\pmod{\mathfrak M_2};\qquad
[a_1(\eta)]^{\nu\kappa}A(\eta)
\equiv[c_1^{(\nu)}(\eta)]^{\nu\kappa}\pmod{\mathfrak M_{\nu\kappa+1}},
\tag{6}
\end{equation}
y de aquí, como $A(\eta)$ comienza con términos de al menos primera dimensión, $c_1^{(\nu)}(\eta)=0$. Así, en completa analogía con (5), para todo $\nu$:
\[
 a_1(\eta)\chi_\nu(\eta)\equiv0\pmod{\mathfrak M_2};\qquad
 [a_1(\eta)]^{\nu\kappa}A(\eta)\equiv0\pmod{\mathfrak M_{2\nu\kappa}},
\]
y por tanto $a_1(\eta)=0$. Prosiguiendo así, la aplicación alternada de (5) y (6) da
\[
 a_0=0,\qquad a_1(\eta)=0,\quad \ldots,\quad a_{\sigma-1}(\eta)=0.
\]
Como la consideración análoga vale también respecto de $\mathfrak G_\tau$: \emph{Bajo toda relación isomorfa entre $\mathfrak G_\sigma$ y $\mathfrak G_\tau$ $(\sigma\ne\tau)$, los módulos $\mathfrak M_\sigma$ y $\mathfrak M_\tau$ corresponden isomórficamente.}

Supóngase, por ejemplo, que $\sigma>\tau$. A las indeterminadas $\xi$ correspondían las cantidades $f(\eta)$ de $\mathfrak G_\sigma$; y, como todas las cantidades de $\mathfrak G_\tau$ pueden representarse por polinomios en los $\xi$ mod. $\mathfrak M_\tau$, la correspondencia de $\mathfrak M_\tau$ y $\mathfrak M_\sigma$ implica que todas las cantidades de $\mathfrak G_\sigma$ se obtienen como polinomios en los $f(\eta)$ mod. $\mathfrak M_\sigma$. Puesto que $\sigma>\tau\ge1$, las indeterminadas $\eta$ ciertamente no pertenecen a $\mathfrak M_\sigma$, y por tanto, al ser expresables entera y racionalmente por los $f(\eta)$ mod. $\mathfrak M_\sigma$, debe haber un $f(\eta)$ con términos lineales no nulos. Sea, pues,
\[
 \xi:f(\eta)\equiv a_0+a_1(\eta)\pmod{\mathfrak M_2}\quad[a_1(\eta)\ne0],
\]
\[
 \xi^\tau:[f(\eta)]^\tau\equiv a_0^\tau\pmod{\mathfrak M_1}\quad\text{para }a_0\ne0,
\]
\[
 \hphantom{\xi^\tau:[f(\eta)]^\tau}\equiv [a_1(\eta)]^\tau
 \pmod{\mathfrak M_{\tau+1}}\quad\text{para }a_0=0.
\]
Ahora $\xi^\tau$ pertenece a $\mathfrak M_\tau$, mientras que $[f(\eta)]^\tau$ comienza con términos no nulos de dimensión cero o $\tau$-ésima y, como $\tau<\sigma$, no puede pertenecer a $\mathfrak M_\sigma$, en contradicción con la correspondencia de $\mathfrak M_\sigma$ y $\mathfrak M_\tau$ deducida arriba. El caso $\tau>\sigma$ se trata al mismo tiempo intercambiando $\eta$ y $\xi$. Queda demostrado, por tanto, que los dominios $\mathfrak G_\sigma$ y $\mathfrak G_\tau$ son \emph{no isomorfos, o esencialmente diferentes}. Así, \emph{los dominios $\mathfrak G_i$ $(i=1,2,\ldots\text{ in inf.})$ dan un ejemplo de infinitas clases numerables de dominios $\mathfrak G$.}

\emph{Observación.} Aunque cualesquiera dos dominios $\mathfrak G_i$ no son isomorfos, puede muy bien suceder, para dos dominios, que cada uno sea isomorfo a una parte del otro.\footnote{Este comportamiento también puede darse para cuerpos; cf. Steinitz, loc. cit., conclusión.} Sea, por ejemplo, $\tau=1$ y $\sigma$ arbitrario:
\[
\mathfrak G_1:a_0+\mathfrak M_1(\xi);\qquad
\mathfrak G_\sigma:a_0+a_1(\eta)+\cdots+a_{\sigma-1}(\eta)+\mathfrak M_\sigma(\eta).
\]
La asignación $\xi=\eta$ hace de $\mathfrak G_\sigma$ un divisor propio de $\mathfrak G_1$; recíprocamente, la asignación biyectiva entre $\mathfrak G_1$ y $\mathfrak G_\sigma$ dada por $\xi=\eta^\sigma$, $\eta=\sqrt[\sigma]{\xi}$ hace de $\mathfrak G_1$ un divisor propio de $\mathfrak G_\sigma$. Así, todo $\mathfrak G_i$ es también isomorfo a un divisor propio de sí mismo. En contraste con la noción de intersección, la noción de una clase de intersección ---es decir, una clase tal que cada uno de sus dominios fuese isomorfo a un divisor propio de todo dominio de toda otra clase--- no existe, al menos no como noción unívocamente determinada.

\section*{§ 10. Las cantidades enteras de un cuerpo arbitrario.}

Los desarrollos de los párrafos precedentes se han apoyado, para los dominios $\mathfrak G$ ---siempre que los números algebraicos no se incluyan en las propiedades definitorias III y IV de los dominios $\mathfrak G$---, sólo en el hecho de que la totalidad de todos los números complejos forma un cuerpo;\footnote{Sólo en § 7 se utilizó además la hipótesis de que el ``cuerpo primo'' de los números racionales contiene infinitos elementos; en el caso de un cuerpo primo con sólo finitos elementos, este párrafo simplemente desaparece, como veremos.} para los dominios $\mathfrak H$, se han apoyado además en que este cuerpo es algebraicamente cerrado. Por ello, los resultados obtenidos admiten directamente la \emph{generalización al cuerpo más general}, que, siguiendo a Weber y E. Steinitz,\footnote{loc. cit., introducción.} se define abstractamente como ``un sistema de elementos con dos operaciones, adición y multiplicación, que están sujetas a las leyes asociativa y conmutativa, están enlazadas por la ley distributiva, y admiten operaciones inversas irrestrictas y únicas''.\footnote{Sólo debe excluirse la división por cero.}

Todo cuerpo de tal clase contiene un elemento unidad, y por consiguiente el ``cuerpo primo'' que se deriva de él mediante las operaciones de adición y multiplicación y sus inversas. Este cuerpo primo es o bien del tipo de los números racionales (característica $0$), o bien del tipo del sistema de clases de residuos de un número primo $p$ (característica $p$).\footnote{Steinitz, § 4.} Las \emph{cantidades enteras del cuerpo primo} están entonces bien definidas como las cantidades que surgen del elemento unidad por adición y sustracción; para cuerpos primos de característica $p$, las cantidades enteras ya agotan el cuerpo primo, y no hay cantidades fraccionarias del cuerpo primo. Todo cuerpo algebraicamente cerrado contiene un ``cuerpo absolutamente algebraico'', que surge del cuerpo primo adjuntando todos los elementos algebraicos sobre el cuerpo primo; para cuerpos de característica $0$, éste es por tanto del tipo del cuerpo de todos los números algebraicos. Las \emph{cantidades algebraicamente enteras del cuerpo absolutamente algebraico} están entonces bien definidas como todas las cantidades algebraicamente enteras sobre las cantidades enteras del cuerpo primo (o, lo que es lo mismo, sobre la unidad); para cuerpos de característica $p$ tampoco hay, por tanto, cantidades algebraicamente fraccionarias del cuerpo absolutamente algebraico. Después de estas observaciones preliminares, la definición de las cantidades ``enteras'' y ``algebraicamente enteras'' puede transferirse a cuerpos arbitrarios, respectivamente a cuerpos algebraicamente cerrados: \emph{Las cantidades enteras de un cuerpo $\mathfrak R$ se definen como un dominio íntegro $\mathfrak G$ de cantidades de $\mathfrak R$ cuyo cuerpo de cocientes\footnote{Es decir, el cuerpo formado por todos los cocientes de pares de cantidades de $\mathfrak G$.} agota $\mathfrak R$, y que contiene el elemento unidad pero ninguna cantidad fraccionaria del cuerpo primo.}

\emph{Las cantidades algebraicamente enteras de un cuerpo algebraicamente cerrado $\mathfrak R$ se definen como un dominio íntegro algebraicamente enteramente cerrado $\mathfrak H$ de cantidades de $\mathfrak R$ cuyo cuerpo de cocientes agota $\mathfrak R$, y que contiene el elemento unidad pero ninguna cantidad algebraicamente fraccionaria del cuerpo absolutamente algebraico.}

Para cuerpos de característica cero, los resultados obtenidos en los párrafos precedentes se aplican literalmente, siempre que se sustituyan los ``números algebraicos'' por las cantidades del cuerpo primo, respectivamente del cuerpo absolutamente algebraico. Para cuerpos de característica $p$, los resultados se simplifican aún más, porque el cuerpo primo no contiene cantidades fraccionarias, de modo que simplemente desaparece la condición lateral para la base racional y para los sistemas que se adjuntan.\footnote{Cf. la primera nota de este párrafo.} Así resulta: \emph{La definición abstracta permite como cantidades enteras de un cuerpo de característica $p$: todo dominio íntegro derivado de una base racional arbitraria, el cuerpo mismo, y todo dominio íntegro situado entre estos dominios y el cuerpo. Correspondientemente, como cantidades algebraicamente enteras de un cuerpo algebraicamente cerrado de característica prima se tienen: todo dominio íntegro algebraicamente enteramente cerrado situado entre el dominio derivado de una base algebraica y el propio cuerpo, incluidos los dos dominios frontera.} Para los cuerpos más generales de característica prima, por tanto, como para los correspondientes cuerpos primos, se borra la distinción entre entero y fraccionario.

\medskip
\noindent Erlangen, 30 de marzo de 1915.

\clearpage
\section*{10. Las ecuaciones funcionales de la aplicación isomorfa}
\begin{center}
Math. Ann. 77 (1916), pp. 536--545
\end{center}

Según Dedekind,\footnote{Cf. por ejemplo Dirichlet-Dedekind, \emph{Vorlesungen über Zahlentheorie} (4.ª ed.), suplemento XI, \S\ 161. La designación ``aplicación isomorfa'' o ``isomorfismo'', modelada sobre la teoría de grupos, sólo aparece en la literatura posterior; Dedekind habla de las ``permutaciones de un cuerpo''.} se entiende por una \emph{aplicación isomorfa del cuerpo} $\mathfrak A$ \emph{sobre un sistema} $\mathfrak B$ ---que después resulta ser igualmente un cuerpo--- una correspondencia unívoca tal que a cada elemento de $\mathfrak A$ corresponda uno y sólo un elemento de $\mathfrak B$, y tal que a la suma, la diferencia, el producto y el cociente de dos elementos cualesquiera de $\mathfrak A$ correspondan siempre de nuevo la suma, la diferencia, el producto y el cociente de los elementos inequívocamente correspondientes de $\mathfrak B$.

Sea tal aplicación mediada por una función $f(z)$ ---unívoca según lo anterior---. Si $z$ recorre todos los elementos $x,y,\ldots$ del cuerpo $\mathfrak A$, entonces $f(z)$ queda caracterizada por las ecuaciones funcionales:
\[
\begin{array}{ll}
(1)\quad f(x+y)=f(x)+f(y); & (2)\quad f(x-y)=f(x)-f(y);\\[0.45em]
(3)\quad f(x\cdot y)=f(x)\cdot f(y); & (4)\quad f\!\left(\dfrac{x}{y}\right)=\dfrac{f(x)}{f(y)}.
\end{array}
\]

\emph{En lo que sigue se darán las soluciones unívocas más generales} $f(z)$ \emph{de estas ecuaciones funcionales cuando} $x,y,\ldots$ \emph{recorren todos los números reales y complejos; es decir, la función más general} $f(z)$ \emph{que realiza una aplicación isomorfa del cuerpo de todos los números complejos}.\footnote{Esta cuestión me fue planteada ocasionalmente por el señor Landau. Como advierto posteriormente, una parte de las soluciones se encuentra ya en H. Lebesgue, \emph{Sur les transformations ponctuelles, transformant les plans en plans}, Atti Acad. Torino, 1906/07; e implícitamente también en A. Ostrowski, \emph{Über einige Fragen der allgemeinen Körpertheorie}, Crelle's Journal 143 (1913), \S\ 2.} De manera completamente análoga se obtiene la solución para cuerpos arbitrarios.

G. Hamel\footnote{G. Hamel, \emph{Eine Basis aller Zahlen und die unstetigen Lösungen der Funktionalgleichung:} $f(x+y)=f(x)+f(y)$, Math. Ann. 60, p. 459 (1905).} construyó las soluciones más generales de la ecuación funcional (1) mediante una base \emph{lineal} de todos los números. La construcción que aquí se da para las soluciones de (1) a (4) ---por tanto, para la función de aplicación respecto de la cual son invariantes las operaciones racionales y algebraicas--- utiliza la base \emph{racional} y \emph{algebraica} de todos los números. Después de examinar con más detalle la aplicación isomorfa en 1, se dan en 2 condiciones necesarias para $f(z)$; en 3 se reconoce que éstas son suficientes y conducen a la construcción efectiva.\footnote{Cf. las consideraciones paralelas de mi trabajo: \emph{Die allgemeinsten Bereiche aus ganzen transzendenten Zahlen}. Math. Ann. 77, p. 103 (1915), en especial \S\S\ 8 y 10.} Las soluciones de (1) a (4), salvo las conocidas excepciones $f(z)=z$ o $\overline z$, son discontinuas; en 4 se muestra además que son ``extremadamente discontinuas'', hecho que G. Hamel demostró para las soluciones reales de (1), pero que no tiene por qué cumplirse en el dominio complejo.\footnote{La expresión de Hamel ``totalmente discontinuo'' se utiliza en el resto de la literatura en otro sentido.}

\subsection*{1.}

Extraigamos primero algunas consecuencias de las ecuaciones funcionales (1) a (4), directamente para cuerpos generales $\mathfrak A$. Por (1), de la univocidad se sigue $f(0)=0$. Si $f(y)$ no se anula, entonces tampoco puede anularse el $y$ originario, y las ecuaciones funcionales (1) a (4) muestran, por tanto, que \emph{el sistema de imagen} $\mathfrak B$ \emph{de todos los valores} $f(z)$ \emph{forma de nuevo un cuerpo}. Recíprocamente, si se prescribe la condición inicial $f(0)=0$, entonces (2) implica la univocidad de la función $f(z)$: \emph{la exigencia de univocidad y la condición inicial} $f(0)=0$ \emph{son, por tanto, equivalentes}. De (4) se sigue que $f(z)$ no puede anularse idénticamente; y de (3), además, que $f(z)$ sólo se anula para $z=0$. Pues de $f(y)=0$ y $y\ne0$ se seguiría, ya que $x/y$ pertenece también a $\mathfrak A$, que
\[
f(x)=f\!\left(\frac{x}{y}\cdot y\right)=f\!\left(\frac{x}{y}\right)\cdot f(y)=0,
\]
es decir, en contradicción con (4), la anulación idéntica de $f(z)$. De $f(x)=f(y)$ se sigue así $x=y$; esto es, a cada valor de $f(z)$ corresponde también uno y sólo un valor de $z$: \emph{$f(z)$ es una función biyectiva de} $z$. Como muestran las ecuaciones funcionales, la aplicación mediada por la función inversa es de nuevo isomorfa. Puesto que toda operación racional está compuesta por un número finito de adiciones, multiplicaciones y sus inversas, también puede decirse:

\emph{Mediante la aplicación isomorfa queda establecida una relación biunívoca entre los elementos de $\mathfrak A$ y los de $\mathfrak B$, de tal modo que todas las relaciones racionales existentes entre un número finito de elementos de $\mathfrak A$,
\[
F(x,y,\ldots,t)=0,
\]
se conservan en $\mathfrak B$, y recíprocamente.}

Conviene observar todavía que, para un cuerpo, $f(z)$ queda ya caracterizada por las ecuaciones funcionales (1) y (3), por la exigencia de que $f(z)$ no se anule idénticamente, y por la condición inicial $f(0)=0$. En efecto, (2) se sigue de (1) reemplazando $x$ por el elemento $(x-y)$ perteneciente a $\mathfrak A$; después, (2) y $f(0)=0$ dan la univocidad de $f(z)$. Como se mostró arriba, de (3) se sigue además, reemplazando $x$ por el elemento $x/y$ perteneciente a $\mathfrak A$, que $f(y)$ no puede anularse para $y\ne0$, y por consiguiente la validez de (4) y al mismo tiempo la biyectividad. Si en lugar de un cuerpo se toma como base un dominio íntegro $\mathfrak J$, entonces $x/y$ no tiene por qué pertenecer a $\mathfrak J$, y de (3) no se puede concluir la biyectividad. Pero aquí basta la formulación más conocida: \emph{una aplicación isomorfa es una correspondencia biunívoca entre los elementos de $\mathfrak J$ y de $\mathfrak J'$, tal que a la suma y al producto siempre correspondan la suma y el producto de los elementos correspondientes}.\footnote{Cf. sobre núm. 1: Dedekind, loc. cit., \S\ 161.}

\subsection*{2.}

Desde ahora, por simplicidad, tomaremos como base, en lugar de $\mathfrak A$, el cuerpo $\mathfrak C$ de todos los números reales y complejos. Se trata ante todo de discutir los sistemas de valores que toma $f(z)$. De (3) o (4) se sigue $f(1)=1$, y de ahí, por las ecuaciones funcionales, para todo \emph{número racional} $\alpha$: $f(\alpha)=\alpha$; por tanto, bajo la aplicación, todo número racional se corresponde consigo mismo. Como un \emph{número algebraico} $\beta$ satisface una ecuación irreducible con coeficientes numéricos racionales $F(\beta,\alpha)=0$, por 1 y por lo que acabamos de observar se obtiene para $f(\beta)$ la ecuación $F(f(\beta),\alpha)=0$; \emph{todo número algebraico pasa, por tanto, a sí mismo o a uno de sus conjugados}, y por la biyectividad, la totalidad de los números algebraicos corresponde de nuevo al cuerpo de todos los números algebraicos.

Para reconocer el comportamiento de los números \emph{trascendentes}, y al mismo tiempo separar los valores conjugados, tomamos como base una \emph{base racional de todos los números},\footnote{Cf. mi trabajo citado sobre los ``números trascendentes enteros'', \S\ 6.} es decir, un sistema $\Theta$ de números $\vartheta$ que permite expresar todos los números racionalmente ---con coeficientes numéricos racionales\footnote{Por ``función racional'' entenderé siempre una cuyos coeficientes son también números racionales.}---, mientras que, bajo al menos un buen orden, ningún número de base es expresable racionalmente por un número finito de los precedentes. La existencia de esta base se sigue del teorema del buen orden exactamente de modo análogo a las bases lineales y algebraicas.

Sométase $\Theta$ a un buen orden $\Omega$; entonces cada número es o bien expresable racionalmente por un número finito de los precedentes, o bien es racionalmente independiente de los precedentes. Estos últimos números forman la base $\Theta$, y el argumento de inducción muestra que efectivamente todo número es racionalmente expresable por un número finito de $\vartheta$.

Para cada número de base $\vartheta$ está definido el ``cuerpo de segmento $\mathfrak K(\vartheta)$'': consta de todas las combinaciones racionales de aquellos números de base que preceden a $\vartheta$ en el buen orden. Como cuerpo de segmento del primer número de base debe considerarse el cuerpo $\mathfrak R$ de todos los números racionales. El número $\vartheta$ puede presentar dos comportamientos respecto de $\mathfrak K(\vartheta)$: puede ser algebraicamente dependiente de $\mathfrak K(\vartheta)$ o algebraicamente independiente de él. Puesto que, por 1, todas las relaciones válidas en $\mathfrak C$ valen también en el cuerpo de imagen y recíprocamente, la totalidad de los valores $f(\vartheta)$ correspondientes a los $\vartheta$ forma una base del dominio de imagen, bien ordenada mediante el buen orden de $\Theta$ mismo. En particular, al cuerpo de segmento $\mathfrak K(\vartheta)$ corresponde el cuerpo de segmento $\mathfrak K(f(\vartheta))$; y según que $\vartheta$ sea algebraicamente dependiente o independiente de $\mathfrak K(\vartheta)$, lo mismo vale para $f(\vartheta)$ respecto de $\mathfrak K(f(\vartheta))$; y en el caso de dependencia corresponden también las ecuaciones irreducibles de $\vartheta$ y de $f(\vartheta)$.

La totalidad de los valores $\vartheta$ que son, en cada caso, algebraicamente independientes de $\mathfrak K(\vartheta)$ forman ---como de nuevo muestra la inducción--- una \emph{base algebraica de todos los números},\footnote{Cf. E. Zermelo, \emph{Über ganze transzendente Zahlen}, \S\ 1, Math. Ann. 75 (1914).} es decir, un sistema $H$ de números $\eta$ que permite expresar algebraicamente todos los números, mientras que entre ellos no existen relaciones algebraicas. A esta base algebraica $H$ corresponde en el dominio de imagen de nuevo una base algebraica $Z$, que por la biyectividad tiene la misma cardinalidad que $H$, a saber, la cardinalidad del continuo, pues es sabido que una extensión algebraica no aumenta la cardinalidad.\footnote{Cf. E. Steinitz, \emph{Algebraische Theorie der Körper}, Crelle's Journal 137 (1910), \S\ 24.}

Del conocimiento de los valores $f(\vartheta)$ se obtiene, sin embargo, inmediatamente $f(z)$ para todos los sistemas de valores, ya que de $z=\Psi(\vartheta)$ se sigue siempre $f(z)=\Psi(f(\vartheta))$.

\subsection*{3.}

Mostramos ahora que las condiciones necesarias encontradas realizan efectivamente también la construcción de $f(z)$. La base $Z$ que corresponde biunívocamente a la base algebraica $H$ se construye y se asigna a $H$ del modo siguiente. Se toma un segundo buen orden $\Omega'$, que proporciona una base algebraica $Z'$ de todos los números; sea $Z$, con elementos $\zeta$, un subconjunto de $Z'$ de la misma cardinalidad que $Z'$, y por tanto que $H$. A cada elemento $\eta_i$ de $H$ se le asigna entonces biunívocamente un $\zeta_i$ (con la misma numeración) por la regla: $\zeta_i$ es el \emph{primero}, en el buen orden $\Omega'$, de todos los elementos de $Z$ que no están asignados a ningún elemento de $H$ precedente a $\eta_i$.

Defínase ahora $f(z)$:
\begin{enumerate}
\item[a)] para todos los números racionales $\alpha$, por $f(0)=0$, $f(\alpha)=\alpha$;
\item[b)] para todas las cantidades de la base algebraica $H$, por $f(\eta_i)=\zeta_i$;
\item[c)] para todas las demás cantidades $\vartheta$ de la base racional $\Theta$ ---las que dependen algebraicamente del cuerpo de segmento $\mathfrak K(\vartheta)$ y por consiguiente satisfacen una ecuación $F(\vartheta;\vartheta_1,\ldots,\vartheta_r)=0$ irreducible en $\mathfrak K(\vartheta)$---, como raíz de la ecuación $F(f(\vartheta);f(\vartheta_1),\ldots,f(\vartheta_r))=0$;
\item[d)] para todos los demás valores $z$, que por la definición de $\Theta$ son representables como $z=\psi(\vartheta_{k_1},\ldots,\vartheta_{k_r})$, por $f(z)=\psi(f(\vartheta_{k_1}),\ldots,f(\vartheta_{k_r}))$.\footnote{En (d) está contenido (a) como caso especial.}
\end{enumerate}

Para mostrar que el $f(z)$ así definido satisface efectivamente las ecuaciones funcionales (1) a (4), basta por núm. 1 probarlo para (1) y (3), ya que $\mathfrak C$ es un cuerpo, $f(z)$ no puede anularse idénticamente a causa de (a) y (b), y además satisface la condición inicial $f(0)=0$.

Mediante la base racional $\Theta$, exprésense $x$, $y$, $x+y$ por:
\[
x=\psi_1(\vartheta_1,\ldots,\vartheta_t),\quad
 y=\psi_2(\vartheta_1,\ldots,\vartheta_t),\quad
 x+y=\psi_3(\vartheta_1,\ldots,\vartheta_t);
\]
el hecho de que $f(z)$ satisfaga la ecuación funcional (1),
\[
f(x)+f(y)-f(x+y)=0,
\]
es, por (d), idéntico a que de
\[
\psi_1(\vartheta)+\psi_2(\vartheta)-\psi_3(\vartheta)=\frac{H_1(\vartheta)}{H_2(\vartheta)}=0
\]
se siga siempre:
\[
\psi_1(f(\vartheta))+\psi_2(f(\vartheta))-\psi_3(f(\vartheta))
=\frac{H_1(f(\vartheta))}{H_2(f(\vartheta))}=0,
\]
donde $H_1,H_2$ significan funciones racionales enteras de sus argumentos. Exactamente a la misma forma de la condición conduce también la ecuación funcional (3); y la satisfacción de todas las ecuaciones funcionales es, por consiguiente, idéntica a la existencia de la condición ---necesaria también por núms. 1 y 2--- de que, para toda función racional entera $H(\vartheta)$, se siga siempre: de $H(\vartheta)=0$, también $H(f(\vartheta))=0$; de $H(\vartheta)\ne0$, también $H(f(\vartheta))\ne0$, esto último a causa de la aparición del denominador.

Que esto es efectivamente así lo muestra el argumento de inducción. Supongamos que en $H(\vartheta_1,\ldots,\vartheta_n)$ el número de base $\vartheta_n$ es el último, en el buen orden, entre $\vartheta_1,\ldots,\vartheta_n$;\footnote{Las expresiones $H(\vartheta)$ de grado cero en los $\vartheta$ pasan a sí mismas bajo la aplicación y, por tanto, no necesitan investigarse más.} entonces $H(\vartheta_1,\ldots,\vartheta_n)=0$ puede considerarse como una relación que satisface $\vartheta_n$ respecto del cuerpo de segmento $\mathfrak K(\vartheta_n)$, e igualmente $H(f(\vartheta_1),\ldots,f(\vartheta_n))=0$ como una relación de $f(\vartheta_n)$ respecto de $\mathfrak K(f(\vartheta_n))$.

Si no todas las relaciones $H(\vartheta)=0$ se satisficieran también para $f(\vartheta)$, y recíprocamente, tendría que haber un primer número de base ---digamos $\vartheta_n$--- para el cual al menos una relación válida respecto de $\mathfrak K(\vartheta_n)$ no se conservara en el cuerpo de imagen, o recíprocamente, mientras que los cuerpos de segmento precedentes $\mathfrak K(\vartheta_n)$ y $\mathfrak K(f(\vartheta_n))$ fueran isomorfos.

En particular, por (a), el cuerpo de segmento del primer número de base de $\Theta$, el cuerpo $\mathfrak R$ de todos los números racionales, es isomorfo a su cuerpo de imagen. Supóngase ahora que $H(\vartheta_n)$ tiene la forma:
\[
H(\vartheta_n)=\vartheta_n^{\mu}+a_1(\vartheta)\vartheta_n^{\mu-1}+\cdots+a_\mu(\vartheta),
\]
donde $a_i(\vartheta)$ son cantidades de $\mathfrak K(\vartheta_n)$. Entonces hay dos posibilidades:

1) $\vartheta_n$ es algebraicamente \emph{independiente} de $\mathfrak K(\vartheta_n)$. Entonces $H(\vartheta_n)=0$ se satisface idénticamente en $\vartheta_n$,\footnote{Por la expresabilidad de $x,y,\ldots$ mediante los $\vartheta$, esta forma de relación puede muy bien aparecer.} se tiene $a_i(\vartheta)=0$ para todos los coeficientes y, por tanto, debido al isomorfismo de $\mathfrak K(\vartheta_n)$ y $\mathfrak K(f(\vartheta_n))$, también $a_i(f(\vartheta))=0$; de $H(\vartheta_n)=0$ se sigue, pues, siempre $H(f(\vartheta_n))=0$. Recíprocamente, sea $H(f(\vartheta_n))=0$; como $\vartheta_n$, por ser algebraicamente independiente de $\mathfrak K(\vartheta_n)$, pertenece a la base algebraica $H$, $f(\vartheta_n)$ pertenece por (b) a la base algebraica $Z$ y es algebraicamente independiente de $\mathfrak K(f(\vartheta_n))$. Por tanto se anulan todos los $a_i(f(\vartheta))$, y por el isomorfismo todos los $a_i(\vartheta)$; de $H(f(\vartheta_n))=0$ se sigue $H(\vartheta_n)=0$, o de $H(\vartheta_n)\ne0$ también $H(f(\vartheta_n))\ne0$.

2) $\vartheta_n$ es algebraicamente \emph{dependiente} de $\mathfrak K(\vartheta_n)$. Entonces $H(\vartheta_n)$ es divisible por la ecuación $F(\vartheta_n)$ irreducible respecto de $\mathfrak K(\vartheta_n)$, es decir, idénticamente en $t$:
\[
H(t;a_i)=F(t;\vartheta_n)\cdot Q(t;a_i).
\]
Como todos los coeficientes que aparecen pertenecen a $\mathfrak K(\vartheta_n)$, en el cuerpo isomorfo $\mathfrak K(f(\vartheta_n))$ se satisface la relación correspondiente:
\[
H(t;a_i(f))=F(t;f(\vartheta_n))\cdot Q(t;a_i(f)).
\]
Pero por (c) $f(\vartheta_n)$ fue elegido de tal modo que es raíz de $F(t;f(\vartheta_n))$; de $H(\vartheta_n)=0$ se sigue, por tanto, $H(f(\vartheta_n))=0$.

A la inversa, si se tiene $H(f(\vartheta_n))=0$, de la divisibilidad de $H(t;a_i(f(\vartheta)))$ por la función $F(t;f(\vartheta_n))$ irreducible respecto de $\mathfrak K(f(\vartheta_n))$ se sigue igualmente, para el cuerpo isomorfo $\mathfrak K(\vartheta_n)$, la divisibilidad de $H(t;a_i(\vartheta))$ por $F(t;\vartheta_n)$; y como $F(t;f(\vartheta_n))$ surgió de la ecuación irreducible de $\vartheta_n$ por la relación isomorfa, y esta relación es biunívoca, $F(t;\vartheta_n)$ representa necesariamente la función irreducible cuya raíz es $\vartheta_n$. De $H(f(\vartheta_n))=0$ se sigue, pues, $H(\vartheta_n)=0$, o de $H(\vartheta_n)\ne0$ también $H(f(\vartheta_n))\ne0$.

De cuerpos de segmento isomorfos $\mathfrak K(\vartheta_n)$ y $\mathfrak K(f(\vartheta_n))$ surgen, por tanto, al adjuntar $\vartheta_n$ y $f(\vartheta_n)$ respectivamente, también cuerpos isomorfos; la suposición de que esto falle para algún $\vartheta_n$ conduce a una contradicción.

Con ello queda probado que \emph{la función} $f(z)$ \emph{definida por (a) a (d) representa efectivamente una solución de las ecuaciones funcionales (1) a (4), y precisamente, por 2, la más general}.

Todas las consideraciones utilizadas para construir $f(z)$ sólo han usado el hecho de que la totalidad $\mathfrak C$ de todos los números forma un cuerpo, y valen por consiguiente para todo cuerpo $\mathfrak A$ definido abstractamente. Debe observarse que al cuerpo $\mathfrak R$ de los números racionales corresponde el ``cuerpo primo'' de $\mathfrak A$, es decir, el cuerpo derivado del elemento unidad de $\mathfrak A$ por las operaciones de adición y multiplicación y sus inversas. Si, por tanto, en (a) a (d) se sustituyen los números racionales por los elementos del cuerpo primo, la función $f(z)$ así obtenida \emph{realiza la aplicación isomorfa más general de un cuerpo abstracto arbitrariamente definido}.

\subsection*{4.}

Resta mostrar que las funciones $f(z)$ definidas para el cuerpo $\mathfrak C$ de todos los números reales y complejos ---que, salvo $f(z)=z$ o $\overline z$, son notoriamente discontinuas--- son \emph{extremadamente discontinuas}; esto es, que en toda vecindad de cualquier número real o complejo $z_0$ hay valores $z$ para los que $f(z)$ se aproxima arbitrariamente a cualquier valor $Z_0$ prescrito.

Esta extrema discontinuidad pertenece, como G. Hamel demostró en el lugar citado, a las soluciones reales $\varphi(x)$ de la ecuación funcional (1) definidas para todos los números reales; pero ---como muestran los ejemplos que se dan abajo--- no necesariamente a las soluciones definidas en lo complejo. Hamel razona como sigue. Si $\varphi(x)$ es distinta de $C\cdot x$, entonces existen al menos dos valores de la base lineal,\footnote{La base lineal de todos los números (reales) está dada por un sistema de números (reales) que permite expresar todos los números (reales) linealmente, con coeficientes racionales, mientras que entre cada número finito de números de base no existe ninguna relación lineal con coeficientes racionales.} digamos $x_1$ y $x_2$, tales que $\varphi(x_1):x_1\ne\varphi(x_2):x_2$; entonces las dos ecuaciones
\[
\alpha x_1+\beta x_2=a,
\qquad
\alpha\varphi(x_1)+\beta\varphi(x_2)=b
\]
tienen una solución en números reales $\alpha,\beta$; y por consiguiente $a$ y $b$ pueden aproximarse arbitrariamente mediante números racionales $\alpha,\beta$. Pero, por ser racionales $\alpha,\beta$, la cantidad $\alpha\varphi(x_1)+\beta\varphi(x_2)$ representa realmente $\varphi(\alpha x_1+\beta x_2)$. Este argumento falla para sistemas de valores complejos; de hecho, también en lo complejo se pueden indicar soluciones de (1) que son discontinuas pero no extremadamente discontinuas, por ejemplo
\[
\varphi(z)=\varphi(x+iy)=\varphi(x)+i\bigl(y+\varphi(x)\bigr),
\]
donde $\varphi(x)$ denota una solución real extremadamente discontinua en lo real; o, aún más simplemente, $\varphi(z)=\varphi(x)$. En ambos casos la parte real de $\varphi(z)$ es extremadamente discontinua, mientras que la parte imaginaria queda codeterminada por la parte real.

Este comportamiento, y al mismo tiempo el comportamiento de la función $f(z)$, se vuelve completamente claro cuando se introduce el concepto de \emph{rango}. Escribimos $z=x+iy$, $\varphi(z)=X+iY$, y decimos que $\varphi(z)$ tiene rango cuatro, tres, dos o uno respectivamente según que no existan, o existan una, dos o tres relaciones linealmente independientes:
\[
c_1x+c_2y+c_3X+c_4Y=0,
\]
donde, naturalmente, las $c$ son números reales.

1) Sea $\varphi(z)$ de rango cuatro. Entonces existen al menos cuatro valores de la base lineal, digamos $z_1,z_2,z_3,z_4$, tales que el determinante
\[
\begin{vmatrix}
x_1&x_2&x_3&x_4\\
y_1&y_2&y_3&y_4\\
X_1&X_2&X_3&X_4\\
Y_1&Y_2&Y_3&Y_4
\end{vmatrix}
\]
no se anula. Las ecuaciones
\[
\begin{aligned}
\alpha x_1+\beta x_2+\gamma x_3+\delta x_4&=a,\\
\alpha y_1+\cdots+\delta y_4&=b,\\
\alpha X_1+\cdots+\delta X_4&=A,\\
\alpha Y_1+\cdots+\delta Y_4&=B
\end{aligned}
\]
tienen, por tanto, una solución en números reales $\alpha,\ldots,\delta$, y $a,b,A,B$ pueden aproximarse arbitrariamente mediante números racionales $\alpha,\beta,\gamma,\delta$. Además se tiene:
\[
\alpha(X_1+iY_1)+\beta(X_2+iY_2)+\gamma(X_3+iY_3)+\delta(X_4+iY_4)
=\varphi(\alpha z_1+\beta z_2+\gamma z_3+\delta z_4).
\]
``En general'', las soluciones $\varphi(z)$ son, pues, extremadamente discontinuas también en lo complejo; los \emph{valores de discontinuidad} de $\varphi(z)$ en la vecindad de $z_0=a+ib$ forman una variedad lineal bidimensional.

2) Sea $\varphi(z)$ de rango tres. Entonces existen al menos tres sistemas de valores de la base lineal, digamos $z_1,z_2,z_3$, tales que el determinante anterior tiene rango tres. Los valores $a,b,A,B$ que satisfacen la relación
\[
c_1a+c_2b+c_3A+c_4B=0
\]
y sólo éstos pueden aproximarse arbitrariamente; los \emph{valores de discontinuidad} forman una variedad lineal \emph{unidimensional} determinada por $z_0=a+ib$. Como muestran los ejemplos indicados, este caso puede ocurrir realmente.

3) Sea $\varphi(z)$ de rango dos. Puesto que siempre se pueden elegir dos números complejos $z_1,z_2$ tales que
\[
\begin{vmatrix}x_1&x_2\\ y_1&y_2\end{vmatrix}\ne0,
\]
las relaciones tienen la forma:
\[
X=\lambda_1x+\lambda_2y,\qquad Y=\mu_1x+\mu_2y.
\]
La expresión resultante
\[
\varphi(z)=(\lambda_1x+\lambda_2y)+i(\mu_1x+\mu_2y)
\]
representa precisamente la solución \emph{continua} más general de la ecuación funcional (1) en el dominio de los números complejos.

4) Sea $\varphi(z)$ de rango uno. Este caso no puede ocurrir en el dominio de los números complejos; en el dominio de los números reales significa la solución \emph{continua} $\varphi(x)=C\cdot x$.

Mostramos ahora que las \emph{soluciones} $f(z)$ \emph{de las ecuaciones funcionales (1) a (4) definidas para todos los números reales y complejos sólo pueden tener rango dos o rango cuatro}. El rango dos da sólo las dos funciones \emph{continuas} $f(z)=z$ o $\overline z$, como muestra la especialización $f(1)=1$ y $f(i)=\pm i$ en conexión con 3). Como el rango uno ya fue excluido arriba, queda por excluir el rango tres. Supongamos entonces que existe al menos una relación:
\[
c_1x+c_2y+c_3X+c_4Y=0,
\]
de modo que $f(z)$ sea de rango tres, o eventualmente de rango menor, y mostremos que bajo esta hipótesis siempre se da este último caso. Como de nuevo muestra la especialización $f(1)=1$, $f(i)=\pm i$, la relación se transforma en:
\[
c_3(X-x)+c_4(Y\mp y)=0,
\]
donde $c_3$ y $c_4$ no se anulan simultáneamente. Supongamos primero, por ejemplo, que $c_4=0$. La relación $(X-x)=0$ significa entonces que a un $z$ puramente imaginario corresponde también un $f(z)$ puramente imaginario. Así, para todo $y$ real se tiene $f(iy)=i\cdot\mathfrak Y$, donde $\mathfrak Y$ es de nuevo real. Pero de esto, por $f(iy)=\pm i f(y)$, se sigue también $f(y)=\pm\mathfrak Y$; por tanto a todo $z$ real corresponde también un $f(z)$ real, y a causa de $(X-x)$ todos los valores reales pasan a sí mismos. Se tiene así sólo $f(z)=z$ o $\overline z$, luego efectivamente rango dos. De modo completamente análogo se concluye si $c_3=0$, $c_4\ne0$.

Sean ahora $c_3$ y $c_4$ distintos de cero, de modo que la relación pueda escribirse en la forma
\[
(Y\mp y)=c\cdot(X-x)
\]
con $c$ real, no nulo y finito. Esto significa que para $y=\pm cx$ se tiene también $Y=c\cdot X$. Por tanto, teniendo en cuenta las ecuaciones funcionales, para todo $r$ real vale:
\[
f\{r\cdot(1\pm ic)\}=X(1+ic)=f(r)\cdot(1+i f(c)),
\]
donde $X$ es de nuevo real. Pero aquí, por núm. 1, el factor de $f(r)$ no puede anularse idénticamente, y la división da:
\[
f(r)=X(\sigma+i\tau)
\]
con $\sigma$ y $\tau$ reales e independientes de $r$ y de $X$.

En particular, también para $r=1$ corresponde un $X_0$ real, y la condición $1=X_0(\sigma+i\tau)$ muestra que necesariamente $\tau=0$, de modo que $f(r)=\sigma\cdot X$. Así, a todo $z$ real corresponde también un $f(z)$ real; equivalentemente, de $y=0$ se sigue necesariamente $Y=0$. Con ello la relación lineal para $z$ real pasa a ser $X-x=0$;\footnote{Aquí se utiliza el hecho de que $c\ne0$, de modo que ni $c_3$ ni $c_4$ se anulan, mientras que antes sólo se utilizó $c_4\ne0$; por eso hubo que tratar de antemano el caso en que uno de estos dos se anula.} cada valor real corresponde a sí mismo. Se tiene de nuevo sólo $f(z)=z$ o $\overline z$, luego efectivamente rango dos. Así queda excluido el rango tres en todos los casos. Pero, puesto que ---como muestran las investigaciones de los primeros números--- existen efectivamente soluciones discontinuas y éstas deben tener por consiguiente necesariamente rango cuatro, se obtiene por 1) el teorema: \emph{``Todas las soluciones} $f(z)$ \emph{de las ecuaciones funcionales (1) a (4) ---con excepción de las soluciones continuas} $f(z)=z$ \emph{o} $\overline z$--- \emph{son extremadamente discontinuas en lo real y en lo complejo}.''\footnote{Lebesgue muestra loc. cit. que la parte real y la parte imaginaria de $f(z)$ son ambas funciones ``no medibles'' de $x$ e $y$; como muestra el ejemplo del comienzo de núm. 4, $\varphi(z)=\varphi(x)+i(y+\varphi(x))$, de ello no se sigue todavía la extrema discontinuidad en lo complejo.}

\medskip
\noindent Erlangen, 30 de octubre de 1915.

\newpage

\begin{center}
{\Large\bfseries 11. Ecuaciones con grupo prescrito}\par
\vspace{1.2em}
\emph{Math. Ann.} 78 (1918), pp. 221--229
\end{center}

\vspace{1.5em}

El problema de la construcción de ecuaciones con grupo prescrito puede atacarse en dos direcciones, que pueden caracterizarse brevemente como la ``irracional'', o la que caracteriza las raíces, y la ``racional'', o la que caracteriza los coeficientes.

En la dirección ``irracional'', que trabaja de manera teórico-funcional y aritmética, se sitúa el teorema de Kronecker según el cual todos los cuerpos abelianos en el dominio de los números racionales son cuerpos ciclotómicos, y los teoremas correspondientes para cuerpos relativamente abelianos con respecto a cuerpos cuadráticos. Estos teoremas dan, por una parte, la realizabilidad de todos los grupos abelianos con respecto a esos campos especiales de racionalidad; por otra parte, proporcionan la totalidad de las ecuaciones y al mismo tiempo dan una comprensión más profunda de la estructura de los cuerpos numéricos definidos por ellas. Pero cada teorema particular está restringido a un campo especial de racionalidad, y cada nuevo campo de racionalidad exige un tratamiento completamente nuevo.

La dirección ``racional'', que trabaja algebraicamente, se apoya en el teorema de irreducibilidad de Hilbert, según el cual en la representación de los coeficientes puede uno limitarse a una representación paramétrica. A ella pertenece, por ejemplo, la existencia de arbitrariamente muchas ecuaciones con grupo alternante respecto de cualquier campo de racionalidad. Aquí falta, naturalmente, la comprensión indicada arriba de la estructura de los cuerpos definidos por las ecuaciones; pero la ventaja de esta dirección ``racional'' consiste en que la realizabilidad de los grupos se demuestra simultáneamente para todos los campos de racionalidad, aunque las representaciones paramétricas conocidas hasta ahora no dan, en general, la totalidad de las ecuaciones.\footnote{Para ecuaciones resolubles, la representación paramétrica de los coeficientes puede reemplazarse por la de las raíces. Recientemente F. Mertens ha dado representaciones generales de raíces para ciertos grupos de grado 8: ``Gleichungen mit Quaternionengruppen'', Sitzb. d. Ak. d. Wiss., Wien, Abt. IIa, vol. 125, p. 735. Según tomo de esta nota, ya antes G. Bucht había indicado las expresiones generales de raíces para formas normales siempre producibles de las ecuaciones de grados 3 y 4 y de las ecuaciones de grado 8 con los grupos mencionados: ``Über einige algebraische Körper achten Grades'', Arkiv for Math., Astron. och Physik, vol. 6, núm. 30. [Adición en la corrección.]}

Lo que sigue es una contribución a la dirección ``racional'': se dan condiciones suficientes para la existencia de representaciones paramétricas que, para un campo de racionalidad arbitrario, proporcionan la totalidad de las ecuaciones con grupo prescrito. Estas condiciones consisten en que el cuerpo de invariantes correspondiente al grupo ---el \emph{Gattungsbereich} lagrangiano--- pueda representarse racionalmente mediante funciones independientes del campo, es decir, que posea una ``base mínima''; o, dicho de otro modo, que sea isomorfo al cuerpo de todas las funciones racionales de $n$ indeterminadas.

Esto se cumple, como se mostrará aún en el \S 2, por ejemplo para todos los grupos que aparecen en las ecuaciones de grados 3 y 4. La construcción efectiva y la discusión más precisa de estas ecuaciones de grados 3 y 4 se encuentran en la disertación de F. Seidelmann,\footnote{\emph{Die Gesamtheit der kubischen und biquadratischen Gleichungen mit Affekt bei beliebigem Rationalitätsbereich}. Erlangen, 1916.} de la cual sigue inmediatamente un extracto a esta comunicación.

La cuestión de la base mínima de los dominios de género lagrangianos fue planteada ---como ya se mencionó en el trabajo ``Körper und Systeme rationaler Funktionen'' (Math. Ann. 76)--- independientemente de la conexión expuesta, por E. Fischer en conversación conmigo, y dio el impulso a las investigaciones presentes.\footnote{Compárese la parte 2 de la comunicación preliminar sobre ``Rationale Funktionenkörper'', Jahresb. d. d. Math. Vereinig., vol. 22, 1913.}

\subsection*{\S 1. Condiciones suficientes para la representación paramétrica.}

\textbf{1.} La posibilidad de reducir la construcción de ecuaciones con grupo prescrito a una representación paramétrica se sigue del teorema de irreducibilidad de Hilbert, que dice lo siguiente. Supóngase que la ecuación
\begin{equation}\tag{1}
f(x)=x^n+F_1(\lambda_1,\ldots,\lambda_r)x^{n-1}+\cdots+F_n(\lambda_1,\ldots,\lambda_r)=0,
\end{equation}
donde $F_1(\lambda_1,
\ldots,\lambda_r),\ldots,F_n(\lambda_1,\ldots,\lambda_r)$ significan funciones racionales de los parámetros independientes $\lambda_1,\ldots,\lambda_r$, con coeficientes en un cuerpo numérico\footnote{En lugar del cuerpo numérico podría figurar también un campo general de racionalidad.} $\Omega$, tiene, con respecto al cuerpo $\Omega(\lambda_1,
\ldots,\lambda_r)$ de todas las funciones racionales de $\lambda_1,
\ldots,\lambda_r$ con coeficientes en $\Omega$, el grupo $\Gamma$. Entonces se pueden reemplazar los parámetros $\lambda$ por arbitrariamente muchos números racionales de tal modo que la ecuación numérica resultante
\begin{equation}\tag{2}
g(x)=x^n+a_1x^{n-1}+\cdots+a_n=0
\end{equation}
tenga precisamente el grupo $\Gamma$ respecto de $\Omega$. Tal $g(x)$ se designará más exactamente por $g_\Gamma$.

Se pregunta ahora si puede darse una representación paramétrica \emph{tal} (1) que la ecuación numérica resultante (2) recorra la totalidad de todos los $g_\Gamma$ cuando $\lambda_1,
\ldots,\lambda_r$ recorren todas las magnitudes de $\Omega$, excluidas las magnitudes de $\Omega$ que producen una reducción del grupo.\footnote{Aquí debe contarse también la aparición de una raíz doble de $g(x)=0$.} En otras palabras: el sistema no lineal de ecuaciones
\begin{equation}\tag{3}
F_1(\lambda_1,
\ldots,\lambda_r)=a_1;\quad \ldots;\quad F_n(\lambda_1,\ldots,\lambda_r)=a_n
\end{equation}
debe poseer, para cada sistema de valores $a_1,
\ldots,a_n$ que corresponda a un $g_\Gamma$, una solución, y precisamente en magnitudes $\lambda_1,
\ldots,\lambda_r$ de $\Omega$. Como el sistema (3) no es lineal, desde el comienzo debe imponerse una restricción a la exigencia de obtener la totalidad de los $g_\Gamma$ por una representación paramétrica. En efecto, en tales sistemas no lineales siempre pueden aparecer sistemas de valores singulares $a_1,
\ldots,a_n$, que satisfacen una relación algebraica $H(a_1,
\ldots,a_n)=0$, para los cuales no existe solución;\footnote{Esto se desarrollará próximamente de modo general en un trabajo sobre la ``alternativa en sistemas de ecuaciones no lineales''.} la representación paramétrica (1) falla entonces y se vuelve necesaria una representación complementaria. Sin embargo, en el caso presente esos valores singulares de $a_1,
\ldots,a_n$ se pueden caracterizar de manera simple.

\textbf{2.} Para obtener una representación paramétrica de esa clase, exigimos además que los $\lambda_i$ sean idénticamente, como funciones de las raíces consideradas como indeterminadas, irracionalidades naturales pertenecientes al grupo. Con ello entendemos lo siguiente. Si en (1) se reemplazan los parámetros $\lambda_1,
\ldots,\lambda_r$ por funciones racionales adecuadamente elegidas $\varphi_1(x),\ldots,\varphi_r(x)$ de las indeterminadas $x_1,
\ldots,x_n$ ---con coeficientes en $\Omega$---, entonces (1) debe transformarse en
\begin{equation}\tag{4}
h(x)=x^n-\sigma_1(x)x^{n-1}+\sigma_2(x)x^{n-2}\cdots \pm \sigma_n(x),
\end{equation}
donde $\sigma_1(x),\ldots,\sigma_n(x)$ significan las funciones simétricas elementales de $x_1,
\ldots,x_n$; y $h(x)$ debe tener el grupo $\Gamma$ con respecto al campo de racionalidad $\Omega(\varphi_1(x),\ldots,\varphi_r(x))$ surgido así de $\Omega(\lambda_1,
\ldots,\lambda_r)$. Además, a causa de la independencia de los parámetros, $\varphi_1(x),\ldots,\varphi_r(x)$ deben ser funciones algebraicamente independientes.

Para precisar esta exigencia hay que aclarar la conexión de $\Omega(\varphi_1(x),\ldots,\varphi_r(x))$ con el grupo $\Gamma$. Para ello determinamos primero el campo de racionalidad más general $K$ respecto del cual $h(x)$ se convierte en $h_\Gamma$. Denótese por $\Omega_\Gamma$ el cuerpo de todas las funciones racionales, con coeficientes racionales enteros, de $x_1,
\ldots,x_n$ que admiten a $\Gamma$ ---también llamado el \emph{Gattungsbereich} lagrangiano perteneciente a $\Gamma$, o el cuerpo de invariantes de $\Gamma$---; y por $\Omega(\Omega_\Gamma)$ el cuerpo correspondiente que se obtiene al adjuntar $\Omega_\Gamma$ a $\Omega$.

Así, $\Omega_\Gamma$ es un divisor del cuerpo $\mathfrak R(x_1,
\ldots,x_n)$ de todas las funciones racionales, con coeficientes racionales enteros, de $x_1,
\ldots,x_n$. Entonces, por la definición del grupo, el campo más general $K$ viene dado por todo cuerpo que contenga $\Omega_\Gamma$ como divisor y cuya intersección con $\mathfrak R(x_1,
\ldots,x_n)$ sea exactamente $\Omega_\Gamma$. De modo correspondiente, para todo campo $K$ que contenga a $\Omega$, la intersección con $\Omega(x_1,
\ldots,x_n)$ está dada por $\Omega(\Omega_\Gamma)$. En particular, como por nuestra exigencia $\Omega(\varphi_1(x),\ldots,\varphi_r(x))$ se convierte en un campo $K$, la intersección de $\Omega(\varphi_1(x),\ldots,\varphi_r(x))$ con $\Omega(x_1,
\ldots,x_n)$ está dada por $\Omega(\Omega_\Gamma)$; por otro lado, como $\Omega(\varphi_1(x),\ldots,\varphi_r(x))$ es un divisor de $\Omega(x_1,
\ldots,x_n)$, se tiene exactamente
\[
\Omega(\varphi_1(x),\ldots,\varphi_r(x))=\Omega(\Omega_\Gamma).
\]
La exigencia significa, por consiguiente, que el \emph{Gattungsbereich} lagrangiano $\Omega(\Omega_\Gamma)$ debe surgir de $\Omega$ por adjunción de las funciones algebraicamente independientes $\varphi_1(x),\ldots,\varphi_r(x)$. Aquí necesariamente debe ser $r=n$, porque $\Omega_\Gamma$ contiene las $n$ funciones elementales algebraicamente independientes, mientras que no pueden ser algebraicamente independientes más de $n$ funciones. Diremos que $\varphi_1(x),\ldots,\varphi_n(x)$ forman una \emph{base mínima} de $\Omega(\Omega_\Gamma)$; o también que $\Omega_\Gamma$ posee una base mínima con coeficientes en $\Omega$.

\textbf{3.} Es ahora esencial que toda esta consideración pueda invertirse, y que así conduzca efectivamente a condiciones suficientes\footnote{La exigencia fue ampliada en 2 respecto de 1, de modo que las condiciones no tienen por qué ser necesarias.} para la representación paramétrica buscada. Supongamos, pues, que $\varphi_1(x),\ldots,\varphi_n(x)$ es una base mínima de $\Omega(\Omega_\Gamma)$; y sea $\Omega$ el cuerpo más pequeño que contiene los coeficientes de $\varphi_1(x),\ldots,\varphi_n(x)$. Entonces $\Omega$ es necesariamente un cuerpo numérico algebraico; en particular puede ser el cuerpo $\mathfrak R$ de todos los números racionales. Como $\sigma_1(x),\ldots,\sigma_n(x)$ pertenecen a $\Omega(\Omega_\Gamma)$, existen identidades en $x_1,
\ldots,x_n$ ---donde los $F_i$ significan funciones racionales de sus argumentos---:
\begin{equation}\tag{5}
\sigma_1(x)=F_1(\varphi_1(x),\ldots,\varphi_n(x));\quad\ldots;\quad
\sigma_n(x)=F_n(\varphi_1(x),\ldots,\varphi_n(x)).
\end{equation}
Inversamente, $\varphi_1(x),\ldots,\varphi_n(x)$, y por tanto también
$\varphi(x)=\mu_1\varphi_1(x)+\cdots+\mu_n\varphi_n(x)$, dependen algebraicamente de $\sigma_1(x),\ldots,\sigma_n(x)$ mediante la ecuación irreducible relativa a $\Omega(\sigma_1(x),\ldots,\sigma_n(x))$:
\begin{equation}\tag{6}
G(\varphi(x);\sigma_1(x),\ldots,\sigma_n(x))=0,
\end{equation}
que de nuevo representa una identidad en los $x$.

En virtud de (5), esta identidad (6) pasa a
\begin{equation}\tag{7}
G\bigl(\varphi(x);F_1(\varphi_1,
\ldots,\varphi_n),\ldots,F_n(\varphi_1,
\ldots,\varphi_n)\bigr)
=H(\varphi_1(x),\ldots,\varphi_n(x))=0.
\end{equation}
Y por la independencia algebraica de $\varphi_1(x),\ldots,\varphi_n(x)$, (7) no sólo se satisface idénticamente en $x$, sino también en los $\varphi_i$; es decir, vale idénticamente en parámetros independientes $\lambda$:
\begin{equation}\tag{8}
H(\lambda_1,
\ldots,\lambda_n)=G\bigl(\mu_1\lambda_1+\cdots+\mu_n\lambda_n;
F_1(\lambda_1,
\ldots,\lambda_n),\ldots,F_n(\lambda_1,
\ldots,\lambda_n)\bigr)=0.
\end{equation}
De la identidad (8) se sigue que la ecuación
\begin{equation}\tag{9}
f(x)=x^n-F_1(\lambda_1,
\ldots,\lambda_n)x^{n-1}+F_2(\lambda_1,
\ldots,\lambda_n)x^{n-2}\cdots \pm F_n(\lambda_1,
\ldots,\lambda_n)=0
\end{equation}
tiene el grupo $\Gamma$ respecto de $\Omega(\lambda_1,
\ldots,\lambda_n)$. En efecto, esta identidad muestra que el elemento racionalmente conocido $\lambda=\mu_1\lambda_1+
\cdots+
\mu_n\lambda_n$ es una función de las raíces perteneciente a $\Gamma$; pero no puede ser racionalmente conocida ninguna otra función perteneciente a un subgrupo de $\Gamma$, como muestra la sustitución $\lambda_i=\varphi_i(x)$; mientras que, por la misma sustitución, $\lambda$ es también distinto de todos sus conjugados. Además, si $\Omega^*$ es un cuerpo numérico arbitrario que contiene a $\Omega$, entonces $f(x)$ posee el grupo $\Gamma$ respecto de $\Omega^*(\lambda_1,
\ldots,\lambda_n)$, pues por 2, al sustituir $\lambda_i=\varphi_i(x)$, el grupo no se reduce tampoco tras adjuntar $\Omega^*$.

Queda todavía demostrar que, en el sentido modificado de 1, $f(x)$ recorre efectivamente la totalidad de todos los $g_\Gamma$. Para ello observamos que, según (5), los $F_i(\lambda)$ son funciones algebraicamente independientes de sus argumentos $\lambda$. Por tanto, el sistema de ecuaciones
\begin{equation}\tag{10}
F_1(\lambda_1,
\ldots,\lambda_n)=-a_1;\quad F_2(\lambda_1,
\ldots,\lambda_n)=a_2;\quad\ldots;\quad F_n(\lambda_1,
\ldots,\lambda_n)=\pm a_n
\end{equation}
posee soluciones para cada sistema de valores arbitrario $a_1,
\ldots,a_n$, salvo los singulares que todavía se indicarán. Para todo sistema de valores $a_1,
\ldots,a_n$ que corresponda a un $g_\Gamma$ respecto de $\Omega^*$, las $\lambda_i$ correspondientes, como irracionalidades naturales pertenecientes al grupo,\footnote{Compárese también más abajo.} se convierten en magnitudes de $\Omega^*$. Si $\lambda_1,
\ldots,\lambda_n$ recorren por tanto todos los sistemas de valores, $g(x)$ recorre la totalidad de todas las ecuaciones (en el sentido modificado); dentro de ella está contenida la totalidad de los $g_\Gamma$ a los que corresponden valores de $\Omega^*$.

Recíprocamente, como por el teorema de irreducibilidad de Hilbert a los valores $\lambda$ de $\Omega^*$ les corresponden ``en general'' también ecuaciones $g_\Gamma$, se obtiene en el sentido modificado de 1 que la representación paramétrica (9) da efectivamente la totalidad de todos los $g_\Gamma$ respecto de $\Omega^*$.

Resta ahora determinar los sistemas de valores singulares para los que falla la representación paramétrica. Supóngase que, separadas en numerador y denominador, las funciones de la base mínima están dadas por
\[
\varphi_i(x)=\frac{\psi_i(x)}{\chi_i(x)}.
\]
Al sustituir esto en las identidades (5), supóngase que éstas se transforman (sin cancelación) en
\begin{equation}\tag{11}
\sigma_k(x)=\frac{G'_k(x)}{H_k(x)}\quad\text{o}\quad H_k(x)\cdot\sigma_k(x)=G_k(x),
\end{equation}
de modo que el producto de todos los $H_k(x)$ es divisible por el producto de todos los denominadores $\chi_i(x)$. Para todos los sistemas de raíces $\alpha_1,
\ldots,\alpha_n$ para los que el producto
\[
H_1(\alpha)H_2(\alpha)\cdots H_n(\alpha)\ne0
\]
y para los que, por consiguiente, tampoco se anula ningún denominador $\chi_i(\alpha)$, se puede dividir en (11) por $H_k(\alpha)$ y se obtiene
\begin{equation}\tag{12}
\sigma_k(\alpha)=\frac{G_k(\alpha)}{H_k(\alpha)}=F_k(\varphi_1(\alpha),\ldots,\varphi_n(\alpha)),
\end{equation}
donde las $\varphi_i(\alpha)$ toman valores finitos por no anularse los denominadores, y para ecuaciones $g_\Gamma$ se convierten en magnitudes de $\Omega^*$. (12) representa el sistema de soluciones de (10) tan pronto como se determinan $\alpha_1,
\ldots,\alpha_n$ de tal modo que $a_k=(-1)^k\sigma_k(\alpha)$.\footnote{La resolución de la ecuación queda así reducida a la resolución del sistema formado por funciones independientes:
$\varphi_1(\alpha_1,
\ldots,\alpha_n)=\lambda_1;\ldots;\varphi_n(\alpha_1,
\ldots,\alpha_n)=\lambda_n$.}

Singulares sólo pueden llegar a ser, por tanto, aquellos sistemas de coeficientes $a_k=(-1)^k\sigma_k(\alpha)$ para los cuales
\[
H_1(\alpha)H_2(\alpha)\cdots H_n(\alpha)=0;
\]
es decir, aquellos para los que (11) pasa a la representación $0=0$. Entonces se requiere una investigación especial para decidir si entre los $a_1,
\ldots,a_n$ singulares así definidos se encuentran también algunos a los que correspondan ecuaciones $g_\Gamma$, y para qué cuerpos numéricos.\footnote{Por ejemplo, como ha mostrado F. Seidelmann loc. cit., sólo para el grupo alternante de las ecuaciones de grados 3 y 4 aparecen tales ecuaciones ---correspondientes a sistemas de valores singulares---, y sólo para aquellos cuerpos numéricos $\Omega^*$ que contienen el cuerpo de las terceras raíces de la unidad. En cada caso puede darse aquí una representación complementaria simple.} En resumen se obtiene el siguiente

\smallskip
\noindent\textbf{Teorema.} \emph{Si el \emph{Gattungsbereich} lagrangiano correspondiente al grupo $\Gamma$ posee una base mínima\footnote{De la existencia de una base mínima se sigue inmediatamente la existencia de arbitrariamente muchas, que pueden deducirse de la original mediante transformaciones racionales invertibles.} con coeficientes en $\Omega$, entonces para todo cuerpo numérico $\Omega^*$ que contenga a $\Omega$ la totalidad de las ecuaciones con grupo prescrito $\Gamma$ respecto de $\Omega^*$ puede construirse racionalmente mediante la representación paramétrica (2), eventualmente con excepción de las singulares respecto de la representación paramétrica, que pueden darse separadamente. Los parámetros deben recorrer todos los valores de $\Omega^*$ que no produzcan una reducción del grupo. La construcción es posible para cualquier cuerpo numérico si el campo de coeficientes $\Omega$ de la base mínima es igual al cuerpo $\mathfrak R$ de todos los números racionales.}

Además, por el teorema de irreducibilidad de Hilbert, al excluir de la variedad paramétrica aquellos valores de $\Omega^*$ que producen una reducción del grupo no se reduce la variedad. Por tanto, vale todavía: \emph{de la existencia de una base mínima se sigue que la variedad de las ecuaciones con el grupo $\Gamma$ respecto de $\Omega^*$ es igual a la variedad de las ecuaciones sin afecto; la variedad no queda reducida por el afecto.}

\subsection*{\S 2. Aplicación a ecuaciones especiales.}

Demostramos ahora la existencia de la base mínima para todos los grupos que aparecen en las ecuaciones de grados 3 y 4; para ello mostramos primero, de manera completamente general, que para ecuaciones de grado $n$ el problema se reduce a uno de $n-2$ indeterminadas. La primera reducción corresponde a la transformación lineal de Tschirnhaus que elimina el segundo término de mayor grado. Tomamos como primera función de base $\varphi_1(x)=\sigma_1(x)$, y observamos además que $\Omega_\Gamma$ se puede derivar racionalmente de todas las formas homogéneas de $x_1,
\ldots,x_n$ que admiten a $\Gamma$. Sea $p(x_1,
\ldots,x_n)$ una forma de ese tipo; entonces también pertenece a $\Omega_\Gamma$
\[
q(x)=p\left(x_1-\frac{\sigma_1(x)}{n},\ldots,x_n-\frac{\sigma_1(x)}{n}\right)=p\left(x_i-\frac{\sigma_1(x)}{n}\right),
\]
y además se tiene
\[
p(x)\equiv p\left(x_i-\frac{\sigma_1(x)}{n}\right)\pmod{\sigma_1(x)},
\quad\text{o bien}\quad
p(x)=q(x)+\sigma_1(x)\cdot p_1(x),
\]
con $p_1(x)$ perteneciente a $\Omega_\Gamma$; éste es además de menor dimensión que $p(x)$. La continuación del procedimiento muestra, por consiguiente, que todas las formas homogéneas de $\Omega_\Gamma$ ---y, por tanto, en general todas las funciones racionales de $\Omega_\Gamma$--- pueden expresarse racionalmente mediante $\varphi_1(x)=\sigma_1(x)$ y mediante formas homogéneas
\[
q(x)=p\left(x_i-\frac{\sigma_1(x)}{n}\right)=p(y_i).
\]
Entre esas combinaciones $y_i$ existe, sin embargo, la relación lineal $\sum y_i=0$; así, las $q(x)$ dependen sólo de $(n-1)$ argumentos.

La segunda reducción se basa en que las $q(x)$ son formas homogéneas de sus argumentos. Tomamos como segunda función de base
\[
\varphi_2(x)=\frac{\tau_3(x)}{\tau_2(x)},
\]
donde
\[
\tau_3(x)=\sigma_3\left(x_i-\frac{\sigma_1(x)}{n}\right)=\sigma_3(y_i),
\qquad
\tau_2(x)=\sigma_2\left(x_i-\frac{\sigma_1(x)}{n}\right)=\sigma_2(y_i).
\]
Así, $\varphi_2(x)$ es una función homogénea-fraccionaria de primer grado de los $(n-1)$ argumentos, a saber, de los $y_i$ con $\sum y_i=0$. Por medio de $\varphi_2(x)$, todas las $q(x)$ se pueden expresar racionalmente mediante las funciones, también pertenecientes a $\Omega_\Gamma$,
\[
r(x)=\frac{q(x)}{\varphi_2(x)^\lambda}
=\frac{q(x)\tau_2(x)^\lambda}{\tau_3(x)^\lambda}
=\frac{p(y_i)\sigma_2(y_i)^\lambda}{\sigma_3(y_i)^\lambda},
\]
donde $\lambda$ designa el grado de $q(x)$. Entonces $r(x)$ es homogénea de grado cero y depende, por tanto, sólo de $(n-2)$ argumentos, las razones $y_1:y_2:\cdots:y_n$ con $\sum y_i=0$. Así, $\Omega_\Gamma$ se puede expresar racionalmente mediante
\[
\varphi_1(x)=\sigma_1(x),
\qquad
\varphi_2(x)=\frac{\tau_3(x)}{\tau_2(x)},
\]
y el cuerpo de todas las funciones $r(x)$ de $\Omega_\Gamma$ que dependen de esos $(n-2)$ argumentos,\footnote{Para la ecuación de grado $n$ sin segundo término
\[
f(x)=x^n+a_2x^{n-2}+a_3x^{n-3}+\cdots+a_n=0
\]
corresponde a esta segunda reducción la sustitución, siempre posible para $a_2\ne0$, $a_3\ne0$,
\[
y=\frac{x a_2}{a_3},
\]
por la cual $f(x)$, salvo un factor, pasa a
\[
g(y)=y^n+\frac{a_2^3}{a_3^2}y^{n-2}+\frac{a_2^3}{a_3^2}y^{n-3}
+\frac{a_4a_2^4}{a_3^4}y^{n-4}+\cdots+a_n\frac{a_2^n}{a_3^n}=0,
\]
es decir, a una ecuación que contiene sólo $n-2$ parámetros:
\[
g(y)=y^n+c\,y^{n-2}+c\,y^{n-3}+c_4y^{n-4}+\cdots+c_n=0.
\]} con lo cual queda realizada la reducción deseada.

Para ecuaciones de grados 3 y 4 se tiene por ello, en particular, una reducción a cuerpos de funciones de uno, respectivamente dos, argumentos. Para éstos siempre existe una base mínima, por los teoremas de Lüroth y Castelnuovo.\footnote{J. Lüroth: ``Beweis eines Satzes über rationale Kurven'', Math. Ann. 9 (1875); G. Castelnuovo: ``Sulla razionalità delle involuzioni piane'', Math. Ann. 44 (1893). F. Enriques mostró, mediante un contraejemplo, que para cuerpos de más de dos indeterminadas no tiene por qué existir una base mínima: Rend. Acc. Linc., vol. 21, 21 de enero de 1912.} En el caso de una indeterminada, la base mínima puede deducirse mediante operaciones racionales, y por tanto contiene, en particular para $\Omega_\Gamma$, sólo coeficientes numéricos racionales. En dos indeterminadas podría tratarse todavía, según Castelnuovo, de una base con coeficientes en un cuerpo numérico algebraico $\Omega$; la investigación efectiva (véase F. Seidelmann loc. cit.) muestra que también aquí se obtiene para cada $\Omega_\Gamma$ una base con coeficientes racionales. Esto se ve casi sin cálculo ya por el hecho de que para el grupo de Klein de orden cuatro puede elegirse una base con coeficientes racionales, a saber
\[
\varphi_1(x)=\sigma_1(x);\quad
\varphi_2(x)=\frac{vw}{u};\quad
\varphi_3(x)=\frac{uw}{v};\quad
\varphi_4(x)=\frac{uv}{w},
\]
donde
\[
u=x_1+x_2-x_3-x_4;\quad v=x_1-x_2+x_3-x_4;\quad w=x_1-x_2-x_3+x_4,
\]
y porque el grupo alternante pasa al cíclico de las tres funciones de base $\varphi_2,\varphi_3,\varphi_4$, y cada grupo octavo pasa al simétrico de dos de esas funciones, con lo cual se efectúa una reducción a los grupos de las ecuaciones de grados 3 y 2. Para el grupo cíclico, una base con coeficientes racionales ya era conocida por Abel, aunque no formulada como tal.\footnote{Véase Weber, \emph{Algebra} (edición pequeña), \S 93, fórmula (12).} Así queda demostrado: \emph{la totalidad de las ecuaciones de grados 3 y 4 con grupo prescrito ---eventualmente con excepción de las singulares respecto de la representación paramétrica--- puede construirse racionalmente por representación paramétrica para cualquier campo de racionalidad; su variedad es igual a la variedad de las ecuaciones sin afecto.}

La investigación efectiva muestra (véase F. Seidelmann loc. cit.) que sólo para el grupo alternante, respecto de campos de racionalidad que contienen el cuerpo de las terceras raíces de la unidad, se necesita una representación complementaria, mientras que en todos los demás casos la representación paramétrica proporciona la totalidad sin restricción.

Obsérvese aún que la base mínima es conocida también para todos los grupos abelianos.\footnote{E. Fischer: ``Die Isomorphie der Invariantenkörper der endlichen Abelschen Gruppen linearer Transformationen'', Gött. Nachr. 1915, y ``Zur Theorie der endlichen Abelschen Gruppen'', Math. Ann. 77 (1915).} Pero aquí se trata de una base con coeficientes en el cuerpo ciclotómico derivado de los caracteres del grupo. Por tanto, en este caso no se obtienen nuevos resultados para la construcción de las ecuaciones con grupo prescrito.

\medskip
\noindent Gotinga, julio de 1916.

\clearpage

\begin{center}
{\Large\bfseries 12. Invariantes de expresiones diferenciales arbitrarias}\par
\vspace{1.5em}
Nachr. v. d. Ges. d. Wiss. zu Göttingen 1918, pp. 37--44
\end{center}

\vspace{3em}
\begin{center}
Presentado en la sesión del 25 de enero de 1918 por F. Klein.
\end{center}

Por una expresión diferencial entiendo una función
\[
f(x,dx)=f(x_1,\ldots,x_n;\,dx_1,\ldots,dx_n)
\]
de las \(n\) variables \(x_1,\ldots,x_n\) y de sus diferenciales
\(dx_1,\ldots,dx_n\), que se supone analítica en los dos sistemas de
\(n\) argumentos, pero no necesariamente real. Tomo ahora como base el grupo
de todas las transformaciones analíticas de las variables, al cual corresponde
el grupo de todas las transformaciones lineales de los diferenciales:
\begin{equation}
x_i=x_i(y_1,\ldots,y_n);\qquad
dx_i=\sum_k\frac{\partial x_i}{\partial y_k}\,dy_k;\qquad
\delta x_i=\sum_k\frac{\partial x_i}{\partial y_k}\,\delta y_k,
\tag{1}
\end{equation}
por las cuales \(f(x,dx)\) debe pasar a \(g(y,dy)\). Por una invariante de
\(f(x,dx)\) entiendo una invariante con respecto a este grupo y al grupo que
éste induce para los diferenciales superiores \(d^2x,d\delta x,\ldots\) y para
las derivadas de \(f(x,dx)\); es decir, una función \(J\) (analítica) tal que,
para \(f\) arbitraria y en virtud de (1), se tenga identicamente en las
variables y en todos los diferenciales que aparezcan:
\[
\begin{aligned}
&J\!\left(f,\frac{\partial f}{\partial dx}\cdots
\frac{\partial^{\rho+\sigma}f}{\partial x^\rho\partial dx^\sigma}
\cdots dx,\delta x,d^2x,\ldots\right)\\
&\qquad =
J\!\left(g,\frac{\partial g}{\partial dy}\cdots
\frac{\partial^{\rho+\sigma}g}{\partial y^\rho\partial dy^\sigma}
\cdots dy,\delta y,d^2y,\ldots\right).\footnotemark
\end{aligned}
\]
\footnotetext{\(\displaystyle
\frac{\partial^{\rho+\sigma}f}{\partial x^\rho\partial dx^\sigma}\)
se usará en todo momento como abreviatura de
\[
\frac{\partial^{\rho+\sigma}f}
{\partial x_{i_1}\cdots\partial x_{i_\rho}
 \partial dx_{k_1}\cdots\partial dx_{k_\sigma}},
\]
donde los índices significan cualesquiera de los números de \(1\) a \(n\).}
De manera enteramente análoga se define la invariante de un sistema simultáneo
de expresiones diferenciales. Si, en particular, una invariante de este tipo
contiene sólo primeros diferenciales \(dx,\delta x\), y sólo derivadas con
respecto a esos diferenciales y no con respecto a las variables, entonces no
entra en consideración la aparición de las variables en las funciones ni la
transformación lineal de los diferenciales; esas invariantes especiales son
invariantes respecto del grupo de todas las transformaciones lineales de los
diferenciales \(dx,\delta x\), con coeficientes indeterminados, y se llamarán
invariantes proyectivas del sistema simultáneo.

Para formas diferenciales cuadráticas homogéneas, Christoffel (Crelle 70) y
Ricci (Math. Annalen 54) han establecido un sistema de formaciones invariantes
tal que todas las invariantes se convierten en invariantes proyectivas de ese
sistema simultáneo; con ello las cuestiones sobre la totalidad de las invariantes
y sobre la equivalencia quedan reducidas a cuestiones de la teoría lineal de
invariantes. A esto quisiera llamarlo el ``teorema de reducción''. El sistema
completo consta de infinitas formas numerables; pero, para invariantes de
cualquier orden finito \(\rho\), es decir, aquellas que no contienen derivadas
respecto de las variables de orden superior al \(\rho\)-ésimo, consta sólo de
un número finito de formas.

En lo que sigue demuestro la validez del ``teorema de reducción'' para
expresiones diferenciales arbitrarias; de nuevo se trata de un sistema completo
de formaciones invariantes infinitas numerables, pero sólo de un número finito
para cada orden finito. Sin embargo, mientras Christoffel y Ricci se apoyan,
para generar las invariantes y demostrar el teorema de reducción, en
eliminaciones difíciles de abarcar, yo genero las invariantes directamente por
procesos invariantes de diferenciación y de variación, estos últimos
interpretables sencillamente como procesos de diferenciación respecto de otros
parámetros. El algoritmo de estos procesos invariantes de diferenciación fue
desarrollado, siguiendo a Lagrange (\emph{Mécanique analytique}, parte 2,
sección IV), por Riemann (Obras XXII) y Lipschitz (Crelle 70, 72) para formas
diferenciales cuadráticas, aunque sin avanzar hasta el teorema de reducción. El
medio auxiliar para probar el teorema de reducción lo proporcionan las
``coordenadas normales'' introducidas por Riemann para las formas cuadráticas
(Obras XIII), que transforman en rectas las extremales que salen de un punto en
el problema variacional correspondiente, pero que existen sólo para expresiones
diferenciales homogéneas \(f(x,\varkappa\cdot dx)=\Phi(\varkappa)\cdot
f(x,dx)\footnote{Se ve fácilmente que \(\Phi(\varkappa)=\varkappa^c\), donde
\(c\) denota una cantidad real o compleja cualquiera.}\). El caso inhomogéneo,
sin embargo, se reduce a éste.

\begin{center}
\textbf{I. Generación de invariantes.}
\end{center}

Una sucesión, por lo general no terminante, de invariantes proyectivas de
\(f(x,dx)\) viene dada por las polares. Por la linealidad de la transformación
de los diferenciales, de \(f(x,dx)=g(y,dy)\) se sigue también
\(f(x,dx+\lambda\delta x)=g(y,dy+\lambda\delta y)\); de aquí, al desarrollar
en \(\lambda\), resultan las relaciones características de la invariancia:
\begin{equation}
\begin{gathered}
f(dx)=g(dy),\\
f_\delta=\sum\frac{\partial f(dx)}{\partial dx}\,\delta x
       =\sum\frac{\partial g(dy)}{\partial dy}\,\delta y,\\
\vdots\\
f_{\delta^\sigma}=
\sum\frac{\partial^\sigma f(dx)}{\partial dx^\sigma}\,\delta x^\sigma
=
\sum\frac{\partial^\sigma g(dy)}{\partial dy^\sigma}\,\delta y^\sigma,\\
\vdots
\end{gathered}
\tag{2}
\end{equation}

Cada polar, mediante la aplicación de un proceso de variación \(\bdelta\), da
lugar a una sucesión no terminante de invariantes generales:
\begin{equation}
\begin{gathered}
\bdelta^\rho f(x,dx)=\bdelta^\rho g(y,dy),\\
\vdots\\
\bdelta^\rho f_{\delta^\sigma}
=\bdelta^\rho\sum\frac{\partial^\sigma f(dx)}{\partial dx^\sigma}\,
\delta x^\sigma
=\bdelta^\rho\sum\frac{\partial^\sigma g(dy)}{\partial dy^\sigma}\,
\delta y^\sigma,\\
\vdots\\[0.2em]
(\rho=0,1,2,\ldots).
\end{gathered}
\tag{3}
\end{equation}
Las relaciones (2) y (3) permiten leer, por medio de (1), las leyes de
transformación de las derivadas de \(f(x,dx)\); esto, sin embargo, no se usará
en lo que sigue. A partir de las invariantes (3) se puede obtener por
combinación lineal la ``forma normal de la \(\rho\)-ésima variación'', que para
\(\rho=1\) se encuentra ya en Lagrange y para \(\rho=2\) en Riemann, y que se
caracteriza porque los diferenciales mixtos de orden \(\rho+1\), a saber
\(d^\rho\delta,d^{\rho-1}\delta^2,\ldots\), quedan eliminados. Para ella se
obtiene:
\begin{equation}
\begin{gathered}
\Omega_1=\delta f(dx)-d f_\delta(dx),\\
\Omega_2=\delta^2 f-\delta d f_\delta+\frac12 d^2 f_{\delta^2},\\
\vdots\\
\Omega_\rho=\delta^\rho f-\delta^{\rho-1}d f_\delta
+\frac12\delta^{\rho-2}d^2 f_{\delta^2}
+\cdots+\frac{(-1)^\rho}{\rho!}\,d^\rho f_{\delta^\rho},\\
\vdots
\end{gathered}
\tag{4}
\end{equation}

A partir de las invariantes (4), mediante la generalización de un planteamiento
de Riemann, se llega ahora a invariantes que contienen sólo primeros
diferenciales; tales invariantes se denominarán ``funciones fundamentales''. En
efecto, si en \(\Omega_1\) se reemplazan los diferenciales por los cocientes
diferenciales, entonces \(\Omega_1=0\), idénticamente en \(\delta x\), da
precisamente las ecuaciones de Lagrange del problema variacional correspondiente
a
\[
f\!\left(\frac{dx}{dt}\right).
\]
Impongo ahora la condición invariante de que la dirección \((d+\lambda\delta)\)
satisfaga esas ecuaciones de Lagrange:
\begin{equation}
\Omega_1(d+\lambda\delta)=
\bdelta f(dx+\lambda\delta x)
-(d+\lambda\delta)f_{\bdelta}(dx+\lambda\delta x)=0
\quad(\text{idénticamente en }\delta x),
\tag{5}
\end{equation}
de donde, por la sustitución \(\bdelta=d+\lambda\delta\), para \(f(dx)\)
homogénea se sigue además
\begin{equation}
(d+\lambda\delta)f(dx+\lambda\delta x)=0,
\tag{6}
\end{equation}
con excepción del caso de homogeneidad de primer orden. En este último caso,
(6) debe agregarse como nueva condición, pues entonces el sistema de ecuaciones
(5) representa sólo \(n-1\) condiciones independientes. Si en (5),
respectivamente en (5) y (6), se igualan a cero los coeficientes de la potencia
cero, primera y segunda de \(\lambda\), se obtienen tres sistemas invariantes de
ecuaciones \(\varphi=0\), \(\psi=0\), \(\chi=0\), idénticamente en \(\delta x\),
por medio de los cuales \(d^2x,d\delta x,\delta^2x\) pueden expresarse mediante
primeros diferenciales. Los sistemas invariantes de ecuaciones ulteriores
\begin{equation}
\begin{gathered}
d^\sigma\varphi=0,\qquad d^\sigma\psi=0,\qquad
d^\sigma\chi=0,\qquad d^{\sigma-1}\delta\chi=0,\\
d^{\sigma-2}\delta^2\chi=0,\ldots,\delta^\sigma\chi=0
\qquad(\sigma=0,1,2,\ldots)
\end{gathered}
\tag{7}
\end{equation}
bastan entonces exactamente para expresar todos los diferenciales superiores
mediante los primeros.\footnote{Sólo la forma lineal \(\sum A_i\,dx_i\) requiere
un tratamiento especial, pues aquí el sistema de ecuaciones (5) no contiene
segundos diferenciales.} Las funciones (4), junto con el sistema invariante de
ecuaciones (7), conducen por tanto a funciones fundamentales \([\Omega_\rho]\)
que contienen sólo primeros diferenciales; además \([\Omega_1]\) desaparece
idénticamente.

Del mismo modo, a partir del diferencial \(\delta h(dx,\delta x)\) de una
función fundamental \(h\), puede obtenerse de nuevo, mediante (7), una función
fundamental, que se llamará la ``derivada covariante'' \([h^{(1)}]\) de \(h\);
en general, \([h^{(\nu)}]\) denotará la derivada covariante de
\([h^{(\nu-1)}]\). Estos dos procesos conducen así a una sucesión doblemente
infinita de funciones fundamentales:
\begin{equation}
\begin{array}{cccccc}
[\Omega_2], & [\Omega_2^{(1)}], & [\Omega_2^{(2)}], & \ldots,
& [\Omega_2^{(\nu)}], & \ldots\\
\vdots &&&& \vdots&\\
{}[\Omega_\rho], & [\Omega_\rho^{(1)}], & \ldots, & \ldots,
& [\Omega_\rho^{(\nu)}], & \ldots\\
\vdots & \vdots &&& \vdots&
\end{array}
\tag{8}
\end{equation}

Se muestra además que los miembros izquierdos de las ecuaciones que expresan
los diferenciales superiores mediante los primeros son cogredientes con los
primeros diferenciales, es decir, están sometidos a las mismas transformaciones
lineales. Si, por tanto, en lugar de los diferenciales superiores se introducen
esos miembros izquierdos \(p,q,r,\ldots\) como argumentos de la invariante,
entonces ésta, aparte de las derivadas, depende sólo de magnitudes cogredientes
entre sí. La formación de las funciones (8) indicada arriba consiste simplemente
en introducir estas nuevas magnitudes; toda invariante pasa con ello a una suma
de invariantes, y de esa suma se toma precisamente la que contiene
\(dx,\delta x\), pero no \(p,q,r,\ldots\).

A continuación se demuestra que, bajo la hipótesis de homogeneidad
\[
f(\varkappa\cdot dx)=\Phi(\varkappa)\cdot f(dx),
\]
las funciones fundamentales (8) agotan el sistema completo buscado.

\begin{center}
\textbf{II. Coordenadas normales, equivalencia y teorema de reducción.}
\end{center}

Llego a las coordenadas normales integrando el problema variacional asociado a
\(f\!\left(\frac{dx}{dt}\right)\):
\begin{equation}
\begin{gathered}
\Omega_1=\delta f\!\left(\frac{dx}{dt}\right)
-\frac{d}{dt}f_\delta\!\left(\frac{dx}{dt}\right)=0
\quad(\text{idénticamente en }\delta x),\\
\frac{d}{dt}f\!\left(\frac{dx}{dt}\right)=0
\quad(\text{como consecuencia o condición adicional}).
\end{gathered}
\tag{9}
\end{equation}
Por la homogeneidad supuesta de \(f(dx)\), este sistema de ecuaciones se
transforma en sí mismo bajo la sustitución de parámetro \(t=\varkappa\cdot\tau\).
Así, si mediante (9) se expresa \(\frac{d^\rho x_i}{dt^\rho}\) como función
\(\varphi_\rho^{(i)}\!\left(\frac{dx}{dt}\right)\) de las primeras derivadas,
entonces \(\varphi_\rho^{(i)}\) es homogénea de orden \(\rho\):
\begin{equation}
\varphi_\rho^{(i)}\!\left(\varkappa\cdot\frac{dx}{dt}\right)
=\varkappa^\rho\varphi_\rho^{(i)}\!\left(\frac{dx}{dt}\right).
\tag{10}
\end{equation}
Si ahora, al integrar (9), se prescriben para \(t=t_0\) los valores iniciales
\[
x=(x)_0,\qquad \frac{dx}{dt}=\left(\frac{dx}{dt}\right)_0
\]
y se pone
\begin{equation}
u_i=(t-t_0)\left(\frac{dx_i}{dt}\right)_0,
\tag{11}
\end{equation}
donde, en virtud de \(f\!\left(\frac{dx}{dt}\right)=\mathrm{const.}\), el
parámetro \((t-t_0)\) se vuelve proporcional a la ``longitud de la extremal'',
entonces el desarrollo en potencias de \((t-t_0)\), por (10), da
\begin{equation}
x_i-(x_i)_0=u_i+\frac12\varphi_2^{(i)}(u)+\cdots+
\frac{1}{\rho!}\varphi_\rho^{(i)}(u)+\cdots .
\tag{12}
\end{equation}
Este desarrollo puede considerarse como una transformación de variables
\(x_i=x_i(u)\), donde las cantidades \(u\) son precisamente las coordenadas
normales de Riemann. Según (11) y (12), estas coordenadas transforman en rectas
las extremales que pasan por \((x)_0\). Además, a una transformación arbitraria
de variables \(x=x(y)\) corresponde, por (11), una transformación lineal
\(u=u(v)\) de las coordenadas normales asociadas, cuyos coeficientes dependen
sólo del sistema arbitrario de valores \((x)_0\); por tanto, los diferenciales
\(du,\delta u\) están sometidos a la misma transformación que las propias
\(u\).

Por medio de (12), \(f(x,dx)\) pasa a \(F(u,du)\); o bien, desarrollada en
potencias de las \(u\),
\begin{equation}
F(u,du)=F_0(u,du)+F_1(u,du)+\cdots+F_\rho(u,du)+\cdots,
\tag{13}
\end{equation}
donde \(F_\rho(u,du)\) designa una forma homogénea de dimensión \(\rho\) en
\(u\). Por la linealidad de la transformación \(u=u(v)\), en virtud de
\(F(u,du)=G(v,dv)\), también los componentes homogéneos individuales pasan a
los correspondientes:
\begin{equation}
F_\rho(u,du)=G_\rho(v,dv).
\tag{14}
\end{equation}
Si se toma, por tanto, para \((x)_0\) cualquier sistema fijo y constante de
valores, (13) y (14) muestran que la equivalencia de \(f(dx)\) con \(g(dy)\)
equivale a la equivalencia, respecto de una transformación lineal, de los
sistemas de funciones correspondientes \(F_\rho(u,du)\). Además, las
\(F_\rho(u,du)\), o las correspondientes \(F_\rho(\delta u,du)\), representan
invariantes de \(f(dx)\), tomadas en el punto \(x=(x)_0\); de hecho, forman un
sistema completo de funciones fundamentales para ese punto \(x=(x)_0\). En
efecto, se tiene
\[
F_\rho(\delta u,du)=\frac1{\rho!}\sum
\frac{\partial F_\rho}{\partial \delta u^\rho}\,\delta u^\rho;
\qquad
\frac{\partial F_\rho}{\partial \delta u^\rho}
=\left(\frac{\partial F(du)}{\partial u^\rho}\right)_0,
\]
y esta última derivada puede expresarse, mediante (12), exactamente de la misma
manera por derivadas de \(f(dx)\), tomadas para \(x=(x)_0\), que la derivada
correspondiente formada a partir de \(G(dv)\) puede expresarse por derivadas de
\(g(dy)\). Puesto que
\[
\left(\frac{\partial^{\rho+\sigma}F(du)}
{\partial u^\rho\partial du^\sigma}\right)_0
=
\frac{\partial^{\rho+\sigma}F_\rho(\delta u,du)}
{\partial\delta u^\rho\partial du^\sigma},
\]
toda invariante de \(f(dx)=F(du)\), tomada en el punto \(x=(x)_0\), se convierte
en una invariante proyectiva de las \(F_\rho(\delta u,du)\) tan pronto como se
reemplazan los diferenciales superiores por los \(p,q,r,\ldots\) cogredientes
con \(dx,\delta x\). Pero ahora \((x)_0\) puede considerarse como parámetro; a
cada invariante en el punto \(x=(x)_0\) corresponde una invariante de \(f(dx)\)
si se reemplaza de nuevo \((x)_0\) por \(x\). Así, si las
\(F_\rho(\delta u,du)\) se convierten en invariantes
\(\Psi_\rho(\delta x,dx)\), tomadas para \(x=(x)_0\) ---pues para \(x=(x)_0\)
\(du,\delta u\) pasan a \(dx,\delta x\)---, las propias \(\Psi_\rho\) forman las
funciones fundamentales de un sistema completo. Queda sólo expresar este sistema
de funciones fundamentales mediante invariantes explícitas de \(f(dx)\), a saber,
mediante las funciones (8). Que esto es posible lo muestra su generación por
procesos de diferenciación. En los términos de orden más alto estos procesos
pasan, en efecto, a procesos polares simples, y una transformación análoga al
desarrollo en serie de Clebsch--Gordan, unida a un razonamiento inductivo por
los términos despreciados, prueba la afirmación. Si se trata de invariantes de
un sistema simultáneo, a las funciones fundamentales (8) deben añadirse todavía
las derivadas covariantes de las demás expresiones diferenciales, como muestra
la introducción, en esas expresiones, de las coordenadas normales pertenecientes
a \(f(dx)\).

La equivalencia de \(f(dx)\) con \(g(dy)\), que había sido reducida a (14),
queda por tanto también equivalente a la equivalencia de las \(\Psi_\rho\), o
igualmente de las funciones (8), respecto de la transformación lineal de los
diferenciales, sin que se requieran condiciones especiales de integrabilidad;
esto se sigue de la posibilidad de reducción a (14). Con ello el teorema de
reducción queda probado en todas sus partes. En particular, la anulación
idéntica de la sucesión de funciones \([\Omega_2],[\Omega_3],\ldots,
[\Omega_\rho],\ldots\) es necesaria y suficiente para que \(f(dx)\) pueda
transformarse en una expresión con coeficientes constantes.

Finalmente, si en lugar del grupo de todas las transformaciones analíticas se
toma como base un subgrupo, entonces también el grupo de las transformaciones
lineales correspondientes de las \(u\) pasa a ser un subgrupo del grupo
proyectivo; las invariantes de \(f(dx)\) vuelven a ser invariantes de las
funciones (8) respecto de transformación lineal, pero ahora respecto de este
subgrupo. Así, por homogenización, el caso de funciones no homogéneas \(f(dx)\)
se reduce al grupo afín; aquí el sistema completo puede deducirse de las
funciones (8) formadas para una variable adicional.

Del teorema de reducción se obtiene el caso especial de las formas cuadráticas,
y más generalmente de las formas de dimensión \(p\), por el hecho de que aquí el
sistema de funciones se interrumpe con \([\Omega_2]\), más generalmente con
\([\Omega_p]\) y sus derivadas covariantes, puesto que todas las funciones
fundamentales posteriores se convierten en invariantes proyectivas de éstas.
Así se explica la posición distinguida de la forma de curvatura riemanniana
\([\Omega_2]\) para las formas cuadráticas. La cuestión de la equivalencia de
formas cuadráticas, respectivamente de formas de dimensión \(p\), se reduce
precisamente a la equivalencia de \([\Omega_2]\), respectivamente de
\([\Omega_2]\ldots[\Omega_p]\), y de sus derivadas covariantes, respecto de una
transformación lineal; y la anulación idéntica de \([\Omega_2]\), respectivamente
de \([\Omega_2]\ldots[\Omega_p]\), es necesaria y suficiente para la
transformación en formas con coeficientes constantes, con lo cual se enuncia,
para las formas cuadráticas, uno de los resultados más conocidos.

Una exposición más detallada aparecerá próximamente en los \emph{Mathematische Annalen}.

\clearpage

\begin{center}
{\Large\bfseries 13. Problemas variacionales invariantes}\par
\vspace{1.5em}
\emph{Nachrichten von der Gesellschaft der Wissenschaften zu Göttingen} 1918, pp. 235--257
\end{center}

\vspace{3em}
\begin{center}
Presentado por F. Klein en la sesión del 26 de julio de 1918.\footnote{La versión definitiva del manuscrito no fue entregada sino hasta fines de septiembre.}
\end{center}

Se trata de problemas variacionales que admiten un grupo continuo (en el sentido de Lie); las consecuencias que de ello resultan para las ecuaciones diferenciales correspondientes encuentran su expresión más general en los teoremas formulados en el \S{}1 y demostrados en los parágrafos siguientes. Para estas ecuaciones diferenciales que nacen de problemas variacionales pueden darse afirmaciones mucho más precisas que para ecuaciones diferenciales arbitrarias que admiten un grupo, las cuales forman el objeto de las investigaciones de Lie. Lo que sigue se apoya, por tanto, en una combinación de los métodos del cálculo variacional formal con los de la teoría de grupos de Lie. Para grupos y problemas variacionales especiales esta combinación de métodos no es nueva; menciono a Hamel y Herglotz para grupos finitos especiales, y a Lorentz y sus discípulos (por ejemplo Fokker), Weyl y Klein para grupos infinitos especiales.\footnote{Hamel: \emph{Math. Ann.} vol. 59 y \emph{Zeitschrift f. Math. u. Phys.} vol. 50. Herglotz: \emph{Ann. d. Phys.} (4) vol. 36, especialmente \S{}9, p. 511. Fokker, Verslag d. Amsterdamer Akad., 27 ene. 1917. Para la bibliografía ulterior compárese la segunda nota de Klein: Göttinger Nachrichten, 19 de julio de 1918. En un trabajo recién aparecido de Kneser (\emph{Math. Zeitschrift} vol. 2), se trata de la construcción de invariantes por un método semejante.} En particular, la segunda nota de Klein y las presentes exposiciones se han influido mutuamente; a este respecto puedo remitir a las observaciones finales de la nota de Klein.

\section*{\S{}1. Observaciones preliminares y formulación de los teoremas}

Todas las funciones que aparecen en lo que sigue deben suponerse analíticas, o al menos continuas y finitamente muchas veces continuamente diferenciables, y unívocas en el dominio considerado.

Por ``grupo de transformaciones'' se entiende, como es usual, un sistema de transformaciones tal que para cada transformación existe una transformación inversa contenida en el sistema, y tal que la composición de cualesquiera dos transformaciones del sistema pertenece de nuevo al sistema. El grupo se llama una \(\G_\rho\) continua finita cuando sus transformaciones están contenidas en una transformación más general que depende analíticamente de \(\rho\) parámetros esenciales \(\eps\) (es decir, los \(\rho\) parámetros no han de poder representarse como \(\rho\) funciones de un número menor de parámetros). Correspondientemente, por una \(\G_{\infty\rho}\) continua infinita se entiende un grupo cuyas transformaciones más generales dependen de \(\rho\) funciones arbitrarias esenciales \(p(x)\) y de sus derivadas, analíticamente o al menos de modo continuo y con un número finito de derivadas continuas. Entre ambos casos se encuentra, como caso intermedio, el grupo que depende de infinitos parámetros pero no de funciones arbitrarias. Finalmente, se denomina grupo mixto a uno que depende tanto de funciones arbitrarias como de parámetros.\footnote{Lie define, en \emph{Grundlagen für die Theorie der unendlichen kontinuierlichen Transformationsgruppen} (Berichte der Königl. Sächs. Gesellschaft der Wissenschaften, 1891), el grupo continuo infinito como un grupo de transformaciones cuyas transformaciones están dadas por las soluciones más generales de un sistema de ecuaciones diferenciales parciales, siempre que estas soluciones no dependan sólo de un número finito de parámetros. Así se obtiene uno de los tipos indicados arriba, distinto del grupo finito; inversamente, sin embargo, el caso límite de infinitos parámetros no necesita satisfacer necesariamente un sistema de ecuaciones diferenciales.}

Sean \(x_1,\ldots,x_n\) variables independientes, y sean \(u_1(x),\ldots,u_\mu(x)\) funciones que dependen de ellas. Si se someten las \(x\) y las \(u\) a las transformaciones de un grupo, entonces, por la invertibilidad supuesta de las transformaciones, entre las magnitudes transformadas deben aparecer de nuevo exactamente \(n\) independientes, denotadas \(y_1,\ldots,y_n\); las restantes magnitudes dependientes de éstas se denotan \(v_1(y),\ldots,v_\mu(y)\). Las transformaciones pueden contener también las derivadas de las \(u\) con respecto a las \(x\), es decir \(\p u/\p x,\ \p^2u/\p x^2,\ldots\).\footnote{Suprimo los índices en lo posible, también en las sumas; así, por ejemplo, \(\p^2u/\p x^2\) denota \(\p^2u_\nu/\p x_i\p x_k\), etc.} Una función se llama un invariante del grupo si existe una relación
\[
P\!\left(x,u,\frac{\p u}{\p x},\frac{\p^2u}{\p x^2},\ldots\right)
 =
P\!\left(y,v,\frac{\p v}{\p y},\frac{\p^2v}{\p y^2},\ldots\right).
\]
En particular, un integral \(I\) será por tanto un invariante del grupo si existe una relación
\begin{equation}
\begin{aligned}
I&=\int\!\cdots\!\int f\!\left(x,u,\frac{\p u}{\p x},\frac{\p^2u}{\p x^2},\ldots\right)\dd x\\
 &=\int\!\cdots\!\int f\!\left(y,v,\frac{\p v}{\p y},\frac{\p^2v}{\p y^2},\ldots\right)\dd y,
\end{aligned}
\tag{1}
\end{equation}
donde la integración se efectúa sobre un dominio real arbitrario de las \(x\) y sobre el dominio correspondiente de las \(y\).\footnote{Escribo abreviadamente \(\dd x,\dd y\) por \(\dd x_1\cdots \dd x_n,\dd y_1\cdots \dd y_n\). Todos los argumentos \(x,u,\eps,p(x)\) que aparecen en las transformaciones han de tomarse reales, mientras que los coeficientes pueden ser complejos. Pero como en los resultados finales se trata de identidades en las \(x,u\), en los parámetros y en las funciones arbitrarias, éstas valen también para valores complejos tan pronto como todas las funciones que aparecen se supongan analíticas. Además, una gran parte de los resultados puede fundarse sin integrales, de modo que la restricción a lo real no es necesaria ni siquiera para la justificación. En cambio, las consideraciones del final del \S{}2 y del comienzo del \S{}5 no parecen ser realizables sin integrales.}

Por otra parte, para un integral arbitrario \(I\), no necesariamente invariante, formo la primera variación \(\delta I\) y la transformo por integración parcial según las reglas del cálculo variacional. Como es sabido, si \(\delta u\), junto con todas las derivadas que aparecen en la frontera, se toma nula en la frontera, pero por lo demás arbitraria, se obtiene:
\begin{equation}
\delta I=\int\!\cdots\!\int \delta f\,\dd x
=\int\!\cdots\!\int\left(\sum_i\psi_i\!\left(x,u,\frac{\p u}{\p x},\ldots\right)\delta u_i\right)\dd x,
\tag{2}
\end{equation}
donde las \(\psi\) denotan las expresiones lagrangianas, es decir los lados izquierdos de las ecuaciones de Lagrange del problema variacional correspondiente \(\delta I=0\). A esta relación integral corresponde una identidad sin integrales en \(\delta u\) y sus derivadas, que surge al escribir también los términos de frontera. Como muestra la integración parcial, estos términos de frontera son integrales de divergencias, es decir de expresiones
\[
\Div A=\frac{\p A_1}{\p x_1}+\cdots+\frac{\p A_n}{\p x_n},
\]
donde \(A\) es lineal en \(\delta u\) y sus derivadas. Así se obtiene:
\begin{equation}
\sum_i\psi_i\,\delta u_i=\delta f+\Div A.
\tag{3}
\end{equation}
En particular, si \(f\) contiene sólo primeras derivadas de las \(u\), entonces, en el caso del integral simple, la identidad (3) coincide con la llamada ``ecuación central lagrangiana'':
\begin{equation}
\sum_i \psi_i\,\delta u_i
=\delta f-\frac{\dd}{\dd x}\left(\sum_i\frac{\p f}{\p u_i'}\,\delta u_i\right),
\qquad \left(u_i'=\frac{\dd u_i}{\dd x}\right),
\tag{4}
\end{equation}
mientras que para el integral \(n\)-múltiple (3) pasa a ser:
\begin{equation}
\sum_i\psi_i\,\delta u_i
=
\delta f
-\frac{\p}{\p x_1}\left(\sum_i\frac{\p f}{\p \left(\frac{\p u_i}{\p x_1}\right)}\,\delta u_i\right)
-\cdots
-\frac{\p}{\p x_n}\left(\sum_i\frac{\p f}{\p \left(\frac{\p u_i}{\p x_n}\right)}\,\delta u_i\right).
\tag{5}
\end{equation}
Para el integral simple y \(\kappa\) derivadas de las \(u\), (3) está dada por:
\begin{equation}
\begin{aligned}
\sum_i\psi_i\,\delta u_i
={}&\delta f
-\frac{\dd}{\dd x}\Bigg\{\sum_i\left[
\binom{1}{1}\frac{\p f}{\p u_i^{(1)}}\delta u_i
+\binom{2}{1}\frac{\p f}{\p u_i^{(2)}}\delta u_i^{(1)}
+\cdots+
\binom{\kappa}{1}\frac{\p f}{\p u_i^{(\kappa)}}\delta u_i^{(\kappa-1)}
\right]\Bigg\}\\
&+\frac{\dd^2}{\dd x^2}\Bigg\{\sum_i\left[
\binom{2}{2}\frac{\p f}{\p u_i^{(2)}}\delta u_i
+\binom{3}{2}\frac{\p f}{\p u_i^{(3)}}\delta u_i^{(1)}
+\cdots+
\binom{\kappa}{2}\frac{\p f}{\p u_i^{(\kappa)}}\delta u_i^{(\kappa-2)}
\right]\Bigg\}\\
&+\cdots
+(-1)^\kappa\frac{\dd^\kappa}{\dd x^\kappa}
\left\{\sum_i\binom{\kappa}{\kappa}
\frac{\p f}{\p u_i^{(\kappa)}}\delta u_i\right\}.
\end{aligned}
\tag{6}
\end{equation}
Una identidad correspondiente vale para el integral \(n\)-múltiple; en particular, \(A\) contiene \(\delta u\) hasta la derivada de orden \((\kappa-1)\). Que (4), (5) y (6) definan efectivamente las expresiones lagrangianas \(\psi_i\) se sigue de que, en las combinaciones de los lados derechos, se eliminan todas las derivadas superiores de \(\delta u\), mientras se cumple la relación (2), a la que conduce de manera única la integración parcial.

En lo que sigue se trata de los dos teoremas siguientes:

I. Si el integral \(I\) es invariante con respecto a una \(\G_\rho\), entonces \(\rho\) combinaciones linealmente independientes de las expresiones lagrangianas se convierten en divergencias; recíprocamente, de ello se sigue la invariancia de \(I\) con respecto a una \(\G_\rho\). El teorema vale también en el caso límite de infinitos parámetros.

II. Si el integral \(I\) es invariante con respecto a una \(\G_{\infty\rho}\) en la cual las funciones arbitrarias aparecen hasta la derivada \(\sigma\)-ésima, entonces existen \(\rho\) relaciones idénticas entre las expresiones lagrangianas y sus derivadas hasta el orden \(\sigma\); también aquí vale el recíproco.\footnote{Para ciertos casos excepcionales triviales compárese el \S{}2, segunda nota.}

Para grupos mixtos rigen las afirmaciones de ambos teoremas; por consiguiente aparecen a la vez dependencias y relaciones de divergencia independientes entre sí.

Al pasar de estas identidades al problema variacional correspondiente, esto es, al poner \(\psi=0\),\footnote{Un poco más generalmente se puede poner también \(\psi_i=T_i\); compárese el \S{}3, primera nota.} el Teorema I dice, en el caso unidimensional --- donde la divergencia se vuelve una diferencial total ---, que existen \(\rho\) primeras integrales, aunque puedan presentarse dependencias no lineales entre ellas;\footnote{Compárese el final del \S{}3.} en el caso multidimensional se obtienen las ecuaciones de divergencia que recientemente se han llamado a menudo ``leyes de conservación''. El Teorema II dice que \(\rho\) de las ecuaciones de Lagrange son consecuencias de las restantes.

El ejemplo más simple del Teorema II --- sin el recíproco --- es la representación paramétrica de Weierstrass; aquí, bajo homogeneidad de primer orden, se sabe que el integral es invariante si se reemplaza la variable independiente \(x\) por una función arbitraria de \(x\), dejando inalterada \(u\), \((y=p(x);\ v_i(y)=u_i(x))\). Así aparece una función arbitraria, pero sin derivadas; y a ello corresponde la conocida relación lineal entre las expresiones lagrangianas mismas:
\[
\sum_i \psi_i\,\frac{\dd u_i}{\dd x}=0.
\]
Otro ejemplo lo ofrece la ``teoría general de la relatividad'' de los físicos. Aquí se trata del grupo de todas las transformaciones de las \(x\), \(y_i=p_i(x)\), mientras las \(u\) (denotadas \(g_{\mu\nu}\) y \(q\)) se someten a las transformaciones inducidas de ese modo para los coeficientes de una forma diferencial cuadrática y de una lineal, transformaciones que contienen las primeras derivadas de las funciones arbitrarias \(p(x)\). A ello corresponden las conocidas \(n\) dependencias entre las expresiones lagrangianas y sus primeras derivadas.\footnote{Compárese, por ejemplo, la exposición de Klein.}

Si, en particular, se especializa el grupo admitiendo que en las transformaciones no aparezcan derivadas de \(u(x)\), y además haciendo que las magnitudes independientes transformadas dependan sólo de las \(x\) y no de las \(u\), entonces, como se mostrará en el \S{}5, la invariancia de \(I\) implica la invariancia relativa de \(\sum_i\psi_i\delta u_i\),\footnote{Es decir, \(\sum_i\psi_i\delta u_i\) adquiere un factor bajo las transformaciones.} y asimismo de las divergencias que aparecen en el Teorema I, tan pronto como los parámetros se someten a transformaciones apropiadas. Se sigue además que las primeras integrales mencionadas también admiten el grupo. Para el Teorema II se obtiene de igual modo la invariancia relativa de los lados izquierdos de las dependencias, combinados por medio de las funciones arbitrarias; y, como consecuencia, también una función cuya divergencia se anula idénticamente y que admite el grupo: la función que en la teoría física de la relatividad media la conexión entre dependencias y teorema de energía.\footnote{Compárese la segunda nota de Klein.} Finalmente, en forma grupo-teórica, el Teorema II da la demostración de una afirmación conexa de Hilbert sobre el fallo de los teoremas de energía propios en la ``relatividad general''. Con estas observaciones suplementarias, el Teorema I contiene todos los teoremas sobre primeras integrales conocidos en mecánica, etc., mientras que el Teorema II puede describirse como la generalización grupo-teórica más amplia posible de la ``teoría general de la relatividad''.

\section*{\S{}2. Relaciones de divergencia y dependencias}

Sea \(\G\) un grupo continuo, finito o infinito. Siempre puede disponerse que la transformación idéntica corresponda al valor cero de los parámetros \(\eps\), respectivamente de las funciones arbitrarias \(p(x)\).\footnote{Compárese, por ejemplo, Lie, \emph{Grundlagen}, p. 331. Si intervienen funciones arbitrarias, los valores especiales \(a^\sigma\) de los parámetros han de reemplazarse por funciones fijas \(p^\sigma,\ \p p^\sigma/\p x,\ldots\); y correspondientemente los valores \(a^\sigma+\eps\) por \(p^\sigma+p(x),\ \p p^\sigma/\p x+\p p/\p x\), etc.} Así, la transformación más general tendrá la forma
\[
y_i=A_i\!\left(x,u,\frac{\p u}{\p x},\ldots\right)=x_i+\D x_i+\cdots,\qquad
v_i(y)=B_i\!\left(x,u,\frac{\p u}{\p x},\ldots\right)=u_i+\D u_i+\cdots,
\]
donde \(\D x_i,\D u_i\) denotan los términos de menor dimensión en \(\eps\), respectivamente en \(p(x)\) y sus derivadas; han de suponerse lineales en estas magnitudes. Como se mostrará más adelante, esto no restringe la generalidad.

Supóngase ahora que el integral \(I\) es invariante con respecto a \(\G\), de modo que se cumple la relación (1). En particular, entonces \(I\) también es invariante con respecto a la transformación infinitesimal contenida en \(\G\):
\[
y_i=x_i+\D x_i,\qquad v_i(y)=u_i+\D u_i.
\]
Para ésta, la relación (1) se convierte en:
\begin{equation}
\begin{aligned}
0=\D I
&=\int\!\cdots\!\int f\!\left(y,v(y),\frac{\p v}{\p y},\ldots\right)\dd y\\
&\qquad-\int\!\cdots\!\int f\!\left(x,u(x),\frac{\p u}{\p x},\ldots\right)\dd x.
\end{aligned}
\tag{7}
\end{equation}
El primer integral debe extenderse sobre el dominio \(x+\D x\) correspondiente al dominio \(x\). Pero esta integración puede transformarse también en una integración sobre el dominio \(x\) mediante la transformación válida para \(\D x\) infinitesimal:
\begin{equation}
\begin{aligned}
\int\!\cdots\!\int f\!\left(y,v(y),\frac{\p v}{\p y},\ldots\right)\dd y
&=\int\!\cdots\!\int f\!\left(x,v(x),\frac{\p v}{\p x},\ldots\right)\dd x\\
&\quad+\int\!\cdots\!\int \Div(f\cdot \D x)\,\dd x.
\end{aligned}
\tag{8}
\end{equation}
Si, por tanto, se introduce, en lugar de la transformación infinitesimal \(\D u\), la variación
\begin{equation}
\bd u_i=v_i(x)-u_i(x)=\D u_i-\sum_\lambda\frac{\p u_i}{\p x_\lambda}\D x_\lambda,
\tag{9}
\end{equation}
entonces (7) y (8) se convierten en:
\begin{equation}
0=\int\!\cdots\!\int\{\bd f+\Div(f\cdot\D x)\}\dd x.
\tag{10}
\end{equation}
Como la relación (10) se cumple al integrar sobre todo dominio arbitrario, el integrando debe anularse idénticamente. Por tanto, las ecuaciones diferenciales de Lie para la invariancia de \(I\) pasan a la relación
\begin{equation}
\bd f+\Div(f\cdot\D x)=0.
\tag{11}
\end{equation}
Si aquí se expresa \(\bd f\), por (3), en términos de las expresiones lagrangianas, se obtiene
\begin{equation}
\sum_i\psi_i\bd u_i=\Div B,\qquad (B=A-f\cdot\D x),
\tag{12}
\end{equation}
y esta relación representa por consiguiente, para todo integral invariante \(I\), una identidad en todos los argumentos que aparecen; es la forma buscada de las ecuaciones diferenciales de Lie para \(I\).\footnote{(12) se convierte en \(0=0\) para el caso trivial --- que sólo puede ocurrir si \(\D x,\D u\) dependen también de derivadas de \(u\) --- en que \(\Div(f\cdot\D x)=0,\ \bd u=0\). Estas transformaciones infinitesimales deben, por tanto, separarse siempre de los grupos, y en la formulación de los teoremas sólo debe contarse el número de los parámetros o funciones arbitrarias restantes. Debe quedar abierta la cuestión de si las transformaciones infinitesimales restantes todavía forman siempre un grupo.}

Supóngase primero que \(\G\) es un grupo continuo finito \(\G_\rho\). Como, por hipótesis, \(\D u\) y \(\D x\) son lineales en los parámetros \(\eps_1,\ldots,\eps_\rho\), por (9) lo mismo vale para \(\bd u\) y sus derivadas; por tanto \(A\) y \(B\) son lineales en las \(\eps\). Pongo, por consiguiente,
\[
B=B^{(1)}\eps_1+\cdots+B^{(\rho)}\eps_\rho,\qquad
\bd u=\bd u^{(1)}\eps_1+\cdots+\bd u^{(\rho)}\eps_\rho,
\]
donde \(\bd u^{(1)},\ldots\) son funciones de \(x,u,\p u/\p x,\ldots\). De (12) se siguen las relaciones de divergencia buscadas:
\begin{equation}
\sum_i\psi_i\bd u_i^{(1)}=\Div B^{(1)},\ldots,\qquad
\sum_i\psi_i\bd u_i^{(\rho)}=\Div B^{(\rho)}.
\tag{13}
\end{equation}
Así, \(\rho\) combinaciones linealmente independientes de las expresiones lagrangianas se convierten en divergencias. La independencia lineal se sigue de que, por (9), \(\bd u=0,\D x=0\) implicaría \(\D u=0,\D x=0\), y por tanto una dependencia entre las transformaciones infinitesimales. Pero por hipótesis no existe tal dependencia para ningún valor de los parámetros; de lo contrario, la \(\G_\rho\) que surge de nuevo por integración a partir de las transformaciones infinitesimales dependería de menos de \(\rho\) parámetros esenciales. La posibilidad adicional \(\bd u=0,\ \Div(f\D x)=0\) ha sido excluida. Estas conclusiones siguen siendo válidas en el caso límite de infinitos parámetros.

Sea ahora \(\G\) un grupo continuo infinito \(\G_{\infty\rho}\). Entonces también \(\bd u\) y sus derivadas, y por tanto también \(B\), son lineales en las funciones arbitrarias \(p(x)\) y en sus derivadas.\footnote{Que no es una restricción suponer que las \(p\) son independientes de \(u,\p u/\p x,\ldots\) se muestra por el recíproco.} La sustitución de los valores de \(\bd u\), todavía independientemente de (12), da:
\[
\sum_i\psi_i\bd u_i
=\sum_{\lambda,i}\psi_i\left\{
a_i^{(\lambda)}(x,u,\ldots)p^{(\lambda)}(x)
+b_i^{(\lambda)}(x,u,\ldots)\frac{\p p^{(\lambda)}}{\p x}
+\cdots+c_i^{(\lambda)}(x,u,\ldots)
\frac{\p^\sigma p^{(\lambda)}}{\p x^\sigma}\right\}.
\]
Ahora, de modo análogo a la fórmula de integración parcial, mediante la identidad
\[
\varphi(x,u,\ldots)\frac{\p^\tau p(x)}{\p x^\tau}
=(-1)^\tau\frac{\p^\tau\varphi}{\p x^\tau}\,p(x)
\quad \bmod \text{ divergencias},
\]
las derivadas de \(p\) pueden reemplazarse por \(p\) mismo y por divergencias que son lineales en \(p\) y sus derivadas. Así:
\begin{equation}
\begin{aligned}
\sum_i\psi_i\bd u_i
={}&\sum_\lambda\left\{
a_i^{(\lambda)}\psi_i-\frac{\p}{\p x}\bigl(b_i^{(\lambda)}\psi_i\bigr)
+\cdots+(-1)^\sigma
\frac{\p^\sigma}{\p x^\sigma}\bigl(c_i^{(\lambda)}\psi_i\bigr)
\right\}p^{(\lambda)}+\Div\Gamma,
\end{aligned}
\tag{14}
\end{equation}
con suma sobre \(i\) dentro de la llave. En combinación con (12) esto da:
\begin{equation}
\sum_\lambda\left\{
a_i^{(\lambda)}\psi_i-\frac{\p}{\p x}\bigl(b_i^{(\lambda)}\psi_i\bigr)
+\cdots+(-1)^\sigma
\frac{\p^\sigma}{\p x^\sigma}\bigl(c_i^{(\lambda)}\psi_i\bigr)
\right\}p^{(\lambda)}
=\Div(B-\Gamma).
\tag{15}
\end{equation}
Formo ahora el integral \(n\)-múltiple de (15), extendido sobre un dominio arbitrario, y elijo las \(p(x)\) de modo que ellas, junto con todas las derivadas que aparecen en \(B-\Gamma\), se anulen en la frontera. Como el integral de una divergencia se reduce a un integral de frontera, también el integral del lado izquierdo de (15) se anula para \(p(x)\) arbitrarias que se anulan en la frontera con un número suficiente de derivadas. De los argumentos usuales se sigue que el integrando se anula para toda \(p(x)\), y se obtienen las \(\rho\) relaciones
\begin{equation}
\sum_i\left\{
a_i^{(\lambda)}\psi_i-\frac{\p}{\p x}\bigl(b_i^{(\lambda)}\psi_i\bigr)
+\cdots+(-1)^\sigma
\frac{\p^\sigma}{\p x^\sigma}\bigl(c_i^{(\lambda)}\psi_i\bigr)
\right\}=0,\qquad(\lambda=1,2,\ldots,\rho).
\tag{16}
\end{equation}
Éstas son las dependencias buscadas entre las expresiones lagrangianas y sus derivadas cuando \(I\) es invariante con respecto a \(\G_{\infty\rho}\). La independencia lineal se sigue como antes, pues el recíproco conduce de nuevo a (12); y porque de las transformaciones infinitesimales se puede volver a concluir las finitas, como se lleva a cabo con más detalle en el \S{}4. En consecuencia, para una \(\G_{\infty\rho}\), ya en las transformaciones infinitesimales aparecen \(\rho\) transformaciones arbitrarias. De (15) y (16) se sigue además que \(\Div(B-\Gamma)=0\).

Si, correspondiendo a un ``grupo mixto'', se toman \(\D x\) y \(\D u\) lineales en las \(\eps\) y en las \(p(x)\), entonces, poniendo una vez las \(p(x)\) y otra vez las \(\eps\) iguales a cero, se ve que existen tanto las relaciones de divergencia (13) como las dependencias (16).

\section*{\S{}3. Recíproco en el caso del grupo finito}

Para demostrar el recíproco, las consideraciones precedentes deben recorrerse primero, en lo esencial, en orden inverso. De la validez de (13), multiplicando por las \(\eps\) y sumando, resulta la validez de (12); y por medio de la identidad (3) se obtiene una relación
\[
\bd f+\Div(A-B)=0.
\]
Así, si se pone \(\D x=\frac1f(A-B)\), se llega con ello a (11); por integración esto da finalmente (7), \(\D I=0\), y por tanto la invariancia de \(I\) con respecto a la transformación infinitesimal determinada por \(\D x,\D u\), donde las \(\D u\) se determinan a partir de \(\D x\) y \(\bd u\) mediante (9), y \(\D x,\D u\) se vuelven lineales en los parámetros. Pero \(\D I=0\) implica, de la manera conocida, la invariancia de \(I\) con respecto a las transformaciones finitas que surgen por integración del sistema simultáneo:
\begin{equation}
\frac{\dd x_i}{\dd t}=\D x_i,\qquad
\frac{\dd u_i}{\dd t}=\D u_i,\qquad
\left\{\begin{array}{l}
x_i=y_i,\\
u_i=v_i
\end{array}\right\}
\quad\text{para }t=0.
\tag{17}
\end{equation}
Estas transformaciones finitas contienen \(\rho\) parámetros \(a_1,\ldots,a_\rho\), a saber, las combinaciones \(t\eps_1,\ldots,t\eps_\rho\). De la hipótesis de que haya \(\rho\), y sólo \(\rho\), relaciones de divergencia linealmente independientes (13), se sigue además que las transformaciones finitas, tan pronto como no contienen las derivadas \(\p u/\p x\), forman siempre un grupo. En caso contrario, al menos una transformación infinitesimal surgida por el proceso de corchetes de Lie no sería una combinación lineal de las \(\rho\) restantes; y como \(I\) admite también esta transformación, habría más de \(\rho\) relaciones de divergencia linealmente independientes. O bien esta transformación infinitesimal sería de la forma especial \(\bd u=0,\ \Div(f\D x)=0\). Pero entonces \(\D x\) o \(\D u\) dependerían de derivadas, contra la hipótesis. Debe quedar abierta la cuestión de si este caso puede ocurrir cuando aparecen derivadas en \(\D x\) o \(\D u\); entonces todas las funciones \(\D x\) para las cuales \(\Div(f\D x)=0\) deben añadirse todavía a las \(\D x\) determinadas arriba, para obtener de nuevo la propiedad de grupo, pero por convenio los parámetros añadidos de esta manera no han de contarse. Con esto queda demostrado el recíproco.

De este recíproco se sigue también que \(\D x\) y \(\D u\) pueden suponerse efectivamente lineales en los parámetros. Si \(\D u\) y \(\D x\) fuesen formas de grado superior en \(\eps\), entonces, por la independencia lineal de los productos de potencias de las \(\eps\), se seguirían relaciones enteramente correspondientes a (13), sólo que en mayor número; por el recíproco, éstas implican la invariancia de \(I\) con respecto a un grupo cuyas transformaciones infinitesimales contienen los parámetros linealmente. Si este grupo ha de contener exactamente \(\rho\) parámetros, entonces deben existir dependencias lineales entre las relaciones de divergencia obtenidas originalmente de los términos de grado superior en \(\eps\).

Debe señalarse también que, en el caso en que \(\D x\) y \(\D u\) contengan asimismo derivadas de las \(u\), las transformaciones finitas pueden depender de infinitamente muchas derivadas de las \(u\); pues en este caso la integración de (17), al determinar
\[
\frac{\dd^2x_i}{\dd t^2},\quad \frac{\dd^2u_i}{\dd t^2},
\]
conduce a
\[
\D\!\left(\frac{\p u}{\p x_\lambda}\right)
=\frac{\p\,\D u}{\p x_\lambda}
-\sum_\nu\frac{\p u}{\p x_\nu}
\frac{\p\,\D x_\nu}{\p x_\lambda},
\]
de modo que, en general, el número de derivadas de \(u\) crece en cada paso. Un ejemplo es:
\[
f=\frac12 u'^2,\qquad
\psi=-u'',\qquad
\psi\cdot x=\frac{\dd}{\dd x}\{(u-u'x)\},\qquad
\bd u=x\eps,
\]
\[
\D x=-\frac{2u}{u'}\,\eps,\qquad
\D u=\left(x-\frac{2u}{u'}\right)\eps.
\]
Como las expresiones lagrangianas de una divergencia se anulan idénticamente, el recíproco muestra finalmente lo siguiente: si \(I\) admite una \(\G_\rho\), entonces todo integral que difiera de \(I\) sólo en un integral de frontera, es decir en un integral de una divergencia, también admite una \(\G_\rho\) con las mismas \(\bd u\), mientras que sus transformaciones infinitesimales contendrán en general derivadas de las \(u\). Así, correspondiendo al ejemplo precedente,
\[
f^*=\frac12\left\{u'^2-\frac{\dd}{\dd x}\left(\frac{u^2}{x}\right)\right\}
\]
admite la transformación infinitesimal \(\D u=x\eps,\ \D x=0\); en cambio, en las transformaciones infinitesimales correspondientes a \(f\) aparecen derivadas de \(u\).

Si se pasa al problema variacional, es decir, si se pone \(\psi_i=0\),\footnote{\(\psi_i=0\), o algo más generalmente \(\psi_i=T_i\), donde las \(T_i\) son funciones recientemente adjuntadas, se llaman en física ``ecuaciones de campo''. En el caso \(\psi_i=T_i\), las identidades (13) pasan a las ecuaciones \(\Div B^{(\lambda)}=\sum_i T_i\bd u_i^{(\lambda)}\), que en física se siguen llamando leyes de conservación.} entonces (13) se convierte en las ecuaciones
\[
\Div B^{(1)}=0,\ldots,\qquad \Div B^{(\rho)}=0,
\]
que a menudo se llaman ``leyes de conservación''. En el caso unidimensional esto da:
\[
B^{(1)}=\mathrm{const.},\ldots,\qquad B^{(\rho)}=\mathrm{const.}.
\]
Aquí las \(B\) contienen a lo sumo derivadas de \(u\) de orden \((2\kappa-1)\) (por (6)), tan pronto como \(\D u\) y \(\D x\) no contienen derivadas superiores a las \(\kappa\)-ésimas que aparecen en \(f\). Como \(\psi\) contiene, en general, derivadas \(2\kappa\)-ésimas,\footnote{Tan pronto como \(f\) es no lineal en las derivadas \(\kappa\)-ésimas.} se tiene así la existencia de \(\rho\) primeras integrales. Que pueden existir dependencias no lineales entre ellas lo muestra de nuevo la \(f\) anterior. A las transformaciones linealmente independientes \(\D u=\eps_1,\ \D x=\eps_2\) corresponden las relaciones linealmente independientes
\[
u''=\frac{\dd}{\dd x}u',\qquad
u''u'=\frac12\frac{\dd}{\dd x}(u')^2,
\]
mientras que entre las primeras integrales \(u'=\mathrm{const.},\ u'^2=\mathrm{const.}\) existe una dependencia no lineal. Este es el caso elemental en que \(\D u,\D x\) no contienen derivadas de \(u\).\footnote{En caso contrario se tiene además \(u''u'^{\lambda-1}=\mathrm{const.}\) para todo \(\lambda\), correspondiente a
\[
u''(u')^{\lambda-1}=\frac1\lambda\frac{\dd}{\dd x}(u')^\lambda.
\]}

\section*{\S{}4. Recíproco en el caso del grupo infinito}

Primero debe mostrarse que la hipótesis de linealidad de \(\D x\) y \(\D u\) no constituye restricción. Aquí esto se sigue, aun sin el recíproco, del hecho de que \(\G_{\infty\rho}\) depende formalmente de \(\rho\), y sólo de \(\rho\), funciones arbitrarias. En efecto, en el caso no lineal, al componer transformaciones se suman los términos de menor orden y aumentaría el número de funciones arbitrarias. Supóngase, por ejemplo, que
\[
\begin{aligned}
y&=A\!\left(x,u,\frac{\p u}{\p x},\ldots;p\right)\\
&=x+\sum a(x,u,\ldots)p^\nu
+b(x,u,\ldots)p^{\nu-1}\frac{\p p}{\p x}
+c\,p^{\nu-2}\left(\frac{\p p}{\p x}\right)^2+\cdots
+d\left(\frac{\p p}{\p x}\right)^\nu+\cdots,
\end{aligned}
\]
donde
\[
p^\nu=(p^{(1)})^{\nu_1}\cdots(p^{(\rho)})^{\nu_\rho},
\]
y correspondientemente
\[
v=B\!\left(x,u,\frac{\p u}{\p x},\ldots;p\right).
\]
Entonces, por composición con
\[
z=A\!\left(y,v,\frac{\p v}{\p y},\ldots;q\right),
\]
los términos de menor orden dan
\[
z=x+\sum a(p^\nu+q^\nu)+
b\left\{p^{\nu-1}\frac{\p p}{\p x}+q^{\nu-1}\frac{\p q}{\p x}\right\}
+c\left\{p^{\nu-2}\left(\frac{\p p}{\p x}\right)^2
+q^{\nu-2}\left(\frac{\p q}{\p x}\right)^2\right\}+\cdots.
\]
Si aquí algún coeficiente distinto de \(a\) y \(b\) es diferente de cero, de modo que un término
\[
p^{\nu-\sigma}\left(\frac{\p p}{\p x}\right)^\sigma
+q^{\nu-\sigma}\left(\frac{\p q}{\p x}\right)^\sigma
\]
aparece efectivamente para \(\sigma>1\), entonces no puede escribirse como el cociente diferencial de una única función, ni como producto de potencias de tal cociente. Así, el número de funciones arbitrarias ha aumentado contra la hipótesis. Si todos los coeficientes distintos de \(a\) y \(b\) se anulan, entonces, según los valores de los exponentes \(\nu_1,\ldots,\nu_\rho\), el segundo término es el cociente diferencial del primero (como, por ejemplo, siempre para una \(\G_{\infty1}\)), de modo que aparece efectivamente la linealidad; o bien también aquí aumenta el número de funciones arbitrarias. Por tanto, debido a la linealidad en \(p(x)\), las transformaciones infinitesimales satisfacen un sistema de ecuaciones diferenciales parciales lineales; y, como se cumple la propiedad de grupo, forman un ``grupo infinito de transformaciones infinitesimales'' en el sentido de Lie (\emph{Grundlagen}, \S{}10).

El recíproco resulta ahora de manera semejante al caso del grupo finito. La existencia de las dependencias (16), después de multiplicar por \(p^{(\lambda)}(x)\) y sumar, conduce por medio de la transformación idéntica (14) a
\[
\sum_i\psi_i\bd u_i=\Div\Gamma.
\]
De aquí, como en el \S{}3, se obtiene la determinación de \(\D x\) y \(\D u\), y la invariancia de \(I\) con respecto a estas transformaciones infinitesimales, que en efecto dependen linealmente de \(\rho\) funciones arbitrarias y de sus derivadas hasta el orden \(\sigma\). Que estas transformaciones infinitesimales, cuando no contienen derivadas \(\p u/\p x,\ldots\), forman ciertamente un grupo se sigue como en el \S{}3, pues de lo contrario aparecerían más funciones arbitrarias por composición, mientras que por hipótesis sólo deben existir \(\rho\) dependencias (16). Así forman un ``grupo infinito de transformaciones infinitesimales''. Pero tal grupo consiste (\emph{Grundlagen}, Teorema VII, p. 391) en las transformaciones infinitesimales más generales de un cierto ``grupo infinito \(\G\) de transformaciones finitas'' definido de ese modo, en el sentido de Lie. Cada transformación finita se genera de este modo a partir de infinitesimales (\emph{Grundlagen}, \S{}7),\footnote{De esto se sigue en particular que el grupo \(\G\) generado a partir de las transformaciones infinitesimales \(\D x,\D u\) de una \(\G_{\infty\rho}\) conduce de nuevo a \(\G_{\infty\rho}\). Pues \(\G_{\infty\rho}\) no contiene transformaciones infinitesimales dependientes de funciones arbitrarias y distintas de \(\D x,\D u\), y tampoco puede contener transformaciones independientes dependientes de parámetros, ya que entonces sería un grupo mixto. Pero por el argumento precedente las transformaciones finitas están determinadas por las infinitesimales.} y por tanto surge por integración del sistema simultáneo:
\[
\frac{\dd x_i}{\dd t}=\D x_i,\qquad
\frac{\dd u_i}{\dd t}=\D u_i,\qquad
\left\{\begin{array}{l}
x_i=y_i,\\
u_i=v
\end{array}\right\}
\quad\text{para }t=0.
\]
Puede ser necesario, sin embargo, permitir que las \(p(x)\) arbitrarias dependan todavía de \(t\). Así \(\G\) depende efectivamente de \(\rho\) funciones arbitrarias. En particular, si basta tomar \(p(x)\) independiente de \(t\), entonces esta dependencia es analítica en las funciones arbitrarias \(q(x)=t\,p(x)\).\footnote{La cuestión de si quizá este último caso se presenta siempre fue planteada por Lie en otra formulación (\emph{Grundlagen}, \S{}7 y \S{}13, final).} Si aparecen derivadas \(\p u/\p x,\ldots\), puede ser necesario adjuntar transformaciones infinitesimales \(\bd u=0,\ \Div(f\D x)=0\) antes de extraer las mismas conclusiones.

Siguiendo un ejemplo de Lie (\emph{Grundlagen}, \S{}7), indico ahora un caso bastante general en el que se llega hasta fórmulas explícitas, las cuales muestran también que las derivadas de las funciones arbitrarias aparecen hasta el orden \(\sigma\); en este caso, por tanto, el recíproco es completo. Se trata de grupos de transformaciones infinitesimales a los que corresponde el grupo de todas las transformaciones de las \(x\) y las transformaciones de las \(u\) ``inducidas'' por aquéllas; es decir, transformaciones de las \(u\) en las que \(\D u\), y por tanto \(u\), depende sólo de las funciones arbitrarias que aparecen en \(\D x\). Se sigue suponiendo que las derivadas \(\p u/\p x,\ldots\) no aparecen en \(\D u\). Así:
\[
\D x_i=p^{(i)}(x),\qquad
\D u_i=
\sum_{\lambda=1}^{n}\left\{
a_i^{(\lambda)}(x,u)p^{(\lambda)}
+b_i^{(\lambda)}\frac{\p p^{(\lambda)}}{\p x}
+\cdots+
c_i^{(\lambda)}\frac{\p^\sigma p^{(\lambda)}}{\p x^\sigma}
\right\}.
\]
Como la transformación infinitesimal \(\D x=p(x)\) genera toda transformación \(x=y+g(y)\) con \(g(y)\) arbitraria, se puede en particular determinar \(p(x)\) como función de \(t\) de manera que se genere el grupo de un parámetro:
\begin{equation}
x_i=y_i+t\,g_i(y),
\tag{18}
\end{equation}
que para \(t=0\) se convierte en la identidad y para \(t=1\) en la \(x=y+g(y)\) deseada. Al diferenciar (18) se obtiene
\begin{equation}
\frac{\dd x_i}{\dd t}=g_i(y)=p^{(i)}(x,t),
\tag{19}
\end{equation}
donde \(p(x,t)\) se determina a partir de \(g(y)\) por inversión de (18); recíprocamente, (18) surge de (19) mediante la condición lateral \(x_i=y_i\) para \(t=0\), por la cual el integral queda fijado de manera única. Mediante (18), las \(x\) en \(\D u\) pueden reemplazarse por las constantes de integración \(y\) y por \(t\); entonces las \(g(y)\) aparecen exactamente hasta sus derivadas \(\sigma\)-ésimas, ya que en
\[
\frac{\p p}{\p x}=\sum_k\frac{\p g}{\p y_k}\frac{\p y_k}{\p x}
\]
las \(\p y/\p x\) se expresan por \(\p x/\p y\), y en general \(\p^\sigma p/\p x^\sigma\) se reemplaza por su valor en
\[
\frac{\p g}{\p y},\ldots,\frac{\p^\sigma g}{\p y^\sigma},\ldots,\frac{\p^\sigma x}{\p y^\sigma}.
\]
Para la determinación de las \(u\) se obtiene por tanto el sistema de ecuaciones
\[
\frac{\dd u_i}{\dd t}
=
F_i\!\left(g(y),\frac{\p g}{\p y},\ldots,\frac{\p^\sigma g}{\p y^\sigma},u,t\right),
\qquad (u_i=v_i\ \text{para }t=0),
\]
en el que sólo \(t\) y \(u\) son variables, mientras que \(g(y),\ldots\) pertenecen al dominio de coeficientes; por tanto la integración da
\[
u_i=v_i+
B_i\!\left(v,g(y),\frac{\p g}{\p y},\ldots,\frac{\p^\sigma g}{\p y^\sigma},t\right)_{t=1}.
\]
Así se obtienen transformaciones que dependen exactamente de \(\sigma\) derivadas de las funciones arbitrarias. La identidad está contenida en esto para \(g(y)=0\), por (18); y la propiedad de grupo se sigue del hecho de que el procedimiento recién descrito produce toda transformación \(x=y+g(y)\), con lo cual la transformación inducida de las \(u\) queda determinada unívocamente, de modo que el grupo \(\G\) queda agotado.

Del recíproco se sigue incidentalmente que no es restricción suponer que las funciones arbitrarias dependan sólo de las \(x\), y no de \(u,\p u/\p x,\ldots\). En este último caso, en la transformación idéntica (14), y por tanto también en (15), además de las \(p^{(\lambda)}\) aparecerían las derivadas \(\p p^{(\lambda)}/\p u,\ \p p^{(\lambda)}/\p(\p u/\p x),\ldots\). Si se toman ahora sucesivamente las \(p^{(\lambda)}\) de grado cero, primero y superior en \(u,\p u/\p x,\ldots\), con funciones arbitrarias de \(x\) como coeficientes, entonces surgen de nuevo dependencias (16), sólo que en mayor número. Pero por el recíproco precedente pueden reconducirse al caso anterior combinándolas con funciones arbitrarias que dependen sólo de \(x\). De manera similar se muestra que a los grupos mixtos corresponde la aparición simultánea de dependencias y relaciones de divergencia independientes.\footnote{Como en el \S{}3, también aquí se sigue del recíproco que, además de \(I\), todo integral \(I^*\) que difiera de él por un integral de una divergencia admite asimismo un grupo infinito con las mismas \(\bd u\); aquí, sin embargo, \(\D x\) y \(\D u\) contendrán en general derivadas de las \(u\). Einstein introdujo un tal integral \(I^*\) en la relatividad general para obtener una formulación más simple de los teoremas de energía. Doy las transformaciones infinitesimales admitidas por este \(I^*\), usando exactamente la notación de la segunda nota de Klein. El integral \(I=\int\!\cdots\!\int K\,\dd\omega=\int\!\cdots\!\int \mathfrak{R}\,\dd S\) admite el grupo de todas las transformaciones de las \(w\) y el grupo inducido por él para las \(g_{\mu\nu}\). A esto corresponden las dependencias ((30) en Klein):
\[
\sum \mathfrak{R}_{\mu\nu}g_{\mu\nu}^{\alpha\beta}
+2\sum\frac{\p g_{\mu\nu}^{\alpha\beta}\mathfrak{R}_{\mu\nu}}{\p w^\sigma}=0.
\]
Ahora \(I^*=\int\!\cdots\!\int \mathfrak{R}^*\,\dd S\), donde \(\mathfrak{R}^*=\mathfrak{R}+\Div\), y por consiguiente \(\mathfrak{R}_{\mu\nu}^*=\mathfrak{R}_{\mu\nu}\), donde \(\mathfrak{R}_{\mu\nu}^*,\mathfrak{R}_{\mu\nu}\) denotan respectivamente las expresiones lagrangianas. Las dependencias dadas son, por tanto, también dependencias para \(\mathfrak{R}_{\mu\nu}^*\); y después de multiplicar por \(p^\tau\) y sumar, invirtiendo las transformaciones de la diferenciación de productos, se obtiene
\[
\sum \mathfrak{R}_{\mu\nu}p^{\mu\nu}
+2\Div\!\left(\sum g^{\mu\sigma}\mathfrak{R}_{\mu\tau}p^\tau\right)=0,
\]
\[
\delta\mathfrak{R}^*+\Div\!\left(\sum 2g^{\mu\sigma}\mathfrak{R}_{\mu\tau}p^\tau
-\frac{\p\mathfrak{R}^*}{\p g_{\mu\nu}^{\sigma}}p^{\mu\nu}\right)=0.
\]
La comparación con la ecuación diferencial de Lie
\(\delta\mathfrak{R}^*+\Div(\mathfrak{R}^*\D w)=0\) da
\[
\D w^\sigma=\frac1{\mathfrak{R}^*}\left(\sum 2g^{\mu\sigma}\mathfrak{R}_{\mu\tau}p^\tau
-\frac{\p\mathfrak{R}^*}{\p g_{\mu\nu}^{\sigma}}p^{\mu\nu}\right),
\qquad
\D g^{\mu\nu}=p^{\mu\nu}+\sum g^{\mu\nu}_{\sigma}\D w^\sigma,
\]
como transformaciones infinitesimales admitidas por \(I^*\). Estas transformaciones infinitesimales dependen, por tanto, de las primeras y segundas derivadas de las \(g^{\mu\nu}\), y contienen las \(p\) arbitrarias hasta la primera derivada.}

\section*{\S{}5. Invariancia de los componentes individuales de las relaciones}

Si el grupo \(\G\) se especializa al caso más simple, el que ordinariamente se considera, no admitiendo derivadas de las \(u\) en las transformaciones y haciendo que las variables independientes transformadas dependan sólo de las \(x\) y no de las \(u\), entonces se puede inferir la invariancia de los componentes individuales en las fórmulas. Primero, por argumentos familiares, se obtiene la invariancia de
\[
\int\!\cdots\!\int\left(\sum_i\psi_i\delta u_i\right)\dd x,
\]
y por tanto la invariancia relativa de \(\sum_i\psi_i\delta u_i\), donde \(\delta\) denota una variación cualquiera. En efecto, por una parte,
\[
\delta I
=\int\!\cdots\!\int \delta f\!\left(x,u,\frac{\p u}{\p x},\ldots\right)\dd x
=\int\!\cdots\!\int \delta f\!\left(y,v,\frac{\p v}{\p y},\ldots\right)\dd y.
\]
Por otra parte, para \(\delta u,\delta(\p u/\p x),\ldots\) que se anulan en la frontera, a lo cual, por la transformación lineal homogénea de \(\delta u,\delta(\p u/\p x),\ldots\), corresponden \(\delta v,\delta(\p v/\p y),\ldots\) que también se anulan en la frontera, se tiene
\[
\int\!\cdots\!\int \delta f\!\left(x,u,\frac{\p u}{\p x},\ldots\right)\dd x
=
\int\!\cdots\!\int\left(\sum_i\psi_i(u,\ldots)\delta u_i\right)\dd x,
\]
\[
\int\!\cdots\!\int \delta f\!\left(y,v,\frac{\p v}{\p y},\ldots\right)\dd y
=
\int\!\cdots\!\int\left(\sum_i\psi_i(v,\ldots)\delta v_i\right)\dd y.
\]
Así, para \(\delta u,\delta(\p u/\p x),\ldots\) que se anulan en la frontera,
\[
\int\!\cdots\!\int\left(\sum_i\psi_i(u,\ldots)\delta u_i\right)\dd x
=
\int\!\cdots\!\int\left(\sum_i\psi_i(v,\ldots)\delta v_i\right)\dd y
\]
\[
=
\int\!\cdots\!\int\left(\sum_i\psi_i(v,\ldots)\delta v_i\right)
\left|\frac{\p y_i}{\p x_k}\right|\dd x.
\]
Si en el tercer integral se expresan \(y,v,\delta v\) por \(x,u,\delta u\), y se iguala al primero, se obtiene una relación
\[
\int\!\cdots\!\int\left(\sum_i\chi_i(u,\ldots)\delta u_i\right)\dd x=0
\]
para \(\delta u\) nula en la frontera pero por lo demás arbitraria. De aquí se sigue, como es sabido, que el integrando se anula para \(\delta u\) arbitraria; por tanto, la relación
\[
\sum_i\psi_i(u,\ldots)\delta u_i
=
\left|\frac{\p y_i}{\p x_k}\right|
\left(\sum_i\psi_i(v,\ldots)\delta v_i\right)
\]
vale idénticamente en \(\delta u\). Esto expresa la invariancia relativa de \(\sum_i\psi_i\delta u_i\), y por consiguiente la invariancia de
\[
\int\!\cdots\!\int\left(\sum_i\psi_i\delta u_i\right)\dd x.
\]
\footnote{Estos argumentos fallan si \(y\) depende también de las \(u\), porque entonces \(\delta f(y,v,\p v/\p y,\ldots)\) contiene también términos \(\sum(\p f/\p y)\delta y\), de modo que la transformación de divergencia no conduce a las expresiones lagrangianas. Fallan asimismo si se admiten derivadas de las \(u\); entonces las \(\delta v\) se convierten en combinaciones lineales de \(\delta u,\delta(\p u/\p x),\ldots\), y por tanto sólo después de una nueva transformación de divergencia conducen a una identidad \(\int\!\cdots\!\int(\sum\chi_i(u,\ldots)\delta u_i)\dd x=0\), de modo que en el lado derecho tampoco aparecen las expresiones lagrangianas. La pregunta de si la invariancia de \(\int\!\cdots\!\int(\sum\psi_i\delta u_i)\dd x\) permite ya concluir la existencia de relaciones de divergencia es, por el recíproco, equivalente a la pregunta de si implica la invariancia de \(I\) con respecto a un grupo que no necesita conducir a las mismas \(\D u,\D x\), pero sí a las mismas \(\bd u\). En el caso especial del integral simple y sólo primeras derivadas en \(f\), para el grupo finito se puede inferir la existencia de primeras integrales de la invariancia de las expresiones lagrangianas (compárese, por ejemplo, Engel, Gött. Nachr. 1916, p. 270).}

Para aplicar esto a las relaciones de divergencia y dependencias obtenidas, hay que mostrar primero que la \(\bd u\) derivada de \(\D u,\D x\) satisface efectivamente las leyes de transformación de la variación \(\delta u\), con tal sólo de que los parámetros, respectivamente las funciones arbitrarias, en \(\bd v\) se determinen como corresponden al grupo similar de transformaciones infinitesimales en \(y,v\). Sea \(T_q\) la transformación que lleva \(x,u\) a \(y,v\), y sea \(T_p\) una transformación infinitesimal en \(x,u\). Entonces la similar en \(y,v\) está dada por \(T_r=T_qT_pT_q^{-1}\), donde los parámetros, respectivamente las funciones arbitrarias, \(r\) están determinados por \(p\) y \(q\). En fórmulas, esto se expresa así:
\[
T_p:\ \xi=x+\D x(x,p),\qquad u^*=u+\D u(x,u,p);
\]
\[
T_q:\ y=A(x,q),\qquad v=B(x,u,q);
\]
\[
T_qT_p:\ \eta=A(x+\D x(x,p),q),\quad
v^*=B(x+\D x(p),\,u+\D u(p),q).
\]
Pero de esto surge \(T_r=T_qT_pT_q^{-1}\), luego
\[
\eta=y+\D y(r),\qquad v^*=v+\D v(r),
\]
donde, por medio de la inversa de \(T_q\), se consideran las \(x\) como funciones de las \(y\) y se retienen sólo los términos infinitesimales. Así se tiene la identidad
\begin{equation}
\eta=y+\D y(r)=y+\sum\frac{\p A(x,q)}{\p x}\D x(p),
\tag{20a}
\end{equation}
\begin{equation}
v^*=v+\D v(r)=v+\sum\frac{\p B(x,u,q)}{\p x}\D x(p)
+\sum\frac{\p B(x,u,q)}{\p u}\D u(p).
\tag{20b}
\end{equation}
Si en esto se reemplaza \(\xi=x+\D x\) por \(\xi-\D \xi\), de modo que \(\xi\) vuelve a \(x\) y por tanto \(\D x\) se anula, entonces por la primera fórmula (20) también \(\eta\) vuelve a \(y=\eta-\D\eta\). Si bajo esta sustitución \(\D u(p)\) pasa a \(\bd u(p)\), entonces también \(\D v(r)\) pasa a \(\bd v(r)\), y la segunda fórmula (20) da
\[
v+\bd v(y,v,\ldots,r)=v+\sum\frac{\p B(x,u,q)}{\p u}\,\bd u(p),
\]
\[
\bd v(y,v,\ldots,r)=\sum\frac{\p B}{\p u_\kappa}\,\bd u_\kappa(x,u,p).
\]
Así, las fórmulas de transformación para variaciones se satisfacen efectivamente, con tal sólo de que \(\bd v\) se tome dependiente de los parámetros, respectivamente funciones arbitrarias, \(r\).\footnote{Se ve de nuevo que \(y\) debe suponerse independiente de \(u\), etc., si han de valer las conclusiones. Como ejemplo pueden citarse las \(\delta g^{\mu\nu}\) y \(\delta q_\sigma\) indicadas por Klein, que satisfacen las transformaciones para variaciones tan pronto como las \(p\) se someten a una transformación vectorial.}

Se sigue en particular que \(\sum_i\psi_i\bd u_i\) es relativamente invariante; por tanto también, por (12), puesto que las relaciones de divergencia se cumplen igualmente en \(y,v\), la invariancia relativa de \(\Div B\); y además, por (14) y (13), la invariancia relativa de \(\Div\Gamma\) y de los lados izquierdos de las dependencias combinados con las \(p^{(\lambda)}\), donde en las fórmulas transformadas las \(p(x)\) arbitrarias (respectivamente los parámetros) deben reemplazarse por las \(r\) en todas partes. De esto se sigue también la invariancia relativa de \(\Div(B-\Gamma)\), es decir, de la divergencia de un sistema de funciones \(B-\Gamma\) que no se anula idénticamente, pero cuya divergencia se anula idénticamente.

De la invariancia relativa de \(\Div B\), en el caso unidimensional y para un grupo finito, se puede extraer una conclusión adicional sobre la invariancia de las primeras integrales. La transformación de los parámetros correspondiente a la transformación infinitesimal se vuelve, por (20), lineal y homogénea; y como todas las transformaciones son invertibles, las \(\eps\) son también lineales y homogéneas en los parámetros transformados \(\eps^*\). Esta invertibilidad se conserva ciertamente cuando se pone \(\psi=0\), ya que en (20) no aparecen derivadas de \(u\). Igualando los coeficientes de las \(\eps^*\) en
\[
\Div B(x,u,\ldots,\eps)=\frac{\dd y}{\dd x}\,\Div B(y,v,\ldots,\eps^*)
\]
las magnitudes \(\frac{\dd}{\dd y}B^{(\lambda)}(y,v,\ldots)\) se convierten también en funciones lineales homogéneas de las \(\frac{\dd}{\dd x}B^{(\lambda)}(x,u,\ldots)\). Así, de
\[
\frac{\dd}{\dd x}B^{(\lambda)}(x,u,\ldots)=0
\quad\text{o}\quad B^{(\lambda)}(x,u)=\mathrm{const.}
\]
se sigue que
\[
\frac{\dd}{\dd y}B^{(\lambda)}(y,v,\ldots)=0
\quad\text{o}\quad B^{(\lambda)}(y,v)=\mathrm{const.}
\]
Las \(\rho\) primeras integrales correspondientes a una \(\G_\rho\) admiten, por tanto, cada una el grupo, de modo que la integración ulterior también se simplifica. El ejemplo más simple de esto es que \(f\) no contenga \(x\) o no contenga una \(u\), correspondiente respectivamente a las transformaciones infinitesimales \(\D x=\eps,\D u=0\) y \(\D x=0,\D u=\eps\). Entonces \(\bd u=-\eps\,\dd u/\dd x\), respectivamente \(\eps\); y como \(B\) se deriva de \(f\) y \(\bd u\) por diferenciaciones y operaciones racionales, tampoco contiene \(x\), respectivamente \(u\), y admite los grupos correspondientes.\footnote{En aquellos casos en que la existencia de primeras integrales ya se sigue de la invariancia de \(\int(\sum\psi_i\delta u_i)\dd x\), éstas no admiten el grupo completo \(\G_\rho\). Por ejemplo, \(\int(u''\delta u)\dd x\) admite la transformación infinitesimal \(\D x=\eps_2,\D u=\eps_1+x\eps_3\); mientras que la primera integral \(u-u'x=\mathrm{const.}\), correspondiente a \(\D x=0,\D u=x\eps_3\), no admite las otras dos transformaciones infinitesimales, ya que contiene explícitamente tanto \(u\) como \(x\). A esta primera integral corresponden transformaciones infinitesimales para \(f\) que contienen derivadas. Así se ve que la invariancia de \(\int\!\cdots\!\int(\sum\psi_i\delta u_i)\dd x\) logra en cualquier caso menos que la invariancia de \(I\), punto relevante para la cuestión planteada en una nota precedente.}

\section*{\S{}6. Una afirmación de Hilbert}

De lo anterior resulta finalmente la demostración de una afirmación de Hilbert sobre la conexión entre el fallo de los teoremas de energía propios y la ``relatividad general'' (primera nota de Klein, Göttinger Nachr. 1917, Respuesta, primer párrafo), en una formulación grupo-teórica generalizada.

Supóngase que el integral \(I\) admite una \(\G_{\infty\rho}\), y sea \(\G_\sigma\) un grupo finito cualquiera que surge por especialización de las funciones arbitrarias, por tanto un subgrupo de \(\G_{\infty\rho}\). Al grupo infinito \(\G_{\infty\rho}\) corresponden dependencias (16), al grupo finito \(\G_\sigma\) relaciones de divergencia (13); recíprocamente, de la existencia de cualesquiera relaciones de divergencia se sigue la invariancia de \(I\) con respecto a un grupo finito, que coincide con \(\G_\sigma\) si y sólo si las \(\bd u\) son combinaciones lineales de las que surgen de \(\G_\sigma\). Por consiguiente, la invariancia con respecto a \(\G_\sigma\) no puede conducir a relaciones de divergencia distintas de (13). Pero como de la existencia de (16) se sigue la invariancia de \(I\) con respecto a las transformaciones infinitesimales \(\D u,\D x\) de \(\G_{\infty\rho}\) para \(p(x)\) arbitraria, se sigue en particular que \(I\) ya es invariante con respecto a las transformaciones infinitesimales de \(\G_\sigma\) que resultan de éstas por especialización, y por tanto con respecto a \(\G_\sigma\). Las relaciones de divergencia
\[
\sum_i\psi_i\bd u_i^{(\lambda)}=\Div B^{(\lambda)}
\]
deben, por tanto, ser consecuencias de las dependencias (16), que también pueden escribirse
\[
\sum_i\psi_i a_i^{(\lambda)}=\Div\chi^{(\lambda)},
\]
donde las \(\chi^{(\lambda)}\) son combinaciones lineales de las expresiones lagrangianas y sus derivadas. Como las \(\psi\) entran linealmente tanto en (13) como en (16), las relaciones de divergencia deben ser en particular combinaciones lineales de las dependencias (16). Así,
\[
\Div B^{(\lambda)}=\Div\left(\sum \alpha_\mu\chi^{(\mu)}\right),
\]
y las \(B^{(\lambda)}\) mismas están compuestas linealmente de las \(\chi\), es decir, de las expresiones lagrangianas y sus derivadas, y de funciones cuya divergencia se anula idénticamente, como por ejemplo el \(B-\Gamma\) que aparece al final del \S{}2, para el cual \(\Div(B-\Gamma)=0\), y donde la divergencia tiene al mismo tiempo el carácter de un invariante. Llamaré ``impropias'' a las relaciones de divergencia en las que las \(B^{(\lambda)}\) pueden componerse de la manera indicada a partir de las expresiones lagrangianas y sus derivadas, y ``propias'' a todas las restantes.

Recíprocamente, si las relaciones de divergencia son combinaciones lineales de las dependencias (16), y por tanto ``impropias'', entonces la invariancia con respecto a \(\G_\sigma\) se sigue de la invariancia con respecto a \(\G_{\infty\rho}\); \(\G_\sigma\) se convierte en un subgrupo de \(\G_{\infty\rho}\). Así, las relaciones de divergencia correspondientes a un grupo finito \(\G_\sigma\) se vuelven impropias si y sólo si \(\G_\sigma\) es subgrupo de un grupo infinito respecto del cual \(I\) es invariante.

Por especialización de los grupos se sigue de esto la afirmación original de Hilbert. Por ``grupo de traslaciones'' se entiende el grupo finito
\[
y_i=x_i+\eps_i,\qquad v_i(y)=u_i(x);
\]
por tanto
\[
\D x_i=\eps_i,\qquad \D u_i=0,\qquad
\bd u_i=-\sum_\lambda\frac{\p u_i}{\p x_\lambda}\eps_\lambda.
\]
La invariancia con respecto al grupo de traslaciones dice, como es sabido, que en
\[
I=\int\!\cdots\!\int f\!\left(x,u,\frac{\p u}{\p x},\ldots\right)\dd x
\]
las \(x\) no aparecen explícitamente en \(f\). Las \(n\) relaciones de divergencia asociadas
\[
\sum_i\psi_i\frac{\p u_i}{\p x_\lambda}=\Div B^{(\lambda)}
\qquad(\lambda=1,2,\ldots,n)
\]
han de llamarse ``relaciones de energía'', porque las ``leyes de conservación'' \(\Div B^{(\lambda)}=0\) correspondientes al problema variacional corresponden a los ``teoremas de energía'', y las \(B^{(\lambda)}\) a las ``componentes de energía''. Así se obtiene: si \(I\) admite el grupo de traslaciones, entonces las relaciones de energía se vuelven impropias si y sólo si \(I\) es invariante con respecto a un grupo infinito que contiene al grupo de traslaciones como subgrupo.\footnote{Los teoremas de energía de la mecánica clásica, y asimismo los de la antigua ``teoría de la relatividad'' (donde \(\sum \dd x^2\) pasa a sí misma), son ``propios'', puesto que aquí no aparecen grupos infinitos.}

Un ejemplo de tales grupos infinitos lo da el grupo de todas las transformaciones de las \(x\) y de las transformaciones inducidas de las \(u(x)\) en las que sólo aparecen derivadas de las funciones arbitrarias \(p(x)\). El grupo de traslaciones surge por la especialización \(p^{(i)}(x)=\eps_i\). Debe quedar indeciso, sin embargo, si por este medio --- y por los grupos que surgen al cambiar \(I\) por un integral de frontera --- se han dado ya los grupos más generales de este tipo. Las transformaciones inducidas de la clase indicada surgen, por ejemplo, al someter las \(u\) a las transformaciones de coeficientes de una ``forma diferencial total'', es decir, de una forma
\[
\sum a\,\dd^2x_i+\sum b\,\dd x_i\dd x_k+\cdots,
\]
que contiene, además de las \(\dd x\), también diferenciales superiores. Transformaciones inducidas más especiales, en las que las \(p(x)\) aparecen sólo en primeras derivadas, están dadas por las transformaciones de coeficientes de formas diferenciales ordinarias \(\sum c\,\dd x_{i_1}\cdots \dd x_{i_\lambda}\); éstas son las que se han considerado usualmente.

Otro grupo de la clase indicada --- que, por la aparición del término logarítmico, no puede ser una transformación de coeficientes --- es, por ejemplo, el siguiente:
\[
y=x+p(x),\qquad
v_i=u_i+\lg(1+p'(x))=u_i+\lg\frac{\dd y}{\dd x};
\]
\[
\D x=p(x),\qquad
\D u_i=p'(x),\qquad
\bd u_i=p'(x)-u_i'p(x).
\]
Aquí las dependencias (16) se convierten en
\[
\sum_i\left(\psi_i u_i'+\frac{\dd\psi_i}{\dd x}\right)=0,
\]
y las relaciones de energía impropias se convierten en
\[
\sum_i\left(\psi_i u_i'+
\frac{\dd(\psi_i+\mathrm{const.})}{\dd x}\right)=0.
\]
Un integral invariante simplicísimo del grupo es
\[
I=\int \frac{e^{-2u_1}}{u_1'-u_2'}\,\dd x.
\]
El \(I\) más general se determina por integración de la ecuación diferencial de Lie (11)
\[
\bd f+\frac{\dd}{\dd x}(f\cdot \D x)=0,
\]
que, tras sustituir los valores de \(\D x\) y \(\bd u\), se convierte, tan pronto como se supone que \(f\) depende sólo de las primeras derivadas de las \(u\), en
\[
\frac{\p f}{\p x}p(x)
+\left\{\sum_i\frac{\p f}{\p u_i}
-\sum_i\frac{\p f}{\p u_i'}u_i'+f\right\}p'(x)
+\left\{\sum_i\frac{\p f}{\p u_i'}\right\}p''(x)=0
\]
(identidad en \(p(x),p'(x),p''(x)\)). Ya para dos funciones \(u(x)\), este sistema de ecuaciones tiene soluciones que contienen realmente las derivadas, a saber
\[
f=(u_1'-u_2')\,
\Phi\!\left(u_1-u_2,\frac{e^{-u_1}}{u_1'-u_2'}\right),
\]
donde \(\Phi\) denota una función arbitraria de los argumentos indicados.

Hilbert formula su afirmación diciendo que el fallo de los teoremas de energía propios es una característica de la ``teoría general de la relatividad''. Para que esta afirmación valga literalmente, la denominación ``relatividad general'' debe entenderse por tanto en un sentido más amplio que el usual, y extenderse también a los grupos precedentes que dependen de \(n\) funciones arbitrarias.\footnote{Con esto se confirma de nuevo la corrección de una observación de Klein: que el término ``relatividad'', usual en física, debe sustituirse por ``invariancia relativa a un grupo''. (Über die geometrischen Grundlagen der Lorentzgruppe. \emph{Jhrber. d. d. Math. Vereinig.} vol. 19, p. 287, 1910; reimpreso en la \emph{Physikalische Zeitschrift}.)}

% --- Paper 14 local macro support ---
\providecommand{\p}{\partial}
\providecommand{\frS}{\mathfrak S}
\providecommand{\frp}{\mathfrak p}
\providecommand{\fra}{\mathfrak a}
\providecommand{\frb}{\mathfrak b}
\providecommand{\frc}{\mathfrak c}
\providecommand{\frf}{\mathfrak f}
\providecommand{\frr}{\mathfrak r}
\providecommand{\frm}{\mathfrak m}
\providecommand{\frD}{\mathfrak D}
\providecommand{\calA}{\mathcal A}
\providecommand{\calB}{\mathcal B}
\providecommand{\calN}{\mathcal N}
% --- end Paper 14 local macro support ---

\section*{14. La teoría aritmética de las funciones algebraicas de una variable en su relación con las demás teorías y con la teoría de los cuerpos numéricos}
\begin{center}
\emph{Jahresbericht der Deutschen Mathematiker-Vereinigung} 38 (1919), pp. 182--203
\end{center}

El informe siguiente nació del deseo de contar con un eslabón de unión entre las distintas teorías de las funciones algebraicas; al mismo tiempo puede considerarse como un complemento del informe de Brill--Noether,\footnote{El desarrollo de la teoría de las funciones algebraicas en la época antigua y en la reciente. \emph{Jahresbericht der Deutschen Mathematiker-Vereinigung} vol. 3 (1894) [citado como Brill--Noether].} en el cual falta esencialmente la discusión de la teoría aritmética.\footnote{Con excepción de una discusión de la investigación de Kronecker sobre el discriminante. (\emph{J. f. M.} vol. 91.)}

Para la discusión de las teorías particulares, además de Brill--Noether, entran en consideración los siguientes artículos de la \emph{Enzyklopädie}, junto con la bibliografía allí indicada:

Para la teoría de Riemann: Wirtinger (II B 2), núms. 1--19;

para la teoría de Weierstrass: Wirtinger (II B 2), núms. 23, 32, 33, 34;

para la teoría algebraico-geométrica de Brill--Noether: Wirtinger (II B 2), núms. 21--31; Berzolari (III C 4), núms. 23--38;

para la teoría aritmética de Dedekind--Weber y de Hensel--Landsberg: Hensel (II C 5), [citado como Hensel];\footnote{Las partes relativas a Dedekind--Weber proceden de A. Ostrowski.}

más generalmente, para la teoría aritmética de las cantidades algebraicas: Landsberg (I C 5), y para la teoría aritmética de las funciones algebraicas de varias variables: Jung (II C 6).

Para la teoría de cuerpos numéricos compárese el ``Zahlbericht'' de Hilbert\footnote{La teoría de los cuerpos de números algebraicos. \emph{Jahresbericht der Deutschen Mathematiker-Vereinigung} vol. 4 (1894/95).} y las lecciones de teoría de números de Dirichlet--Dedekind [citadas como Dedekind--teoría de números].\footnote{Braunschweig, Vieweg, 4.ª ed. 1894.}

\subsection*{Índice}
\begin{enumerate}
\item Fundamentos de las distintas teorías.
\item Paralelismo entre números algebraicos y funciones algebraicas (elementos conjugados, teoremas de bases, concepto de módulo).
\item Paralelismo y diferencias en la teoría de ideales de los números y de las funciones algebraicas.
\item Relación de la fundamentación aritmética del desarrollo en serie con la teoría de ideales.
\item Comparación de la teoría aritmética de las funciones algebraicas con la teoría geométrica. Concepto distinto de invariancia.
\item Continuación de la comparación. Puntos singulares.
\item Comparación entre el teorema del residuo y los teoremas de la teoría de ideales.
\item Cuestiones trascendentes para funciones algebraicas y números algebraicos.
\end{enumerate}

\subsection*{1. Fundamentos de las distintas teorías}
La teoría de Riemann es la única puramente trascendente, apoyada en teoremas de existencia (principio de Dirichlet, métodos de Schwarz--Neumann). En ella los integrales de las funciones algebraicas son lo primario; las funciones algebraicas mismas, lo secundario. Sobre la superficie de Riemann las funciones se caracterizan por sus ceros y polos, lo que corresponde a los polígonos (respectivamente divisores) de la teoría aritmética; Weierstrass, en cambio, habla también de las funciones mismas junto con sus ceros y polos. El problema principal, aparte de los teoremas de existencia, es la construcción de todas las funciones con polos dados, lo cual se realiza mediante integrales de primera y segunda especie; además, la pregunta por el número de constantes arbitrarias que todavía intervienen cuando se prescriben los puntos de infinito. Esta última pregunta la responde el teorema de Riemann--Roch; a ella corresponde, en la teoría aritmética, la pregunta por la dimensión de una clase de polígonos (respectivamente de divisores), y en la teoría algebraico-geométrica la pregunta por la dimensión de una serie lineal completa, cuya respuesta viene dada también aquí por el teorema de Riemann--Roch. La formación efectiva de las funciones se hace, en la teoría aritmética, mediante los teoremas de bases; en la teoría algebraico-geométrica, mediante el teorema del residuo. En particular, a los ceros de los integrandos de primera especie corresponde la clase diferencial en la teoría aritmética, y la doctrina de los grupos especiales (es decir, la serie de grupos de puntos cortada por las curvas adjuntas de orden \((n-3)\), las curvas \(\varphi\)) en la teoría algebraico-geométrica.

Así como la teoría de Riemann precedió históricamente a las demás teorías, aunque la demostración rigurosa de los teoremas de existencia sólo se obtuvo más tarde, también posteriormente, con sus auxilios trascendentes, se resolvieron cuestiones algebraicas que hasta hoy no son accesibles a un tratamiento puramente algebraico o aritmético; debe mencionarse aquí ante todo la teoría de Hurwitz de las correspondencias singulares.\footnote{A. Hurwitz, Über algebraische Korrespondenzen und das erweiterte Korrespondenzprinzip. \emph{Math. Ann.} vol. 28. (Compárese Brill--Noether, sección 10.)}

En el último apartado se discutirán analogías con las cuestiones trascendentes en la teoría de cuerpos de números algebraicos.

Para las demás teorías, las funciones algebraicas son lo primario y los integrales lo secundario. A este respecto, la teoría de Weierstrass descansa sobre los fundamentos trascendentes del desarrollo en series de potencias y del teorema del residuo; pero después prosigue algebraicamente, mediante la indicación explícita de todas las funciones que entran en consideración.

También la teoría aritmética en la forma de Hensel--Landsberg utiliza los fundamentos trascendentes del desarrollo en serie. Pero, al asignar divisores a los puntos individuales sobre esa base, continúa entonces de manera puramente aritmética y puede considerarse, por tanto, como una fusión de Weierstrass y Dedekind--Weber. En contraste con estos últimos, suministra a la vez métodos para la formación efectiva de todas las funciones de base de que se trata. En detalle presenta una serie de simplificaciones frente a Dedekind--Weber, esencialmente por la introducción de divisores ``fraccionarios''. Recientemente Hensel\footnote{\emph{Mathematische Zeitschrift} vol. 4; la fundamentación sirve también de base a su informe.} ha fundamentado aritméticamente la teoría del desarrollo en series (salvo por un método de mayorantes para la demostración de convergencia), con lo cual la teoría se ha vuelto casi puramente aritmética.

La teoría algebraico-geométrica de Brill--Noether presupone sólo el concepto de ``puntos sucesivos'' de una rama, que se reduce al desarrollo en serie en un punto ordinario\footnote{Para puntos ``singulares'' sucesivos compárese Brill--Noether, sección 6, especialmente núms. 10 y 11. Los puntos ordinarios sucesivos son, por ejemplo, los puntos de intersección de la tangente con la curva, o puntos de contacto de orden superior (puntos de ramificación).} y corresponde al concepto de elemento de Weierstrass; después continúa con métodos algebraico-racionales, esencialmente los de la teoría ternaria de la eliminación; también ella llega hasta representaciones explícitas.

La teoría originaria, puramente aritmética, de Dedekind--Weber\footnote{\emph{J. f. M.} 92 (citado como D.--W.).} es afín a la de Riemann en cuanto sus teoremas vuelven a tener el carácter de demostraciones de existencia; conserva la plena pureza del método.\footnote{Por esta razón es también más adecuada para compararla con las demás teorías que, por ejemplo, la teoría de Hensel--Landsberg, que aparecerá sólo ocasionalmente en las comparaciones siguientes. --- La teoría de los integrales abelianos y de su inversión, en cuanto presupone continuidad, no se trata, sin embargo, ni en D.--W. ni en el ulterior desarrollo de la teoría algebraica (M. Noether, \emph{Ann.} 37).} La fundamentación aritmética de la teoría, donde las variables desempeñan sólo el papel de indeterminadas, corre paralelamente a la teoría de los números algebraicos, como se explicará más precisamente en 2.\footnote{Una comparación de la teoría de Hensel--Landsberg con la teoría de números, especialmente con la teoría de los números \(p\)-ádicos, se encuentra en Hensel.} Esta fundamentación proporciona el medio de introducir aritméticamente el concepto de punto (compárese Hensel núm. 10a) y así establecer el concepto de la superficie de Riemann absoluta sin ninguna consideración de continuidad. La teoría adquiere con ello un carácter formal; por ejemplo, también el concepto de diferencial se introduce de manera puramente formal. Sobre la base obtenida por el concepto de punto puede demostrarse el estrecho parentesco con los métodos y conceptos de la teoría algebraico-geométrica, que, aparte del concepto de elemento, está también libre de auxilios trascendentes; esto se hará en los apartados siguientes.

\subsection*{2. Paralelismo entre números algebraicos y funciones algebraicas. (Elementos conjugados, teoremas de bases, concepto de módulo)}
Un cuerpo de números algebraicos se define como un cuerpo de grado \(n\) sobre el cuerpo \(R\) de los números racionales; un cuerpo de funciones algebraicas de una variable, como un cuerpo de grado \(n\) sobre el cuerpo \(K(z)\) de todas las funciones racionales de una indeterminada \(z\) con coeficientes complejos arbitrarios. Por tanto, para ambos cuerpos valen en común todos los teoremas que abstraen completamente de la naturaleza especial del cuerpo base: los teoremas sobre bases, normas, trazas y discriminantes.\footnote{D.--W., § 1 y 2; Hensel núm. 4 (con nota 5).} Dedekind--Weber introduce todos estos conceptos sin utilizar elementos conjugados. No obstante, debe observarse que el concepto de cuerpo conjugado puede definirse de manera enteramente formal, como mostró Steinitz siguiendo a Kronecker.\footnote{Steinitz, Algebraische Theorie der Körper. (\emph{J. f. M.} 137, § 6 y 8.)} En efecto, si \(K\) es un cuerpo arbitrario y \(\varphi(x)\) un polinomio en \(x\) irreducible sobre \(K\), entonces el sistema de clases de restos módulo \(\varphi(x)\) forma un cuerpo. Si se asigna a la clase de restos de \(x\) un elemento \(\xi\), entonces, bajo la asignación isomorfa, a la clase de restos representada por el polinomio \(f(x)\) corresponde el elemento \(f(\xi)\). Puesto que todos los polinomios divisibles por \(\varphi(x)\) representan la clase cero, \(\varphi(\xi)\) queda asignado a la clase cero; así, el elemento \(\xi\) puede considerarse simbólicamente como una raíz de \(\varphi(x)=0\) y genera el cuerpo algebraico \(K(\xi)\). La repetición del procedimiento conduce a los \(n\) cuerpos conjugados \(K(\xi),K(\xi'),\ldots,K(\xi^{(n-1)})\). Por consiguiente, también para las funciones algebraicas puede hablarse en particular de elementos conjugados sin abandonar la teoría puramente aritmética.

Los dos cuerpos base \(R\) y \(K(z)\) tienen en común la propiedad de que en ellos está definido el sistema de todas las cantidades enteras: los enteros racionales, respectivamente las funciones racionales enteras de una indeterminada. Estas cantidades enteras tienen la propiedad característica de que dos cualesquiera \(a\) y \(b\) poseen un máximo común divisor \(d\) representable en la forma \(ma+nb\), donde también \(d,m,n\) son cantidades enteras de \(R\), respectivamente de \(K(z)\); y \(d\) es precisamente el número absolutamente más pequeño, respectivamente el polinomio de grado mínimo, representable en esa forma. Sólo en esta propiedad, que trae consigo la descomposición única en factores primos para los enteros racionales y para las funciones de una indeterminada, descansan los teoremas sobre las cantidades enteras del cuerpo algebraico o cuerpo superior. El razonamiento que demuestra la existencia del máximo común divisor da también la existencia de un sistema fundamental o de una base de las cantidades enteras del cuerpo algebraico. Si se consideran sólo aquellas bases del cuerpo que consisten en cantidades enteras, su discriminante se convierte en un entero racional, respectivamente en una función racional entera de una indeterminada. Toda base que corresponde al número absolutamente más pequeño, respectivamente a la función de grado mínimo, forma un sistema fundamental.

Una comprensión más precisa de las cuestiones de bases, también para cuerpos más generales, se obtiene por la introducción del concepto de módulo; lo formularemos de manera que incluya las definiciones que aparecen de forma diversa en la literatura según la naturaleza de los elementos. Un módulo --- con respecto a un dominio básico de integridad \(\frS\) --- es un sistema de elementos tal que la diferencia de dos elementos pertenece de nuevo al sistema, y lo mismo el producto de un elemento por una cantidad arbitraria de \(\frS\). (Si \(\frS\) consiste en los enteros racionales, entonces la segunda exigencia es consecuencia de la primera.)\footnote{El concepto de módulo se debe a Dedekind, quien introdujo por primera vez el módulo numérico (\(\frS=\) enteros racionales) en la segunda edición de la teoría de números; también el módulo de formas lineales enteras está implícitamente contenido allí. En Dedekind--Weber aparecen los ``módulos de funciones'', cuyos elementos consisten en funciones algebraicas, mientras que \(\frS\) consiste en todas las funciones racionales enteras de una indeterminada. --- El módulo de formas homogéneas en \(n\) indeterminadas, más generalmente de polinomios, aparece primero en Hilbert (\emph{Annalen} 36); \(\frS\) consiste en todos los polinomios, y como los elementos son asimismo polinomios, este módulo cae bajo nuestra definición de ideal. Los ``sistemas de módulos'' de polinomios de Kronecker parten de la base, no de la definición abstracta.}

Todo módulo que posee una base de finitos elementos por medio de la cual todo elemento puede representarse linealmente con coeficientes de \(\frS\) se llama finito (finitamente generado); en particular, todo módulo de un término es por tanto de la forma \(c\cdot\alpha\), donde \(\alpha\) denota el elemento de base y \(c\) recorre todas las cantidades de \(\frS\). \emph{Un módulo cuyos elementos pertenecen ellos mismos a \(\frS\) se llama un ideal en \(\frS\);} de ahí se sigue el concepto de ideal finito, y en particular el de ideal de un término (ideal principal).

Ahora bien, todos los teoremas de bases para cantidades enteras e ideales descansan en el siguiente teorema:

\emph{Si \(M\) es un módulo con respecto a \(\frS\), cuyos elementos consisten en formas lineales en \(n\) indeterminadas con coeficientes de \(\frS\), y si todo ideal en \(\frS\) es finito, entonces \(M\) es también un módulo finito. Si en particular todo ideal en \(\frS\) es un ideal principal, entonces \(M\) posee una base de formas linealmente independientes.}\footnote{El teorema aparece en la literatura separadamente para los distintos dominios de coeficientes \(\frS\). Para enteros racionales compárese, por ejemplo, Hilbert, Zahlbericht § 3 y 4; para polinomios en varias indeterminadas: Hilbert, Über die vollen Invariantensysteme, final del § 2 (\emph{Math. Ann.} 42).}

En efecto, si los elementos de \(M\) están dados por \(l(x)=a_1x_1+\cdots+a_nx_n\), entonces la totalidad de los coeficientes \(a_1\) recorre un ideal de \(\frS\), que por hipótesis es de la forma \(b_1a^{(1)}+\cdots+b_\rho a^{(\rho)}\). La forma lineal perteneciente a \(M\),
\[
m(x)=l(x)-b_1l^{(1)}(x)-\cdots-b_\rho l^{(\rho)}(x),
\]
depende por tanto sólo de \((n-1)\) indeterminadas, de donde la afirmación se sigue por un número finito de repeticiones. Si en cada caso \(\rho=1\), entonces la base consiste en \(n\) formas en, respectivamente, \(n,n-1,\ldots,1\) indeterminadas, lo que da la independencia lineal de las formas no nulas.

El teorema sobre el sistema fundamental se deduce de aquí en los distintos casos como sigue:

Si el cuerpo algebraico se obtiene por adjunción del elemento algebraico entero \(\delta\) al cuerpo base, y si \(\frS\) denota la totalidad de todas las cantidades enteras del cuerpo base, entonces toda cantidad entera del cuerpo superior admite la representación (compárese, por ejemplo, Zahlbericht § 3)
\[
\omega=\frac{a_1\delta^{\,n-1}+a_2\delta^{\,n-2}+\cdots+a_n}{\Delta(\delta)},
\]
donde los \(a_i\) son cantidades de \(\frS\) y \(\Delta(\delta)\) denota el discriminante de \(\delta\). Mediante la sustitución
\[
x_i=\frac{\delta^{\,n-i}}{\Delta(\delta)}
\]
el módulo de todas las cantidades enteras se transforma en un módulo de formas lineales; de ahí que la existencia del sistema fundamental se siga en todos los casos en que todo ideal de \(\frS\) es finito. El mismo razonamiento muestra más generalmente que todo ideal en el dominio de las cantidades enteras del cuerpo superior es finito. Ahora bien, en particular para los enteros racionales y para las funciones de una indeterminada todo ideal es principal; de aquí se sigue que los sistemas fundamentales en cuerpos de números algebraicos y en cuerpos de funciones algebraicas de una indeterminada consisten exactamente en \(n\) cantidades cuando \(n\) indica el grado. Y puesto que en estos cuerpos todo ideal es finito por lo anterior, se obtiene al mismo tiempo el sistema fundamental para cuerpos relativos, que en general constará de más cantidades de las que indica el grado relativo.\footnote{Según las investigaciones de Steinitz sobre módulos de formas lineales, donde \(\frS\) consiste en los enteros de un cuerpo de números algebraicos (finito), bastan \(\nu+1\) elementos de base cuando \(\nu\) denota el grado relativo (\emph{Math. Ann.} 71, § 4).} Además, si \(\frS\) denota el dominio de todos los polinomios en varias indeterminadas, todo ideal en \(\frS\) es también finito;\footnote{Por el teorema de Hilbert sobre la base de módulos (\emph{Math. Ann.} 36). Los módulos de polinomios aquí llamados ideales por coherencia se denominan habitualmente ``módulos de formas'' o simplemente ``módulos''. El teorema de Hilbert de que tal módulo \(N\) es finito puede reducirse también al teorema sobre formas lineales. Sea, en efecto, \(f\) un polinomio de grado \(n\) contenido en \(N\) en las \(r+1\) indeterminadas \(x_1,\ldots,x_r,x_{r+1}\), que contenga el término \(x_{r+1}^{\,n}\) (lo cual puede lograrse siempre mediante una transformación lineal); entonces todos los polinomios de \(N\) son congruentes módulo \(f\) con los representables en la forma \(a_1(x)x_{r+1}^{\,n-1}+\cdots+a_n(x)\); por tanto, mediante \(x_{r+1}^{\,n-i}=X_i\), forman un módulo de formas lineales para el cual puede suponerse finito todo ideal en el dominio de \(r\) indeterminadas, ya que esto ocurre para una indeterminada.} por tanto se obtiene también la existencia del sistema fundamental para cuerpos de funciones algebraicas de varias indeterminadas, y éste constará de nuevo de más cantidades de las que indica el grado.

\subsection*{3. Paralelismo y diferencias en la teoría de ideales de los números y de las funciones algebraicas}
La demostración del teorema fundamental de la teoría de ideales, según el cual todo ideal puede representarse de manera única como un producto de potencias de ideales primos, puede llevarse a cabo conjuntamente para números y funciones algebraicas, únicamente sobre la base de la propiedad de que, para los enteros racionales y para las funciones de una indeterminada, todo ideal es principal.\footnote{Compárese una observación en la introducción a Hurwitz: Über die Theorie der Ideale (\emph{Gött. Nachr.} 1894). En el manual de álgebra de Weber, la demostración se lleva a cabo en paralelo para ambos cuerpos, siguiendo a Kronecker, mediante la introducción de funcionales. [La conexión entre la teoría de Kronecker y la teoría de ideales está mediada por el ``teorema de Gauss ampliado''; Hurwitz, loc. cit.]} El punto principal de la demostración reside en mostrar que, de la divisibilidad de un ideal \(\frc\) por un ideal \(\fra\) (es decir, de que todo elemento de \(\frc\) esté contenido en \(\fra\)), se sigue la representación producto \(\frc=\fra\frb\).\footnote{Subrayado especialmente por Dedekind (teoría de números), Suplemento XI, teorema VII, § 177; compárese la nota en p. 554. El teorema usa el hecho de que se trata de cuerpos algebraicos; ya no vale para ideales (módulos) de polinomios, donde el mínimo común múltiplo ocupa el lugar del producto.} En Hurwitz, esta demostración descansa en el ``teorema de Gauss ampliado'', según el cual una cantidad entera \(\omega\) que divide a cada coeficiente del producto de dos polinomios \(\varphi\) y \(\psi\) divide también el producto de todos los coeficientes de \(\varphi\) y \(\psi\); en Dedekind descansa en un teorema de módulos equivalente a ese teorema.\footnote{Teoría de números, Supl. XI; teorema VI, § 173. Dedekind señala la equivalencia de los dos teoremas (Über die Begründung der Idealtheorie, \emph{Göttinger Nachr.} 1895).}

Si se continúa, el concepto de módulo complementario (Hensel núm. 12), y a partir de él el ideal de ramificación (ideal fundamental, diferente) \(\frD\), puede definirse conjuntamente (su norma se convierte en el discriminante del cuerpo \(D\)); y vale el teorema fundamental de que el ideal de ramificación \(\frD\) es el \emph{máximo común divisor de todos los ideales principales \(F'(\delta)\)}; aquí \(\delta\) recorre todas las cantidades enteras del cuerpo, y \(F(\delta)=0\) denota en cada caso la ecuación irreducible correspondiente. De aquí se sigue que el \emph{discriminante del cuerpo contiene todas y sólo aquellas primas racionales, respectivamente factores lineales en \(z\), que son divisibles por el cuadrado de un ideal primo}. Para \(F'(\delta)\) mismo se obtiene la representación
\[
F'(\delta)=\frf\cdot\frD,
\]
y por paso a las normas
\[
\Delta(\delta)=R^2\cdot D;
\]
aquí \(\frf\) denota el conductor del anillo derivado de \(\delta\) (orden, dominio íntegro), es decir, \(\frf\) consiste en la totalidad de las cantidades (enteras) que, después de multiplicarse por cualquier cantidad entera, son divisibles por el módulo \([1,\delta,\ldots,\delta^{n-1}]\). \(R^2=\operatorname{Norm}\frf\) representa el cuadrado del determinante de sustitución que lleva la base \(1,\delta,\ldots,\delta^{n-1}\) a una base de las cantidades enteras.\footnote{Dedekind, Über die Diskriminante endlicher Körper (\emph{Abhandl. der Gött. Ges. d. Wissensch.} vol. 29, 1882). Las definiciones dadas allí para números algebraicos pueden trasladarse directamente a funciones algebraicas; también las demostraciones pueden trasladarse casi exactamente. El orden de la demostración en D.--W. es distinto: véase más abajo. Para los teoremas citados compárese también Zahlbericht, § 31 y 32. (Allí \(F'(\delta)\) se llama la ``diferente'' del entero \(\delta\).)}

Las \emph{diferencias} en el desarrollo ulterior están condicionadas por propiedades especiales de los cuerpos base. Así, en el caso de los números algebraicos, el hecho de que los \emph{elementos del cuerpo base sean al mismo tiempo exponentes} conduce a la posibilidad de que el ideal de ramificación sea divisible por una potencia más alta que la potencia \((e-1)\)-ésima de un ideal primo \(\frp\) que aparece en \(p\) a la \(e\)-ésima potencia (a saber, cuando \(e\) es divisible por \(p\)); esto conduce a los conceptos adicionales de cuerpo de inercia y cuerpo de ramificación,\footnote{Compárese Zahlbericht, § 39--45.} que no tienen análogo en la teoría de las funciones algebraicas. Pero sobre todo, debido a la \emph{existencia de infinitas constantes} en el cuerpo base, para las funciones algebraicas desaparece la finitud del número de clases de ideales, y con ella todas las cuestiones conectadas con la determinación del número de clases. No obstante, en cierto sentido el dominio de todas las formas primas (trascendentes)\footnote{Klein, Zur Theorie der Abelschen Funktionen (\emph{Ann.} 36, 1890).} puede considerarse como un análogo trascendente del cuerpo de clases; pues aquí cada punto es generado por una sola forma, un ideal principal.

Por otra parte, la existencia de infinitas constantes en el cuerpo base \(K(z)\) conduce, para las funciones algebraicas, a teoremas y cuestiones que no tienen análogo para los números algebraicos. Primero se sigue de ello cierta simplificación en la disposición de la demostración; por ejemplo en la demostración dada por D.--W. de la descomposición única de los ideales en ideales primos, donde el ``teorema de Gauss ampliado'', o un teorema equivalente, puede evitarse usando el hecho de que toda forma lineal \((z-c)\) tiene infinitas clases de restos, a saber, todas las constantes (mientras que el número de clases de restos módulo un número primo es finito).\footnote{D.--W. § 9. Compárense las notas en pp. 212 y 213.}

El hecho ulterior de que todo polinomio con funciones racionales enteras de \(z\) como coeficientes se descompone en factores lineales módulo \((z-\alpha)\) (lo que no sucede para un polinomio entero módulo \(p\)) conduce a resultados genuinamente nuevos; este hecho sigue valiendo si, en lugar de todas las constantes, se toman sólo todos los números algebraicos.\footnote{D.--W. señalan en la introducción que en este caso toda la teoría permanece intacta.} De aquí se sigue, a saber, que todo ideal primo es un ideal de primer grado;\footnote{D.--W., § 9, núm. 7.} y además que el discriminante del cuerpo \(D\) se convierte en la intersección de todos los discriminantes \(\Delta(\delta)\), donde \(\delta\) recorre todas las cantidades enteras;\footnote{D.--W., p. 224.} ambas cosas están en contraste con los números algebraicos. De esto se sigue a su vez que el ideal de ramificación es la intersección de todos los ideales principales \(F'(\delta)\).

Otras diferencias están condicionadas por el hecho de que, para los números algebraicos, el cuerpo base \(R\) de los números racionales es algo absolutamente distinguido, a saber, el cuerpo primo contenido en el cuerpo (es decir, el cuerpo derivado de la unidad), mientras que el cuerpo base \(K(z)\) es algo relativo. Toda función no constante \(\xi\) puede elegirse como variable independiente; el cuerpo se convierte entonces en un cuerpo algebraico sobre \(K(\xi)\). El teorema de que todo ideal primo \(\frp\) es de primer grado, y por tanto toda función es congruente módulo \(\frp\) con una constante, unido a la posibilidad de tomar una función arbitraria como variable independiente, conduce al concepto más importante que va más allá de los números algebraicos: el concepto de punto en forma puramente aritmética (compárese Hensel 10a). ``Cada punto está representado aquí por un cuerpo de constantes isomorfo al cuerpo de funciones.''\footnote{Como es sabido, dos cuerpos se llaman ``isomorfos'' si pueden ponerse en correspondencia biunívoca de tal modo que también correspondan la suma, la diferencia, el producto y el cociente.} La totalidad de los puntos así definidos forma la ``superficie de Riemann absoluta''. Esta definición de punto depende sólo del cuerpo, no de una variable especial. La demostración de existencia, sin embargo, se lleva a cabo distinguiendo alguna función como variable independiente \(z\); un punto \(P\) es generado por un ideal primo \(\frp\) al asignar a cada función su resto constante módulo \(\frp\); así, en particular, las funciones divisibles por \(\frp\) se anulan en \(P\), y recíprocamente. Cuando \(\frp\) recorre todos los ideales primos en \(z\), se obtienen con ello todos los puntos en los que \(z\) es finito; los puntos restantes se obtienen introduciendo la variable independiente \(\xi=1/z\) y corresponden a los ideales primos que dividen el ideal principal \(\xi\). Naturalmente, tampoco los conceptos ligados al concepto de punto pueden tener análogo para los números algebraicos, como el concepto de polígono, de clase de polígonos, de dimensión de una clase de polígonos, etc.

Debe señalarse asimismo que la libre elección de la variable independiente conduce a otra interpretación de la descomposición \(F'(\vartheta)=\frf\cdot\frD\). En efecto, si se considera \(\vartheta\) como variable independiente, entonces se obtiene correspondientemente
\[
\left(\frac{\p F}{\p z}\right)=\frf\cdot\fra,
\]
puesto que el anillo derivado de \(\vartheta\) estaba definido simétricamente respecto de \(z\) y \(\vartheta\); y \(\frf\) se vuelve el máximo común divisor de los dos ideales principales \(\frac{\p F}{\p z}\) y \(\frac{\p F}{\p\vartheta}\); el conductor \(\frf\) se convierte al mismo tiempo en el ``ideal de puntos dobles'' respecto de \(\vartheta,z\).\footnote{D.--W. p. 263.}

\subsection*{4. Relación de la fundamentación aritmética del desarrollo en serie con la teoría de ideales}
En Hensel--Landsberg, el auxilio trascendente del desarrollo en serie y el concepto de punto tomado de la teoría de funciones ocupan el lugar de la teoría de ideales. La fundamentación aritmética del desarrollo en serie para números y funciones algebraicas dada recientemente por Hensel\footnote{Eine neue Theorie der algebraischen Zahlen (\emph{Math. Zeitschr.} vol. 2, 1918) y Neue Begründung der arithmetischen Theorie der algebraischen Funktionen einer Variablen (\emph{Math. Zeitschr.} vol. 4, 1919). Para una exposición más detallada compárese Hensel núm. 44 ss.} puede considerarse por tanto como un equivalente de la teoría de ideales. De hecho, esta fundamentación aritmética consiste en pasar, para cada ideal primo, a un campo de extensión (trascendente) en el que ese ideal primo se vuelve principal, de modo que valen los teoremas usuales de descomposición; y basta realizar esta consideración para todos los ideales primos que aparecen en un número primo \(p\), respectivamente en una función lineal \(z-a\).

Primero se introduce un cuerpo de extensión trascendente formado por todos los desarrollos en serie en potencias enteras de \(z-a\), respectivamente de \(p\) (números \(p\)-ádicos), con a lo sumo finitas potencias negativas en cada caso. En este cuerpo, la ecuación irreducible que define el cuerpo de números algebraicos o el cuerpo de funciones se descompone en el producto de \(\rho\) factores irreducibles, correspondientes a la descomposición de \(p\), respectivamente de \(z-a\), en el producto de los \(\rho\) ideales ``primarios'' \(\mathfrak q_1,\ldots,\mathfrak q_\rho\). La representación ulterior de cada ideal primario como potencia de un ideal primo,
\[
\mathfrak q_i=\frp_i^{\,\ell_i},
\]
corresponde en Hensel a la introducción de otro cuerpo, algebraico respecto del cuerpo de extensión trascendente. En el caso de las funciones algebraicas este cuerpo puede definirse por una ecuación pura; en el caso de los números algebraicos, más generalmente, por una ecuación algebraicamente resoluble. La simplificación para las funciones algebraicas se debe a que aquí no aparecen cuerpos de inercia ni cuerpos de ramificación (compárese núm. 3).

La ``serie terminante'' en un elemento primo \(\pi\)\footnote{Neue Begründung \(\ldots\), fórmula (11).} que aparece en estas investigaciones (y que todavía no requiere una demostración trascendente de convergencia)
\[
v=c_0+c_1\pi+\cdots+c_{\sigma-1}\pi^{\sigma-1}-v_\sigma\pi^\sigma,
\]
donde \(c_0,\ldots,c_{\sigma-1}\) son constantes y \(v_\sigma\) denota un elemento entero del cuerpo, se encuentra en forma exactamente correspondiente en D.--W.\footnote{D.--W., p. 231 y § 15.} Allí se deriva de la teoría de ideales y da la base para definir el orden de anulación (respectivamente de hacerse infinita) de una función en un punto.

\subsection*{5. Comparación de la teoría aritmética de las funciones algebraicas con la teoría geométrica. Concepto distinto de invariancia}
``Invariante'' significa cualquier concepto caracterizado únicamente por el cuerpo. En la teoría aritmética esta invariancia está dada por lo general ya en la definición, que se enuncia respecto de todos los elementos del cuerpo, no respecto de alguna base o variable distinguida; el concepto de punto dado arriba es un ejemplo. En cambio, las teorías restantes parten de alguna variable o base especial y muestran después la independencia del concepto respecto de esa elección especial; la invariancia se convierte aquí en invariancia bajo transformación birracional.

En efecto, supóngase que el cuerpo surge, por una parte, adjuntando el elemento algebraico \(s\) a \(K(z)\), y por otra adjuntando \(\sigma\) a \(K(\xi)\), donde \(f(z,s)=0\), respectivamente \(F(\xi,\sigma)=0\), son las ecuaciones irreducibles de \(s\), respectivamente de \(\sigma\); o, todavía más generalmente, adjuntando un número finito de elementos a \(K(\eta)\), con las ecuaciones irreducibles correspondientes
\[
g_1(\eta,\tau_1)=0,\qquad g_2(\eta,\tau_1,\tau_2)=0,\ldots .
\]
Entonces, por definición, todos los elementos pueden expresarse racionalmente por \(z\) y \(s\), así como por \(\xi\) y \(\sigma\), y también por \(\eta,\tau_1,\tau_2,\ldots\). En particular, \(z\) y \(s\) son por tanto expresables racionalmente por \(\xi\) y \(\sigma\), y recíprocamente; asimismo \(z\) y \(s\) racionalmente por \(\eta,\tau_1,\tau_2,\ldots\), y recíprocamente. La curva plana \(f(z,s)=0\) pasa entonces por ``transformación birracional'' a la curva plana \(F(\xi,\sigma)=0\), o también a la curva en el espacio de las variables \(\eta,\tau_1,\tau_2,\ldots\) definida por \(g_1(\eta,\tau_1)=0; g_2(\eta,\tau_1,\tau_2)=0,\ldots\). La invariancia aparece por tanto en el hecho de que una definición enunciada para una curva especial permanece válida para todas las curvas obtenidas de ella por transformación birracional (en un espacio de un número arbitrario de dimensiones), para todas las curvas de la ``clase'' en el sentido de Riemann. Como ``punto'' se define aquí un sistema coherente de valores \(z=a,s=b\), de modo que \(f(a,b)=0\); bajo una transformación birracional este punto pasa en general de nuevo a un sistema coherente de valores \(\alpha,\beta\) de \(F(\xi,\sigma)=0\), pero en casos especiales, cuando \(a,b\) era un punto ``singular'', puede corresponder a varios ``puntos'' sobre \(F(\xi,\sigma)=0\); y siempre pueden indicarse curvas sobre las que un punto ``singular'' de \(f(z,s)=0\) corresponde a un número finito de puntos simples, situados separadamente o sucesivamente (resolución de singularidades).\footnote{Compárese, por ejemplo, Brill--Noether, sección 6. Información más precisa sobre puntos ``singulares'' aparece en la sección siguiente.} Del mismo modo, un grupo de puntos pasa de nuevo a un grupo de puntos; y el número de sus puntos es invariante tan pronto como se cuenta un punto singular como tantos puntos cuantos puntos simples le corresponden bajo una transformación adecuada. Como toda curva puede transformarse birracionalmente en una con puntos ``múltiples ordinarios'',\footnote{Ibid.} la teoría geométrica desarrolla sus conceptos sólo para tales curvas especiales y después demuestra la invariancia bajo toda transformación birracional, incluida una que conduzca a los puntos singulares más generales.

En particular, puede introducirse como curva distinguida de la ``clase'' --- excluido el caso hiperelíptico --- la llamada ``curva normal de las \(\varphi\)'' en el espacio de \((p-1)\) dimensiones, que no tiene puntos singulares; aquí las \(\varphi\) son las \(p\) ``curvas adjuntas de orden \((n-3)\)'' para la ecuación \(f(z,s)=0\) hecha homogénea. A una transformación birracional de \(f(z,s)=0\) en \(F(\xi,\sigma)=0\) corresponde una transformación lineal de las \(\varphi\), de modo que se trata de invariancia en sentido proyectivo. La representación de todas las funciones del cuerpo como funciones racionales de las \(\varphi\) se denomina ``representación invariante'' en la teoría geométrica.\footnote{Compárese, por ejemplo, Brill--Noether, sección 8. Para funciones algebraicas de varias variables no se ha tratado la cuestión de si la invariancia puede reducirse a tal invariancia bajo transformación lineal.}

El sistema lineal de las \(\varphi\), o también de los ``grupos especiales'' cortados por ellas, pasa por tanto a sí mismo bajo toda transformación birracional, y en consecuencia queda caracterizado sólo por el cuerpo. Corresponde a la clase diferencial en la teoría aritmética, donde la invariancia se muestra de nuevo directamente en la definición, dada por propiedades respecto de cada función del cuerpo como variable independiente.

\subsection*{6. Continuación de la comparación. Puntos singulares}
Con la distinta concepción del concepto de punto se relaciona el hecho de que los ``puntos singulares'' de las curvas, que en la teoría algebraico-geométrica presentan dificultades especiales, no aparezcan en absoluto en la fundamentación de la teoría aritmética; y, por consiguiente, que tampoco las ``funciones adjuntas'' sean necesarias para construir la teoría aritmética, como sí sucede en la teoría geométrica, sino que pertenezcan a partes especiales y superiores de la teoría. Por otra parte, en la fundamentación de la teoría geométrica no aparecen los puntos de ramificación; de hecho, éstos desaparecen también en la representación de los integrales en la forma normal de Aronhold (Hensel núm. 27, fórmula 80).

La definición de punto singular dada en 5 puede formularse así: \(f(z,s)=0\) (respectivamente una curva en un espacio superior) posee en \(z=a\), \(s=b\) (respectivamente \(z=a\), \(s_i=b_i\)) un punto singular si este sistema de valores es asumido por más de un punto de la superficie de Riemann absoluta (punto en sentido aritmético), o si es asumido en orden superior en uno o varios puntos. En este último caso se trata de un punto del polígono de ramificación tanto de \(z\) como de \(s\); corresponde a los ``puntos sucesivos'' de la teoría geométrica; el punto se convierte en un punto de retroceso (de orden superior) de \(f(z,s)=0\) (respectivamente de la curva en el espacio superior).

Como toda función arbitraria \(\eta\) del cuerpo puede expresarse racionalmente, por hipótesis, mediante \(z\) y \(s\), todo punto de la superficie de Riemann absoluta en el que \(z=a\), \(s=b\) se obtiene, por la isomorfía, asignando a cada función \(\eta\), representada mediante \(z\) y \(s\), su valor para \(z=a\), \(s=b\). Para que el punto sea singular, debe haber, por tanto, al menos una función \(\eta\) que para \(z=a\), \(s=b\) tome varios valores o tome un sistema de valores varias veces. Y, debido a la expresabilidad racional mediante \(z\) y \(s\), esto sólo puede ocurrir si en esa representación numerador y denominador se anulan simultáneamente para \(z=a\), \(s=b\). De aquí se sigue que la definición anterior es idéntica a la usual, según la cual \(f\), \(\frac{\partial f}{\partial z}\) y \(\frac{\partial f}{\partial s}\) se anulan para \(z=a\), \(s=b\). En efecto, en este caso \(\eta=\frac{z-a}{s-b}\) es una función de la clase indicada que es multivaluada para \(z=a\), \(s=b\), mientras que en el caso contrario toda función, aun cuando se vuelva \(\frac{0}{0}\), toma un solo valor.

Con el concepto aritmético de punto no puede aparecer el punto singular, pues dos puntos se consideran distintos tan pronto como a una sola función se le asignan sistemas de valores distintos. Pero incluso en la teoría de ideales que sirve para fundamentar el concepto aritmético de punto no aparece el concepto de punto singular, precisamente porque siempre se considera a la vez la totalidad de todas las cantidades enteras del cuerpo. Todos los puntos singulares en la parte finita de \(F(z,\delta)=0\), donde \(\delta\) es algebraicamente entero con respecto a \(z\), son generados (según la segunda definición y el final del núm. 3) por el ``ideal de puntos dobles'' \(\frf\). Ahora bien, como por el teorema fundamental el ideal de ramificación \(\frD\) es el máximo común divisor de todos los ideales principales \(F'(\delta)\), para cada ideal primo \(\frp\) puede indicarse un \(\delta\) tal que el correspondiente \(\frf\) no sea divisible por \(\frp\), y por consiguiente --- puesto que \(\frf\) es el máximo común divisor de \(F'(\delta)\) y \(F'(z)\) --- al menos uno de los dos ideales principales \(F'(\delta)\) y \(F'(z)\) no sea divisible por \(\frp\). Por tanto, \emph{no hay ningún sistema de valores para el cual se anulen simultáneamente todos los \(F'(\delta)\) y \(F'(z)\)}, y en consecuencia el sistema de valores \(z=a,\delta_i=b_i\) se asume en un único punto de la superficie de Riemann absoluta; la \emph{totalidad de las funciones algebraicas enteras del cuerpo} (que puede interpretarse como una curva en un espacio de infinitas dimensiones) \emph{no posee ningún punto singular en la parte finita}; pero en la teoría de ideales sólo entran en consideración valores finitos.

Tampoco las funciones de base \(\omega_1,\ldots,\omega_n\) de todas las cantidades enteras poseen puntos singulares en la parte finita. Por lo anterior, cada punto en sentido aritmético ya está determinado unívocamente por un sistema de constantes asignado isomórficamente a todas las funciones enteras. Si, por tanto, una ``curva espacial'' definida por funciones enteras cualesquiera \(s_i\) posee un punto singular en la parte finita para \(s_i=\alpha_i\), entonces debe haber al menos una función entera \(\eta\) que tome varios sistemas de valores para \(s_i=\alpha_i\). Si ahora \(\omega_1,\ldots,\omega_n\) denota una base de todas las cantidades enteras, entonces, por \(\delta=a_1\omega_1+\cdots+a_n\omega_n\) con \(a_i\) racionales enteros, a cada sistema de valores de las \(\omega\) corresponde un único sistema de valores de \(\delta\); de ahí que la curva espacial definida por \(\omega_1,\ldots,\omega_n\) no pueda poseer ningún punto singular en la parte finita; y \emph{todo punto de la superficie de Riemann absoluta en el que \(z\) es finito queda ya definido unívocamente por un sistema de constantes asignado isomórficamente a las \(\omega\). Las funciones de base son precisamente las funciones \(\eta\) que se vuelven multivaluadas para \(s_i=\alpha_i\).} Así, al introducir la base de las cantidades enteras, por ejemplo en lugar de una base \(1,\delta,\ldots,\delta^{n-1}\), se evita la dificultad de los puntos singulares.

En la teoría geométrica, la tarea de generar unívocamente todos los puntos de la superficie de Riemann absoluta --- más generalmente, de formar todas las funciones que se vuelven infinitas en puntos dados --- se realiza mediante las formas adjuntas y sus cocientes. Una forma \(\psi\) se llama adjunta si \(\psi=0\) posee al menos un punto \((i-1)\)-ple en cada punto \(i\)-ple de la curva base \(f(z,s)=0\); un punto singular superior debe resolverse primero en puntos múltiples ordinarios. Arriba se mostró que, para la función más general \(\eta=\frac{\psi}{\chi}\), numerador y denominador deben anularse en todo punto singular de \(f(z,s)=0\); que deben ser adjuntos lo muestra la consideración de los diferenciales finitos en todas partes. Que, recíprocamente, las condiciones de adjunción son también suficientes para formar la función más general lo muestra el ``teorema del residuo'', que se discutirá con más detalle más abajo. Así, las funciones adjuntas ocupan el lugar de las funciones de base de la teoría aritmética.

Aritméticamente, a una forma adjunta corresponde el numerador de una función del cuerpo divisible por el ``ideal de puntos dobles'' \(\frf\), tan pronto como se supone que los puntos singulares --- lo cual siempre puede lograrse por transformación --- están todos en la parte finita. De aquí resulta formalmente la representación de las cantidades de base \(\omega_1,\ldots,\omega_n\) como cocientes de formas adjuntas, y de ahí la conexión entre las distintas maneras de representarlas, como sigue:

Sea \(R(z)\) el determinante de sustitución de una base \(1,\delta,\ldots,\delta^{n-1}\) en una base \(\omega_1,\ldots,\omega_n\). Entonces, recíprocamente, todo \(\omega\) puede expresarse como cociente de una función entera de \(z\) y \(\delta\) por \(R(z)\):
\[
\omega_i=\frac{G_i(z,\delta)}{R(z)}.
\]
Como \(R(z)\omega_i\) pertenece al anillo derivado de \(\delta\), por la definición del conductor \(R(z)\) es divisible por \(\frf\) [de \((R(z))^2=\operatorname{Norm}\frf\) esta divisibilidad se sigue sólo para \((R(z))^2\)]. Como ahora \(\omega_i\), en cuanto función entera, permanece finita para todo \(z\) finito, también la función numerador \(G_i(z,\delta)\) debe ser divisible por \(\frf\); por tanto, numerador y denominador de todas las \(\omega_i\) son adjuntos. De esta representación de las \(\omega\) se sigue, sin embargo, la representación de todas las funciones del cuerpo por cocientes de adjuntas.

Según lo anterior, la teoría geométrica puede caracterizarse como la \emph{teoría del anillo derivado de una cantidad \(\delta\)} (orden), en contraste con la teoría aritmética, que considera simultáneamente \emph{todas las cantidades enteras}. Así es claro que el conductor del anillo --- los puntos singulares --- desempeña un papel decisivo. Pueden darse además otras analogías entre la teoría aritmética de los anillos (órdenes) y la teoría geométrica. Así, la teoría geométrica no cuenta en ningún lugar los puntos de intersección que caen en los puntos singulares al cortar con curvas adjuntas, sino que considera sólo los grupos de puntos distintos de ellos. A esto corresponde que Dedekind\footnote{Über die Anzahl der Idealklassen in den verschiedenen Ordnungen eines endlichen Körpers. (Festschrift Braunschweig 1877); § 3--5. (Los teoremas enunciados en esos parágrafos para números algebraicos valen literalmente para funciones algebraicas de una variable.)} considera sólo aquellos ideales de anillo que son primos con el ideal conductor \(\frf\), y para este dominio de ideales establece todas las leyes de divisibilidad, incluida la descomposición única en ideales primos.

También el ``teorema fundamental de las funciones algebraicas''\footnote{M. Noether, Über einen Satz aus der Theorie der algebraischen Funktionen. \emph{Math. Ann.} vol. 6 (1873).} que está en la base de la teoría geométrica, o al menos una consecuencia inmediata de él, tiene en cierto sentido un análogo en la teoría aritmética de los anillos. Este teorema da las condiciones para una representabilidad
\[
\psi(z,s)=A(z,s)f(z,s)+B(z,s)\varphi(z,s),
\]
donde todas las cantidades que aparecen son racionales enteras en \(s,z\) (más generalmente enteras y homogéneas en \(x_1,x_2,x_3\)). De estas condiciones se sigue,\footnote{Brill--Noether, núm. 55.} en virtud de \(f(z,s)=0\), que el cociente \(\frac{\psi}{\varphi}\) es siempre igual a una función racional entera \(B(z,s)\) cuando es adjunto a \(f\). A esto corresponde el teorema del núm. 3, que subyace a todas estas comparaciones, según el cual el conductor \(\frf\) de un anillo derivado de \(\delta\) es igual al ideal de puntos dobles de \(f(z,\delta)=0\). Pues, por la definición del conductor, toda cantidad algebraica entera divisible por \(\frf\) pertenece al anillo, y por tanto es entera y racionalmente representable mediante \(z\) y \(\delta\); mientras que el teorema citado dice que tal cantidad es una función adjunta.\footnote{Siguiendo a Dedekind, aquí se ha hecho siempre la hipótesis de que \(\delta\) es un entero algebraico (lo cual siempre puede lograrse por transformación). El caso que corresponde más exactamente a la geometría, \(f(z,s)=0\), donde \(s\) puede depender también algebraicamente de modo fraccionario de \(z\), lo tratan Hensel--Landsberg; compárese Hensel núms. 26 y 29.}

\subsection*{7. Comparación entre el teorema del residuo y los teoremas de la teoría de ideales}
La consecuencia recién mencionada del ``teorema fundamental de las funciones algebraicas'' conduce en la teoría geométrica directamente al ``teorema del residuo'', y con ello a la base sobre la que se edifica toda la teoría geométrica. En cierto respecto, sin embargo, el teorema del residuo tiene de nuevo un carácter más elemental que los restantes fundamentos, en cuanto puede formularse aritméticamente como consecuencia directa de la descomposición única de los ideales en ideales primos, o más bien de los hechos que están en la base de este teorema de descomposición, y por tanto no descansa en la teoría superior de los anillos. Se llega a esta formulación aritmética de la manera siguiente:

Supóngase que la curva base \(f(z,\delta)=0\) --- como siempre puede lograrse mediante una transformación lineal fraccionaria de las variables --- es tal que los puntos singulares y los finitos grupos de puntos considerados yacen en la parte finita; y además que \(\delta\) depende algebraicamente de modo entero de \(z\), lo cual siempre puede lograrse mediante una transformación lineal posterior entera en las variables. Entonces, a un ideal primo \(\frp\) que genera un punto \(P\) de la superficie de Riemann absoluta en el que \(z=a,\delta=b\), corresponde el punto \(z=a,\delta=b\) de \(f(z,\delta)=0\); más generalmente, un ideal \(\fra=\frp_1^{\,e_1}\cdots\frp_\rho^{\,e_\rho}\) genera el grupo de puntos formado por los puntos \(z=a_i,\delta=b_i\) con la multiplicidad correspondiente; y así se generan todos los puntos de \(f(z,\delta)=0\) que yacen en la parte finita. Como en un punto \(P\) se anulan todas las funciones \(g_1,g_2,\ldots\) divisibles por \(\frp\), y como recíprocamente esta totalidad de funciones \(g\) genera el ideal \(\frp\), el punto correspondiente sobre \(f\) es la intersección de
\[
g_1(z,\delta)=0,\quad g_2(z,\delta)=0,\ldots
\]
con \(f=0\); lo mismo vale para el grupo de puntos correspondiente a un ideal \(\fra\). De hecho, por el teorema de que todo ideal es el máximo común divisor de dos ideales principales, dos curvas de este tipo bastan siempre para cortar un grupo de puntos. En particular, el grupo de puntos correspondiente a un ideal principal está cortado por una curva \(g(z,\delta)=0\), y recíprocamente; el ideal principal corresponde a la intersección completa.

Dos grupos de puntos \(A\) y \(B\) se llaman \emph{corresiduales} en la teoría geométrica si pueden completarse con el mismo grupo de puntos \(R\) hasta una intersección completa. Si estos grupos de puntos son generados respectivamente por los ideales \(\fra,\frb,\frr\), entonces entre estos ideales se tiene la relación
\[
\fra\cdot\frr=(\alpha),\qquad \frb\cdot\frr=(\beta),
\]
lo que muestra que los ideales \(\fra\) y \(\frb\) son equivalentes. Recíprocamente, de la equivalencia de dos ideales, \(\fra=\eta\frb\), se sigue siempre esta relación sobre la base del teorema de que todo ideal puede transformarse en un ideal principal multiplicándolo por otro ideal;\footnote{Compárese Dedekind, teoría de números § 181; en efecto, si \(\frr\) se elige de modo que \(\fra\frr\) se convierta en un ideal principal \((\alpha)\), entonces \(\frb\frr=\eta^{-1}\cdot(\alpha)\); y como \(\frb\frr\) contiene sólo cantidades enteras, lo mismo vale para \(\eta^{-1}\cdot(\alpha)\), que por tanto es también un ideal principal \((\beta)\).} \emph{los ideales equivalentes y los grupos de puntos corresiduales son, por tanto, conceptos idénticos.}

El teorema del residuo dice ahora: \emph{si \(A\) y \(B\) son corresiduales, entonces pueden cortarse por curvas adjuntas con el mismo residuo (fuera de los puntos singulares).} Como los ceros y polos de una función son siempre corresiduales, esto demuestra al mismo tiempo que las funciones más generales del cuerpo son representables como cocientes de formas adjuntas. En el lenguaje de la teoría de ideales, el teorema del residuo dice simplemente que, además de \(\fra\) y \(\frb\), también \(\fra\frf\) y \(\frb\frf\) son equivalentes, donde \(\frf\) denota el ideal de puntos dobles. Pues entonces, por lo anterior, existe siempre un ideal \(\frm\) tal que
\[
\fra\frf\frm=(\alpha),\qquad \frb\frf\frm=(\beta)
\]
se vuelven ideales principales; \(\alpha=0\) y \(\beta=0\) son las curvas adjuntas que cortan los grupos de puntos \(A\) y \(B\), respectivamente sus subgrupos \(A'\) y \(B'\) distintos de los puntos singulares.\footnote{Si \(\fra=\frf\fra_1,\ \frb=\frf\frb_1,\ \frr=\frf\frr_1\), donde \(\fra_1,\frb_1,\frr_1\) son primos entre sí, entonces ya \(\fra_1\frf\mathfrak n,\ \frb_1\frf\mathfrak n\) se vuelven ideales principales, donde \(\mathfrak n\) es primo con \(\frf\). Pues todo ideal puede transformarse en un ideal principal multiplicándolo por un ideal que sea primo con uno dado (D.--W. p. 214; o teoría de números, § 178, teorema XI).} La demostración del teorema del residuo está por tanto contenida aritméticamente en la posibilidad de asignar puntos a ideales primos, y en la demostración de que los ideales equivalentes pueden transformarse en ideales principales multiplicándolos por uno y el mismo ideal.

Debe observarse que dos grupos de puntos corresiduales no necesitan consistir en el mismo número de puntos, del mismo modo que ideales equivalentes no necesitan tener el mismo grado. Así, por ejemplo, todos los ideales principales son equivalentes, y todos los grupos de puntos generados por intersecciones completas son corresiduales. Esta diferencia respecto de polígonos (divisores) equivalentes\footnote{Un polígono \(A\) consiste en finitos puntos de la superficie de Riemann absoluta; dos polígonos \(A\) y \(B\) se llaman equivalentes si son ceros y polos de una función \(\eta\); entonces \(\eta\) admite la representación por cocientes de polígonos \(\eta=\frac{A}{B}\). (D.--W. § 17 y 18.)} se explica por el hecho de que aquellos puntos en los que \(z\) se vuelve infinito no son generados por ideales primos en \(z\). En la representación de una función como cociente de ideales \(\eta=\frac{\fra}{\frb}\), que expresa la equivalencia de \(\fra\) y \(\frb\), no aparecen por tanto los ceros o polos de \(\eta\) en los que \(z\) se vuelve infinito. Por ejemplo, el ideal principal \((\alpha)\) se convierte en el producto de los ideales primos correspondientes a los ceros de la función entera \(\alpha\); los polos de \(\alpha\) no aparecen, puesto que \(\alpha\) se vuelve infinita sólo en puntos en los que \(z\) se vuelve infinito. En este respecto, la teoría de ideales es el análogo exacto de la teoría de formas en la teoría geométrica, donde igualmente sólo hay ceros y no polos.

El paso de formas a funciones se realiza geométricamente formando cocientes de formas del mismo orden; aritméticamente, esto corresponde al paso de las clases de ideales a las clases de polígonos (divisores) definidas invariantemente. En efecto, como aquí ya no se distingue ninguna variable, la representación de una función por cocientes de polígonos da todos los ceros y polos. La relación de la clase de ideales (clase de grupos de puntos corresiduales) con la clase de polígonos es la siguiente: si \(\calA\) y \(\calB\) son los polígonos generados por los ideales \(\fra\) y \(\frb\), y \(\calN\) es el polígono de los puntos en los que \(z\) se vuelve infinito (el polígono denominador de \(z\) y de las funciones enteras de \(z\)), entonces de la equivalencia de \(\fra\) y \(\frb\) se sigue la equivalencia de \(\calA\) y \(\calB\) si y sólo si \(\fra\) y \(\frb\) tienen el mismo grado; en general sólo se sigue la equivalencia de \(\calA\calN^\nu\) con \(\calB\calN^{\nu'}\) (\(\nu\ne\nu'\)). La división de los ideales (respectivamente de los grupos de puntos corresiduales) en clases es, por tanto, una reunión de infinitas clases de polígonos en una sola clase; también puede considerarse como una división de los polígonos en clases módulo una potencia de \(\calN\).\footnote{Hensel--Landsberg, 25.ª lección, p. 438. --- Si se representan las dos variables \(z\) y \(s\) por cocientes de polígonos, esto puede entenderse como una homogeneización binaria de cada una de estas variables; si se eligen en particular \(z\) y \(s\) de modo que tengan el mismo polígono denominador, surge la homogeneización ternaria usual en geometría.}

El teorema del residuo para el caso de grupos de puntos corresiduales que consisten en números distintos de puntos sólo aparece en cuestiones puramente geométricas. En la aplicación a funciones algebraicas aparece de hecho sólo el caso correspondiente a clases de polígonos, el de grupos de puntos corresiduales con el mismo número. El ``sistema lineal completo'' generado por un grupo de puntos, es decir, la totalidad de los grupos de puntos corresiduales a ese grupo y con el mismo número de puntos, corresponde exactamente a la clase de polígonos generada por el polígono correspondiente; también coinciden los conceptos de ``suma de dos sistemas completos'' y ``producto de dos clases de polígonos''. Que tal sistema completo contiene sólo finitos grupos de puntos linealmente independientes es consecuencia directa del teorema del residuo, que reduce la dimensión del sistema (el número de grupos de puntos linealmente independientes) a la dimensión de todas las curvas adjuntas de orden arbitrario pero fijo que pasan por el grupo residual. Para llegar a la determinación efectiva de esta dimensión, el teorema de Riemann--Roch, la teoría geométrica deriva primero del teorema del residuo el ``teorema de reducción'' (teorema sobre los puntos fijos);\footnote{Brill--Noether, núm. 58.} y de éste, el teorema de Riemann--Roch.

De manera enteramente correspondiente, la teoría aritmética muestra que toda clase de polígonos es de dimensión finita, y determina su dimensión mediante una base especial de las cantidades enteras, una base normal definida por medio del módulo complementario; así obtiene el teorema de Riemann--Roch como consecuencia de la conexión del módulo complementario con el polígono de ramificación. En la exposición de D.--W. aparece también el teorema de reducción,\footnote{D.--W., § 30, núm. 2.} aunque no como fundamento, como en la teoría geométrica. Como consecuencia de la correspondencia mencionada en el núm. 6 entre funciones de base y funciones adjuntas, también aquí puede establecerse cierta correspondencia entre los medios auxiliares de ambas teorías. En resumen, puede decirse que las formaciones de conceptos en las dos teorías coinciden esencialmente, aunque, aparte de la formulación, difieran en su disposición. Ambas teorías surgieron independientemente una de otra; la geométrica, una década antes. Debe recordarse además que el teorema de Riemann--Roch aparece originariamente, en Riemann y Roch y también en Weierstrass, como teorema de rango: se trata del rango de la matriz \((\varphi_i(x_j))\), donde las \(\varphi\) recorren las \(p\) curvas adjuntas de orden \(n-3\), y las \(x\), los puntos del grupo residual.

\subsection*{8. Cuestiones trascendentes para funciones algebraicas y números algebraicos}
Hilbert señaló en particular analogías entre cuestiones trascendentes para funciones algebraicas y para números algebraicos;\footnote{Mathematische Probleme, problema 12. (\emph{Göttinger Nachrichten} 1900).} se trata esencialmente de cuestiones de existencia. Así, la cuestión de Riemann sobre la existencia de una función algebraica para una superficie de Riemann dada (por tanto también para puntos de ramificación dados) corresponde a la cuestión de la existencia de cuerpos de números algebraicos con grupo y discriminante prescritos.\footnote{El tratamiento de funciones algebraicas sugerido por ello, partiendo del grupo, fue llevado a cabo por R. König como caso especial de cuestiones más generales; compárese especialmente: Riemannsche Funktionen- und Differentialsysteme in der Ebene (\emph{J. f. M.} vol. 148) y \emph{Math. Ann.} 78, 79. En lugar del cuerpo algebraico, König tiene un módulo finito especial respecto de \(K(z)\): sólo hay ``propiedad de clase''.} Mientras que en el caso de las funciones algebraicas la demostración de existencia puede llevarse a cabo en general, en el dominio de los cuerpos de números algebraicos se ha logrado sólo para cuerpos base especiales (números racionales, cuerpos cuadráticos) y sólo para grupos abelianos.\footnote{Sólo el caso simple de los grupos que aparecen para ecuaciones de tercer y cuarto grado, y algunos grupos muy especiales para ecuaciones de grado 8, pueden tratarse para cuerpos base arbitrarios: G. Bucht, Über einige algebraische Körper 8. Grades (\emph{Arkiv f. Math., Astron. och Physik}, vol. 6, núm. 30). Compárese también E. Noether, Gleichungen mit vorgeschriebener Gruppe y la tesis de Seidelmann citada allí. (\emph{Math. Ann.} 78).} Esto concierne al teorema de Kronecker según el cual todo cuerpo numérico abeliano sobre los números racionales puede componerse de cuerpos de raíces de la unidad, y a los teoremas correspondientes de Fueter\footnote{\emph{Math. Ann.} 75.} y Hecke\footnote{\emph{Math. Ann.} 76.} para cuerpos numéricos relativamente abelianos. Así, los cuerpos numéricos deseados son suministrados realmente por medios trascendentes: funciones exponenciales, funciones modulares; en contraste con la pura demostración de existencia para funciones algebraicas.

A la demostración de existencia de funciones algebraicas la precede la existencia de integrales finitos en todas partes. Un análogo de esto es la existencia de los ``números primarios singulares'' de un cuerpo de números algebraicos definidos con respecto a un número primo \(\ell\), cuyas raíces \(\ell\)-ésimas son relativamente no ramificadas y suministran los primeros bloques de construcción del cuerpo de clases.\footnote{Furtwängler, Reziprozitätsgesetze für Potenzreste mit Primzahlexponenten in algebraischen Zahlkörpern. \emph{Math. Ann.} 67 y 72.} Esta demostración de existencia se lleva a cabo con ayuda esencial del teorema según el cual en el cuerpo numérico existen siempre ideales primos con caracteres residuales prescritos; por tanto, este teorema puede considerarse como análogo del problema de valores de frontera sobre el que descansa la existencia de los integrales finitos en todas partes.

Los integrales finitos en todas partes proporcionan un criterio trascendente para determinar si dos polígonos pertenecen a la misma clase (dos grupos de puntos del mismo número son corresiduales), a saber, mediante el teorema de Abel, que dice que los integrales de primera especie, extendidos sobre ``caminos corresiduales'', se anulan; o mediante el correspondiente teorema diferencial junto con su recíproco.\footnote{Compárese, por ejemplo, la formulación en M. Noether (\emph{Ann.} 37). Como el teorema del residuo nació históricamente del teorema de Abel, la teoría geométrica habla de la ``suma'' de grupos de puntos y de sistemas lineales, siguiendo las sumas integrales; en contraste con el ``producto'' de las clases en la teoría aritmética.}

Como análogo del teorema de Abel para cuerpos de números algebraicos debe tomarse la ``ley de reciprocidad para residuos de potencias \(\ell\)-ésimas'', que da, para el cuerpo mismo o al menos para un cuerpo superior adecuado, un criterio de que dos ideales pertenezcan a la misma clase. Su contenido puede formularse también diciendo que en este cuerpo superior los números primarios singulares tienen el mismo carácter residual respecto de todos los ideales de la misma clase. Este último hecho vale también en el cuerpo mismo, pero allí no coincide con la ley de reciprocidad.\footnote{Furtwängler, loc. cit.}

Finalmente, la teoría de Hensel de los números \(p\)-ádicos suministra el análogo del desarrollo en series de potencias de las funciones algebraicas; para ello, véase el informe de Hensel.


% Paper 15 current invariant-theory macros
\providecommand{\frH}{\mathfrak H}
\providecommand{\frK}{\mathfrak K}
\providecommand{\frM}{\mathfrak M}
\providecommand{\frN}{\mathfrak N}
\providecommand{\frG}{\mathfrak G}
\providecommand{\frS}{\mathfrak S}
\providecommand{\tmod}[1]{\;(\operatorname{mod} #1)}

\clearpage
\begin{center}
{\Large\bfseries 15. La finitud del sistema de invariantes enteros de las formas binarias}\par
\vspace{1em}
\emph{Nachrichten von der Gesellschaft der Wissenschaften zu Göttingen} 1919, pp. 138--156
\end{center}

\vspace{2em}
\begin{center}
Presentado por F. Klein en la sesión del 27 de marzo de 1919.
\end{center}

En lo que sigue se demuestra, como ampliación del conocido teorema de finitud, el siguiente teorema:
\emph{todos los invariantes polinómicos enteros de un sistema binario de formas fundamentales son funciones polinómicas racionales, con coeficientes enteros, de un número finito de ellos.}
Aquí, como de costumbre, por ``invariante'' se entiende siempre un invariante polinómico racional de una forma fundamental o de un sistema de formas fundamentales, es decir, un invariante respecto de la transformación de coeficientes inducida por el grupo de todas las transformaciones lineales.

La demostración se apoya en un refinamiento de integridad de la demostración de finitud de Mertens-Hilbert.\footnote{F. Mertens, Wiener Sitzungsberichte Jan. 1889; D. Hilbert, \emph{Math. Ann.} vol. 33 (1889) [fechado en marzo de 1888]. Las dos demostraciones, surgidas independientemente una de otra siguiendo a Mertens, Crelle vol. 100, coinciden en contenido.}
Esa demostración de finitud puede caracterizarse como sigue. Se descomponen las formas fundamentales en sus factores lineales - dicho con más precisión, se asocia a cada forma fundamental un producto de formas lineales -, con lo cual todo invariante \(I\) se transforma en un invariante de esas formas lineales, y por tanto en una función de las raíces homogeneizadas. Esta función debe satisfacer tres condiciones: 1) la condición de invariancia, 2) condiciones de peso, y 3) condiciones de simetría. Las condiciones 2) y 3) son necesarias para que \(I\) llegue a ser polinómico y racional en los coeficientes de las formas fundamentales. La condición 1) se satisface representando \(I\) como una función polinómica racional de los determinantes de las formas lineales; la condición 2), tomando, en lugar de los determinantes individuales, ciertos productos de potencias de ellos como nuevos argumentos. La condición 3) dice entonces simplemente que \(I\) se convierte en una función polinómica racional simétrica de filas de cantidades y que, recíprocamente, toda función simétrica de tal tipo de filas de cantidades se convierte en un invariante de las formas fundamentales. Las funciones simétricas elementales de esas filas de cantidades proporcionan por tanto el sistema completo.

Si se consideran solamente invariantes enteros, la condición 1) puede satisfacerse exactamente igual tan pronto como, para los invariantes de formas lineales, se haya probado el teorema más preciso de que \emph{todos los invariantes enteros de formas lineales binarias son funciones polinómicas enteras de los determinantes de esas formas lineales}. Muestro que esto es efectivamente así en el § 3, en la forma de que la totalidad de todas las relaciones módulo cualquier primo entre esos determinantes coincide con la totalidad de todas las relaciones idénticas; esto es, el módulo correspondiente a esas relaciones sigue siendo un módulo primo módulo todo número primo. Esta verificación se efectúa normalizando las clases residuales respecto del módulo de formas cuadráticas correspondiente a las relaciones cuadráticas entre los determinantes; esas formas cuadráticas resultan entonces formar una base entera del módulo, módulo todo primo, mientras que antes se las conocía como base de módulo sólo cuando se prescindía de la integridad.

Para satisfacer la condición 2) no hace falta modificación alguna, pues la introducción de productos de potencias no afecta a los coeficientes enteros. La condición 3) dice ahora que \(n_1!\cdots n_\mu! I\) (donde \(n_1,\ldots,n_\mu\) son los grados de las formas fundamentales individuales) se convierte en una función simétrica entera de filas de cantidades y, recíprocamente, toda función simétrica entera de esas filas se convierte en un invariante entero. Pero, como se demuestra fácilmente, esas funciones simétricas enteras de filas poseen una base entera, para la cual las funciones elementales por sí solas ya no bastan. La aplicación del teorema de Hilbert sobre la base entera de módulos proporciona entonces también una base entera para los propios \(I\).

El § 1 reúne los teoremas especiales de finitud correspondientes a las condiciones 1), 2) y 3), de los cuales se deduce en el § 2 la demostración de la finitud. Los teoremas especiales de finitud correspondientes a la condición 3) dan también la finitud de los invariantes enteros de grupos finitos de sustituciones enteras o algebraico-enteras, como se desarrolla brevemente en el § 4.

\section*{§ 1. Teoremas especiales de finitud}

En primer lugar enuncio los siguientes teoremas especiales de finitud, que servirán para satisfacer las condiciones de invariancia, peso y simetría mencionadas en la introducción.

\textbf{I.} Todo invariante entero de un sistema de formas lineales binarias es una función polinómica racional, con coeficientes enteros, de los determinantes de esas formas lineales. La demostración se dará en el § 3.

\textbf{II.} Todo sistema \(\frS\) de productos de potencias de \(n\) variables con exponentes no negativos que, junto con cualesquiera dos productos de potencias \(f\) y \(g\), contiene también su cociente siempre que dicho cociente sea polinómico en las variables, posee una base multiplicativa finita \(f_1,\ldots,f_k\), formada por funciones de \(\frS\), tal que para todo \(f\) de \(\frS\) se tiene una representación
\[
        f=f_1^{\lambda_1}\cdots f_k^{\lambda_k}
\]
con exponentes no negativos \(\lambda_1,\ldots,\lambda_k\).

En efecto, por el teorema de Hilbert sobre la base de módulos existe un número finito de funciones de \(\frS\), digamos \(f_1,\ldots,f_k\), cada una de las cuales contiene realmente las \(x\), tales que toda \(f\) no constante de \(\frS\) pertenece al módulo generado por ellas:
\[
        f\equiv0\pmod{(f_1,\ldots,f_k)}.
\]
De aquí se sigue que para toda \(f=x_1^{i_1}\cdots x_n^{i_n}\) existe al menos una \(f_\nu=x_1^{\nu_1}\cdots x_n^{\nu_n}\) tal que \(\nu_1\le i_1,\ldots,\nu_n\le i_n\). En la representación
\[
        f=x_1^{i_1-\nu_1}\cdots x_n^{i_n-\nu_n} f_\nu = A f_\nu
\]
el factor \(A\) pertenece a \(\frS\) por hipótesis, pero tiene grado menor en las variables que \(f\). Por consiguiente, repitiendo el argumento un número finito de veces, se obtiene la afirmación. Para que la representación valga también para \(f=1\), basta poner todos los exponentes \(\lambda\) iguales a cero.\footnote{Esta representación puede demostrarse también directamente, y recíprocamente de ella se deduce el teorema de Hilbert: compárese P. Gordan, Neuer Beweis des Hilbertschen Satzes über homogene Funktionen, Gött. Nachr. 1899. Otra demostración directa de esta representación, de la que también se deriva el Teorema II, se encuentra en A. Ostrowski, Über die Existenz einer endlichen Basis bei gewissen Funktionensystemen, \emph{Math. Ann.} 78 (1916), § 2,1 y § 3,2. El caso especial del Teorema II usado en el § 2, donde los exponentes recorren todas las soluciones no negativas de un sistema diofántico de ecuaciones, se halla ya en Gordan-Kerschensteiner, \emph{Vorlesungen über Invariantentheorie}, vol. I, p. 199.}

\textbf{III.} Todas las funciones simétricas polinómicas enteras de \(n\) filas de cantidades son funciones polinómicas enteras de un número finito de ellas.

Sean \(\tau_1(x),\ldots,\tau_n(x)\) las funciones simétricas elementales de las indeterminadas \(x_1,\ldots,x_n\). Entonces para todo exponente \(k\) existe una identidad
\[
        x_i^k = G_0^{(k)}(\tau)+x_iG_1^{(k)}(\tau)+\cdots+x_i^{n-1}G_{n-1}^{(k)}(\tau),
        \qquad (i=1,2,\ldots,n),
\]
donde \(G_0^{(k)},\ldots,G_{n-1}^{(k)}\) son funciones polinómicas enteras de las \(\tau(x)\). Si \(x_i,y_i,\ldots,z_i\) son los elementos de la \(i\)-ésima fila, entonces, mediante estas identidades y las correspondientes para \(y,\ldots,z\), todas las funciones simétricas especiales, de una sola fila,
\[
        \sum_{i=1}^n x_i^{a}y_i^{b}\cdots z_i^{c}
\]
se convierten en funciones polinómicas enteras de las finitamente muchas
\[
        \sum_{i=1}^n x_i^{a}y_i^{b}\cdots z_i^{c}
        \qquad
        (0\le a<n,\;0\le b<n,\;\ldots,\;0\le c<n)
\]
y de las funciones simétricas elementales \(\tau(x),\tau(y),\ldots,\tau(z)\).\footnote{Hilbert, loc. cit., incluye en la base las sumas de potencias en lugar de las \(\tau(x),\ldots,\tau(z)\); con ello la base consta sólo de funciones de una fila, pero se pierde la integridad de la representación.}
Si se tienen las funciones simétricas más generales de las \(n\) filas - caso que no aparece más abajo -, se demuestra del mismo modo que las funciones
\[
        \sum x_{i_1}^{a_1}y_{i_1}^{b_1}\cdots z_{i_1}^{c_1}
              x_{i_2}^{a_2}y_{i_2}^{b_2}\cdots z_{i_2}^{c_2}
              \cdots
              x_{i_n}^{a_n}y_{i_n}^{b_n}\cdots z_{i_n}^{c_n}
\]
\[
        (0\le a_\nu<n,\;0\le b_\nu<n,\;\ldots,\;0\le c_\nu<n),
\]
donde \(i_1,i_2,\ldots,i_n\) recorre todas las permutaciones, junto con \(\tau(x),\tau(y),\ldots,\tau(z)\), forman una base entera.

\textbf{IV.} Sean \(F_1(x),\ldots,F_k(x)\) funciones polinómicas enteras de \(x_1,\ldots,x_n\), y sea \(a\) un entero fijo. Entonces todas las funciones
\[
        \frac{G(F)}{a},
\]
donde \(G\) denota una función polinómica entera, que llegan a ser enteras en las \(x\), son funciones polinómicas enteras de un número finito de ellas.

En efecto, por el teorema de Hilbert sobre la base entera de módulos, para cada tal \(G(F)/a\) se tiene una representación
\[
 \frac{G(F)}{a}
 =
 A_1(F_1\ldots F_k)\frac{G_1(F)}{a}
 +\cdots+
 A_\rho(F_1\ldots F_k)\frac{G_\rho(F)}{a},
\]
donde \(A_1,\ldots,A_\rho\) son funciones polinómicas enteras de las \(F\), y
\[
        \frac{G_1(F)}{a},\ldots,\frac{G_\rho(F)}{a}
\]
pertenecen al sistema. Las funciones
\[
        F_1=\frac{aF_1}{a},\ldots,F_k=\frac{aF_k}{a};
        \quad
        \frac{G_1(F)}{a},\ldots,\frac{G_\rho(F)}{a}
\]
forman por tanto la base requerida.

\setcounter{equation}{0}
\section*{§ 2. Finitud de los invariantes binarios enteros}

Sean \(x,y\) variables binarias sometidas a una transformación con coeficientes indeterminados
\begin{equation}
        x=s_{11}x'+s_{12}y',\qquad y=s_{21}x'+s_{22}y',
\end{equation}
y sea \(f\) una forma fundamental binaria cuyos coeficientes son indeterminadas. Entonces, por (1),
\begin{equation}
        f=a_0x^n+a_1x^{n-1}y+\cdots+a_ny^n
        =a_0'x'^n+a_1'x'^{\,n-1}y'+\cdots+a_n'y'^n.
\end{equation}
Por invariante entero de \(f\) se entiende, como es sabido, una función polinómica entera \(I(a)\) tal que
\begin{equation}
        I(a')=|s_{ik}|^\rho I(a),\qquad \text{idénticamente en \(a,s_{ik}\)}.
\end{equation}
Los invariantes enteros de un sistema de formas fundamentales se definen de manera correspondiente; se puede incluir en la definición la homogeneidad en los coeficientes de las formas fundamentales individuales, puesto que de tales invariantes individualmente homogéneos se componen los restantes de manera lineal entera.

Además de (2), considérese una forma binaria especial \(g\), a saber, el producto de \(n\) formas lineales cuyos coeficientes son de nuevo indeterminadas:\footnote{Si, como se hace a menudo en teoría de invariantes, se escribe \(f\) con coeficientes binomiales, \(f=b_0x^n+\binom n1b_1x^{n-1}y+\cdots+b_ny^n\), entonces el dominio de todos los invariantes enteros \(I(b)\) se amplía con respecto al de los \(I(a)\), y los argumentos de finitud empleados aquí fallan; \(I(b)\) ya no es una función entera de las raíces.}
\begin{align}
        g&=(\alpha_1x+\beta_1y)(\alpha_2x+\beta_2y)\cdots(\alpha_nx+\beta_ny)\notag\\
         &=\sigma_0(\alpha,\beta)x^n+\sigma_1(\alpha,\beta)x^{n-1}y+\cdots+\sigma_n(\alpha,\beta)y^n\notag\\
         &=(\alpha_1'x'+\beta_1'y')(\alpha_2'x'+\beta_2'y')\cdots(\alpha_n'x'+\beta_n'y')\notag\\
         &=\sigma_0'x'^n+\sigma_1'x'^{\,n-1}y'+\cdots+\sigma_n'y'^n.
\end{align}
Así, \(\sigma_0(\alpha,\beta),\sigma_1(\alpha,\beta),\ldots,\sigma_n(\alpha,\beta)\) denotan las funciones simétricas elementales homogeneizadas, y los coeficientes transformados \(\sigma_i'\) se vuelven iguales a \(\sigma_i(\alpha',\beta')\).

Por la independencia algebraica de estas funciones elementales, a toda relación
\begin{equation}
        G(\sigma_0,\sigma_1,\ldots,\sigma_n)=0
        \quad \text{idénticamente en \(\alpha,\beta\)}
\end{equation}
corresponde una relación ulterior
\begin{equation}
        G(\sigma_0,\sigma_1,\ldots,\sigma_n)=0
        \quad \text{idénticamente en \(\sigma_0,\ldots,\sigma_n\)},
\end{equation}
y por tanto, para indeterminadas \(a_0,\ldots,a_n\),
\begin{equation}
        G(a_0,a_1,\ldots,a_n)=0
        \quad \text{idénticamente en \(a_0,\ldots,a_n\)}.
\end{equation}

Ahora bien, todo invariante entero \(I(a)\) de \(f\), mediante la sustitución \(a_i=\sigma_i(\alpha,\beta)\), da un invariante entero \(I(\sigma)=H(\alpha,\beta)\) de \(g\), donde \(H\) se convierte en una función entera de \(\alpha,\beta\). Para este invariante, la relación (3) se convierte en
\begin{equation}
        I(\sigma')=|s_{ik}|^\rho I(\sigma)
        \quad \text{idénticamente en \(\sigma,s_{ik}\)},
\end{equation}
y como \(\sigma'=\sigma(\alpha',\beta')\), (8) da
\begin{equation}
        H(\alpha',\beta')=|s_{ik}|^\rho H(\alpha,\beta)
        \quad \text{idénticamente en \(\alpha,\beta;s_{ik}\)}.
\end{equation}
Así \(I(\sigma)=H(\alpha,\beta)\) se convierte en un invariante entero de las \(n\) formas lineales \((\alpha_i x+\beta_i y)\). Recíprocamente, supóngase que \(H(\alpha,\beta)\) es un invariante entero de estas \(n\) formas lineales, de modo que satisface (9), y que al mismo tiempo es igual a una función entera \(I(\sigma)\) de las \(\sigma\). Puesto que entonces
\[
        H(\alpha',\beta')=I(\sigma(\alpha',\beta'))=I(\sigma'),
\]
la ecuación (9) puede escribirse también
\begin{equation*}
        I(\sigma')=|s_{ik}|^\rho I(\sigma)
        \quad \text{idénticamente en \(\alpha,\beta;s_{ik}\)}.\tag{9a}
\end{equation*}
De (5) y (6) se sigue que vale (8), y de (7) que vale (3). Por tanto, a todo invariante entero \(I(a)\) le corresponde un invariante entero \(H(\alpha,\beta)=I(\sigma)\) de las \(n\) formas lineales, y recíprocamente. Además, si \(I_1(a),\ldots,I_k(a)\) es una base entera de todos los \(I(a)\), entonces \(H_1(\alpha,\beta),\ldots,H_k(\alpha,\beta)\) es una base entera de todos los \(H(\alpha,\beta)\) que al mismo tiempo son funciones enteras \(I(\sigma)\); también aquí vale la recíproca por (5), (6) y (7). Así, para probar la existencia de una base entera de todos los \(I(a)\), basta probarla para todos los invariantes enteros \(H(\alpha,\beta)\) que al mismo tiempo son funciones enteras \(I(\sigma)\); y la base
\[
        H_1(\alpha,\beta)=I_1(\sigma),\ldots,H_k(\alpha,\beta)=I_k(\sigma)
\]
proporciona entonces la base \(I_1(a),\ldots,I_k(a)\) para los \(I(a)\).

Hágase recorrer a \(H(\alpha,\beta)=I(\sigma)\) la totalidad de los invariantes enteros de (4). Como por (3) todo invariante es homogéneo en las \(\sigma\), y por (4) cada \(\sigma\) es homogénea y lineal en los coeficientes de cada forma lineal, \(H(\alpha,\beta)\) es homogénea en cada fila \(\alpha_i,\beta_i\) y de la misma dimensión en cada fila; además, por supuesto, es una función simétrica de las \(n\) filas. Recíprocamente, por el teorema fundamental sobre las funciones simétricas, toda función simétrica entera \(H(\alpha,\beta)\) de las \(n\) filas, separadamente homogénea y de la misma dimensión en cada fila, es una función polinómica entera de las \(\sigma\). Así, por el argumento precedente, es un invariante \(I(\sigma)\) tan pronto como \(H(\alpha,\beta)\) es invariante. La totalidad de los invariantes enteros de (4) queda por tanto caracterizada como la totalidad de las funciones enteras \(H(\alpha,\beta)\) con las siguientes propiedades:
\[
\begin{array}{ll}
1) & \text{invariante entero de las \(n\) formas lineales \(\alpha_i x+\beta_i y\);}\\
2) & \text{separadamente homogénea y de la misma dimensión en cada una de las \(n\) filas \(\alpha_i,\beta_i\);}\\
3) & \text{simétrica en las \(n\) filas \(\alpha_i,\beta_i\).}
\end{array}
\]
Por el teorema especial de finitud I, \(H(\alpha,\beta)\), a causa de 1), se convierte en una función polinómica entera de los determinantes
\[
        (ik)=\alpha_i\beta_k-\beta_i\alpha_k;
\]
por tanto
\[
        H(\alpha,\beta)=G((ik)).
\]
En esta representación la simetría también puede expresarse formalmente aplicando todas las permutaciones de las \(n\) filas a \(G((ik))\):
\begin{equation}
        n!\,H(\alpha,\beta)
        =
        G^{(1)}((ik))+G^{(2)}((ik))+\cdots+G^{(n!)}((ik))
        =G^*((ik)).
\end{equation}
Junto con cada producto de potencias \(P\) de los \((ik)\), \(G^*\) contiene por tanto también la suma \(Q=P^{(1)}+\cdots+P^{(n!)}\) sobre todas las permutaciones de \(P\). De hecho, \(G^*\) se convierte en una combinación lineal entera de finitamente muchas \(Q\), \(G^*=\sum cQ\). Como las \(Q\) satisfacen las condiciones 1), 2) y 3), y por consiguiente se convierten ellas mismas en invariantes, basta considerar las \(Q\) en lugar de \(G^*((ik))\).

Hágase ahora recorrer a \(P\) el sistema \(\frS\) de todos los productos de potencias de los \((ik)\) que aparecen en las \(Q\), es decir, que satisfacen la condición 2). Si el cociente de dos \(P\) es polinómico en los \((ik)\), entonces también satisface la condición 2), y por tanto pertenece a \(\frS\). Así, por el Teorema II, hay finitamente muchas de tales \(P\), digamos \(P_1,\ldots,P_\rho\), tales que toda \(P\) es representable en la forma
\[
        P=P_1^{\lambda_1}\cdots P_\rho^{\lambda_\rho}
\]
con exponentes no negativos \(\lambda\). Aplicando todas las permutaciones se obtiene
\[
        Q=
        P_1^{(1)\lambda_1}\cdots P_\rho^{(1)\lambda_\rho}
        +\cdots+
        P_1^{(n!)\lambda_1}\cdots P_\rho^{(n!)\lambda_\rho}.
\]
Así \(Q\) se convierte en una función simétrica entera de las \(n!\) filas de cantidades, cada una con \(\rho\) elementos:
\[
        P_1^{(1)}\ldots P_\rho^{(1)};
        \quad
        P_1^{(2)}\ldots P_\rho^{(2)};
        \quad \ldots;\quad
        P_1^{(n!)}\ldots P_\rho^{(n!)};
\]
y por tanto, por el Teorema III, es una función polinómica entera de finitamente muchas funciones simétricas de estas filas. Estas funciones de base son ellas mismas cantidades \(Q\) o funciones simétricas elementales de los \(n!\) elementos \(P_i^{(1)},P_i^{(2)},\ldots,P_i^{(n!)}\), y por consiguiente satisfacen las condiciones 1), 2) y 3) y se convierten en invariantes enteros
\begin{equation}
        H_1(\alpha,\beta)=I_1(\sigma);\ldots;\quad
        H_l(\alpha,\beta)=I_l(\sigma).
\end{equation}
Por (10), ya que \(G^*=\sum cQ\), todo \(H(\alpha,\beta)=I(\sigma)\) es por tanto de la forma
\[
        H(\alpha,\beta)=I(\sigma)=\frac1{n!}\Phi(I_1,\ldots,I_l),
\]
donde \(\Phi\) denota una función polinómica entera; recíprocamente, \(\frac1{n!}\Phi\) es un invariante entero tan pronto como es polinómico en \(\alpha,\beta\). Por consiguiente, por el Teorema IV, pueden adjuntarse a (11) finitamente muchos invariantes adicionales \(I_{l+1}(\sigma),\ldots,I_k(\sigma)\), de tal modo que
\[
        H_1(\alpha,\beta)=I_1(\sigma);\ldots;H_l(\alpha,\beta)=I_l(\sigma);
\]
\[
        H_{l+1}(\alpha,\beta)=I_{l+1}(\sigma);\ldots;H_k(\alpha,\beta)=I_k(\sigma)
\]
formen una base entera de todos los invariantes \(H(\alpha,\beta)=I(\sigma)\). Esto prueba el teorema para el caso de una sola forma fundamental.

Si se trata de un sistema de formas fundamentales, entonces, como se señaló arriba, se puede suponer la homogeneidad de los invariantes en los coeficientes de cada forma fundamental. De nuevo se reemplazan las formas fundamentales generales por otras que son, respectivamente, productos de formas lineales. Sus invariantes se convierten entonces en: 1) invariantes de esas formas lineales; 2) separadamente homogéneos en las filas formadas por los coeficientes de las formas lineales, y de igual dimensión en las filas pertenecientes a cada forma fundamental individual; 3) simétricos en las filas pertenecientes a cada forma fundamental individual. Recíprocamente, todo invariante entero de todas las formas lineales que satisfaga estas condiciones vuelve a corresponder a un invariante entero del sistema de formas fundamentales. La demostración es absolutamente paralela a la acabada de dar. A cada producto de potencias \(P=P_1^{\lambda_1}\cdots P_\rho^{\lambda_\rho}\) de los determinantes individuales que satisface las condiciones de peso (Teoremas I y II) hay que aplicarle todas las
\[
        N=n_1!n_2!\cdots n_\mu!
\]
permutaciones, donde \(n_1,\ldots,n_\mu\) denotan los grados de las formas fundamentales. De esta manera \(N\cdot I\) se convierte en una función simétrica entera de \(N\) filas, cada una con \(\rho\) elementos. De ello se obtiene la base entera de todos los \(N\cdot I\), y por consiguiente, mediante el Teorema IV, la base entera de todos los \(I\).

\setcounter{equation}{0}
\section*{§ 3. La base de los invariantes enteros de formas lineales binarias}

Queda por dar la demostración del teorema especial de finitud I. Por un teorema conocido, probado por primera vez por Clebsch, todos los invariantes polinómicos racionales de las \(n\) formas lineales \(\alpha_i x+\beta_i y\) son funciones polinómicas racionales de los determinantes \((ik)\). En particular, por tanto, todos los invariantes enteros de estas formas lineales se convierten en funciones polinómicas racionales, aunque no necesariamente enteras, de esos determinantes.

Se mostrará a continuación que, en vista de las relaciones idénticas existentes entre los \((ik)\), esta representación siempre puede hacerse de manera entera.

En fórmulas, y en el lenguaje de la teoría de módulos, esto se expresa como sigue. Sean \(\xi_{ik}\), \(i<k\), indeterminadas. Entonces la totalidad de todos los polinomios enteros \(F(\xi_{ik})\) con la propiedad de que \(F((ik))\) se anula idénticamente en \(\alpha,\beta\) forma un módulo \(\frM\); la suma de dos tales \(F\), y también el producto de una tal \(F\) por un polinomio entero arbitrario en las \(\xi_{ik}\), vuelve a pertenecer a esa totalidad. En fórmulas:
\[
        F((ik))=0\quad \text{idénticamente en \(\alpha,\beta\)}
\]
implica
\[
        F(\xi_{ik})\equiv0\pmod{\frM}\quad \text{idénticamente en \(\xi_{ik}\)},
\]
y recíprocamente.

Muestro ahora que vale la afirmación completamente análoga siguiente. De
\[
        F((ik))\equiv0\pmod p\quad \text{idénticamente en \(\alpha,\beta\)}
\]
\begin{equation}
        \text{se sigue que}\qquad
        F(\xi_{ik})\equiv0\pmod{(\frM,p)}
        \quad \text{idénticamente en \(\xi_{ik}\)},
\end{equation}
y recíprocamente, para todo número primo \(p\). Por tanto \(\frM\) es también idéntico al módulo \(\frM_p\) correspondiente a la totalidad de todas las relaciones módulo \(p\) entre los \((ik)\); este módulo consta de todos los polinomios enteros \(F(\xi_{ik})\) para los cuales \(F((ik))\) se anula idénticamente módulo \(p\) en \(\alpha,\beta\).

Que (1) afirma la representabilidad entera se ve como sigue. También puede escribirse en la forma: de \(F((ik))=p\cdot H(\alpha,\beta)\), idénticamente en \(\alpha,\beta\), se sigue que
\[
        F(\xi_{ik})-p\cdot G(\xi_{ik})\equiv0\pmod{\frM}
        \quad \text{idénticamente en \(\xi_{ik}\)},
\]
y por tanto, por la definición de \(\frM\),
\[
        F((ik))=p\cdot G((ik))\quad \text{idénticamente en \(\alpha,\beta\)}.\footnote{Esta última conclusión prueba simultáneamente la recíproca de (1).}
\]
Así \(p\cdot H(\alpha,\beta)=p\cdot G((ik))\), o
\[
        H(\alpha,\beta)=G((ik)).
\]
Por consiguiente, si se da una representación
\[
        H(\alpha,\beta)
        =
        \frac1{p_1^{\lambda_1}\cdots p_\rho^{\lambda_\rho}}
        F((ik)),
\]
donde \(p_1,\ldots,p_\rho\) son números primos, entonces también se sigue que
\[
        H(\alpha,\beta)
        =
        \frac1{p_1^{\lambda_1-1}\cdots p_\rho^{\lambda_\rho}}
        G((ik)),
\]
y repitiendo esto un número finito de veces se obtiene la integridad de la representación.

Para la demostración de (1), denótense por
\[
        f_{ik,rs}=\xi_{ik}\xi_{rs}-\xi_{ir}\xi_{ks}+\xi_{is}\xi_{kr},
        \qquad i<k<r<s,
\]
las formas en las indeterminadas \(\xi_{ik}\) correspondientes a las relaciones cuadráticas entre los \((ik)\),
\[
        (ik)(rs)-(ir)(ks)+(is)(kr)=0,
\]
y denótese por \(\frN=(f_{ik,rs})\) el módulo de polinomios enteros generado por ellas. De
\[
        F(\xi_{ik})\equiv0\pmod{\frN}
\]
se sigue que
\[
        F((ik))=0\quad \text{idénticamente en \(\alpha,\beta\)},
\]
y
\begin{equation}
\begin{array}{rcl}
        \text{de }F(\xi_{ik})&\equiv&0\pmod{(\frN,p)}\\
        \text{se sigue que }F((ik))&\equiv&0\pmod p\quad\text{idénticamente en \(\alpha,\beta\)}.
\end{array}
\end{equation}
Ahora mostraré, normalizando las clases residuales respecto de \(\frN\), que también vale la recíproca de (2). Entonces \(\frN=\frM=\frM_p\), y (1) quedará probado. La igualdad \(\frN=\frM\) se obtiene como el caso especial \(\frN=\frM_p\) para \(p=0\); el caso \(p=0\) está siempre incluido en las consideraciones que siguen. En otras palabras, los \(f_{ik,rs}\) forman una base entera de módulo de \(\frM_p\), y en particular también de \(\frM_0=\frM\).

\subsection*{a) El caso de dos o tres filas \(\alpha_i,\beta_i\)}

Como no aparecen cuatro índices distintos, \(\frN\) es el módulo nulo. Por tanto hay que mostrar que de
\begin{equation}
        F((ik))\equiv0\pmod p\quad \text{idénticamente en \(\alpha,\beta\)}
\end{equation}
se sigue que
\[
        F(\xi_{ik})\equiv0\pmod p\quad \text{idénticamente en \(\xi_{ik}\)}.
\]
Esto prueba (1) con \(\frM=\frM_p=0\).

Como la condición \(F((ik))=p\cdot H(\alpha,\beta)\) debe satisfacerse por separado para las componentes homogéneas en las filas individuales, basta en todo momento suponer que \(F((ik))\) es homogénea en cada fila individual, y por tanto que \(F(\xi_{ik})\) es homogénea en cada índice. En el caso de dos filas sólo existe el determinante \((12)\) y la indeterminada correspondiente \(\xi_{12}\). Por la homogeneidad supuesta,
\[
        F=c\,\xi_{12}^\nu.
\]
De
\[
        c\,(12)^\nu\equiv0\pmod p\quad \text{idénticamente en \(\alpha,\beta\)}
\]
se sigue, puesto que \((12)\not\equiv0\pmod p\) idénticamente en \(\alpha,\beta\), que \(c\equiv0\pmod p\); de ahí
\[
        F(\xi_{ik})\equiv0\pmod p\quad \text{idénticamente en \(\xi_{ik}\)}.
\]
Esto prueba (3), y por tanto (1).

En el caso de tres filas aparecen los determinantes \((12),(13),(23)\), correspondientes a las indeterminadas \(\xi_{12},\xi_{13},\xi_{23}\). Así
\[
        F=\sum c\,\xi_{12}^{\tau_{12}}\xi_{13}^{\tau_{13}}\xi_{23}^{\tau_{23}}.
\]
Sea \(F\) de grados respectivos \(\sigma_1,\sigma_2,\sigma_3\) en los índices \(1,2,3\). Entonces
\[
        \sigma_1=\tau_{12}+\tau_{13},\qquad
        \sigma_2=\tau_{12}+\tau_{23},\qquad
        \sigma_3=\tau_{13}+\tau_{23},
\]
con la inversa única
\[
        \tau_{12}=\frac12(\sigma_1+\sigma_2-\sigma_3),\quad
        \tau_{13}=\frac12(\sigma_1-\sigma_2+\sigma_3),\quad
        \tau_{23}=\frac12(-\sigma_1+\sigma_2+\sigma_3).
\]
Así cada \(F\) separadamente homogénea contiene a lo sumo un término,
\[
        F=c\,\xi_{12}^{\tau_{12}}\xi_{13}^{\tau_{13}}\xi_{23}^{\tau_{23}},
\]
del cual se sigue, como antes, (3), y por consiguiente (1).

\subsection*{b) El caso de cuatro filas \(\alpha_i,\beta_i\)}

Aparecen las seis indeterminadas \(\xi_{12},\xi_{13},\xi_{14},\xi_{23},\xi_{24},\xi_{34}\) y los determinantes correspondientes; y \(\frN=(f_{12,34})\). Puesto que
\begin{equation}
        \xi_{14}\xi_{23}\equiv \xi_{13}\xi_{24}-\xi_{12}\xi_{34}\pmod{\frN},
\end{equation}
todo producto de potencias de las \(\xi_{ik}\) es congruente módulo \(\frN\) con una combinación lineal entera de tales productos de potencias en los que no aparece el producto \(\xi_{14}\xi_{23}\). Así, para todo polinomio \(F(\xi_{ik})\),
\begin{equation}
        F(\xi_{ik})\equiv F_0(\xi_{ik})\pmod{\frN},
\end{equation}
donde \(F_0\) no contiene ningún producto \(\xi_{14}\xi_{23}\). El polinomio \(F_0\) se llamará el polinomio normalizado de la clase residual de \(F\). Como (4) es separadamente homogéneo, a una \(F\) separadamente homogénea le corresponde de nuevo una \(F_0\) separadamente homogénea.

Pongamos ahora
\begin{equation}
        F_0(\xi_{ik})=G(\xi_{ik})+\xi_{34}F_1(\xi_{ik}),
\end{equation}
donde \(G\) no contiene ningún producto de potencias divisible por \(\xi_{34}\), y además pongamos
\begin{equation}
        [G]_{\xi_{14}=\xi_{13}}=K(\xi_{12},\xi_{13},\xi_{23}).
\end{equation}
De
\[
        F((ik))\equiv0\pmod p\quad \text{idénticamente en \(\alpha,\beta\)}
\]
se sigue de (5), teniendo en cuenta (2), que
\begin{equation}
        F_0((ik))\equiv0\pmod p\quad \text{idénticamente en \(\alpha,\beta\)}.
\end{equation}
Especializando \(\alpha_4=\alpha_3,\;\beta_4=\beta_3\), (6) y (7) dan
\[
        K((ik))\equiv0\pmod p\quad \text{idénticamente en \(\alpha,\beta\)}.
\]
Por el resultado probado en a), por tanto,
\[
        K(\xi_{ik})\equiv0\pmod p\quad \text{idénticamente en \(\xi_{ik}\)}.
\]
De esto se sigue lo que todavía queda por mostrar abajo:
\[
        G(\xi_{ik})\equiv0\pmod p\quad \text{idénticamente en \(\xi_{ik}\)}.
\]
Así (6) da
\[
        F_0(\xi_{ik})\equiv \xi_{34}F_1(\xi_{ik})\pmod p
        \quad \text{idénticamente en \(\xi_{ik}\)}.
\]
Como \((34)\not\equiv0\pmod p\) idénticamente en \(\alpha,\beta\), se sigue de (8) que
\[
        F_1((ik))\equiv0\pmod p\quad \text{idénticamente en \(\alpha,\beta\)},
\]
donde \(F_1(\xi_{ik})\) vuelve a estar normalizado y tiene grado menor en \(\xi_{34}\). Repitiendo el procedimiento un número finito de veces se obtiene
\[
        F_0(\xi_{ik})\equiv\xi_{34}^{\lambda}F_\lambda(\xi_{ik})\pmod p,
\]
donde \(F_\lambda\) es de grado cero en \(\xi_{34}\), y por tanto se anula módulo \(p\) por lo que acaba de probarse. Por consiguiente
\begin{equation}
\begin{array}{rcl}
        F_0(\xi_{ik})&\equiv&0\pmod p\quad \text{idénticamente en \(\xi_{ik}\)},\quad\text{o}\\
        F(\xi_{ik})&\equiv&0\pmod{(\frN,p)}\quad \text{idénticamente en \(\xi_{ik}\)}.
\end{array}
\end{equation}
Esto prueba la recíproca de (2), y por tanto (1), al mismo tiempo en la afirmación más precisa para \(F_0(\xi_{ik})\).\footnote{Esta afirmación más precisa significa que cada clase residual contiene un único polinomio normalizado.}

Sólo resta añadir que \(G(\xi_{ik})\equiv0\pmod p\) se sigue de \(K(\xi_{ik})\equiv0\pmod p\); equivalentemente, de \(G(\xi_{ik})\not\equiv0\pmod p\) se sigue que \(K(\xi_{ik})\not\equiv0\pmod p\). La normalización sólo se necesita para esta parte de la inducción. Bajo la sustitución \(\xi_{14}=\xi_{13}\), los productos de potencias individuales de \(G\) no se anulan. Una anulación de \(K\) sólo podría ocurrir porque productos de potencias distintos de \(G\) correspondieran a productos iguales en \(K\); y esto, a su vez, sólo puede suceder cuando tales productos de potencias difieren meramente por un intercambio de los índices 3 y 4. Supóngase que uno de tales productos de potencias normalizados es, por ejemplo,
\[
        \xi_{12}^{\tau_{12}}\xi_{13}^{\tau_{13}}\xi_{14}^{\tau_{14}}\xi_{24}^{\tau_{24}}.
\]
Debido a la homogeneidad separada supuesta en cada índice, una sustitución de 3 por 4 en una \(\xi_{ik}\) debe ir acompañada de una sustitución de 4 por 3 en otra; un término relevante contendría por tanto, en lugar de \(\xi_{13}\xi_{24}\), el producto \(\xi_{14}\xi_{23}\), que está excluido por la normalización. Exactamente del mismo modo, contra
\[
        \xi_{12}^{\tau_{12}}\xi_{13}^{\tau_{13}}\xi_{14}^{\tau_{14}}\xi_{24}^{\tau_{24}}
\]
sólo podría cancelar un término que contuviera el factor \(\xi_{14}\xi_{23}\), nuevamente excluido por la normalización. De ahí que \(K(\xi_{ik})\equiv0\pmod p\) implique \(G(\xi_{ik})\equiv0\pmod p\).

\subsection*{c) El caso de \(n\) filas \(\alpha_i,\beta_i\)}

El caso de \(n\) filas queda ahora resuelto por inducción completa, como generalización exacta del caso de cuatro filas tratado acaba de tratarse.\footnote{La demostración sigue siendo aplicable para \(n=3\).}
Primero muestro de nuevo que para todo polinomio \(F\) existe una congruencia
\begin{equation}
        F(\xi_{ik})\equiv F_0(\xi_{ik})\pmod{\frN}
        \quad \text{idénticamente en \(\xi_{ik}\)},
\end{equation}
donde \(\frN=(f_{ik,rs})\) y \(F_0(\xi_{ik})\) está normalizado. La normalización debe consistir en lo siguiente: si \(i,k,r,s\) son cuatro índices distintos y \(i<k<r<s\), entonces \(F_0\) no debe contener el producto \(\xi_{is}\xi_{kr}\); es decir, los índices de una \(\xi\) no deben quedar ambos entre los de la otra. Esta normalización no impone restricción en el caso de dos y tres filas, donde todo polinomio ya está normalizado; en el caso de cuatro filas, como se mostró, también existe un polinomio normalizado para cada clase residual. Por tanto puede suponerse probada, para \((n-1)\) filas, la existencia de un polinomio normalizado para cada clase residual.

Escribamos un producto arbitrario de potencias de \(F\) en la forma
\[
        \xi_{i_1n}\cdots\xi_{i_\nu n}\,P,
\]
donde \(P\) ya no contiene el índice \(n\). Por la hipótesis de inducción, \(P\) es congruente módulo \(\frN\) con un polinomio entero normalizado \(P_0\):
\[
        \xi_{i_1n}\cdots\xi_{i_\nu n}P
        \equiv
        \xi_{i_1n}\cdots\xi_{i_\nu n}P_0
        \pmod{\frN}.
\]
Si \(P_0\) no contiene ninguna \(\xi_{rs}\) tal que, para al menos un índice \(i_\lambda\), los dos índices \(r,s\) estén entre \(i_\lambda\) y \(n\), es decir, \(i_\lambda<r<s<n\), entonces ya se ha alcanzado la normalización, pues las hipótesis se cumplen para el producto de cualesquiera dos \(\xi\) de \(P_0\). En caso contrario, sea \(i_\lambda\) un índice de ese tipo, \(i_\lambda<r<s<n\), y sea \(P_0^*\) el producto de potencias de \(P_0\) que contiene \(\xi_{rs}\). Como
\[
        \xi_{i_\lambda n}\xi_{rs}\equiv
        \xi_{i_\lambda s}\xi_{rn}-\xi_{i_\lambda r}\xi_{sn}
        \pmod{\frN},
\]
se obtiene
\begin{align*}
 &\xi_{i_1n}\cdots\xi_{i_\lambda n}\cdots\xi_{i_\nu n}P_0^*\\
 &\equiv
 \xi_{i_1n}\cdots\xi_{rn}\cdots\xi_{i_\nu n}Q
 +
 \xi_{i_1n}\cdots\xi_{sn}\cdots\xi_{i_\nu n}R
 \pmod{\frN}\\
 &\equiv
 \xi_{i_1n}\cdots\xi_{rn}\cdots\xi_{i_\nu n}Q_0
 +
 \xi_{i_1n}\cdots\xi_{sn}\cdots\xi_{i_\nu n}R_0
 \pmod{\frN},
\end{align*}
donde \(Q_0\) y \(R_0\) denotan de nuevo los polinomios normalizados de las clases residuales de \(Q\) y \(R\). Estos existen por la hipótesis de inducción, ya que \(Q\) y \(R\) no contienen el índice \(n\), y son ciertamente polinomios enteros, puesto que \(Q\) y \(R\) son simplemente productos de potencias. Como \(i_\lambda<r<s<n\), se tiene además
\[
 ((n-i_1)+\cdots+(n-r)+\cdots+(n-i_\nu))
 <
 ((n-i_1)+\cdots+(n-i_\lambda)+\cdots+(n-i_\nu)),
\]
\[
 ((n-i_1)+\cdots+(n-s)+\cdots+(n-i_\nu))
 <
 ((n-i_1)+\cdots+(n-i_\lambda)+\cdots+(n-i_\nu)).
\]
Así, al pasar de \(P_0^*\) a \(Q_0,R_0\), esta suma de diferencias ha disminuido; como la suma es necesariamente \(\ge\sigma\), el procedimiento debe terminar después de un número finito de pasos para cada producto de potencias \(P_0^*\). Por tanto
\[
        F(\xi_{ik})
        =
        \sum \xi_{i_1n}\cdots\xi_{i_\nu n}P^{(1)}
        \equiv
        \sum \xi_{j_1n}\cdots\xi_{j_\sigma n}S_0^{(j)}
        =
        F_0(\xi_{ik})\pmod{\frN},
\]
donde el polinomio entero \(F_0(\xi_{ik})\) está necesariamente normalizado, pues de lo contrario la suma de diferencias podría reducirse todavía más. Esto prueba (10). Como, por la hipótesis de inducción para \((n-1)\) índices, toda \(F\) separadamente homogénea posee una \(F_0\) normalizada separadamente homogénea, el procedimiento recién dado muestra que esto también vale para \(n\) índices. De ahí que en lo sucesivo sólo deban considerarse polinomios separadamente homogéneos.

Como en b), pongamos
\begin{equation}
        F_0(\xi_{ik})=G(\xi_{ik})+\xi_{n-1,n}F_1(\xi_{ik}),\footnote{La identidad correspondiente a (10) y (11),
        \[
          F((ik))=F_0((ik))=G((ik))+(n-1,n)F_1((ik)),
        \]
        es un análogo del desarrollo en serie de Clebsch-Gordan tan pronto como las dos filas \((n-1)\) y \((n)\) se sustituyen por filas de variables \(x,y\), de modo que \((n-1,n)\) se sustituye por \((xy)\). Pero mientras que Clebsch-Gordan implica un desarrollo mediante polares, la normalización de \(F_0\) corresponde a ``términos iniciales'' de polares, y de ese modo fuerza la integridad que no vale para Clebsch-Gordan. La demostración conocida de que, prescindiendo de la integridad, los \(f_{ik,rs}\) forman una base de \(\frM\), se lleva a cabo precisamente mediante el desarrollo en serie de Clebsch-Gordan (Gordan-Kerschensteiner, \emph{Vorlesungen über Invariantentheorie}, p. 132).}
\end{equation}
donde \(G\) no contiene ningún producto de potencias divisible por \(\xi_{n-1,n}\), y pongamos además
\begin{equation}
        [G]_{\xi_{in}=\xi_{i,n-1}}=K(\xi_{ik}).
\end{equation}
Como muestra la definición de normalización, \(K\), al igual que \(G\), está normalizado.

De nuevo existe una correspondencia biunívoca entre \(G\) y \(K\). Productos de potencias distintos de \(G\) pueden corresponder al mismo producto en \(K\) sólo si difieren meramente por un intercambio de los índices \(n\) y \((n-1)\). Supongamos, por ejemplo, que \(G\) es de grado \(\rho\) en el índice \((n-1)\) y de grado \(\sigma\) en el índice \(n\), y que
\[
        \chi_1\le\chi_2,\ldots,\le\chi_\rho\le\chi_{\rho+1},\ldots,\le\chi_{\rho+\sigma}
\]
son los índices conectados con \((n-1)\) y \(n\) en un producto de potencias de \(G\). Su número debe ser \(\rho+\sigma\), pues por (11) la indeterminada \(\xi_{n-1,n}\) no aparece en \(G\). Por la normalización, el producto de potencias es por tanto de la forma
\[
        \xi_{\chi_1,n-1}\cdots\xi_{\chi_\rho,n-1}
        \xi_{\chi_{\rho+1},n}\cdots\xi_{\chi_{\rho+\sigma},n}\,P(\xi_{ik}),
\]
donde \(P(\xi_{ik})\) ya no contiene los índices \(n\) y \((n-1)\). Todo intercambio de \(n\) por \((n-1)\) daría, para \(i\ne j\), el factor \(\xi_{in}\xi_{j,n-1}\) con \(i<j\); entonces \(j,n-1\) queda entre \(i\) y \(n\), y el término no estaría normalizado. Así, para \(\chi_1,\chi_2,\ldots,\chi_{\rho+\sigma}\) y \(P\) fijos, \(G\) contiene sólo un producto de potencias. Esto prueba que productos de potencias distintos de \(G\) corresponden a productos distintos en \(K\). Por tanto
\begin{equation}
        K(\xi_{ik})\equiv0\pmod p\quad\Longrightarrow\quad
        G(\xi_{ik})\equiv0\pmod p.
\end{equation}

De
\[
        F((ik))\equiv0\pmod p\quad \text{idénticamente en \(\alpha,\beta\)}
\]
se sigue ahora de (10), teniendo en cuenta (2), que
\begin{equation}
        F_0((ik))\equiv0\pmod p\quad \text{idénticamente en \(\alpha,\beta\)}.
\end{equation}
Especializando \(\alpha_n=\alpha_{n-1},\;\beta_n=\beta_{n-1}\), (11) y (12) dan
\begin{equation}
        K((ik))\equiv0\pmod p\quad \text{idénticamente en \(\alpha,\beta\)}.
\end{equation}
Como \(K(\xi_{ik})\) está normalizado y contiene sólo \((n-1)\) índices, se sigue de lo probado en a) y en b), fórmula (9), que
\[
        K(\xi_{ik})\equiv0\pmod p\quad \text{idénticamente en \(\xi_{ik}\)}.
\]
Luego, por (13),
\[
        G(\xi_{ik})\equiv0\pmod p\quad \text{idénticamente en \(\xi_{ik}\)}.
\]
De aquí (11) da
\[
        F_0(\xi_{ik})\equiv \xi_{n-1,n}F_1(\xi_{ik})\pmod p
        \quad \text{idénticamente en \(\xi_{ik}\)}.
\]
Como \((n,n-1)\not\equiv0\pmod p\) idénticamente en \(\alpha,\beta\), se sigue de (14) que
\[
        F_1((ik))\equiv0\pmod p\quad \text{idénticamente en \(\alpha,\beta\)},
\]
donde \(F_1(\xi_{ik})\) está de nuevo normalizado y tiene menor grado en \(\xi_{n-1,n}\) que \(F_0\). Repitiendo el proceso un número finito de veces se obtiene
\[
        F_0(\xi_{ik})\equiv\xi_{n-1,n}^{\lambda}F_\lambda(\xi_{ik})\pmod p,
\]
donde \(F_\lambda\) es de grado cero en \(\xi_{n-1,n}\), y por tanto se anula módulo \(p\) por lo que acaba de probarse. Así se obtiene
\begin{equation}
\begin{array}{rcl}
        F_0(\xi_{ik})&\equiv&0\pmod p\quad \text{idénticamente en \(\xi_{ik}\)},\quad\text{o}\\
        F(\xi_{ik})&\equiv&0\pmod{(\frN,p)}\quad \text{idénticamente en \(\xi_{ik}\)}.
\end{array}
\end{equation}
Esto prueba la recíproca de (2), y por tanto (1). Al mismo tiempo prueba, en el caso de \(n\) índices, y por tanto en general, la afirmación más precisa usada en el paso de inducción para \(F_0(\xi_{ik})\).

\setcounter{equation}{0}
\section*{§ 4. Finitud de los invariantes enteros de grupos finitos}

La demostración elemental de finitud dada en mi nota ``El teorema de finitud para invariantes de grupos finitos''\footnote{\emph{Math. Ann.} 77, p. 89 (1915).} admite, mediante los teoremas especiales de finitud III y IV del § 1, el refinamiento a la integridad, mientras que la demostración usual basada en el teorema de Hilbert sobre la base de módulos no proporciona esa integridad. Incluso el uso de la base entera de módulos sólo da una representación en la que aparece en el denominador una potencia arbitrariamente alta del entero \(h\), el orden del grupo.

Uso en todo momento la notación de la nota recién citada. Sea el grupo finito \(\frH\) formado por las \(h\) transformaciones lineales \(A_1,\ldots,A_h\) de determinante no nulo,\footnote{Bastaría suponer que \(\frH\) es simplemente isomorfo a un grupo finito abstracto; esto es, si se anula un determinante de transformación de los \(A\), entonces \(\frH\) debería ser semejante a un grupo \(\begin{pmatrix}\frS&0\\0&0\end{pmatrix}\), donde los determinantes de las sustituciones de \(\frS\) no se anulan. En este caso los invariantes de \(\frH\) coinciden con los invariantes de \(\frS\), de modo que la condición de no anulación de los determinantes no es en realidad una restricción de generalidad.} donde \(A_k\) denota la transformación lineal
\[
        x_1^{(k)}=\sum_{\nu=1}^{n}a_{1\nu}^{(k)}x_\nu,
        \quad\ldots,\quad
        x_n^{(k)}=\sum_{\nu=1}^{n}a_{n\nu}^{(k)}x_\nu,
\]
o, brevemente, \((x^{(k)})=A_k(x)\). Así el grupo \(\frH\) envía la fila \((x)\), con elementos \(x_1,\ldots,x_n\), a las filas \((x^{(k)})\), con elementos \(x_1^{(k)},\ldots,x_n^{(k)}\). Se supone que los coeficientes \(a_{i\nu}^{(k)}\) son enteros algebraicos; esto no restringe la generalidad, puesto que por un teorema de J. Schur\footnote{Über Gruppen linearer Substitutionen mit Koeffizienten aus einem algebraischen Zahlkörper. Satz VII. \emph{Math. Ann.} 71, p. 555 (1912).} todo grupo finito de transformaciones lineales es semejante a uno para el cual los \(a_{i\nu}^{(k)}\) se vuelven enteros algebraicos.

Por invariante absoluto entero del grupo se entiende una función polinómica racional de \(x_1,\ldots,x_n\), con enteros algebraicos del cuerpo numérico \(\frK\) generado por los \(a_{i\nu}^{(k)}\) como coeficientes, que permanece idénticamente inalterada bajo la aplicación de \(A_1,\ldots,A_h\). Para tal invariante \(f(x)\) se tiene
\begin{equation}
        f(x)=f(x^{(1)})=\cdots=f(x^{(h)})
        =\frac1h\sum_{k=1}^{h}f(x^{(k)}).
\end{equation}
Recíprocamente, toda función simétrica entera de las \(h\) filas de cantidades \((x^{(k)})\) se convierte también en un invariante entero del grupo.

Sean \(\omega_1,\omega_2,\ldots,\omega_\rho\) una base de todos los enteros de \(\frK\). Entonces todo invariante puede representarse como
\[
        f(x)=\omega_1 f_1(x)+\cdots+\omega_\rho f_\rho(x),
\]
donde \(f_1(x),\ldots,f_\rho(x)\) tienen coeficientes enteros racionales. Así, por (1),
\[
        h\,f(x)
        =
        \omega_1\sum_{k=1}^{h}f_1(x^{(k)})
        +\cdots+
        \omega_\rho\sum_{k=1}^{h}f_\rho(x^{(k)}).
\]
Aquí las sumas individuales son funciones simétricas de las \(h\) filas de cantidades \((x^{(k)})\), y además con coeficientes enteros racionales. Por el Teorema III, por tanto, pueden expresarse polinómica y racionalmente, con coeficientes enteros racionales, en términos de finitamente muchas funciones simétricas entero-racionales de esas filas. Por el argumento precedente, estas funciones de base \(I_1(x),\ldots,I_\sigma(x)\) se convierten en invariantes enteros del grupo, en general con coeficientes algebraico-enteros, ya que \(\frH\) no tiene por qué contener, junto con \(A\), todas sus transformaciones algebraicamente conjugadas.

Existe por consiguiente una representación
\begin{equation}
        h\cdot f(x)=\omega_1G_1(I)+\cdots+\omega_\rho G_\rho(I),
\end{equation}
donde \(G_1,\ldots,G_\rho\) tienen coeficientes enteros racionales.

Sean ahora \(w_1,\ldots,w_\rho\) indeterminadas, y hágase recorrer a
\[
        L=w_1G_1(I)+\cdots+w_\rho G_\rho(I)
\]
todas las formas lineales en las \(w\) que, bajo la sustitución \(w_i=\omega_i/h\), corresponden a invariantes \(f(x)\) en la representación (2). El teorema de Hilbert sobre la base entera de módulos da
\begin{equation}
        L=A_1(I)L_1+\cdots+A_\tau(I)L_\tau,
\end{equation}
donde los \(A_i\) tienen coeficientes enteros racionales y donde, puesto que todos los \(L\) son lineales en las \(w\), los \(A_i\) ya no contienen las \(w\). Si la sustitución \(\omega_i=w_i/h\) envía \(L_1,\ldots,L_\tau\) a los invariantes \(\mathfrak I_1(x),\ldots,\mathfrak I_\tau(x)\), entonces (3) se convierte en
\begin{equation}
        f(x)=A_1(I)\mathfrak I_1(x)+\cdots+A_\tau(I)\mathfrak I_\tau(x).
\end{equation}
Esta representación muestra que
\[
        I_1(x),\ldots,I_\sigma(x);\mathfrak I_1(x),\ldots,\mathfrak I_\tau(x)
\]
forman una base de todos los invariantes enteros del grupo \(\frH\), de tal manera que todo invariante es representable como una función polinómica racional, con coeficientes enteros racionales, de esos finitamente muchos invariantes.

Si, en particular, el grupo \(\frH\) tiene la propiedad de que junto con cada transformación \(A\) contiene también todas las \(A',A'',\ldots\) con coeficientes numéricos algebraicamente conjugados, entonces las funciones simétricas entero-racionales de las \(h\) filas \((x^{(k)})\) se convierten en funciones de \(x\) con coeficientes enteros racionales. Esto vale en particular también para las funciones de base de las funciones simétricas. Así, si se restringe uno a invariantes \(f(x)\) con coeficientes enteros racionales, entonces las funciones de base
\[
        I_1(x),\ldots,I_\sigma(x);\mathfrak I_1(x),\ldots,\mathfrak I_\tau(x)
\]
tienen también coeficientes enteros racionales. Bajo esta hipótesis especial, que en general no satisfacen los grupos \(\frH\), la finitud entera vale por tanto también para el subdominio de invariantes con coeficientes enteros racionales.





% ---- Paper 16 appended in current round ----
\clearpage
\begin{center}
{\Large\bfseries 16. Sobre el desarrollo en serie en la teoría de formas}\par
\vspace{1em}
\emph{Math. Ann.} 81 (1920), pp. 25--30
\end{center}

Por ``desarrollo en serie de la teoría de formas'' se entiende, como es sabido, por una parte la generalización del desarrollo en serie de Clebsch--Gordan, es decir, del desarrollo de una forma que contiene dos filas binarias según potencias del determinante de esas filas, cuyos coeficientes son polares de formas de una sola fila de variables o, lo que es idéntico, desaparecen al aplicar el proceso \(\Omega\), a formas de un número arbitrario de series de \(n\) variables cada una; y por otra parte, como generalización de las formas normales que aparecen para las ``coordenadas complejas'' \(t_{ik}\), un desarrollo según formas normales para formas que contienen simultáneamente magnitudes de distintos niveles \(t_i,t_{ik},t_{ikl},\ldots\). Este segundo tipo de desarrollo puede obtenerse, sin embargo, del primero mediante procesos invariantes de diferenciación, como lo llevaron a cabo ante todo Mertens y Deruyts.\footnote{F. Mertens, Über invariante Gebilde quaternärer Formen (Sitzber. d. Ak. d. Wiss. Wien 98, Abt. IIa (1889)) para el caso cuaternario. Las formas normales generales se encuentran en J. Deruyts: Sur la réduction des fonctions invariantes dans le système des variables géométriques, Bulletins de l'Ac. R. de Belgique 24 (1892). Sin conocer ese trabajo, y en estrecha dependencia de Mertens, di el desarrollo general según formas normales en el § 7 del trabajo: Zur Invariantentheorie der Formen von \(n\) Variabeln, J. f. M. 139 (1910).}

Los distintos desarrollos de la primera especie resultan todos, como se bosqueja brevemente, por ejemplo, en \emph{Math. Ann.} 77, loc. cit. III, como consecuencia del desarrollo de una forma de \(n\) filas de \(n\) variables cada una según potencias del determinante de esas filas, donde los coeficientes son polares de formas de sólo \((n-1)\) filas, derivadas ellas mismas de la forma inicial por formación de polares y por el proceso \(\Omega\).\footnote{Compárense las referencias bibliográficas loc. cit. Desde entonces apareció, también basado en el método formal: Schouten, Over reeksontwikkelingen van algebraische vormen met verschillende rijen van variabelen van verschillenden graad, Verslag Ak. Amsterdam 27 (1919), p. 1481.} El desarrollo original de Clebsch--Gordan \((n=2)\) había ido todavía un paso más allá: los procesos polares especiales que allí aparecen están dados explícitamente, incluso con los coeficientes numéricos.

Lo que ahora quiero añadir a la nota indicada, donde el desarrollo en serie se concibe más conceptualmente como un problema de familias lineales de formas, es una representación explícita hasta coeficientes numéricos, obtenida por especialización de las investigaciones de E. Fischer sobre sistemas generales de módulos.\footnote{E. Fischer, Über die Differentiationsprozesse der Algebra, J. f. M. 148 (1917).}

A esta clasificación conduce la interpretación de la forma normal en las \(t_{ik}\). En efecto, si se sustituyen las ``coordenadas complejas'' \(t_{ik}\) por las ``coordenadas de recta''
\[
        p_{ik}=(x_i y_k-x_k y_i),
\]
entonces, a causa de las relaciones idénticas existentes entre las \(p_{ik}\), la representación de una forma \(F(p_{ik})\) ya no es unívoca. La univocidad puede conseguirse exigiendo que entre los coeficientes de \(F\) rijan las mismas relaciones lineales que entre los productos de potencias de las \(p_{ik}\) que aparecen en \(F\).\footnote{Compárese Clebsch, Über eine Fundamentalaufgabe der Invariantentheorie, § 4, Abh. d. K. G. d. Wissensch. zu Göttingen 17 (1872).} Para esta forma normal \(\Phi\) vale por tanto:
\begin{equation}
\begin{aligned}
        F(p_{ik})&=\Phi(p_{ik}) &&[\text{idénticamente en }x,y],\quad\text{o}\notag\\
        F(t_{ik})&\equiv\Phi(t_{ik}) &&\modu{M},
\end{aligned}
\tag{1}
\end{equation}
donde \(M\) designa el módulo de todas las relaciones idénticas entre las \(p_{ik}\), es decir, consta de la totalidad de las formas \(G(t_{ik})\) para las cuales \(G(p_{ik})\) desaparece idénticamente en \(x,y\). \(\Phi(t_{ik})\) es la ``forma normal'' en coordenadas complejas, mientras que las funciones de base de \(M\) se convierten en formas cuadráticas en las \(t_{ik}\). Por iteración del procedimiento con respecto a los factores de estas funciones de base surge el desarrollo en serie completo. Consideraciones correspondientes valen cuando aparecen simultáneamente las \(t_i,t_{ik},t_{ikl},\ldots\).

A una línea de pensamiento análoga puede reducirse el desarrollo según polares. Si, por ejemplo, en una forma \(F(x,y,z)\) se sustituyen las tres filas ternarias \(x=(x_1,x_2,x_3),y,z\) por filas \(\xi,\eta,\zeta\) que se componen linealmente de sólo dos de ellas, por ejemplo \(\xi=x,\eta=y;\zeta=\lambda_1x+\lambda_2y\), entonces la representación de \(F(\xi,\eta,\zeta)\) vuelve a no ser unívoca. Puede prescribirse de nuevo para los coeficientes las mismas relaciones lineales que para los productos de potencias de \(\xi,\eta,\zeta\) que aparecen en \(F\). Como aquí el módulo \(M\) de todas las relaciones idénticas tiene como función de base el determinante \(\Delta=(xyz)\) de las tres filas (loc. cit. III, lema auxiliar a)), aquí en lugar de [1] aparece:
\begin{equation}
\begin{aligned}
        F(\xi,\eta,\zeta)&=\Phi(\xi,\eta,\zeta)&&[\text{idénticamente en }x,y],\quad\text{o}\notag\\
        F(x,y,z)&\equiv\Phi(x,y,z)&&\modu{\Delta}.
\end{aligned}
\tag{2}
\end{equation}
Se ve que \(\Phi\) surge por formación de polares (compárese la fórmula (2) loc. cit.); el desarrollo en serie completo se obtiene de nuevo por iteración.

La formación de las formas normales \(\Phi\) queda así subordinada en ambos casos al siguiente problema: si en \(F(z)\) se sustituyen las indeterminadas \(z_i\) por polinomios \(f_i(\xi)\) en parámetros \(\xi\), de modo que la representación \(F(f_i(\xi))\) ya no sea unívoca, entonces la univocidad debe fijarse exigiendo que entre los coeficientes rijan las mismas relaciones lineales que entre los productos de potencias de las \(f_i(\xi)\) que aparecen en \(F\). Así se tiene
\begin{equation}
\begin{aligned}
        F(f_i(\xi))&=\Phi(f_i(\xi))&&[\text{idénticamente en }\xi],\quad\text{o}\notag\\
        F(z)&\equiv\Phi(z)&&\modu{M},
\end{aligned}
\tag{3}
\end{equation}
donde \(M\) designa el módulo de todas las relaciones idénticas entre las \(f_i(\xi)\).

Para esta forma normal \(\Phi(z)\), E. Fischer dio en el § 11, I, de su trabajo citado, un desarrollo explícito hasta coeficientes numéricos, a saber, mediante operaciones lineales. Al mismo tiempo, sin embargo, en la Introducción III y en el § 6 II dio también, para la diferencia \(F-\Phi\) perteneciente a \(M\), y bajo el supuesto de conocer una base de módulo, una expresión explícita en el mismo sentido; y al combinar estos desarrollos en el § 11 II sustituyó así la fórmula [3] por una representación explícita. La cuestión todavía abierta de los coeficientes numéricos es, según el § 12 del mismo trabajo, idéntica a la búsqueda de los valores propios y funciones propias de las operaciones lineales (fórmula 82).

Especializo ahora, en 1., las fórmulas de Fischer al caso del desarrollo en serie según polares y obtengo de ellas por iteración el desarrollo en serie. La operación para determinar \(\Phi\) pasa así a procesos polares determinados, lo cual es considerablemente más preciso que la formulación conocida hasta ahora y mencionada arriba. En la operación para determinar \(F-\Phi\) intervienen mezclados la multiplicación por el determinante \(\Delta\) y el proceso \(\Omega\). Para llegar a la forma usual, no explícita, hay que observar además que \(\Omega\Delta\) puede expresarse mediante procesos polares.

En este punto debe añadirse también una corrección. En el lugar citado (p. 101, línea 6 desde arriba) supuse inadvertidamente como válida en general la transformación idéntica
\[
        \Omega(\Delta\psi)=c_1\psi+c_2\Delta\cdot\Omega(\psi),
\]
que sólo vale en el dominio binario, e hice con ello el paso de la fórmula fundamental (2), que da la expresión de la forma normal \(\Phi\) en [2], al desarrollo en serie (3). Por tanto, las consideraciones de allí sólo dan la fórmula (2), y como caso especial de ella el teorema de reducción (1); pero, al faltar una expresión explícita para la diferencia \(F-\Phi\), no bastan para fundamentar el desarrollo en serie.

En 2. obtengo igualmente, por especialización, la forma normal \(\Phi(t_{ik})\) de [1]; más generalmente la forma normal \(\Phi(t_i,t_{ik},t_{ikl},\ldots)\), lo cual conduce a fórmulas explícitas conocidas. Pero como el establecimiento del desarrollo en serie completo exige conocer una base de módulo de \(M\), aquí es preferible el camino de la teoría de invariantes (compárese la literatura citada), que no exige ese conocimiento sino que, por el contrario, proporciona una base de módulo mediante el desarrollo en serie.

1. Según Fischer, § 11 (fórmula (19) y líneas siguientes), la forma normal \(\Phi(z)\) de [3], allí denotada \(N(\zeta)\), está dada por:
\begin{equation}
        \Phi(z)=(a_1S+a_2S^2+\cdots+a_rS^r)F(z),
\tag{4}
\end{equation}
donde las constantes \(a_1\ldots a_r\) pertenecen al dominio de coeficientes de las \(f(\xi)\) y dependen sólo del grado de \(F(z)\); el operador \(S\) está definido por
\begin{equation}
\begin{gathered}
        S=AB;\qquad BF(z)=F(f(\xi));\\
        AF(f(\xi))=\left\{\sum z_i f_i\!\left(\frac{\p}{\p\xi}\right)\right\}^{p}F(f(\xi))\footnotemark.
\end{gathered}
\tag{5}
\end{equation}
\footnotetext{Compárese la fórmula (75). En el caso de [3], la suma \(\alpha\) recorre todos los productos de potencias de dimensión \(p\)-ésima, y (75) se especializa por tanto en [5].}
Si en [2] se considera \(F\) como una forma de dimensión \(p\)-ésima sólo en \(z\), entonces \(f_i(\xi)\) se especializa en \(\lambda_1x_i+\lambda_2y_i\); por consiguiente se obtiene:
\[
\begin{gathered}
        BF=F(x,y,\lambda_1x+\lambda_2y),\\
        SF=\left\{\sum z_i\left(\frac{\p}{\p\lambda_1}\frac{\p}{\p x_i}+
        \frac{\p}{\p\lambda_2}\frac{\p}{\p y_i}\right)\right\}^{p}
        F(x,y,\lambda_1x+\lambda_2y),
\end{gathered}
\]
o bien, escrito en procesos polares (compárese loc. cit. II, 2 y 3; también para la notación abreviada):
\begin{equation}
        SF(x,y,z)=\frac1{p!}\left[z\left(\frac{\p}{\p\lambda_1}\frac{\p}{\p x}+
        \frac{\p}{\p\lambda_2}\frac{\p}{\p y}\right)\right]^p
        \left[(\lambda_1x+\lambda_2y)\frac{\p}{\p z}\right]^p F(x,y,z).
\tag{6}
\end{equation}
\emph{Por medio de} [4] \emph{y} [6] \emph{queda dada la expresión explícita buscada para} \(\Phi(x,y,z)\); al mismo tiempo \(\Phi(x,y,z)\) se subordina a la forma general mencionada al comienzo (\(\sum PH\) en (2), loc. cit.), esto es, polares de expresiones que dependen sólo de 2 filas y se obtienen de \(F\) por procesos polares. Los coeficientes aún no determinados \(a_i\) son números racionales, ya que aquí las \(f_i(\xi)\) tienen coeficientes numéricos racionales enteros. Si [2] se escribe en la forma
\begin{equation}
        F=\Phi+\Delta F_1,
\tag{2a}
\end{equation}
entonces, por especialización de la fórmula dada por Fischer en la Introducción III y en el § 6 II (p. 41), se obtiene
\begin{equation}
        F_1=(c_1\Omega+c_2\Omega\Delta\Omega+\cdots+c_i\Omega\Delta\Omega\cdots\Delta\Omega)F\footnotemark,
        \qquad \left[\Omega=\left(\frac{\p}{\p x}\frac{\p}{\p y}\frac{\p}{\p z}\right)\right],
\tag{7}
\end{equation}
\footnotetext{Kerim usa esta fórmula especializada en su disertación de Erlangen, de próxima aparición, en la aplicación a las ``formas de inercia''.}
donde de nuevo los \(c_i\) son números racionales que dependen sólo del grado de \(F\).

La iteración conduce a:
\[
        F_1=\Phi_1+\Delta F_2;\quad F_2=\Phi_2+\Delta F_3;\quad\ldots;\quad
        F_i=\Phi_i+\Delta F_{i+1};\quad\ldots,
\]
donde en cada caso las \(\Phi_i\) son las formas normales correspondientes a las \(F_i\). Como las \(F_i\) son de grado \((p-i)\) en \(z\), de [4] y [6] se obtiene:
\begin{equation}
\begin{gathered}
        \Phi_i=(a_{i1}S_i+a_{i2}S_i^2+\cdots)F_i,\\
        S_iF_i=\frac1{(p-i)!}
        \left[z\left(\frac{\p}{\p\lambda_1}\frac{\p}{\p x}+
        \frac{\p}{\p\lambda_2}\frac{\p}{\p y}\right)\right]^{p-i}
        \left[(\lambda_1x+\lambda_2y)\frac{\p}{\p z}\right]^{p-i}F_i.
\end{gathered}
\tag{8}
\end{equation}
A la fórmula [7] corresponde:
\[
        F_i=(\alpha_{i1}\Omega+\alpha_{i2}\Omega\Delta\Omega+\cdots)F_{i-1},
        \quad\text{y por iteración}
\]
\begin{equation}
        F_i=(\beta_1\Omega^i+\beta_2\Omega^i\Delta\Omega+\cdots
        +\beta_\rho\Omega^{k_1}\Delta\Omega^{k_2}\cdots\Delta\Omega^{k_\lambda}+\cdots)F,
        \quad(k_1+k_2+\cdots+k_\lambda-\lambda+1=i).
\tag{9}
\end{equation}
\emph{Mediante} [8] \emph{y} [9], \emph{en el desarrollo en serie}
\[
        F=\Phi+\Delta\Phi_1+\Delta^2\Phi_2+\cdots+\Delta^i\Phi_i+\cdots+\Delta^p\Phi_p
\]
\emph{se ha alcanzado la representación explícita en el sentido indicado.}

Si se aplica a [7] la transformación idéntica \(\Omega\Delta F=PF\), donde \(P\) designa procesos polares,\footnote{Compárese, por ejemplo, F. Mertens, Über ganze Funktionen von \(m\) Systemen von je \(n\) Unbestimmten, Monatsh. f. Math. u. Phys. 4, fórmula (5).} entonces, por la conmutabilidad de esos procesos con el proceso \(\Omega\), se obtiene la fórmula conocida:
\[
        F_1=P\Omega F;\quad\ldots;\quad F_i=P\Omega^iF;
\]
y de ella, en conexión con [8] o con las consideraciones loc. cit., se sigue la formulación conocida del desarrollo en serie, como por ejemplo en (3) loc. cit. Si se trata de \(n\) filas de \(n\) variables, \(x^{(1)},x^{(2)},\ldots,x^{(n-1)},z\), entonces correspondientemente
\[
        f_i(\xi)=\lambda_1x_i^{(1)}+\lambda_2x_i^{(2)}+\cdots+\lambda_{n-1}x_i^{(n-1)};
\]
y asimismo \(\Delta\) y \(\Omega\) deben definirse para \(n\) filas.

2. Para la forma normal de \(F(t_i,t_{ik},t_{ikl},\ldots)\), las funciones \(f_i(\xi),f_{ik}(\xi),\ldots\) están dadas por los determinantes de una, dos, hasta \((n-1)\) filas
\[
        \xi_i^{(1)},\quad (\xi^{(1)}\xi^{(2)})_{ik},\quad\ldots,\quad
        (\xi^{(1)}\xi^{(2)}\cdots\xi^{(n-1)})_{i_1i_2\cdots i_{n-1}}.
\]
Así, como especialización de [5], se obtiene:
\[
\begin{gathered}
        S=AB;\qquad
        BF=F\bigl(\xi_i^{(1)},(\xi^{(1)}\xi^{(2)})_{ik},\ldots\bigr),\\
        ABF=\left(\sum t_i\frac{\p}{\p\xi_i^{(1)}}\right)^{\lambda_1}
        \left(\sum t_{ik}\left(\frac{\p}{\p\xi^{(1)}}\frac{\p}{\p\xi^{(2)}}\right)_{ik}\right)^{\lambda_2}
        \cdots F\bigl(\xi_i^{(1)},(\xi^{(1)}\xi^{(2)})_{ik},\ldots\bigr),
\end{gathered}
\]
donde \(\lambda_1,\lambda_2,\ldots\) designan los grados en \(t_i,t_{ik},\ldots\).

De consideraciones de la teoría de invariantes (en efecto, todos los invariantes de \(F\) que son lineales en los coeficientes pueden componerse linealmente a partir de \(\Phi\) y de formas de \(M\), por tanto en particular también el \(SF\) que no pertenece a \(M\))\footnote{Esto se sigue, por ejemplo, de los §§ 6 y 7, 2., de mi trabajo citado ``Zur Invariantentheorie der Formen von \(n\) Variabeln''.} se sigue aquí la congruencia \(\Phi\equiv a_1SF\modu{M}\), y por la propiedad de forma normal resulta de ella, en lugar de [4], la relación simple:
\begin{equation}
        \Phi=a_1SF=a_1S\Phi.
\tag{10}
\end{equation}
La segunda igualdad vale porque la aplicación de \(S\) hace desaparecer toda forma de \(M\). La fórmula [10] muestra que la forma normal \(\Phi\) es ella misma al mismo tiempo una función propia; y, de hecho, por la consideración de la teoría de invariantes precedente no hay otras funciones propias que sean a la vez invariantes de \(F\) (compárese, sobre [10], Deruyts, fórmulas (10) y (10')).

\begin{center}
\textbf{Correcciones.}
\end{center}

1. En el trabajo ``Die allgemeinsten Bereiche aus ganzen transzendenten Zahlen'' (\emph{Math. Ann.} 77), en la página 121, línea 8 desde arriba, línea 9 desde arriba y línea 11 desde abajo, el cuerpo \((H,Y)\) debe sustituirse por el ``cuerpo \((H,Y,Y',\ldots,Y^{(\lambda_0-1)})\), donde \(Y',\ldots,Y^{(\lambda_0-1)}\) designan los conjugados de \(Y\) respecto de \(H\)''. Correspondientemente, en p. 120, línea 5 desde arriba, y p. 121, línea 9 desde arriba, \((H,y)\) debe sustituirse por \((H,y,y',\ldots,y^{(\lambda_0-1)})\).

2. En el trabajo ``Gleichungen mit vorgeschriebener Gruppe'' (\emph{Math. Ann.} 78), en la formulación del teorema de irreducibilidad de Hilbert, p. 222, línea 3 desde abajo, en lugar de ``cuerpo numérico \(\Omega\)'' debe decirse más precisamente ``cuerpo numérico algebraico finito \(\Omega\)'', como me señala el señor Ostrowski. En efecto, respecto del cuerpo de todos los números algebraicos, por ejemplo, el grupo de toda ecuación con coeficientes numéricos algebraicos es el grupo identidad. Correspondientemente, en la introducción, por ``dominio arbitrario de racionalidad'' debe entenderse un cuerpo numérico algebraico finito al que todavía pueden adjuntarse finitamente muchos parámetros, o uno de los cuerpos intermedios que surgen así.

\begin{center}
(Aceptado en septiembre de 1919.)
\end{center}


% Paper 17 continuation through Section 9

\clearpage
\providecommand{\mM}{\mathfrak{M}}
\providecommand{\mN}{\mathfrak{N}}
\providecommand{\mL}{\mathfrak{L}}
\providecommand{\mG}{\mathfrak{G}}
\providecommand{\mA}{\mathfrak{A}}
\providecommand{\mB}{\mathfrak{B}}
\providecommand{\mX}{\mathfrak{X}}
\providecommand{\mY}{\mathfrak{Y}}
\providecommand{\mU}{\mathfrak{U}}
\providecommand{\mP}{\mathfrak{P}}
\providecommand{\mQ}{\mathfrak{Q}}
\providecommand{\mS}{\mathfrak{S}}
\providecommand{\mD}{\mathfrak{D}}
\providecommand{\mR}{\mathfrak{R}}
\providecommand{\mV}{\mathfrak{V}}
\providecommand{\ma}{\mathfrak{a}}
\providecommand{\mb}{\mathfrak{b}}
\providecommand{\mx}{\mathfrak{x}}
\providecommand{\my}{\mathfrak{y}}
\providecommand{\mz}{\mathfrak{z}}
\providecommand{\me}{\mathfrak{e}}
\providecommand{\mbarA}{\overline{\mathfrak{A}}}
\providecommand{\mbarB}{\overline{\mathfrak{B}}}
\providecommand{\bara}{\overline{\mathfrak{a}}}
\setcounter{footnote}{0}
\section*{17. En colaboración con W. Schmeidler: Módulos en dominios no conmutativos, en particular procedentes de expresiones diferenciales y de diferencias}
\begin{center}
\emph{Math. Zs.} 8 (1920), pp. 1--35
\end{center}

\begin{center}
\textbf{Índice.}
\end{center}
\begin{center}
\begin{minipage}{0.82\linewidth}
Introducción.\\
\S~1. Observaciones formales sobre expresiones diferenciales y de diferencias.\\
\S~2. El dominio polinómico general.\\
\S~3. Módulos unilaterales y sus grupos residuales.\\
\S~4. Reducibilidad de grupos residuales en dos subgrupos y su relación con el módulo.\\
\S~5. Cálculo con unidades y reducibilidad en $k$ subgrupos.\\
\S~6. Teoremas de finitud.\\
\S~7. Un caso de descomposición única.\\
\S~8. Descomposición aditiva de grupos completamente reducibles en grupos primos. El teorema de isomorfía.\\
\S~9. Conexión entre los conceptos ``isomorfo'' y ``del mismo tipo''.\\
\S~10. Existencia de infinitas descomposiciones de un grupo.\\
\S~11. Relaciones con sistemas de ecuaciones diferenciales y sus integrales.\\
\S~12. Ejemplos.
\end{minipage}
\end{center}

\begin{center}\textbf{Introducción.}\end{center}

La teoría formal de las expresiones diferenciales en una variable conduce a una descomposición en factores irreducibles que es única en cierto sentido.\footnote{Cf. E. Landau, Ein Satz über die Zerlegung homogener linearer Differentialausdrücke in irreduzible Faktoren, \emph{Journal für die reine und angewandte Mathematik} 124 (1902), pp. 115--120, y después A. Loewy, Über reduzible lineare homogene Differentialgleichungen, \emph{Mathematische Annalen} 56 (1903), pp. 549--584.} Pero mientras el teorema correspondiente de álgebra puede trasladarse a polinomios de varias variables, esto ya no es posible para las expresiones diferenciales parciales.\footnote{Cf. un contraejemplo de Landau publicado en H. Blumberg, Über algebraische Eigenschaften von linearen homogenen Differentialausdrücken, disertación, Göttingen 1912, 52 pp., pp. 51--52.} Sin embargo, para expresiones diferenciales de una variable, junto a la representación como producto existe también una representación como mínimo común múltiplo, en la que valen leyes análogas; pues, en dos descomposiciones distintas, los constituyentes individuales pueden asignarse por pares: son ``del mismo tipo''.\footnote{Cf., además de los trabajos citados, A. Loewy, Über vollständig reduzible lineare homogene Differentialgleichungen, \emph{Mathematische Annalen} 62 (1906), pp. 89--117; citado como Loewy.}

Este concepto de mínimo común múltiplo puede interpretarse ahora como concepto de módulos y así trasladarse a $n$ variables; en el presente trabajo esto da la generalización natural de los teoremas indicados. Más allá de ello observamos, sin embargo, que para expresiones diferenciales sólo se utiliza el siguiente hecho: pueden interpretarse como polinomios con multiplicación no conmutativa, para los cuales el grado de un producto es igual a la suma de los grados de los factores. Estas condiciones las satisfacen, por ejemplo, también las expresiones de diferencias; por tanto, los teoremas valen también para ellas, donde hasta ahora eran conocidos igualmente sólo para una variable.\footnote{Cf. G. Wallenberg--A. Guldberg, \emph{Theorie der linearen Differenzengleichungen}, Leipzig y Berlín 1911.}

Desarrollamos, pues, en general, una teoría de módulos cuyos elementos son polinomios con multiplicación no conmutativa; al hacerlo reducimos el estudio de los módulos al de las clases residuales.\footnote{Para el caso de dominios conmutativos, cf. W. Schmeidler, Über Moduln und Gruppen hyperkomplexer Größen, \emph{Mathematische Zeitschrift} 3 (1919), pp. 29--42.} A diferencia del caso especial conmutativo, aquí se puede hablar ciertamente de la suma de clases residuales, pero no de su producto; sólo se puede hablar del producto de una clase residual por un polinomio. El concepto de ``grupo residual'' debe, por ello, modificarse respecto del caso conmutativo. En el caso especial de polinomios en una variable, además, el grupo residual es ``finito'', es decir, posee sólo un número finito de clases residuales linealmente independientes. En general no existe tal ``orden'' finito; los grupos son ``infinitos''. Para nuestros métodos, sin embargo, esta diferencia es irrelevante.

Como en el caso conmutativo, a la representación de un módulo como mínimo común múltiplo de módulos relativamente primos corresponde el procedimiento más sencillo de la descomposición aditiva del grupo residual. En este punto la isomorfía de los grupos residuales resulta esencial; al especializar a expresiones diferenciales de una variable, esta isomorfía se convierte en el concepto de mismo tipo para los módulos asociados, mientras que al especializar a módulos conmutativos se convierte en identidad. De paso, el concepto de mismo tipo también puede trasladarse a $n$ variables, y puede probarse su concordancia con el concepto de isomorfía. El hecho de que en la descomposición aditiva del grupo residual aparezcan siempre sólo finitos sumandos se obtiene trasladando el teorema de Hilbert sobre la base de módulos, con el método de demostración de Gordan; el método original de Hilbert exigiría restricciones sobre las leyes de multiplicación (en la fórmula (6), $a$ a la derecha en lugar de $a_i$, condición cumplida para expresiones diferenciales, pero no para expresiones de diferencias).

Para llegar a un teorema de unicidad debemos especializar los módulos -- como también hace Loewy -- a los ``completamente reducibles'', es decir, a aquellos representables como mínimos comunes múltiplos de módulos primos. Dadas dos descomposiciones distintas, los constituyentes son entonces o bien idénticos por pares, o por lo menos a cada constituyente de una descomposición corresponde un constituyente isomorfo de la otra, y recíprocamente; al mismo tiempo, cada descomposición contiene entonces al menos dos constituyentes isomorfos. Además, de la aparición de dos constituyentes isomorfos en una misma descomposición se deduce la existencia de infinitas descomposiciones; esto vale incluso sin la restricción a módulos primos, cosa que hasta ahora no era conocida ni siquiera en el caso especial de expresiones diferenciales en una variable.

Finalmente, para módulos provenientes de expresiones diferenciales, discutimos la conexión con las integrales y mostramos que, para grupos residuales finitos, el orden es igual al mayor número de integrales linealmente independientes. Así, las funciones arbitrarias aparecen en la integral general precisamente cuando el grupo residual no es finito. El concepto de que dos módulos sean del mismo tipo significa aquí, como en el caso de una variable, una relación ``birracional'' entre las integrales. Algunos ejemplos cierran el trabajo.

El primer impulso para este trabajo fue dado hace varios años por una pregunta de E. Landau a E. Noether sobre la generalización del teorema del producto. Entonces E. Noether sólo observó que la cuestión era de teoría de módulos en el dominio polinómico no conmutativo, pero aún no podía abordarse. Los teoremas de módulos de Lasker conocidos entonces no eran directamente trasladables a este dominio, porque el método de demostración de Lasker descansa en el teorema de que un polinomio se descompone de manera única en factores irreducibles.\footnote{Recientemente E. Noether observó que también los teoremas de Lasker pueden establecerse sin usar este teorema, mediante métodos emparentados con los empleados aquí.} Sólo los métodos de Schmeidler ofrecieron la posibilidad de una transferencia, que luego llevamos a cabo conjuntamente. En detalle, señalamos que la generalización del grupo residual, la transferencia y aplicación del teorema de finitud de Hilbert, y la existencia de infinitas descomposiciones se deben a E. Noether; la generalización del concepto de reducibilidad completa de Loewy, el concepto de isomorfía y su conexión con el concepto de mismo tipo se deben a W. Schmeidler.

\subsection*{\S~1. Observaciones formales sobre expresiones diferenciales y de diferencias.}

1. Para expresiones diferenciales en una variable se ha introducido una formación simbólica de productos. Sea
\[
 a_0(x)\frac{\dd^n y}{\dd x^n}+a_1(x)\frac{\dd^{n-1}y}{\dd x^{n-1}}+\cdots+a_n(x)y=A(y),
\]
\[
 b_0(x)\frac{\dd^m y}{\dd x^m}+b_1(x)\frac{\dd^{m-1}y}{\dd x^{m-1}}+\cdots+b_m(x)y=B(y).
\]
Entonces se define el producto simbólico $AB(y)$ como la expresión diferencial
\[
 a_0(x)\frac{\dd^n B(y)}{\dd x^n}+a_1(x)\frac{\dd^{n-1}B(y)}{\dd x^{n-1}}+\cdots+a_n(x)B(y).
\]
Esta formación de productos es en realidad una multiplicación de dos operadores; se escribe del modo más simple sustituyendo, como de costumbre, $\frac{\dd^n}{\dd x^n}$ por $\left(\frac{\dd}{\dd x}\right)^n$. Entonces $AB$ se convierte en
\[
 \left(a_0\left(\frac{\dd}{\dd x}\right)^n+\cdots+a_n\right)
 \left(b_0\left(\frac{\dd}{\dd x}\right)^m+\cdots+b_m\right)y,
\]
donde al efectuar el producto rigen las reglas ordinarias para diferenciar sumas y productos. Cada uno de tales operadores puede considerarse ahora como un polinomio en la única variable $\frac{\dd}{\dd x}=\xi$, omitiendo la función sobre la que actúa el operador y fijando para la variable $\xi$ la regla de multiplicación correspondiente a la derivación de un producto:
\[
 \frac{\dd}{\dd x}\bigl(a(x)y\bigr)=a(x)\frac{\dd}{\dd x}(y)+a'(x)y,
\]
o, en nuestra notación $(n=1,m=0)$,
\begin{equation}
 \xi\cdot a=a\cdot\xi+a'.
\tag{1}
\end{equation}

Se procede exactamente del mismo modo para expresiones diferenciales parciales lineales
\[
 \left(\sum a_{\alpha_1\ldots\alpha_n}(x_1,\ldots,x_n)
 \frac{\p^{\alpha_1+\cdots+\alpha_n}}
 {\p x_1^{\alpha_1}\cdots\p x_n^{\alpha_n}}\right)y
\]
de una función $y$. Para los $n$ operadores $\frac{\p}{\p x_1},\ldots,\frac{\p}{\p x_n}$ introducimos las variables $\xi_1,\ldots,\xi_n$ y obtenemos así el polinomio
\[
 \sum a_{\alpha_1\ldots\alpha_n}(x_1,\ldots,x_n)
 \xi_1^{\alpha_1}\cdots\xi_n^{\alpha_n}.
\]
El producto de dos polinomios se calcula entonces según las leyes
\begin{equation}
 \xi_i a=a\xi_i+a_i \qquad (i=1,\ldots,n),
\tag{2}
\end{equation}
donde $a_i$ denota la derivada parcial de $a$ respecto de $x_i$; los productos de las $\xi$ entre sí son conmutativos. Salvo la ley conmutativa de la multiplicación, todas las leyes de la suma y de la multiplicación valen para los polinomios así definidos como en el álgebra ordinaria. En particular, para el producto de dos de tales polinomios el grado total, y el grado en cada variable individual, es igual a la suma de los grados correspondientes de los dos factores.

2. Para expresiones de diferencias (cf. Wallenberg, loc. cit.) de la forma
\[
 a_x^{(0)}y_{x+n}+a_x^{(1)}y_{x+n-1}+\cdots+a_x^{(n)}y_x=A_x,
\]
\[
 b_x^{(0)}y_{x+m}+b_x^{(1)}y_{x+m-1}+\cdots+b_x^{(m)}y_x=B_x
\]
se tiene una definición correspondiente del producto
\[
 (AB)_x=a_x^{(0)}B_{x+n}+a_x^{(1)}B_{x+n-1}+\cdots+a_x^{(n)}B_x,
\]
que puede ejecutarse con ayuda de la regla de multiplicación
\[
 (ay)_{x+1}=a_{x+1}y_{x+1}.
\]
Si introducimos el operador $\xi$ para el paso de $x$ a $x+1$, obtenemos $\xi\{ay\}=a^*\xi\{y\}$, donde $a^*=a_{x+1}$. De nuevo se puede omitir la función $y$ sobre la que actúa el operador y se obtiene la ley
\begin{equation}
 \xi a=a^*\xi,
\tag{3}
\end{equation}
donde, junto con $a$, tampoco $a^*$ se anula idénticamente. Exactamente del mismo modo, para $n$ variables,
\begin{equation}
 \xi_i a=a_i\xi_i\qquad (i=1,\ldots,n),
\tag{4}
\end{equation}
donde $\xi_i$ representa el paso de $x_i$ a $x_i+1$, y $a_i$ denota la función obtenida de $a$ sustituyendo $x_i$ por $x_i+1$; por tanto no se anula idénticamente si $a$ no lo hace. Con estas convenciones vuelven a valer todas las reglas de la suma y de la multiplicación, incluida la regla sobre el grado de un producto, salvo la ley conmutativa.

Ambos casos se subordinarán, en el párrafo siguiente, a un dominio polinómico definido de modo completamente general con coeficientes arbitrarios.

\subsection*{\S~2. El dominio polinómico general.}

Partimos de un cuerpo $P$ definido abstractamente de modo arbitrario (cuyos elementos denotamos por letras latinas minúsculas), para el cual valen todas las reglas del álgebra, en particular también la ley conmutativa de la multiplicación con inversión única. A este cuerpo añadimos $n$ indeterminadas $\xi_1,\ldots,\xi_n$; es decir, consideramos el dominio íntegro $J$ formado por los productos $a\xi_1^{\nu_1}\xi_2^{\nu_2}\cdots\xi_n^{\nu_n}$ y por sus sumas finitas, en el que deben valer todas las leyes de la suma y de la multiplicación, con excepción de la ley conmutativa de la multiplicación. En particular, llamamos ``polinomios'' (que en adelante se denotarán siempre con letras latinas mayúsculas) a las sumas finitas de la forma
\[
 F(\xi_1,\ldots,\xi_n)=\sum a_{\nu_1\ldots\nu_n}\xi_1^{\nu_1}\cdots\xi_n^{\nu_n}.
\]
En lugar de la ley conmutativa, sometemos ahora los elementos de $J$ a leyes de producto tales que \emph{el producto de dos polinomios vuelva a ser un polinomio}; así, el dominio de todos los polinomios de $J$ agota ya a $J$. Para esto es evidentemente necesario y suficiente que, para los elementos $\xi_1,\ldots,\xi_n$ y para un elemento arbitrario $a$ de $P$, valgan reglas de producto de la forma
\begin{equation}
 \xi_i a=F_i(\xi_1,\ldots,\xi_n)\qquad (i=1,\ldots,n),
\tag{5}
\end{equation}
donde $F_i(\xi_1,\ldots,\xi_n)$ denota algún polinomio que depende naturalmente de $\xi_i$ y de $a$. (En particular, las reglas (2) y (4) para expresiones diferenciales y de diferencias, respectivamente, son de esta forma.) En efecto, puesto que $\xi_i a$ no es por sí mismo un polinomio, mientras que por hipótesis pertenece al dominio, la ecuación es necesaria. Recíprocamente, por aplicación finitamente repetida de la fórmula (5), todo elemento del dominio $J$ puede transformarse en un polinomio.

Para la unidad del cuerpo, que denotamos por 1, exigiremos además $\xi_i\cdot1=1\cdot\xi_i=\xi_i$. También estipularemos que la multiplicación de $\xi_1,\ldots,\xi_n$ entre sí sea conmutativa.\footnote{Esto se necesita sólo desde \S~6 en adelante, para el teorema de finitud y sus consecuencias.} Entonces todo polinomio de $J$ se escribe exactamente como un polinomio ordinario en $\xi_1,\ldots,\xi_n$ con coeficientes de $P$, y dos polinomios coinciden sólo cuando coinciden todos sus coeficientes.

Más especialmente, en lo que sigue\footnote{También esto se usa sólo desde \S~6 en adelante.} exigimos que, en lugar de (5), rija la regla
\begin{equation}
 \xi_i a=a_i\xi_i+b_i\qquad (a_i\ne0)\quad(i=1,\ldots,n).
\tag{6}
\end{equation}
Aquí también quedan incluidas, en particular, las leyes (2) y (4). Entonces el grado de un producto es igual a la suma de los grados de los factores; esto vale tanto para el grado total como para el grado en cada variable individual, de modo que el producto de dos polinomios que no se anulan idénticamente vuelve a ser distinto de 0.

\subsection*{\S~3. Módulos unilaterales y sus grupos residuales.}

Para nuestro dominio $J$, por ser la multiplicación no conmutativa, se distinguen los siguientes tipos de módulos\footnote{Denotamos los módulos por letras alemanas mayúsculas.}:\footnote{Los ideales unilaterales ya fueron considerados por A. Hurwitz; cf. \emph{Vorlesungen über die Zahlentheorie der Quaternionen} (Berlín, Springer 1919), lección 6. Allí, sin embargo, se trata esencialmente de ideales principales; lo mismo vale para los trabajos de Du Pasquier citados allí en la nota 1.}

1. El \emph{módulo por la derecha}, que, junto con $M$, contiene también $FM$ para un polinomio arbitrario $F$, y junto con $M_1$ y $M_2$ contiene también $M_1+M_2$.

2. El \emph{módulo por la izquierda}, que, junto con $M$, contiene también $MF$, y junto con $M_1$ y $M_2$ contiene también $M_1+M_2$.

A estos dos tipos los designamos con el nombre común de \emph{módulos unilaterales}; los módulos que son a la vez por la derecha y por la izquierda se llamarán \emph{bilaterales}. En lo que sigue nos restringimos a módulos por la derecha y observamos que la teoría de los módulos por la izquierda puede construirse exactamente del mismo modo.

Llamamos congruentes módulo $\mM$ a dos polinomios $F$ y $G$, y escribimos $F\equiv G(\mM)$, si su diferencia es divisible por $\mM$, es decir, si pertenece a $\mM$. La totalidad de todos los polinomios congruentes con un polinomio dado forma una \emph{clase residual} módulo $\mM$. (Las clases residuales se denotarán por letras alemanas minúsculas.) Como de
\[
 F_1\equiv F_2(\mM)\qquad\text{y}\qquad G_1\equiv G_2(\mM)
\]
se deduce $F_1+G_1\equiv F_2+G_2(\mM)$, la suma de dos clases residuales $\mx$ y $\my$ está definida y es de nuevo una clase residual $\mx+\my$. En cambio, como el módulo es unilateral, no se puede hablar del producto de dos clases residuales. En efecto, de
\[
 F_1=F_2+M_1,\qquad G_1=G_2+M_2
\]
se sigue sólo
\[
 F_1G_1=F_2G_2+F_2M_2+M_1G_2+M_1M_2;
\]
pero como $M_1G_2$ no tiene por qué pertenecer a $\mM$ cuando $M_1$ pertenece a $\mM$, $F_1G_1$ y $F_2G_2$ no tienen por qué ser congruentes en general. Por otra parte, si $M_1=0$, es decir, $F_1=F_2=F$, se ve que $G_1\equiv G_2$ implica $FG_1\equiv FG_2$. Así una clase residual arbitraria $\my$ puede multiplicarse por delante por un polinomio $F$, y esto da una clase residual bien definida $F\my=\mz$.

Para el cálculo con clases residuales valen las leyes conmutativa y asociativa de la suma, así como la ley de reversión única de la suma; además, para la multiplicación de clases residuales por polinomios y la suma, vale la ley distributiva $F(\my_1+\my_2)=F\my_1+F\my_2$, como se sigue directamente del cálculo con polinomios módulo $\mM$.

La clase residual a la que pertenece la unidad se llamará clase unidad, o unidad, y se denotará por $\me$. Entonces $F\me=\mx$ es la clase residual a la que pertenece el polinomio $F$. A diferencia de las clases residuales generales, sin embargo, aquí también está bien definido el producto $\mx\me$ y es igual a $\mx$; pues de
\[
 F_1=F_2+M_1,\qquad G_1=1+M_2
\]
y $M_1\cdot1=M_1$ se deduce
\[
 F_1G_1=F_2+M.
\]
La totalidad de las clases residuales módulo $\mM$ se llama \emph{grupo residual de} $\mM$. Un grupo residual posee por tanto las siguientes propiedades: 1. Junto con cada clase residual $\mx$ y un polinomio arbitrario $F$, contiene la clase residual $F\mx$; y junto con las clases residuales $\mx_1$ y $\mx_2$ contiene la clase residual $\mx_1+\mx_2$.\footnote{Esta propiedad podría llamarse, en un sentido más general, la ``propiedad modular'' del grupo residual.} 2. Posee una unidad $\me$, de modo que $F\me$ representa, para cada polinomio $F$, la clase residual $\mx$ de $F$. \emph{Así, cuando} $F$ \emph{recorre todos los polinomios,} $F\me$ \emph{recorre todo el grupo residual.}

\subsection*{\S~4. Reducibilidad del grupo residual en dos subgrupos y su relación con el módulo.}

Para los módulos unilaterales, los conceptos de \emph{divisibilidad}, \emph{máximo común divisor} y \emph{mínimo común múltiplo} permanecen como en el caso bilateral, como muestran las definiciones siguientes.

Se dice que un módulo $\mM$ es divisible por otro módulo $\mN$ si de $M\equiv0(\mM)$ se sigue también $M\equiv0(\mN)$.

Para dos módulos $\mM$ y $\mN$, el \emph{máximo común divisor} $(\mM,\mN)$ se define como la totalidad de todos los polinomios representables en la forma $M+N$, donde $M\equiv0(\mM)$ y $N\equiv0(\mN)$; de nuevo es un módulo. Lo mismo vale para varios módulos.

La totalidad de todos los polinomios divisibles tanto por $\mM$ como por $\mN$ forma un módulo $[\mM,\mN]$, el \emph{mínimo común múltiplo} de $\mM$ y $\mN$. Esta definición se aplica también de manera correspondiente a varios módulos, incluso a un conjunto de módulos de cardinalidad arbitraria.

Dos módulos se llaman relativamente primos cuando su máximo común divisor es el módulo unidad, es decir, la totalidad de todos los polinomios.

Partimos ahora de un módulo $\mM$ que es representable como mínimo común múltiplo de dos módulos relativamente primos $\mN$ y $\mL$, ambos supuestos divisores propios:
\begin{equation}
 \mM=[\mN,\mL].
\tag{7}
\end{equation}
Investigamos el grupo residual $\mG$ de $\mM$. Por hipótesis,
\begin{equation}
 1\equiv N+L(\mM)
\tag{8}
\end{equation}
para dos polinomios $N\equiv0(\mN)$ y $L\equiv0(\mL)$. A partir de esto mostramos
\begin{equation}
 \mL=(\mM,L),\qquad \mN=(\mM,N),
\tag{9}
\end{equation}
donde, por ejemplo, $(\mM,L)$ significa el módulo compuesto por todos los polinomios de la forma $M+FL$. En efecto, primero $(\mM,L)\equiv0(\mL)$. Recíprocamente, de (8) se sigue, para un polinomio arbitrario $F$, que
\begin{equation}
 F\equiv FN(\mL).
\tag{10}
\end{equation}
Ahora bien, si $F\equiv0(\mL)$, entonces $FN\equiv0(\mL)$, luego $\equiv0(\mM)$, y por consiguiente $F\equiv FL(\mM)$. Del mismo modo se prueba $\mN=(\mM,N)$. De paso se deduce, puesto que $\mN$ y $\mL$ son divisores propios de $\mM$, que no puede valer ni $L\equiv0$ ni $N\equiv0(\mM)$.

Sea $\ma$ la clase residual de $N$ módulo $\mM$, y $\mb$ la clase residual de $L$ módulo $\mM$. Entonces, por (8),
\[
 \me=\ma+\mb\qquad(\ma\ne0,\ \mb\ne0).
\]
Como toda clase residual es de la forma $F\me$, la ley distributiva da, para toda clase residual,
\begin{equation}
 F\me=F\ma+F\mb\qquad(\ma\ne0,\ \mb\ne0).
\tag{11}
\end{equation}
Esta representación de una clase residual arbitraria como suma de dos clases residuales de la forma $G\ma$ y $H\mb$ es, sin embargo, \emph{única}. Pues de una relación $0=G\ma+H\mb$, si $K$ es un polinomio de la clase residual $G\ma=-H\mb$, se tiene claramente $K\equiv GN(\mM)$, $K\equiv-HL(\mM)$, luego $K\equiv0(\mN)$, $K\equiv0(\mL)$ y por tanto $K\equiv0(\mM)$; es decir, $G\ma=0$, $H\mb=0$.

Recíprocamente, supóngase ahora que la relación (11) se cumple para todo polinomio $F$, y de modo único en el sentido anterior. Sea entonces $N$ un polinomio de $\ma$ y $L$ un polinomio de $\mb$. Formamos los dos módulos (9) y afirmamos (7) y (8). Esta última sigue de (11) para $F=1$. Además, por la definición de $\mN$ y $\mL$, $\mM$ es un múltiplo común de estos dos módulos, y precisamente el menor. Pues de $F\equiv0(\mL)$ y $F\equiv0(\mN)$ se sigue $F\equiv HL\equiv GN(\mM)$; por tanto $G\ma=H\mb$, $G\ma-H\mb=0$, y de la unicidad de la representación, $G\ma=0$, $H\mb=0$, $F\equiv0(\mM)$. Queda así mostrado que la \emph{representación de un módulo} $\mM$ \emph{como mínimo común múltiplo de dos módulos relativamente primos equivale a la descomposición aditiva única del grupo residual} $\mG$ \emph{de} $\mM$, \emph{tal como la da la fórmula} (11).

Investigamos ahora el sistema $\mA$ de clases residuales de la forma $F\ma$, que queremos relacionar con el grupo residual $\mbarA$ de $\mL$. Por (11), todo polinomio $F$ de $\mL$ pertenece a la clase residual $F\ma$ cuando $\bara$ es la clase residual de $N$ módulo $\mL$ y, por tanto, la clase unidad de $\mbarA$. Si asociamos ahora las clases residuales $F\ma$ de $\mA$ con las clases residuales $F\bara$ de $\mbarA$, a cada clase residual de $\mA$ se asigna una clase residual de $\mbarA$, de tal manera que la suma de dos clases residuales de $\mA$ corresponde a la suma de las clases residuales asignadas de $\mbarA$, y el producto de una clase residual de $\mA$ por un polinomio corresponde al producto de la clase residual asignada de $\mbarA$ por el mismo polinomio. Esta asignación es además biunívoca: de $F\ma=0$ se sigue $F\equiv0(\mL)$ por (11), luego $F\bara=0$; recíprocamente, $F\bara=0$, luego $F\equiv0(\mL)$, dice también por (10) que $FN\equiv0(\mL)$, luego $FN\equiv0(\mM)$, es decir, $F\ma=0$.

Esto nos conduce a un concepto fundamental, que definimos así:

\emph{Una correspondencia biunívoca entre dos sistemas $\mA$ y $\mbarA$ de clases residuales se llama isomorfa si a la suma de dos clases residuales de $\mA$ se asigna la suma de las correspondientes clases residuales de $\mbarA$, y al producto de una clase residual arbitraria de $\mA$ por un polinomio se asigna el producto de la correspondiente clase residual de $\mbarA$ por el mismo polinomio.}

Un subsistema $\mA$ de clases residuales en $\mG$ que es isomorfo al grupo residual de un módulo se llama \emph{subgrupo} de $\mG$. Como muestra la demostración anterior, la asignación aquí es precisamente tal que cada clase residual de $\mA$ surge de la correspondiente clase residual de $\mbarA$ por adjunción de ciertos polinomios, a saber, todos los polinomios que pertenecen a $\mL$ pero no a $\mM$.

Del mismo modo se muestra que la totalidad de todas las clases residuales de la forma $F\mb$ forma un subgrupo $\mB$ de $\mG$ isomorfo al grupo residual $\mbarB$ de $\mN$.

Un grupo residual $\mG$ cuyas clases residuales admiten una descomposición aditiva única (11) se llamará \emph{reducible} en los dos subgrupos $\mA$ y $\mB$, y se escribirá $\mG=\mA+\mB$; un grupo que no es reducible se llama \emph{irreducible}. Ocasionalmente llamaremos también reducibles o irreducibles a los módulos asociados. Por tanto vale lo siguiente.

\textbf{Teorema I.}\footnote{Cf. Schmeidler, loc. cit., p. 39.} \emph{Un módulo $\mM$ es representable como mínimo común múltiplo de dos divisores propios relativamente primos $\mL$ y $\mN$ si y sólo si su grupo residual $\mG$ es reducible en dos subgrupos $\mA$ y $\mB$. En este caso $\mA$ y $\mB$ son respectivamente isomorfos a los grupos residuales de $\mL$ y $\mN$.}

\subsection*{\S~5. Cálculo con unidades y reducibilidad en $k$ subgrupos.}

Las clases residuales $\ma$ y $\mb$ de $\mG$ tienen en común con la unidad $\me$ la propiedad de que el producto de $\ma$ y $\mb$ por una clase residual arbitraria $\mx$ está definido; por tanto llamamos también a $\ma$ y $\mb$ unidades de $\mG$. En efecto, de (11), para $F=M$, se sigue
\[
 0=M\me=M\ma+M\mb,
\]
y por tanto, por unicidad, $M\ma=0$, $M\mb=0$, y entonces
\[
 (F+M)\ma=F\ma=\mx\ma,
\]
donde $\mx$ denota la clase residual de $F$. En lugar de la relación (11) podemos escribir, por consiguiente, la relación equivalente entre clases
\[
 \mx\me=\mx\ma+\mx\mb.
\]
En particular, tomando $\mx=\ma$ se obtiene
\[
 \ma=\ma^2+\ma\mb,
\]
y por unicidad $\ma^2=\ma$, $\ma\mb=0$. De modo análogo se obtiene $\mb^2=\mb$, $\mb\ma=0$.

Correspondientemente a la representación de un grupo residual como suma de dos constituyentes, podemos definir ahora también tal representación como suma de $k$ constituyentes. Supóngase dada una representación única
\[
 F\me=F\ma_1+F\ma_2+\cdots+F\ma_k\qquad(\ma_1\ne0,\ldots,\ma_k\ne0),
\]
donde de $0=G_1\ma_1+\cdots+G_k\ma_k$ se sigue también $G_i\ma_i=0$. Entonces, en particular para $F=M$, se obtiene la relación $M\ma_i=0$, de modo que la representación también puede escribirse en la forma\footnote{Para la descomposición aditiva de los grupos residuales valen las leyes conmutativa y asociativa; a diferencia del caso conmutativo, sin embargo, de $\mU_1+\mU_2=\mU_1+\mU_3$ no se sigue $\mU_2=\mU_3$; cf. ejemplo 1 en \S~12.}
\begin{equation}
 \mx\me=\mx\ma_1+\cdots+\mx\ma_k\qquad(\ma_1\ne0,\ldots,\ma_k\ne0).
\tag{12}
\end{equation}
Esta representación también puede considerarse como surgida de una descomposición sucesiva en dos constituyentes. Si esto puede continuarse sin límite, con todos los $\ma_i\ne0$, hablamos de una suma de infinitos constituyentes.

Al poner $\mx=\ma_i$ en (12) se obtiene
\begin{equation}
 \ma_i^2=\ma_i,
 \qquad \ma_i\ma_j=0\quad(i\ne j),
 \qquad i,j=1,\ldots,k,
\tag{13}
\end{equation}
con $\ma_i\ne0$. Recíprocamente, supóngase presente en $\mG$ un sistema de clases residuales $\ma_1,\ldots,\ma_k$ para el cual se cumplen las condiciones de ortogonalidad (13) y para el cual vale
\begin{equation}
 \me=\ma_1+\cdots+\ma_k.
\tag{14}
\end{equation}
Mostraremos que esto da una representación del grupo como suma de $k$ constituyentes. Primero, de la hipótesis se sigue, para todo polinomio $A_i$ de $\ma_i$, ya que $A_i\ma_i=(A_i+M)\ma_i$, que $M\ma_i=0$; por tanto $\ma_i$ puede multiplicarse por toda clase. Así (14) implica la representación (12), y esta representación es única. Pues si $0=\mx_1\ma_1+\cdots+\mx_k\ma_k$, la multiplicación a la derecha por $\ma_i$ da, por (13), $0=\mx_i\ma_i$, lo que prueba la unicidad y por tanto la afirmación. Las clases residuales $\ma_1,\ldots,\ma_k$ que aparecen en la descomposición de la clase unidad se llamarán las unidades pertenecientes a la descomposición; equivalentemente están caracterizadas por las relaciones de ortogonalidad (13) y la condición (14). Los desarrollos posteriores consistirán esencialmente en calcular con unidades.

Queda por establecer la conexión entre la descomposición aditiva del grupo en $k$ constituyentes y la correspondiente descomposición del módulo. Como en el caso $k=2$, resultará por inducción una representación del módulo como mínimo común múltiplo de $k$ módulos relativamente primos.

Sea dada la descomposición (12), que escribimos en la forma
\begin{align}
 \mx\me&=\mx\mb+\mx\ma_k,\tag{15}\\
 \mx\mb&=\mx\ma_1+\cdots+\mx\ma_{k-1}.\tag{16}
\end{align}
De (15) se sigue por el Teorema I que las clases residuales $\mx\overline{\mb}$ son isomorfas al grupo residual del módulo
\begin{equation}
 \mN=(\mM,A_k).
\tag{17}
\end{equation}
Para sus clases residuales $\mx\overline{\mb}$ es válida por tanto también la representación correspondiente a (16):
\begin{equation}
 \mx\overline{\mb}=\mx\overline{\ma}_1+\cdots+\mx\overline{\ma}_{k-1}.
\tag{18}
\end{equation}
Supóngase ahora ya probado que
\begin{equation}
 \mN=[\mM_1,\ldots,\mM_{k-1}],
\tag{19}
\end{equation}
donde
\begin{equation}
\left\{
\begin{aligned}
 \mM_1&=(\mN,A_2+\cdots+A_{k-1}),\\
 \mM_2&=(\mN,A_1+A_3+\cdots+A_{k-1}),\\
 &\hspace{1.8cm}\vdots\\
 \mM_{k-1}&=(\mN,A_1+\cdots+A_{k-2}).
\end{aligned}\right.
\tag{20}
\end{equation}
Para dos sumandos esta afirmación concuerda con el resultado anterior. Además, los polinomios $A_i$ pueden elegirse aquí concretamente como polinomios de $\overline{\ma}_i$ que al mismo tiempo están contenidos en $\ma_i$, pues por \S~4 los polinomios de esta última clase residual están contenidos en la primera. Entonces, por (15),
\begin{equation}
 \mM=[\mN,\mM_k],
\tag{21}
\end{equation}
donde
\[
 \mM_k=(\mM,A_1+\cdots+A_{k-1}).
\]
De (19) y (21) se obtiene
\[
 \mM=[[\mM_1,
 \ldots,\mM_{k-1}],\mM_k]=[\mM_1,\ldots,\mM_{k-1},\mM_k]
\]
por la definición de mínimo común múltiplo; y de (17) y (20) se obtiene
\[
 \mM_1=(\mM,A_k,A_2+\cdots+A_{k-1}).
\]
Se sigue que
\[
 \mM_1=(\mM,A_2+\cdots+A_k),
\]
pues este último módulo contiene, por (13), también el polinomio
\[
 A_k=A_k(A_2+\cdots+A_k)\pmod{\mM}.
\]
Se obtienen expresiones correspondientes para $\mM_2,\ldots,\mM_{k-1}$. Además, los módulos $\mM_1,\ldots,\mM_k$ son relativamente primos, pues por (14)
\[
 \frac{1}{k-1}\Bigl((A_2+\cdots+A_k)+(A_1+A_3+\cdots+A_k)+\cdots+(A_1+\cdots+A_{k-1})\Bigr)\equiv1\pmod{\mM}.
\]
Forman además un sistema totalmente relativamente primo, es decir, el mínimo común múltiplo de cualesquiera de ellos es relativamente primo al mínimo común múltiplo de los restantes.\footnote{Esto es una generalización de un concepto introducido por Loewy.} Esto se sigue directamente agrupando los sumandos de la fórmula (12) en dos grupos adecuados.

Queda así mostrado que a una descomposición del grupo residual en $k$ sumandos corresponde una representación del módulo como mínimo común múltiplo de $k$ módulos que forman un sistema totalmente relativamente primo. Probamos ahora el teorema recíproco, también por inducción.

Supóngase, pues, que
\[
 \mM=[\mM_1,\ldots,\mM_k]=[[\mM_1,\ldots,\mM_{k-1}],\mM_k]=[\mN,\mM_k],
\]
donde $\mM_1,\ldots,\mM_k$ forman un sistema totalmente relativamente primo. En combinación con el Teorema I esto da la representación (15) para el grupo residual. Los módulos $\mM_1,\ldots,\mM_{k-1}$ también forman un sistema totalmente relativamente primo. Pues si, bajo alguna división de los módulos de este sistema en dos grupos, el mínimo común múltiplo de un grupo y el mínimo común múltiplo del otro tuviesen un divisor común, este divisor permanecería si el módulo $\mM_k$ se añadiera a uno de los dos grupos y se formara entonces el mínimo común múltiplo.\footnote{Continuando este argumento se ve que entre $\mM_1,\ldots,\mM_k$ cualesquiera dos son relativamente primos; la recíproca, que esto implicaría que $\mM_1,\ldots,\mM_k$ forman un sistema totalmente relativamente primo, sólo es válida en general en el caso conmutativo; cf. el contraejemplo 1 en \S~12.} Para $k=2$, ``totalmente relativamente primo'' significa lo mismo que ``relativamente primo''; por tanto nuestra afirmación puede suponerse probada hasta $k-1$. Así, para la clase residual $\mx\overline{\mb}$ vale la representación única (18), luego, por la isomorfía, también (16), y la afirmación se sigue.

\textbf{Teorema II.} \emph{Un módulo $\mM$ es representable como mínimo común múltiplo de $k$ módulos $\mM_1,\ldots,\mM_k$ que forman un sistema totalmente relativamente primo si y sólo si su grupo residual $\mG$ es reducible en $k$ subgrupos $\mU_1,\ldots,\mU_k$. Con notación adecuada, $\mU_i$ es entonces isomorfo al grupo residual de $\mM_i$.}

\subsection*{\S~6. Teoremas de finitud.}

En este párrafo se mostrará que todo grupo residual puede representarse como suma de un número finito de constituyentes irreducibles. Para ello transferimos primero el teorema de Hilbert sobre la base de módulos.

De ahora en adelante usamos las hipótesis más especiales formuladas en \S~2: la multiplicación de $\xi_1,\ldots,\xi_n$ entre sí es conmutativa, y la ley de multiplicación de los polinomios tiene la forma
\begin{equation}
 \xi_i a=a_i\xi_i+b_i,
 \qquad (a_i\ne0)\quad(i=1,\ldots,n).
\tag{6}
\end{equation}
Basta probar la afirmación sólo para módulos. Pues si $\mS$ es cualquier sistema de polinomios, formamos el módulo $\mM$ derivado de él añadiendo a $\mS$ todos los polinomios que pueden componerse a partir de un número finito de polinomios $G_1,\ldots,G_k$ de $\mS$ en la forma $A_1G_1+\cdots+A_kG_k$. Si $F_1,\ldots,F_l$ es entonces una base modular de $\mM$, y si
\[
 F_i=\sum_j A_{ij}G_j\qquad(i=1,\ldots,l),
\]
entonces los finitos polinomios $G_j$ que aparecen en estas representaciones forman una base de $\mS$.

Probamos el teorema de la base de módulos siguiendo a Gordan (\emph{Göttingen Nachrichten} 1899). La prueba se divide en dos partes. La primera concierne a un sistema $\mP$ de infinitos productos de potencias y permanece esencialmente inalterada a causa de la conmutatividad de $\xi_1,\ldots,\xi_n$. Los productos de potencias
\[
 P=\xi_1^{a_1}\cdots\xi_n^{a_n},
\]
cada uno de los cuales se introduce en el sistema una sola vez, se ordenan por grado creciente y, dentro del mismo grado, de alguna manera fija. Consideramos ahora el subsistema $\mQ$ de $\mP$ formado por todos los productos de potencias $P=Q$ que no son divisibles por ningún $P$ precedente. Entonces ningún $Q$ es divisible tampoco por un $Q$ posterior, ya que los posteriores son de grado mayor o, en el mismo grado, distintos de los anteriores. Por otra parte, todo $P$ de $\mP$ es divisible por al menos un $Q$. Pues supóngase, en contra, que $P_m$ es el primer $P$ no divisible por ningún $Q$. Entonces $P_m$ no puede ser él mismo un $Q$, de modo que es divisible por al menos un $P$ precedente; como por hipótesis este $P$ precedente es divisible por algún $Q$, también $P_m$ es divisible por un $Q$, contradicción.

Afirmamos ahora que el número de los $Q$ es finito. En efecto, sea
\[
 Q_0=\xi_1^{h_1}\cdots\xi_n^{h_n}
\]
el primer $Q$. Como un $Q_\lambda=\xi_1^{h_{1\lambda}}\cdots\xi_n^{h_{n\lambda}}$ arbitrario no es divisible por $Q_0$, para todo $Q$ debe valer al menos una desigualdad $h_{\rho\lambda}<h_\rho$; sea $\lambda$ el menor índice para el cual esto ocurre. Agrupamos todos los $Q$ para los cuales $h_{\rho\lambda}$ tiene el mismo valor fijo $c_\lambda$ en una clase $K(c_\lambda)$. Puesto que $c_\lambda<h_\lambda$, el número de estas clases es finito. Pero cada clase contiene sólo finitos elementos. En efecto, si ponemos un $Q$ de la clase $K(c_\lambda)$ igual a $\xi_\lambda^{c_\lambda}Q'$, entonces los $Q'$ forman también un sistema de productos de potencias, ninguno de los cuales divide a otro, y contienen sólo $n-1$ variables. Para $n=2$ los $Q'$ contienen sólo una variable; su número es entonces 1 y, por tanto, finito, de modo que el número puede suponerse finito también para $n-1$ variables. Esto prueba la afirmación.

En segundo lugar, sea $\mM$ un módulo de polinomios cuyos términos están ordenados lexicográficamente. Sea $\mP$ el sistema de los productos de potencias que aparecen en el término principal, cada uno escrito una sola vez. Entonces al sistema $\mQ=(Q_1,\ldots,Q_\rho)$ corresponde un sistema de finitos polinomios $F_1,\ldots,F_\rho$, obtenido asignando a cada $Q_i$ algún polinomio de $\mM$ cuyo término principal sea $b_iQ_i$, de modo que $F_i=b_iQ_i+\cdots$ con términos inferiores.

Sea ahora $F=bP+\cdots$ un polinomio arbitrario de $\mM$, donde $P=\xi_1^{\beta_1}\cdots\xi_n^{\beta_n}Q_i$. Formemos el producto
\[
 \xi_1^{\beta_1}\cdots\xi_n^{\beta_n}F_i
 =\xi_1^{\beta_1}\cdots\xi_n^{\beta_n}(b_iQ_i+\cdots)
 =c\,\xi_1^{\beta_1}\cdots\xi_n^{\beta_n}Q_i+\cdots
 =cP+\cdots.
\]
Aquí $c\ne0$ por (6), y todos los términos siguientes son términos inferiores en el orden lexicográfico, puesto que la multiplicación es conmutativa en cuanto a los exponentes, salvo términos inferiores.\footnote{Esto muestra que en lugar de la ley (6) también se podría usar la ley algo menos restrictiva
\[
 \xi_i a=a_i\xi_i+a_{i,i+1}\xi_{i+1}+\cdots+a_{in}\xi_n+b_i\qquad(a_i\ne0)
\]
para alguna numeración fija de las $\xi$.} La diferencia
\[
 F-\frac bc\,\xi_1^{\beta_1}\cdots\xi_n^{\beta_n}F_i
\]
contiene entonces sólo términos inferiores y pertenece de nuevo a $\mM$. Por tanto el procedimiento termina después de un número finito de pasos, y $F_1,\ldots,F_\rho$ quedan reconocidos como una base modular.

\textbf{Teorema III.} \emph{Todo sistema $\mS$ de infinitos polinomios tiene una base modular finita $G_1,\ldots,G_k$, de modo que todo polinomio $G$ de $\mS$ admite una representación $G=A_1G_1+\cdots+A_kG_k$.}

Lo siguiente es inmediato.

\textbf{Corolario.} \emph{Todo sistema $\mS$ de infinitas clases residuales mod $\mM$ tiene una base $\mx_1,\ldots,\mx_k$, de modo que todo $\mx$ de $\mS$ es representable en la forma $\mx=A_1\mx_1+\cdots+A_k\mx_k$.}

En efecto, se asigna a cada clase residual algún representante y se considera el sistema de estos representantes junto con todos los polinomios de $\mM$.

Pasamos ahora a la demostración del siguiente teorema.

\textbf{Teorema IV.} \emph{Todo grupo residual puede representarse como suma de finitos constituyentes irreducibles; equivalentemente, por el Teorema II, todo módulo puede representarse como mínimo común múltiplo de finitos módulos irreducibles que forman un sistema totalmente relativamente primo.}

Supóngase, en contra, que $\mG$ es una suma de infinitos constituyentes $\mU_1,\mU_2,\ldots$ (cf. \S~5), y sean $\ma_1,\ma_2,\ldots$ las unidades correspondientes, todas no nulas por definición. Formamos el sistema $\mS$ de todas estas unidades; por el corolario al Teorema III tiene una base finita $\ma_{i_1},\ldots,\ma_{i_v}$ con $i_1<\cdots<i_v$. Sea $\mu>i_v$. Entonces
\[
 \ma_\mu=\mx_1\ma_{i_1}+\cdots+\mx_v\ma_{i_v}.
\]
Esto es una relación lineal entre unidades, excluida por la definición. Por tanto, en toda descomposición aditiva de un grupo residual aparece un constituyente irreducible después de finitos pasos; por la conmutatividad de la descomposición podemos ponerlo al comienzo. Así, sin pérdida de generalidad, podemos suponer que los propios $\mU_i$ son irreducibles. Pero por el argumento precedente sólo pueden aparecer finitos constituyentes $\mU_i$. La afirmación queda probada.

\subsection*{\S~7. Un caso de descomposición única.}

Sea
\[
 \mG=\mU_1+\cdots+\mU_k=\mB_1+\cdots+\mB_l
\]
un grupo residual cuyos constituyentes $\mU_1,\ldots,\mU_k$, $\mB_1,\ldots,\mB_l$ son irreducibles. Entonces
\begin{equation}
 \mx\me=\mx\ma_1+\cdots+\mx\ma_k
 =\mx\mb_1+\cdots+\mx\mb_l,
\tag{22}
\end{equation}
donde $\ma_\chi$ es la clase residual de $\mG$ asignada isomórficamente a la clase unidad de $\mU_\chi$, y $\mb_\lambda$ la asignada a la clase unidad de $\mB_\lambda$. Sustituyendo $\mx$ por $\mx\mb_\lambda$, con $\lambda$ fijo en el rango $1,\ldots,l$, se obtiene
\begin{equation}
 \mx\mb_\lambda=\mx\mb_\lambda\ma_1+\cdots+\mx\mb_\lambda\ma_k.
\tag{23}
\end{equation}
Esta representación del grupo $\mB_\lambda$ es única, pero no es una descomposición aditiva de $\mB_\lambda$, porque las clases residuales individuales de la derecha no tienen por qué pertenecer a $\mB_\lambda$. Puesto que $\mb_\lambda$ no se anula, no todos los productos $\mb_\lambda\ma_\chi$ pueden ser cero. Sin pérdida de generalidad, sea
\[
 \mb_\lambda\ma_1\ne0,
 \ldots,
 \mb_\lambda\ma_\rho\ne0,
 \qquad
 \mb_\lambda\ma_{\rho+1}=0,
 \ldots,
 \mb_\lambda\ma_k=0.
\]
Entonces
\[
 \mx\mb_\lambda=\mx\mb_\lambda\ma_1+\cdots+\mx\mb_\lambda\ma_\rho.
\]

Para un $\chi$ fijo en el rango $1,\ldots,\rho$, considérese la totalidad de aquellas clases residuales $\mx\mb_\lambda$ que son no nulas y para las cuales también $\mx\mb_\lambda\ma_\chi\ne0$. Los dos sistemas $\mx\mb_\lambda$ y $\mx\mb_\lambda\ma_\chi$ son isomorfos. Primero, la asignación es biunívoca: de $\mx\mb_\lambda\ma_\chi=0$ se sigue, por hipótesis, $\mx\mb_\lambda=0$, y de esto último, por unicidad de la representación (23), $\mx\mb_\lambda\ma_\chi=0$. Además, las sumas corresponden a sumas, y los productos por polinomios arbitrarios a los productos correspondientes. Manteniendo fijo $\chi$, el mismo argumento puede hacerse para todo $\lambda$ con $\mb_\lambda\ma_\chi\ne0$. Para cada $\ma_\chi$ debe existir al menos un tal $\mb_\lambda$, ya que $\ma_\chi=(\mb_1+\cdots+\mb_l)\ma_\chi\ne0$.

Primero suponemos que para cada $\chi$ hay sólo un $\lambda$, y para cada $\lambda$ sólo un $\chi$, tal que $\mb_\lambda\ma_\chi\ne0$.

Esta correspondencia implica $k=l$; además, la notación puede escogerse de modo que $\mb_\chi\ma_\chi\ne0$ y $\mb_\lambda\ma_\chi=0$ para $\lambda\ne\chi$. Entonces decimos que los productos $\mb_\lambda\ma_\chi$ forman un esquema diagonal. En este caso, por (23), la unidad es $\mb_\chi=\mb_\chi\ma_\chi$, y como
\begin{equation}
 \ma_\chi=\mb_1\ma_\chi+\cdots+\mb_l\ma_\chi=\mb_\chi\ma_\chi=\mb_\chi,
\tag{24}
\end{equation}
se tiene en general $\mU_\chi=\mB_\chi$. Hay así unicidad de la descomposición, y los productos $\ma_\chi\mb_\lambda$ forman un esquema diagonal.

Si, por ejemplo, la ley conmutativa vale específicamente para las clases residuales $\ma_\chi\mb_\lambda$, entonces la hipótesis anterior siempre se cumple. Pues en (23) cada constituyente $\mx\mb_\lambda\ma_\chi=\mx\ma_\chi\mb_\lambda$ está entonces contenido en $\mB_\lambda$, de modo que (23) da una descomposición aditiva del grupo $\mB_\lambda$. Como $\mB_\lambda$ es irreducible, todos menos uno de los $\mb_\lambda\ma_\chi$ son cero. Por (24), correspondientemente, para $\chi$ fijo sólo un $\mb_\lambda\ma_\chi$ es distinto de cero, pues de lo contrario (24), por intercambiabilidad de $\ma_\chi$ y $\mb_\lambda$, daría una descomposición aditiva de $\mU_\chi$, que se supone irreducible.\footnote{Lo mismo vale en el caso bilateral no conmutativo tratado por Schmeidler, loc. cit., donde el producto de toda clase residual de un subgrupo con toda clase residual del otro subgrupo se anula. Entonces, como $\mb_\lambda\ma_\chi=\mb_\lambda\ma_\chi\mb_1+\cdots+\mb_\lambda\ma_\chi\mb_l$ y $\mb_\lambda(\ma_\chi\mb_\mu)=0$ para $\lambda\ne\mu$, se tiene $\mb_\lambda\ma_\chi=\mb_\lambda\ma_\chi\mb_\lambda$, luego esta clase está contenida en $\mB_\lambda$; por irreducibilidad de $\mB_\lambda$, se sigue que sólo para un, y por tanto exactamente un, $\chi$ el producto $\mb_\lambda\ma_\chi\ne0$. Del mismo modo se muestra que para cada $\ma_\chi$ existe uno y sólo un $\mb_\lambda$ con $\ma_\chi\mb_\lambda\ne0$, de donde $k=l$. Así los $\mb_\lambda\ma_\chi$ forman un esquema diagonal, lo que prueba la unicidad establecida allí.}

Las mismas conclusiones valen cuando las condiciones $\mb_\lambda\ma_\lambda\ne0$, $\mb_\lambda\ma_\chi=0$ para $\lambda\ne\chi$, $\chi=1,\ldots,k$, e igualmente, para cada $\mu=1,\ldots,l$, $\mb_\mu\ma_\chi=0$, $\mb_\mu\ma_\mu\ne0$, $\chi=1,\ldots,\sigma$, se cumplen no para todos, sino sólo para parte de los $\mb_\lambda$ ($\lambda=1,\ldots,\sigma$). Se obtiene $\mU_\chi=\mB_\chi$ para $\chi=1,\ldots,\sigma$; y si se pone
\[
 \mU_{\sigma+1}+\cdots+\mU_k=\mU,
 \qquad
 \mB_{\sigma+1}+\cdots+\mB_l=\mB,
\]
entonces, también para las unidades $\ma$ y $\mb$ de $\mU$ y $\mB$, se cumplen las condiciones $\mb\ma_\chi=0$, $\mb_\chi\ma=0$ ($\chi=1,\ldots,\sigma$), y $\mb\ma\ne0$, de donde $\mU=\mB$.

\subsection*{\S~8. Descomposición aditiva de grupos completamente reducibles en grupos primos. El teorema de isomorfía.}

Consideramos ahora otro caso, en el que no vale la unicidad, pero sí la isomorfía de las distintas descomposiciones. Ésta es la generalización del caso considerado por Loewy.

\textbf{Definición.} \emph{Un módulo se llama módulo primo si no tiene ningún divisor aparte de sí mismo y del módulo unidad que contiene todos los polinomios. El grupo residual de un módulo primo se llama grupo primo.}

Todo grupo primo es, a fortiori, irreducible. También hay grupos primos infinitos (cf. la definición en \S~11 y el ejemplo \S~12, 2).

Sea ahora el grupo residual dado $\mG$ representable de dos maneras como suma de finitos grupos primos $\mU_1,\ldots,\mU_k$ y $\mB_1,\ldots,\mB_l$. Un grupo residual que admite al menos una tal representación como suma de grupos primos, y el módulo correspondiente $\mM$, se llamarán completamente reducibles, siguiendo la terminología de Loewy.

Consideramos los productos $\ma_\chi\mb_\lambda$ ($\chi=1,\ldots,k$, $\lambda=1,\ldots,l$), y mostramos que o bien $\ma_\chi\mb_\lambda=0$, o bien $\mx\ma_\chi\mb_\lambda$ recorre todo el grupo $\mB_\lambda$ cuando $\mx$ recorre todas las clases residuales de $\mG$. En efecto, el módulo correspondiente a $\mB_\lambda$ es
\[
 (\mM,\mb_1+\cdots+\mb_{\lambda-1}+\mb_{\lambda+1}+\cdots+\mb_l)=\mM_{\mb_\lambda};
\]
se obtiene de $\mM$ adjuntando todos los polinomios de la clase residual $\mb_1+\cdots+\mb_{\lambda-1}+\mb_{\lambda+1}+\cdots+\mb_l$. Este módulo tiene el divisor
\[
 (\mM,\mb_1+\cdots+\mb_{\lambda-1}+\mb_{\lambda+1}+\cdots+\mb_l,
 \ma_\chi\mb_\lambda).
\]
Pero, como $\mM_{\mb_\lambda}$ es por hipótesis un módulo primo, el divisor es igual a $\mM_{\mb_\lambda}$ o al módulo unidad. En el primer caso $\ma_\chi\mb_\lambda$ pertenece a $\mM_{\mb_\lambda}$; por tanto existe una relación
\[
 \ma_\chi\mb_\lambda=\mx(\mb_1+\cdots+\mb_{\lambda-1}+\mb_{\lambda+1}+\cdots+\mb_l),
\]
que da $\ma_\chi\mb_\lambda=0$. En el segundo caso hay dos clases $\mx_0$ y $\mx_1$ tales que
\[
 \mx_0\ma_\chi\mb_\lambda
 +\mx_1(\mb_1+\cdots+\mb_{\lambda-1}+\mb_{\lambda+1}+\cdots+\mb_l)
 =\me=\mb_1+\cdots+\mb_{\lambda-1}+\mb_{\lambda+1}+\cdots+\mb_l+\mb_\lambda,
\]
y por tanto, por unicidad, $\mx_0\ma_\chi\mb_\lambda=\mb_\lambda$. Así $F\mx_0\ma_\chi\mb_\lambda$, y por tanto también $\mx\ma_\chi\mb_\lambda$, recorre todo el grupo $\mB_\lambda$ cuando $F$ o $\mx$ recorre todos los polinomios o todas las clases residuales de $\mG$. Esto prueba la afirmación.

Si $\ma_\chi\mb_\lambda\ne0$, entonces $\mU_\chi$ también es isomorfo a $\mB_\lambda$. Pues por \S~7 la totalidad de todas las clases $\mx\ma_\chi\ne0$ para las cuales $\mx\ma_\chi\mb_\lambda\ne0$ es isomorfa al sistema de clases $\mx\ma_\chi\mb_\lambda$; y éstas agotan el subgrupo $\mB_\lambda$. Sólo queda mostrar que las clases indicadas $\mx\ma_\chi$ agotan $\mU_\chi$. Para ello considérese la totalidad de las clases residuales $\mx\ma_\chi$ para las cuales $\mx\ma_\chi\mb_\lambda=0$. Este sistema tiene la propiedad de que la suma de dos de sus clases y el producto por un polinomio arbitrario vuelven a pertenecer al sistema. Si se adjuntan todas estas clases residuales al módulo $\mM_{\ma_\chi}$, se obtiene de nuevo un módulo que es divisor de $\mM_{\ma_\chi}$; por tanto es $\mM_{\ma_\chi}$ o el módulo unidad. En el último caso, sin embargo, $\ma_\chi$ misma estaría contenida en él, de modo que $\ma_\chi\mb_\lambda=0$, contra la hipótesis. Luego no hay clases $\mx\ma_\chi\ne0$ para las cuales $\mx\ma_\chi\mb_\lambda=0$, y queda probada la isomorfía de $\mU_\chi$ con $\mB_\lambda$.

Supóngase ahora, para un $\chi$ fijo, que $\ma_\chi\mb_{\lambda_1}\ne0,\ldots,\ma_\chi\mb_{\lambda_i}\ne0$ con $i>1$. Entonces $\mU_\chi$ es isomorfo a los grupos $\mB_{\lambda_1},\ldots,\mB_{\lambda_i}$, de modo que estos grupos son isomorfos entre sí.

Mostramos ahora que al menos dos de los $\mU$ deben ser entonces también isomorfos. Pues si en los productos $\mb_\lambda\ma_\chi$ correspondiera a cada $\mb_\lambda$ un único $\ma_\chi$ con $\mb_\lambda\ma_\chi\ne0$, de modo que $\mb_\lambda=\mb_\lambda\ma_\chi$, entonces $\mB_\lambda=\mU_\chi$, ya que $\mU_\chi$, como grupo primo, no puede contener un subgrupo propio. Así $k=l$ y, con una numeración correcta, $\ma_\chi=\mb_\chi$; pero entonces el esquema $\ma_\chi\mb_\lambda$ sería diagonal, lo que no ocurre. Queda probado el teorema siguiente.

\textbf{Teorema V.} \emph{Si un grupo completamente reducible $\mG$ se representa de dos maneras como suma de grupos primos,
\[
 \mG=\mU_1+\cdots+\mU_k=\mB_1+\cdots+\mB_l,
\]
entonces todo grupo $\mU$ es isomorfo a por lo menos un grupo $\mB$, y recíprocamente. Si todo $\mU$ es isomorfo exactamente a un $\mB$, entonces las descomposiciones son idénticas. En el otro caso, al menos dos grupos $\mU$ son isomorfos entre sí, y del mismo modo al menos dos grupos $\mB$ son isomorfos entre sí.}

\subsection*{\S~9. Conexión entre los conceptos ``isomorfo'' y ``del mismo tipo''.}

En este párrafo mostramos que el concepto ``del mismo tipo'', conocido para expresiones diferenciales de una variable, al generalizarse de manera natural conduce a la isomorfía de los grupos residuales asociados.

\textbf{Definición.}\footnote{Para expresiones diferenciales de una variable, cf. la definición formal de Blumberg, loc. cit., p. 25.} \emph{Dos módulos $\mM$ y $\mN$ se llaman del mismo tipo si existen dos polinomios $P$ y $Q$, con $P$ relativamente primo con $\mM$ y $Q$ relativamente primo con $\mN$, tales que para todo $M\equiv0(\mM)$ y $N\equiv0(\mN)$}
\begin{align}
 MQ&\equiv0(\mN),\tag{25}\\
 NP&\equiv0(\mM)\tag{26}
\end{align}
\emph{se cumplen. La primariedad relativa significa aquí la existencia de dos polinomios $X$ e $Y$ tales que}
\begin{align}
 YQ&\equiv1(\mN),\tag{27}\\
 XP&\equiv1(\mM)\tag{28}
\end{align}
\emph{se cumplen.}

Tenemos ahora el siguiente teorema.

\textbf{Teorema VI.} \emph{Dos módulos son del mismo tipo si y sólo si sus grupos residuales son isomorfos.}

1. Sea $\mA$ el grupo residual de $\mM$, $\mB$ el grupo residual de $\mN$, y supóngase que $\mA$ es isomorfo a $\mB$. Sean $\ma$ y $\mb$ las clases unidad de $\mA$ y $\mB$. Además, sea $Q\mb$ la clase de $\mB$ correspondiente a la unidad $\ma$ de $\mA$, y sea $P\ma$ la clase de $\mA$ correspondiente a la unidad $\mb$ de $\mB$. Escribimos esta asignación como
\begin{align}
 \ma&\sim Q\mb,\tag{29}\\
 P\ma&\sim\mb.\tag{30}
\end{align}
De ello se sigue que
\[
 0=M\ma\sim MQ\mb,
 \qquad\text{es decir,}\qquad MQ\equiv0(\mN),
\]
lo que prueba (25). Del mismo modo,
\[
 NP\ma\sim N\mb=0,
 \qquad\text{es decir,}\qquad NP\equiv0(\mM),
\]
lo que prueba (26). Además, por (29), $P\ma\sim PQ\mb$; como por (30) se había supuesto $P\ma\sim\mb$, la unicidad de la asignación da $PQ\mb=\mb$, luego $PQ\equiv1(\mN)$, y análogamente $QP\equiv1(\mM)$. Así quedan probadas (27) y (28), con $X=Q$, $Y=P$.

2. Recíprocamente, supóngase que $\mM$ y $\mN$ son del mismo tipo, de modo que valen las relaciones (25)--(28). Asignamos a una clase arbitraria $F\ma$ de $\mA$ la clase $FQ\mb$ de $\mB$. Esto asigna a cada clase de $\mA$ una sola clase de $\mB$: pues de $F\ma=0$, luego $F=M$, se sigue por (25) que $FQ\mb=0$. Además, todo el grupo $\mB$ queda agotado por esta asignación, ya que la clase especial $Y\ma$ corresponde a $YQ\mb$, que es igual a $\mb$ por (27). Tenemos así una aplicación unívoca $\Phi$ de $\mA$ sobre $\mB$; aún no se ha mostrado que sea biunívoca. Del mismo modo, las ecuaciones (26) y (28) dan una aplicación unívoca $\Psi$ de $\mB$ sobre $\mA$, que tampoco se ha mostrado todavía biunívoca.

Si una de estas dos aplicaciones es biunívoca, la isomorfía queda probada, puesto que se satisfacen los requisitos relativos a sumas y productos. Pero si ambas aplicaciones $\Phi$ y $\Psi$ no fuesen biunívocas, habría clases $G\ma$ en $\mA$ correspondientes a la clase $GQ\mb=0$ en $\mB$. Las clases $FQ\mb$ correspondientes a todas las demás clases residuales $F\ma$ agotarían entonces ya el grupo $\mB$, y por consiguiente la aplicación $\Psi$ devolvería todo el grupo $\mA$ en las clases $FQP\ma$. Como sabemos, hay también clases no nulas en $\mB$ que bajo $\Psi$ corresponden a la clase cero en $\mA$; denotamos por $G_1\ma$ todas aquellas clases para las cuales $G_1\ma\ne0$ pero $G_1QP\ma=0$. Todos los polinomios restantes $F$ tienen la propiedad de que $FQP\ma$ agota el grupo $\mA$. Aquellos polinomios $F$ para los cuales $FQP\ma=G_1\ma$ se denotarán por $G_2$; tales polinomios existen siempre que existan polinomios $G_1$. Para cada $G_2$ se tiene $G_2QP\ma\ne0$, pero $G_2(QP)^2\ma=0$. Correspondientemente, para todo entero $\nu$ existen clases $G_\nu\ma$ tales que $G_\nu(QP)^{\nu-1}\ma\ne0$ pero $G_\nu(QP)^\nu\ma=0$.

Considérese ahora el sistema $\mS$ formado por todos los polinomios $G_1,G_2,\ldots$. Tiene una base modular finita $G_{\lambda_1},\ldots,G_{\lambda_e}$. Elíjase $\nu$ mayor que el mayor de los índices $\lambda_1,\ldots,\lambda_e$, que denotamos por $\lambda$. Entonces
\[
 G_\nu=A_1G_{\lambda_1}+\cdots+A_eG_{\lambda_e}.
\]
Multiplicando por $(QP)^\lambda$ se obtiene
\[
 G_\nu(QP)^\lambda\equiv0(\mM),
\]
contradicción con nuestra hipótesis. Por tanto se sigue la isomorfía y, por la parte 1, las relaciones más precisas (27), (28) valen con $X=Q$, $Y=P$.

Por el Teorema VI, el Teorema V puede expresarse también de la siguiente manera.

\textbf{Teorema VII.}\footnote{Cf. Loewy, pp. 100--101.} \emph{Si un módulo completamente reducible $\mM$ es representable de dos maneras como mínimo común múltiplo de módulos primos $\mP_1,\ldots,\mP_k$ y $\mD_1,\ldots,\mD_l$, cada uno de los cuales forma un sistema totalmente relativamente primo, entonces todo $\mP$ es del mismo tipo que al menos un $\mD$, y recíprocamente. Si todo $\mP$ es del mismo tipo que exactamente un $\mD$, entonces las descomposiciones son idénticas. En el otro caso, al menos dos de los módulos $\mP$ son del mismo tipo entre sí, y del mismo modo al menos dos de los módulos $\mD$ son del mismo tipo entre sí.}

\setcounter{footnote}{19}
\subsection*{\S~10. Existencia de infinitas descomposiciones de un grupo.}

Si la representación como suma de un grupo completamente reducible no es única, entonces, por el Teorema V, en cada representación existen al menos dos subgrupos mutuamente isomorfos. Consideramos ahora el subgrupo formado a partir de tal sistema de subgrupos isomorfos,
\[
  \mathfrak H=\mathfrak A_1+\cdots+\mathfrak A_\alpha,
\]
donde, por tanto, $\mathfrak A_1,\mathfrak A_2,\ldots,\mathfrak A_\alpha$ son todos isomorfos, y afirmamos que $\mathfrak H$ admite entonces infinitas descomposiciones distintas de este tipo, siempre que el dominio de racionalidad $P$ contenga infinitas ``constantes'', es decir, magnitudes $a$ para las cuales $\xi_i a=a\xi_i$. Este teorema vale incluso sin imponer ninguna otra hipótesis sobre los grupos $\mathfrak A_i$, ni siquiera la irreducibilidad. Así, se tiene:

\textbf{Teorema VIII.} \emph{Si $\mathfrak G$ puede representarse como suma de finitos subgrupos mutuamente isomorfos, entonces esta representación es posible de infinitas maneras distintas, siempre que el dominio de racionalidad $P$ contenga infinitas constantes. Por constantes se entienden aquí aquellas magnitudes $a$ del dominio para las cuales $\xi_i a=a\xi_i$.}

\textbf{Demostración.} Sea, pues,
\[
  \mathfrak G=\mathfrak A_1+\cdots+\mathfrak A_\alpha,
\]
y sea $\mathfrak A_i$ isomorfo a $\mathfrak A_\varkappa$ para $i,\varkappa=1,\ldots,\alpha$ ($\alpha\ge2$). Para fijar una correspondencia isomorfa determinada, distinguimos provisionalmente un subgrupo, digamos $\mathfrak A_1$. Supongamos que la isomorfía entre $\mathfrak A_1$ y $\mathfrak A_i$ está fijada por
\begin{equation}
 a_1\sim t_{1i}a_i,
 \qquad
 a_i\sim t_{i1}a_1,
 \qquad
 t_{11}a_1=a_1
 \quad(i,\varkappa=1,
\ldots,\alpha),
\tag{31}
\end{equation}
de modo que, en el grupo $\mathfrak A_1$, cada clase residual queda asignada a sí misma. Como $Ma_1=0$ y $Ma_i=0$, se sigue que $Mt_{1i}a_i=0$ y $Mt_{i1}a_1=0$; por tanto, las clases $t_{1i}a_i$ y $t_{i1}a_1$, análogamente a las unidades, admiten multiplicación por clases, y no sólo por polinomios. De la definición de isomorfía se sigue entonces que, para dos clases asignadas $\mathfrak y$ y $\bar{\mathfrak y}$ que admiten esta multiplicación por clases, las clases $\mathfrak r\mathfrak y$ y $\mathfrak r\bar{\mathfrak y}$ vuelven a estar asignadas entre sí. Por consiguiente, las relaciones $a_ra_s=0$ $(r\ne s)$ ($r=1,\ldots,\alpha$; $s=1,\ldots,\alpha$) dan, para $s=1$,
\begin{equation}
\left\{
\begin{aligned}
 a_r t_{1\varkappa}a_\varkappa&=0 &&(\varkappa=1,
\ldots,
\alpha,
\ r\ne1),\\
 \text{y para }s>1\qquad a_r t_{s1}a_1&=0 &&(r\ne s),\\
 \text{de ahí también}\qquad t_{1i}a_i&=a_1t_{1i}a_i,
 &t_{i1}a_1&=a_it_{i1}a_1.
\end{aligned}\right.
\tag{32}
\end{equation}
Además, de (31), generalizando (27) y (28), y de la unicidad de la asignación, se obtiene
\begin{equation}
 t_{1i}t_{i1}a_1=a_1,
 \qquad
 t_{\varkappa1}t_{1\varkappa}a_\varkappa=a_\varkappa.
\tag{33}
\end{equation}
Como isomorfía entre $\mathfrak A_i$ y $\mathfrak A_\varkappa$ fijamos ahora la relación obtenida por composición de (31); tal fijación es necesaria, pues los grupos individuales pueden admitir isomorfías en sí mismos. De (31), (32) y (33) resulta por tanto
\begin{equation}
 a_i\sim t_{i1}t_{1\varkappa}a_\varkappa=t_{i\varkappa}a_\varkappa,
 \qquad
 t_{i\varkappa}a_i=a_i,
 \qquad
 a_rt_{s\varkappa}a_\varkappa=0\quad(r\ne s),
\tag{34}
\end{equation}
donde se ha puesto $t_{i1}t_{1\varkappa}=t_{i\varkappa}$. De (33) y (34) se obtiene además
\[
 t_{i\varkappa}t_{\varkappa l}a_l=t_{i1}t_{1\varkappa}t_{\varkappa1}t_{1l}a_l
 =t_{i1}t_{1l}a_l=t_{il}a_l;
\]
con ello queda suprimida de nuevo la distinción del grupo $\mathfrak A_1$. La relación isomorfa entre $\mathfrak A_i$ y $\mathfrak A_\varkappa$ permanece igual cualquiera que sea el grupo $\mathfrak A_\varkappa$ usado en lugar de $\mathfrak A_1$ para la definición. Por tanto, los $t_{i\varkappa}$ son clases que, en conjunto, satisfacen las relaciones
\begin{equation}
 Mt_{i\varkappa}a_\varkappa=0;
 \quad a_rt_{i\varkappa}a_\varkappa=0\ (r\ne i);
 \quad t_{i\varkappa}t_{\varkappa l}a_l=t_{il}a_l;
 \quad t_{\varkappa i}a_i=t_{\varkappa i}t_{i\varkappa}a_\varkappa=a_\varkappa.
\tag{35}
\end{equation}
Estas relaciones dicen exactamente que los módulos $\mathfrak M_i$ y $\mathfrak M_\varkappa$ pertenecientes a $\mathfrak A_i$ y $\mathfrak A_\varkappa$ son del mismo tipo (cf. la definición del \S~9); en lugar de los $P$ y $Q$ de allí aparecen aquí las clases $t_{i\varkappa}$ y $t_{\varkappa i}$, y las dos primeras relaciones (35), tomadas conjuntamente, corresponden a la fórmula (25), respectivamente (26). Recíprocamente, de estas relaciones se sigue por el Teorema VI la isomorfía.\footnote{Y, de hecho, sin aplicar el teorema de finitud, puesto que las fórmulas exhiben ya respectivamente una aplicación $\varphi_{i\varkappa}$ y su inversa $\varphi_{\varkappa i}$. En efecto, de $ra_i=ra_it_{i\varkappa}a_\varkappa t_{\varkappa i}a_i$ se sigue que $ra_i\ne0$ implica también $ra_it_{i\varkappa}a_\varkappa\ne0$.}

A partir de las clases $t_{i\varkappa}a_\varkappa$ componemos ahora clases $b_\rho$ que volverán a resultar unidades de grupos $\mathfrak B_\rho$, de donde se seguirá una representación
\[
 \mathfrak G=\mathfrak B_1+\cdots+\mathfrak B_\alpha.
\]
Sean $\lambda_{i\varkappa}$ constantes cuyo determinante $|\lambda_{i\varkappa}|$ sea distinto de cero, y sean $\lambda^{\varkappa r}$ los cofactores correspondientes divididos por $|\lambda_{i\varkappa}|$, de modo que valgan las relaciones
\begin{equation}
\left\{
\begin{aligned}
 \sum_{r=1}^{\alpha}\lambda_{ri}\lambda^{r\varkappa}&=\delta_{i\varkappa}=\begin{cases}0,&i\ne\varkappa,\\1,&i=\varkappa,
 \end{cases}\\
 \sum_{r=1}^{\alpha}\lambda_{ir}\lambda^{\varkappa r}&=\delta_{i\varkappa}.
\end{aligned}\right.
\tag{36}
\end{equation}
Definimos entonces
\begin{equation}
 b_\rho=\sum_{i,\varkappa}\lambda_{\rho i}\lambda^{\rho\varkappa}t_{i\varkappa}a_\varkappa
 \qquad(\rho=1,
\ldots,\alpha).
\tag{37}
\end{equation}
De (35), (36) y (37) se sigue $Mb_\rho=0$ y
\begin{equation}
 \sum_{\rho=1}^{\alpha}b_\rho
 =\sum_{i,\varkappa,\rho}\lambda_{\rho i}\lambda^{\rho\varkappa}t_{i\varkappa}a_\varkappa
 =\sum_{i,\varkappa}\delta_{i\varkappa}t_{i\varkappa}a_\varkappa
 =\sum_i t_{ii}a_i=\sum_i a_i=e,
\tag{38}
\end{equation}
y también
\[
\begin{aligned}
 b_\rho b_\sigma
 &=\sum_{\substack{i,\varkappa\\ \mu,\nu}}
  \lambda_{\rho i}\lambda^{\rho\varkappa}\lambda_{\sigma\mu}\lambda^{\sigma\nu}
  t_{i\varkappa}a_\varkappa t_{\mu\nu}a_\nu \\
 &=\sum_{i,\mu,\nu}\lambda_{\rho i}\lambda^{\rho\mu}\lambda_{\sigma\mu}\lambda^{\sigma\nu}t_{i\nu}a_\nu
 =\delta_{\rho\sigma}\sum_{i,\nu}\lambda_{\rho i}\lambda^{\sigma\nu}t_{i\nu}a_\nu
 =\begin{cases}0&(\rho\ne\sigma),\\ b_\rho&(\rho=\sigma),\end{cases}
\end{aligned}
\]
$(\rho,\sigma=1,\ldots,\alpha)$. Así, la representación $\mathfrak r e=\mathfrak r b_1+\cdots+\mathfrak r b_\alpha$ es una descomposición aditiva del grupo $\mathfrak G$. Los grupos $\mathfrak B_\rho$ son distintos de los grupos $\mathfrak A_\sigma$ tan pronto como los $\lambda_{i\varkappa}$ no forman simplemente una matriz diagonal. Pues si formamos $a_\sigma b_\rho$, por (35) obtenemos
\[
 a_\sigma b_\rho
 =\sum_{i,\varkappa}\lambda_{\rho i}\lambda^{\rho\varkappa}a_\sigma t_{i\varkappa}a_\varkappa
 =\lambda_{\rho\sigma}
 \sum_\varkappa\lambda^{\rho\varkappa}t_{\sigma\varkappa}a_\varkappa\ne0
 \quad\text{para }\lambda_{\rho\sigma}\ne0,
\]
puesto que, de otro modo, cada sumando individual de la última suma tendría que anularse, cosa que no ocurre. Así, si para un $\sigma$ fijo existen al menos dos $\rho$ tales que $\lambda_{\rho\sigma}\ne0$, entonces $\mathfrak A_\sigma$ no puede ser idéntico a ningún $\mathfrak B_\rho$.

Pero los grupos $\mathfrak B_\rho$ son también mutuamente isomorfos e isomorfos a los grupos $\mathfrak A_\sigma$. Primero mostramos la existencia de clases $r_{\rho\sigma}a_\sigma$ y $r^{\sigma\rho}b_\rho$ que median la isomorfía de $\mathfrak A_\sigma$ y $\mathfrak B_\rho$ ($\rho,\sigma=1,\ldots,\alpha$). En efecto, mostraremos que para las clases
\[
 r_{\rho\sigma}a_\sigma=\sum_i\lambda_{\rho i}t_{i\sigma}a_\sigma,
 \qquad
 r^{\sigma\rho}=\sum_\varkappa\lambda^{\rho\varkappa}t_{\sigma\varkappa}a_\varkappa
 =r^{\sigma\rho}b_\rho
\]
se satisfacen las relaciones ``del mismo tipo'' entre los módulos
\[
 (\mathfrak M,b_1+\cdots+b_{\rho-1}+b_{\rho+1}+\cdots+b_\alpha)
\]
y
\[
 (\mathfrak M,a_1+\cdots+a_{\sigma-1}+a_{\sigma+1}+\cdots+a_\alpha):
\]
\begin{equation}
\left\{
\begin{aligned}
 b_\varkappa r_{\rho\sigma}a_\sigma&=0 &&(\varkappa\ne\rho),&
 r^{\sigma\rho}r_{\rho\sigma}a_\sigma&=a_\sigma,\\
 a_\varkappa r^{\sigma\rho}b_\rho&=0 &&(\varkappa\ne\sigma),&
 r_{\rho\sigma}r^{\sigma\rho}b_\rho&=b_\rho,
\end{aligned}\right.
\tag{39}
\end{equation}
de donde, como en (35), se sigue la isomorfía mediante la asignación
\begin{equation}
 b_\rho\sim r_{\rho\sigma}a_\sigma,
 \qquad
 a_\sigma\sim r^{\sigma\rho}b_\rho.
\tag{40}
\end{equation}
En efecto,
\[
\begin{aligned}
 b_\varkappa r_{\rho\sigma}a_\sigma
 &=\sum_{\mu,\nu,i}\lambda_{\varkappa\mu}\lambda^{\varkappa\nu}t_{\mu\nu}a_\nu\lambda_{\rho i}t_{i\sigma}a_\sigma\\
 &=\sum_{\mu,i}\lambda_{\varkappa\mu}\lambda^{\varkappa i}\lambda_{\rho i}t_{\mu\sigma}a_\sigma
 =\delta_{\varkappa\rho}\sum_\mu\lambda_{\varkappa\mu}t_{\mu\sigma}a_\sigma
 =\delta_{\varkappa\rho}r_{\rho\sigma}a_\sigma=0\quad(\varkappa\ne\rho),
\end{aligned}
\]
y además
\[
 r^{\sigma\rho}r_{\rho\sigma}a_\sigma
 =\sum_{\nu,i}\lambda^{\rho\nu}t_{\sigma\nu}a_\nu\lambda_{\rho i}t_{i\sigma}a_\sigma
 =\sum_i\lambda^{\rho i}\lambda_{\rho i}a_\sigma=a_\sigma.
\]
Recíprocamente, del mismo modo se encuentra
\[
 a_\varkappa r^{\sigma\rho}b_\rho
 =a_\varkappa\sum_\mu\lambda^{\rho\mu}t_{\sigma\mu}a_\mu=0
 \quad(\varkappa\ne\sigma),
\]
y
\[
 r_{\rho\sigma}r^{\sigma\rho}b_\rho
 =\sum_{i,\mu}\lambda_{\rho i}t_{i\sigma}a_\sigma\lambda^{\rho\mu}t_{\sigma\mu}a_\mu
 =\sum_{i,\mu}\lambda_{\rho i}\lambda^{\rho\mu}t_{i\mu}a_\mu=b_\rho.
\]
De las relaciones del mismo tipo acabadas de probar se sigue la isomorfía de $\mathfrak A_\sigma$ y $\mathfrak B_\rho$ (otra vez sin el teorema de finitud). La isomorfía entre $\mathfrak B_\rho$ y $\mathfrak B_\tau$ se obtiene de ésta por composición; registramos las fórmulas. Póngase
\[
 \ell_{\rho\tau}b_\tau=r_{\rho\sigma}r^{\sigma\tau}b_\tau
 =\sum_{i,j}\lambda_{\rho i}t_{i\sigma}a_\sigma\lambda^{\tau j}t_{\sigma j}a_j
 =\sum_{i,j}\lambda_{\rho i}\lambda^{\tau j}t_{ij}a_j;
\]
entonces la asignación es $b_\rho\sim\ell_{\rho\tau}b_\tau$. También aquí las relaciones del mismo tipo pueden verificarse formalmente; si los $\ell_{\rho\tau}$ se ponen en lugar de los $t_{i\varkappa}$, satisfacen (35). En efecto,
\[
 b_\sigma\ell_{\rho\tau}b_\tau=b_\sigma r_{\rho\sigma}r^{\sigma\tau}b_\tau=0\quad(\sigma\ne\rho),
\]
\[
 \ell_{\rho\sigma}\ell_{\sigma\tau}b_\tau=r_{\rho\varkappa}r^{\varkappa\sigma}r_{\sigma\lambda}r^{\lambda\tau}b_\tau
 =r_{\rho\varkappa}a_\varkappa r^{\varkappa\tau}b_\tau=\ell_{\rho\tau}b_\tau.
\]

Con esto queda demostrado el Teorema VIII, pues sólo hay que elegir los $\lambda_{i\varkappa}$ de infinitas maneras de modo que no formen una matriz diagonal y tampoco se transformen unos en otros por matrices diagonales.

Una consecuencia del Teorema VIII, en conexión con el Teorema V o VII, es:

\textbf{Teorema IX.}\footnote{Cf. Loewy, pp. 106--107.} \emph{Sea el módulo $\mathfrak M$ completamente reducible y el mínimo común múltiplo de módulos primos $\mathfrak P_1,\ldots,\mathfrak P_k$ que forman un sistema totalmente relativamente primo, y supóngase que el dominio de racionalidad $P$ contiene infinitas constantes. Entonces la condición necesaria y suficiente para que $\mathfrak M$ sea representable de infinitas maneras distintas como mínimo común múltiplo de módulos primos es que al menos dos de los módulos $\mathfrak P_1,\ldots,\mathfrak P_k$ sean del mismo tipo.}

Ahora derivaremos criterios generales para que un módulo $\mathfrak M$ sea divisible por infinitos módulos primos, y probamos primero lo siguiente:

\textbf{Teorema auxiliar.}\footnote{Cf. Loewy, p. 96.} \emph{Si un módulo completamente reducible $\mathfrak M=[\mathfrak P_1,\ldots,\mathfrak P_k]$ es divisible por un módulo primo $\mathfrak P'$ distinto de todos los $\mathfrak P_i$, entonces al menos dos de los $\mathfrak P_i$ son del mismo tipo.}

Seleccionamos primero de $\mathfrak P_1,\ldots,\mathfrak P_k$ un sistema totalmente relativamente primo, eliminando sucesivamente cada uno de estos módulos primos que esté contenido en el mínimo común múltiplo de los no eliminados aún. Supongamos, pues, ahora que $\mathfrak P_1,\ldots,\mathfrak P_k$ mismos forman un sistema totalmente relativamente primo. Sea $a'$ una clase de $\mathfrak M$ correspondiente a la clase unidad módulo $\mathfrak P'$. Entonces
\[
 Fa'=Fa'a_1+
\cdots+Fa'a_k
\]
para todo polinomio $F$. Al menos dos de los productos $a'a_i$ son no nulos, digamos $a'a_1$ y $a'a_2$; pues si, por ejemplo, sólo $a'a_1\ne0$, entonces $Fa'=Fa'a_1$, y las clases residuales $Fa'$ formarían un subgrupo de $\mathfrak A_1$, por tanto, como $\mathfrak A_1$ es un grupo primo, serían idénticas con $\mathfrak A_1$, de modo que también los módulos $\mathfrak P'$ y $\mathfrak P_1$ coincidirían. Por el razonamiento del \S~8, se sigue que $\mathfrak A'$ es isomorfo a $\mathfrak A_1$ y a $\mathfrak A_2$.

Sea ahora $\mathfrak M$ un módulo arbitrario. Consideramos el mínimo común múltiplo $\mathfrak N$ de todos los módulos primos por los que $\mathfrak M$ es divisible.\footnote{Si hay infinitos, entonces, por el Teorema IV, ninguna descomposición aditiva del grupo residual puede corresponder a esta representación. Cf. también el ejemplo \S~12, 4.} Son posibles dos casos. O bien $\mathfrak N$ no es representable como mínimo común múltiplo de finitos módulos primos; en tal caso $\mathfrak N$, y por tanto también $\mathfrak M$, tiene infinitos módulos primos como divisores. O bien $\mathfrak N$ es representable como mínimo común múltiplo de finitos módulos primos; entonces $\mathfrak N$ se llama el mayor módulo completamente reducible de $\mathfrak M$ (siguiendo una construcción correspondiente de Loewy). Si entre estos finitos módulos primos dos son del mismo tipo, entonces, por el Teorema IX, el módulo $\mathfrak N$, y por tanto $\mathfrak M$, es divisible por infinitos módulos primos. Recíprocamente, si éste es el caso, entonces hay al menos un divisor primo de $\mathfrak N$ distinto de todos los finitos módulos cuyo mínimo común múltiplo representa a $\mathfrak N$. Por el teorema auxiliar, $\mathfrak N$, y por tanto $\mathfrak M$, posee dos divisores primos del mismo tipo. Así hemos probado:

\textbf{Teorema X.} \emph{Si el dominio de racionalidad $P$ contiene infinitas constantes, entonces la condición necesaria y suficiente para que un módulo sea divisible por infinitos módulos primos distintos es que no exista ningún mayor módulo completamente reducible, o bien, si tal módulo existe, que sea divisible por al menos dos módulos primos distintos del mismo tipo.}

\subsection*{\S~11. Relaciones con sistemas de ecuaciones diferenciales y sus integrales.}

En virtud del teorema de finitud, a cada módulo de expresiones diferenciales puede asignársele como base un sistema de finitamente muchas expresiones diferenciales, de modo que la totalidad de los integrales simultáneos de todas las ecuaciones diferenciales obtenidas igualando a cero una expresión diferencial del módulo coincide con la totalidad de los integrales del sistema finito de ecuaciones diferenciales que se obtiene igualando a cero las expresiones diferenciales de la base. Esto da la conexión con la teoría de sistemas de ecuaciones diferenciales parciales lineales; examinamos esta conexión algo más en el caso en que el grupo residual del módulo es ``finito''.

\textbf{Definición.} \emph{Un grupo residual se llama finito si existe un número $\mu$, el orden del grupo, tal que entre cualesquiera $\mu+1$ clases residuales existe una relación lineal con coeficientes de $P$, mientras que existe al menos un sistema de $\mu$ clases residuales entre las cuales no existe ninguna relación; a éste lo llamaremos una base lineal del grupo residual. En el caso contrario, el grupo residual se llama infinito.}

La suma de dos grupos residuales finitos tiene orden finito, a saber, la suma de sus órdenes. La suma de dos grupos infinitos, o de un grupo finito y uno infinito, es un grupo infinito. Ejemplos de grupos finitos los proporcionan todos los grupos residuales de módulos en una variable; pero para varias variables el orden también puede ser finito, por ejemplo para el módulo con base $\xi_1^2$, $\xi_2^3$, donde todas las clases residuales pueden componerse linealmente a partir de las de $1,\xi_1,\xi_2,\xi_1\xi_2$. Para ejemplos de grupos infinitos véase el \S~12.

Tenemos ahora:

\textbf{Teorema XI.} \emph{Si un módulo tiene un grupo residual finito, entonces todo sistema de ecuaciones diferenciales obtenido al igualar a cero una base del módulo admite exactamente tantos integrales linealmente independientes como indica el orden de su grupo residual. Recíprocamente, para un sistema dado de finitamente muchas ecuaciones diferenciales parciales lineales que poseen sólo un número finito de integrales linealmente independientes, este número es igual al orden del grupo residual del módulo generado por las expresiones diferenciales.}

Aquí el dominio de racionalidad $P$ consiste en funciones analíticas de $n$ variables $x_1,\ldots,x_n$, cuyas singularidades, dentro de cierto dominio $n$-dimensional, forman a lo sumo variedades de dimensiones inferiores, de modo que existen puntos arbitrariamente numerosos con vecindades regulares $n$-dimensionales.

Sea $\mathfrak M$ un módulo con grupo residual finito. Mostramos que existe una base lineal de productos de potencias con la propiedad de que \emph{al expresar una clase residual arbitraria mediante esta base, sólo pueden aparecer en el denominador potencias de finitamente muchas magnitudes distintas de $P$}.

Para ello ordenamos lexicográficamente todos los productos de potencias $\xi_1^{a_1}\cdots\xi_n^{a_n}$ y tomamos como representantes $Q_1,\ldots,Q_m$ de las clases residuales todos aquellos que son linealmente independientes, módulo $\mathfrak M$, de los precedentes. Por hipótesis sólo hay finitamente muchos tales productos. Todos los demás pueden expresarse linealmente mediante éstos.

Sea $\mathfrak S$ el sistema de todos los productos de potencias no admitidos en la base lineal. Entonces, por la primera parte de la demostración del Teorema III, $\mathfrak S$ posee un subsistema finito $\mathfrak T$ de productos de potencias tal que todo producto de potencias de $\mathfrak S$ es divisible por al menos uno de $\mathfrak T$. Sean $P_1,\ldots,P_\rho$ los productos de potencias de $\mathfrak T$, y sean $M_1=0,\ldots,M_\rho=0\;(\mathfrak M)$ las relaciones mediante las cuales $P_1,\ldots,P_\rho$ se expresan cada uno, módulo $\mathfrak M$, como combinación lineal de productos de potencias inferiores; por ejemplo,
\[
 M_i=b_iP_i+\text{términos inferiores}.
\]
Si $P$ es un producto de potencias arbitrario de $\mathfrak S$, por tanto de la forma
\[
 \xi_1^{h_1}\cdots\xi_n^{h_n}P_i,
\]
entonces
\[
\begin{aligned}
 \xi_1^{h_1}\cdots\xi_n^{h_n}M_i
 &=b_i\xi_1^{h_1}\cdots\xi_n^{h_n}P_i+\text{términos inferiores}\\
 &=b_iP+\text{términos inferiores}.
\end{aligned}
\]
Supóngase ahora que ya se ha probado para todos los productos de potencias inferiores a $P$ que sólo potencias de $b_1,\ldots,b_\rho$ aparecen en los denominadores, mientras que los numeradores se forman a partir de los coeficientes de los $M_i$ por adición, multiplicación y diferenciación. Entonces lo mismo vale para $P$ mismo. Esta suposición está permitida, pues se cumple para $P_1$, el primer elemento de $\mathfrak S$ y de $\mathfrak T$.

Por el Teorema III, además, $M_1,\ldots,M_\rho$ es una base modular de $\mathfrak M$. Consideremos ahora un sistema de valores $\beta_1,\ldots,\beta_n$ de $x_1,\ldots,x_n$ para el cual $b_1\ne0,\ldots,b_\rho\ne0$ y los restantes coeficientes de los $M_i$ son regulares, es decir, un punto regular del sistema de ecuaciones diferenciales $M_1=0,\ldots,M_\rho=0$. A los productos de potencias $Q_i=\xi_1^{h_{i1}}\cdots\xi_n^{h_{in}}$ $(i=1,\ldots,m)$ les corresponden las derivadas parciales
\[
 \frac{\partial^{h_{i1}+\cdots+h_{in}}}{\partial x_1^{h_{i1}}\cdots\partial x_n^{h_{in}}}y,
\]
para las cuales prescribimos valores constantes arbitrarios $(\gamma_1,\ldots,\gamma_m)$ en el punto $(\beta_1,\ldots,\beta_n)$. Entonces las ecuaciones diferenciales $M_i=0$ determinan unívocamente, como valores finitos en este punto, los valores de todas las derivadas superiores de $y$, pues sólo los $b_i$ aparecen en los denominadores; y es sabido que esto implica la existencia de $m$ integrales analíticos linealmente independientes.

Por otra parte, bajo nuestras hipótesis no pueden existir más de $m$ integrales linealmente independientes, pues de otro modo en un punto regular podrían prescribirse arbitrariamente los valores de $m+1$ derivadas convenientemente elegidas, contra el hecho de que entre cualesquiera $m+1$ productos de potencias deben existir dependencias lineales módulo $\mathfrak M$.

Recíprocamente, si un módulo tiene un grupo residual infinito, el mismo método muestra que existen infinitos integrales linealmente independientes; por tanto, en el caso de sólo $m$ integrales linealmente independientes, el grupo residual es finito y su orden es, por lo anterior, igual a $m$.

Para concluir este párrafo mostramos que, para dos módulos de expresiones diferenciales, el concepto de isomorfía, o equivalentemente de ser del mismo tipo, tiene para los integrales de los sistemas de ecuaciones diferenciales asociados el mismo significado que el concepto de tipo de Poincaré para una variable (cf. Blumberg, loc. cit., Apéndice I).

\textbf{Teorema XII.} \emph{Si dos módulos $\mathfrak M$ y $\mathfrak N$ de expresiones diferenciales son del mismo tipo, entonces existen dos expresiones diferenciales $P$ y $Q$ tales que $P(y)$ y $Q(z)$ recorren todos los integrales de $\mathfrak N$ y de $\mathfrak M$, respectivamente, cuando $y$ recorre todos los integrales de $\mathfrak M$ y $z$ todos los integrales de $\mathfrak N$.}

Por hipótesis, a saber (por el \S~9), para dos polinomios $P$ y $Q$ convenientemente elegidos valen las relaciones
\[
 MQ=0(\mathfrak N),
 \qquad
 PQ=1(\mathfrak N),
\]
\[
 NP=0(\mathfrak M),
 \qquad
 QP=1(\mathfrak M).
\]
Así, puesto que $NP(y)=0$, la función $P(y)$ es un integral de $\mathfrak N$. Del mismo modo, $Q(z)$ es un integral de $\mathfrak M$. Como $PQ(z)=z$ y $QP(y)=y$, $P(y)$ recorre todos los integrales de $\mathfrak N$, y $Q(z)$ todos los de $\mathfrak M$. De este modo, a cada $y$ no nulo le corresponde un $z$ no nulo, y recíprocamente.

\subsection*{\S~12. Ejemplos.}

1. Como primer ejemplo consideramos una ecuación diferencial lineal ordinaria de segundo orden cuyos coeficientes se representan como funciones simétricas de los integrales. El dominio de racionalidad $P$ consiste aquí en todas las funciones racionales de los integrales y de sus derivadas; nos restringimos a una vecindad suficientemente pequeña de un punto regular de la ecuación diferencial. Definimos el módulo $\mathfrak M$ como la totalidad de todas las expresiones diferenciales que tienen a $y_1$ y $y_2$ como integrales, y para las cuales escribimos de nuevo polinomios en una indeterminada $\xi$. Entonces la base modular consta de un solo miembro y está dada por el polinomio
\[
 M=
 \begin{vmatrix} y_1&y_2\\ y_1'&y_2'\end{vmatrix}\xi^2
 -\begin{vmatrix} y_1&y_2\\ y_1''&y_2''\end{vmatrix}\xi
 +\begin{vmatrix} y_1'&y_2'\\ y_1''&y_2''\end{vmatrix},
\]
que queda determinado unívocamente hasta un factor de $P$. Ahora
\begin{equation}
 \mathfrak M=[\mathfrak M_1,
\mathfrak M_2],
\tag{41}
\end{equation}
donde $\mathfrak M_i$ consiste en todas las expresiones diferenciales que tienen el integral $y_i$, y cuya base está determinada, hasta un factor de $P$, por
\[
 M_i=y_i\xi-y_i'
 \qquad(i=1,2).
\]
En efecto, $\mathfrak M$ es divisible por $\mathfrak M_1$ y $\mathfrak M_2$, pues se anula para los integrales $y_1$ y $y_2$, y por tanto es divisible por $[\mathfrak M_1,\mathfrak M_2]$; y como el orden del polinomio de base de $[\mathfrak M_1,\mathfrak M_2]$ debe ser al menos 2, se tiene $\mathfrak M=[\mathfrak M_1,\mathfrak M_2]$. Además, $\mathfrak M_1$ y $\mathfrak M_2$ son relativamente primos, porque
\begin{equation}
 \frac{y_2}{D}(y_1\xi-y_1')-\frac{y_1}{D}(y_2\xi-y_2')=1,
 \qquad
 D=D(y)=\begin{vmatrix}y_1&y_2\\ y_1'&y_2'\end{vmatrix}.
\tag{42}
\end{equation}
También son módulos primos, ya que su grupo residual tiene orden 1.

Pero, además de $y_1$ y $y_2$, $\mathfrak M$ tiene también todos los integrales
\[
 z_1=\alpha_{11}y_1+\alpha_{12}y_2,
 \qquad
 z_2=\alpha_{21}y_1+\alpha_{22}y_2,
 \qquad
 A=|\alpha_{ik}|\ne0,
\]
y por tanto es idéntico al mínimo común múltiplo de los dos módulos relativamente primos $\mathfrak N_1$ y $\mathfrak N_2$, que tienen respectivamente los integrales $z_1$ y $z_2$, con polinomios de base
\[
 N_i=z_i\xi-z_i'=(\alpha_{i1}y_1+\alpha_{i2}y_2)\xi-(\alpha_{i1}y_1'+\alpha_{i2}y_2').
\]
Así $\mathfrak M$ admite infinitas representaciones distintas como mínimo común múltiplo de módulos primos relativamente primos; por el Teorema VII, los dos módulos $\mathfrak N_1$ y $\mathfrak N_2$ son por tanto del mismo tipo entre sí y del mismo tipo que todos los $\mathfrak M_1$ y $\mathfrak M_2$. En efecto, los grupos residuales de todos estos módulos tienen orden 1 y, por tanto, son isomorfos.

Ahora escribiremos la descomposición aditiva del grupo residual $\mathfrak G=\mathfrak A_1+\mathfrak A_2$ correspondiente a la representación (41). Para ello partimos, como en el \S~4, de la relación (42), que expresa la primariedad relativa de $\mathfrak M_1$ y $\mathfrak M_2$. Para las unidades $a_1$ y $a_2$ obtenemos las clases residuales generadas por los polinomios
\[
 -\frac{y_1}{D}(y_2\xi-y_2')
 \quad\text{y}\quad
 \frac{y_2}{D}(y_1\xi-y_1').
\]
Así
\[
 \mathfrak M_1=\left(\mathfrak M,\frac{y_2}{D}(y_1\xi-y_1')\right),
 \qquad
 \mathfrak M_2=\left(\mathfrak M,-\frac{y_1}{D}(y_2\xi-y_2')\right).
\]
Si ahora formamos del mismo modo las unidades $b_1$ y $b_2$ correspondientes a la descomposición $\mathfrak M=[\mathfrak N_1,\mathfrak N_2]$, sustituyendo $y_i$ por $z_i$, obtenemos
\[
 -\frac{z_1}{D(z)}(z_2\xi-z_2')
 =-\frac{\alpha_{11}y_1+\alpha_{12}y_2}{AD(y)}
 ((\alpha_{21}y_1+\alpha_{22}y_2)\xi-(\alpha_{21}y_1'+\alpha_{22}y_2')),
\]
\[
 \frac{z_2}{D(z)}(z_1\xi-z_1')
 =\frac{\alpha_{21}y_1+\alpha_{22}y_2}{AD(y)}
 ((\alpha_{11}y_1+\alpha_{12}y_2)\xi-(\alpha_{11}y_1'+\alpha_{12}y_2')),
\]
y por tanto
\begin{equation}
\left\{
\begin{aligned}
 b_1&=\frac{\alpha_{11}y_1+\alpha_{12}y_2}{y_1}
 \left(\frac{\alpha_{22}}{A}\right)a_1
 +\frac{\alpha_{11}y_1+\alpha_{12}y_2}{y_2}
 \left(-\frac{\alpha_{21}}{A}\right)a_2
 =\frac{z_1}{y_1}\alpha^{11}a_1+\frac{z_1}{y_2}\alpha^{12}a_2,\\
 b_2&=\frac{\alpha_{21}y_1+\alpha_{22}y_2}{y_1}
 \left(-\frac{\alpha_{12}}{A}\right)a_1
 +\frac{\alpha_{21}y_1+\alpha_{22}y_2}{y_2}
 \left(\frac{\alpha_{11}}{A}\right)a_2
 =\frac{z_2}{y_1}\alpha^{21}a_1+\frac{z_2}{y_2}\alpha^{22}a_2.
\end{aligned}\right.
\tag{43}
\end{equation}
Aquí $(\alpha^{ik})$ es la matriz inversa de $(\alpha_{ik})$.

Recíprocamente, intercambiando $z_i$ con $y_i$ se obtiene
\begin{equation}
\left\{
\begin{aligned}
 a_1&=\frac{y_1}{z_1}\alpha_{11}b_1+\frac{y_1}{z_2}\alpha_{12}b_2,\\
 a_2&=\frac{y_2}{z_1}\alpha_{21}b_1+\frac{y_2}{z_2}\alpha_{22}b_2.
\end{aligned}\right.
\tag{44}
\end{equation}
Estas fórmulas exhiben la isomorfía de los grupos $\mathfrak A$ y $\mathfrak B$. A saber, por nuestras consideraciones generales, si $a_1b_\sigma\ne0$, la asignación correspondiente es $a_1\sim a_1b_\sigma$, y por tanto, por (44),
\[
 a_1\sim\frac{y_1}{z_\sigma}\alpha_{1\sigma}b_\sigma,
 \qquad\text{si }\alpha_{1\sigma}\ne0,
\]
y por (43) asimismo
\[
 b_\sigma\sim\frac{z_\sigma}{y_2}\alpha^{\sigma2}a_2,
 \qquad\text{si }\alpha^{\sigma2}\ne0,
\]
de modo que
\[
 \frac{y_1}{z_\sigma}\alpha_{1\sigma}b_\sigma
 \sim
 \frac{y_1}{z_\sigma}\alpha_{1\sigma}\frac{z_\sigma}{y_2}\alpha^{\sigma2}a_2
 =\frac{y_1}{y_2}\alpha_{1\sigma}\alpha^{\sigma2}a_2.
\]
Por tanto
\[
 a_1\sim\alpha_{1\sigma}\alpha^{\sigma2}\frac{y_1}{y_2}a_2
 \quad\text{para }\alpha_{1\sigma}\alpha^{\sigma2}\ne0,
 \qquad
 \frac{1}{\alpha_{1\sigma}\alpha^{\sigma2}}\frac{y_2}{y_1}a_1\sim a_2,
\]
lo que da la inversa. Si $(\alpha_{ik})$ no es una matriz diagonal, esta condición $\alpha_{1\sigma}\alpha^{\sigma2}\ne0$ se cumple siempre para algún $\sigma$. Debe excluirse el caso de una matriz diagonal, pues entonces los módulos son idénticos. Si, por ejemplo, sólo $\alpha_{12}=0$, entonces $\mathfrak A_1=\mathfrak B_1$, pero no $\mathfrak A_2=\mathfrak B_2$. Así, de $\mathfrak A_1+\mathfrak A_2=\mathfrak B_1+\mathfrak B_2$ y $\mathfrak A_1=\mathfrak B_1$ no se sigue $\mathfrak A_2=\mathfrak B_2$ (cf. nota 12). Los tres módulos $\mathfrak M_1$, $\mathfrak M_2$ y $\mathfrak N_2$ son entonces relativamente primos por pares, pero no forman un sistema totalmente relativamente primo, puesto que $[\mathfrak M_1,\mathfrak M_2]$ es divisible por $\mathfrak N_2$ (cf. nota 15).

Poniendo además
\[
 t_{12}=\alpha_{1\sigma}\alpha^{\sigma2}\frac{y_1}{y_2}a_2,
 \qquad
 t_{21}=\frac{1}{\alpha_{1\sigma}\alpha^{\sigma2}}\frac{y_2}{y_1}a_1,
\]
se verifica por cálculo directo que las relaciones (35) se satisfacen.

2. Todo módulo en más de una variable cuya base consta de un solo polinomio proporciona un ejemplo de grupo infinito. El siguiente ejemplo, algo más complicado, quiere mostrar que también pueden aparecer grupos infinitos en módulos con varios miembros, es decir, en módulos cuya base contiene más de un polinomio.\footnote{Cf. el ejemplo de Landau comunicado por Blumberg (loc. cit., pp. 51--52). Aquí y abajo designamos las variables independientes por $x$ e $y$.}
\[
 \mathfrak M=[(\xi+x\eta),(\xi+1)]=[(P),(Q)]
 \qquad\left(\xi=\frac{\partial}{\partial x},\ \eta=\frac{\partial}{\partial y}\right).
\]
La regla de multiplicación es la de las expresiones diferenciales. El grupo residual es aquí infinito, porque los grupos residuales de $(\xi+x\eta)$ y $(\xi+1)$ son ambos infinitos, mientras que los módulos $(\xi+x\eta)$ y $(\xi+1)$ son relativamente primos, pues
\[
 1=(-x\xi-x^2\eta+1)Q+(x\xi+x-1)P.
\]
Ahora es muy fácil mostrar que $\mathfrak M$ no posee ningún polinomio de segundo grado: se pone
\[
 (u_1\xi+u_2\eta+u_3)(\xi+x\eta)=(v_1\xi+v_2\eta+v_3)(\xi+1),
\]
lo que da $u_1=u_2=u_3=v_1=v_2=v_3=0$. Si se buscan a continuación polinomios de tercer grado en $\mathfrak M$, se encuentran dos linealmente independientes, por ejemplo
\[
 (\xi^2+x\xi\eta+\xi+(2+x)\eta)Q=(\xi^2+2\xi+1)P
\]
y
\[
 (\xi^2+2x\xi\eta+x^2\eta^2+\eta)Q
 =(\xi^2+x\xi\eta+\xi+(x-1)\eta)P.
\]
Si $\mathfrak M$ fuese un módulo de un solo miembro, ambos tendrían que ser divisibles por el polinomio de base y, como $\mathfrak M$ no contiene ningún polinomio de segundo grado, por un polinomio de tercer grado; por tanto tendrían que ser proporcionales, cosa que no son. El carácter de varios miembros de $\mathfrak M$ es la razón interna del fracaso de los teoremas del producto de Landau, que son válidos en el caso de una variable.

3. Otro ejemplo muestra que también pueden aparecer grupos infinitos en módulos primos.\footnote{Nuestra definición de ``módulo primo'' difiere de lo que se llama módulo primo en el caso conmutativo, pues allí existen siempre divisores de dimensión inferior, salvo cuando la dimensión es $0$, es decir, cuando el grupo residual es finito.} Sea $P$ el dominio de todas las funciones racionales de dos variables $x,y$. El módulo $\mathfrak M$ con base
\[
 \xi-\frac{y}{x}=\xi-a
\]
tiene un grupo infinito, puesto que este grupo contiene la clase residual de toda potencia de $y$; por otra parte, $\mathfrak M$ es un módulo primo. En efecto, mostramos que todo divisor
\[
 \mathfrak M_1=(\xi-a,
F(\xi,\eta))
\]
es igual al módulo unidad. Todo polinomio $F(\xi,\eta)$ puede representarse módulo $(\xi-a)$ por un polinomio $G(\eta)$:
\[
 F(\xi,
\eta)=c_0\eta^\nu+c_1\eta^{\nu-1}+\cdots+c_\nu\quad(\mathfrak M).
\]
Sin pérdida de generalidad sea $c_0=1$. Por la ley de multiplicación, $\xi\eta^\nu-\eta^\nu\xi=0$; por tanto, como
\[
 \xi\equiv a(\mathfrak M_1),
 \qquad
 \eta^\nu\equiv F_1(\mathfrak M_1),
\]
donde $F_1=-(c_1\eta^{\nu-1}+\cdots+c_\nu)$, se tiene
\[
 \xi F_1-\eta^\nu a\equiv0(\mathfrak M_1).
\]
Desarrollando
\[
 \eta^\nu a=a\eta^\nu+\nu\frac{\partial a}{\partial y}\eta^{\nu-1}+\cdots,
\]
y reemplazando de nuevo $\eta^\nu$ por $F_1$, se obtiene
\[
 \xi(c_1\eta^{\nu-1}+\cdots+c_\nu)-a(c_1\eta^{\nu-1}+\cdots+c_\nu)
 +\nu\frac{\partial a}{\partial y}\eta^{\nu-1}+\cdots\equiv0(\mathfrak M_1),
\]
o bien
\[
 c_1\eta^{\nu-1}a+\frac{\partial c_1}{\partial x}\eta^{\nu-1}+\cdots
 -ac_1\eta^{\nu-1}+\cdots+
\nu\frac{\partial a}{\partial y}\eta^{\nu-1}+\cdots
 \equiv0(\mathfrak M_1),
\]
y finalmente
\[
 \left(\frac{\partial c_1}{\partial x}+\nu\frac{\partial a}{\partial y}\right)\eta^{\nu-1}+\cdots
 \equiv0(\mathfrak M_1).
\]
En el lado izquierdo aparece ahora un polinomio en $\eta$ de grado exacto $\nu-1$ que no se anula idénticamente, pues el coeficiente principal
\[
 \frac{\partial c_1}{\partial x}+\nu\frac{\partial a}{\partial y}
 =\frac{\partial c_1}{\partial x}+\nu\frac{1}{x}
 \quad\text{es }\ne0,
\]
ya que, de lo contrario, $c_1=-\nu\log x+\varphi(y)$ no sería racional. Así el procedimiento puede continuarse y conduce a un polinomio no idénticamente nulo de grado cero, lo que prueba que la unidad pertenece también a $\mathfrak M_1$.\footnote{Este procedimiento equivale a probar que las condiciones de integrabilidad necesarias para un sistema de ecuaciones diferenciales no se cumplen.}

4. Por otra parte, si se adjunta la función $\log x$ al dominio de racionalidad, entonces $\mathfrak M$ tiene los divisores
\[
 \left(\xi-\frac{y}{x},\ \eta-\log x+\varphi(y)\right),
\]
donde $\varphi(y)$ recorre todas las funciones racionales de $y$. Estos divisores son todos distintos, la condición de integrabilidad se cumple para cada uno, y el integral correspondiente es
\[
 y\log x-\int\varphi(y)\,dy+c.
\]
Por tanto, la unidad no puede estar contenida en el módulo, pues en tal caso la función cero sería el único integral posible. Por otra parte, $\xi$ y $\eta$ son congruentes, con respecto al módulo, a una magnitud del dominio; luego el grupo residual es finito y de orden 1. Además, dos cualesquiera de los infinitos divisores son relativamente primos, y el módulo dado $\mathfrak M$ es igual al mínimo común múltiplo $\mathfrak N$ de todos estos divisores. En efecto, es divisible por todos los divisores, y por tanto también por $\mathfrak N$. Si $\mathfrak N$ fuese un divisor propio de $\mathfrak M$, entonces $\mathfrak N$ tendría que contener todavía al menos un polinomio $F$, digamos de grado $\rho$, que dependiese sólo de $\eta$. Haciendo que $\varphi(y)$ recorra las funciones $y,\ldots,y^{\rho+1}$, no hay ninguna relación lineal con coeficientes independientes de $y$ entre los integrales correspondientes
\[
 y\log x-\frac{y^{i+1}}{i+1}+c_i
 \qquad(i=1,
\ldots,
\rho+1).
\]
Pero la ecuación diferencial de orden $\rho$, $F=0$, no puede poseer $\rho+1$ integrales linealmente independientes.

Sin embargo, el módulo $\mathfrak M=\mathfrak N$ no es completamente reducible.\footnote{Así tenemos un ejemplo del primer caso del Teorema X.} Pues, de lo contrario, sería el mínimo común múltiplo de finitamente muchos de estos módulos primos, y por tanto su grupo residual sería de orden finito; pero el grupo residual de $\xi-y/x$ contiene las clases residuales de todas las potencias de $\eta$, y por ello es infinito.

Gotinga, 1 de agosto de 1919.

\begin{center}
(Recibido el 4 de agosto de 1919.)
\end{center}

\clearpage
\setcounter{footnote}{0}
\begin{center}
{\Large\bfseries 18. Sobre un trabajo de K. Hentzelt, caído en la guerra, sobre la teoría de eliminación}\par
\vspace{1em}
\emph{Jahresbericht der DMV} 30 (1921), p. 101
\end{center}

\begin{center}
\textbf{11.ª sesión, viernes 23 de septiembre de 1921, 11 h}\\
(Presidencia: Loewy.)
\end{center}

\textbf{1. E. Noether: Sobre un trabajo de K. Hentzelt, caído en la guerra, sobre la teoría de eliminación.}

Kurt Hentzelt está desaparecido ante Dixmuiden desde octubre de 1914 y debe contarse entre los muertos. Tengo intención de publicar próximamente, en una redacción libre en los \emph{Mathematische Annalen}, las partes esenciales de su disertación de Erlangen, que dejó construida sin lagunas pero expuesta de modo casi incomprensible.

El asunto concierne al problema general de la \textbf{teoría de eliminación}: determinar los ceros comunes de todos los polinomios de un ideal $\mathfrak M$ formado por polinomios. Si primero se somete a las variables a una transformación lineal con coeficientes indeterminados, entonces el resultado principal puede formularse así:

Al ideal $\mathfrak M$ se le puede asignar unívocamente una forma resultante
\[
 R_{\mathfrak M}=R^{(1)}(x_1,
\ldots,x_n)\cdots R^{(i)}(x_i,
\ldots,x_n)\cdots R^{(m)}(x_n)=0(\mathfrak M),
\]
tal que $R_{\mathfrak M}$ se anula en todos los ceros de $\mathfrak M$ y sólo en éstos. Si $\mathfrak N$ es un divisor de $\mathfrak M$ y $R_{\mathfrak N}$ y $R_{\mathfrak M}$ coinciden, entonces también coinciden los propios $\mathfrak N$ y $\mathfrak M$.

La segunda parte del teorema muestra que la forma resultante proporciona los ceros con multiplicidad característica. Se apoya en el hecho de que $\mathfrak M$ puede considerarse como un módulo de formas lineales en los productos de potencias de $x_1,\ldots,x_{i-1}$; los factores individuales $R^{(i)}(x_i,\ldots,x_n)$ de $R_{\mathfrak M}$ se convierten entonces en las normas del respectivo ``módulo básico'' respecto de este módulo, cuando $x_{i+1},\ldots,x_n$ se adjunta al dominio de coeficientes. Si se recurre a la teoría abstracta de ideales, se obtiene la siguiente interpretación: si
\[
 \mathfrak M=[\mathfrak Q,\mathfrak Q_1,
\ldots,
\mathfrak Q_r]=[\mathfrak Q,\mathfrak L]
\]
es una descomposición de $\mathfrak M$ en ideales primarios $\mathfrak Q_i$, y $\mathfrak L$ es el complemento de $\mathfrak Q$, entonces al ideal $\mathfrak Q$ le corresponde un factor primario $Q^{(i)}(x_i,\ldots,x_n)$ de $R^{(i)}(x_i,\ldots,x_n)$, que en el sentido especificado arriba representa la norma de $\mathfrak L$ respecto de $\mathfrak Q$.



\clearpage

\section*{19. Teoría de ideales en dominios de anillos}
\begin{center}
Math. Ann. 83 (1921), pp. 24--66
\end{center}

\begin{center}\textbf{Índice.}\end{center}
\noindent Introducción.\\
\S~1. Dominio de anillos, ideal, condición de finitud.\\
\S~2. Representación de un ideal como mínimo común múltiplo de un número finito de ideales irreducibles.\\
\S~3. Igualdad del número de componentes en dos descomposiciones distintas en ideales irreducibles.\\
\S~4. Ideales primarios. Unicidad de los ideales primos asociados en dos descomposiciones distintas en ideales irreducibles.\\
\S~5. Representación de un ideal como mínimo común múltiplo de ideales primarios máximos. Unicidad de los ideales primos asociados.\\
\S~6. Representación única de un ideal como mínimo común múltiplo de ideales relativamente-primos irreducibles.\\
\S~7. Unicidad de los ideales aislados.\\
\S~8. Representación única de un ideal como producto de ideales coprimos-irreducibles.\\
\S~9. Extensión de la investigación a módulos. Igualdad del número de componentes en descomposiciones en módulos irreducibles.\\
\S~10. Caso especial del dominio de polinomios.\\
\S~11. Ejemplos de la teoría de números y de la teoría de expresiones diferenciales.\\
\S~12. Ejemplo de la teoría de los divisores elementales.

\subsection*{Introducción.}
El contenido del presente trabajo es la \emph{transferencia de los teoremas de descomposición de los enteros racionales, respectivamente de los ideales en cuerpos algebraicos de números, a ideales en dominios íntegros arbitrarios y, más generalmente, en dominios de anillos}. Para comprender esta transferencia, indiquemos primero los teoremas de descomposición de los enteros racionales en una forma algo distinta de la formulación usual.

Si en
\[
 a=p_1^{\varrho_1}p_2^{\varrho_2}\cdots p_\sigma^{\varrho_\sigma}=q_1q_2\cdots q_\sigma
\]
se consideran las potencias de primos $q_i$ como los componentes de la descomposición, entonces estos componentes poseen las siguientes propiedades características:

1. Son \emph{coprimos dos a dos}; pero ningún $q_i$ puede representarse como producto de números coprimos dos a dos, de modo que en este sentido hay irreducibilidad. De la coprimalidad dos a dos se deduce además que el producto $q_1\cdots q_\sigma$ es igual al mínimo común múltiplo $[q_1\cdots q_\sigma]$.

2. Dos componentes cualesquiera, $q_i$ y $q_k$, son \emph{relativamente primos}; es decir, si $bq_i$ es divisible por $q_k$, entonces $b$ es divisible por $q_k$. También en este sentido hay irreducibilidad.

3. Todo $q$ es \emph{primario}; esto es, si un producto $b\cdot c$ es divisible por $q$, pero $b$ no es divisible, entonces una potencia\footnote{Si esta potencia es siempre la primera, se trata, como es sabido, de números primos.} de $c$ es divisible. La representación es además una representación por \emph{componentes primarios máximos}, puesto que el producto de dos $q$ distintos ya no es primario. También con respecto a la descomposición en componentes primarios máximos los $q$ son irreducibles.

4. Todo $q$ es \emph{irreducible} en el sentido de que no puede representarse como el mínimo común múltiplo de dos divisores propios.

La conexión de estos números primarios $q$ con los números primos $p$ consiste en lo siguiente: a cada $q$ le corresponde un $p$, y salvo el signo sólo uno, que es divisor de $q$ y del cual alguna potencia es divisible por $q$: el número primo asociado. Si $p^\varrho$ es la menor potencia de este tipo -- $\varrho$ el exponente de $q$ --, entonces aquí, en particular, $p^\varrho$ es igual a $q$. El teorema de unicidad puede formularse ahora así:

\emph{En dos descomposiciones distintas de un entero racional en los componentes $q$ irreducibles y primarios máximos, coinciden el número de componentes, los números primos asociados (salvo el signo) y los exponentes. Como $p^\varrho=q$, de ello se deduce también la coincidencia de los propios $q$ (salvo el signo).}

La indeterminación que proviene del signo se elimina, como es sabido, si en lugar de los números se consideran los ideales derivados de ellos (todos los números divisibles por $a$); entonces la formulación vale exactamente igual para la descomposición única de los ideales de los cuerpos algebraicos de números finitos en potencias de ideales primos.

En lo que sigue (\S~1) se toma como base un dominio de anillos general que sólo debe satisfacer la condición de finitud de que cada ideal del dominio posea una base ideal finita. Sin tal condición de finitud, los ideales irreducibles y los ideales primos no tendrían por qué existir en absoluto, como muestra el dominio de todos los enteros algebraicos, en el cual no existe descomposición en ideales primos.

Resulta que -- correspondiendo a las cuatro propiedades características de los componentes $q$ -- en general existen cuatro descomposiciones separadas, cada una de las cuales surge de la precedente por subdivisión. La descomposición en ideales coprimos-irreducibles es una representación como producto; las otras tres descomposiciones son representaciones reducidas (\S~2) como mínimos comunes múltiplos. También se conserva la conexión entre un ideal primario -- también los ideales irreducibles son primarios -- y el ideal primo asociado: a todo ideal primario $\ideal Q$ le queda determinado de modo único un ideal primo asociado $\ideal P$, que es divisor de $\ideal Q$ y del cual alguna potencia es divisible por $\ideal Q$. Si $\ideal P^\varrho$ es la menor potencia de esta clase -- $\varrho$ el exponente de $\ideal Q$ --, aquí $\ideal P^\varrho$ no necesita coincidir con $\ideal Q$. El teorema de unicidad se expresa así:

Las descomposiciones 1 y 2 son únicas; en dos descomposiciones distintas 3 o 4 coinciden el número de componentes y los ideales primos asociados;\footnote{Es de presumir que, más allá de esto, también coincidan los exponentes y, aún más generalmente, que los componentes correspondientes sean isomorfos.} los ideales aislados que aparecen entre los componentes (\S~7) están determinados de modo único.

Para la demostración de los teoremas de descomposición, de la condición de finitud se obtiene primero el ``teorema de la cadena finita'', enunciado por Dedekind para módulos numéricos finitos; de él se deduce la representación 4 de cada ideal como mínimo común múltiplo de un número finito de ideales irreducibles. Transformando el concepto de reducibilidad de un componente se obtiene de ahí el teorema fundamental de unicidad para la descomposición 4 en ideales irreducibles. Combinando cada vez un número finito de componentes se asciende a las demás descomposiciones, cuyos teoremas de unicidad resultan entonces como consecuencia del teorema de unicidad 4.

Finalmente se muestra (\S~9) que la representación por un número finito de componentes irreducibles subsiste bajo hipótesis más débiles; no se requiere la conmutatividad del dominio de anillos, y basta considerar, en lugar de un ideal, un módulo respecto del dominio. En este caso más general sigue valiendo la igualdad del número de componentes en dos descomposiciones distintas, mientras que las nociones de primo y primario están ligadas a la conmutatividad y al concepto de ideal; en cambio, la noción de coprimo se conserva para ideales en dominios no conmutativos.

El dominio de anillos más simple para el cual aparecen efectivamente las cuatro descomposiciones separadas es el dominio de todos los polinomios en $n$ variables con coeficientes complejos arbitrarios. Aquí las descomposiciones individuales pueden interpretarse, irracionalmente, por el comportamiento de los entes algebraicos, y el teorema de unicidad de los ideales primos asociados corresponde al teorema fundamental de la teoría de la eliminación sobre la descomposición única de los entes algebraicos en irreducibles. Otros ejemplos los dan todos los dominios íntegros finitos formados por polinomios (\S~10). Pero incluso el dominio simple de todos los números pares, y más generalmente de todos los números divisibles por un número fijo, proporciona ya un ejemplo de descomposiciones parcialmente separadas (\S~11). Un ejemplo de teoría de ideales en dominios no conmutativos lo proporciona la teoría de los divisores elementales (\S~12), donde vale la descomposición única en ideales, respectivamente clases, irreducibles. Estas clases irreducibles caracterizan por completo los constituyentes irreducibles de los divisores elementales, y acaso puedan verse como su equivalente para dominios en los cuales falla la teoría usual de divisores elementales.

Sobre la literatura existente cabe decir lo siguiente. La descomposición en ideales primarios máximos para el dominio de polinomios con coeficientes complejos arbitrarios, respectivamente enteros, fue dada por Lasker y desarrollada por Macaulay en puntos particulares.\footnote{E. Lasker, Zur Theorie der Moduln und Ideale. Math. Ann. 60 (1905), p. 20, teoremas VII y XIII. -- F. S. Macaulay, On the Resolution of a given Modular System into Primary Systems including some Properties of Hilbert Numbers. Math. Ann. 74 (1913), p. 66.} Ambos se apoyan en la teoría de la eliminación y utilizan por tanto el hecho de que un polinomio puede representarse de modo único como producto de polinomios irreducibles. En realidad, los teoremas de descomposición para ideales son independientes de esta premisa, como sugiere la teoría de ideales en cuerpos algebraicos de números y como muestra el presente trabajo. También el ideal primario es definido por Lasker y Macaulay sobre la base de conceptos de la teoría de la eliminación.

La descomposición en ideales irreducibles y la descomposición en ideales irreducibles relativamente primos no parecen haber sido advertidas en la literatura ni siquiera para el dominio de polinomios; sólo en Macaulay hay una observación sobre la unicidad de los ideales primarios aislados.

La descomposición en ideales coprimos-irreducibles fue dada para el dominio de polinomios por Schmeidler,\footnote{W. Schmeidler, Über Moduln und Gruppen hyperkomplexer Größen. Math. Zeitschr. 3 (1919), p. 29.} usando la teoría de la eliminación para la demostración de la finitud. Allí, sin embargo, el teorema de unicidad se formula sólo para clases de ideales, no para los propios ideales. Este último teorema de unicidad aparece en un trabajo conjunto,\footnote{E. Noether--W. Schmeidler, Moduln in nichtkommutativen Bereichen, insbesondere aus Differential- und Differenzenausdrücken. Math. Zeitschr. 8 (1920), p. 1.} donde intervienen ideales en dominios de polinomios no conmutativos. Allí sólo se utiliza la base ideal finita; por tanto, los teoremas y métodos siguen siendo válidos para dominios de anillos generales, y en el presente trabajo se afinan en lo tocante a la igualdad del número (\S~11). Las presentes investigaciones son una fuerte generalización y un desarrollo ulterior de las formaciones conceptuales que subyacen a esos dos trabajos. El rasgo esencial de ambos trabajos es el paso de la representación como mínimo común múltiplo a una descomposición aditiva del sistema de clases residuales. Aquí, para simplificar la exposición, permanecemos en el mínimo común múltiplo; la descomposición aditiva se corresponde entonces con la transformación de la noción de reducibilidad en una propiedad del complemento (\S~3). Sin embargo, mediante consideraciones que corresponden esencialmente a las del trabajo conjunto, todos los teoremas aquí enunciados pueden entenderse también como teoremas de descomposición aditiva para el sistema de clases residuales y ciertos sistemas parciales. Este sistema de clases residuales forma un anillo de la misma generalidad que el anillo tomado originariamente como base; en efecto, todo anillo puede ser considerado como el sistema de clases residuales del ideal que corresponde a la totalidad de las relaciones idénticas entre los elementos del anillo; o también de un sistema parcial de estas relaciones, si se supone que las restantes relaciones ya se satisfacen en el dominio.

Esta observación también sitúa los trabajos de Fraenkel.\footnote{A. Fraenkel, Über die Teiler der Null und die Zerlegung von Ringen. J. f. M. 145 (1914), p. 139. Über gewisse Teilbereiche und Erweiterungen von Ringen. Habilitationsschrift, Leipzig, Teubner, 1916. Über einfache Erweiterungen zerlegbarer Ringe. J. f. M. 151 (1920), p. 121.} Fraenkel considera descomposiciones aditivas de anillos sometidos a condiciones restrictivas tales (existencia de elementos regulares, división por éstos, condición de descomponibilidad) que, para el ideal correspondiente, las cuatro descomposiciones coinciden. Debido a esta coincidencia, su condición de finitud -- que el ideal tenga sólo un número finito de divisores propios, de la cual él se aparta en parte -- tampoco significa una restricción más fuerte que la nuestra. El punto de partida de Fraenkel está condicionado por los objetivos distintos, esencialmente algebraicos, de sus trabajos; mediante extensión algebraica llega luego a anillos más generales con condiciones menos restrictivas.

\subsection*{\S~1. Dominio de anillos, ideal, condición de finitud.}

\textbf{1.} El dominio $\Sigma$ tomado como base debe ser un anillo (conmutativo) en la definición abstracta;\footnote{La definición está tomada de la tesis de habilitación de Fraenkel, omitiendo sus condiciones restrictivas 6, I y II; para ello hubo que incluir la ley conmutativa de la adición. Son, por tanto, las leyes que definen un cuerpo con la invertibilidad de la multiplicación omitida.} esto es, $\Sigma$ consiste en un sistema de elementos $a,b,c,\ldots,f,g,h,\ldots$ en el que se define, como igualdad, una relación que satisface las condiciones usuales; y en el que, mediante dos operaciones (modos de composición), la adición y la multiplicación, de cualesquiera dos elementos del anillo $a$ y $b$ se obtiene siempre de manera única un tercero, respectivamente como suma $a+b$ y como producto $a\cdot b$. El anillo y las operaciones, por lo demás enteramente arbitrarias, deben satisfacer las leyes siguientes:
\begin{enumerate}
\item Ley asociativa de la adición: $(a+b)+c=a+(b+c)$.
\item Ley conmutativa de la adición: $a+b=b+a$.
\item Ley asociativa de la multiplicación: $(a\cdot b)c=a(b\cdot c)$.
\item Ley conmutativa de la multiplicación: $a\cdot b=b\cdot a$.
\item Ley distributiva: $a(b+c)=ab+ac$.
\item Ley de la sustracción irrestricta y única. En $\Sigma$ hay un único elemento $x$ que satisface la ecuación $a+x=b$. (Se denota $x=b-a$.)
\end{enumerate}
De estas propiedades se sigue la existencia del cero; pero un anillo no tiene por qué poseer unidad; y el producto de dos elementos puede anularse sin que uno de los factores se anule. Los anillos para los cuales de la anulación de un producto se deduce siempre la anulación de un factor, y que además poseen unidad, se llamarán dominios íntegros propios. Para la suma finita $a+a+\cdots+a$ introducimos la abreviatura usual $na$, donde los enteros $n$ deben considerarse meramente como signos abreviados, no como elementos del anillo, y se definen recursivamente por $a=1a$, $na+a=(n+1)a$.

\textbf{2.} Por un ideal $\ideal M$\footnote{Los ideales se denotan por letras góticas mayúsculas. $\ideal M$ recuerda el ejemplo del ideal en polinomios que suele llamarse ``módulo'' o módulo de formas. Por lo demás, \S\S~1--3 utilizan sólo la propiedad de módulo y no la propiedad de ideal; compárese \S~9.} en $\Sigma$ se entiende un sistema de elementos de $\Sigma$ que satisface las dos condiciones:
\begin{enumerate}
\item $\ideal M$ contiene, junto con $f$, también $a\cdot f$, donde $a$ es un elemento arbitrario de $\Sigma$.
\item $\ideal M$ contiene, junto con $f$ y $g$, también la diferencia $f-g$; por tanto, junto con $f$ también $nf$ para todo entero $n$.
\end{enumerate}
Si $f$ es un elemento de $\ideal M$, lo expresamos como de costumbre mediante $f\equiv0(\ideal M)$ y decimos que $f$ es divisible por $\ideal M$. Si todo elemento de $\ideal N$ es al mismo tiempo un elemento de $\ideal M$, y por tanto divisible por $\ideal M$, decimos: $\ideal N$ es divisible por $\ideal M$; en símbolos, $\ideal N\equiv0(\ideal M)$. $\ideal M$ se llama divisor propio de $\ideal N$ cuando contiene elementos distintos de los de $\ideal N$, y por tanto no es recíprocamente divisible por $\ideal N$. De $\ideal N\equiv0(\ideal M)$ y $\ideal M\equiv0(\ideal N)$ se sigue $\ideal N=\ideal M$.

También se conservan literalmente las demás nociones conocidas. Por máximo común divisor de dos ideales $\ideal A$ y $\ideal B$ -- $\ideal D=(\ideal A,\ideal B)$ -- entendemos la totalidad de los elementos que pueden representarse en la forma $a+b$, donde $a$ recorre todos los elementos de $\ideal A$ y $b$ todos los elementos de $\ideal B$; $\ideal D$ es de nuevo un ideal. Asimismo, el máximo común divisor de infinitos ideales -- $\ideal D=(\ideal A_1,\ideal A_2,\ldots,\ideal A_\nu,\ldots)$ -- se define como la totalidad de los elementos $d$ representables como sumas de elementos tomados cada vez de sólo un número finito de los ideales: $d=a_{i_1}+a_{i_2}+\cdots+a_{i_n}$; también aquí $\ideal D$ vuelve a ser un ideal.

Si, en particular, el ideal $\ideal M$ contiene un número finito de elementos $f_1,f_2,\ldots,f_\varrho$ tales que
\[
 \ideal M=(f_1\ldots f_\varrho),\qquad
 f=a_1f_1+\cdots+a_\varrho f_\varrho+n_1f_1+\cdots+n_\varrho f_\varrho
\]
para todo $f\equiv0(\ideal M)$, donde los $a_i$ son magnitudes del dominio de anillos y los $n_i$ son enteros, entonces $\ideal M$ se llama ideal finito; $f_1,\ldots,f_\varrho$ una base ideal.

En lo sucesivo tomamos como base sólo anillos $\Sigma$ que satisfacen la condición de finitud: \emph{Todo ideal de $\Sigma$ es finito, por tanto posee una base ideal.}

\textbf{3.} De la condición de finitud se sigue directamente el resultado en que se apoyan todas las consideraciones posteriores.

\textbf{Teorema I (teorema de la cadena finita).}\footnote{Enunciado por primera vez para módulos numéricos por Dedekind: Zahlentheorie, Suppl. XI, \S~172, teorema VIII (4.ª ed.); nuestra demostración y la designación ``cadena'' están tomadas de allí. Para ideales en polinomios véase Lasker, loc. cit., p. 56 (lema). Pero en ambos casos el teorema sólo encuentra aplicación aislada. Nuestras aplicaciones descansan enteramente en el principio de elección.} \emph{Sea $\ideal M,\ideal M_1,\ideal M_2,\ldots,\ideal M_\nu,\ldots$ un sistema infinito numerable de ideales en $\Sigma$, cada uno de los cuales es divisible por el siguiente. Entonces a partir de algún índice finito $n$ todos los ideales son idénticos, $\ideal M_n=\ideal M_{n+1}=\cdots$. En otras palabras: si $\ideal M,\ideal M_1,\ideal M_2,\ldots,\ideal M_s,\ldots$ forman una cadena simplemente ordenada de ideales tal que todo ideal es divisor propio del inmediatamente precedente, entonces la cadena se interrumpe después de un número finito de pasos.}

En efecto, sea $\ideal D=(\ideal M,\ideal M_1,\ideal M_2,\ldots,\ideal M_\nu,\ldots)$ el máximo común divisor del sistema, y sea $f_1,\ldots,f_\varrho$ una base de $\ideal D$, que existe siempre por la condición de finitud. De la hipótesis de divisibilidad se sigue que todo elemento de $\ideal D$ es al mismo tiempo elemento de uno de los ideales de la cadena; pues de
\[
 f=g+h,\qquad g\equiv0(\ideal M_r),\qquad h\equiv0(\ideal M_s),\qquad (r<s)
\]
se obtiene $g\equiv0(\ideal M_s)$ y por tanto $f\equiv0(\ideal M_s)$. Lo mismo vale cuando $f$ es una suma de varios sumandos. Por consiguiente hay un índice finito $n$ tal que
\[
 f_1\equiv0(\ideal M_n),\quad f_2\equiv0(\ideal M_n),\quad\ldots,\quad f_\varrho\equiv0(\ideal M_n).
\]
Como recíprocamente $\ideal M_n\equiv0(\ideal D)$, se tiene $\ideal M_n=\ideal D$; y como además $\ideal M_n\equiv0(\ideal D)$ y $\ideal D=\ideal M_n\equiv0(\ideal M_r)$, también se tiene $\ideal M_r=\ideal D=\ideal M_n$ para todo $r>n$, con lo que el teorema queda probado.

Observamos que, recíprocamente, de este teorema se sigue de nuevo la existencia de una base ideal, de modo que la condición de finitud también podría haberse formulado de esta forma sin bases.

\subsection*{\S~2. Representación de un ideal como mínimo común múltiplo de un número finito de ideales irreducibles.}

El mínimo común múltiplo $[\ideal B_1,\ideal B_2,\ldots,\ideal B_k]$ de los ideales $\ideal B_1,\ideal B_2,\ldots,\ideal B_k$ se define como de costumbre como la totalidad de los elementos divisibles por $\ideal B_1$, por $\ideal B_2$, $\ldots$, y por $\ideal B_k$; en símbolos,
\[
 \text{de } f\equiv0(\ideal B_i),\quad (i=1,2,\ldots,k),
 \quad\text{se sigue}\quad
 f\equiv0([\ideal B_1,\ideal B_2,\ldots,\ideal B_k]),
\]
y recíprocamente. El mínimo común múltiplo es de nuevo un ideal; los ideales $\ideal B_i$ se llaman también los componentes de la descomposición.

\textbf{Definición I.} Una representación $\ideal M=[\ideal B_1\cdots\ideal B_k]$ se llama representación reducida si ningún $\ideal B_i$ está contenido en el mínimo común múltiplo $\ideal U_i$ de los demás ideales, y si ningún $\ideal B_i$ puede reemplazarse por un divisor propio.\footnote{Un ejemplo de representación no reducida es $(x,xy)=[(x),(x^2,xy,y^\lambda)]$ para todo exponente $\lambda\ge2$; la representación correspondiente a $\lambda=1$, $[(x),(x^2,y)]$, es una reducida asociada.} Si las condiciones se satisfacen sólo para el ideal $\ideal B_i$, se dice que la representación es reducida respecto de $\ideal B_i$. El mínimo común múltiplo $\ideal U_i=[\ideal B_1,\ldots,\ideal B_{i-1},\ideal B_{i+1},\ldots,\ideal B_k]$ se llama el complemento de $\ideal B_i$. Las representaciones en las que sólo se satisface la primera condición se llaman representaciones más cortas.

Para representar un ideal como mínimo común múltiplo basta, por tanto, restringirse a representaciones reducidas, de acuerdo con el siguiente resultado.

\textbf{Lema I.} \emph{Toda representación de un ideal como mínimo común múltiplo de un número finito de ideales puede ser reemplazada, al menos de una manera, por una representación reducida; tal representación puede alcanzarse en particular mediante descomposición sucesiva.}

En efecto, sea $\ideal M=[\ideal B_1\cdots\ideal B_k]$ una representación arbitraria de $\ideal M$. Omitimos, por orden, aquellos $\ideal B_i$ que están contenidos en el mínimo común múltiplo de los ideales que se dejan. Puesto que los ideales restantes siguen dando $\ideal M$, en la representación resultante
\[
 \ideal M=[\ideal B_1\cdots\ideal B_k]=[\ideal U_i,\ideal B_i]
\]
se satisface entonces la primera condición, luego es una representación más corta; y esta condición permanece satisfecha si cualquier $\ideal B_i$ se reemplaza por un divisor propio. La segunda condición, sin embargo, siempre puede alcanzarse mediante el teorema de la cadena finita (Teorema I). Pues si
\[
 \ideal B_i\supset \ideal B_i'\supset \ideal B_i''\supset\cdots\supset \ideal B_i^{(\nu)}\supset\cdots,
 \qquad
 \ideal M=[\ideal U_i,\ideal B_i]=[\ideal U_i,\ideal B_i']=\cdots=[\ideal U_i,\ideal B_i^{(\nu)}]=\cdots,
\]
donde cada $\ideal B_i^{(\nu)}$ es un divisor propio del inmediatamente precedente, entonces por ese teorema la cadena $\ideal B_i,\ideal B_i',\ldots,\ideal B_i^{(\nu)},\ldots$ debe terminar después de un número finito de pasos; por tanto, en la representación $\ideal M=[\ideal U_i,\ideal B_i^{(\nu)}]$ el ideal $\ideal B_i^{(\nu)}$ ya no puede ser reemplazado por un divisor propio; y esto vale a fortiori si se reemplaza $\ideal U_i$ por un divisor propio. Aplicando el procedimiento sucesivamente a cada $\ideal B_i$, formando siempre el complemento con los $\ideal B$ ya reducidos, se obtiene una representación reducida.\footnote{Que tal representación no está definida de manera única por la representación dada lo muestra el ejemplo anterior. Para $(x^2,xy)=[(x),(x^2,xy,y^\lambda)]$, donde $\lambda\ge2$, además de $[(x),(x^2,y)]$ también $[(x),(x^2,ux+y)]$, para $u$ arbitrario, es una representación reducida.}

Para obtener tal representación sucesivamente, queda por mostrar que de las representaciones reducidas individuales
\[
 \ideal M=[\ideal B_1,\ideal C_1],\qquad
 \ideal C_1=[\ideal B_2,\ideal C_2],\quad\ldots,\quad
 \ideal C_{k-1}=[\ideal B_k,\ideal C_k]
\]
se sigue que la representación resultante $\ideal M=[\ideal B_1\cdots\ideal B_k,\ideal C_k]$ es reducida. Para ello basta mostrar que de representaciones reducidas
\[
 \ideal M=[\ideal B,\ideal C],\qquad \ideal C=[\ideal C_1,\ideal C_2]
\]
se obtiene una representación reducida $\ideal M=[\ideal B,\ideal C_1,\ideal C_2]$. En efecto, por hipótesis ningún $\ideal B$ está en su complemento; si esto ocurriera para algún $\ideal C_i$, entonces, contra la hipótesis en la primera representación, $\ideal C$ sería reemplazable por un divisor propio, ya que por la segunda representación reducida $\ideal C_1$ y $\ideal C_2$ son divisores propios de $\ideal C$; por tanto la representación es más corta. Además, por hipótesis ningún $\ideal B$ puede ser reemplazado por un divisor propio; si esto fuera posible para algún $\ideal C_i$, correspondería, contra la hipótesis, a reemplazar $\ideal C$ por un divisor propio, puesto que la representación de $\ideal C$ es reducida. Así queda probado el lema.

\textbf{Definición II.} Un ideal $\ideal M$ se llama reducible si puede representarse como mínimo común múltiplo de dos divisores propios; en caso contrario $\ideal M$ se llama irreducible.

Demostramos ahora, por medio del Teorema I de la cadena finita y usando representaciones reducidas:

\textbf{Teorema II.} \emph{Todo ideal es representable como mínimo común múltiplo de un número finito de ideales irreducibles.}\footnote{Que tal representación no es única en general lo muestra el ejemplo anterior: $(x^2,xy)=[(x),(x^2,ux+y)]$ para $u$ arbitrario. Los dos componentes son irreducibles para $u$ arbitrario. En efecto, todos los divisores de $(x)$ son de la forma $(x,g(y))$, donde $g(y)$ denota un polinomio en $y$; por tanto, el mínimo común múltiplo de dos cualesquiera tiene también esta forma, luego es un divisor propio de $(x)$. El ideal $(x^2,ux+y)$ posee sólo el único divisor $(x,y)$, y por tanto es necesariamente irreducible también.}

Un ideal arbitrario $\ideal M$ es, en efecto, o bien irreducible; entonces $\ideal M=[\ideal M]$ es una representación del tipo requerido por el Teorema II; o bien $\ideal M=[\ideal B_1,\ideal C_1]$, donde $\ideal B_1$ y $\ideal C_1$ son divisores propios de $\ideal M$, y la representación puede suponerse reducida por el Lema I. Para $\ideal C_1$ vale la misma alternativa: o bien es irreducible, o bien existe una representación reducida $\ideal C_1=[\ideal B_2,\ideal C_2]$.

Continuando así, se obtiene la serie de representaciones reducidas
\begin{equation}
 \ideal M=[\ideal B_1,\ideal C_1]=[\ideal B_1,\ideal B_2,\ideal C_2]=\cdots=[\ideal B_1,\ldots,\ideal B_n,\ideal C_n]=\cdots .
\tag{1}
\end{equation}
En la cadena $\ideal C_1,\ideal C_2,\ldots,\ideal C_n,\ldots$ cada $\ideal C_n$ es un divisor propio del inmediatamente precedente, y por tanto la cadena termina después de un número finito de pasos; hay un índice $n$ tal que $\ideal C_n$ es irreducible. Por el Lema I, además, la representación $\ideal M=[\ideal U_n,\ideal C_n]$ es reducida; $\ideal C_n$ no puede, por tanto, estar en su complemento $\ideal U_n$, y en la representación $\ideal M=[\ideal U_n,\ideal C_n]$ el ideal $\ideal U_n$ no puede reemplazarse por un divisor propio. Reemplazando, si es necesario, $\ideal U_n$ por un divisor propio,\footnote{De hecho la representación es también reducida respecto de $\ideal U_n$, como se mostrará en \S~3 (Lema IV) como recíproco del Lema I.} de modo que la representación se vuelva reducida, queda mostrado que todo ideal reducible admite una representación reducida como mínimo común múltiplo de un ideal irreducible y un ideal complementario. En la serie (1), por tanto, todos los $\ideal B_i$ pueden suponerse irreducibles sin pérdida de generalidad; la repetición del argumento precedente da la existencia de un $\ideal C_n$ irreducible, con lo que queda probado el Teorema II.

\subsection*{\S~3. Igualdad del número de componentes en dos descomposiciones distintas en ideales irreducibles.}

Para demostrar la igualdad del número, primero hay que expresar la reducibilidad o irreducibilidad de un ideal por propiedades de su complemento.

\textbf{Teorema III.}\footnote{El Teorema III corresponde al paso de módulos a grupos residuales en los trabajos de Schmeidler y de Noether--Schmeidler (compárese la Introducción). Al grupo residual corresponde $\ideal C$, y a $\ideal N_1$ y $\ideal N_2$ los subgrupos en que se descompone el grupo residual.} \emph{Sea la representación más corta $\ideal M=[\ideal U,\ideal C]$ reducida respecto de $\ideal C$. Entonces la condición necesaria y suficiente para que $\ideal C$ sea reducible es la existencia de dos ideales $\ideal N_1$ y $\ideal N_2$, divisores propios de $\ideal M$, tales que}
\begin{equation}
 \ideal N_1\ideal N_2\equiv0(\ideal M),\qquad
 \ideal N_i\equiv0(\ideal U),\qquad
 [\ideal N_1,\ideal N_2]=\ideal M.
\tag{2}
\end{equation}
\emph{De aquí se sigue además: si las condiciones (2) se satisfacen y $\ideal C$ es irreducible, entonces al menos uno de los $\ideal N_i$ no es divisor propio de $\ideal M$; $\ideal N_i=\ideal M$.}

Sea $\ideal C=[\ideal C_1,\ideal C_2]$, donde $\ideal C_1,\ideal C_2$ son divisores propios de $\ideal C$. Entonces
\[
 \ideal M=[\ideal U,\ideal C]=[\ideal U,\ideal C_1,\ideal C_2]
       =[[\ideal U,\ideal C_1],[\ideal U,\ideal C_2]].
\]
Aquí los ideales $[\ideal U,\ideal C_i]$ son divisores propios de $\ideal M$, pues de otro modo $[\ideal U,\ideal C]$ no sería reducida respecto de $\ideal C$. Puesto que también se cumple la divisibilidad por $\ideal U$, queda probada la necesidad de la condición (2). (La representación (2) no es reducida, ya que uno de los $[\ideal U,\ideal C_i]$ puede ser reemplazado por $\ideal C_i$.)

Supongamos recíprocamente que (2) se satisface. Formamos los ideales
\[
 \ideal C_1=(\ideal C,\ideal N_1),\qquad
 \ideal C_2=(\ideal C,\ideal N_2),\qquad
 \ideal C^*=[\ideal C_1,\ideal C_2].
\]
Entonces $\ideal C$ es divisible por $\ideal C_1$ y por $\ideal C_2$, por tanto también por su mínimo común múltiplo $\ideal C^*$. Para mostrar la divisibilidad de $\ideal C^*$ por $\ideal C$, sea
\[
 f\equiv0(\ideal C^*);\quad f\equiv0(\ideal C_1);\quad f\equiv0(\ideal C_2),
 \quad\text{o también}\quad f=c+n_1=t+n_2,
\]
donde $c,t,n_1,n_2$ son magnitudes de $\ideal C,\ideal C,\ideal N_1,\ideal N_2$, y en particular $n_1$ es divisible por $\ideal N_1$. Por tanto la diferencia
\[
 g=c-t=n_2-n_1
\]
es divisible tanto por $\ideal C$ como por $\ideal U$, y por ello por $\ideal M$. Como además $n_1=n_2+m$, $n_1$ (y del mismo modo $n_2$) es divisible por $\ideal N_1$ y por $\ideal N_2$, luego por $\ideal M$. Así
\[
 f=c+m,
 \qquad f\equiv0(\ideal C),
 \qquad \ideal C^*=\ideal C.
\]
Los ideales $\ideal C_1$ y $\ideal C_2$ son aquí divisores propios de $\ideal C$; pues si $\ideal C_1=(\ideal C,\ideal N_1)=\ideal C$, entonces $\ideal N_1$ sería divisible por $\ideal C$, y por consiguiente, a causa de la divisibilidad por $\ideal U$, igual a $\ideal M$, contra la hipótesis. Así $\ideal C=[\ideal C_1,\ideal C_2]$ queda reconocido como reducible; el Teorema III está probado.

Casi el mismo razonamiento muestra también lo siguiente.

\textbf{Lema II.} \emph{Si en una representación más corta $\ideal M=[\ideal U,\ideal C]$ el ideal $\ideal C$ puede reemplazarse por un divisor propio, entonces $\ideal C$ es reducible.}

Supongamos, pues,
\[
 \ideal M=[\ideal U,\ideal C]=[\ideal U,\ideal C_1],
\]
y pongamos
\[
 \ideal C^*=[\ideal C_1,(\ideal U,\ideal C)].
\]
De nuevo $\ideal C$ es divisible por $\ideal C^*$. De $f\equiv0(\ideal C^*)$ se sigue
\[
 f=c_1=a+c.
\]
La diferencia $a=c_1-c$ es por tanto divisible tanto por $\ideal U$ como por $\ideal C_1$, luego por $\ideal M$; así $c_1=c+m$, $f\equiv0(\ideal C)$, $\ideal C^*=\ideal C$. Puesto que $\ideal C_1$ y $(\ideal U,\ideal C)$ son, por hipótesis, divisores propios de $\ideal C$, el ideal $\ideal C=\ideal C^*$ queda reconocido como reducible.

Un $\ideal C$ irreducible, por tanto, no puede ser reemplazado por un divisor propio.

Sean ahora dadas dos representaciones más cortas distintas de $\ideal M$ como mínimo común múltiplo de un número finito de ideales irreducibles:
\[
 \ideal M=[\ideal B_1\cdots\ideal B_k]=[\ideal D_1\cdots\ideal D_l].
\]
Por la observación posterior al Lema II, estas representaciones son al mismo tiempo reducidas. Primero demostramos:

\textbf{Lema III.} \emph{Para cada complemento $\ideal U_i=[\ideal B_1\cdots\ideal B_{i-1}\ideal B_{i+1}\cdots\ideal B_k]$ existe un ideal $\ideal D_j$ tal que $\ideal M=[\ideal U_i,\ideal D_j]$.}

En efecto, poniendo $\ideal M=[\ideal B_i,\ideal C_1]$, $\ideal C_1=[\ideal D_1,\ideal C_2]$, etc., se obtiene
\[
 \ideal M=[\ideal U_i,\ideal B_i]
       =[\ideal U_i,[\ideal D_1,\ideal C_2]]
       =[\ideal B_i,[\ideal U_i,\ideal D_1],\ideal C_2].
\]
Aquí, para $\ideal B_i=[\ideal U_i,\ideal D_1]$ y $\ideal C_i=[\ideal U_i,\ideal C_2]$, se satisfacen las condiciones (2) del Teorema III, puesto que $\ideal M=[\ideal U_i,\ideal B_i]$ es reducida respecto de $\ideal B_i$ y la representación es más corta. Como $\ideal B_i$ se suponía irreducible, uno de los $\ideal N_i$ debe hacerse necesariamente igual a $\ideal M$.

Si $[\ideal U_i,\ideal D_1]=\ideal M$, el lema está probado; si $[\ideal U_i,\ideal C_2]=\ideal M$, entonces correspondientemente $\ideal M=[[\ideal U_i,\ideal D_1],[\ideal U_i,\ideal C_2]]$, donde por el mismo argumento de nuevo un componente debe hacerse igual a $\ideal M$. Continuando de este modo, o bien $\ideal M=[\ideal U_i,\ideal D_j]$ con $j<l$, o bien $\ideal M=[\ideal U_i,\ideal D_1,\ldots,\ideal D_{l-1}]$; como $\ideal D_1\cdots\ideal D_{l-1}=\ideal D_l$, el lema queda probado.

De ello se sigue:

\textbf{Teorema IV.} \emph{En dos representaciones más cortas distintas de un ideal como mínimo común múltiplo de ideales irreducibles, el número de componentes es el mismo.}

Por el lema, para $i=1$,
\[
 \ideal M=[\ideal B_1\cdots\ideal B_k]
       =[\ideal U_1,\ideal D_{j_1}]
       =[\ideal D_{j_1},\ideal B_2,\ldots,\ideal B_k].
\]
Considerando ahora las dos descomposiciones
\[
 \ideal M=[\ideal D_{j_1},\ideal B_2,\ldots,\ideal B_k]
          =[\ideal D_1,\ideal D_2,\ldots,\ideal D_l]
\]
y repitiendo el argumento anterior respecto del complemento $\ideal U_2=[\ideal D_{j_1},\ideal B_3,\ldots,\ideal B_k]$ de $\ideal B_2$, se obtiene
\[
 \ideal M=[\ideal U_2,\ideal B_2]=[\ideal U_2,\ideal D_{j_2}]
          =[\ideal D_{j_1},\ideal D_{j_2},\ideal B_3,\ldots,\ideal B_k],
\]
y, continuando el procedimiento,
\[
 \ideal M=[\ideal D_{j_1},\ideal D_{j_2},\ldots,\ideal D_{j_k}].
\]
Como por hipótesis la representación $\ideal M=[\ideal D_1\cdots\ideal D_l]$ es más corta, de modo que ningún $\ideal D$ puede omitirse, los ideales distintos entre los $\ideal D_{j_i}$ deben agotar todos los $\ideal D$; por tanto $k\ge l$. Si en el lema y en los argumentos siguientes se intercambian por completo los $\ideal B$ y los $\ideal D$, se obtiene correspondientemente $l\ge k$, y por tanto $k=l$, con lo que queda probada la igualdad del número. También se sigue que los ideales $\ideal D_{j_i}$ son todos mutuamente distintos, pues de otro modo en una representación más corta por los $\ideal D_{j_i}$ aparecerían menos de $k$ componentes; así se puede escoger la notación de modo que $\ideal D_{j_i}=\ideal D_i$. Por la misma razón, todas las representaciones intermedias $\ideal M=[\ideal D_{j_1}\cdots\ideal D_{j_r},\ideal B_{r+1},\ldots,\ideal B_k]$ son más cortas y, por la observación posterior al Lema II, reducidas.

La igualdad del número conduce a un recíproco del Lema I:

\textbf{Lema IV.} \emph{Si en una representación reducida se combinan los componentes en grupos y se forman sus mínimos comunes múltiplos, la representación resultante es reducida. En otras palabras, de una representación reducida}
\[
 \ideal M=[\ideal C_1\cdots\ideal C_k]
\]
\emph{se sigue que $\ideal M=[\ideal N_1,\ideal N_2]$ también es reducida, si $\ideal N_1=[\ideal C_1,\ldots,\ideal C_h]$ y $\ideal N_2=[\ideal C_{h+1},\ldots,\ideal C_k]$.}

Primero, $\ideal N_1$ no puede estar en su complemento $\ideal N_2$, puesto que esto no ocurre para ninguno de sus divisores $\ideal C_i$; por tanto la representación es más corta. Para mostrar que $\ideal N_1$ no puede ser reemplazado por un divisor propio, descompónganse los $\ideal C$ en ideales irreducibles $\ideal B$,\footnote{Esto significa siempre una representación más corta, por tanto reducida, mediante los $\ideal B$.} lo que da, por el Lema I, las representaciones reducidas
\[
 \ideal N_1=[\ideal B_1\cdots\ideal B_\lambda],
 \qquad
 \ideal N_2=[\ideal B_{\lambda+1}\cdots\ideal B_\kappa].
\]
Sea $\ideal M=[\ideal N_1,\ideal N_2]$ reducida respecto de $\ideal N_2$, y sea $\ideal N_1^*$ un divisor propio de $\ideal N_1$. Por el Lema II,
\[
 \ideal N_1=[\ideal N_1^*,(\ideal N_1,\ideal N_2)],
\]
y esta representación es necesariamente reducida respecto de $\ideal N_1^*$, pues de lo contrario $\ideal N_1$ también podría ser reemplazado en $\ideal M$ por un divisor propio. Si es necesario, reemplácese también $(\ideal N_1,\ideal N_2)$ por un divisor propio, de modo que se obtenga una representación reducida de $\ideal N_1$. Al descomponer los dos componentes de $\ideal N_1$ en ideales irreducibles, el número $\lambda$ de ideales irreducibles de $\ideal N_1$ se compone entonces aditivamente de los números de los componentes; por tanto, el número de ideales irreducibles correspondiente al divisor propio $\ideal N_1^*$ es necesariamente menor que $\lambda$. Pero entonces la descomposición de $\ideal M=[\ideal N_1^*,\ideal N_2]$ en ideales irreducibles conduciría a menos de $\kappa$ ideales, en contradicción con la igualdad del número. Como caso especial $\varrho=2$, se sigue también que la representación $\ideal M=[\ideal N_1,\ideal N_2]$ es reducida asimismo respecto del complemento $\ideal N_2$.

\subsection*{\S~4. Ideales primarios. Unicidad de los ideales primos asociados en dos descomposiciones distintas en ideales irreducibles.}

En lo que sigue se trata de la conexión entre ideales primarios e ideales irreducibles.

\textbf{Definición III.} Un ideal $\ideal Q$ se llama primario si de $a\cdot b\equiv0(\ideal Q)$ y $a\not\equiv0(\ideal Q)$ se sigue necesariamente que $b^\varkappa\equiv0(\ideal Q)$, donde el exponente $\varkappa$ es un número finito.

La definición también puede formularse así: si un producto $a\cdot b$ es divisible por $\ideal Q$, entonces o bien uno de los factores es divisible, o bien una potencia de cada factor es divisible. Si en particular $\varkappa$ es siempre igual a 1, el ideal se llama ideal primo.

De la definición de ideal primario (respectivamente primo), en virtud de la existencia de bases, se sigue la formulación que sólo involucra productos de ideales:\footnote{Si $\ideal A=(a_1,\ldots,a_r)$ y $\ideal B=(b_1,\ldots,b_s)$, entonces $\ideal A\ideal B$ se entiende como el ideal que tiene por base los productos $a_i b_k$; que éstos, y sólo éstos, forman el ideal producto se sigue de la representación por bases, usando la conmutatividad.}

\textbf{Definición IIIa.} Un ideal $\ideal Q$ se llama primario si de $\ideal A\ideal B\equiv0(\ideal Q)$ y $\ideal A\not\equiv0(\ideal Q)$ se sigue necesariamente que $\ideal B^\varkappa\equiv0(\ideal Q)$. Si $\varkappa$ es siempre igual a 1, el ideal se llama ideal primo. Para un ideal primo $\ideal P$, por tanto, de $\ideal A\ideal B\equiv0(\ideal P)$ y $\ideal A\not\equiv0(\ideal P)$ se sigue siempre que $\ideal B\equiv0(\ideal P)$.

Como IIIa contiene a la Definición III como caso especial $\ideal A=(a)$, $\ideal B=(b)$, todo ideal primario según IIIa es también primario según III. Recíprocamente, sea $\ideal Q$ primario según III y supóngase cumplida la hipótesis de IIIa: $\ideal A\ideal B\equiv0(\ideal Q)$. Entonces o bien se sigue $\ideal A\equiv0(\ideal Q)$, o bien existe al menos una cantidad $a\not\equiv0(\ideal Q)$ tal que $a\ideal B\equiv0(\ideal Q)$ y $a\not\equiv0(\ideal Q)$. Si $b_1,\ldots,b_s$ es una base ideal de $\ideal B$, entonces por la Definición III, puesto que $ab_i\equiv0(\ideal Q)$,
\[
 b_i^{\varkappa_i}\equiv0(\ideal Q)\qquad (i=1,\ldots,s).
\]
Por consiguiente, como
\[
 b=n_1b_1+\cdots+n_sb_s+a_1b_1+\cdots+a_sb_s,
\]
para $\varkappa=\varkappa_1+\cdots+\varkappa_s$ todo producto de $\varkappa$ cantidades $b$ es divisible por $\ideal Q$, lo que prueba que los ideales primarios según III satisfacen la Definición IIIa. En el caso especial de ideales primos $\ideal P$, de $a\ideal B\equiv0(\ideal P)$, por tanto también de $ab\equiv0(\ideal P)$ para todo $b\equiv0(\ideal B)$, y $a\not\equiv0(\ideal P)$, se sigue $b\equiv0(\ideal P)$ y por ello $\ideal B\equiv0(\ideal P)$. Queda así probada la equivalencia de las dos definiciones.

La conexión entre ideales primarios y primos se establece mediante la observación de que la totalidad $\ideal P$ de todos los elementos $p$ con la propiedad de que alguna potencia de $p$ es divisible por $\ideal Q$ forma un ideal primo. Primero es claro que $\ideal P$ es un ideal, ya que junto con $p_1$ y $p_2$ también los elementos $ap_1$ y $p_1-p_2$ tienen la propiedad indicada. Por el mismo argumento de bases aplicado en la definición de IIIa, se sigue además la existencia de un número $\lambda$ tal que $\ideal P^\lambda\equiv0(\ideal Q)$.

Sea ahora
\[
 ab\equiv0(\ideal P),\qquad a\not\equiv0(\ideal P).
\]
Entonces, por la definición de $\ideal P$,
\[
 a^\lambda b^\lambda\equiv0(\ideal Q),\qquad a^\lambda\not\equiv0(\ideal Q),
\]
y por tanto, por la definición de $\ideal Q$,
\[
 b^{\lambda\varkappa}\equiv0(\ideal Q),\qquad\text{y en consecuencia}\qquad b\equiv0(\ideal P),
\]
con lo que $\ideal P$ queda probado como ideal primo. $\ideal P$ se define también como el máximo común divisor de todos los ideales $\ideal B$ con la propiedad de que alguna potencia de $\ideal B$ es divisible por $\ideal Q$. Cada uno de estos $\ideal B$ es divisible por $\ideal P$, y por tanto también lo es su máximo común divisor $\ideal D$. Recíprocamente, $\ideal P$ mismo es un ideal $\ideal B$, y por consiguiente es divisible por $\ideal D$, lo que prueba $\ideal P=\ideal D$. Así $\ideal P$ es un ideal primo que es divisor de $\ideal Q$ y del cual una potencia es divisible por $\ideal Q$; esta propiedad lo define de manera única. Pues de
\[
 \ideal Q\equiv0(\ideal P),\qquad \ideal P^\varrho\equiv0(\ideal Q),
 \qquad
 \ideal Q\equiv0(\ideal B),\qquad \ideal B^\sigma\equiv0(\ideal Q)
\]
se sigue
\[
 \ideal P^\varrho\equiv0(\ideal B),\qquad \ideal B^\sigma\equiv0(\ideal P),
\]
y por la propiedad de los ideales primos,
\[
 \ideal P\equiv0(\ideal B),\qquad \ideal B\equiv0(\ideal P),\qquad \ideal B=\ideal P.
\]
En resumen:

\textbf{Teorema V.} \emph{A todo ideal primario $\ideal Q$ le corresponde uno y sólo un ideal primo $\ideal P$ que es divisor de $\ideal Q$ y del cual una potencia es divisible por $\ideal Q$; $\ideal P$ se llamará el ``ideal primo asociado''.}\footnote{Que el recíproco no vale lo muestra el ejemplo $\ideal M=(x^2,xy)$. El ideal primo $(x)$ satisface todas las condiciones, pero $\ideal M$ no es primario.} \emph{$\ideal P$ se define como el máximo común divisor de todos los ideales $\ideal B$ con la propiedad de que una potencia de $\ideal B$ es divisible por $\ideal Q$. Si $\varrho$ es el menor número tal que $\ideal P^\varrho\equiv0(\ideal Q)$, entonces $\varrho$ se llama el exponente de $\ideal Q$.}\footnote{Que en general no tiene por qué valer $\ideal P^\varrho=\ideal Q$, como ocurre en el dominio de los enteros racionales, respectivamente algebraicos, lo muestra el ejemplo $\ideal Q=(x^2,y)$; $\ideal P=(x,y)$; $\ideal P^2=(x^2,xy,y^2)\equiv0(\ideal Q)$, pero $\ideal Q\not\equiv0(\ideal P^2)$; por tanto $\ideal Q$ es distinto de $\ideal P^2$.}

Ahora demostramos la conexión entre lo primario y lo irreducible:

\textbf{Teorema VI.} \emph{Todo ideal no primario es reducible; en otras palabras, todo ideal irreducible es primario.}\footnote{Que aquí el recíproco no vale lo muestra, por ejemplo, $\ideal Q=(x^2,xy,y^\lambda)=[(x^2,y),(x,y^\lambda)]$, donde $\lambda\ge2$. Aquí $\ideal Q$ es primario pero reducible. Que $\ideal Q$ es primario se sigue de que contiene todos los productos de potencias de dimensión $\lambda$ en $x,y$; por tanto, para todo polinomio sin término constante alguna potencia es divisible por $\ideal Q$. Pero si en $ab\equiv0(\ideal Q)$ el polinomio $b$ contiene un término constante -- luego $b^\varkappa\not\equiv0(\ideal Q)$ para todo $\varkappa$ --, entonces, puesto que los polinomios de base de $\ideal Q$ son homogéneos y cada componente homogénea de $ab$ es divisible por $\ideal Q$, $a$ debe ser divisible por $\ideal Q$.}

Sea $\ideal K$ un ideal no primario, de modo que por la Definición III existe al menos un par de cantidades $a,b$ tales que
\begin{equation}
 ab\equiv0(\ideal K),\qquad
 a\not\equiv0(\ideal K),\qquad
 b^\varkappa\not\equiv0(\ideal K)\quad\text{para todo }\varkappa.
\tag{3}
\end{equation}
Formamos ahora los dos ideales
\[
 \ideal K_1=(\ideal K,a),
 \qquad
 \ideal N_1=(\ideal K,b),
\]
que por (3) son divisores propios de $\ideal K$ y para los cuales, por (3),
\begin{equation}
 \ideal K_1\ideal N_1\equiv0(\ideal K).
\tag{4}
\end{equation}
Para los elementos $f$ del mínimo común múltiplo $\ideal K_2=[\ideal K_1,\ideal N_1]$ vale la alternativa siguiente.

O bien de
\[
 f\equiv0(\ideal K_1),\qquad f\equiv0(\ideal N_1),
 \qquad\text{es decir}\qquad f\equiv a_1 b\;(\ideal K),
\]
se sigue siempre una representación
\[
 f\equiv l_1b\;(\ideal K),
 \qquad l_1\equiv0(\ideal K_1),
\]
y por tanto, por (4), $f\equiv0(\ideal K)$, de modo que $\ideal K_2\equiv0(\ideal K)$; como también $\ideal K\equiv0(\ideal K_2)$, se tiene $\ideal K=\ideal K_2$, y $\ideal K$ queda reconocido como reducible.

O bien existe al menos un $f\equiv0(\ideal K_2)$ para el cual no existe tal $l_1$. Con el $a_1$ correspondiente formamos entonces
\[
 \ideal K_2=(\ideal K,a,a_1),
 \qquad
 \ideal N_2=(\ideal K,b^2).
\]
Entonces, como $a_1b\equiv0(\ideal K_1)$, (4) da de nuevo
\begin{equation}
 \ideal K_2\ideal N_2\equiv0(\ideal K),
\tag{4'}
\end{equation}
y $\ideal K_2$ es un divisor propio de $\ideal K_1$.

La misma alternativa vale para los elementos $f$ de $\ideal K_3=[\ideal K_2,\ideal N_2]$. O bien de
\[
 f\equiv0(\ideal K_2),
 \qquad f\equiv0(\ideal N_2),
 \qquad\text{es decir}\qquad f\equiv a_2b^2\;(\ideal K),
\]
se sigue siempre
\[
 f\equiv l_2b^2\;(\ideal K),
 \qquad l_2\equiv0(\ideal K_2),
\]
y por tanto, por (4'), $f\equiv0(\ideal K)$.

O bien, para al menos un $f$, no existe tal $l_2$, lo que conduce a la formación de $\ideal K_3=(\ideal K,a,a_1,a_2)$ y $\ideal N_3=(\ideal K,b^3)$, con $\ideal K_3\ideal N_3\equiv0(\ideal K)$, donde $\ideal K_3$ es un divisor propio de $\ideal K_2$. Continuando de este modo, definimos en general
\[
 \ideal K_n=(\ideal K,a,a_1,\ldots,a_{n-1}),
 \qquad
 \ideal N_n=(\ideal K,b^n),
\]
\[
 \ideal K_n\ideal N_n\equiv0(\ideal K),
\]
donde los $a_n$ se definen por el hecho de que existe un $f$ tal que
\[
 f\equiv0(\ideal K_n),\quad f\equiv0(\ideal N_n),\quad
 f\equiv a_n b^n\;(\ideal K),\quad\text{pero}\quad
 f\not\equiv l_n b^n\;(\ideal K),\quad l_n\equiv0(\ideal K_n).
\]
Así, en general, $\ideal K_n\ideal N_n\equiv0(\ideal K)$, $\ideal N_n$ es un divisor propio de $\ideal K$ por (3), y $\ideal K_n$ es un divisor propio de $\ideal K_{n-1}$. Por el Teorema I de la cadena finita, la cadena de los $\ideal K_n$ debe terminar después de un número finito de pasos, digamos en $\ideal K_n$. Entonces, para todo $f\equiv0(\ideal K_n)$ y $f\equiv0(\ideal N_n)$, se tiene $f\equiv l_n b^n\;(\ideal K)$ con $l_n\equiv0(\ideal K_n)$, y en consecuencia, por el argumento precedente, $\ideal K=[\ideal K_n,\ideal N_n]$, de modo que $\ideal K$ queda reconocido como reducible.

La \emph{unicidad de los ideales primos asociados} se sigue de lo que acaba de probarse.

Sean
\[
 \ideal M=[\ideal B_1\cdots\ideal B_k]=[\ideal D_1\cdots\ideal D_k]
\]
dos representaciones más cortas, por tanto reducidas, de $\ideal M$ como mínimo común múltiplo de ideales irreducibles, cuyo número coincide por el Teorema IV. Entonces por ese teorema también son representaciones más cortas las representaciones intermedias que allí aparecen (donde, como allí se observó, puede ponerse el índice $j_i=i$):
\[
 \ideal M=[\ideal D_1\cdots\ideal D_{i-1}\ideal B_i\ideal B_{i+1}\cdots\ideal B_k]
       =[\ideal D_1\cdots\ideal D_{i-1}\ideal D_i\ideal B_{i+1}\cdots\ideal B_k]
       =[\ideal U_i,\ideal B_i]=[\ideal U_i,\ideal D_i].
\]
Así
\[
 \ideal U_i\ideal B_i\equiv0(\ideal D_i),\qquad \ideal U_i\not\equiv0(\ideal D_i),
 \qquad
 \ideal U_i\ideal D_i\equiv0(\ideal B_i),\qquad \ideal U_i\not\equiv0(\ideal B_i).
\]
Como por el Teorema VI los ideales irreducibles $\ideal B_i$ y $\ideal D_i$ son primarios, se sigue la existencia de dos números $\lambda_i$ y $\mu_i$ tales que
\begin{equation}
 \ideal B_i^{\lambda_i}\equiv0(\ideal D_i),
 \qquad
 \ideal D_i^{\mu_i}\equiv0(\ideal B_i).
\tag{5}
\end{equation}
Sean $\ideal P_i$ y $\bar{\ideal P}_i$ los ideales primos asociados de $\ideal B_i$ y $\ideal D_i$, respectivamente; así $\ideal P_i^{\varrho_i}\equiv0(\ideal B_i)$ y $\bar{\ideal P}_i^{\sigma_i}\equiv0(\ideal D_i)$. De (5) se sigue entonces
\[
 \ideal P_i^{\lambda_i\varrho_i}\equiv0(\bar{\ideal P}_i),
 \qquad
 \bar{\ideal P}_i^{\mu_i\sigma_i}\equiv0(\ideal P_i),
\]
y por la propiedad de los ideales primos
\[
 \ideal P_i\equiv0(\bar{\ideal P}_i),
 \qquad
 \bar{\ideal P}_i\equiv0(\ideal P_i),
 \qquad
 \ideal P_i=\bar{\ideal P}_i.
\]
Por tanto:

\textbf{Teorema VII.} \emph{En dos representaciones más cortas distintas de un ideal como mínimo común múltiplo de ideales irreducibles, coinciden los ideales primos asociados; entre ellos pueden aparecer también iguales,\footnote{Lo muestra, por ejemplo, el ejemplo de la nota 19: $(x^2,xy,y^\lambda)=[(x^2,y),(x,y^\lambda)]$, donde $\lambda\ge2$. Los dos ideales primos asociados son aquí $(x,y)$.} y aparecen con la misma frecuencia en toda descomposición. En consecuencia, los propios ideales pueden emparejarse al menos de una manera tal que en cada caso una potencia de un ideal $\ideal B_i$ sea divisible por el ideal $\ideal D_i$ que se le asigna, y recíprocamente. Su número coincide por el Teorema IV.}\footnote{Para la unicidad de los ideales ``aislados'' contenidos entre los ideales irreducibles, compárese \S~7.}

\subsection*{\S~5. Representación de un ideal como mínimo común múltiplo de ideales primarios máximos. Unicidad de los ideales primos asociados.}

\textbf{Definición IV.} \emph{Una representación más corta $\ideal M=[\ideal Q_1\cdots\ideal Q_\alpha]$ se llamará un mínimo común múltiplo de ideales primarios máximos si todos los $\ideal Q$ son primarios, pero el mínimo común múltiplo de dos $\ideal Q$ ya no es primario.}

Que existe siempre al menos una representación de esta clase se sigue de la representación de $\ideal M$ como mínimo común múltiplo de ideales irreducibles. Estos ideales son primarios; o ya se tiene una representación por ideales primarios máximos, o bien el mínimo común múltiplo de dos ideales es de nuevo primario. Como con ello el número de ideales disminuye en uno, la repetición del procedimiento conduce, después de un número finito de pasos, a la representación deseada.

Esta representación es reducida por el Lema IV. Recíprocamente, toda representación reducida por ideales primarios máximos surge de este modo, como muestra la descomposición de los $\ideal Q$ en ideales irreducibles.

Para deducir aquí del Teorema VII un teorema de unicidad correspondiente, debe investigarse la conexión con los ideales primos asociados.

\textbf{Teorema VIII.} \emph{Si los ideales primarios $\ideal N_1,\ideal N_2,\ldots,\ideal N_\lambda$ poseen todos el mismo ideal primo asociado $\ideal P$, entonces su mínimo común múltiplo $\ideal Q=[\ideal N_1,\ideal N_2,\ldots,\ideal N_\lambda]$ también es primario y tiene $\ideal P$ como ideal primo asociado. Recíprocamente, si $\ideal Q=[\ideal N_1\cdots\ideal N_\lambda]$ es una representación reducida del ideal primario $\ideal Q$, entonces todos los $\ideal N_i$ son primarios y tienen, como ideal primo asociado, el ideal primo asociado $\ideal P$ de $\ideal Q$.}

Para demostrar la primera parte, obsérvese primero que de $\ideal P^{\varrho_i}\equiv0(\ideal N_i)$ para todo $i$ se sigue también $\ideal P^\tau\equiv0(\ideal Q)$, donde $\tau$ denota el mayor de los índices $\varrho_i$. Como $\ideal P$ es además un divisor de $\ideal Q$, $\ideal P$ es necesariamente el ideal primo asociado si $\ideal Q$ es primario. De
\[
 \ideal A\ideal B\equiv0(\ideal Q),
 \qquad
 \ideal B^k\not\equiv0(\ideal Q)\quad\text{(para todo }k)
\]
se sigue
\[
 \ideal B\not\equiv0(\ideal P),
 \quad\text{luego}\quad
 \ideal B^k\not\equiv0(\ideal N_i)\quad\text{(para todo }k),
\]
y por consiguiente $\ideal A\equiv0(\ideal N_i)$; de ahí $\ideal A\equiv0(\ideal Q)$. Así $\ideal Q$ queda probado primario, y $\ideal P$ como ideal primo asociado.

Recíprocamente, sea primero
\[
 \ideal Q=[\ideal N_1\cdots\ideal N_\lambda]=[\ideal N_i,\ideal N_i']
\]
una representación más corta de $\ideal Q$ por ideales primarios $\ideal N_i$, y sean $\ideal P_i$ los ideales primos asociados. De
\[
 \ideal N_i\ideal N_i'\equiv0(\ideal Q),
 \qquad
 \ideal N_i'\not\equiv0(\ideal N_i)
\]
(puesto que la representación es más corta) se sigue
\[
 \ideal N_i'\not\equiv0(\ideal P_i),
 \qquad\text{pero}\qquad
 \ideal N_i'\equiv0(\ideal P).
\]
Como $\ideal P_i$ es al mismo tiempo un divisor de $\ideal Q$, $\ideal P_i$ coincide así, para todo $i$, con el ideal primo asociado $\ideal P$ de $\ideal Q$.

Queda por mostrar que en toda representación reducida $\ideal Q=[\ideal N_1\cdots\ideal N_\lambda]$ los $\ideal N_i$ son primarios.\footnote{Que aquí la representación reducida es esencial lo muestra el ejemplo de la representación no reducida $\ideal Q=[x^2,xy,y^\lambda]=[(x^2,xy,y^\lambda,yz),(x,y^\lambda)]$, donde $\lambda\ge2$. Aquí $(x^2,xy,y^\lambda,yz)=[(x^2,y),(x,y^\lambda,z)]$ no es primario por lo que se acaba de demostrar; pues esta última representación es una más corta por ideales primarios, pero los ideales primos asociados $(x,y)$ y $(x,y,z)$ son distintos. ($\ideal Q$ es primario por la nota 19.)} Para ello, descompónganse los $\ideal N_i$ en sus ideales irreducibles $\ideal B$; en la representación más corta resultante $\ideal Q=[\ideal B_1\cdots\ideal B_s]$, todo ideal es primario y, por lo que acaba de probarse, tiene $\ideal P$ como ideal primo asociado. Entonces, por la primera parte del teorema demostrada arriba, lo mismo vale para cada $\ideal N_i$, con lo que el Teorema VIII queda completamente probado.

\textbf{Adición.} También se sigue que un ideal primo es necesariamente irreducible. Pues de la representación reducida $\ideal P=[\ideal N,\ideal K]$ se obtiene por el Teorema VIII $\ideal N\equiv0(\ideal P)$; luego, como $\ideal P\equiv0(\ideal N)$, también $\ideal P=\ideal N$, y análogamente $\ideal P=\ideal K$. La irreducibilidad de $\ideal P$ se sigue asimismo directamente: de $\ideal P=[\ideal N,\ideal K]$ se deduce $\ideal N\ideal K\equiv0(\ideal P)$, $\ideal N\not\equiv0(\ideal P)$, $\ideal K\not\equiv0(\ideal P)$, contra la propiedad definitoria de un ideal primo.

Sean ahora
\[
 \ideal M=[\ideal Q_1\cdots\ideal Q_\alpha]=[\ideal Q'_1\cdots\ideal Q'_\beta]
\]
dos representaciones reducidas de $\ideal M$ como mínimo común múltiplo de ideales primarios máximos. Descomponiendo los $\ideal Q$ en ideales irreducibles $\ideal B$, y correspondientemente los $\ideal Q'$ en ideales irreducibles $\ideal D$, se obtienen dos representaciones reducidas de $\ideal M$ como mínimo común múltiplo de ideales irreducibles; por el Teorema VII coinciden tanto el número de componentes como los ideales primos asociados. Por el Teorema VIII, todos los ideales irreducibles $\ideal B$ que aparecen para un $\ideal Q_i$ fijo tienen el mismo ideal primo asociado $\ideal P_i$; mientras que el $\ideal P_j$ asociado a $\ideal Q_j$ es necesariamente distinto de éste, ya que de otro modo, por el Teorema VIII, no habría representación por ideales primarios máximos. Por consiguiente el número $\alpha$ de los $\ideal Q$ es igual al número de ideales primos asociados distintos $\ideal P$ de los $\ideal B$; estos $\ideal P$ distintos forman los ideales primos asociados de los $\ideal Q$. Lo mismo vale para los $\ideal Q'$ respecto de su descomposición en los $\ideal D$. El Teorema VII da por tanto la igualdad del número para los $\ideal Q$ y $\ideal Q'$, así como la coincidencia de sus ideales primos asociados. Al mismo tiempo se ve que la agrupación de los ideales irreducibles en ideales primarios máximos descrita al comienzo del parágrafo consiste en reunir todos, y sólo aquellos, que tienen el mismo ideal primo asociado. El Teorema VIII muestra además la propiedad de irreducibilidad de los ideales primarios máximos: no admiten ninguna representación reducida como mínimo común múltiplo de ideales primarios máximos.

En resumen:

\textbf{Teorema IX.} \emph{En dos representaciones reducidas de un ideal como mínimo común múltiplo de ideales primarios máximos, coinciden el número de componentes y los ideales primos asociados, todos ellos distintos entre sí. En otras palabras, a cada $\ideal Q$ se puede asignar uno y sólo un $\ideal Q'$ tal que una potencia de $\ideal Q$ sea divisible por $\ideal Q'$, y recíprocamente.}\footnote{Un ejemplo de representaciones distintas es el dado en la nota 12 para el Teorema II: $(x^2,xy)=[(x),(x^2,\mu x+y)]$ para $\mu$ arbitrario. Como los ideales primos asociados $\ideal P_1=(x)$ y $\ideal P_2=(x,y)$ son distintos, éstos son ideales primarios máximos. -- Para la unicidad de los ideales primarios máximos ``aislados'', compárese \S~7.} \emph{Los $\ideal Q$ y $\ideal Q'$ poseen la propiedad de irreducibilidad respecto de la descomposición en ideales primarios máximos.}

\textbf{Adición.} Debe notarse que el Teorema IX permanece esencialmente válido si se supone sólo una representación más corta en lugar de una reducida. Si, por ejemplo,
\[
 \ideal M=[\ideal Q_1\cdots\ideal Q_i^*\cdots\ideal Q_\alpha]
\]
es reducida respecto de $\ideal Q_i^*$, y $\ideal Q_i^*$ es un divisor propio de $\ideal Q_i$, mientras $\ideal U_i$ es el complemento de $\ideal Q_i$, entonces por el Lema IV
\[
 \ideal Q_i=[\ideal Q_i^*,(\ideal U_i,\ideal Q_i)],
\]
y esta representación es reducida respecto de $\ideal Q_i^*$. Por el Teorema VIII -- reemplazando $(\ideal U_i,\ideal Q_i)$ por un divisor propio si es necesario al aplicar el teorema -- $\ideal Q_i^*$ es por tanto primario y tiene el mismo ideal primo asociado $\ideal P_i$ que $\ideal Q_i$. Continuando el procedimiento se muestra que a toda representación de este tipo puede asignársele una representación reducida por ideales primarios máximos de tal modo que coincidan el número de componentes y los ideales primos asociados. Así, incluso para representaciones más cortas vale que en dos representaciones distintas coinciden el número de componentes y los ideales primos asociados.

Los ideales primos asociados, mutuamente distintos y definidos de modo único de esta manera, se llamarán simplemente ``los ideales primos asociados de $\ideal M$''.




% --- Paper 19 continuation: sections 6--12 complete ---
\subsection*{\S~6. Representación única de un ideal como mínimo común múltiplo de ideales relativamente primos irreducibles.}

\textbf{Definición V.} \emph{Un ideal $\ideal R$ se llama relativamente primo respecto de $\ideal S$ si de
\[
 \ideal T\ideal R\equiv0(\ideal S)
\]
se sigue necesariamente $\ideal T\equiv0(\ideal S)$. Si tanto $\ideal R$ es relativamente primo respecto de $\ideal S$ como $\ideal S$ respecto de $\ideal R$, entonces $\ideal R$ y $\ideal S$ se llaman mutuamente relativamente primos.}\footnote{La condición de ser relativamente primo no es simétrica. Por ejemplo, $\ideal R=(x^2,y)$ es relativamente primo respecto de $\ideal S=(x)$; pero $\ideal S$ no es relativamente primo respecto de $\ideal R$, pues $\ideal S^2\equiv0(\ideal R)$, mientras que $\ideal S\not\equiv0(\ideal R)$.}

\emph{Un ideal se llama relativamente-primo irreducible si no puede representarse como mínimo común múltiplo de divisores propios mutuamente relativamente primos.}

En particular, si en lugar de $\ideal T$ se toma el máximo común divisor $\ideal T_0$ de todos los $\ideal T$ para los cuales $\ideal T\ideal R\equiv0(\ideal S)$, entonces también $\ideal T_0\ideal R\equiv0(\ideal S)$ y $\ideal S\equiv0(\ideal T_0)$. Así, $\ideal T_0=\ideal S$ cuando $\ideal R$ es relativamente primo respecto de $\ideal S$; y $\ideal T_0$ es un divisor propio de $\ideal S$ cuando $\ideal R$ no es relativamente primo respecto de $\ideal S$.\footnote{El $\ideal T_0$ así definido coincide con el ``módulo residual'' de Lasker; y, en su extensión a módulos de números en lugar de ideales, con el ``cociente'' de dos módulos de Dedekind. Lasker, loc. cit., p. 49; Dedekind (Zahlentheorie), p. 504.}

La prueba de unicidad se apoya en:

\textbf{Teorema X.} \emph{1. Si $\ideal R$ es relativamente primo respecto de los ideales $\ideal S_1,\ldots,\ideal S_\lambda$, entonces $\ideal R$ es también relativamente primo respecto de su mínimo común múltiplo $\ideal S$.}

\emph{2. Si los ideales $\ideal S_1,\ldots,\ideal S_\lambda$ son relativamente primos respecto de $\ideal R$, entonces su mínimo común múltiplo $\ideal S$ es también relativamente primo respecto de $\ideal R$.}

\emph{3. Si $\ideal R$ es relativamente primo respecto de $\ideal S$ y $\ideal S=[\ideal S_1\cdots\ideal S_\lambda]$ es una representación reducida de $\ideal S$, entonces $\ideal R$ es también relativamente primo respecto de cada $\ideal S_i$.}

\emph{4. Si $\ideal S$ es relativamente primo respecto de $\ideal R$, entonces todo divisor $\ideal S_i$ de $\ideal S$ es también relativamente primo respecto de $\ideal R$.}

1. Puesto que $\ideal T\ideal R\equiv0(\ideal S)$ da necesariamente $\ideal T\ideal R\equiv0(\ideal S_i)$, la hipótesis da $\ideal T\equiv0(\ideal S_i)$, y por tanto $\ideal T\equiv0(\ideal S)$.

2. Pongamos
\[
 \widehat{\ideal S}_1=[\ideal S_2\cdots\ideal S_\lambda],\quad
 \widehat{\ideal S}_{12}=[\ideal S_3\cdots\ideal S_\lambda],\quad \ldots,\quad
 \widehat{\ideal S}_{12\ldots\lambda-1}=\ideal S_\lambda.
\]
De $\ideal T\ideal S\equiv0(\ideal R)$ se sigue $\ideal T\ideal S_1\widehat{\ideal S}_1\equiv0(\ideal R)$; por tanto, por la hipótesis, $\ideal T\widehat{\ideal S}_1\equiv0(\ideal R)$. Aplicando sucesivamente el mismo argumento se obtiene $\ideal T\widehat{\ideal S}_{12}\equiv0(\ideal R)$ y finalmente $\ideal T\widehat{\ideal S}_{12\ldots\lambda-1}=\ideal T\ideal S_\lambda\equiv0(\ideal R)$; de ahí $\ideal T\equiv0(\ideal R)$.

3. Sea $\widehat{\ideal S}_i$ el complemento de $\ideal S_i$. De $\ideal T_0\ideal R\equiv0(\ideal S_i)$, junto con $\widehat{\ideal S}_i\ideal R\equiv0(\ideal S_i)$, se sigue que $[\ideal T_0,\widehat{\ideal S}_i]\ideal R\equiv0(\ideal S)$. Si $\ideal T_0$ fuera un divisor propio de $\ideal S_i$, entonces, como $\ideal S=[\ideal S_i,\widehat{\ideal S}_i]$ es reducida, $[\ideal T_0,\widehat{\ideal S}_i]$ sería un divisor propio de $\ideal S$, contra la hipótesis.

4. De $\ideal T\ideal S_i\equiv0(\ideal R)$, donde $\ideal T$ es un divisor propio de $\ideal R$, se sigue también $\ideal T\ideal S\equiv0(\ideal R)$; por tanto $\ideal T$ sería un divisor propio de $\ideal R$, contra la hipótesis.

La Definición V muestra en particular que todo ideal es relativamente primo respecto del ideal unitario $\ideal D$ formado por todos los elementos de $\Sigma$;\footnote{$\ideal D$ desempeña el papel de unidad sólo con respecto a la divisibilidad y al mínimo común múltiplo, no con respecto a la formación de productos. Por ejemplo, $\ideal D=(x)$ para el dominio de todos los polinomios enteros en $x$ sin término constante; y $\ideal D=(2)$ para el dominio de todos los números pares.} sin embargo, los teoremas siguientes se aplican sólo a ideales distintos de $\ideal D$.

Del Teorema X se obtiene primero:

\textbf{Lema V.} \emph{Toda representación de un ideal como mínimo común múltiplo de ideales mutuamente relativamente primos y distintos de $\ideal D$ es reducida.}

Sea
\[
 \ideal M=[\ideal R_1\ideal R_2\cdots\ideal R_\sigma]=[\ideal R_i,\ideal L_i]
\]
una representación de este tipo. Por el Teorema X, 2, también $\ideal L_i$ es relativamente primo respecto de $\ideal R_i$; por tanto $\ideal R_i$ no puede estar contenido en $\ideal L_i$, y la representación es más corta. Pues de $\ideal L_i\equiv0(\ideal R_i)$, siendo $\ideal L_i$ relativamente primo respecto de $\ideal R_i$, se seguiría de $\ideal D\ideal L_i\equiv0(\ideal R_i)$ que $\ideal D\equiv0(\ideal R_i)$, y por tanto $\ideal R_i=\ideal D$, contra la hipótesis.

Si ahora $\ideal R_i$ pudiera sustituirse por un divisor propio $\ideal R_i^*$, entonces por el Lema II $\ideal R_i$ sería reducible y
\[
 \ideal R_i=[\ideal R_i^*,(\ideal R_i,\ideal L_i)].
\]
Si es necesario, sustitúyase $(\ideal R_i,\ideal L_i)$ por un divisor propio $(\ideal R_i,\ideal L_i)^*$ y asimismo sustitúyase $\ideal R_i^*$ por un divisor propio, de modo que se obtenga una representación reducida de $\ideal R_i$. Entonces el Teorema X, 3 muestra que $\ideal L_i$ es también relativamente primo respecto de $(\ideal R_i,\ideal L_i)^*$; y este ideal es distinto de $\ideal D$, porque la representación original era más corta. Pero $(\ideal R_i,\ideal L_i)^*$ está contenido en $(\ideal R_i,\ideal L_i)$ y $(\ideal R_i,\ideal L_i)$ está contenido en $\ideal L_i$, lo que da la contradicción.

El Teorema X da además el siguiente teorema, que media la conexión con los ideales primos asociados.

\textbf{Teorema XI.} \emph{Si $\ideal R$ es relativamente primo respecto de $\ideal S$ y $\ideal S$ es distinto de $\ideal D$, entonces ningún ideal primo asociado de $\ideal R$\footnote{Los ideales primos asociados de $\ideal R$ se definieron al final de \S~5.} es divisible por un ideal primo asociado de $\ideal S$. Recíprocamente, si no se da tal divisibilidad, entonces $\ideal R$ es relativamente primo respecto de $\ideal S$, y desde luego $\ideal S$ es distinto de $\ideal D$.}

Sean
\[
 \ideal R=[\ideal Q_1\cdots\ideal Q_\alpha],
 \qquad
 \ideal S=[\ideal Q_1^*\cdots\ideal Q_\beta^*]
\]
representaciones reducidas de $\ideal R$ y $\ideal S$ por ideales primarios máximos, y sean $\ideal P_1,\ldots,\ideal P_\alpha$ y $\ideal P_1^*,\ldots,\ideal P_\beta^*$ sus ideales primos asociados. Probamos la afirmación en la forma siguiente: si algún $\ideal P$ es divisible por algún $\ideal P^*$, entonces $\ideal R$ no puede ser relativamente primo respecto de $\ideal S$, y recíprocamente.

Supongamos entonces que $\ideal P_\mu\equiv0(\ideal P_\nu^*)$. Entonces también $\ideal Q_\mu^{\varrho}\equiv0(\ideal P_\nu^*)$, y por la definición de $\ideal P_\nu^*$ esto da $\ideal Q_\mu^{\varrho}\equiv0(\ideal Q_\nu^*)$, por tanto $\ideal R^{\varrho}\equiv0(\ideal Q_\nu^*)$. Sea $\ideal R^\tau$ la menor potencia de $\ideal R$ divisible por $\ideal Q_\nu^*$. Si $\tau=1$, entonces $\ideal R$ es divisible por $\ideal Q_\nu^*$; pero, puesto que $\ideal S\ne\ideal D$, $\ideal R$ no es relativamente primo respecto de $\ideal Q_\nu^*$. Si $\tau\ge2$, entonces $\ideal R^{\tau-1}\ideal R\equiv0(\ideal Q_\nu^*)$ mientras $\ideal R^{\tau-1}\not\equiv0(\ideal Q_\nu^*)$, de modo que de nuevo $\ideal R$ no es relativamente primo respecto de $\ideal Q_\nu^*$;\footnote{$\ideal R^0$ no está definido, porque $\Sigma$ no necesita contener una unidad; por eso había que tratar primero el caso $\tau=1$. El caso $\tau=0$ queda también excluido, si $\Sigma$ tiene unidad, por la hipótesis sobre $\ideal S$.} en ambos casos el Teorema X, 3 implica entonces que $\ideal R$ no es relativamente primo respecto de $\ideal S$.

Recíprocamente, si $\ideal R$ no es relativamente primo respecto de $\ideal S$, entonces por el Teorema X, 1 no es relativamente primo respecto de al menos un $\ideal Q_\nu^*$. Por tanto
\[
 \ideal T_0\ideal R\equiv0(\ideal Q_\nu^*),
 \qquad
 \ideal T_0\not\equiv0(\ideal Q_\nu^*),
\]
y como $\ideal Q_\nu^*$ es primario,
\[
 \ideal R^\varrho\equiv0(\ideal Q_\nu^*),
 \qquad
 \ideal Q_1^\varrho\cdots\ideal Q_\alpha^\varrho\equiv0(\ideal Q_\nu^*).
\]
Para los ideales primos asociados esto implica
\[
 \ideal P_1^{\varrho_1}\cdots\ideal P_\alpha^{\varrho_\alpha}\equiv0(\ideal P_\nu^*),
\]
y por tanto, por la propiedad de los ideales primos, $\ideal P_\nu^*$ está contenido en al menos uno de los $\ideal P$. Esto prueba el Teorema XI.

Los Teoremas X y XI dan ahora la existencia\footnote{La existencia de la descomposición también puede probarse directamente, en exacta analogía con la prueba de \S~2 de la existencia de una descomposición en un número finito de ideales irreducibles.} y la unicidad de la descomposición en ideales relativamente-primos irreducibles de la manera siguiente.

Sea $\ideal M=[\ideal Q_1\cdots\ideal Q_\alpha]$ una representación reducida (o al menos más corta) de $\ideal M$ por ideales primarios máximos, y sean $\ideal P_1,\ldots,\ideal P_\alpha$ los ideales primos asociados. Reunimos los $\ideal P$ en grupos de modo que ningún ideal de un grupo sea divisible por un ideal de otro grupo, mientras que un solo grupo no pueda descomponerse en dos subgrupos que ambos tengan esta propiedad.

Para construir tal agrupación, obsérvese que por definición un solo grupo $G$ debe contener, junto con cada ideal $\ideal P$, todos los divisores y múltiplos de $\ideal P$ que aparezcan entre $\ideal P_1,\ldots,\ideal P_\alpha$. Sean, por ejemplo, $\ideal P_j^{(i)}$ todos los múltiplos de $\ideal P$, $\ideal P_{j_1}^{(i_1)}$ todos los divisores de $\ideal P_j^{(i)}$, $\ideal P_{j_1}^{(i_1 i_2)}$ todos los múltiplos de $\ideal P_{j_1}^{(i_1)}$, y así sucesivamente. Como en total aparecen sólo un número finito de ideales $\ideal P$, el procedimiento debe detenerse tras un número finito de pasos. El conjunto de ideales así obtenido es entonces un grupo $G$ con las propiedades requeridas: contiene todos los divisores y todos los múltiplos prescritos por la construcción, y ningún ideal exterior a él puede ser divisor ni múltiplo de un ideal interior a él.

El grupo $G$ satisface también la condición de irreducibilidad. Si se lo dividiera en dos subgrupos $G^{(1)}$ y $G^{(2)}$, y si $G^{(1)}$ contuviera uno de los ideales que aparecen en alguna etapa de la construcción alternada, entonces contendría todos los anteriores, pues éstos son alternativamente múltiplos y divisores (o divisores y múltiplos), y por tanto también $\ideal P$ y con ello el grupo entero $G$. Repitiendo la construcción con los ideales que no están contenidos en $G$, se obtiene una descomposición en grupos $G_1,\ldots,G_\sigma$ de todos los $\ideal P$, y el mismo argumento muestra que esta agrupación es única.

Sea $\ideal Q_{i\mu}$, donde $\mu$ recorre de $1$ a $\lambda_i$, los ideales $\ideal Q$ cuyos ideales primos asociados están reunidos en un grupo $G_i$. Como los $\ideal P$ son todos distintos, los ideales primarios $\ideal Q_{i\mu}$ que pertenecen a una representación más corta están determinados de modo único. Pongamos
\[
 \ideal R_i=[\ideal Q_{i1}\cdots\ideal Q_{i\lambda_i}],
 \qquad
 \ideal M=[\ideal R_1\cdots\ideal R_\sigma].
\]
El Teorema XI sigue siendo aplicable aunque $\ideal R_i$ sea sólo una representación más corta, pues por la adición al Teorema IX sus ideales primos asociados ya están bien definidos. Puesto que ningún ideal primo asociado de $\ideal R_i$ es divisible por uno de $\ideal R_j$, ni recíprocamente, el Teorema XI muestra que $\ideal R_i$ y $\ideal R_j$ son mutuamente relativamente primos; y cada $\ideal R_i$ es relativamente-primo irreducible por la propiedad de irreducibilidad del grupo $G_i$. Por el Lema V la representación es reducida. Recíprocamente, toda descomposición en ideales relativamente-primos irreducibles conduce, por el Teorema XI, a la agrupación que acaba de describirse.

Sea ahora
\[
 \ideal M=[\bar{\ideal R}_1\cdots\bar{\ideal R}_\tau]
\]
una segunda representación de $\ideal M$ por ideales relativamente-primos irreducibles. Es reducida por el Lema V. Al descomponer los $\ideal R$ en ideales primarios máximos se ve que coinciden los ideales primos asociados, y por tanto coinciden las agrupaciones de estos ideales primos. Así $\tau=\sigma$, y puede elegirse la notación de modo que $\ideal R_i$ y $\bar{\ideal R}_i$ pertenezcan al mismo grupo. Si
\[
 \ideal M=[\ideal R_i,\ideal L_i]=[\bar{\ideal R}_i,\bar{\ideal L}_i]
\]
es la representación por ideal y complemento, entonces, puesto que $\bar{\ideal R}_i$ está asignado al mismo grupo que $\ideal R_i$, por el Teorema XI también $\ideal L_i$ es relativamente primo respecto de $\bar{\ideal R}_i$, y $\bar{\ideal L}_i$ relativamente primo respecto de $\ideal R_i$. De
\[
 \ideal R_i\ideal L_i\equiv0(\bar{\ideal R}_i),
 \qquad
 \bar{\ideal R}_i\bar{\ideal L}_i\equiv0(\ideal R_i)
\]
se sigue por tanto
\[
 \ideal R_i\equiv0(\bar{\ideal R}_i),
 \qquad
 \bar{\ideal R}_i\equiv0(\ideal R_i),
 \qquad
 \ideal R_i=\bar{\ideal R}_i.
\]
Así:

\textbf{Teorema XII.} \emph{Todo ideal puede representarse de modo único como mínimo común múltiplo de un número finito de ideales mutuamente relativamente primos y relativamente-primo irreducibles.}

\subsection*{\S~7. Unicidad de los ideales aislados.}

\textbf{Definición VI.} \emph{Si la representación más corta $\ideal M=[\ideal R,\ideal L]$ es reducida respecto de $\ideal L$, entonces $\ideal R$ se llama un ideal aislado si ningún ideal primo asociado de $\ideal R$ está contenido en un ideal primo asociado de $\ideal L$; equivalentemente, si $\ideal L$ es relativamente primo respecto de $\ideal R$.}

Entonces la representación $\ideal M=[\ideal R,\ideal L]$ satisface las condiciones de la representación $\ideal M=[\ideal R_i,\ideal L_i]$ del Lema V, y el Lema V prueba:

\textbf{Lema VI.} \emph{Si $\ideal R$ es un ideal aislado de la representación más corta $\ideal M=[\ideal R,\ideal L]$, reducida respecto de $\ideal L$, entonces la representación es también reducida respecto de $\ideal R$.}

Así, los ideales aislados aparecen sólo en representaciones reducidas. Los ideales primos asociados que surgen en las descomposiciones de $\ideal R$ y de $\ideal L$ en ideales irreducibles se complementan entre sí hasta dar los ideales primos asociados, determinados de modo único, que surgen en la correspondiente descomposición de $\ideal M$. Por tanto, ningún ideal primo asociado perteneciente a la descomposición de $\ideal R$ en ideales irreducibles está contenido en los restantes ideales primos asociados de $\ideal M$. Recíprocamente, si esta condición se cumple y $\ideal R$ aparece en al menos una representación $\ideal M=[\ideal R,\ideal L]$ reducida respecto de $\ideal R$, y por tanto también en una representación reducida $\ideal M=[\ideal R,\ideal L^*]$, entonces $\ideal R$ es aislado por la Definición VI. Esto da la siguiente definición, independiente del complemento particular $\ideal L$.

\textbf{Definición VIa.} \emph{$\ideal R$ se llama un ideal aislado si los ideales primos pertenecientes a su descomposición en ideales irreducibles no están contenidos en los demás ideales primos asociados que pertenecen a la correspondiente descomposición de $\ideal M$, y si $\ideal R$ aparece en al menos una representación $\ideal M=[\ideal R,\ideal L]$ reducida respecto de $\ideal R$.}\footnote{Si se introdujeran ideales primos asociados ordinarios (final de \S~5), habría que añadir la condición especial de que los de $\ideal L$ fueran todos distintos de los de $\ideal R$. Con la Definición VIa, ya no hay que suponer que la representación sea reducida respecto del complemento.}

Supongamos
\[
 \ideal M=[\ideal R,\ideal L]=[\bar{\ideal R},\bar{\ideal L}]
\]
que son dos representaciones de $\ideal M$ por ideales aislados $\ideal R$ y $\bar{\ideal R}$, con complementos $\ideal L$ y $\bar{\ideal L}$, tales que coinciden los ideales primos asociados de $\ideal R$ y de $\bar{\ideal R}$. Sustitúyanse $\ideal L$ y $\bar{\ideal L}$ por divisores $\ideal L^*$ y $\bar{\ideal L}^{*}$ de modo que las representaciones se vuelvan reducidas. Entonces también coinciden los ideales primos asociados de $\ideal L^*$ y $\bar{\ideal L}^{*}$; por el Teorema XI, $\ideal L^*$ es relativamente primo respecto de $\bar{\ideal R}$ y $\bar{\ideal L}^{*}$ es relativamente primo respecto de $\ideal R$. Puesto que
\[
 \ideal R\ideal L^*\equiv0(\bar{\ideal R}),
 \qquad
 \bar{\ideal R}\bar{\ideal L}^{*}\equiv0(\ideal R),
\]
obtenemos
\[
 \ideal R\equiv0(\bar{\ideal R}),
 \qquad
 \bar{\ideal R}\equiv0(\ideal R),
 \qquad
 \ideal R=\bar{\ideal R}.
\]
Así, los ideales aislados están determinados de modo único por sus ideales primos asociados. Esto afina los Teoremas VII y IX sobre descomposiciones en ideales irreducibles, respectivamente en ideales primarios máximos; por la adición al Teorema IX basta suponer una representación más corta. En resumen:

\textbf{Teorema XIII.} \emph{En toda representación más corta de un ideal como mínimo común múltiplo de ideales irreducibles, respectivamente de ideales primarios máximos, los ideales aislados irreducibles, respectivamente primarios máximos, están determinados de modo único; la ambigüedad afecta sólo a los ideales irreducibles, respectivamente primarios máximos, no aislados.}\footnote{Para ideales en dominios de polinomios, este teorema ya fue enunciado sin prueba por Macaulay en el caso de la descomposición en ideales primarios máximos; su definición de ideales primarios aislados y no aislados (incrustados) puede considerarse una forma irracional de la dada aquí.} \emph{En general, los ideales aislados están determinados de modo único por sus ideales primos asociados.}

Si en tal representación más corta por ideales irreducibles, respectivamente por ideales primarios máximos, los ideales $\ideal B_i$, respectivamente $\ideal Q_j$, son no aislados, entonces por definición sus complementos $\ideal U_i$, respectivamente $\ideal L_j$, son divisibles por $\ideal P_i$, respectivamente $\ideal P_j$. Por tanto
\[
 \ideal U_i^{\varrho_i}\equiv0(\ideal B_i)
 \qquad \text{respectivamente} \qquad
 \ideal L_j^{\sigma_j}\equiv0(\ideal Q_j).
\]
Recíprocamente, si estas relaciones se cumplen, entonces $\ideal P_i$, respectivamente $\ideal P_j$, está contenido en al menos un ideal primo asociado del complemento y por tanto, por la adición al Teorema IX, en un ideal primo asociado del divisor $\ideal L_j^*$ de $\ideal L_j$ que suministra una representación reducida respecto de $\ideal L_j^*$. Así $\ideal B_i$, respectivamente $\ideal Q_j$, es no aislado. En particular, los ideales irreducibles $\ideal B_i$ cuyo ideal primo asociado aparece más de una vez en la descomposición de $\ideal M$ son siempre no aislados. \emph{Los ideales primarios no aislados se caracterizan, por tanto, también por el hecho de que una potencia de todo complemento es divisible por ellos; los ideales primarios aislados, por el hecho de que esto no puede ocurrir.}

\subsection*{\S~8. Representación única de un ideal como producto de ideales coprimos-irreducibles.}

Si el dominio de anillos subyacente $\Sigma$ tiene una unidad, es decir, un elemento $\varepsilon$ tal que $\varepsilon a=a$ para todo elemento de $\Sigma$,\footnote{Como la multiplicación es conmutativa, $\Sigma$ puede tener por supuesto a lo sumo una unidad: para dos unidades cualesquiera $\varepsilon_1$ y $\varepsilon_2$ se tiene $\varepsilon_1\varepsilon_2=\varepsilon_2=\varepsilon_1$.} entonces los ideales coprimos pueden definirse como sigue.

\textbf{Definición VIII.} \emph{Dos ideales $\ideal R$ y $\ideal S$ se llaman coprimos si su máximo común divisor es el ideal unitario $\ideal D=(\varepsilon)$ formado por todos los elementos de $\Sigma$. Un ideal se llama coprimo-irreducible si no puede representarse como mínimo común múltiplo de ideales dos a dos coprimos.}

Obsérvese que dos ideales coprimos son siempre mutuamente relativamente primos. Por definición hay elementos $r\equiv0(\ideal R)$ y $s\equiv0(\ideal S)$ con $\varepsilon=r+s$. De $\ideal T\ideal R\equiv0(\ideal S)$ se sigue $\ideal T r\equiv0(\ideal S)$, de ahí $\ideal T\varepsilon=\ideal T\equiv0(\ideal S)$; correspondientemente, de $\ideal T\ideal S\equiv0(\ideal R)$ se sigue $\ideal T\equiv0(\ideal R)$.\footnote{La recíproca es falsa: por ejemplo, $\ideal R=(x)$ y $\ideal S=(y)$ son mutuamente relativamente primos, pero no coprimos.} Por tanto el Lema V implica que toda representación por ideales dos a dos coprimos es también reducida.

La prueba de unicidad se apoya en el análogo del Teorema X:

\textbf{Teorema XIV.} \emph{Si $\ideal R$ es coprimo con cada uno de los ideales $\ideal S_1,\ldots,\ideal S_\lambda$, entonces $\ideal R$ es también coprimo con $\ideal S=[\ideal S_1\cdots\ideal S_\lambda]$. Recíprocamente, si $\ideal R$ y $\ideal S$ son coprimos, entonces $\ideal R$ es coprimo con cada $\ideal S_j$. Si $\ideal R=[\ideal R_1\cdots\ideal R_\mu]$ y todo $\ideal R_i$ es coprimo con todo $\ideal S_j$, entonces $\ideal R$ y $\ideal S$ son coprimos; también aquí vale la recíproca.}

En efecto, si $\ideal R$ es coprimo con cada $\ideal S_j$, entonces existen elementos $s_j$ tales que
\[
 s_j\equiv0(\ideal S_j),
 \qquad
 s_j\equiv\varepsilon(\ideal R).
\]
En consecuencia
\[
 s_1s_2\cdots s_\lambda\equiv0(\ideal S),
 \qquad
 s_1s_2\cdots s_\lambda\equiv\varepsilon(\ideal R),
 \qquad
 (\ideal R,\ideal S)=(\varepsilon).
\]
Como $(\ideal R,\ideal S)$ es divisible por cada $(\ideal R,\ideal S_j)$, se obtiene también la recíproca. La aplicación repetida del mismo argumento da la segunda parte. En efecto, si $\ideal R_i$ es coprimo con $\ideal S_1,\ldots,\ideal S_\lambda$ para un $i$ fijo, entonces $\ideal R_i$ es coprimo con $\ideal S$; si esto vale para todo $i$, entonces, como la coprimalidad es mutua, $\ideal S$ es coprimo con $\ideal R$. Recíprocamente, de la coprimalidad de $\ideal R$ y $\ideal S$ se sigue la coprimalidad de $\ideal S$ con $\ideal R_i$, y por tanto de $\ideal R_i$ con $\ideal S$.

La prueba de existencia y unicidad\footnote{La existencia de la descomposición puede probarse de nuevo directamente en analogía con \S~2; también la prueba de unicidad puede realizarse directamente (compárese lo dicho en la introducción sobre Schmeidler y Noether--Schmeidler). La prueba dada aquí proporciona además una visión de la estructura de los ideales coprimo-irreducibles.} de la descomposición en ideales coprimo-irreducibles consiste, como en la prueba correspondiente para ideales relativamente primos, en una agrupación única. Sin embargo, como los ideales relativamente-primos irreducibles $\ideal R_1,\ldots,\ideal R_\sigma$ de $\ideal M$ están definidos de modo único por el Teorema XII, aquí no es necesario volver a los ideales primos asociados.

Reunimos los ideales relativamente-primos irreducibles $\ideal R_1,\ldots,\ideal R_\sigma$ de $\ideal M$, definidos de modo único, en grupos de tal manera que \emph{todo ideal de un grupo sea coprimo con todo ideal de cualquier grupo distinto, mientras que un solo grupo no pueda escindirse en dos subgrupos tales que todo ideal de un subgrupo sea coprimo con todo ideal del otro subgrupo}. Tal agrupación se obtiene como sigue. Por definición, un grupo $G$ debe contener, junto con todo ideal $\ideal R$, todos los ideales no coprimos con $\ideal R$. Sean éstos $\ideal R_{i_1}$; sean $\ideal R_{i_1i_2}$ los ideales no coprimos con éstos; en general, sean $\ideal R_{i_1\ldots i_\lambda}$ no coprimos con $\ideal R_{i_1\ldots i_{\lambda-1}}$. Como sólo intervienen finitamente muchos ideales, este proceso termina tras un número finito de pasos. El conjunto de ideales así obtenido forma un grupo $G$ con las propiedades deseadas: todos los ideales exteriores a $G$ son, por construcción, coprimos con todos los ideales de $G$. Si $G$ se escindiera en dos subgrupos $G^{(1)}$ y $G^{(2)}$, y si $\ideal R_{i_1\ldots i_\lambda}$ perteneciera a $G^{(1)}$, entonces, como la coprimalidad es mutua, $G^{(1)}$ tendría que contener también $\ideal R_{i_1\ldots i_{\lambda-1}},\ldots,\ideal R_{i_1}$, de ahí también $\ideal R$ y por tanto el grupo entero $G$. Esto prueba la irreducibilidad. Repitiendo la construcción con los ideales restantes se obtiene una agrupación $G_1,\ldots,G_\tau$ de todos los $\ideal R$.

La agrupación es única: si $G'_1,\ldots,G'_\tau$ es una segunda agrupación y $\ideal R_{i_1\ldots i_\lambda}$ es un elemento de $G'_i$, entonces $G'_i$ contiene $\ideal R$ y por tanto $G_1$, y, por la condición de irreducibilidad, no puede contener elementos distintos de $G_1$; por tanto $G'_i=G_1$.

Sea $\ideal T_i$ el mínimo común múltiplo de los ideales $\ideal R$ reunidos en un grupo $G_i$. Mostramos que
\[
 \ideal M=[\ideal T_1\cdots\ideal T_\tau]
\]
es una representación de $\ideal M$ por ideales coprimo-irreducibles. Descomponer primero los $\ideal T$ en los $\ideal R$ muestra que $[\ideal T_1\cdots\ideal T_\tau]$ representa realmente $\ideal M$. Por el Teorema XIV los $\ideal T$ son dos a dos coprimos, y cada $\ideal T$ es coprimo-irreducible. Con esto queda probada la existencia de tal representación.

Para la unicidad, sea $\ideal M=[\bar{\ideal T}_1\cdots\bar{\ideal T}_\tau]$ una segunda representación de este tipo. Al descomponer los $\bar{\ideal T}$ en ideales relativamente-primos irreducibles $\ideal R$, el Teorema XIV muestra que los ideales $\ideal R$ que aparecen en distintos $\bar{\ideal T}$ son coprimos entre sí, y por tanto también mutuamente relativamente primos; por ello coinciden con los ideales relativamente-primos irreducibles $\ideal R$ de $\ideal M$, definidos de modo único. Los ideales $\bar{\ideal T}_i$ inducen también, por el Teorema XIV, una agrupación $G'_i$ de los $\ideal R$ con las propiedades indicadas. Como esta agrupación es única y cada $\bar{\ideal T}_i$ está determinado de modo único por el grupo $G'_i=G_i$, se tiene $\bar{\ideal T}_i=\ideal T_i$; queda probada la unicidad.

Para ideales dos a dos coprimos, el mínimo común múltiplo es además igual al producto. Pues por el Teorema XIV, si $\ideal M=[\ideal T_1\cdots\ideal T_\tau]$, el complemento $\ideal L_i$ es también coprimo con $\ideal T_i$; por tanto hay elementos
\[
 t_i\equiv0(\ideal T_i),
 \qquad
 l_i\equiv0(\ideal L_i),
 \qquad
 \varepsilon=t_i+l_i.
\]
De $f\equiv0(\ideal T_i)$ y $f\equiv0(\ideal L_i)$ se sigue que
\[
 f=f\varepsilon=ft_i+fl_i,
 \qquad
 f\equiv0(\ideal L_i\ideal T_i).
\]
Recíprocamente, $\ideal L_i\ideal T_i$ es divisible por $[\ideal L_i,\ideal T_i]$, de donde $[\ideal L_i,\ideal T_i]=\ideal L_i\ideal T_i$; al continuar el procedimiento con $\ideal L_i$ se obtiene
\[
 \ideal M=[\ideal T_1\cdots\ideal T_\tau]
       =\ideal T_1\ideal T_2\cdots\ideal T_\tau.
\]
Por tanto:

\textbf{Teorema XV.} \emph{Todo ideal puede representarse de modo único como producto de un número finito de ideales dos a dos coprimos y coprimo-irreducibles.}

\subsection*{\S~9. Extensión de la investigación a módulos. Igualdad del número de componentes en descomposiciones en módulos irreducibles.}

Mostramos ahora que el contenido de los tres primeros parágrafos, que se refiere a ideales irreducibles y no a ideales primarios ni primos, permanece válido bajo hipótesis más débiles. En efecto, esos parágrafos no usan la ley conmutativa de la multiplicación y se refieren sólo a la propiedad de los ideales de ser módulos. Por tanto permanecen válidos para módulos respecto de dominios no conmutativos, que ahora deben definirse. La definición de los módulos ha de fundarse en un dominio doble $(\Sigma,T)$ con las propiedades siguientes.

$\Sigma$ es un anillo no conmutativo definido abstractamente, es decir, $\Sigma$ es un sistema de elementos $a,b,c,\ldots$ para los cuales están definidos dos modos de combinación, la adición de anillo $(+)$ y la multiplicación de anillo $(\times)$, que satisfacen las leyes indicadas en \S~1, con excepción de la ley conmutativa 4 de la multiplicación de anillo.

$T$ es un sistema de elementos $\alpha,\beta,\gamma,\ldots$ para el cual, en conexión con $\Sigma$, se definen asimismo dos combinaciones: la adición, que de dos elementos cualesquiera $\alpha,\beta$ produce un tercer elemento único $\alpha+\beta$; y la multiplicación de un elemento $\alpha$ de $T$ por un elemento $c$ de $\Sigma$, que produce de modo único un elemento $c\alpha$ de $T$.\footnote{Así, aquí se trata de una multiplicación ``por la derecha'', de un dominio $T$ ``por la derecha'' y, por consiguiente, de módulos e ideales ``por la derecha''. Si se basara $T$ en cambio en una multiplicación por la izquierda $\alpha c$, resultaría una teoría correspondiente de módulos e ideales por la izquierda; $M$ contendría $\alpha c$ junto con $\alpha$. La ley asociativa tendría entonces la forma $(\gamma b)a=\gamma(b\times a)$.}

Para estas combinaciones valen las leyes siguientes:
\begin{enumerate}
\item La ley asociativa de la adición: $(\alpha+\beta)+\gamma=\alpha+(\beta+\gamma)$.
\item La ley conmutativa de la adición: $\alpha+\beta=\beta+\alpha$.
\item La ley de la sustracción ilimitada y única: existe en $T$ un único elemento $\xi$ que satisface la ecuación $\alpha+\xi=\beta$ (se escribe $\xi=\beta-\alpha$).
\item La ley asociativa de la multiplicación: $a(b\gamma)=(a\times b)\gamma$.
\item La ley distributiva:
\[
 (a+b)\gamma=a\gamma+b\gamma,
 \qquad
 c(\alpha+\beta)=c\alpha+c\beta.
\]
\end{enumerate}
De estas condiciones se sigue, como es sabido, la existencia del elemento cero y la validez de la ley distributiva también para la combinación de sustracción y multiplicación:
\[
 (a-b)\gamma=a\gamma-b\gamma,
 \qquad
 c(\alpha-\beta)=c\alpha-c\beta,
\]
donde $(-)$ designa la sustracción en $\Sigma$. Si $\Sigma$ contiene una unidad $\varepsilon$, debe cumplirse $\varepsilon\alpha=\alpha$ para todo elemento $\alpha$ de $T$.

Por un módulo $M$ en $(\Sigma,T)$ entendemos un sistema de elementos de $T$ que satisface las dos condiciones:
\begin{enumerate}
\item junto con $\alpha$, $M$ contiene $c\alpha$, donde $c$ es cualquier elemento de $\Sigma$;
\item junto con $\alpha$ y $\beta$, $M$ contiene la diferencia $\alpha-\beta$, y por tanto también $n\alpha$ para todo entero $n$.\footnote{Los enteros han de entenderse otra vez como signos abreviados, no como elementos del anillo.}
\end{enumerate}
Por esta definición, $T$ mismo forma un módulo en $(\Sigma,T)$. Si en particular el dominio $T$ y las combinaciones establecidas allí coinciden con el dominio $\Sigma$ y las combinaciones válidas allí, entonces el módulo $M$ se convierte en un ideal derecho $M$ en $\Sigma$. Si además $\Sigma$ se supone conmutativo, se obtiene el concepto ordinario de ideal, que es por tanto un caso especial del concepto de módulo.\footnote{El ejemplo más sencillo de un módulo es el módulo de las formas lineales enteras: aquí $\Sigma$ consta de todos los enteros racionales y $T$ de todas las formas lineales enteras. Un módulo algo más general se obtiene si en $\Sigma$ y $T$ se toman enteros algebraicos en lugar de enteros racionales, o por ejemplo todos los enteros pares. Si, en lugar de las formas lineales, se considera cada complejo de coeficientes como un elemento, entonces las combinaciones en $\Sigma$ y $T$ son de hecho distintas. Los ideales en dominios de polinomios no conmutativos son objeto del trabajo conjunto de Noether--Schmeidler. Los ideales en otros dominios no conmutativos especiales se tratan en las lecciones de Hurwitz sobre la teoría de números de los cuaterniones (Berlín, Springer 1919) y en los trabajos de Du Pasquier citados allí.}

Para los módulos siguen siendo válidas todas las definiciones de \S~1. Así, $\alpha\eqzero{M}$, respectivamente $N\eqzero{M}$, significa que $\alpha$, respectivamente todo elemento de $N$, es un elemento de $M$; con otras palabras, que $\alpha$, respectivamente $N$, es divisible por $M$. $M$ es un divisor propio de $N$ si $M$ contiene elementos distintos de los de $N$; de $N\eqzero{M}$ y $M\eqzero{N}$ se sigue $M=N$. Las definiciones de máximo común divisor y mínimo común múltiplo permanecen asimismo literalmente iguales. Si el módulo $M$ contiene, en particular, un número finito de elementos $\alpha_1,\ldots,\alpha_\rho$ tales que
\[
 M=(\alpha_1,\ldots,\alpha_\rho),
\]
es decir,
\[
 \alpha=c_1\alpha_1+\cdots+c_\rho\alpha_\rho+n_1\alpha_1+\cdots+n_\rho\alpha_\rho
\]
para todo $\alpha\eqzero{M}$, donde los $c_i$ son elementos de $\Sigma$ y los $n_i$ son enteros, entonces $M$ se llama un módulo finito y $\alpha_1,\ldots,\alpha_\rho$ una base del módulo.

En lo que sigue suponemos, igual que en \S~1, sólo aquellos dominios $(\Sigma,T)$ que satisfacen la condición de finitud: todo módulo en $(\Sigma,T)$ es finito, y por tanto tiene una base de módulo.

Para tal dominio $(\Sigma,T)$, el Teorema I sobre cadenas finitas vale también para módulos, como muestra la prueba dada allí, y por consiguiente se satisfacen todas las hipótesis para \S\S~2 y 3. La Definición I y el Lema I sobre representaciones más cortas y reducidas se trasladan directamente; igualmente el Teorema II sobre la representabilidad de todo módulo como mínimo común múltiplo de un número finito de módulos irreducibles, donde el Lema II muestra que toda representación más corta de esta clase es al mismo tiempo reducida. También permanece válido el Teorema III, que expresa la reducibilidad de un módulo por medio de propiedades de su complemento; de él se sigue el Lema III y finalmente el Teorema IV, que afirma la igualdad del número de componentes en dos representaciones más cortas distintas de un módulo como mínimo común múltiplo de módulos irreducibles. El Teorema IV da también el Lema IV como recíproco del Lema I sobre representaciones reducidas.

\emph{Adición.} El mismo argumento muestra que todos estos teoremas y definiciones permanecen válidos si, para dominios biláteros $T$, es decir, dominios que son tanto derechos como izquierdos, se entiende en todas partes por módulo un módulo bilateral: uno que, junto con $\alpha$, contiene también $c\alpha$ y $\alpha c$, y junto con $\alpha$ y $\beta$ contiene también $\alpha-\beta$.

Mientras que todos estos teoremas descansan sólo en el concepto de divisibilidad y en el de mínimo común múltiplo, los teoremas de unicidad posteriores descansan esencialmente en el concepto de producto y por tanto no admiten una transferencia directa. Así, la definición de ideal primario y de ideal primo no puede trasladarse a módulos, pues el producto de dos elementos de $T$ no está definido. La transferencia formalmente posible a anillos no conmutativos pierde su sentido, pues aquí no puede probarse la existencia del ideal primo asociado\footnote{De $p_1^\rho\eqzero{\ideal Q}$ y $p_2^\sigma\eqzero{\ideal Q}$ no se sigue aquí que $(p_1-p_2)^{\rho+\sigma}\eqzero{\ideal Q}$; del mismo modo, de $(ab)^\rho\eqzero{\ideal Q}$ no se sigue que $a^\rho b^\rho\eqzero{\ideal Q}$. Así, no puede probarse ni que $\ideal P$ sea un ideal ni que tenga la propiedad de un ideal primo. Para ideales biláteros $\ideal P$ sí puede probarse que es un ideal, pero no que sea primo.} y falla también la prueba de que un ideal irreducible sea primario. Estos dos hechos constituyen la base de los teoremas de unicidad siguientes. En cambio, si el anillo no conmutativo tiene unidad, pueden definirse ideales coprimos y coprimo-irreducibles, y los argumentos utilizados en la prueba del Teorema II muestran que todo ideal puede representarse como mínimo común múltiplo de un número finito de ideales dos a dos coprimos y coprimo-irreducibles.\footnote{Para ideales especiales, ``completamente reducibles'', también aquí pueden establecerse teoremas de unicidad; compárese el trabajo conjunto Noether--Schmeidler.}

Finalmente mencionamos un criterio suficiente para que se cumpla la condición de finitud en $(\Sigma,T)$: si $\Sigma$ contiene una unidad y satisface la condición de finitud, y si $T$ mismo es un módulo finito en $(\Sigma,T)$, entonces todo módulo en $(\Sigma,T)$ es finito.

En efecto, de la existencia de la unidad en $\Sigma$, para
\[
 \ideal M=(f_1,\ldots,f_e),
 \qquad
 f=\bar b_1f_1+\cdots+\bar b_ef_e+n_1f_1+\cdots+n_ef_e,
\]
donde los $\bar b_i$ son elementos de $\Sigma$ y los $n_i$ son enteros, se sigue también una representación
\[
 f=b_1f_1+\cdots+b_ef_e
\]
con los $b_i$ en $\Sigma$. Pues, como $f_i=\varepsilon f_i$, se tiene $n_if_i=n_i\varepsilon f_i$, y como $n_i\varepsilon=(\varepsilon+\cdots+\varepsilon)$ pertenece a $\Sigma$, $b_i=\bar b_i+n_i\varepsilon$ es un elemento de $\Sigma$. La hipótesis sobre $T$ dice que todo elemento $\alpha$ de $T$ admite una representación
\[
 \alpha=\bar a_1\tau_1+\cdots+\bar a_k\tau_k+n_1\tau_1+\cdots+n_k\tau_k,
\]
y por tanto
\[
 \alpha=a_1\tau_1+\cdots+a_k\tau_k,
\]
la segunda representación siguiéndose de la primera exactamente como antes, porque $\varepsilon\tau_i=\tau_i$.

Cuando $\alpha$ recorre los elementos de un módulo $M$ en $(\Sigma,T)$, los coeficientes $a_k$ de $\tau_k$ recorren un ideal $\ideal M_k$ de $\Sigma$; por lo anterior, para todo $a_k\eqzero{\ideal M_k}$ se tiene
\[
 a_k=b_1a_k^{(1)}+\cdots+b_ea_k^{(e)}.
\]
Sea $\alpha^{(i)}$ un elemento de $M$ cuyo coeficiente de $\tau_k$ es $a_k^{(i)}$. Entonces $\alpha-b_1\alpha^{(1)}-\cdots-b_e\alpha^{(e)}$ es un elemento perteneciente a $M$ que depende sólo de $\tau_1,\ldots,\tau_{k-1}$. El procedimiento puede repetirse para la totalidad de estos elementos, que ya no contienen $\tau_k$ y que forman un módulo $M'$, y la afirmación queda probada tras un número finito de repeticiones.

\subsection*{\S~10. Caso especial del dominio de polinomios.}

1. Sea el dominio de anillos subyacente $\Sigma$ el conjunto de todos los polinomios en $x_1,\ldots,x_n$ con coeficientes complejos arbitrarios. Por el teorema de la base de Hilbert (Ann. 36), la condición de finitud se cumple para él. Se trata de la relación de nuestros teoremas con los teoremas conocidos de la teoría de eliminación y de la teoría de módulos.

Esta relación se establece mediante el siguiente caso especial de un teorema conocido de Hilbert:\footnote{Über die vollen Invariantensysteme. Math. Ann. 42 (1893), \S~3, p. 313.}

Si $f$ se anula para todo sistema finito de valores de $x_1,\ldots,x_n$ que sea un cero de todos los polinomios de un ideal primo $\ideal P$ (un cero de $\ideal P$), entonces $f$ es divisible por $\ideal P$. En otras palabras, un ideal primo $\ideal P$ consta de la totalidad de los polinomios que se anulan en sus ceros.\footnote{Como mostró Lasker [Math. Ann. 60 (1905), p. 607], este caso especial también puede probarse directamente para formas homogéneas; recíprocamente, de él se sigue de nuevo el teorema de Hilbert en los casos homogéneo e inhomogéneo. Podemos expresar ese teorema como sigue: si un ideal $\ideal R$ se anula en todos los ceros de $M$, entonces una potencia de $\ideal R$ es divisible por $M$. Este teorema, y también el caso especial, vale sólo si el dominio de valores de los $x$ es algebraicamente cerrado; por tanto no puede seguirse de nuestros teoremas solos, sino que debe usar la existencia de raíces, por ejemplo en el punto en que se usa que un ideal cuyos ceros consisten sólo en $x_1=0,\ldots,x_n=0$ contiene todos los productos de potencias de los $x$ a partir de alguna dimensión. El resto de la prueba puede simplificarse algo en comparación con Lasker usando nuestros teoremas. Pues Lasker debe apelar también al teorema de Hilbert para probar la descomposición de un ideal en ideales primarios máximos.}

Así, si un producto $fg$ se anula en todos los ceros de $\ideal P$, al menos uno de los factores se anula; los ceros forman una variedad algebraica irreducible. Recíprocamente, si se adopta esta definición de variedad irreducible, entonces la totalidad de los polinomios que se anulan en una variedad irreducible forma un ideal primo. Ideal primo y variedad irreducible se corresponden por tanto biunívocamente. Como $\ideal P^\rho\eqzero{\ideal Q}$ y $\ideal Q\eqzero{\ideal P}$, los ceros de un ideal primario coinciden con los de su ideal primo asociado.\footnote{Macaulay (compárese la introducción) toma esta propiedad de un ideal primario, a saber, poseer una variedad irreducible, como definición; Lasker incluye en la definición sólo el concepto de la multiplicidad de una variedad y por lo demás define abstractamente. Los ideales primarios que se anulan sólo para $x_1=0,\ldots,x_n=0$ ocupan para Lasker una posición especial.} Por tanto la representación de un ideal como mínimo común múltiplo de ideales primarios máximos da una resolución de todos los ceros del ideal en variedades irreducibles; y, como mostró Lasker, también vale la recíproca. La prueba de unicidad de los ideales primos asociados corresponde aquí al teorema fundamental de la teoría de eliminación sobre la descomponibilidad única de una variedad algebraica en variedades irreducibles, y para dominios especiales de polinomios, en los que no existe representación única de los polinomios como producto de polinomios irreducibles del dominio y por consiguiente no existe teoría de eliminación, puede servir como equivalente de aquel teorema de eliminación.

Las variedades irreducibles que corresponden a los ideales primarios aislados son exactamente aquellas que aparecen en la ``resolvente mínima'';\footnote{Compárese, por ejemplo, J. König, Einleitung in die allgemeine Theorie der algebraischen Größen (Leipzig, Teubner, 1903), p. 235.} pues los ceros de todo ideal primario máximo no aislado son al mismo tiempo ceros de al menos uno aislado, a saber, uno cuyo ideal primo asociado es divisible por el del ideal no aislado. La unicidad de los ideales primarios aislados da por tanto, en los exponentes, nuevos números invariantes de multiplicidad. Los teoremas de unicidad para la resolución de ideales primarios en irreducibles pueden considerarse igualmente suplementos de la teoría de eliminación en el sentido de la multiplicidad.

Después de estas observaciones, las distintas descomposiciones pueden interpretarse por su comportamiento respecto de las variedades algebraicas. A los ideales dos a dos coprimos corresponden variedades que no tienen ningún cero común; en los ideales mutuamente relativamente primos, ninguna variedad irreducible de un ideal es al mismo tiempo un cero común; los ideales primarios máximos se anulan sólo en variedades irreducibles todas distintas entre sí; en la descomposición en ideales irreducibles, las mismas variedades irreducibles pueden aparecer con multiplicidad.

Debe señalarse además que, en lugar del dominio general de polinomios, puede tomarse también como base el dominio de todas las formas homogéneas; se verifica fácilmente que los teoremas generales permanecen válidos para las combinaciones allí vigentes, donde la adición está definida sólo para formas de la misma dimensión.\footnote{Que en el caso ambiguo pueden existir descomposiciones inhomogéneas junto a homogéneas lo muestra el ejemplo $(x^3,xy,y^3)=[(x^3,y),(y^3,x)]=[(xy,x^3,y^3,x+y^3),(xy,x^3,y^3,y+x^2)]$.}

Un ejemplo más sencillo de las cuatro descomposiciones diferentes -- cuyas fórmulas siguen abajo -- viene dado, de acuerdo con lo anterior, por una recta y dos rectas que la cortan y son oblicuas a ella, una de estas últimas conteniendo, con multiplicidad mayor, un punto distinto del punto de intersección. La descomposición en ideales coprimo-irreducibles corresponde a la descomposición en la recta y la variedad oblicua a ella; en la descomposición en ideales relativamente-primos irreducibles esta variedad se descompone en las dos rectas; la descomposición en ideales primarios máximos corresponde a separar el punto de mayor multiplicidad, mientras que la descomposición en ideales irreducibles fuerza una resolución de ese punto.

Tomando este punto como origen, la recta que pasa por él como eje $y$, la recta que la corta paralela al eje $x$, y la recta oblicua a ella paralela al eje $z$, tal configuración está representada, por ejemplo, por los siguientes ideales irreducibles:\footnote{De éstos, los tres primeros son irreducibles como ideales primos; $\ideal B_4$ porque sólo tiene los divisores $(x^2,y,z)$ y $(x,y,z)$; $\ideal B_5$ porque todo divisor contiene el polinomio $xy$.}
\[
\begin{gathered}
 \ideal B_1=(x-1,y),\quad
 \ideal B_2=(y-1,z),\quad
 \ideal B_3=(x,z),\quad
 \ideal B_4=(x^3,y,z),\quad
 \ideal B_5=(x^2,y^2,z).
\end{gathered}
\]
Los ideales primos asociados son
\[
 \ideal P_1=\ideal B_1,
 \quad \ideal P_2=\ideal B_2,
 \quad \ideal P_3=\ideal B_3,
 \quad \ideal P_4=\ideal P_5=(x,y,z).
\]
Los ideales primarios máximos son
\[
 \ideal Q_1=\ideal B_1,
 \quad \ideal Q_2=\ideal B_2,
 \quad \ideal Q_3=\ideal B_3,
 \quad \ideal Q_4=[\ideal B_4,\ideal B_5]=(x^3,y^2,x^2y,z).
\]
Los ideales relativamente-primos irreducibles son
\[
 \ideal R_1=\ideal Q_1,
 \quad \ideal R_2=\ideal Q_2,
 \quad \ideal R_3=[\ideal Q_3,\ideal Q_4]=(x^3,x^2y,xy^2,z).
\]
Los ideales coprimo-irreducibles son
\[
 \ideal S_1=\ideal R_1,
 \quad
 \ideal S_2=[\ideal R_2,\ideal R_3]=((y-1)x^3,(y-1)x^2y,(y-1)xy^2,z).
\]
Así el ideal total es
\[
\begin{aligned}
 \ideal M&=[\ideal S_1,\ideal S_2]=\ideal S_1\ideal S_2 \\
 &=((x-1)(y-1)x^3,(y-1)x^2y,(y-1)xy^2,(x-1)z,y(y-1)x^3,yz),
\end{aligned}
\]
donde
\[
 1=-(y-1)x^3+(y-1)(x^3-1)+y,
\]
\[
 (y-1)(x^3-1)+y\eqzero{\ideal S_1},
 \qquad -(y-1)x^3\eqzero{\ideal S_2}.
\]
Aquí los ideales $\ideal B_1$, $\ideal B_2$, $\ideal B_3$ son aislados, y por tanto están determinados de modo único también en las descomposiciones en ideales irreducibles y primarios máximos. Los ideales $\ideal B_4$ y $\ideal B_5$, respectivamente $\ideal Q_4$, son no aislados; no están determinados de modo único, pero pueden, por ejemplo, reemplazarse por
\[
 \ideal D_4=(x^3,y+\lambda x^2,z),
 \qquad
 \ideal D_5=(x^2+\mu xy,y^2,z),
\]
y $\ideal Q_4$ puede reemplazarse igualmente por
\[
 \overline{\ideal Q}_4=(x^3,x^2y,y^2+\lambda xy,z).
\]

2. Así como los dominios generales, y los dominios íntegros, de polinomios satisfacen la condición de finitud, también lo hace todo dominio íntegro finito de polinomios, por el teorema de Hilbert sobre bases de módulos;\footnote{Recíprocamente, si un dominio de polinomios satisface la condición de finitud y todo polinomio tiene al menos una representación en la que los multiplicadores son de grado menor en $x$, entonces el dominio es un dominio íntegro finito.} los coeficientes pueden tomarse en un cuerpo arbitrario. Damos otro ejemplo, para el dominio de todos los polinomios pares, el dominio más simple en el cual, por $x^2y^2=(xy)^2$, no existe representación única de los polinomios como producto de polinomios irreducibles del dominio. Es la misma configuración que en el ejemplo precedente, dada ahora por los ideales irreducibles
\[
\begin{gathered}
 \ideal B_1=(x^2-1,xy,y^2,yz),\quad
 \ideal B_2=(y^2-1,xz,yz,z^2),\\
 \ideal B_3=(x^2,xy,xz,yz,z^2),\quad
 \ideal B_4=(x^4,x^3y,y^2,xz,yz,z^2),\\
 \ideal B_5=(x^2,y^2,xz,yz,z^2).
\end{gathered}
\]
Los ideales primos asociados son
\[
 \ideal P_1=\ideal B_1,
 \quad \ideal P_2=\ideal B_2,
 \quad \ideal P_3=\ideal B_3,
 \quad \ideal P_4=\ideal P_5=(x^2,xy,y^2,xz,yz,z^2).
\]
Que $\ideal P_1$ es un ideal primo se sigue de que todo polinomio del dominio tiene la forma
\[
 f=\varphi(z^2)+xz\,\psi(z^2)\pmod{\ideal P_1}.
\]
Así, si
\[
 f_1=\varphi_1(z^2)+xz\,\psi_1(z^2),
 \qquad
 f_2=\varphi_2(z^2)+xz\,\psi_2(z^2),
\]
entonces de $f_1f_2\eqzero{\ideal P_1}$ se obtienen las ecuaciones
\[
 \varphi_1\varphi_2+z^2\psi_1\psi_2=0,
 \qquad
 \varphi_2\psi_1+\varphi_1\psi_2=0,
\]
y de ahí $f_1\eqzero{\ideal P_1}$ o $f_2\eqzero{\ideal P_1}$. Exactamente del mismo modo se ve que $\ideal P_2$ es un ideal primo; $\ideal P_3$ es un ideal primo porque todo polinomio del dominio módulo $\ideal P_3$ es congruente con un polinomio en $y^2$; $\ideal P_4$ consta de todos los polinomios del dominio. Así, $\ideal B_1$, $\ideal B_2$ y $\ideal B_3$ son irreducibles como ideales primos, mientras que $\ideal B_4$ y $\ideal B_5$ tienen cada uno sólo el único divisor propio $\ideal P_4$, y por tanto son necesariamente irreducibles.

De los ideales irreducibles se obtienen los ideales primarios máximos
\[
\begin{gathered}
 \ideal Q_1=\ideal B_1,
 \quad \ideal Q_2=\ideal B_2,
 \quad \ideal Q_3=\ideal B_3,
 \quad \ideal Q_4=[\ideal B_4,\ideal B_5]=(x^4,x^3y,y^2,xz,yz,z^2),
\end{gathered}
\]
los ideales relativamente-primos irreducibles
\[
 \ideal R_1=\ideal Q_1,
 \quad \ideal R_2=\ideal Q_2,
 \quad
 \ideal R_3=[\ideal Q_3,\ideal Q_4]=(x^4,x^3y,x^2y^2,xy^3,xz,yz,z^2),
\]
y los ideales coprimo-irreducibles
\[
 \ideal S_1=\ideal R_1,
 \quad \ideal S_2=[\ideal R_2,\ideal R_3]
 =((y^2-1)x^4,(y^2-1)x^3y,(y^2-1)x^2y^2,
 (y^2-1)xy^3,xz,yz,z^2).
\]
Como en el primer ejemplo,
\[
 [\ideal B_4,\ideal B_5]=[\ideal D_4,\ideal D_5],
 \qquad
 [\ideal Q_3,\ideal Q_4]=[\ideal Q_3,\overline{\ideal Q}_4],
\]
donde
\[
 \ideal D_4=(x^4,xy+\lambda x^2,\ldots),
 \qquad
 \ideal D_5=(x^2+\mu xy,\ldots),
 \qquad \lambda\mu\ne1,
\]
y
\[
 \overline{\ideal Q}_4=(x^4,x^3y,y^2+\lambda xy,\ldots).
\]
Los restantes ideales irreducibles, respectivamente primarios máximos, $\ideal B_1$, $\ideal B_2$, $\ideal B_3$, están determinados de modo único como ideales aislados.

\subsection*{\S~11. Ejemplos de la teoría de números y de la teoría de las expresiones diferenciales.}

1. Sea el dominio subyacente $\Sigma$ el conjunto de todos los enteros racionales pares. Así $\Sigma$ puede ponerse en correspondencia biunívoca con el dominio de todos los enteros racionales, asignando a cada número $2a$ de $\Sigma$ el número $a$. De ello se sigue inmediatamente que todo ideal de $\Sigma$ es un ideal principal $(2a)$, donde en la representación de base $2c=n\cdot2a$ para todo elemento $2c$ del ideal, los enteros impares $n$ son sólo abreviaturas de sumas finitas.

Los ideales primos del dominio son $\ideal P_0=\ideal D=(2)$ y $\ideal P=(2p)$, donde $p$ designa un número primo impar; por tanto todo ideal primo es divisible por $\ideal P_0$, pero por ningún otro ideal primo. Los ideales primarios son
\[
 \ideal Q_0=(2\cdot2^{e_0})
 \qquad\text{y}\qquad
 \ideal Q_p=(2p^e).
\]
Son al mismo tiempo ideales irreducibles, y por lo dicho sobre los ideales primos, dos ideales que pertenecen a números primos impares distintos son mutuamente relativamente primos, pero ningún $\ideal Q$ es relativamente primo respecto de ningún $\ideal Q_0$.\footnote{En efecto, de $2b\cdot2p_1^{e_1}\eqzero{2p_2^{e_2}}$ para $p_1\ne p_2$ impares se obtiene siempre $2b\eqzero{2p_2^{e_2}}$; pero de $2b\cdot2p^e\eqzero{2\cdot2^{e_0}}$ se obtiene sólo $2b=2\cdot2^{e_0-1}\pmod{2\cdot2^{e_0}}$.} Así, los $\ideal Q_0$ son los únicos ideales primarios no aislados. La descomposición única de $a$ en potencias de primos corresponde a la representación única del ideal $(2a)$ por ideales primarios máximos, a la vez irreducibles:
\[
 (2a)=[(2\cdot2^{e_0}),(2p_1^{e_1}),
\ldots,(2p_\mu^{e_\mu})].
\]
Por tanto aquí, a diferencia de los ejemplos del dominio de polinomios, incluso los ideales primarios máximos no aislados están determinados de modo único. Si $e_0=0$, ésta es al mismo tiempo una representación por ideales mutuamente relativamente primos, mientras que si $e_0>0$ el ideal es relativamente-primo irreducible.

Así, mientras en el dominio de todos los enteros coinciden las cuatro descomposiciones distintas, aquí esto ocurre sólo para las dos descomposiciones en ideales primarios máximos y en ideales irreducibles, por una parte, y, para $e_0>0$, para los ideales coprimo-irreducibles (todo ideal es coprimo-irreducible, pues el dominio no tiene unidad) y relativamente-primos irreducibles, por otra; para $e_0=0$, en cambio, las descomposiciones coprimo-irreducible y relativamente-primo irreducible se vuelven distintas. Al mismo tiempo se obtiene ya aquí un ejemplo en el cual un ideal primo puede ser divisible por otro sin ser idéntico a él; más generalmente, se ve que la representación por producto no se sigue de la divisibilidad. Este último hecho -- consecuencia de la ausencia de unidad en el dominio -- es también la razón por la cual no existe representación única por producto de los números del dominio mediante números irreducibles del dominio, aunque todo ideal sea principal; la introducción de mínimos comunes múltiplos se revela por tanto necesaria. Obsérvese además que la situación permanece exactamente igual si, en lugar de todos los enteros pares, se toman todos los enteros divisibles por un número primo fijo o por una potencia de un primo.

Sin embargo, los ideales irreducibles y primarios se vuelven distintos si $\Sigma$ consta de todos los números divisibles por un número compuesto $g=p_1^{\nu_1}\cdots p_s^{\nu_s}$. De nuevo todo ideal es un ideal principal $(g\cdot a)$, y los ideales primos vuelven a ser $\ideal P_0=\ideal D=(g)$ y $\ideal P=(g\cdot p)$, donde $p$ es un primo distinto de los primos que aparecen en $g$. Pero los ideales irreducibles son $\ideal B_{\lambda_i}=(g\cdot p_i^{\lambda_i})$ y $\ideal B_e=(g\cdot p^e)$; los ideales primarios son $\ideal Q_e=\ideal B_e$ y los ideales, distintos de los irreducibles,
\[
 \ideal Q_{\lambda_1\ldots\lambda_s}=(g\cdot p_1^{\lambda_1}\cdots p_s^{\lambda_s}),
\]
donde los $\ideal B_{\lambda_i}$ y $\ideal Q_{\lambda_1\ldots\lambda_s}$ tienen todos el mismo ideal primo asociado $\ideal P_0=(g)$. También aquí la descomposición en ideales irreducibles es única, y por consiguiente lo es la descomposición en ideales primarios máximos; de nuevo, por tanto, los ideales no aislados están determinados de modo único.

2. Un ejemplo de dominio de anillos no conmutativo lo suministra la teoría de ideales en dominios de polinomios no conmutativos tratada en el trabajo de Noether--Schmeidler. En particular se tienen ideales ``completamente reducibles'', es decir, ideales cuyos componentes de descomposición son dos a dos coprimos y no tienen divisores propios; por consiguiente los componentes son \emph{a fortiori} irreducibles. Así, además del isomorfismo probado allí según \S~9, se obtiene la igualdad del número de componentes en dos descomposiciones distintas. Para la descomposición de sistemas de expresiones diferenciales lineales parciales u ordinarias, que aparece como caso especial de aquel trabajo, esto da un resultado que no parece haber sido advertido ni siquiera en el caso conocido de una sola expresión diferencial lineal ordinaria.

Al mismo tiempo, el sistema $T$ de todas las clases residuales de un ideal fijo $M$, junto con el dominio de polinomios no conmutativo $\Sigma$, suministra un dominio doble $(\Sigma,T)$, donde $T$ tiene la propiedad de módulo respecto de $\Sigma$. En efecto, la diferencia de dos clases residuales es de nuevo una clase residual, e igualmente el producto de una clase residual por un polinomio arbitrario; en cambio, el producto de dos clases residuales no existe (loc. cit., \S~3). Así, los sistemas de clases residuales llamados allí ``subgrupos'' proporcionan ejemplos de módulos en dominios dobles $(\Sigma,T)$, donde el dominio de anillos subyacente $\Sigma$ es no conmutativo.

\subsection*{\S~12. Ejemplo de la teoría de los divisores elementales.}

Se trata de una concepción de la teoría de los divisores elementales dictada por los desarrollos generales; la teoría misma, sin embargo, se presupone conocida.

Sea $\Sigma$ el dominio de todas las matrices enteras con $n^2$ elementos, con adición y multiplicación definidas en el sentido usual para matrices. Entonces $\Sigma$ es un dominio de anillos no conmutativo; por consiguiente los ideales son en general uniláteros, y los ideales biláteros sólo son un caso especial.\footnote{La teoría de ideales de estos dominios es el objeto de los trabajos de Du Pasquier: Zahlentheorie der Tettarionen, disertación, Zürich, Vierteljahrsschr. d. Naturf. Ges. Zürich, 51 (1906); Zur Theorie der Tettarionenideale, ibid., 52 (1907). El contenido del segundo trabajo es la prueba de que todo ideal es principal.}

Primero mostramos que todo ideal es principal. Para ello, en el caso de ideales derechos, a cada matriz $A=(a_{ik})$ le asignamos el módulo
\[
 A=(a_{11}\xi_1+\cdots+a_{1n}\xi_n,
 \ldots,
 a_{n1}\xi_1+\cdots+a_{nn}\xi_n)
\]
de formas lineales enteras. Recíprocamente, a este módulo corresponde toda matriz que suministra una base de $A$, y por tanto, junto a $A$, también $UA$, donde $U$ es unimodular. Más generalmente, al producto $PA$ corresponde un módulo $B$ que es un múltiplo de $A$. Una sola forma lineal de $A$ viene dada por matrices $P$ que contienen sólo una fila no nula.

Sean ahora $A_1,A_2,\ldots,A_\nu,\ldots$ todos los elementos de un ideal $\ideal M$; sean $A_1,A_2,\ldots,A_\nu,\ldots$ los módulos asignados a ellos; sea $A$ su máximo común divisor y $UA$ la matriz más general asignada a este módulo $A$. A cada forma lineal individual de $A$ corresponde entonces, por la definición de máximo común divisor, una matriz $P_1A_{i_1}+\cdots+P_\sigma A_{i_\sigma}$, donde, como antes, los $P$ tienen sólo una fila no nula. De ello se sigue que la matriz $A$ correspondiente a una base de $A$ admite también una representación de este tipo, ahora con $P$ generales, y por tanto se convierte en elemento de $\ideal M$. Como, además, todo módulo $A_i$ es divisible por $A$, toda matriz $A_i$ es divisible por $A$; $A$, y en general $UA$, forma una base de $\ideal M$. Si se trata de ideales izquierdos, entonces de modo correspondiente deben considerarse las columnas de cada matriz como la base de un módulo; todo ideal es un ideal principal, para el cual, junto a $A$, también $AV$ es una base, con $V$ una matriz unimodular arbitraria.

Para la conexión con la teoría de divisores elementales, nos basamos ahora en ideales biláteros; por tanto, en particular, $PAQ$ pertenece al ideal siempre que $A$ pertenezca. Por lo anterior, la base más general de un ideal de este tipo es $UAV$, donde $U$ y $V$ son unimodulares. Los elementos de base agotan por tanto una clase de matrices equivalentes,\footnote{Los elementos de base de ideales uniláteros corresponden a clases derechas o izquierdas.} y hay una relación biunívoca entre ideal y clase, y por tanto también entre ideal y sistema de divisores elementales $(a_1\mid a_2\mid\cdots\mid a_n)$ de la clase, donde los $a_i$ son, como es sabido, enteros no negativos, cada uno de los cuales divide al siguiente. Así, la matriz de la clase que aparece en la forma normal determinada por los divisores elementales puede considerarse como una base especial del ideal; la divisibilidad de ideales, respectivamente de clases, implica la divisibilidad de los divisores elementales, y recíprocamente.

Pero Du Pasquier ha mostrado, loc. cit. \S~11, que para todo ideal bilateral el rango es $n$ y todos los divisores elementales coinciden. Para poder considerar el caso de divisores elementales generales, debemos partir por tanto no de los ideales, sino directamente de las clases biláteras (que asimismo se denotarán con letras alemanas mayúsculas). Una clase $\mathfrak A=UAV$ es divisible por otra clase $\mathfrak B$ si $A=PBQ$.

En general: \emph{el mínimo común múltiplo (el máximo común divisor) de dos clases se obtiene formando el mínimo común múltiplo (el máximo común divisor) de los sistemas correspondientes de divisores elementales}.\footnote{En el trabajo Zur Theorie der Moduln, Math. Ann. 52 (1899), p. 1, E. Steinitz define el mínimo común múltiplo (el máximo común divisor) de clases mediante los mínimos comunes múltiplos y máximos comunes divisores de los sistemas elementales. Independientemente, el mínimo común múltiplo de clases aparece como ``composición de congruencias'' en H. Brandt, Komposition der binären quadratischen Formen relativ einer Grundform, J. f. M. 150 (1919), p. 1.} En efecto, sean
\[
 (a_1\mid a_2\mid\cdots\mid a_n),\quad
 (b_1\mid b_2\mid\cdots\mid b_n),\quad
 (c_1\mid c_2\mid\cdots\mid c_n),
\]
donde $c_i=[a_i,b_i]$, los sistemas de divisores elementales de $\mathfrak A$, $\mathfrak B$ y $\mathfrak C^*$, respectivamente, y sea $\mathfrak C=[\mathfrak A,\mathfrak B]$. Entonces $\mathfrak C^*$ es divisible por $\mathfrak A$ y $\mathfrak B$, luego por $\mathfrak C$; recíprocamente, los divisores elementales de $\mathfrak C$ son divisibles por los de $\mathfrak C^*$, de modo que $\mathfrak C$ es divisible por $\mathfrak C^*$ y por tanto $\mathfrak C=\mathfrak C^*$. La prueba para el máximo común divisor se obtiene de modo correspondiente.

La representación única de los divisores elementales $a_i$ como mínimos comunes múltiplos de potencias de primos corresponde por tanto a una representación de $\mathfrak A$ como mínimo común múltiplo de clases $\mathfrak Q$ cuyos divisores elementales están dados por potencias de un solo número primo, simbólicamente
\[
 \mathfrak Q\sim(p^{r_1}\mid p^{r_2}\mid\cdots\mid p^{r_\rho}\mid0\mid\cdots\mid0),
 \qquad r_1\le r_2\le\cdots\le r_\rho.
\]
Si en particular el rango $\rho$ es igual a $n$, entonces se tiene aquí una descomposición en clases coprimas y coprimo-irreducibles, que es única a pesar del dominio no conmutativo.

Las clases $\mathfrak Q$ pueden descomponerse más en clases que corresponden a los sistemas de divisores elementales
\[
\begin{gathered}
 \mathfrak B_1\sim(p^{r_1}\mid\cdots\mid p^{r_1}),\quad
 \mathfrak B_2\sim(1\mid p^{r_2}\mid\cdots\mid p^{r_2}),\quad \ldots,\\
 \mathfrak B_\nu\sim(1\mid\cdots\mid1\mid p^{r_\nu}\mid\cdots\mid p^{r_\nu}),\quad
 \mathfrak B_{\rho+1}\sim(1\mid\cdots\mid1\mid0\mid\cdots\mid0),
\end{gathered}
\]
donde en $\mathfrak B_\nu$ el número $1$ aparece $(\nu-1)$ veces. Si, por ejemplo, $r_1=r_2=\cdots=r_\mu=0$, entonces $\mathfrak B_1,\ldots,\mathfrak B_\mu$ son iguales a la clase unidad y deben omitirse de la descomposición; lo mismo vale para $\mathfrak B_{\rho+1}$ si $\rho=n$. Si además, digamos, $r_\nu=r_{\nu+1}=\cdots=r_{\nu+\lambda}$, entonces $\mathfrak B_{\nu+1},\ldots,\mathfrak B_{\nu+\lambda}$ son divisores propios de $\mathfrak B_\nu$, y también deben descartarse. Denotemos por $\mathfrak B_{i_1},\ldots,\mathfrak B_{i_k}$ las que quedan y dan ahora una representación más corta. Entonces
\[
 \mathfrak Q=[\mathfrak B_{i_1},\ldots,\mathfrak B_{i_k}]
\]
es la descomposición única de $\mathfrak Q$ en clases irreducibles. En efecto, supóngase que
\[
 \mathfrak B_\nu\sim(1\mid\cdots\mid1\mid p^{r_\nu}\mid\cdots\mid p^{r_\nu})
\]
puede representarse como mínimo común múltiplo de
\[
 \mathfrak C\sim(1\mid\cdots\mid1\mid p^{s_\nu}\mid\cdots\mid p^{s_i})
 \quad\text{y}\quad
 \mathfrak D\sim(1\mid\cdots\mid1\mid p^{t_\nu}\mid\cdots\mid p^{t_i}).
\]
Entonces, en los sistemas de divisores elementales de $\mathfrak C$ y $\mathfrak D$, el número $1$ debe estar en los primeros $(\nu-1)$ lugares; en el lugar $\nu$-ésimo, uno de los exponentes $s_\nu$ o $t_\nu$, digamos $s_\nu$, debe hacerse igual a $r_\nu$. Pero como
\[
 r_\nu=s_\nu\le s_{\nu+1}\le\cdots\le s_i\le r_\nu,
\]
se obtiene $\mathfrak C=\mathfrak B_\nu$, y por tanto $\mathfrak B_\nu$ es irreducible.\footnote{Las $\mathfrak B$ siguen siendo reducibles en clases uniláteras; aquí ya no hay una relación única entre divisores elementales y clase, y por tanto no hay descomposición única en clases uniláteras irreducibles. El siguiente ejemplo, que me fue comunicado por H. Brandt, muestra esto para la descomposición en clases derechas (donde las clases se representan por una matriz de base, o por el módulo correspondiente):
\[
\begin{gathered}
 \mathfrak B=[\mathfrak C_1,\mathfrak C_2]=[\mathfrak D_1,\mathfrak D_2],\quad
 \mathfrak B\sim\begin{pmatrix}p&0\\0&p\end{pmatrix},\quad
 \mathfrak C_1\sim\begin{pmatrix}1&0\\0&p\end{pmatrix},\quad
 \mathfrak C_2\sim\begin{pmatrix}p&0\\0&1\end{pmatrix},\\
 \mathfrak D_1\sim\begin{pmatrix}1&0\\0&p\end{pmatrix},
 \quad
 \mathfrak D_2\sim\begin{pmatrix}p&0\\(p-1)&1\end{pmatrix}.
\end{gathered}
\]
En efecto, los módulos $(\xi,p\eta)$, $(p\xi,\eta)$ y $(p\xi,(p-1)\xi+\eta)$ tienen cada uno como único divisor propio el módulo $(\xi,\eta)$, y son por tanto irreducibles y mutuamente distintos.} Lo mismo vale para $\mathfrak B_{\rho+1}$, donde $p$ debe reemplazarse por $0$. Pero, como muestra la formación del mínimo común múltiplo, cada una de estas clases irreducibles da un exponente determinado y el lugar en que este exponente aparece por primera vez en el sistema de divisores elementales de $\mathfrak Q$, respectivamente de $\mathfrak A$, mientras que $\mathfrak B_{\rho+1}$ da el rango. Puesto que estos números están determinados de modo único por el sistema de divisores elementales de $\mathfrak Q$, respectivamente de $\mathfrak A$, y puesto que la relación entre divisores elementales y clase es biunívoca, la descomposición de $\mathfrak Q$, y asimismo de cualquier clase, en clases irreducibles es única.

En resumen: \emph{Toda clase bilátera $\mathfrak A$ de matrices enteras con número acotado de entradas puede representarse de modo único como mínimo común múltiplo de un número finito de clases biláteras irreducibles. Cada clase irreducible representa un divisor primo fijo del sistema de divisores elementales de $\mathfrak A$, un exponente asociado y el lugar en que este exponente aparece por primera vez. La clase irreducible correspondiente al divisor $0$ da el rango de $\mathfrak A$.}

\begin{flushleft}
Erlangen, octubre de 1920.
\end{flushleft}
\begin{center}
(Recibido el 16. 10. 1920.)
\end{center}


\clearpage
\setcounter{footnote}{0}
\section*{20. Un criterio algebraico para la irreducibilidad absoluta}
\begin{center}
Math. Ann. 85 (1922), pp. 26--33
\end{center}

Un polinomio -- con coeficientes en un cuerpo cualquiera, definido abstractamente -- se llama \emph{absolutamente irreducible} cuando sigue siendo irreducible en el cuerpo algebraicamente cerrado al que puede ampliarse el dominio de coeficientes.\footnote{E. Steinitz mostró, en su \emph{Algebraische Theorie der Körper} (J. f. M. 137 (1910), p.~167), que todo cuerpo puede ampliarse, esencialmente de manera única, a uno algebraicamente cerrado, es decir, a uno en el que todo polinomio de una variable se descompone en factores lineales; con ello dio el equivalente racional del teorema fundamental del álgebra.} Muestro en lo que sigue que la irreducibilidad absoluta, a diferencia de la irreducibilidad con respecto a un cuerpo previamente dado, puede aprehenderse algebraicamente. Más precisamente, vale el teorema siguiente:

\emph{A todo polinomio en $n\ge2$ variables con coeficientes indeterminados puede asignársele una función entera racional, con coeficientes enteros, de estos coeficientes y de otras indeterminadas -- la forma de reducibilidad -- de tal modo que, para todo polinomio especial del mismo grado, la condición necesaria y suficiente para la irreducibilidad absoluta esté dada por el no anulamiento de la forma de reducibilidad. En otras palabras: para todo sistema especial de valores de los coeficientes para el cual la forma de reducibilidad se anula idénticamente en las indeterminadas y que no produzca una disminución de grado, el polinomio especial es reducible, y recíprocamente.}

Si en lugar de polinomios se consideran formas homogéneas en una variable más, la condición sobre el grado se sustituye por la más simple de que el sistema de coeficientes no se anule idénticamente.

Como consecuencia directa del teorema se obtiene además el siguiente resultado, formulado primero por A. Ostrowski:\footnote{Zur arithmetischen Theorie der algebraischen Größen. Gött. Nachr. 1919, p.~279 (lema, p.~296). Para el caso de cuerpos compuestos por números ordinarios, Ostrowski, según me consta, llegó también al teorema principal anterior; publicará próximamente sus investigaciones al respecto. Que, en la disposición de mi demostración, el teorema principal valga para cuerpos arbitrarios definidos abstractamente se debe a que la forma de reducibilidad construida para indeterminadas tiene coeficientes numéricos independientes del cuerpo particular tomado como cuerpo base.}

\emph{Todo polinomio absolutamente irreducible con números algebraicos como coeficientes permanece irreducible módulo todo ideal primo de un cuerpo algebraico finito $\Omega$ que contenga el dominio de coeficientes, salvo a lo sumo para un número finito de ideales primos de $\Omega$. En particular, $\Omega$ puede ser también el cuerpo de los números racionales.}

Pienso tratar en otro lugar otras aplicaciones a la teoría de las funciones relativamente enteras.\footnote{D. Hilbert, Mathematische Probleme. Gött. Nachr. 1900, p.~253, problema 14.}

1. Mostremos primero que \emph{todo polinomio en $n\ge2$ variables con coeficientes indeterminados es absolutamente irreducible}. En efecto, pongamos
\begin{equation}
\begin{aligned}
F(x)&=\sum A_{i_1\ldots i_n}x_1^{i_1}\cdots x_n^{i_n} &&(i_1+\cdots+i_n\le l),\\
G(x)&=\sum B_{j_1\ldots j_n}x_1^{j_1}\cdots x_n^{j_n} &&(j_1+\cdots+j_n\le m),\\
H(x)&=F(x)G(x)=\sum C_{k_1\ldots k_n}(A,B)x_1^{k_1}\cdots x_n^{k_n} &&(k_1+\cdots+k_n\le l+m),
\end{aligned}
\tag{1}
\end{equation}
donde $A$ y $B$ designan indeterminadas, y los $C(A,B)$ son combinaciones bilineales enteras de los $A$ y $B$. Por tanto, los cocientes
\[
D=C(A,B)/C_{0\ldots0}(A,B)
\]
son funciones racionales de los cocientes $A/A_{0\ldots0}$ y $B/B_{0\ldots0}$; pero el número de estas funciones es mayor que el número de sus argumentos, pues esos números son, respectivamente,
\[
\binom{l+m+n}{n}-1,
\qquad
\binom{l+n}{n}-1,
\qquad
\binom{m+n}{n}-1,
\]
y se tiene
\[
\binom{l+m+n}{n}+1>
\binom{l+n}{n}+\binom{m+n}{n}
\quad\text{para } n\ge2,\ l>0,\ m>0.
\]
Así, existe al menos \emph{una} relación algebraica entre los $D(A,B)$ y, por consiguiente, también entre los $C(A,B)$:
\begin{equation}
 \Phi(Z)\ne0\quad[\text{idénticamente en }Z_{k_1\ldots k_n}],
 \qquad
 \Phi(C(A,B))=0\quad[\text{idénticamente en }A,B],
\tag{2}
\end{equation}
donde $\Phi$ denota un polinomio entero.

Por tanto, un polinomio con indeterminadas $Z$ como coeficientes,
\begin{equation}
 E(x)=\sum Z_{k_1\ldots k_n}x_1^{k_1}\cdots x_n^{k_n}
 \qquad(k_1+\cdots+k_n\le l+m),
\tag{3}
\end{equation}
es necesariamente absolutamente irreducible para $n\ge2$.\footnote{Pues, en la terminología de 2, las indeterminadas $Z$ no son ceros de $\mathfrak P$.}

2. La totalidad de estos polinomios $\Phi(Z)$ que se anulan idénticamente en $A,B$ bajo la sustitución $Z=C(A,B)$ forma un \emph{ideal de polinomios},\footnote{A tales totalidades se las suele llamar módulos; pero, de hecho, su propiedad característica -- que junto con $\Phi_1$ y $\Phi_2$ contienen también $\Phi_1-\Phi_2$, y junto con $\Phi$ también $\Psi\Phi$, donde $\Psi$ es un polinomio entero arbitrario en $Z$ -- es la propiedad ideal.} más exactamente un ideal primo $\mathfrak P$; pues, si por la sustitución se anula un producto $\Phi_1\Phi_2$, entonces se anula al menos uno de los factores. Como (2) se cumple idénticamente en $A,B$, sigue siendo válida para todo sistema especial de valores $A,B$, donde $A,B$ y, por consiguiente, también $\Gamma=C(A,B)$ designan magnitudes de algún cuerpo. \emph{Los coeficientes $\Gamma$ de todo polinomio reducible en un cuerpo de extensión satisfacen, por tanto, las condiciones $\Phi(\Gamma)=0$; son ceros de $\mathfrak P$,} o también de los finitos polinomios de base de $\mathfrak P$. Queda por demostrar el recíproco.

Para ello debe observarse que todas las relaciones entre $A,B,C$ se conservan cuando, introduciendo una nueva variable $y$, los polinomios (1) se transforman en formas homogéneas $F(x,y)$, $G(x,y)$, $H(x,y)$ de dimensiones $l,m,l+m$. Así, también ciertos coeficientes $\Gamma$ se convierten en ceros de $\mathfrak P$ cuando, por ejemplo, $F(x,y)$\footnote{Los polinomios y formas que se obtienen sustituyendo en $F,G,\ldots,f,g,\ldots$ las indeterminadas por sistemas especiales de valores se denotarán en todo momento con una barra: $\bar F,\bar G,\ldots,\bar f,\bar g,\ldots$.} se convierte en una potencia de $y$; para ellos, el paso a polinomios no homogéneos ocasiona una disminución del grado de $\bar H(x)$ y no una factorización. Se mostrará que, con excepción de esta disminución de grado, el recíproco siempre vale; en particular, que para \emph{formas homogéneas el recíproco vale sin excepción, siempre que no se anulen todos los $\Gamma$; en otras palabras, todo cero de $\mathfrak P$ es suministrado por la representación paramétrica $\Gamma=C(A,B)$ con magnitudes $A,B$ de un cuerpo de extensión}.\footnote{Que ``en general'' los ceros de $\mathfrak P$ son suministrados por la representación paramétrica se sigue de la propiedad del ideal primo $\mathfrak P$ de definir una variedad algebraica irreducible. Pero, como es sabido, de esto no se sigue -- cosa que yo había pasado por alto originalmente y sobre la que A. Ostrowski llamó mi atención hace ya tiempo -- que la representación paramétrica no pueda fallar para ningún sistema especial de valores. La demostración que se da aquí se apoya en nociones más elementales.}

Llevo a cabo esta demostración por construcción directa de un número finito de funciones $\Phi(Z)$: mediante la sustitución de Kronecker
\[
 x_i=\xi^{d^{\,n-i}}
\]
asigno a los polinomios (1) polinomios de \emph{una} variable; muestro que aquí aparecen lagunas en los exponentes, y, formando la norma de los coeficientes correspondientes, llego a las funciones $\Phi(Z)$ y al criterio buscado.

3. Pongamos:\footnote{Festschrift, p.~11.}
\begin{equation}
 d=l+m+1,\qquad
 i=i_1d^{n-1}+i_2d^{n-2}+\cdots+i_n,
\quad
 j=j_1d^{n-1}+\cdots+j_n,
\quad
 k=k_1d^{n-1}+\cdots+k_n.
\tag{4}
\end{equation}
Entonces, por la sustitución de Kronecker, a los polinomios $F,G,H$ de (1) les corresponden, respectivamente, los polinomios
\begin{equation}
\begin{aligned}
 f(\xi)&=\sum A_{i_1\ldots i_n}\xi^i &&(i_1+\cdots+i_n\le l),\\
 g(\xi)&=\sum B_{j_1\ldots j_n}\xi^j &&(j_1+\cdots+j_n\le m),\\
 h(\xi)&=\sum C_{k_1\ldots k_n}(A,B)\xi^k &&(k_1+\cdots+k_n\le l+m).
\end{aligned}
\tag{5}
\end{equation}
Esta correspondencia es, como se sabe, biunívoca, ya que para todos los exponentes inducidos que aparecen en (5), de $k=0$ se sigue $k_1=0,\ldots,k_n=0$; y, por tanto, de $i+j=i'+j'$ se sigue también $i_1+j_1=i'_1+j'_1;\ldots;i_n+j_n=i'_n+j'_n$. Así, productos distintos $\xi^i\xi^j$ sólo pueden producir el mismo exponente $\xi^k$ cuando los productos correspondientes $x_1^{i_1}\cdots x_n^{i_n}x_1^{j_1}\cdots x_n^{j_n}$ conducen también al mismo exponente. Por tanto se obtiene
\begin{equation}
 h(\xi)=f(\xi)g(\xi),
\tag{6}
\end{equation}
y, de la validez de una relación (6), de tal modo que en $f(\xi)$ y $g(\xi)$ no aparezcan exponentes distintos de los inducidos, se recupera a la inversa
\begin{equation}
 H(x)=F(x)G(x).
\tag{7}
\end{equation}
Aquí \emph{aparecen efectivamente lagunas en los exponentes inducidos de $f(\xi)$ y $g(\xi)$}. El exponente más alto de $f(\xi)$ es $ld^{n-1}$, mientras que el segundo más alto es $(l-1)d^{n-1}+d^{n-2}$; su diferencia es $d^{n-2}(d-1)>1$, pues $d=l+m+1\ge3$ y $n\ge2$; y lo mismo vale para $g(\xi)$.

Las relaciones (6) y (7), válidas para las indeterminadas $A,B$, permanecen válidas para todo sistema especial de valores $A,B$. Para conservar los grados, homogenizo (6) y (7) mediante las variables $\eta$ e $y$. \emph{Una forma homogénea $H(x,y)$ se descompone, por tanto, en un cuerpo de extensión algebraico si y sólo si, en dicho cuerpo, la forma binaria correspondiente $h(\xi,\eta)$ puede descomponerse en dos factores de tal modo que en los factores no aparezcan exponentes distintos de los inducidos.}

4. Para efectuar la formación de la norma mencionada al final de 2, sean
\[
 a_1,b_1;\ldots; a_p,b_p;\quad a_{p+1},b_{p+1};\ldots; a_{p+q},b_{p+q}
\]
nuevas indeterminadas. Pongamos
\begin{equation}
\begin{aligned}
 \varphi(\xi,\eta)&=(a_1\xi+b_1\eta)\cdots(a_p\xi+b_p\eta)
     =r_p\xi^p+r_{p-1}\xi^{p-1}\eta+\cdots+r_0\eta^p,\\
 \psi(\xi,\eta)&=(a_{p+1}\xi+b_{p+1}\eta)\cdots(a_{p+q}\xi+b_{p+q}\eta)
     =s_q\xi^q+s_{q-1}\xi^{q-1}\eta+\cdots+s_0\eta^q,\\
 \chi(\xi,\eta)&=\varphi(\xi,\eta)\psi(\xi,\eta)
     =t_{p+q}\xi^{p+q}+t_{p+q-1}\xi^{p+q-1}\eta+\cdots+t_0\eta^{p+q},
\end{aligned}
\tag{8}
\end{equation}
de modo que $r,s,t$ denotan las funciones simétricas elementales homogenizadas; es decir, las funciones simétricas de las filas $a,b$ que son lineales en cada fila individual y a partir de las cuales todas las funciones simétricas enteras, homogéneas y de la misma dimensión en cada fila individual, se componen enteramente, con coeficientes enteros, y de una sola manera.

Formo ahora, con otras indeterminadas $u_{\mu\nu}$, la forma lineal
\begin{equation}
 \zeta(u,a,b)=\sum u_{\mu\nu}r_\mu s_\nu,
\tag{9}
\end{equation}
donde la suma se extiende sobre índices prescritos $\mu$ y $\nu$. El producto de $\zeta(u)$ por todos sus conjugados producidos por permutación de la fila $a,b$, distintos de $\zeta(u)$ y distintos entre sí -- la norma $N(\zeta(u))$ de $\zeta(u)$ respecto del cuerpo de los $t(a,b)$, cuando $\zeta(u)$ se considera como función del cuerpo de los $r_\mu(a,b)$ y $s_\nu(a,b)$ -- es entonces una función simétrica entera de todas las filas $a,b$, homogénea y de la misma dimensión en cada fila. Por lo dicho arriba, es por tanto una función entera con coeficientes enteros, determinada unívocamente, de los $t(a,b)$ y $u_{\mu\nu}$:
\begin{equation}
 N(\zeta(u,a,b))=T(t(a,b),u)\qquad[\text{idénticamente en }a,b,u].
\tag{10}
\end{equation}
El mismo razonamiento muestra que también
\[
N(z-\zeta(u))=z^\delta+T_{\delta-1}(t,u)z^{\delta-1}+\cdots+T(t,u)
\qquad[\text{idénticamente en }a,b,u,z],
\]
donde todos los $T$ son funciones enteras, con coeficientes enteros, de $t$ y $u$. Ambos lados de esta identidad se anulan para $z=\zeta(u)$; y, en efecto, por la independencia algebraica de las funciones simétricas elementales $r(a,b)$ y $s(a,b)$, se anulan idénticamente en $r,s$ cuando $t(a,b)$ se pone, de acuerdo con (8), igual a una combinación bilineal entera $w(r,s)$ de los $r$ y $s$, y $\zeta(u)$ se considera como función de $r$ y $s$. Así resulta
\begin{equation}
\zeta(u)^\delta+T_{\delta-1}(w(r,s),u)\zeta(u)^{\delta-1}+\cdots+T(w(r,s),u)=0
\quad[\text{idénticamente en }r,s,u].
\tag{11}
\end{equation}\footnote{Si la suma de (9) se extiende a todos los índices, (11) representa la extensión kroneckeriana del teorema de Gauss, esencialmente en la disposición de la demostración de Kronecker (Zur Theorie der Formen höherer Stufen, Berliner Ber. 1883, Werke, 2, p.~417).}

Las identidades (10) y (11) permanecen válidas para todos los sistemas especiales de valores $\alpha,\beta$ de $a,b$, respectivamente $\varrho,\sigma$ de $r,s$. Si $\varrho,\sigma$ se eligen en particular de modo que se anulen todos los productos $\varrho_\mu\sigma_\nu$ -- donde $\mu,\nu$ denotan los índices que aparecen en (9) --, de modo que $\zeta(u)$ se anule bajo la sustitución, entonces, poniendo $w(\varrho,\sigma)=\tau$, (11) da
\begin{equation}
 T(\tau,u)=0\qquad[\text{idénticamente en }u].
\tag{12}
\end{equation}

Recíprocamente, sean $\tau_0,\ldots,\tau_{p+q}$ sistemas de valores no todos nulos para los cuales se satisfaga (12). Como en un cuerpo de extensión algebraico existe una descomposición
\begin{equation}
 \bar\chi(\xi,\eta)=\tau_{p+q}\xi^{p+q}+\cdots+\tau_0\eta^{p+q}
=(\alpha_1\xi+\beta_1\eta)\cdots(\alpha_{p+q}\xi+\beta_{p+q}\eta),
\tag{12'}
\end{equation}
en la que $\alpha$ y $\beta$ no se anulan simultáneamente en ningún factor,\footnote{Este ``teorema fundamental del álgebra'' racional es el teorema de existencia sobre el que se apoya la validez sin excepción del recíproco enunciado al final de 2.} aunque algunos $\alpha$ o $\beta$ puedan anularse, se tiene $\tau=t(\alpha,\beta)$. Sustituyendo estos valores en (10), y usando (12), se ve que $\zeta(u,\alpha,\beta)$ o uno de sus factores conjugados se anula idénticamente en $u$. Por tanto, después de una numeración adecuada de $\alpha,\beta$, poner $r(\alpha,\beta)=\varrho$ y $s(\alpha,\beta)=\sigma$ significa precisamente que $\bar\chi(\xi,\eta)$ admite al menos una descomposición
\begin{equation}
\bar\chi(\xi,\eta)=\bar\varphi(\xi,\eta)\bar\psi(\xi,\eta)
=\sum \varrho_\mu\xi^\mu\eta^{p-\mu}\cdot\sum \sigma_\nu\xi^\nu\eta^{q-\nu},
\tag{13}
\end{equation}
de tal manera que se anulen todos los productos $\varrho_\mu\sigma_\nu$ que aparecen en $\zeta(u)$, pero no todos los $\varrho$ ni todos los $\sigma$.

5. Los grados $p,q$ de (8) se toman ahora iguales a $ld^{n-1}$ y $md^{n-1}$, es decir, a los grados de $f(\xi)$ y $g(\xi)$ en (5). Los índices $\mu,\nu$ se descomponen en $\mu^{(1)},\nu^{(1)},\mu^{(2)},\nu^{(2)}$: aquí $\mu^{(1)}$ recorre todos los exponentes que faltan en $f(\xi)$, $\nu^{(1)}$ todos los exponentes de $\xi$ en $\psi(\xi,\eta)$; correspondientemente, $\nu^{(2)}$ recorre todos los exponentes que faltan en $g(\xi)$, y $\mu^{(2)}$ todos los exponentes de $\xi$ en $\varphi(\xi,\eta)$. Además, sea
\[
 e(\xi)=\sum Z_{k_1\ldots k_n}\xi^k\qquad(k_1+\cdots+k_n\le l+m)
\]
el polinomio de una variable que corresponde al polinomio (3), y sea
\[
 S(Z,u)
\]
el polinomio entero en $Z$ y $u$ que se obtiene de $T(t,u)$ -- el miembro derecho, unívocamente determinado, de (10), donde la suma de (9) se extiende a los índices indicados $\mu^{(1)},\nu^{(1)},\mu^{(2)},\nu^{(2)}$ -- reemplazando los $t$ por los coeficientes $Z$ y por los ceros correspondientes de $e(\xi)$.

Si ahora $r,s$ se reemplazan por los coeficientes $A,B$ y por los ceros correspondientes de $f(\xi)$ y $g(\xi)$, entonces por esta sustitución $\zeta(u)$ se anula para la suma indicada. Además, $t=w(r,s)$ pasa a los coeficientes $C(A,B)$ y a los ceros correspondientes de $h(\xi)$; por tanto $T(t,u)$ pasa a $S(C(A,B),u)$. Por la identidad (11) obtenemos así
\begin{equation}
 S(C(A,B),u)=0\qquad[\text{idénticamente en }A,B,u].
\tag{14}
\end{equation}
Como (14) se conserva para todo sistema especial de valores, los coeficientes $\Gamma=C(A,B)$ de tales $H(x)$ especiales, o más generalmente de tales $H(x,y)$, que se descomponen en un cuerpo de extensión en dos factores, satisfacen la condición
\begin{equation}
 S(\Gamma,u)=0\qquad[\text{idénticamente en }u].
\tag{15}
\end{equation}

Recíprocamente, supóngase que para los coeficientes $\Gamma$, no todos nulos, de un polinomio $\bar H(x)$, respectivamente de una forma $\bar H(x,y)$, se satisface la condición (15). Por lo anterior esto equivale a $T(\tau,u)=0$, donde $\tau$ denota todos los coeficientes de $\bar h(\xi,\eta)$. Por (13) existe entonces, en un cuerpo de extensión de los $\Gamma$, una descomposición
\[
 \bar h(\xi,\eta)=\bar f(\xi,\eta)\cdot\bar g(\xi,\eta),
\]
en la que, por la suma prescrita, no aparecen en $\bar f$ ni en $\bar g$ exponentes distintos de los inducidos; por consiguiente existe una relación
\[
 \bar H(x,y)=\bar F(x,y)\cdot\bar G(x,y).
\]
Si, para $y=1$, ningún factor pasa a ser de grado cero -- en particular si los $\Gamma$ no ocasionan una disminución del grado de $H(x)$ -- se sigue además
\[
 \bar H(x)=\bar F(x)\cdot\bar G(x).
\]

\emph{La condición (15) es, por tanto, necesaria y suficiente para que $\bar H(x,y)$, y, excluida la disminución de grado, también $\bar H(x)$, se descomponga en dos factores en un cuerpo de extensión.}

De aquí se sigue además que $S(Z,u)$ no puede anularse idénticamente en $Z$ y $u$, pues en 1 se mostró que para $n\ge2$ los polinomios con indeterminadas como coeficientes, y por tanto también las correspondientes formas homogéneas en una variable más, son absolutamente irreducibles. Así, si $U$ denota productos de potencias de los $u$,
\[
 S(Z,u)=\sum \Phi_\lambda(Z)U_\lambda,
\]
donde los $\Phi_\lambda(Z)$ son, por (14), polinomios pertenecientes a $\mathfrak P$. Con ello se muestra que $\mathfrak P$ no tiene otros ceros que los dados por la representación paramétrica $\Gamma=C(A,B)$, donde $A$ y $B$ pertenecen a un cuerpo de extensión algebraico de los $\Gamma$. El recíproco enunciado al final de 2 queda así demostrado.

6. \emph{La condición de irreducibilidad absoluta} puede formularse ahora directamente. Basta descomponer el grado de $E(x)$ en (3), de todas las maneras posibles, en dos sumandos positivos $l+m$, y formar para cada descomposición la forma $S(Z,u)$. El producto de estas formas,
\begin{equation}
 R(Z,u)=\prod S(Z,u),
\tag{16}
\end{equation}
constituye entonces la forma de reducibilidad de $E(x)$; pues por lo anterior, a toda descomposición de $E(x)$, respectivamente de $E(x,y)$, le corresponde el anulamiento de un factor de (16), y recíprocamente. Queda así demostrado el teorema principal formulado en la introducción y, al mismo tiempo, se muestra cómo puede establecerse efectivamente la forma de reducibilidad en un número finito de pasos.

7. De la existencia de la forma de reducibilidad (16) se sigue inmediatamente el teorema de Ostrowski mencionado en la introducción. Sea
\[
 \bar E(x)=\sum \Gamma_{k_1\ldots k_n}x_1^{k_1}\cdots x_n^{k_n}
 \qquad(k_1+\cdots+k_n\le v)
\]
un polinomio absolutamente irreducible con enteros algebraicos $\Gamma$ como coeficientes; sea $\Omega$ un cuerpo numérico algebraico finito que contiene todos los $\Gamma$, y sea $\mathfrak p$ un ideal primo de $\Omega$.

Entonces la forma de reducibilidad $R(Z,u)$, construida para el grado exacto de $\bar E(x)$, permanece distinta de cero bajo la sustitución $Z=\Gamma$; esto es,
\[
 R(\Gamma,u)=\sum P_\lambda U_\lambda\ne0\qquad[\text{idénticamente en }u].
\]
El ideal $\mathfrak r=(P_1,P_2,\ldots)$ es, por tanto, divisible sólo por un número finito de ideales primos de $\Omega$. Si $\mathfrak p$ es distinto de esos finitos ideales, entonces $R(\Gamma,u)$ no se anula idénticamente módulo $\mathfrak p$; y, como las clases residuales módulo $\mathfrak p$ forman de nuevo un cuerpo, se sigue que $\bar E(x)$ módulo $\mathfrak p$ es irreducible, más precisamente absolutamente irreducible. Lo mismo vale para $\bar E(x,y)$, de modo que tampoco ocurre disminución de grado módulo $\mathfrak p$.

Si los coeficientes de $\bar E(x)$ son números algebraicamente fraccionarios, entonces $\mathfrak p$ debe tomarse también distinto de los finitos ideales primos que dividen el denominador común $\Lambda$; módulo todos los ideales primos restantes, $E(x)$ puede reemplazarse por $\Lambda\bar E(x)$, de modo que se satisfacen las hipótesis anteriores. Con ello queda demostrado el teorema.

\begin{flushleft}
Erlangen, agosto de 1921.
\end{flushleft}
\begin{center}
(Recibido el 28. 8. 1921.)
\end{center}

\clearpage
\section*{21. Cálculo variacional formal e invariantes diferenciales}
\begin{center}
Encyklopädie d. math. Wiss. II, 3 (1922), pp. 68--71 (en: R. Weitzenböck, Differentialinvarianten)
\end{center}

\setcounter{footnote}{148}
\subsection*{28. Cálculo variacional formal e invariantes diferenciales.\footnote{Este número se debe a E. Noether.}}

La generación de todos los invariantes diferenciales y la derivación de los teoremas principales pueden llevarse a cabo de la manera más uniforme mediante los métodos formales del cálculo de variaciones. Este principio se remonta a Riemann\textsuperscript{72)} y recibió su principal desarrollo teórico-formal por obra de Lipschitz\textsuperscript{73)}.

Se parte del problema variacional perteneciente a una forma homogénea $f(dx)$ -- donde se supone
\[
\left|\frac{\partial^2 f}{\partial dx_i\,\partial dx_k}\right|\ne0
\]
--, a saber,
\begin{equation}
 \delta\int_{t_0}^{t_1} f(x')\,dt=0,
 \qquad \left(x_i'=\frac{dx_i}{dt}\right).
\tag{140}
\end{equation}
Por integración parcial se llega al sistema invariante de ecuaciones
\[
 2\psi_i=\frac{d}{dt}\frac{\partial f}{\partial x_i'}-\frac{\partial f}{\partial x_i}=0,
\]
donde las expresiones de Lagrange $\psi_i$ son contragredientes a los $dx$. Formalmente esto se ve de la manera más simple si las $\psi_i$ se definen por la identidad formal sin integral que corresponde a la integración parcial (ecuación central de Lagrange\footnote{Para $f$ especial, la denominación se debe a Heun; véase su artículo IV, 1 11, p. 447.})
\[
 \delta f-df_\delta=-2\sum\psi_i(d,d)\,\delta x_i,
 \qquad
 \left(f_\delta=\sum\frac{\partial f}{\partial dx_i}\,\delta x_i\right),
\]
donde tanto $f$ como $df_\delta$ son invariantes, y por consiguiente también lo es el miembro derecho. En particular, para una forma cuadrática se obtiene
\[
 \delta\sum g_{ik}dx_i dx_k-2d\sum g_{ik}dx_i\delta x_k
 =-2\sum\psi_\mu(d,d)\,\delta x_\mu.
\]
Reemplazando aquí $d$ por el cogrediente $(d+\lambda\delta)$ y comparando las potencias de $\lambda$, se obtiene el correspondiente sistema de identidades:
\begin{equation}
\left\{
\begin{aligned}
D\sum g_{ik}dx_i dx_k-2d\sum g_{ik}dx_iDx_k
   &=-2\sum\psi_\mu(d,d)\,Dx_\mu,\\
D\sum g_{ik}dx_i\delta x_k-\delta\sum g_{ik}dx_iDx_k
   -d\sum g_{ik}\delta x_iDx_k
   &=-2\sum\psi_\mu(d,\delta)\,Dx_\mu,\\
D\sum g_{ik}\delta x_i\delta x_k-2\delta\sum g_{ik}\delta x_iDx_k
   &=-2\sum\psi_\mu(\delta,\delta)\,Dx_\mu.
\end{aligned}
\right.
\tag{141}
\end{equation}
De aquí se sigue que $\psi_\mu(d,\delta)$, y en particular también $\psi_\mu(d,d)$ y $\psi_\mu(\delta,\delta)$, son vectores contragredientes. De (141) se calcula $\psi_\mu(d,\delta)$ como
\[
 \psi_\mu(d,\delta)=\sum g_{i\mu}\,d\delta x_i+
 \sum\left[\begin{smallmatrix}ik\\ \mu\end{smallmatrix}\right]dx_i\delta x_k,
\]
y de esto se obtiene el vector cogrediente
\begin{equation}
 p^\sigma(d,\delta)=\sum g^{\sigma\mu}\psi_\mu
 =d\delta x_\sigma+
 \sum\left\{\begin{smallmatrix}ik\\ \sigma\end{smallmatrix}\right\}dx_i\delta x_k.
\tag{142}
\end{equation}
Poner $p^\nu(d,d)=0$ da evidentemente las ecuaciones de las líneas geodésicas.

El conocimiento del vector cogrediente $p^\sigma(d,\delta)$ conduce directamente a la derivada covariante (compárese n.º 19). Si, por ejemplo, $h(d,\delta)$ es una forma de las dos filas $dx$ y $\delta x$, entonces la expresión
\begin{equation}
 h^{(1)}(dx,\delta x)=\delta h-
 \sum\frac{\partial h}{\partial dx_\sigma}p^\sigma(d,\delta)-
 \sum\frac{\partial h}{\partial \delta x_\sigma}p^\sigma(\delta,\delta)
\tag{143}
\end{equation}
es ella misma un invariante, por ser diferencia de invariantes, y está formada de modo que se cancelen los segundos diferenciales. Así, $h^{(1)}(dx,dx)$ es una forma que contiene sólo primeros diferenciales: la derivada covariante de $h$. Lo correspondiente vale cuando $h$ contiene más filas $d^{(1)}x,\ldots,d^{(\lambda)}x$.\footnote{Lipschitz, J. f. Math. 72 (1870), p. 17.}

En el mismo principio se basa la formación de la forma de curvatura por Riemann y Lipschitz. Se pasa de $\delta^2f$ a la ``forma normal de la segunda variación''
\begin{equation}
 \Omega=\delta^2\sum g_{ik}dx_i dx_k
 -2d\delta\sum g_{ik}dx_i\delta x_k
 +d^2\sum g_{ik}\delta x_i\delta x_k,
\tag{144}
\end{equation}
que, como agregado de invariantes, es ella misma de nuevo un invariante, y en la cual, generalizando ahora la ecuación central de Lagrange, se han eliminado los terceros diferenciales. La eliminación de los segundos diferenciales se efectúa, en analogía con la diferenciación covariante, formando la diferencia
\begin{equation}
 K=\Omega-2\sum\{p^\sigma(d,\delta)\psi_\sigma(d,\delta)-p^\sigma(\delta,\delta)\psi_\sigma(d,d)\},
\tag{145}
\end{equation}
que también puede escribirse como\footnote{Lipschitz, J. f. Math. 82 (1876), § 1.}
\begin{equation}
\begin{aligned}
 \frac12 K={}&d\sum\psi_\sigma(d,\delta)\delta x_\sigma
 -\delta\sum\psi_\sigma(d,d)\delta x_\sigma\\
 &-\sum\{p^\sigma(d,\delta)\psi_\sigma(d,\delta)-p^\sigma(\delta,\delta)\psi_\sigma(d,d)\}.
\end{aligned}
\tag{146}
\end{equation}
La ley de formación de $K$ puede expresarse también diciendo que en $\Omega$ los segundos diferenciales se eliminan poniendo $p(d,\delta)$ y $p(\delta,\delta)$ -- respectivamente los miembros derechos de (141) -- iguales a cero. Ésta es la definición de Riemann de la forma de curvatura (compárese n.º 21), pero debe notarse expresamente -- como (146) lo deja especialmente claro -- que esta puesta a cero sólo puede tener lugar después de haber efectuado la diferenciación; el modo formal de formación de Lipschitz evita esta dificultad. Correspondientemente, la derivada covariante $h^{(1)}$ (143) podría definirse también eliminando en $\delta h$ los segundos diferenciales por la puesta a cero de $p(d,\delta),\ldots$.

De la forma de curvatura se obtiene, por formación sucesiva de derivadas covariantes, una sucesión de formas $[\Phi_4,\Phi_5,\ldots$ en Christoffel], cuyos invariantes proyectivos, tras adjuntar $f$, agotan según Christoffel y Ricci la totalidad de los invariantes diferenciales de la forma cuadrática, y cuya equivalencia bajo transformaciones lineales, sin ninguna otra condición de integrabilidad, expresa la equivalencia de dos formas diferenciales cuadráticas (compárese n.º 22). La razón interna de este teorema de reducción descansa en la existencia de las coordenadas normales riemannianas que surgen del problema variacional (140) (compárese n.º 20): éstas transforman las líneas geodésicas que pasan por un punto en rectas y, en consecuencia, se transforman linealmente bajo transformaciones arbitrarias de los $x$. La formación de todos los invariantes y el problema de equivalencia quedan por ello reducidos a la cuestión correspondiente para los componentes homogéneos individuales en el desarrollo de $f$ en coordenadas normales, y se muestra que estos componentes se componen linealmente de $f,\Phi_4,\Phi_5,\ldots$. Para formas cuadráticas esto fue llevado a cabo en detalle por Vermeil\textsuperscript{92)}, siguiendo una nota de E. Noether\textsuperscript{93)}. Según esa nota, los teoremas y métodos permanecen válidos cuando se toma como base una expresión diferencial arbitraria (compárense los n.os 22 y 23); así se hace patente el alcance de los métodos del cálculo variacional formal frente a los de la pura teoría de eliminación.\footnote{Compárense estas observaciones, y para otras interpretaciones geométricas, las lecciones de seminario de F. Klein indicadas bajo ``Literatura''.}

Levi-Civita\textsuperscript{134)} interpretó geométricamente la puesta $p(d,\delta)=0$ como ``desplazamiento paralelo de un vector'' (compárese también Hessenberg\textsuperscript{133)}). Este concepto, formulado axiomáticamente por Weyl, está esencialmente restringido a formas cuadráticas, pero admite una generalización de otro tipo. Sigue siendo válido cuando se amplía el grupo (multiplicación de los $g_{ik}$ por una función arbitraria), tan pronto como, además de la forma cuadrática, se toma también una forma lineal como base. Entonces $p(d,\delta)$ queda determinado unívocamente, pueden realizarse construcciones invariantes correspondientes a las anteriores, y pueden definirse coordenadas geodésicas; pero $p(d,d)=0$ ya no procede ahora de un problema variacional.\footnote{H. Weyl, Reine Infinitesimalgeometrie, Math. Ztschr. 2 (1918), p. 384--411, y Raum, Zeit, Materie (3.ª y 4.ª ed.).} Las cuestiones sobre el teorema de reducción y la equivalencia todavía no han sido abordadas aquí en general; sólo para invariantes de orden mínimo ha dado R. Weitzenböck un teorema de reducción.\footnote{Wien. Ber. 129 (1920), p. 683 y 697.}

Finalmente, conviene llamar la atención sobre los ``problemas variacionales invariantes'' generales, es decir, aquellos en los que la integral simple o múltiple $J$ es invariante con respecto a algún grupo en el sentido de Lie. Casos especiales de interés físico son tratados en trabajos de Hamel\footnote{Math. Ann. 59 (1904), p. 416--434; Ztschr. Math. Phys. 50 (1904), p. 1--54.}, Herglotz\footnote{Ann. d. Phys. (4) 36 (1911), p. 511, esp. § 9.}, Lorentz\footnote{Verslag der Amsterd. Akad., abril, sept. 1916, oct. 1916, mayo 1917.}, Fokker\footnote{Ibid., enero 1917.}, Weyl\footnote{Ann. d. Phys. (4) 54 (1917), p. 117, § 2.} y Klein\footnote{Gött. Nachr. 1917, p. 469--482; ibid. 1918, p. 171--189.}. La formulación de principio dada por E. Noether\footnote{Gött. Nachr. 1918, p. 235--257.} muestra que a la invariancia de $J$ bajo un $G_\rho$ (grupo de parámetros finitos con $\rho$ parámetros esenciales) corresponden $\rho$ divergencias linealmente independientes; a la invariancia bajo un grupo infinito que contiene $\rho$ funciones arbitrarias hasta la derivada $\sigma$-ésima corresponden $\rho$ dependencias entre las expresiones de Lagrange y sus derivadas hasta el orden $\sigma$. En ambos casos vale el recíproco. Puesto que las expresiones de Lagrange se convierten en invariantes relativos del grupo, se tiene al mismo tiempo un proceso generador de invariantes.

\begin{center}
(Concluido en marzo de 1921.)
\end{center}

\clearpage
\providecommand{\Amod}{\mathfrak A}
\providecommand{\Bmod}{\mathfrak B}
\providecommand{\Gmod}{\mathfrak G}
\providecommand{\Mmod}{\mathfrak M}
\providecommand{\Nmod}{\mathfrak N}
\providecommand{\Rmod}{\mathfrak R}
\providecommand{\Smod}{\mathfrak S}
\providecommand{\mideal}{\mathfrak m}
\providecommand{\nideal}{\mathfrak n}
\providecommand{\gideal}{\mathfrak g}
\providecommand{\hideal}{\mathfrak h}
\providecommand{\jideal}{\mathfrak j}
\providecommand{\tideal}{\mathfrak t}
\providecommand{\aideal}{\mathfrak a}
\providecommand{\bideal}{\mathfrak b}
\providecommand{\bb}{\mathfrak b}
\providecommand{\cc}{\mathfrak c}
\providecommand{\zero}[1]{\equiv 0\pmod{#1}}
\providecommand{\Norm}{N}

\setcounter{footnote}{0}

\section*{22. Reelaboración de un trabajo de K. Hentzelt:\\ sobre la teoría de los ideales de polinomios y las resultantes}
\begin{center}
Math. Ann. 88 (1923), pp. 53--79
\end{center}

En lo que sigue se trata del \emph{problema general de la teoría de la eliminación, a saber, determinar los ceros comunes de todos los polinomios de un ideal de polinomios}; o, lo que es idéntico, los ceros de cualquier número finito de polinomios que formen una base del ideal. Al mismo tiempo se obtendrá una interpretación conceptual de las multiplicidades que aparecen. Los coeficientes de los polinomios pueden pertenecer a un cuerpo cualquiera; los ceros pertenecen entonces al cuerpo algebraicamente cerrado derivado de éste. Además, sin restricción de generalidad, el ideal puede suponerse transformado (§~3). Entonces el resultado esencial es el siguiente:

\emph{A cada ideal $\mideal$ de polinomios se le puede asociar una forma resultante:}
\[
 R_{\mideal}=R^{(1)}(x_1,\ldots,x_n)\cdots
 R^{(i)}(x_i,\ldots,x_n)\cdots R^{(n)}(x_n)\zero{\mideal},
\]
\emph{de tal modo que $R_{\mideal}$ se anula para todos los ceros de $\mideal$, y sólo para ellos. Si $\nideal$ es un divisor de $\mideal$, y si las formas resultantes $R_{\nideal}$ y $R_{\mideal}$ coinciden, entonces también coinciden los ideales $\nideal$ y $\mideal$.}\footnote{Kurt Hentzelt está desaparecido ante Dixmuiden desde octubre de 1914 y debe contarse entre los muertos. El presente artículo es una reelaboración completamente libre de la parte más esencial de su disertación ``Zur Theorie der Formenmoduln und Resultanten'', con la cual se doctoró en Erlangen bajo E. Fischer en el verano de 1914. Esa disertación, compuesta enteramente sobre la base de ideas propias, está construida sin lagunas; pero un lema sigue a otro lema, todos los conceptos están circunscritos por fórmulas con cuatro y cinco índices, y el texto falta casi por completo, de modo que la comprensión encuentra las mayores dificultades. Él mismo no llegó a realizar la reelaboración prevista. Presento aquí el trabajo en una formulación puramente conceptual, con lo cual se obtiene una gran simplificación de las demostraciones, cuyos pensamientos fundamentales proceden en todas partes de Hentzelt, y, según espero, se hace visible la belleza del trabajo. -- Las partes de la disertación que, para una base dada, se refieren a la cuestión de formar en finitos pasos las funciones que aparecen, quedarán reservadas para una publicación separada. (E. N.)}

La segunda parte del teorema principal muestra que la forma resultante proporciona los ceros con multiplicidad característica. Esta segunda parte es análoga al hecho de que dos ideales de un cuerpo algebraico de números, de los cuales uno es divisor del otro, coinciden si y sólo si coinciden sus \emph{normas}; también la razón interna de ambos teoremas es la misma.

En efecto, $\mideal$ puede concebirse como un módulo de formas lineales $\mathfrak M_{i-1}$, considerando cada polinomio como una forma lineal en los productos de potencias de $x_1,\ldots,x_{i-1}$ y adjuntando $x_{i+1},\ldots,x_n$ al dominio de coeficientes. El factor resultante $R^{(i)}$ resulta entonces igual a la \emph{norma} del módulo fundamental correspondiente (§~1) $\Gmod_{i-1}$ respecto de $\mathfrak M_{i-1}$; de ahí se obtiene asimismo inmediatamente una \emph{interpretación de la multiplicidad -- producto del grado y del exponente de un factor de $R^{(i)}$ -- por el número de clases residuales linealmente independientes}\footnote{Estos últimos resultados muestran su verdadero significado cuando se introduce la descomposición de ideales en primarios; la multiplicidad correspondiente a un ideal primario se define entonces como el número de clases residuales linealmente independientes del complemento módulo ese ideal. Esto se tratará en otro lugar. (E. N.)}, hecho que hasta ahora sólo era conocido en el caso especial en que la resultante puede definirse para coeficientes indeterminados\footnote{Cf. Macaulay, \emph{The algebraic theory of modular systems}, Cambridge Tracts 19 (Cambridge University Press, 1916), núm. 67; también Lasker, Zur Theorie der Moduln und Ideale, Math. Ann. 60 (1905), pp. 20--116; p. 98. Que en este caso coinciden la resultante y la forma resultante se demuestra en los teoremas aún no publicados de Hentzelt mencionados en 1). (E. N.)}.

Para deducir estos resultados, §§~1 y 2 dan teoremas sobre módulos de formas lineales cuyos coeficientes son enteros racionales o polinomios en una indeterminada. La conexión con los ideales de polinomios introducidos en §~3 la proporciona el teorema de isomorfía de §~4. La sección~5 da la segunda parte, antes mencionada, del teorema principal; §~7 da el teorema sobre los ceros, al que en §~6 precede una teoría general de la eliminación. Allí se demuestra al mismo tiempo, al introducir nuevas indeterminadas, la descomposición de la resolvente en formas lineales de esas indeterminadas; y con ello se da, para la eliminación presentada aquí, una prueba irreprochable de una afirmación de Kronecker\footnote{Es sabido que no tuvo éxito el intento de demostración en König, \emph{Algebraische Größen} (Leipzig 1903, Teubner), V, §~4, para la teoría de eliminación de Kronecker. Que tampoco para las multiplicidades introducidas por König en la teoría de eliminación de Kronecker vale la segunda parte del teorema principal de Hentzelt, lo muestra Macaulay, loc. cit., p. 23; allí se muestra además que la teoría de eliminación de Kronecker no corresponde a la descomposición en ideales primarios. (E. N.)}.

\section*{§ 1. Módulos de formas lineales.}

Por magnitudes enteras $a,b,c,\ldots$ se entenderán en los dos primeros parágrafos enteros racionales o polinomios en una indeterminada con coeficientes en un cuerpo $P$ arbitrario, definido abstractamente; por formas lineales $a(\xi),b(\xi),\ldots$, tales formas en finitas indeterminadas $\xi_1,\ldots,\xi_k$ con magnitudes enteras como coeficientes.

Un \emph{módulo $\Amod$ de formas lineales} se define por las condiciones siguientes: junto con $a(\xi)$ y $b(\xi)$, $\Amod$ contiene también la diferencia $a(\xi)-b(\xi)$; junto con $a(\xi)$ contiene también $c\cdot a(\xi)$, donde $c$ designa una magnitud entera arbitraria.

Por \emph{rango de $\Amod$} entendemos, como de costumbre, el número máximo de formas lineales linealmente independientes de $\Amod$; por el $\lambda$-ésimo \emph{divisor determinantal} $d_\lambda$ de $\Amod$, el máximo común divisor de todos los determinantes de orden $\lambda$ de $\Amod$ (es decir, de todos los determinantes de orden $\lambda$ que se forman con la matriz de coeficientes de cualesquiera $\lambda$ formas lineales de $\Amod$). Si $\rho$ es el rango de $\Amod$, de modo que $d_1,\ldots,d_\rho$ son los divisores determinantales no nulos, los \emph{divisores elementales} de $\Amod$ se definen por $e_1=d_1$, $e_2=d_2/d_1,\ldots,e_\rho=d_\rho/d_{\rho-1}$, $e_{\rho+i}=0$. Es sabido que $\Amod$ es un módulo finito que posee una base modular de exactamente $\rho$ formas lineales $a_1(\xi),\ldots,a_\rho(\xi)$, mediante la cual toda forma lineal de $\Amod$ puede representarse linealmente, con magnitudes enteras como coeficientes: $\Amod=(a_1(\xi),\ldots,a_\rho(\xi))$. Si $A$ denota la matriz de coeficientes de estas formas lineales, entonces los divisores determinantales y elementales de $\Amod$ coinciden con los de $A$; de esto se sigue primero que \emph{cada divisor elemental $e_i$ divide al siguiente $e_{i+1}$}. La ecuación matricial suministrada por la teoría de divisores elementales,
\[
 P A Q=\begin{pmatrix}
 e_1&&&0\\
 &e_2&&\\
 &&\ddots&\\
 0&&&e_\rho
 \end{pmatrix},
\]
muestra además -- si mediante la transformación unimodular $\xi=Q(\eta)$ se introducen nuevas indeterminadas $\eta_1,\ldots,\eta_k$ -- que vale la ecuación de módulos
\begin{equation}
 \Amod=(e_1\eta_1,e_2\eta_2,\ldots,e_\rho\eta_\rho).
\tag{1}
\end{equation}

El concepto más esencial para lo que sigue es el de módulo fundamental\footnote{Cf. Steinitz, Rechteckige Systeme und Moduln in algebraischen Zahlkörpern II, Math. Ann. 72 (1912), pp. 297--345, núm. 36. Hentzelt parece haber llegado, en todo caso independientemente de Steinitz -- con quien también se hallan puntos de contacto en otros lugares --, para su caso más simple, al concepto en la forma de la fórmula (4). (E. N.)}, según la siguiente definición.

\medskip
\noindent\textbf{Definición I.} \emph{Un módulo $\Gmod$ se llama módulo fundamental si no posee ningún divisor propio del mismo rango}\footnote{Divisor entendido en el sentido modular: $\Amod$ es divisible por $\Bmod$, $\Amod\zero{\Bmod}$, si todo elemento de $\Amod$ está contenido en $\Bmod$; $\Bmod$ se llama divisor propio si contiene elementos distintos de los de $\Amod$.}.

Así, $\Gmod$ es módulo fundamental si y sólo si de
\begin{equation}
 c\cdot g(\xi)\zero{\Gmod};\quad c\ne0 \quad\text{se sigue siempre:}\quad g(\xi)\zero{\Gmod}.
\tag{2}
\end{equation}

La conexión entre un módulo arbitrario y un módulo fundamental viene dada por

\medskip
\noindent\textbf{Teorema I.} \emph{Todo módulo $\Amod$ es divisible por uno y sólo un módulo fundamental $\Gmod$ del mismo rango; $\Gmod$ se define como el menor divisor de $\Amod$ que tiene el mismo rango, y está representado por}
\begin{equation}
 \Gmod=(\eta_1,\ldots,\eta_\rho);
\tag{3}
\end{equation}
\emph{$\Gmod$ se llamará el módulo fundamental de $\Amod$.}

La condición de que $\Gmod$ sea el menor divisor del mismo rango se expresa así: $\Gmod$ contiene todas y sólo aquellas formas lineales $g(\xi)$ que satisfacen una relación
\begin{equation}
 b\cdot g(\xi)\zero{\Amod};\quad b\ne0.
\tag{4}
\end{equation}
Primero se sigue que $\Gmod$ es un módulo, pues junto con
\[
 b_1g_1(\xi)\zero{\Amod};\quad b_1\ne0;
 \qquad
 b_2g_2(\xi)\zero{\Amod};\quad b_2\ne0
\]
se cumple también
\[
 b_1b_2\bigl(g_1(\xi)-g_2(\xi)\bigr)\zero{\Amod};\quad b_1b_2\ne0,
\]
y, por supuesto, también $b_1\cdot c g_1(\xi)\zero{\Amod}$. Pero $\Gmod$ es asimismo un módulo fundamental; pues de
\[
 c\cdot g(\xi)\zero{\Gmod};\quad c\ne0
\]
se sigue, por (4),
\[
 b c g(\xi)\zero{\Amod};\quad bc\ne0,
\]
y de nuevo por (4), $g(\xi)\zero{\Gmod}$.

Además, no existe un módulo fundamental $\Gmod_1$ distinto del menor divisor $\Gmod$ del mismo rango que satisfaga las condiciones. En efecto, $\Gmod_1$ tendría que ser divisible por $\Gmod$; pero, como módulo fundamental, no posee divisor propio del mismo rango, y por tanto es idéntico a $\Gmod$. Finalmente, el módulo representado por el miembro derecho de (3) es un módulo fundamental, pues todo divisor propio debería contener una indeterminada $\eta$ adicional y, por ello, tendría rango mayor. Así, (3) da el módulo fundamental de $\Amod$ determinado de modo único, y \emph{el Teorema I queda demostrado en todas sus partes}.

Otro concepto esencial para lo que sigue es el concepto de Dedekind del cociente de dos módulos\footnote{En Dedekind se trata de módulos de números, para los cuales también está definida la multiplicación, de modo que el cociente de Dedekind no coincide con el cociente aquí definido. (E. N.)}:

\medskip
\noindent\textbf{Definición II.} \emph{El cociente $\cc=\Amod/\Bmod$ de dos módulos de formas lineales se define como la totalidad de las magnitudes enteras $c$ que satisfacen la condición}
\[
 c\cdot\Bmod\zero{\Amod}.
\]
\emph{$\cc$ es un ideal de magnitudes enteras, por tanto un ideal principal. Correspondientemente, el cociente $\mathfrak C=\Amod/\bb$ de un módulo $\Amod$ de formas lineales por un ideal $\bb$ de magnitudes enteras se define como la totalidad de las formas lineales $c(\xi)$ que satisfacen la condición}
\[
 c(\xi)\cdot\bb\zero{\Amod}.
\]
\emph{$\mathfrak C$ es de nuevo un módulo de formas lineales y, tan pronto como $\bb$ es distinto del ideal cero, un módulo del mismo rango que $\Amod$.}

Sólo hay que observar que, por definición, es claro que la propiedad modular pertenece en cada caso al cociente; pero los sistemas de magnitudes enteras con la propiedad modular son ideales.

El concepto de cociente conduce a una conexión entre el módulo fundamental y el divisor elemental más alto, a saber:

\medskip
\noindent\textbf{Teorema II.} \emph{Entre el módulo fundamental y el divisor elemental más alto de $\Amod$ existe la relación recíproca}
\begin{equation}
 (e_\rho)=\Amod/\Gmod;\qquad \Gmod=\Amod/(e_\rho),
\tag{5}
\end{equation}
\emph{donde $(e_\rho)$ denota el ideal principal derivado de $e_\rho$.}

Como todo divisor elemental divide a $e_\rho$, de (1) y (3) se tiene
\begin{equation}
 e_\rho\Gmod\zero{\Amod}\quad \text{y por consiguiente}\quad (e_\rho)\cdot\Gmod\zero{\Amod};
\tag{6}
\end{equation}
por tanto $(e_\rho)$ es divisible por el cociente $\Amod/\Gmod$. Pero también vale la inversa: de
\[
 b\Gmod\zero{\Amod}
 \quad\text{se sigue}\quad b\eta_\rho\zero{\Amod}
 \quad\text{y por tanto}\quad b\zero{(e_\rho)},
\]
lo que demuestra la primera fórmula (5)\footnote{Póngase $\Amod_0=\Amod,\ldots,\Amod_\lambda=(\Amod,\eta_\rho,\ldots,\eta_{\rho-\lambda+1})$, donde cada vez $\Amod_i$ tiene como divisor elemental más alto $e_{\rho-i}$. Entonces, por las mismas conclusiones, se obtiene
\[
 (e_\rho)=\Amod_0/\Amod_1,\ldots,(e_{\rho-\lambda})=\Amod_\lambda/\Amod_{\lambda+1},\ldots,(e_1)=\Amod_{\rho-1}/\Amod_\rho.
\]
Más generalmente, póngase $\zeta_\lambda=b_{\lambda1}\eta_1+\cdots+b_{\lambda\lambda}\eta_\lambda$, supóngase $(b_{\lambda\lambda},e_\lambda)=1$, y sea
\[
 \Bmod_0=\Amod,\ldots,\Bmod_\lambda=(\Amod,\zeta_\rho,\ldots,\zeta_{\rho-\lambda+1}).
\]
Entonces $\Bmod_i$ tiene también divisor elemental más alto $e_{\rho-\lambda}$, y por tanto
\[
 (e_{\rho-\lambda})=\Bmod_\lambda/\Bmod_{\lambda+1}.
\]}. La fórmula (6) muestra además que $\Gmod$ es divisible por el cociente $\Amod/(e_\rho)$; este cociente es un divisor de $\Amod$ del mismo rango, luego, por la definición del módulo fundamental de $\Amod$, es recíprocamente divisible por $\Gmod$. Con ello queda demostrada también la segunda fórmula (5).

Si $d$ denota cualquier determinante de orden $\rho$ de $\Amod$ que no se anula idénticamente, entonces $d$ es divisible por el $\rho$-ésimo divisor determinantal $d_\rho$, y por tanto por $e_\rho$; en consecuencia $d\Gmod\zero{\Amod}$ y, por lo anterior, $\Gmod=\Amod/(d)$. Para lo que sigue, sin embargo, es esencial que esta última representación por cociente de $\Gmod$ pueda demostrarse sin remontarse a los divisores elementales y tenga por ello validez más general, conforme a

\medskip
\noindent\textbf{Teorema III.} \emph{Si por magnitudes enteras se entienden polinomios en varias indeterminadas}\footnote{De hecho sólo se utiliza esto: que las magnitudes enteras se reproducen bajo suma, resta y multiplicación, con validez de las reglas usuales de cálculo, de modo que forman un anillo; y además, que en ese anillo un producto sólo se anula cuando se anula un factor, de modo que el anillo puede extenderse a un cuerpo por formación de cocientes (adjunción de pares de elementos). En cambio no se usa ninguna representación de base del módulo; así, (7) vale por ejemplo también para el dominio de todos los enteros algebraicos como coeficientes. (E. N.)}, \emph{entonces sigue siendo válida la representación por cociente}
\begin{equation}
 \Gmod=\Amod/(d),
\tag{7}
\end{equation}
\emph{donde $d$ denota cualquier determinante de orden $\rho$ de $\Amod$ que no se anula idénticamente.}

Primero es claro que los conceptos de rango, módulo fundamental y cociente de módulos -- que ahora ya no es necesariamente un ideal principal -- se conservan bajo esta ampliación, y también el Teorema I, con excepción de la representación por base.

Sean ahora
\[
 a_1(\xi)=a_{11}\xi_1+\cdots+a_{1k}\xi_k,\ldots,
 a_\rho(\xi)=a_{\rho1}\xi_1+\cdots+a_{\rho k}\xi_k
\]
$\rho$ formas lineales linealmente independientes de $\Amod$, que en general no formarán una base modular. Designemos las indeterminadas $\xi$ de modo que
\[
 d=|a_{ij}|\ne0;\qquad i,j=1,2,\ldots,\rho.
\]
Entonces pertenecen a $\Amod$ $\rho$ formas lineales de forma especial:
\begin{equation}
 d\xi_\lambda+b_{\lambda,\rho+1}\xi_{\rho+1}+\cdots+b_{\lambda,k}\xi_k\zero{\Amod}
 \quad(\lambda=1,2,\ldots,\rho).
\tag{8}
\end{equation}
Pero $\Amod$ no puede contener ninguna forma lineal que dependa sólo de $\xi_{\rho+1},\ldots,\xi_k$, digamos $c_{\rho+1}\xi_{\rho+1}+\cdots+c_k\xi_k$, pues de lo contrario al menos un determinante de orden $\rho+1$, a saber $d\cdot c_\alpha$, sería no nulo.

Ahora bien, el cociente $\Amod/(d)$, como divisor de $\Amod$ del mismo rango, es divisible por $\Gmod$; pero también vale la inversa. En efecto, todo $g(\xi)\zero{\Gmod}$ satisface, por definición,
\[
 c\cdot g(\xi)\zero{\Amod};\quad c\ne0
 \quad\text{y por consiguiente}\quad d\cdot c\cdot g(\xi)\zero{\Amod}.
\]
Por (8), sin embargo,
\[
 d\cdot g(\xi)=g_{\rho+1}\xi_{\rho+1}+\cdots+g_k\xi_k\zero{\Amod},
\]
y en consecuencia
\[
 c g_{\rho+1}\xi_{\rho+1}+\cdots+c g_k\xi_k\zero{\Amod}.
\]
Por lo que se acaba de observar, esto implica $c\cdot g_{\rho+1}=0,\ldots,c\cdot g_k=0$, y como $c\ne0$, también $g_{\rho+1}=0,\ldots,g_k=0$; por tanto $d\cdot g(\xi)\zero{\Amod}$, y (7) queda demostrado.

Debe notarse que $d\cdot g(\xi)\zero{\Amod}$ se sigue también directamente de que los dos módulos $(a_1(\xi),\ldots,a_\rho(\xi))$ y $(g(\xi),a_1(\xi),\ldots,a_\rho(\xi))$ tienen el mismo rango $\rho$, de modo que -- escribiendo $g(\xi)=c_1\xi_1+\cdots+c_k\xi_k$ -- se anula el determinante de orden $\rho+1$
\[
\begin{vmatrix}
 a_1(\xi)&a_{11}&\cdots&a_{1\rho}\\
 \vdots&\vdots&&\vdots\\
 a_\rho(\xi)&a_{\rho1}&\cdots&a_{\rho\rho}\\
 g(\xi)&c_1&\cdots&c_\rho
\end{vmatrix}
\]
La demostración dada arriba, sin embargo, reaparece en §~4 después de la fórmula (30).

\section*{§ 2. Teoremas de descomposición para normas y módulos.}

Volvemos ahora a los módulos de formas lineales respecto de enteros racionales o de polinomios en una indeterminada con coeficientes en $P$.

\medskip
\noindent\textbf{Definición III.} \emph{Si, como de costumbre, todas las formas lineales en $\xi$ congruentes entre sí módulo $\Amod$ se reúnen en una clase residual, el símbolo $\Gmod\mid\Amod$ designará el sistema de clases residuales de $\Gmod$ módulo $\Amod$, es decir, el sistema de todas las clases residuales módulo $\Amod$ que constan de elementos de $\Gmod$. La norma de $\Gmod$ respecto de $\Amod$ -- en símbolos $\Norm(\Gmod\mid\Amod)$ -- se define como el determinante de la sustitución de paso de $\Gmod$ a $\Amod$, es decir, como el determinante de la sustitución que expresa una base modular linealmente independiente cualquiera de $\Amod$ mediante tal base del módulo fundamental $\Gmod$ de $\Amod$.}

De (1) y (3), respectivamente de la observación al Teorema II, se obtiene:
\begin{equation}
 \Norm(\Gmod\mid\Amod)=e_1e_2\cdots e_\rho=d_\rho;
 \qquad
 \Gmod=\Amod/\bigl(\Norm(\Gmod\mid\Amod)\bigr).
\tag{9}
\end{equation}

De (1), (3) y (9) se sigue de la manera conocida que, en el caso de enteros racionales, $\Norm(\Gmod\mid\Amod)$ representa el número de clases residuales de $\Gmod$ módulo $\Amod$, mientras que en el caso de polinomios en una indeterminada -- donde el número de clases residuales es en general infinito -- el grado de $\Norm(\Gmod\mid\Amod)$ se vuelve igual al \emph{número de clases residuales linealmente independientes respecto de $P$}. Si $\Bmod$ es un divisor de $\Amod$ del mismo rango, entonces, junto con un representante de cualquiera de esas clases residuales, $\Bmod$ contiene también todos los elementos de la clase residual; por tanto $\Bmod$ se obtiene de $\Amod$ adjuntando finitas clases residuales, o combinaciones lineales de ellas. De aquí se sigue inmediatamente:

\medskip
\noindent\textbf{Teorema IV.} \emph{Si $\Bmod$ es un divisor de $\Amod$ del mismo rango, de modo que los módulos fundamentales coinciden, y si además $\Norm(\Gmod\mid\Bmod)=\Norm(\Gmod\mid\Amod)$, entonces también $\Bmod=\Amod$.}

Sobre la conexión entre la descomposición de normas y de módulos se tiene lo siguiente.

\medskip
\noindent\textbf{Teorema V.} \emph{Sea}
\begin{equation}
 \Norm(\Gmod\mid\Amod)=r\cdot s,\qquad (r,s)=1;
\tag{10}
\end{equation}
\emph{entonces existe uno y sólo un módulo $\Rmod$, y asimismo uno y sólo un módulo $\Smod$, ambos divisores de $\Amod$ del mismo rango, tales que}
\[
 r=\Norm(\Gmod\mid\Rmod),\qquad s=\Norm(\Gmod\mid\Smod)
\]
\emph{se cumplen. $\Rmod$ y $\Smod$ están definidos por}
\[
 \Rmod=\Amod/(s);\qquad \Smod=\Amod/(r),
\]
\emph{y $\Amod$ es igual al mínimo común múltiplo de $\Rmod$ y $\Smod$.}\footnote{Mínimo común múltiplo $[\Rmod,\Smod]$ entendido en el sentido modular: la totalidad de las formas lineales que son divisibles tanto por $\Rmod$ como por $\Smod$. Correspondientemente, el máximo común divisor $(\Rmod,\Smod)$ en el sentido modular debe entenderse como la totalidad de las formas lineales que pueden representarse como suma de un elemento de $\Rmod$ y un elemento de $\Smod$.}
\emph{Si se pone $r_\rho=(r,e_\rho)$, $s_\rho=(s,e_\rho)$, entonces valen las reciprocidades}
\begin{equation}
 (s_\rho)=\Amod/\Rmod;\quad \Rmod=\Amod/(s_\rho);
 \qquad
 (r_\rho)=\Amod/\Smod;\quad \Smod=\Amod/(r_\rho).
\tag{11}
\end{equation}

En efecto, poniendo en general $r_i=(r,e_i)$, $s_i=(s,e_i)$, se obtiene de (9) y (10):
\[
 (r_i,s_i)=1;\qquad e_i=r_is_i;\qquad r=r_1\cdots r_\rho;\qquad s=s_1\cdots s_\rho,
\]
y cada $r_i$, respectivamente $s_i$, divide al siguiente $r_{i+1}$, respectivamente $s_{i+1}$. Así, si se pone
\[
 \Rmod=(r_1\eta_1,\ldots,r_\rho\eta_\rho);\qquad
 \Smod=(s_1\eta_1,\ldots,s_\rho\eta_\rho),
\]
entonces $\Rmod$ y $\Smod$ son divisores de $\Amod$ del mismo rango; por la primalidad relativa de $r_i$ y $s_i$, $\Amod$ se convierte en el mínimo común múltiplo de $\Rmod$ y $\Smod$, y además
\[
 \Norm(\Gmod\mid\Rmod)=r_1\cdots r_\rho=r;\qquad
 \Norm(\Gmod\mid\Smod)=s_1\cdots s_\rho=s.
\]
Además
\[
 s_\rho\Rmod\zero{\Amod};\qquad r_\rho\Smod\zero{\Amod};
\]
y de $c\Rmod\zero{\Amod}$ se sigue $cr_\rho\eta_\rho\zero{\Amod}$, de donde $c\zero{(s_\rho)}$, lo que prueba $(s_\rho)=\Amod/\Rmod$; de modo análogo, $(r_\rho)=\Amod/\Smod$. De
\[
 b(\eta)=b_1\eta_1+\cdots+b_\rho\eta_\rho\zero{\Amod/(s_\rho)}
\]
se sigue además, puesto que todos los $r_i$ son relativamente primos con $s_\rho$, que $b_i\zero{(r_i)}$; por tanto $\Rmod=\Amod/(s_\rho)$ y, análogamente, $\Smod=\Amod/(r_\rho)$. Así quedan demostradas las reciprocidades (11); en el caso especial $r=1$ pasan a (5).

Queda mostrar la \emph{determinación única} de $\Rmod$ y $\Smod$. Sean $\overline{\Rmod}$ y $\overline{\Smod}$ módulos que satisfacen las condiciones del Teorema V, y sean $\bar r_i$ y $\bar s_i$ sus divisores elementales. Puesto que $\bar r_i$ y $\bar s_i$ son divisores de $e_i$, se sigue, por
\[
 r=\bar r_1\cdots\bar r_\rho,
 \qquad
 s=\bar s_1\cdots\bar s_\rho,
 \qquad
 (\bar r_i,\bar s_i)=1,
\]
que también
\[
 e_i=\bar r_i\bar s_i,
 \qquad
 \bar r_i=(r,e_i)=r_i,
 \qquad
 \bar s_i=(s,e_i)=s_i.
\]
Así $\overline{\Rmod}$ y $\overline{\Smod}$ coinciden con $\Rmod$ y $\Smod$ en divisores elementales y en módulo fundamental. Si por tanto, mediante una sustitución unimodular $\eta=Q(\xi)$, se introducen nuevas indeterminadas $\xi$, entonces
\[
 \overline{\Rmod}=(r_1\xi_1,\ldots,r_\rho\xi_\rho);
\]
\[
 \Amod=\bigl(r_1s_1(q_{11}\xi_1+\cdots+q_{1\rho}\xi_\rho),\ldots,
 r_\rho s_\rho(q_{\rho1}\xi_1+\cdots+q_{\rho\rho}\xi_\rho)\bigr).
\]
De $\Amod\zero{\Rmod}$ se obtiene por ello
\[
 r_i s_i(q_{i1}\xi_1+\cdots+q_{i\rho}\xi_\rho)\zero{\Rmod};
\]
por consiguiente $r_is_iq_{ij}\zero{(r_i)}$. Como $(r,s)=1$, esto da $r_iq_{ij}\zero{(r_i)}$, o sea $\overline{\Rmod}\zero{\Rmod}$. Pero, puesto que $\Norm(\Gmod\mid\Rmod)=\Norm(\Gmod\mid\overline{\Rmod})$, el Teorema IV da $\Rmod=\overline{\Rmod}$, y correspondientemente $\Smod=\overline{\Smod}$. El Teorema V queda así demostrado en todas sus partes.

\section*{§ 3. Ideales de polinomios.}

Sea $\overline{\mideal}$ un \emph{ideal de polinomios} en las $n$ indeterminadas $y_1,\ldots,y_n$, con coeficientes en un cuerpo $P$ arbitrario, definido abstractamente; esto es, junto con $\Phi_1(y)$ y $\Phi_2(y)$, $\overline{\mideal}$ contiene también la diferencia $\Phi_1(y)-\Phi_2(y)$, y junto con $\Phi(y)$ contiene también $A(y)\cdot\Phi(y)$, donde $A(y)$ denota un polinomio arbitrario con coeficientes en $P$\footnote{La designación conceptual que resulta, ``ideal'' en lugar de las antes generalmente usuales ``módulo'' o ``módulo de formas'', se hace ya necesaria aquí para distinguir éstos de los módulos de formas lineales. (E. N.)}.

Sométanse los $y$ a una transformación con indeterminadas como coeficientes,
\begin{equation}
 y=U(x);\quad \text{por ejemplo}\quad
 \begin{cases}
 y_1=x_1,\\
 y_2=u_{21}x_1+x_2,\\
 \cdots\\
 y_n=u_{n1}x_1+\cdots+u_{n,n-1}x_{n-1}+x_n,
 \end{cases}
\tag{12}
\end{equation}
y adjúntense las indeterminadas $u_{\mu\nu}$ al dominio de coeficientes de los polinomios, que pasa así a un cuerpo $P(u)$. Por esta adjunción, $\overline{\mideal}$ pasa a $\overline{\mideal}_{(u)}$; así, $\overline{\mideal}_{(u)}$ contiene, además de los polinomios $\Phi_\lambda(y)$ de $\overline{\mideal}$, todas las combinaciones lineales $\sum \alpha_\lambda(u)\Phi_\lambda(y)$ de finitos $\Phi_\lambda$, donde $\alpha_\lambda(u)$ denota funciones racionales de los $u_{\mu\nu}$ con coeficientes en $P$. Multiplicando por un polinomio adecuado en $u$, puede suponerse sin restricción de generalidad que los $\alpha_\lambda(u)$ son productos de potencias en $u$. Así se tiene
\begin{equation}
 \sum \alpha_\lambda(u)\Phi_\lambda(y)\zero{\overline{\mideal}_{(u)}};
 \qquad
 \Phi_\lambda(y)\zero{\overline{\mideal}}\quad\text{y recíprocamente}.
\tag{13}
\end{equation}
Mediante (12), $\overline{\mideal}_{(u)}$ pasa a un ideal $\mideal$ de polinomios en $x_1,\ldots,x_n$ con coeficientes en $P(u)$:
\begin{equation}
 \overline{\mideal}_{(u)}=\mideal\quad \text{mediante}\quad y=U(x).
\tag{14}
\end{equation}
La combinación de las relaciones (14) y (13) dice lo siguiente:
\begin{equation}
 \text{De } F(x)\zero{\mideal};\quad
 F(x)=\Phi(y)=\sum \alpha_\lambda(u)\Phi_\lambda(y)=\sum \alpha_\lambda(u)F_\lambda(x)
 \text{ se sigue } F_\lambda(x)\zero{\mideal};
\tag{15}
\end{equation}
mientras que (14) y (15), recíprocamente, vuelven a dar (13).

\medskip
\noindent\textbf{Definición IV.} \emph{Los ideales $\mideal$ de polinomios en $x_1,\ldots,x_n$ con coeficientes en $P(u)$ que satisfacen las relaciones (15), y que por tanto surgen de $\overline{\mideal}_{(u)}$ mediante $y=U(x)$, se llamarán ideales transformados.}

Los ideales transformados se distinguen por la existencia de polinomios regulares; \emph{un polinomio de grado $r$ se llama regular respecto de $x_i$ si contiene el término $x_i^r$ con coeficiente no nulo}. Se ve que los polinomios $F_i(x)$ son regulares respecto de $x_1$. En lo que sigue, el punto principal será considerar estos ideales transformados como módulos de formas lineales en ciertos productos de potencias de los $x$. Para ello debe fijarse el concepto de ideal fundamental; más tarde pasará al módulo fundamental. Primero damos una notación que se mantendrá de ahora en adelante:

\medskip
\noindent\textbf{Notación.} \emph{Por $F^{(i)},f^{(i)},a^{(i)},\ldots$ entenderemos siempre polinomios en los $x$ que son libres de $x_1,\ldots,x_{i-1}$.}

\medskip
\noindent\textbf{Definición V.} \emph{A cada ideal $\mideal$ se le definen $n$ ideales fundamentales $\gideal_0,\gideal_1,\ldots,\gideal_{n-1}$ por la siguiente estipulación: el ideal fundamental de la etapa $(i-1)$, $\gideal_{i-1}$, contiene todos y sólo los polinomios $G(x)$ para los cuales existe un polinomio $b^{(i)}$ -- que en general varía con $G(x)$ -- tal que}
\begin{equation}
 b^{(i)}G(x)\zero{\mideal};\qquad b^{(i)}\ne0.
\tag{16}
\end{equation}

Que los $\gideal$ sean efectivamente ideales corresponde exactamente a la demostración de la propiedad modular para $\Gmod$ que sigue a la fórmula (4), con (16) en el papel análogo. El ideal $\gideal_i$ es divisible por $\gideal_{i-1}$, pues todo polinomio $b^{(i+1)}$ es al mismo tiempo un $b^{(i)}$.

Para los ideales fundamentales se tiene:

\medskip
\noindent\textbf{Teorema VI.} \emph{Los ideales fundamentales de ideales transformados son ideales transformados}\footnote{El Teorema VI se usa sólo en §~7.}.

La demostración se apoya en un teorema conocido de Dedekind--Mertens\footnote{Dedekind: Über einen arithmetischen Satz von Gauß, Mitt. d. deutsch. math. Ges. zu Prag 1892. F. Mertens: Über einen algebraischen Satz, Ber. d. Ak. d. Wissensch. Wien 101 (1892), pp. 1560--1566. -- La extensión por Kronecker del teorema de Gauss a polinomios con coeficientes indeterminados es consecuencia inmediata de este teorema.}: sean $a$ y $b$ indeterminadas, y póngase
\[
 \sum a_{i_1,\ldots,i_s}z_1^{i_1}\cdots z_s^{i_s}\cdot
 \sum b_{j_1,\ldots,j_s}z_1^{j_1}\cdots z_s^{j_s}
 =
 \sum c_{k_1,\ldots,k_s}z_1^{k_1}\cdots z_s^{k_s},
\]
\[
 \sum i\le \nu,
 \qquad
 \sum j\le \mu,
 \qquad
 \sum k\le \nu+\mu.
\]
Si $\Amod,\Bmod,\mathfrak C$ denotan los módulos derivados respectivamente de los coeficientes $a,b,c$ mediante enteros racionales, entonces existe un exponente $q$ tal que vale la identidad
\begin{equation}
 \Amod^q\Bmod=\Amod^{q-1}\mathfrak C.
\tag{17}
\end{equation}
Aquí $\Amod^q$ consta de los productos de potencias de dimensión $q$ en los $a$ y de sus combinaciones lineales enteras.

Para la demostración del Teorema VI, póngase ahora
\begin{equation}
\begin{cases}
 G(x)=\Gamma(y)=\sum \alpha_\lambda(u)\Gamma_\lambda(y)=\sum \alpha_\lambda(u)G_\lambda(x),\\
 b^{(i)}(x)=\varphi(y)=\sum \beta_\mu(u)\varphi_\mu(y),
\end{cases}
\tag{18}
\end{equation}
donde $\alpha_\lambda(u)$ y $\beta_\mu(u)$ son productos de potencias en los $u$. Entonces, por (16), (14) y (18),
\begin{equation}
 \sum\beta_\mu(u)\varphi_\mu(y)\cdot
 \sum\alpha_\lambda(u)\Gamma_\lambda(y)\zero{\overline{\mideal}_{(u)}};
 \qquad
 \sum\beta_\mu(u)\varphi_\mu(y)\ne0.
\tag{19}
\end{equation}
En el miembro izquierdo está el producto de dos polinomios en $u$ cuyos coeficientes son los polinomios $\varphi_\mu(y)$ y $\Gamma_\lambda(y)$; los coeficientes del polinomio producto en $u$ son, por el miembro derecho de (19), polinomios de $\overline{\mideal}$. Si, por tanto, en (17) se sustituyen los $a,b,c$ por estos polinomios en $y$, se obtiene
\[
 \left\{\sum\beta_\mu(u)\varphi_\mu(y)\right\}^q\cdot \Gamma_\lambda(y)\zero{\overline{\mideal}_{(u)}},
\]
lo que, por (18), equivale a
\begin{equation}
 b^{(i)}(x)^q\cdot G_\lambda(x)\zero{\mideal};
 \qquad
 G_\lambda(x)\zero{\gideal_{i-1}}.
\tag{20}
\end{equation}
Las relaciones (16), (18) y (20) muestran que los \emph{ideales fundamentales $\gideal_{i-1}$ son efectivamente ideales transformados}.

\section*{§ 4. Conexión entre ideales de polinomios y módulos de formas lineales.}

Para considerar un ideal $\mideal$ de polinomios, que siempre se supone transformado, como un módulo $\Mmod_{i-1}$ $(i=1,\ldots,n)$ de formas lineales, los polinomios individuales $F(x)$ deben considerarse como formas lineales en los productos de potencias $\xi_\lambda$ de $x_1\ldots x_{i-1}$:
\[
 F(x)=\sum a^{(i)}_\lambda\xi_\lambda,
\]
donde, conforme a la notación de §~3, los $a^{(i)}_\lambda$ son polinomios en $x_i\ldots x_n$. De la definición de ideal se sigue inmediatamente que $\Mmod_{i-1}$ tiene la propiedad modular respecto de estas magnitudes enteras $a^{(i)},b^{(i)},\ldots$; y, puesto que los infinitos productos de potencias $\xi$ de $x_1\ldots x_{i-1}$ aparecen sólo linealmente, por su independencia lineal desempeñan para $\Mmod_{i-1}$ el papel de indeterminadas. Cada módulo $\Mmod_{i-1}$ contiene infinitas indeterminadas $\xi$, pero en cada forma lineal individual sólo aparecen finitas indeterminadas. Análogamente, cada ideal fundamental $\gideal_{i-1}$ debe considerarse como un módulo $\Gmod_{i-1}$ de formas lineales, donde las indeterminadas son de nuevo los productos de potencias de $x_1\ldots x_{i-1}$. Para $i=1$ sólo hay una indeterminada $\xi_0$, correspondiente a la unidad.

Ahora se puede, y ésta es la base de todo lo que sigue, restringirse, a pesar de estas infinitas indeterminadas $\xi$, a módulos de formas lineales en finitas indeterminadas, según el teorema siguiente.

\medskip
\noindent\textbf{Teorema VII.} \emph{El sistema de clases residuales $\Gmod_{i-1}\mid\Mmod_{i-1}$ es isomorfo a un sistema de clases residuales $\Gmod^*_{i-1}\mid\Mmod^*_{i-1}$, donde $\Mmod^*_{i-1}$ es un módulo en finitas indeterminadas y $\Gmod^*_{i-1}$ es su módulo fundamental. El módulo $\Mmod_{i-1}$ tiene -- después de adjuntar $x_{i+1}\ldots x_n$ a $P(u)$ -- sólo finitos divisores elementales distintos de $1$, a saber, los divisores elementales de $\Mmod^*_{i-1}$ que son distintos de $1$.}

Los divisores elementales de $\Mmod_{i-1}$ deben definirse como en la nota 8). El isomorfismo que aparece en el Teorema VII se apoya en la definición siguiente.

\medskip
\noindent\textbf{Definición V.} \emph{Dos sistemas de clases residuales se llaman isomorfos si pueden ponerse en correspondencia biunívoca de tal modo que a la diferencia de dos clases se le asigne la diferencia de las clases correspondientes, y al producto de una clase por una magnitud entera el producto de la clase correspondiente por la misma magnitud entera.}

La demostración del Teorema VII se obtiene por inducción completa. Para $i=1$ el teorema es evidente: $\Mmod_0$ es un módulo de formas lineales en una indeterminada $\xi_0$, correspondiente al producto de potencias de dimensión cero de los $x$, es decir, a la unidad; las magnitudes enteras son todos los polinomios en $x_1\ldots x_n$. El ideal fundamental $\gideal_0$ es el ideal unidad, compuesto por todos los polinomios en $x_1\ldots x_n$, y por tanto $\Gmod_0$ es el módulo $(\xi_0)$, el módulo fundamental de $\Mmod_0$; aquí, por consiguiente, $\Gmod^*_0\mid\Mmod^*_0$ coincide con $\Gmod_0\mid\Mmod_0$. Finalmente, como $\Mmod_0$ tiene rango $1$, puede tener a lo sumo un divisor elemental distinto de $1$.

Pasamos primero de $i=1$ a $i=2$, para utilizar la comprensión así obtenida de la estructura de $\Gmod_1\mid\Mmod_1$ en el paso inductivo general. Para ello mostramos primero que para $\Gmod_0\mid\Mmod_0$ puede elegirse un sistema de representantes acotado en $x_1$; por la divisibilidad de $\gideal_1$ por $\gideal_0$, esto conducirá después a las finitas indeterminadas de $\Gmod^*_1$.

Como ideal transformado, $\mideal$ posee por §~3 al menos un polinomio $C^{(1)}(x)$ regular en $x_1$, digamos de dimensión $k$; por tanto $\Mmod_0$ contiene la forma lineal $C^{(1)}\xi_0$. Por la regularidad de $C^{(1)}(x)$, todo polinomio $H(x)$ admite una representación
\begin{equation}
 H(x)=b^{(2)}_0+b^{(2)}_1x_1+\cdots+b^{(2)}_{k-1}x_1^{k-1}\quad(C^{(1)}(x));
\tag{21}
\end{equation}
por ello, puesto que $C^{(1)}\xi_0$ pertenece a $\Mmod_0$, se puede efectivamente escoger para $\Gmod_0\mid\Mmod_0$ un sistema de representantes que no alcance la dimensión $k$ en $x_1$.

Si ahora, en particular, para los polinomios $G(x)$ de $\gideal_1$ y $F(x)$ de $\mideal$ se escribe
\begin{equation}
 \left\{\begin{aligned}
 G(x)&=g^{(2)}_0+g^{(2)}_1x_1+\cdots+g^{(2)}_{k-1}x_1^{k-1} &&(C^{(1)}(x)),\\
 F(x)&=a^{(2)}_0+a^{(2)}_1x_1+\cdots+a^{(2)}_{k-1}x_1^{k-1} &&(C^{(1)}(x)),
 \end{aligned}\right.
\tag{22}
\end{equation}
y si $\Gmod^*_1$ denota el sistema de todas las formas lineales $g(\xi)=g^{(2)}_0\xi_0+\cdots+g^{(2)}_{k-1}\xi_{k-1}$, y correspondientemente $\Mmod^*_1$ el sistema de las formas lineales $a(\xi)=a^{(2)}_0\xi_0+\cdots+a^{(2)}_{k-1}\xi_{k-1}$, entonces la propiedad de ideal de $\gideal_1$ y $\mideal$ muestra que $\Gmod^*_1$ y $\Mmod^*_1$ son módulos de formas lineales en $\xi_0\ldots\xi_{k-1}$ respecto de las magnitudes enteras $b^{(2)},c^{(2)},\ldots$. Asimismo el ideal principal derivado de $C^{(1)}(x)$ da un módulo respecto de estas magnitudes enteras,
\[
 \Gmod_1=(\xi_0,\xi_1,\ldots,\xi_r,\ldots),
\]
donde $\xi_r=x_1^rC^{(1)}(x)$. Los $\xi_r$ son linealmente independientes respecto de estas magnitudes enteras $b^{(2)},c^{(2)},\ldots$, tanto entre sí como respecto de los finitos $\xi$. De la divisibilidad de $\Gmod_1$ por $\Mmod_1$, y por tanto por $\Gmod_1$, se sigue que $\Gmod^*_1$ es divisible por $\Gmod_1$ y $\Mmod^*_1$ por $\Mmod_1$. Teniendo en cuenta (22), se obtiene
\begin{equation}
 \Gmod_1=(\Gmod^*_1,\Gmod_1),\qquad
 \Mmod_1=(\Mmod^*_1,\Gmod_1).
\tag{23}
\end{equation}

De (23) y de la independencia lineal de los $\xi$ y $\xi$, el Teorema VII para $i=2$ se sigue así. Como $\Gmod_1$ es divisible por $\Mmod_1$, primero se tiene $\Gmod_1\mid\Mmod_1=\Gmod^*_1\mid\Mmod_1$. Para demostrar el isomorfismo con $\Gmod^*_1\mid\Mmod^*_1$, tómense ciertas formas lineales $g^*(\xi)$ de $\Gmod^*_1$ que den un sistema de representantes de $\Gmod^*_1\mid\Mmod^*_1$. Por la independencia lineal de los $\xi$ y $\xi$, de $g^*(\xi)\equiv0(\Mmod_1)$ se sigue también $g^*(\xi)\equiv0(\Mmod^*_1)$; las $g^*(\xi)$ son por tanto incongruentes también módulo $\Mmod_1$, forman un sistema de representantes de $\Gmod_1\mid\Mmod_1$, y la correspondencia de $\Gmod_1\mid\Mmod_1$ con $\Gmod^*_1\mid\Mmod^*_1$ mediada por este sistema de representantes es isomorfa en el sentido de la Definición V.

Exactamente de la misma manera se obtiene: de $b^{(2)}g(\xi)\equiv0(\Mmod_1)$, $b^{(2)}\ne0$, se sigue $b^{(2)}g(\xi)\equiv0(\Mmod^*_1)$, $b^{(2)}\ne0$. Como esto se satisface para todos los $g(\xi)$ de $\Gmod^*_1$, y sólo para éstos -- pues por (23) $\Gmod^*_1$ consta de todos los polinomios del ideal fundamental $\gideal_1$ que son al mismo tiempo formas lineales en los $\xi$ --, $\Gmod^*_1$ queda así caracterizado como el módulo fundamental de $\Mmod^*_1$.

Adjuntando finalmente $x_3\ldots x_n$ a $P(u)$, se obtiene
\begin{equation}
\left\{\begin{aligned}
 \Gmod^*_1&=(\eta_1,\ldots,\eta_\rho),&
 \Mmod^*_1&=(e_1\eta_1,\ldots,e_\rho\eta_\rho),\\
 \Gmod_1&=(\eta_\rho,\ldots,\eta_1,\xi_0,\ldots,\xi_r,\ldots),&
 \Mmod_1&=(e_\rho\eta_\rho,\ldots,e_1\eta_1,\xi_0,\ldots,\xi_r,\ldots),
\end{aligned}\right.
\tag{24}
\end{equation}
y de ahí
\[
 (e_\rho)=\Mmod_1/\Gmod_1=\Mmod_1/(\Mmod_1,\eta_\rho),\qquad
 (e_{\rho-1})=(\Mmod_1,\eta_\rho)/\Gmod_1\quad\hbox{etc.}
\]
Así, teniendo en cuenta la nota 8), se reconocen $e_\rho,e_{\rho-1},\ldots,e_1,1,\ldots,1,\ldots$ como los divisores elementales de $\Mmod_1$. Con esto se prueban todas las partes del Teorema VII para $i=2$.

Debe notarse que (22) y (23) sólo utilizan que $C^{(1)}(x)$ es regular en $x_1$. Estas fórmulas, y los argumentos unidos a ellas que conducen al Teorema VII, siguen siendo válidos si en vez de (12) se toma como base la transformación especial
\begin{equation}
 y_1=x_1;\quad y_2=u_{21}x_1+z_2;\quad\ldots;\quad y_n=u_{n1}x_1+z_n
\tag{25}
\end{equation}
por composición con
\begin{equation}
\left\{\begin{array}{l}
 z=U_1(x)\quad\hbox{o}\\[2mm]
 z_2=x_2;\quad z_3=u_{32}x_2+x_3;\quad\ldots;\quad z_n=u_{n2}x_2+\cdots+u_{n,n-1}x_{n-1}+x_n,
\end{array}\right.
\tag{26}
\end{equation}
lo que da de nuevo (12). Si $\widetilde{\mideal}_{(u)}$ pasa mediante (25) a $\widetilde{\mideal}$, y si $\widetilde{\gideal}_1,\widetilde{\Gmod}_1,\ldots$ tienen para $\widetilde{\mideal}$ el mismo significado que $\gideal_1,\Gmod_1,\ldots$ tienen para $\mideal$, entonces
\begin{equation}
 \widetilde{\Gmod}_1=(\widetilde{\Gmod}^{*}_1,\widetilde{\Gmod}_1),\qquad
 \widetilde{\Mmod}_1=(\widetilde{\Mmod}^{*}_1,\widetilde{\Gmod}_1).
\tag{27}
\end{equation}
Aquí (27) surge de (23) simplemente por inversión de (26). Esto es evidente para $\mideal$ y $C^{(1)}(x)$, y por tanto también para $\Gmod_1$ y $\Mmod_1$. También es cierto para $\gideal_1$, pues
\[
 b^{(2)}(x)G(x)\equiv0(\mideal),\quad b^{(2)}\ne0,
 \qquad
 \widetilde b^{(2)}(z)\widetilde G(x,z)\equiv0(\widetilde\mideal),\quad \widetilde b^{(2)}\ne0
\]
se condicionan mutuamente por $z=U_1(x)$. Así $\widetilde\gideal_1=\gideal_1$ por medio de (26); de ahí también $\widetilde\Gmod_1=\Gmod_1$. Por consiguiente, por la definición dada mediante (22), lo mismo vale para $\Gmod_1^*$ y $\Mmod_1^*$ que para $\widetilde\Gmod_1^*$ y $\widetilde\Mmod_1^*$, y por tanto para toda la representación (27).

Para el paso inductivo general supóngase que valen las hipótesis siguientes:
\begin{equation}
\left\{\begin{aligned}
 \Gmod_{i-1}&=(\Gmod_{i-1}^*,\Gmod_{i-1}),&
 \Mmod_{i-1}&=(\Mmod_{i-1}^*,\Gmod_{i-1}),\\
 \Gmod_{i-1}&=(\xi_0,\xi_1,\ldots,\xi_r,\ldots),
\end{aligned}\right.
\tag{28}
\end{equation}
donde $\Gmod^*_{i-1}$ y $\Mmod^*_{i-1}$ son módulos, respecto de las magnitudes enteras $b^{(i)},c^{(i)},\ldots$, de formas lineales en finitas indeterminadas $\xi$, ciertos productos de potencias de $x_1,\ldots,x_{i-1}$; y donde los $\xi$ son linealmente independientes, tanto entre sí como respecto de los $\xi$, con respecto a esas magnitudes enteras. Una representación correspondiente a (28), y la independencia lineal de los $\xi$ y $\xi$, debe seguir valiendo si en vez de (12) se utiliza la transformación especial
\begin{equation}
\begin{aligned}
 y_1&=x_1;\quad \ldots;\quad y_{i-1}=u_{i-1,1}x_1+\cdots+x_{i-1};\\
 y_i&=u_{i,1}x_1+\cdots+u_{i,i-1}x_{i-1}+z_i;\quad\ldots;\quad
 y_n=u_{n,1}x_1+\cdots+u_{n,i-1}x_{i-1}+z_n,
\end{aligned}
\tag{29}
\end{equation}
que puede componerse para dar (12) mediante
\[
 z=U_{i-1}(x)\quad\hbox{o}\quad
 z_i=x_i;\quad z_{i+1}=u_{i+1,i}x_i+x_{i+1};\quad\ldots;\quad
 z_n=u_{n,i}x_i+\cdots+u_{n,n-1}x_{n-1}+x_n.
\]
Además, esta representación especial debe surgir de (28) simplemente por inversión de $z=U_{i-1}(x)$.

De las hipótesis (28) y de la independencia lineal de los $\xi$ y $\xi$, el isomorfismo de $\Gmod_{i-1}\mid\Mmod_{i-1}$ con $\Gmod^*_{i-1}\mid\Mmod^*_{i-1}$ se sigue, mediante la asignación al mismo sistema de representantes, exactamente por los mismos argumentos que antes acompañaban a (23). Igualmente se sigue que $\Gmod^*_{i-1}$ es el módulo fundamental de $\Mmod^*_{i-1}$, y que $\Mmod_{i-1}$ tiene sólo finitos divisores elementales distintos de $1$, a saber, aquellos divisores elementales de $\Mmod^*_{i-1}$ que son distintos de $1$.

Para efectuar el paso inductivo queda por mostrar de nuevo, primero, que para $\Gmod^*_{i-1}\mid\Mmod^*_{i-1}$, y por tanto también para $\Gmod_{i-1}\mid\Mmod_{i-1}$, puede elegirse un sistema de representantes acotado en $x_i$. Esto se sigue de la representación por cociente demostrada en (7), Teorema III:
\begin{equation}
 \Gmod^*_{i-1}=\Mmod^*_{i-1}/(C^{(i)}),
\tag{30}
\end{equation}
donde $C^{(i)}$ significa cualquier determinante no idénticamente nulo de grado $\rho$ de $\Mmod^*_{i-1}$, siendo $\rho$ el rango de $\Mmod^*_{i-1}$. De la parte de las hipótesis referente a la transformación especial (29), se sigue que $C^{(i)}$ puede suponerse siempre regular respecto de $x_i$. En efecto, el módulo $\widetilde\Mmod^*_{i-1}$ correspondiente a $\Mmod^*_{i-1}$ tiene el mismo rango; si $\widetilde C^{(i)}$ denota cualquier determinante no idénticamente nulo de grado $\rho$ de $\widetilde\Mmod^*_{i-1}$ que por $z=U_{i-1}(x)$ pasa a un $C^{(i)}$ regular, entonces por hipótesis $C^{(i)}$ es un determinante de grado $\rho$ de $\Mmod^*_{i-1}$.

De la existencia de $C^{(i)}$ se sigue, como en (8), con una numeración adecuada de los $\xi$, que pertenecen a $\Mmod^*_{i-1}$ $\rho$ formas lineales de la forma especial
\[
 L_\lambda(\xi)=C^{(i)}\xi_\lambda+c^{(i)}_{\lambda,1}\xi_{\rho+1}+\cdots+c^{(i)}_{\lambda,s-\rho}\xi_s
 \qquad(\lambda=1\ldots\rho).
\]
Toda forma lineal en $\xi$ puede llevarse entonces a la forma
\begin{equation}
 H(\xi)=a^{(i)}_1\xi_1+\cdots+a^{(i)}_\rho\xi_\rho+b^{(i)}_1\xi_{\rho+1}+\cdots+b^{(i)}_{s-\rho}\xi_s\quad (L_1(\xi),\ldots,L_\rho(\xi)),
\tag{31}
\end{equation}
donde $a^{(i)}_1\ldots a^{(i)}_\rho$ están acotados en $x_i$ y no alcanzan el grado $k$ de $C^{(i)}$. Esta representación es única; pues de
\[
 a^{(i)}_1\xi_1+\cdots+a^{(i)}_\rho\xi_\rho+b^{(i)}_1\xi_{\rho+1}+\cdots+b^{(i)}_{s-\rho}\xi_s\equiv0\quad (L_1(\xi),\ldots,L_\rho(\xi))
\]
se sigue, comparando coeficientes en $\xi_1\ldots\xi_\rho$ y considerando los grados en $x_i$, que todos los $a^{(i)}$ y $b^{(i)}$ se anulan idénticamente.

Sea ahora en particular
\[
 H(\xi)\equiv0(\Gmod^*_{i-1});\qquad\hbox{por tanto}\qquad C^{(i)}H(\xi)\equiv0(\Mmod^*_{i-1})\quad\hbox{por (30).}
\]
Como en la demostración del Teorema III se obtiene
\[
 C^{(i)}H(\xi)-\sum_{\tau=1}^{s-\rho}\{C^{(i)}b^{(i)}_\tau-\sum_\lambda a^{(i)}_\lambda c^{(i)}_{\lambda,\tau}\}\xi_{\rho+\tau}
 \equiv0(\Mmod^*_{i-1}),
\]
y por tanto, puesto que como en el Teorema III los coeficientes de $\xi_{\rho+1}\ldots\xi_s$ deben anularse,
\[
 C^{(i)}b^{(i)}_\tau=\sum_\lambda a^{(i)}_\lambda c^{(i)}_{\lambda,\tau}\qquad(\tau=1\ldots s-\rho).
\]
Así también los $b^{(i)}_\tau$ están acotados y no alcanzan el grado máximo de los $c^{(i)}_{\lambda,\tau}$. Con ello queda demostrada la existencia de un sistema de representantes acotado en $x_i$ para $\Gmod^*_{i-1}\mid\Mmod^*_{i-1}$, es decir, para $\Gmod_{i-1}\mid\Mmod_{i-1}$, que no alcanza cierto grado fijo $r$.

Sean $\vartheta_1\ldots\vartheta_t$ los finitos productos de potencias
\[
 x_i^\alpha\xi_\lambda\quad(\alpha=0,1,\ldots,k-1;\ \lambda=1,2,\ldots,\rho),\qquad
 x_i^\beta\xi_{\rho+\tau}\quad(\beta=0,1,\ldots,r-1;\ \tau=1,2,\ldots,s-\rho)
\]
y dispónganse las expresiones numerablemente infinitas
\[
 x_i^\gamma L_\lambda\quad(\gamma=0,1,\ldots\text{ hasta el infinito};\ \lambda=1,2,\ldots,\rho),
 \qquad
 x_i^\gamma\xi_\nu\quad(\gamma,\nu=0,1,\ldots\text{ hasta el infinito})
\]
en una sucesión simplemente infinita $\omega_0,\omega_1,\ldots,\omega_\mu,\ldots$. Entonces los $\vartheta$ y los $\omega$ son linealmente independientes, tanto cada familia entre sí como unas respecto de otras, con respecto a las magnitudes enteras $b^{(i+1)},c^{(i+1)},\ldots$. En efecto, de
\[
 \sum c^{(i+1)}_\alpha x_i^\alpha\xi_\lambda+
 \sum c^{(i+1)}_\beta x_i^\beta\xi_{\rho+\tau}+
 \sum d^{(i+1)}_\gamma x_i^\gamma L_\lambda(\xi)+
 \sum e^{(i+1)}_{\gamma\nu}x_i^\gamma\xi_\nu=0,
\]
donde cada suma se extiende sólo sobre finitos términos, la independencia lineal supuesta de los $\xi$ y $\xi$, y de los $\xi$ entre sí respecto de las magnitudes enteras $b^{(i)}$, da la anulación de todos los $e^{(i+1)}_{\gamma\nu}$. La unicidad de la representación (31) da además la anulación de los $c^{(i+1)}_\alpha$ y $c^{(i+1)}_\beta$, y finalmente, por la independencia lineal de los $L_\lambda(\xi)$, la anulación de los $d^{(i+1)}_\gamma$.

Si el módulo $(\Gmod_{i-1},L_1(\xi),\ldots,L_\rho(\xi))$ se considera ahora como un módulo $\Gmod_i$ respecto de las magnitudes enteras $b^{(i+1)}$, entonces $\Gmod_i$ es divisible por $\Mmod_i$, porque $\Gmod_{i-1},L_1(\xi),\ldots,L_\rho(\xi)$ son divisibles por $\Mmod_{i-1}$, y
\[
 \Gmod_i=(\omega_0,\omega_1,\ldots,\omega_\mu,\ldots).
\]
Para cada $H(x)\equiv0(\gideal_{i-1})$, (28), (31) y lo que se acaba de demostrar dan
\[
 H(x)=b^{(i+1)}_1\vartheta_1+\cdots+b^{(i+1)}_t\vartheta_t\quad(\Gmod_i)
\]
como ecuación de módulos respecto de las magnitudes enteras $b^{(i+1)},c^{(i+1)}$. Puesto que $\gideal_i$ es divisible por $\gideal_{i-1}$ y $\mideal$ es divisible por $\gideal_i$, si ahora se pone, en particular, para todo $G(x)$ de $\Gmod_i$ y todo $F(x)$ de $\Mmod_i$,
\[
 G(x)=g^{(i+1)}_1\vartheta_1+\cdots+g^{(i+1)}_t\vartheta_t\quad(\Gmod_i),
 \qquad
 F(x)=a^{(i+1)}_1\vartheta_1+\cdots+a^{(i+1)}_t\vartheta_t\quad(\Gmod_i),
\]
y se denota por $\Gmod_i^*$ el módulo compuesto por todos los $g(\vartheta)$, y por $\Mmod_i^*$ el módulo compuesto por todos los $a(\vartheta)$, entonces exactamente los argumentos que conducían a (23) dan
\begin{equation}
 \Gmod_i=(\Gmod_i^*,\Gmod_i),\qquad
 \Mmod_i=(\Mmod_i^*,\Gmod_i),\qquad
 \Gmod_i=(\omega_0,\omega_1,\ldots,\omega_\mu,\ldots).
\tag{32}
\end{equation}
En la deducción de (32) se utilizó de nuevo sólo la regularidad de $C^{(i)}$ en $x_i$, de modo que todo el argumento sigue siendo válido si se toma la transformación especial (29) con el índice aumentado en uno. Las mismas consideraciones añadidas a (27) muestran entonces que esta representación especial surge de (32) simplemente por inversión de $z=U_i(x)$.

Así quedan demostradas para el índice siguiente todas las hipótesis (28), etc., del paso inductivo general; y como, según se mostró allí, el Teorema VII se sigue directamente de estas hipótesis, el Teorema VII queda demostrado en general. Al mismo tiempo se ha mostrado que la arbitrariedad en $(\Gmod^*_{i-1},\Mmod^*_{i-1})$ causada por la elección arbitraria de $C^{(i)}$ queda de nuevo eliminada por el isomorfismo del sistema de clases residuales $\Gmod^*_{i-1}\mid\Mmod^*_{i-1}$.

En particular, si el cuerpo $P$ utilizado en §~3 como dominio de coeficientes tiene característica cero -- es decir, si el cuerpo primo contenido en $P$ y derivado de la unidad es del tipo de los números racionales --, las indeterminadas $u_{\mu\nu}$ pueden especializarse a magnitudes $\bar u_{\mu\nu}$ de $P$ de modo que para $i=1\ldots n$ cada $C^{(i)}$ permanezca regular en $x_i$. Los argumentos anteriores muestran que el Teorema VII sigue siendo válido también bajo esta especialización.\footnote{Hentzelt toma, en lugar de un $P$ arbitrario, sólo el cuerpo de todos los números complejos y trata principalmente este último caso, mientras que usa indeterminadas sólo en la teoría de eliminación propiamente dicha. Hentzelt muestra después, esencialmente por los argumentos aquí dados, que para formar la forma resultante basta, en cada $\Mmod_{i-1}$, llegar sólo hasta algún grado finito en los $x$ que exceda una cota fija. Lo que falta es el isomorfismo de $\Gmod_{i-1}\mid\Mmod_{i-1}$ con $\Gmod^*_{i-1}\mid\Mmod^*_{i-1}$ dado en el Teorema VII, que explica la independencia respecto de los números de grado. (E. N.)} El módulo fundamental y los divisores elementales no necesitan, sin embargo, provenir necesariamente del caso general por la especialización $u_{\mu\nu}=\bar u_{\mu\nu}$.\footnote{Para $\mideal=(x^2,ux+y)$ se tiene $\gideal_0=(1)$, $\gideal_1=(1)$. Si se escoge $C^{(1)}=x^2$, entonces $\Gmod_1^*=(1,x)$, $\Mmod_1^*=(ux+y,xy)$; este $\Mmod_1^*$ tiene divisores elementales $y^2$ y $1$, y puede escogerse $C^{(2)}=y^2$. Para $u=0$, $C^{(1)}$ y $C^{(2)}$ siguen siendo regulares en $x$ e $y$, respectivamente; pero $\Mmod_1^*=(y,xy)$ tiene divisores elementales $y,y$. Así $\Gmod_1\mid\Mmod_1$ sigue siendo isomorfo al sistema de clases residuales de un módulo finito, pero no es isomorfo a $\Gmod_1\mid\Mmod_1$. Si en cambio se hubiera escogido $C^{(1)}=ux+y$ y se hubiera especializado de modo que $C^{(1)}$ permaneciera regular, también se habría conservado el isomorfismo. (E. N.)}

\section*{§ 5. Forma resultante y forma de divisores elementales de un ideal de polinomios.}

Adjúntense $x_{i+1}\ldots x_n$ a $P(u)$, de modo que $\Gmod_{i-1}$ y $\Mmod_{i-1}$, y por tanto también $\Gmod^*_{i-1}$ y $\Mmod^*_{i-1}$, pasen a módulos de formas lineales en los que las magnitudes enteras pueden considerarse como polinomios en una indeterminada $x_i$. Por el Teorema VII, en este caso los divisores elementales distintos de $1$ de $\Mmod_{i-1}$, junto con a lo sumo finitos divisores elementales iguales a $1$, coinciden con los de $\Mmod^*_{i-1}$. El producto de estos divisores elementales, el divisor determinantal $\rho$-ésimo de $\Mmod^*_{i-1}$, era igual al determinante de la sustitución de paso de $\Gmod^*_{i-1}$ a $\Mmod^*_{i-1}$, por tanto a $N(\Gmod^*_{i-1}\mid\Mmod^*_{i-1})$ según la Definición III. La fórmula (24), o la fórmula análoga que se sigue de (28) para $i$ general, muestra que el producto de estos divisores elementales es también igual al determinante de la sustitución de paso de $\Gmod_{i-1}$ a $\Mmod_{i-1}$, y por ello debe denotarse por $N(\Gmod_{i-1}\mid\Mmod_{i-1})$. Así
\[
 N(\Gmod_{i-1}\mid\Mmod_{i-1})=N(\Gmod^*_{i-1}\mid\Mmod^*_{i-1})=R^{(i)}(x),
\]
donde el polinomio $R^{(i)}(x)$, es decir, el máximo común divisor, en el sentido polinómico, de todos los determinantes de $\rho$ filas de $\Mmod^*_{i-1}$, puede suponerse entero y primitivo en $x_{i+1}\ldots x_n$, y en consecuencia, como divisor de $C^{(i)}$, se vuelve regular en $x_i$. Asimismo el divisor elemental más alto $E^{(i)}(x)$ de $\Mmod_{i-1}$, respectivamente de $\Mmod^*_{i-1}$, puede suponerse entero y primitivo en $x_{i+1}\ldots x_n$ y por tanto regular en $x_i$. Esto conduce a la definición siguiente.

\medskip
\noindent\textbf{Definición VI.} \emph{El polinomio $R^{(i)}(x)$ se llama la resultante de la $i$-ésima etapa de $\mideal$, y $E^{(i)}(x)$ el divisor elemental de la $i$-ésima etapa. Así, la resultante de la $i$-ésima etapa es la norma del ideal fundamental $\gideal_{i-1}$ respecto de $\mideal$, interpretada como la norma del módulo $\Gmod_{i-1}$ respecto de $\Mmod_{i-1}$, donde $x_{i+1}\ldots x_n$ se adjuntan a $P(u)$; asimismo $E^{(i)}(x)$, bajo la misma adjunción, se vuelve igual al cociente $\Mmod_{i-1}/\Gmod_{i-1}$. La forma resultante y la forma de divisores elementales de $\mideal$ son los productos}
\[
 R_\mideal=R^{(1)}R^{(2)}\cdots R^{(n)},\qquad
 E_\mideal=E^{(1)}E^{(2)}\cdots E^{(n)}.
\]

Para la forma resultante y la forma de divisores elementales se tiene:

\medskip
\noindent\textbf{Teorema VIII.} \emph{La forma resultante es divisible por la forma elemental; recíprocamente, una potencia de $E_\mideal$ es divisible por $R_\mideal$. Las formas $E_\mideal$ y $R_\mideal$ son divisibles por $\mideal$; más generalmente, $E^{(i)}\cdots E^{(n)}\gideal_{i-1}\equiv0(\mideal)$ y, en consecuencia, $R^{(i)}\cdots R^{(n)}\gideal_{i-1}\equiv0(\mideal)$.}

En efecto, puesto que la norma es divisible por el divisor elemental más alto y, recíprocamente, una potencia de este divisor elemental más alto es divisible por la norma, esta divisibilidad se sigue para $R^{(i)}(x)$ y $E^{(i)}(x)$ después de adjuntar $x_{i+1}\ldots x_n$. Pero $R^{(i)}(x)$ y $E^{(i)}(x)$ se suponen polinomios primitivos respecto de $x_{i+1}\ldots x_n$; por tanto la divisibilidad vale respecto de todas las indeterminadas $x_i,x_{i+1},\ldots,x_n$. Así $R_\mideal$ es divisible por $E_\mideal$, y alguna potencia de $E_\mideal$ por $R_\mideal$, conforme a la definición de estas expresiones en todas las indeterminadas $x$.

Además, por los Teoremas VII y II, después de adjuntar $x_{i+1}\ldots x_n$, el divisor elemental más alto $E^{(i)}(x)$ está definido como el cociente $\Mmod_{i-1}/\Gmod_{i-1}$. Volviendo entonces a los polinomios en todas las indeterminadas $x$, para todo
\begin{equation}
 G(x)\equiv0(\gideal_{i-1})
 \quad\hbox{existe un}\quad b^{(i+1)}\ne0
 \quad\hbox{tal que}\quad b^{(i+1)}E^{(i)}(x)G(x)\equiv0(\mideal).
\tag{33}
\end{equation}
En particular, $\gideal_0$ es el ideal unidad, de donde
\[
 b^{(2)}E^{(1)}(x)\equiv0(\mideal),\quad b^{(2)}\ne0,
 \qquad\hbox{y por tanto}\qquad E^{(1)}(x)\equiv0(\gideal_1).
\]
En general, de (33), aplicado a
\[
 E^{(1)}(x)\cdots E^{(i-1)}(x)\equiv0(\gideal_{i-1}),
\]
existe $b^{(i+1)}\ne0$ tal que $b^{(i+1)}E^{(1)}\cdots E^{(i)}\equiv0(\mideal)$; de ahí $E^{(1)}\cdots E^{(i)}\equiv0(\gideal_i)$. Como $\gideal_n=\mideal$, se obtiene
\begin{equation}
 E_\mideal=E^{(1)}\cdots E^{(n)}\equiv0(\mideal)
 \quad\hbox{y en consecuencia}\quad
 R_\mideal=R^{(1)}\cdots R^{(n)}\equiv0(\mideal).
\tag{34}
\end{equation}

Si en (33) $G(x)$ recorre los finitos polinomios de una base ideal de $\gideal_{i-1}$, y si $c^{(i+1)}(x)$ denota el producto de los correspondientes $b^{(i+1)}(x)$, entonces
\[
 c^{(i+1)}E^{(i)}(x)\gideal_{i-1}\equiv0(\mideal);
 \qquad c^{(i+1)}\ne0,
 \quad\hbox{y de ahí}\quad E^{(i)}(x)\gideal_{i-1}\equiv0(\gideal_i),
\]
y de esto, como antes, por repetición finita,
\[
 E^{(n)}(x)\cdots E^{(i)}(x)\gideal_{i-1}\equiv0(\mideal)
\]
y por consiguiente también
\[
 R^{(n)}(x)\cdots R^{(i)}(x)\gideal_{i-1}\equiv0(\mideal).
\]
Esto demuestra todas las partes del Teorema VIII.

El Teorema VIII depende esencialmente de propiedades de la forma de divisores elementales; la significación característica de la forma resultante para $\mideal$ la muestra

\medskip
\noindent\textbf{Teorema IX.} \emph{Si $\nideal$ es un divisor de $\mideal$ -- divisor entendido en el sentido de ideales -- y si las formas resultantes $R_\nideal$ y $R_\mideal$ coinciden, entonces los ideales $\nideal$ y $\mideal$ coinciden.}

Sean $\gideal_0\ldots\gideal_{n-1},\gideal_n=\mideal$ y $\hideal_0\ldots\hideal_{n-1},\hideal_n=\nideal$ los ideales fundamentales de $\mideal$ y de $\nideal$. Primero $\gideal_0$ y $\hideal_0$ son ambos el ideal unidad, luego $\gideal_0=\hideal_0$. Supóngase ahora que $\gideal_{i-1}=\hideal_{i-1}$, de modo que $\Mmod_{i-1}$ y $\Nmod_{i-1}$ tienen el mismo módulo fundamental $\Gmod_{i-1}$. La hipótesis $R_\nideal=R_\mideal$, después de adjuntar $x_{i+1}\ldots x_n$, puede escribirse como
\[
 N(\Gmod_{i-1}\mid\Nmod_{i-1})=N(\Gmod_{i-1}\mid\Mmod_{i-1})
 \quad\hbox{o}\quad
 N(\Gmod^*_{i-1}\mid\Nmod^*_{i-1})=N(\Gmod^*_{i-1}\mid\Mmod^*_{i-1}).
\]
Aquí, de acuerdo con (32), $\Nmod_{i-1}=(\Nmod^*_{i-1},\Gmod_{i-1})$, de modo que $\Nmod^*_{i-1}$ es igualmente un módulo de formas lineales en $\xi$, con $\Gmod^*_{i-1}$ como módulo fundamental. Puesto que $\Mmod_{i-1}$ es divisible por $\Nmod_{i-1}$, lo cual implica la divisibilidad de $\Mmod^*_{i-1}$ por $\Nmod^*_{i-1}$, el Teorema IV muestra que, bajo esta adjunción, los dos módulos $\Mmod^*_{i-1}$ y $\Nmod^*_{i-1}$ coinciden, y en consecuencia también $\Nmod_{i-1}$ y $\Mmod_{i-1}$. Pero adjuntar $x_{i+1}\ldots x_n$ significa que se añaden a $\Mmod_{i-1}$, respectivamente a $\mideal$, todos los polinomios $G(x)$ para los cuales
\[
 b^{(i+1)}(x)G(x)\equiv0(\mideal),\qquad b^{(i+1)}\ne0.
\]
Así, por la adjunción, $\Mmod_{i-1}$, respectivamente $\mideal$, pasa a $\gideal_i$, y correspondientemente $\nideal$ a $\hideal_i$; por tanto $\gideal_i=\hideal_i$. Finalmente, $\gideal_n=\hideal_n$, es decir, $\mideal=\nideal$.

Por último, el mismo principio de transferencia, aplicado al Teorema V, da lo siguiente.

\medskip
\noindent\textbf{Teorema X.} \emph{Sea $R^{(i)}=S^{(i)}T^{(i)}$ una descomposición de $R^{(i)}$ en polinomios relativamente primos $S^{(i)}$ y $T^{(i)}$ en el sentido polinómico. Entonces existe uno y sólo un ideal $\jideal$, y asimismo uno y sólo un ideal $\tideal$, ambos divisores de $\mideal$ con el mismo ideal fundamental de la etapa $(i-1)$, $\gideal_{i-1}$, tales que $S^{(i)}=N(\Gmod_{i-1}\mid\Smod_{i-1})$ y $T^{(i)}=N(\Gmod_{i-1}\mid\mathfrak T_{i-1})$. Los ideales $\jideal$ y $\tideal$ se definen como la totalidad de los polinomios $S(x)$ y $T(x)$, respectivamente, tales que}
\[
 s^{(i+1)}(x)T^{(i)}(x)S(x)\equiv0(\mideal),\quad s^{(i+1)}\ne0,
\]
\[
 t^{(i+1)}(x)S^{(i)}(x)T(x)\equiv0(\mideal),\quad t^{(i+1)}\ne0,
\]
\emph{y $\gideal_i$ se convierte en el mínimo común múltiplo de $\jideal$ y $\tideal$,}
\[
 \gideal_i=[\jideal,\tideal].
\]

Sólo queda observar que $s^{(i+1)}$ y $t^{(i+1)}$ pueden elegirse fijos para $\jideal$ y $\tideal$, como producto de los $s^{(i+1)}$ y $t^{(i+1)}$ correspondientes a los elementos de base de $\jideal$ y $\tideal$, y que, como se observó arriba, al adjuntar $x_{i+1}\ldots x_n$ el ideal $\mideal$ se transforma en $\gideal_i$. En particular, la descomposición de $R^{(i)}$ puede continuarse hasta potencias de polinomios irreducibles -- factores primarios --, lo que conlleva una representación correspondiente de $\gideal_i$ como mínimo común múltiplo.

\section*{§ 6. Teoría propia de la eliminación.}

Por el Teorema VIII, respectivamente por la fórmula (34), la forma resultante y la forma de divisores elementales son divisibles por $\mideal$, y por ello se anulan en todos los ceros de $\mideal$. Por ``ceros de $\mideal$'' se entienden aquí sistemas de valores de $x_1\ldots x_i$ pertenecientes al cuerpo algebraicamente cerrado derivado de $P(u;x_{i+1}\ldots x_n)$\footnote{Steinitz mostró, en su Algebraische Theorie der Körper (J. f. M. 137 (1910), pp. 167--309), que todo cuerpo puede extenderse, esencialmente de modo único, a uno algebraicamente cerrado, dando así el equivalente racional del teorema fundamental del álgebra. De hecho, para cada $i=1,\ldots,n$ se trata de un cuerpo de extensión finito de $P(u,x_{i+1},\ldots,x_n)$. (E. N.)}, tales que todo polinomio de $\mideal$ se anula en esos sistemas de valores, mientras que $x_{i+1}\ldots x_n$ se piensan adjuntados al cuerpo de coeficientes $(i=1,\ldots,n)$. Estos $x_{i+1}\ldots x_n$ adjuntados pueden, como se mostrará, reemplazarse también por magnitudes de $P(u)$. El contenido de la teoría de la eliminación, recíprocamente, es la deducción de los ceros de $\mideal$ a partir de los de la forma resultante. Esta recíproca descansa en la eliminación sucesiva que se desarrollará en esta sección; como en la sección siguiente se demostrará la divisibilidad de la forma $D^{(i)}(x)$ que aparece en ella por $E^{(i)}(x)$, la recíproca deseada se sigue de la divisibilidad de una potencia de $E^{(i)}(x)$ por $R^{(i)}(x)$.

La teoría de la eliminación sucesiva que se dará aquí descansa en el resultado siguiente.

\medskip
\noindent\textbf{Teorema XI.} \emph{Sea $\bideal$ un ideal de polinomios en una indeterminada $t$ con coeficientes en $P$, por tanto un ideal principal; sea $h(t)$ un polinomio de $\bideal$ de grado $k$, y sea $\Bmod$ el módulo de formas lineales en $1,t,\ldots,t^{k-1}$ al que pasa $\bideal$ módulo $h(t)$. Entonces $\Bmod$ tiene rango $k-p$, donde $p$ es el grado del polinomio de base $f(t)$ de $\bideal$.}

Por hipótesis $h(t)=h_1(t)f(t)$, donde $h_1(t)$ tiene grado $k-p$. Las clases residuales de $\bideal$ módulo $h(t)$ representadas por $f(t),tf(t),\ldots,t^{k-p-1}f(t)$ son por tanto linealmente independientes; y toda clase residual de $\bideal$ módulo $h(t)$ puede representarse linealmente por estas $k-p$ clases especiales, con coeficientes en $P$. Pero módulo $h(t)$ el ideal $\bideal$ pasa a un módulo de formas lineales en $1,t,\ldots,t^{k-1}$, cuyo rango es por tanto exactamente $k-p$.

Sea ahora $\aideal=\aideal_1$ un ideal transformado de polinomios, y sea $A^{(1)}(x)$ un polinomio de $\aideal_1$, regular en $x_1$, de grado $k_1$. Sea $\mathfrak A_1$ el módulo de formas lineales en $1,x_1,\ldots,x_1^{k_1-1}$, con polinomios $a^{(2)}(x)$ como coeficientes, al que pasa $\aideal_1$ módulo $A^{(1)}(x)$. Denótese por $\aideal_2$ el ideal en $x_2\ldots x_n$ derivado de todos los determinantes de $k_1$ filas de $\mathfrak A_1$. Por el Teorema XI, $\aideal_2$ es el ideal cero si y sólo si los polinomios de $\aideal_1$ tienen un divisor común, en el sentido polinómico, de grado $p_1>0$ en $x_1$. Si $\aideal_2$ es distinto del ideal cero, entonces, como $\aideal$ se supone transformado, contiene un polinomio $A^{(2)}(x)$ regular en $x_2$, de grado $k_2$, como se ve directamente componiendo la transformación (12) a partir de las transformaciones especiales (25) y (26). Sea $\mathfrak A_2$ de nuevo el módulo de formas lineales en $1,x_2,\ldots,x_2^{k_2-1}$ al que pasa $\aideal_2$ módulo $A^{(2)}(x)$, y $\aideal_3$ el ideal en $x_3\ldots x_n$ derivado de todos los determinantes de $k_2$ filas de $\mathfrak A_2$; este ideal debe contener un polinomio $A^{(3)}(x)$ regular en $x_3$.

De este modo se continúa el proceso hasta alcanzar un ideal cero, o hasta que $\aideal_{n+1}$ se vuelva el ideal unidad; pues, al estar libre de los $x$, $\aideal_{n+1}$ sólo puede ser el ideal cero o el ideal unidad. Se ve que los ideales $\aideal_1,\aideal_2,\ldots$ están determinados unívocamente por $\aideal$, independientemente de los polinomios $A^{(\lambda)}(x)$ escogidos. Esto es claro para $\aideal_1=\aideal$, y por tanto puede suponerse para $\aideal_1,\ldots,\aideal_i$. Si $\aideal_i$ pasa módulo $A^{(i)}(x)$ al módulo de formas lineales $\mathfrak A_i$, entonces $\mathfrak A_i$ puede caracterizarse también como la totalidad de los polinomios de $\aideal_i$ que no alcanzan el grado $k_i$ en $x_i$; así, $\mathfrak A_i$ depende sólo del grado $k_i$ de $A^{(i)}(x)$. Sea $\bar A^{(i)}(x)$ un polinomio de $\aideal_i$, regular en $x_i$, de grado $\bar k_i>k_i$, sea $\overline{\mathfrak A}_i$ el módulo correspondiente de formas lineales, y sea $\bar\aideal_{i+1}$ el ideal derivado de los determinantes de $\bar k_i$ filas de $\overline{\mathfrak A}_i$. Entonces vale la ecuación de módulos
\[
 \overline{\mathfrak A}_i=(\mathfrak A_i,A^{(i)},x_iA^{(i)},\ldots,x_i^{\bar k_i-k_i-1}A^{(i)}).
\]
Como, por la regularidad de $A^{(i)}(x)$ en $x_i$, los polinomios $A^{(i)},x_iA^{(i)},\ldots$ pueden introducirse como nuevas indeterminadas de las formas lineales, se sigue que los determinantes de $k_i$ filas de $\mathfrak A_i$ coinciden con los determinantes de $\bar k_i$ filas de $\overline{\mathfrak A}_i$, y por tanto $\bar\aideal_{i+1}=\aideal_{i+1}$.\footnote{Esto es en principio la misma inferencia en la que descansaba el isomorfismo de $\Gmod^*_{i-1}\mid\Mmod^*_{i-1}$ con $\Gmod_{i-1}\mid\Mmod_{i-1}$.}

Además, $\aideal_{i+1}$ es siempre divisible por $\aideal_i$. Esto es claro si $\aideal_{i+1}$ es el ideal cero; en el caso contrario, $\mathfrak A_i$ tiene rango $k_i$, de modo que $\aideal_{i+1}$ es el ideal que por definición consta de los determinantes de $k_i$ filas de $\mathfrak A_i$, y es divisible por $\mathfrak A_i$ y por tanto por $\aideal_i$. Así todo $\aideal_{i+1}$ es divisible por $\aideal$; si por tanto $\aideal_{n+1}$ se vuelve el ideal unidad, entonces $\aideal$ mismo es el ideal unidad.

Sea ahora $\aideal$ distinto del ideal unidad, con $\aideal_1\ne0,\ldots,\aideal_i\ne0$, y $\aideal_{i+1}=0$. Sea $D^{(i)}(x)$ el máximo común divisor, existente por el Teorema XI, de todos los polinomios de $\aideal_i$, donde divisor se entiende en el sentido polinómico; como divisor de $A^{(i)}(x)$, $D^{(i)}(x)$ es él mismo regular en $x_i$ de grado $p_i>0$. Adjúntese $x_{i+1}\ldots x_n$ a $P(u)$; entonces $D^{(i)}(x)$ tiene, en un cuerpo de extensión algebraico de $P(u;x_{i+1}\ldots x_n)$, una descomposición en $p_i$ factores lineales. Sea $x_i-\bar x_i$ uno de esos factores lineales, de modo que todos los polinomios de $\aideal_i$ se anulan para $x_i=\bar x_i$; entonces, por el Teorema XI, los polinomios de $\aideal_{i-1}$ adquieren bajo esta especialización un máximo común divisor $D^{(i-1)}(x)$. Como divisor de $A^{(i-1)}(x)$, $D^{(i-1)}(x)$ es de nuevo regular en $x_{i-1}$ de grado $p_{i-1}>0$; en particular no puede anularse idénticamente, y en un cuerpo de extensión de $P(u,\bar x_i,x_{i+1},\ldots,x_n)$ se descompone en $p_{i-1}$ factores lineales. Si $x_{i-1}-\bar x_{i-1}$ es uno de esos factores lineales, a él corresponde un máximo común divisor de $\aideal_{i-2}$. Continuando así, se llega a finitos ceros comunes de $\aideal$ situados en un cuerpo de extensión algebraico de $P(u;x_{i+1}\ldots x_n)$. Recíprocamente, como $\aideal_i$ es divisible por $\aideal$, todo cero de $\aideal$ es también cero de $\aideal_i$. Si, en particular, $x_1=\bar x_1,\ldots,x_i=\bar x_i,x_{i+1}=x_{i+1},\ldots,x_n=x_n$ es un cero de $\aideal$, entonces se sigue $\aideal_{i+1}=0$, y los polinomios de $\aideal_i$ adquieren un divisor común con el factor lineal $x_i-\bar x_i$. En resumen:

\medskip
\noindent\textbf{Teorema XII.} \emph{Sea $\aideal$ distinto del ideal unidad, con $\aideal_1\ne0,\ldots,\aideal_i\ne0$ y $\aideal_{i+1}=0$. Entonces, después de adjuntar $x_{i+1}\ldots x_n$ a $P(u)$, existen finitos sistemas de valores asociados $\bar x_1\ldots\bar x_i$, situados en un cuerpo de extensión algebraico de $P(u,x_{i+1},\ldots,x_n)$, tales que todos los polinomios de $\aideal$ se anulan para $x_1=\bar x_1,\ldots,x_i=\bar x_i,x_{i+1}=x_{i+1},\ldots,x_n=x_n$. Recíprocamente, si existe tal cero de $\aideal$, entonces se sigue $\aideal_{i+1}=0$, así como la aparición de un divisor común $D^{(i)}(x)$ de todos los polinomios de $\aideal_i$ que tiene el factor lineal $x_i-\bar x_i$.}

El Teorema XII muestra en particular que todo ideal $\aideal$ que no tenga ceros es el ideal unidad.

Para efectuar una descomposición que exhiba explícitamente también la asociación de los sistemas de valores, introdúzcanse nuevas indeterminadas $v$, que pueden adjuntarse a $P(u,x_{i+1},\ldots,x_n)$. Póngase
\begin{equation}
 z_i=-v_1x_1-\cdots-v_{i-1}x_{i-1}+x_i.
\tag{35}
\end{equation}
Entonces la composición de (12) con (35) es de nuevo una sustitución del tipo (12), en la cual ahora sólo los $u_{ik}$ se reemplazan por $w_{ik}=u_{ik}-v_k$, y más generalmente $u_{i+h,k}$ por $w_{i+h,k}=u_{i+h,k}-u_{i+h,i}v_k$. Si la sustitución (35) lleva el ideal, transformado por hipótesis, $\aideal$ a $\bideal$, entonces $\bideal$ se obtiene de $\aideal$ simplemente reemplazando los $u_{\mu\nu}$ por $w_{\mu\nu}$ y adjuntando los $v$; y por el mismo proceso los ideales $\bideal_1,\bideal_2,\ldots$ definidos por $\bideal$ surgen de $\aideal_1,\aideal_2,\ldots$. Recíprocamente, $\aideal_i$ se obtiene de $\bideal_i$ por $v=0$; por tanto los ideales $\aideal_i$ y $\bideal_i$ son simultáneamente ideales cero o simultáneamente distintos del ideal cero.

Sea ahora de nuevo $\aideal_1\ne0;\ldots;\aideal_i\ne0;\aideal_{i+1}=0$, y por tanto también $\bideal_1\ne0;\ldots;\bideal_i\ne0;\bideal_{i+1}=0$. Sea $\bar x_1\ldots\bar x_i,x_{i+1}\ldots x_n$ un sistema de ceros asociado de $\aideal$. Defínase $\bar z_i=\bar x_i+v_1\bar x_1+\cdots+v_{i-1}\bar x_{i-1}$. Entonces $\bar x_1\ldots\bar x_{i-1},\bar z_i,x_{i+1}\ldots x_n$ es, por (35), un cero de $\bideal$. Por consiguiente, si $H^{(i)}(z,x)$ denota el máximo común divisor de todos los polinomios de $\bideal_i$, que puede suponerse entero y primitivo en los $v$, entonces por la última parte del Teorema XII, $H^{(i)}(z,x)$ tiene el factor
\[
 z_i-\bar z_i=z_i-(\bar x_i+v_1\bar x_1+\cdots+v_{i-1}\bar x_{i-1}),
\]
y a cada cero de $\aideal$ que aparece después de adjuntar $x_{i+1}\ldots x_n$ le corresponde tal factor. Debe mostrarse que con esto se agota la factorización de $H^{(i)}(z,x)$, es decir, que no puede separarse de $H^{(i)}$ ningún factor $H^{(i)}_1$ que no se descomponga en factores lineales en $z$ y $v$. Si tal cosa ocurriera, entonces, por la primitividad supuesta en $v$, $H^{(i)}$ contendría al menos un factor lineal $z_i-\bar z_i$, donde $\bar z_i$ pertenece a un cuerpo de extensión algebraico de $P(u,v,x_{i+1},\ldots,x_n)$. Si $\bar x_1\ldots\bar x_{i-1},\bar z_i,x_{i+1}\ldots x_n$ es el cero correspondiente de $\bideal$, entonces debe coincidir con uno de los finitos ceros de $\aideal$ que aparecen después de adjuntar $x_{i+1}\ldots x_n$, y por tanto debe ser independiente de $v$. De ahí que $\bar z_i$ contenga necesariamente los $v$ linealmente, no de manera algebraica ni racional no lineal; contra la suposición, de $H^{(i)}_1$ puede separarse otro factor lineal $z_i-(\bar x_i+v_1\bar x_1+\cdots+v_{i-1}\bar x_{i-1})$ en $z$ y $v$. Así, la descomposición explícita es
\[
 H^{(i)}(z,x)=\prod \{z_i-(\bar x_i+v_1\bar x_1+\cdots+v_{i-1}\bar x_{i-1})\}^{\alpha_\lambda}.
\]
Como el grado de $H^{(i)}$ y los exponentes individuales $\alpha_\lambda$ deben coincidir con los correspondientes de $D^{(i)}(x)$ -- pues $H^{(i)}$ surge de $D^{(i)}$ por la especialización $u_{\mu\nu}=w_{\mu\nu}$, y $D^{(i)}$ de $H^{(i)}$ por la especialización $v=0$ --, esta descomposición muestra también que, debido a que las indeterminadas $u_{\mu\nu}$ están adjuntadas a $P$, todo cero $\bar x_i$ de $D^{(i)}$ que aparece en la eliminación a partir de $\aideal$ corresponde a un máximo común divisor de $\aideal_{i-1}\ldots\aideal_1$ que es respectivamente lineal, o a una potencia de tal divisor; por tanto en cada caso se obtiene un único sistema de ceros asociado $\bar x_1\ldots\bar x_i,x_{i+1}\ldots x_n$ de $\aideal$.\footnote{Debe señalarse que la eliminación sucesiva dada aquí, a diferencia del método de Kronecker, depende sólo del ideal dado $\aideal$ y es independiente de toda base ideal; en particular, esto vale también para los exponentes $\alpha_\lambda$. Una ejecución completa de la eliminación por este método requeriría, según Hentzelt, la continuación del procedimiento sobre el cociente $\aideal/D^{(i)}$, lo cual puede omitirse en vista de la sección siguiente. (E. N.)}

\section*{§ 7. Forma resultante y teoría de la eliminación.}

Para establecer la conexión entre la forma resultante y la teoría de la eliminación, aplíquese la eliminación sucesiva de §~6 específicamente al cociente $\aideal=\mideal/\gideal_{i-1}$ del ideal por el ideal fundamental de la etapa $(i-1)$. Primero debe mostrarse que este cociente es un ideal transformado. Ahora bien, $\mideal$ es transformado y, por el Teorema VI de §~3, también lo es $\gideal_{i-1}$. De
\[
 H(x)\equiv0(\mideal/\gideal_{i-1});\qquad
 H(x)=\overline H(y)=\sum\alpha_i(u)\overline H_i(y)=\sum\alpha_i(u)H_i(x)
\]
se obtiene
\[
 \overline H(y)\overline\gideal_{i-1,(u)}\equiv0(\overline\mideal_{(u)}),
 \quad\hbox{por tanto también}\quad
 \overline H(y)\overline\gideal_{i-1}\equiv0(\overline\mideal_{(u)}),
\]
o bien
\[
 \overline H_i(y)\overline\gideal_{i-1}\equiv0(\overline\mideal),
 \quad\hbox{por tanto}\quad H_i(x)\gideal_{i-1}\equiv0(\mideal).
\]
Esto prueba la divisibilidad de $H_i(x)$ por $\mideal/\gideal_{i-1}$ y, por tanto, que este ideal es transformado, de modo que se aplica el método de eliminación de §~6.

Pero (30), teniendo en cuenta (28) -- §~4 --, muestra que $C^{(i)}(x)$ es divisible por $\mideal/\gideal_{i-1}$, donde $C^{(i)}(x)$ denota cualquier determinante no idénticamente nulo de grado $\rho$ de $\Mmod^*_{i-1}$. Puesto que para $i>1$ este $C^{(i)}$ tiene grado cero en $x_1$, el módulo $\mathfrak A_1$ correspondiente al ideal $\aideal=\aideal_1$ contiene los polinomios $C^{(i)},x_1C^{(i)},\ldots,x_1^{k_1-1}C^{(i)}$, y por tanto $(C^{(i)})^{k_1}$ es divisible por $\aideal_2$; análogamente $(C^{(i)})^{k_1k_2}$ es divisible por $\aideal_3$, y finalmente $(C^{(i)})^{k_1k_2\cdots k_{i-1}}$ por $\aideal_i$. Así $\aideal_1,\ldots,\aideal_i$ son distintos del ideal cero, como también ocurre para $i=1$. Sea ahora $D^{(i)}(x)$ el máximo común divisor de todos los polinomios de $\aideal_i$, que tiene grado $p_i>0$ sólo cuando $\aideal_{i+1}$ es el ideal cero. Por la definición de $D^{(i)}(x)$, y teniendo en cuenta la divisibilidad de $\aideal_i$ por $\aideal$, se obtiene
\[
 d^{(i+1)}D^{(i)}(x)\equiv0(\mideal/\gideal_{i-1});\qquad d^{(i+1)}\ne0.
\]
Así $d^{(i+1)}D^{(i)}$ es un polinomio, libre de $x_1\ldots x_{i-1}$, que pertenece a $\mideal/\gideal_{i-1}$. Pero $E^{(i)}(x)$ fue definido (Teorema II y §~4), como divisor elemental más alto de $\Mmod_{i-1}$ o de $\Mmod^*_{i-1}$, como el máximo común divisor, en el sentido polinómico, de todos los polinomios de $\mideal/\gideal_{i-1}$ que son libres de $x_1\ldots x_{i-1}$. Por tanto $d^{(i+1)}D^{(i)}$, y debido a la regularidad de $E^{(i)}(x)$ en $x_i$ también $D^{(i)}(x)$, es divisible por $E^{(i)}(x)$.

El mismo argumento puede efectuarse si la transformación (12) se hubiera compuesto desde el principio con (35), es decir, si las indeterminadas $u_{\mu\nu}$ hubieran sido reemplazadas por $w_{\mu\nu}$; esto da la divisibilidad de $H^{(i)}(z,x)$ por el correspondiente $\bar E^{(i)}(z,x)$. Puesto que, así como $D^{(i)}(x)$ y $H^{(i)}(z,x)$ pasan uno en otro por especialización, también lo hacen $E^{(i)}(x)$ y $\bar E^{(i)}(z,x)$, e igualmente $R^{(i)}(x)$ y $\bar R^{(i)}(z,x)$, los números de multiplicidad en la descomposición en factores lineales deben coincidir mutuamente. Tomando además en cuenta que una potencia de $E^{(i)}(x)$ es divisible por $R^{(i)}(x)$, se obtiene el siguiente resumen.

\medskip
\noindent\textbf{Teorema XIII.} \emph{Si $R^{(i)}(x)$ se descompone en un cuerpo de extensión algebraico de $P(u,x_{i+1},\ldots,x_n)$ en factores lineales $x_i-\bar x_i$, entonces $\bar x_i$ puede completarse de una y sólo una manera a un sistema de valores asociado $\bar x_1\ldots\bar x_i,x_{i+1}\ldots x_n$ que es cero de $\mideal/\gideal_{i-1}$ y, por consiguiente, de $\mideal$. Los finitos ceros de $\mideal$ que así surgen de los factores lineales de $R^{(i)}$ -- finitos después de adjuntar $x_{i+1}\ldots x_n$ -- quedan reunidos por la descomposición explícita}
\begin{equation}
 \bar R^{(i)}(z,x)=\prod\{z_i-(\bar x_i+v_1\bar x_1+\cdots+v_{i-1}\bar x_{i-1})\}^{\alpha_\lambda}.
\tag{36}
\end{equation}
\emph{Debido a la regularidad de $\bar R^{(i)}(z,x)$ en $z_i$, esta descomposición se conserva si los $x_{i+1}\ldots x_n$ adjuntados a $P(u)$ se reemplazan por sistemas de valores especiales arbitrarios $\bar x_{i+1}\ldots\bar x_n$ de $P(u)$ o del cuerpo algebraicamente cerrado derivado de él.}

Con esto se consigue también una separación de los ceros de $\mideal$ según los factores individuales de la forma resultante; al número de indeterminadas $x$ adjuntadas a $P(u)$ se le llama también la dimensión del objeto algebraico correspondiente.

Si, en la descomposición de $\bar R^{(i)}(z,x)$ o del correspondiente $R^{(i)}(x)$, se agrupan los factores conjugados sobre $P(u,x_{i+1}\ldots x_n)$ -- factores que para un $P$ general pueden ser total o parcialmente idénticos --, se obtiene una descomposición de $R^{(i)}(x)$ en potencias de polinomios irreducibles sobre $P(u,x_{i+1}\ldots x_n)$:
\[
 R^{(i)}(x)=S^{(i)}_1(x)^{\beta_1}\cdots S^{(i)}_\nu(x)^{\beta_\nu}.
\]
Por el Teorema X, a esta descomposición corresponde una representación $\gideal_i=[\jideal_1\ldots\jideal_\nu]$ tal que $S^{(i)}_\mu(x)^{\beta_\mu}$ es igual a la norma de $\gideal_{i-1}$ respecto del ideal $\jideal_\mu$, respectivamente a la norma del módulo $\Gmod_{i-1}$ respecto del módulo correspondiente a $\jideal_\mu$. Esto aclara el significado de los exponentes $\beta_\mu$: en particular, si $\gamma_\mu$ es el grado de $S^{(i)}_\mu$, entonces la multiplicidad $\beta_\mu\gamma_\mu$ es igual al número de clases residuales de $\gideal_{i-1}$ módulo $\jideal_\mu$ que son linealmente independientes sobre $P(u,x_{i+1}\ldots x_n)$.

Finalmente considérese brevemente el caso especial en que $P$ tiene característica cero. Si los $u_{\mu\nu}$ se especializan a magnitudes $\bar u_{\mu\nu}$ de $P$ de modo que los $n$ determinantes $C^{(i)}(x)$ $(i=1,2,\ldots,n)$ permanezcan regulares en $x_i$,\footnote{Hentzelt presupone esta especialización de los $u_{\mu\nu}$ desde el comienzo y por ello debe recurrir a consideraciones mucho más complicadas para probar la descomponibilidad de $H^{(i)}(z,x)$ y $\bar R^{(i)}(z,x)$ en factores lineales en $z$ y $v$. Los exponentes que aparecen allí no tienen por qué coincidir necesariamente con los anteriores. (E. N.)} entonces se conserva también la regularidad de $R^{(i)}$. Bajo esta especialización, $R^{(i)}$ y asimismo $\bar R^{(i)}(z,x)$ se convierten en un divisor común de todos los determinantes de $\rho$ filas de $\Mmod^*_{i-1}$, aunque no necesariamente en el máximo común divisor. Pero la especialización siempre puede escogerse de modo que también esta última propiedad se conserve. Pues, como muestra la existencia de la base modular, $R^{(i)}(z,x)$ puede caracterizarse también como el máximo común divisor de finitos determinantes de $\rho$ filas de $\Mmod^*_{i-1}$. La descomposición del $\bar R^{(i)}(z,x)$ así especializado se obtiene entonces simplemente reemplazando los $u_{\mu\nu}$ por $\bar u_{\mu\nu}$ en (36), de modo que toda la teoría de la eliminación se conserva.

\begin{center}
(Recibido el 17 de marzo de 1922.)
\end{center}


\clearpage
\section*{23. Invariantes algebraicos y diferenciales}
\begin{center}
\emph{Jahresbericht der Deutschen Mathematiker-Vereinigung} 32 (1923), pp. 177--184
\end{center}

Si debo informar sobre el desarrollo de los invariantes algebraicos y diferenciales, quisiera limitarme, en ambos dominios, al período crítico que, según una observación de Hilbert, sigue al período ingenuo y formal. Este período crítico está caracterizado, para los invariantes algebraicos, por el nombre del propio Hilbert; para los invariantes diferenciales, por el nombre de Riemann, o, en términos sustantivos: para los invariantes algebraicos, por los métodos aritméticos del álgebra, que se desarrollaron esencialmente en torno a estas cuestiones y pudieron mostrar aquí toda su agudeza; para los invariantes diferenciales, por los métodos del cálculo variacional formal. Por ello quiero informar sobre los métodos aritméticos, con su significación más allá del objeto particular, mientras que en la segunda parte puedo limitarme a la subordinación de los invariantes diferenciales a los algebraicos, subordinación que se lleva a cabo precisamente en conexión con los métodos del cálculo variacional.

\textbf{1.} Los \emph{invariantes algebraicos} se basan, como es sabido, en el \emph{grupo lineal general} en $x_1,\ldots,x_n$:
\[
\tag{1}
x_i=\sum_k s_{ik}x'_k
\qquad\text{o, abreviadamente:}\qquad
x=S(x'),
\]
donde los $s_{ik}$ designan indeterminadas, de modo que $\D=|s_{ik}|$ es distinto de cero. Si ahora $f(a,x), g(b,x),\ldots$ son formas en $x$, de dimensiones $\alpha,\beta,\ldots$, con indeterminadas $a,b,\ldots$ como coeficientes, entonces por medio de (1) surge la \emph{transformación inducida} de los $a,b,\ldots$; pues de
\[
\begin{aligned}
f(a,x)&=\sum a_{i_1\ldots i_n}x_1^{i_1}\cdots x_n^{i_n}
      =\sum a'_{j_1\ldots j_n}{x'_1}^{j_1}\cdots {x'_n}^{j_n}
      =f(a',x'),\\
g(b,x)&=\sum b_{r_1\ldots r_n}x_1^{r_1}\cdots x_n^{r_n}
      =\sum b'_{s_1\ldots s_n}{x'_1}^{s_1}\cdots {x'_n}^{s_n}
      =g(b',x')
\end{aligned}
\]
se sigue que los $a',b',\ldots$ se convierten en funciones lineales homogéneas de los $a,b,\ldots$, cuyos coeficientes son de grados $\alpha,\beta,\ldots$ en los $s_{ik}$:
\[
\tag{2}
a'=A_\alpha(a);\qquad b'=B_\beta(b)\ldots .
\]
Por \emph{invariante} se entiende entonces un invariante respecto de las transformaciones (1) y (2); es decir, una función racional entera de las indeterminadas $a,b,\ldots,x,y,\ldots$ ---donde también $y=S(y')$--- que, al aplicar las transformaciones, se reproduce salvo un factor, una potencia del determinante de sustitución:
\[
\tag{3}
I(a',b',\ldots,x',y',\ldots)=\D^\rho I(a,b,\ldots,x,y,\ldots).
\]
Los coeficientes de $I$ pueden suponerse racionales, incluso enteros racionales, pues todo invariante con coeficientes en un cuerpo numérico arbitrario $P$ puede expresarse linealmente, con coeficientes de $P$, por medio de un número finito de tales invariantes especiales. Del mismo modo no hay restricción en suponer que $I$ sea separadamente homogéneo en cada serie, ya que también aquí todo invariante puede componerse como suma de un número finito de tales invariantes especiales.

Debe observarse que, recíprocamente, (1) y (2) quedan de nuevo determinadas por la exigencia (3). Como Ostrowski y Schur han mostrado recientemente,\footnote{A. Ostrowski e I. Schur, \emph{Über eine fundamentale Eigenschaft der Invarianten einer binären Form}, Math. Zeitschr. 15 (1922), y trabajos relacionados, todavía inéditos, de Ostrowski.} las transformaciones inducidas pueden caracterizarse también como la totalidad de las transformaciones lineales, separadas en cada serie individual, que ---fuera de ciertas excepciones especificables--- admiten un invariante arbitrario $I$ libre de las $x,y,\ldots$.

El producto de dos invariantes es de nuevo un invariante; la suma lo es si y sólo si coincide el peso, es decir, el exponente $\rho$ en (3). Pero si uno se restringe a transformaciones de determinante uno, entonces toda suma vuelve a ser un invariante y agota también todos los invariantes respecto del grupo así restringido. Se obtiene por tanto un dominio íntegro; de hecho, un \emph{dominio íntegro homogéneo}, esto es, un dominio en el cual la pertenencia de un polinomio implica la pertenencia separada de sus componentes homogéneas, lo que aquí equivale de nuevo a los invariantes separadamente homogéneos según (3).

Aquí surge el problema fundamental en torno al cual se desarrolló la teoría de invariantes, y el álgebra aritmética en general: \emph{¿es finito este dominio íntegro?}; esto es, ¿posee una \emph{base entera finita} $I_1,\ldots,I_k$ tal que todo invariante pueda representarse de modo racional entero mediante $I_1,\ldots,I_k$, $I=G(I_1,\ldots,I_k)$? La solución de Hilbert ---la demostración de Gordan para el caso binario no admite extensión por sus cálculos simbólicos inmanejables, y la demostración de Mertens--Hilbert tampoco, por su uso de la descomposición de una forma binaria en formas lineales--- descansa en reducir la cuestión de la base entera a la de la base de ideales. Quisiera esbozar aquí el curso de las ideas, subrayando las ideas aritméticas fundamentales, y pasar después a tres círculos de cuestiones conectadas: la construcción efectiva de la base mediante un número finito de pasos; las cuestiones de finitud para subgrupos; y las cuestiones de finitud que tienen en cuenta la integridad.

\textbf{2.} Recuerdo la \emph{definición de ideal} para cuerpos algebraicos numéricos finitos. Un sistema $\mathfrak a$ de números del cuerpo se llama ideal si
\begin{enumerate}
\item $\mathfrak a$ pertenece al dominio $\mathfrak o$ de todos los enteros del cuerpo;
\item junto con $\alpha$ y $\beta$, también la diferencia $\alpha-\beta$ pertenece a $\mathfrak a$;
\item junto con $\alpha$, también $\lambda\alpha$ pertenece a $\mathfrak a$, donde $\lambda$ es cualquier elemento de $\mathfrak o$.
\end{enumerate}
Y, como es sabido, vale el teorema de que todo ideal posee una \emph{base ideal}: $\mathfrak a=(\alpha_1,\ldots,\alpha_k)$; esto es, existen finitos elementos $\alpha_1,\ldots,\alpha_k$ de $\mathfrak a$ tales que, mediante 2. y 3., $\mathfrak a$ se deriva de $\alpha_1,\ldots,\alpha_k$; por tanto, $\alpha=\lambda_1\alpha_1+\cdots+\lambda_k\alpha_k$ agota todos los elementos de $\mathfrak a$.\footnote{Que $k$ pueda reducirse a dos es irrelevante para lo que sigue.}

Esta definición de ideal descansa sólo en el hecho de que $\mathfrak o$ forma un dominio íntegro; por tanto, el concepto de ideal, y asimismo el de base ideal, permanece literalmente igual cuando $\mathfrak o$ se reemplaza por un dominio íntegro arbitrario.

La demostración de finitud de Hilbert se descompone ahora en los tres teoremas siguientes:
\begin{enumerate}
\item En el dominio de todos los polinomios en $n$ indeterminadas con coeficientes en un cuerpo de racionalidad, el teorema de la base ideal vale para todo ideal y, en consecuencia, también para todo sistema arbitrario $\mathfrak S$.
\item Un dominio íntegro homogéneo $\mathfrak J$ de polinomios es finito si y sólo si el teorema de la base ideal vale en $\mathfrak J$.
\item Para el dominio de invariantes $\mathfrak J$, el teorema de la base ideal vale en la forma reforzada de que toda base ideal formada respecto del dominio de todos los polinomios es al mismo tiempo una base ideal en $\mathfrak J$, de modo que, junto a
\[
 I=A_1I_1+\cdots+A_kI_k,
\]
se tiene siempre
\[
 I=i_1I_1+\cdots+i_kI_k,
\]
donde las $A$ designan polinomios en los $a,b,\ldots,x,y,\ldots$, y las $i$ designan invariantes.
\end{enumerate}

La demostración de 1. descansa, en principio, en los mismos argumentos que en el cuerpo algebraico numérico; en ambos casos la existencia de la base ideal se sigue directamente de un teorema general sobre módulos de formas lineales.\footnote{Compárese mi informe comparativo sobre la teoría aritmética de las funciones algebraicas, \emph{Jahresber.} 28 (1920), p. 187 y nota 2), p. 188.} De 1. se sigue directamente la existencia de la base ideal para todo dominio íntegro finito de polinomios; para dominios homogéneos, también la recíproca es fácil de ver. Sólo en 3. se emplean propiedades especiales de los invariantes; en efecto, en Hilbert 3. se sigue de un teorema conocido sobre el proceso $\Omega$, mientras que recientemente E. Fischer ha mostrado que 3. puede obtenerse ya ---sobre la base de otros teoremas generales--- de la propiedad del grupo lineal general de contener, junto con cada transformación, también la contragrediente, respectivamente la conjugada-compleja.\footnote{E. Fischer, \emph{Über die Endlichkeit der Invarianten}. Göttinger Nachrichten 1915, y conferencia de Leipzig.}

\textbf{3.} La \emph{construcción efectiva de la base} descansa en una transformación aritmética adicional de la noción de finitud, recurriendo a la teoría de los \emph{cuerpos de funciones}. Vale lo siguiente:
\begin{enumerate}
\item[4.] Un dominio íntegro $\mathfrak J$ de polinomios es finito si y sólo si en $\mathfrak J$ está contenido un subdominio finito $\mathfrak J'$ tal que $\mathfrak J$ depende de $\mathfrak J'$ de modo algebraicamente entero; entonces toda base entera de $\mathfrak J'$ se completa, mediante un sistema fundamental de estas cantidades algebraicamente enteras, hasta una base de $\mathfrak J$.
\end{enumerate}
\footnote{Hilbert muestra (\emph{Über die vollen Invariantensysteme}, Math. Annalen 42 (1893), §§ 1 y 2) que estos hechos se siguen de la finitud supuesta; recíprocamente, que la finitud se sigue de la dependencia algebraicamente entera se deduce de la existencia de la base racional (E. Noether, \emph{Körper und Systeme rationaler Funktionen}, Math. Annalen 76 (1915), § 12).}
Si, de manera especial, como en el caso de los invariantes, se trata de dominios homogéneos $\mathfrak J$ para los cuales el teorema de la base ideal vale en la forma reforzada 3., entonces tal dominio $\mathfrak J'$ puede derivarse de ciertos polinomios siempre existentes y finitos $I_1,\ldots,I_s$ como base entera, de cuya anulación se sigue la anulación de todos los polinomios de $\mathfrak J$. Este hecho se basa de nuevo en un teorema general de la teoría de ideales:
\begin{enumerate}
\item[5.] A todo ideal $\mathfrak a$ de polinomios corresponde un entero racional ordinario $r$ tal que, para todo ideal $\mathfrak b$ que se anula en todos los ceros de $\mathfrak a$, $\mathfrak b^r$ es divisible por $\mathfrak a$.
\end{enumerate}

En el caso de los invariantes se puede determinar una \emph{cota superior para los grados} de $I_1,\ldots,I_s$, y por tanto también para su número $s$, que depende sólo de los grados de las formas fundamentales consideradas. En este punto se usan nuevamente propiedades especiales de los invariantes; el teorema correspondiente dice que un sistema numéricamente especializado de formas fundamentales posee un invariante no nulo si y sólo si $\D=|s_{ik}|$ depende algebraicamente de manera entera de los coeficientes transformados.

A partir de esta cota superior, K. Hentzelt\footnote{Publicaré próximamente el trabajo procedente de su legado.} ha determinado una \emph{cota superior para los grados de los invariantes de la base entera}, dependiente de nuevo sólo de los grados de las formas fundamentales. Lo consigue dando, para el exponente $r$ de 5., una cota superior que depende sólo del número y del grado más alto de los polinomios de una base ideal de $\mathfrak a$, independiente de los coeficientes de estos polinomios, y computable a partir de los números dados mediante un número finito de pasos. En principio, con ello el problema planteado queda completamente resuelto; ciertamente, la cota dada es bastante elevada. Así, para una forma ternaria cuadrática, donde el discriminante ---esto es, un invariante de grado 3--- constituye ya la base de todos los invariantes libres de $x$, según Hilbert y Hentzelt se llega hasta cifras de millones.

\textbf{4.} En la cuestión de finitud para invariantes de \emph{subgrupos} del grupo lineal general se han alcanzado hasta ahora sólo resultados parciales. El método descansa casi exclusivamente en reducir los invariantes de subgrupos a invariantes simultáneos del grupo lineal general, según el procedimiento usado por Study para el grupo ortogonal. Esto fue realizado por Weitzenböck para los invariantes del movimiento y los invariantes afines, y por Deruyts para los semiinvariantes y los invariantes de cizalladura;\footnote{Para referencias bibliográficas, compárese el artículo enciclopédico de Weitzenböck sobre teoría de invariantes; III E 1.} la cuestión del alcance de este método sigue indecisa. La observación de Fischer mencionada al final de 2. resuelve asimismo el problema de finitud para una serie de subgrupos; en particular, proporciona con ello la demostración más simple para el grupo ortogonal. Pero el ejemplo de los semiinvariantes muestra ---como Fischer ha señalado recientemente\footnote{Conferencia de Leipzig y Math. Zeitschrift.}--- que, para subgrupos, pese a la finitud, no tiene por qué cumplirse la forma reforzada 3. de la base ideal.

Interpretada en el sentido de la teoría de cuerpos, la cuestión de finitud que aquí se presenta conduce al problema de Hilbert de las \emph{funciones relativamente enteras}, es decir, a la cuestión de si son siempre finitos aquellos dominios íntegros que constan de la totalidad de los polinomios de un cuerpo de funciones racionales. En esta dirección se sitúa el hecho, demostrable elementalmente, de que para los invariantes de grupos finitos puede darse directamente una base entera distinguida: el sistema de coeficientes de la resolvente de Galois.\footnote{E. Noether, \emph{Der Endlichkeitssatz der Invarianten endlicher Gruppen}. Math. Ann. 77 (1916).}

\textbf{5.} La cuestión de si el teorema de finitud para invariantes puede reforzarse respecto de la \emph{integridad} tampoco está aún resuelta en general. Es verdad que, como mostró Hilbert, el teorema de la base ideal vale también para el dominio de todos los polinomios enteros; y asimismo permanece válida la conexión 2. entre base ideal y base entera;\footnote{De esta conexión se sigue en particular que los teoremas de descomposición de la teoría general de ideales valen para todos los dominios íntegros finitos, como desarrollé en Math. Ann. 83 (1921).} pero falla la demostración de la versión reforzada 3., que en realidad no es necesaria. Para el caso de los invariantes binarios puede demostrarse la finitud entera;\footnote{E. Noether, \emph{Die Endlichkeit des Systems der ganzzahligen Invarianten binärer Formen}. Göttinger Nachrichten 1919.} la demostración, que representa un reforzamiento de integridad de la demostración de finitud binaria de Mertens--Hilbert, usa, además del teorema sobre la base ideal entera, también la descomposición de formas binarias en formas lineales, y por tanto no permite transferencia a más indeterminadas. Sobre fundamentos semejantes se sigue, en cambio, la finitud entera para los invariantes de grupos finitos como reforzamiento de la demostración antes mencionada, aunque ahora de nuevo sólo resulta una prueba de existencia y no una especificación efectiva de la base.\footnote{§ 4 del trabajo citado bajo 11).} También aquí sólo un desarrollo ulterior de la teoría de cuerpos de funciones y de la teoría general de ideales conducirá presumiblemente a la meta.

\textbf{6.} Paso ahora a los \emph{invariantes diferenciales} y a la cuestión planteada en la introducción sobre su reducción a la teoría de invariantes algebraicos. El grupo subyacente está definido, como es sabido, por medio de $x_i=x_i(y)$:
\[
\tag{4}
\dd x_i=\sum_k \frac{\partial x_i}{\partial y_k}\dd y_k;\qquad
\dd\delta x_i=\sum_k \frac{\partial x_i}{\partial y_k}\dd\delta y_k+
\sum_{k,l}\frac{\partial^2 x_i}{\partial y_k\partial y_l}\dd y_k\,\delta y_l;\ \ldots .
\]
Para pasar a las transformaciones inducidas por (4), se puede partir de una función arbitraria $f(x,\dd x)=\varphi(y,\dd y)$; al exigir la invariancia de $\delta f,\delta^2f,\ldots$ respecto de (4), se obtienen las transformaciones de las derivadas de $f$ respecto de $x$ y $\dd x$. Si uno se restringe, por ejemplo, a formas homogéneas en las $\dd x$:
\[
 f(x,\dd x)=\sum a_{i_1\ldots i_n}(x)\dd x_1^{i_1}\cdots\dd x_n^{i_n}
 =\sum a'_{i_1\ldots i_n}(y)\dd y_1^{i_1}\cdots\dd y_n^{i_n}
 =\varphi(y,\dd y),
\]
entonces las $a'$ son lineales en las $a$, las $\frac{\partial a'}{\partial y}$ son lineales en las $a$ y $\frac{\partial a}{\partial x}$, y así sucesivamente:
\[
\tag{5}
a'=A_0(a),\qquad
\frac{\partial a'}{\partial y}=A_1\!\left(a,\frac{\partial a}{\partial x}\right),\qquad
\frac{\partial^2 a'}{\partial y^2}=A_2\!\left(a,\frac{\partial a}{\partial x},\frac{\partial^2 a}{\partial x^2}\right),\ldots
\]
y se trata de la cuestión de la invariancia respecto de las transformaciones (4) y (5). Todas las cantidades que aparecen aquí,
\[
a,\ \frac{\partial a}{\partial x},\ldots,\ \dd x,\ \delta x,\ \dd\delta x,\ldots,\ 
\frac{\partial x}{\partial y},\ \frac{\partial^2 x}{\partial y^2},\ldots,
\]
han de considerarse indeterminadas, de modo que el determinante de cada transformación individual es ciertamente distinto de cero. Sólo cuando la teoría formalmente completada se aplica a funciones especiales surge la cuestión de restringirse a tales funciones para las que las transformaciones permanezcan reversibles de modo unívoco en cierto dominio y exista un número suficiente de cocientes diferenciales, exactamente igual que en el caso algebraico, donde bajo especialización numérica debe suponerse distinto de cero el determinante $|s_{ik}|$.

Con este tratamiento de todos los argumentos que aparecen como indeterminadas, se trata por tanto de un grupo lineal especial, no accesible directamente a los métodos algebraicos ordinarios, en el cual las transformaciones de $\dd x$ y de $a'=A_0(a)$ coinciden con (1) y con la inducida (2). Surge la cuestión de si, quizá de modo enteramente general, mediante la \emph{introducción de nuevos argumentos}, las transformaciones (4) y (5) pueden reducirse al grupo lineal general (1) y a las transformaciones (2) \emph{inducidas} por él.

Esto es, en efecto, lo que ocurre. Los diferenciales superiores $\dd^2x,\delta\dd x,\ldots$ deben reemplazarse por las derivadas de Lagrange que surgen del problema variacional correspondiente a $f(x,\dd x)$, y por otras combinaciones formadas a partir de ellas mediante polarización ---conexiones o transmisiones en el sentido de Weyl y Schouten---, todas las cuales son cogredientes con las $\dd x$, de modo que su transformación está dada por el grupo lineal general (1). Las cantidades $\frac{\partial a}{\partial x},\frac{\partial^2a}{\partial x^2},\ldots$ (o, más generalmente, las derivadas de la $f$, supuesta homogénea, respecto de $x$ y $\dd x$), en cambio, deben reemplazarse por los coeficientes de ciertas formas que surgen de las ``formas normales de las variaciones superiores'' mediante la eliminación indicada de las $\dd^2x,\delta\dd x,\ldots$, y de sus ``derivadas covariantes''
\[
 [\Omega_i],\quad [\Omega_i^{(1)}],\quad [\Omega_i^{(2)}],\ldots\qquad (i=1,2,\ldots);
\]
y como los $[\Omega_i^{(k)}]$ son invariantes respecto de (4) y (5) que dependen sólo de las $\dd x$ y $\delta x$, estos coeficientes satisfacen las transformaciones algebraicamente inducidas (2).\footnote{E. Noether, \emph{Invarianten beliebiger Differentialausdrücke}. Göttinger Nachrichten 1918.} Por este \emph{teorema de reducción}, la teoría de los invariantes diferenciales queda subordinada a la teoría algebraica de invariantes.

La serie de los $[\Omega_i]$ se interrumpe con $[\Omega_\alpha]$ si $f(x,\dd x)$ es una forma homogénea de dimensión $\alpha$ en las $\dd x$. Para las formas cuadráticas, $[\Omega_2]$ se hace idéntica a la forma de curvatura de Riemann, y el proceso de formación indicado aquí es exactamente el que Riemann dio en el ensayo premiado latino, y que Lipschitz, independientemente de esto y en conexión con la lección de habilitación de Riemann, desarrolló también.\footnote{Compárese también mi exposición sobre Riemann y Lipschitz en el núm. 28 de la segunda parte del artículo enciclopédico de Weitzenböck mencionado bajo 7).} La demostración del teorema de reducción, que Christoffel y Ricci habían realizado mediante cálculo para las formas diferenciales cuadráticas, descansa en la introducción de las ``coordenadas normales de Riemann'', que transforman las extremales que parten de un punto en líneas rectas; como de este modo a una transformación arbitraria de las $x$ corresponde una transformación lineal de las coordenadas normales, éstas determinan el paso al grupo lineal general.

Surge la cuestión de si existe también tal teorema de reducción cuando, siguiendo a Weyl y Schouten, el grupo no se define por un problema variacional sino sólo por una conexión; es decir, cuando, en correspondencia con las derivadas de Lagrange, se establecen expresiones $\psi_i(\dd\delta x,\dd x,\delta x)$ por medio de las cuales se efectúa la eliminación de los diferenciales superiores, añadiéndose la restricción ---innecesaria en sí misma--- de que las $\psi_i$ sean lineales en las $\dd x$, por analogía con las formas diferenciales cuadráticas (para una $f(x,\dd x)$ homogénea arbitraria, las derivadas de Lagrange polarizadas son homogéneas de primer orden sólo en $\dd x$). Por la exigencia de que $\psi$ sea cogrediente con las $\dd x$, queda inducida la transformación de sus coeficientes y, por tanto, en analogía con (5), queda fijado el grupo. Para invariantes de los órdenes más bajos, Weitzenböck ha confirmado por cálculo el teorema de reducción en el caso de Weyl; falta todavía un planteamiento general.\footnote{Compárense: Weyl, \emph{Raum, Zeit, Materie}; Schouten, Math. Zeitschr. 13 (1922); Weitzenböck, Sitzungsber. d. Wiener Akademie, 1920 y 1921.}

\begin{center}
(Recibido el 26 de octubre de 1922.)
\end{center}


\clearpage
\setcounter{footnote}{0}


\clearpage

\providecommand{\frakS}{\mathfrak S}
\providecommand{\frakM}{\mathfrak M}
\providecommand{\frakG}{\mathfrak G}
\providecommand{\frakm}{\mathfrak m}
\providecommand{\frakn}{\mathfrak n}
\providecommand{\frakp}{\mathfrak p}
\providecommand{\frakq}{\mathfrak q}
\providecommand{\frakr}{\mathfrak r}
\providecommand{\fraka}{\mathfrak a}
\providecommand{\frakb}{\mathfrak b}
\providecommand{\frakc}{\mathfrak c}
\providecommand{\frakd}{\mathfrak d}
\providecommand{\frake}{\mathfrak e}
\providecommand{\frakg}{\mathfrak g}
\providecommand{\fraks}{\mathfrak s}
\providecommand{\frako}{\mathfrak o}
\providecommand{\frf}{\mathfrak f}
\providecommand{\frK}{\mathfrak K}
\providecommand{\frR}{\mathfrak R}
\providecommand{\barfraka}{\overline{\mathfrak a}}
\providecommand{\barfrakb}{\overline{\mathfrak b}}
\providecommand{\barfrakc}{\overline{\mathfrak c}}
\providecommand{\barfrakm}{\overline{\mathfrak m}}
\providecommand{\calP}{\mathcal P}
\providecommand{\calR}{\mathcal R}
\providecommand{\calN}{\mathcal N}
\providecommand{\calG}{\mathcal G}
\providecommand{\calH}{\mathcal H}
\providecommand{\barr}{\overline}
\providecommand{\dd}{\mathrm d}
\providecommand{\tagg}[1]{\tag{#1}}


\setcounter{footnote}{0}

\section*{24. Teoría de la eliminación y teoría general de ideales}
\begin{center}
Math. Ann. 90 (1923), pp. 229--261
\end{center}

En lo que sigue se trata de situar la teoría de la eliminación ---en la forma aritmética que di a la exposición de Hentzelt\footnote{Kurt Hentzelt, \emph{Zur Theorie der Polynomideale und Resultanten}. Editado por E. Noether. Math. Ann. 88 (1922), pp. 53--79; citado como H.-N.} y que puede describirse como una teoría de ideales en el dominio polinómico--- dentro de la teoría general de ideales,\footnote{E. Noether, \emph{Idealtheorie in Ringbereichen}, Math. Ann. 83 (1921), pp. 24--66; citado como \emph{Idealtheorie}.} y, al mismo tiempo, de dar una nueva fundamentación de las partes referentes a los ceros. Los conceptos básicos de ambas teorías se reúnen brevemente en el § 1, junto con los de la teoría de cuerpos,\footnote{E. Steinitz, \emph{Algebraische Theorie der Körper}, J. f. M. 137 (1910), pp. 167--309; citado como Steinitz.} para poder remitir a ellos aquí.

La clasificación y la nueva fundamentación se agrupan en torno a cuatro cuestiones: la formulación aritmética del concepto de dimensión; la teoría de los ceros; la conexión entre los teoremas de descomposición para normas y divisores elementales y los de la teoría general de ideales; y la caracterización de los ideales primos y primarios por medio de la norma y de la forma de divisores elementales.

La formulación aritmética del \emph{concepto de dimensión} (§ 4) viene dada por cadenas de ideales primos, el equivalente aritmético del hecho de que sobre las superficies hay curvas y sobre las curvas hay puntos. Esta formulación aritmética, que se prueba idéntica a la definición por parámetros de la teoría de la eliminación, puede trasladarse con una leve modificación a dominios de anillos arbitrarios y, junto con la transferencia de ciertas partes de las otras cuestiones, proporciona allí una visión exacta de la estructura de los ideales, como mostraré en otra parte.

En H.-N. la cuestión de los \emph{ceros} se aborda, como de costumbre, directamente para el ideal dado, a saber, remontándose a la eliminación sucesiva, con lo cual no se evita por completo el carácter de verificación. Frente a ello está el camino que se sigue siempre para polinomios en una variable: se representa el polinomio dado como producto de potencias de funciones primas (funciones primarias) y se construye, para cada función prima individual, el cuerpo de ceros como un cuerpo isomorfo al cuerpo de clases residuales. El grado de la función prima da el grado de este cuerpo; el grado de la función primaria, en cambio, cuyos ceros coinciden con los de la función prima, da el grado del anillo de clases residuales módulo esta función primaria. Si, para cada función prima, se pasa al cuerpo de Galois, se obtiene la descomposición en factores lineales; y con ello queda resuelta, para el polinomio inicialmente dado, la cuestión de los ceros, de la descomposición en factores lineales y de la multiplicidad. De modo completamente correspondiente, aquí (§ 3) se amplía el sistema de clases residuales módulo un ideal primo, mediante formación de cocientes, al cuerpo de clases residuales, y se construye un cuerpo de ceros isomorfo a él. Este cuerpo de ceros contiene un subcuerpo generado por adjunción de indeterminadas ---digamos $x_{i+1},\ldots,x_n$---; el grado de la norma da el grado con respecto a este subcuerpo; al mismo tiempo, al pasar al cuerpo de Galois, se obtiene la descomposición de la forma de divisores elementales y por tanto de la norma en factores lineales. Para el ideal primario correspondiente, sin embargo, los ceros son los mismos; también queda dada la descomposición, puesto que la norma se convierte en una potencia de la forma de divisores elementales del ideal primo (§ 2); el grado de la norma da el grado del anillo de clases residuales módulo el ideal primario con respecto al subcuerpo derivado de las indeterminadas.

Así, la cuestión de los ceros de un ideal arbitrario se reduce a la conexión entre los teoremas de descomposición para normas y divisores elementales y los de la teoría general de ideales (§ 5). El hecho de que los ceros del ideal se compongan de los de sus ideales primos asociados individuales corresponde a la descomposición de la norma en los mayores factores primarios, que resultan iguales a potencias de la forma de divisores elementales de los ideales primos asociados individuales;\footnote{Este paralelismo entre teoría de la eliminación y teoría general de ideales no se cumple en la teoría de la eliminación de Kronecker; compárese H.-N., nota *).} con ello se efectúa la descomposición de la norma en factores lineales en general. La multiplicidad, por su parte, recibe en H.-N. su interpretación por medio de los ideales fundamentales y sus componentes; el grado de un mayor factor primario de la norma se hace igual al número de clases residuales linealmente independientes ---después de adjuntar $x_{i+1},\ldots,x_n$--- del ideal fundamental de etapa $(i-1)$ módulo una componente del ideal fundamental de etapa $i$. Ahora se prueba que el ideal fundamental de etapa $(i-1)$ es idéntico a la componente aislada de la descomposición definida unívocamente por todos los ideales primos de dimensión superior a la $(n-i)$-ésima; y que una componente del ideal fundamental de etapa $i$ es idéntica a la componente aislada de la descomposición definida por el ideal primo correspondiente al factor de la norma y por todos los ideales primos de dimensión superior. El grado del factor correspondiente de la norma se vuelve ---después de adjuntar $x_{i+1},\ldots,x_n$--- igual al grado de aquel subanillo de las clases residuales de esta componente que consiste sólo en clases del ideal fundamental de etapa $(i-1)$.

La caracterización de los ideales primos y de los ideales primarios propios (§ 6) resulta como consecuencia directa de la consideración de los cuerpos de clases residuales. Si el dominio de coeficientes original es un cuerpo perfecto, entonces las funciones primas corresponden a los ideales primos como norma y forma de divisores elementales, las funciones primarias propias a los ideales primarios propios, y recíprocamente. Para un cuerpo imperfecto como dominio de coeficientes sólo pueden darse condiciones que son necesarias o suficientes; como muestran los ejemplos, la norma de un ideal primo puede hacerse propiamente primaria, y la forma de divisores elementales de un ideal primario propio puede hacerse una función prima. Más precisamente, en característica $p$ la norma de un ideal primo se convierte en la $p^f$-ésima potencia ($f>0$) de su forma de divisores elementales; mientras que, para ideales primarios propios cuya forma de divisores elementales es una función prima, la norma puede convertirse en una potencia arbitraria, como muestran de nuevo los ejemplos. Finalmente se consideran los ideales primos absolutos (§ 7), es decir, aquellos que permanecen ideales primos al ampliar el dominio de coeficientes a un cuerpo algebraicamente cerrado. Para ideales primos absolutos con coeficientes de un cuerpo numérico algebraico (finito), la caracterización conduce a trasladar un teorema enunciado primero por A. Ostrowski para funciones primas absolutas: la propiedad de ser un ideal primo se conserva módulo todo ideal primo del cuerpo numérico, salvo a lo sumo finitas excepciones.

Señalemos también la analogía de la última cuestión con la teoría de ideales en cuerpos de números algebraicos. Como forma de divisores elementales de tal ideal hay que considerar el menor número racional entero divisible por el ideal (nota 10), que para un ideal primo se convierte por tanto siempre en un número primo; mientras que, recíprocamente, un ideal se convierte en ideal primo tan pronto como su norma es un número primo. Las demostraciones discurren también de modo enteramente paralelo. Los ejemplos mencionados muestran así que, sobre un cuerpo imperfecto, aparece el análogo de ideales primos de grado superior y de ideales de ramificación, mientras que sobre un cuerpo perfecto sólo hay ideales primos de primer grado y no hay ideales de ramificación.

\subsection*{§ 1. Conceptos básicos de la teoría de la eliminación, de la teoría general de ideales y de la teoría de cuerpos}

\textbf{1. El dominio.} Sea $\frakS$ el dominio íntegro de todos los polinomios en $y_1,\ldots,y_n$ con coeficientes de un cuerpo $P$ definido abstractamente. Sometamos las $y$ a una transformación cuyos coeficientes son indeterminadas $u_{\mu\nu}$:\footnote{Con vistas al § 3 y siguientes, la transformación es algo más general que en H.-N.; todos los resultados de allí se conservan con mayor razón.}
\begin{align}
 y_1&=u_{11}x_1+\cdots+u_{1n}x_n,\quad \ldots,\quad
 y_n=u_{n1}x_1+\cdots+u_{nn}x_n, \tag{1}
\end{align}
con inversa
\begin{align}
 x_1&=t_{11}y_1+\cdots+t_{n1}y_n,\quad \ldots,\quad
 x_n=t_{1n}y_1+\cdots+t_{nn}y_n. \tag{1'}
\end{align}
Las $u_{\mu\nu}$, o ---lo que es equivalente por la expresabilidad racional mutua--- las $t_{\mu\nu}$, se adjuntan a $P$, de modo que surge el dominio de coeficientes $P(u)=P(t)$; sea $\frakS'$ el dominio íntegro de todos los polinomios en $x_1,\ldots,x_n$ con coeficientes de $P(u)$. Los dominios $\frakS$ y $\frakS'$ son la base de la teoría de la eliminación; se consideran ideales $\fraka,\frakb,\ldots$ en $\frakS$ y $\fraka',\frakb',\ldots$ en $\frakS'$. Un ideal en un anillo arbitrario se define aquí de la manera usual por la exigencia de que, junto con $\alpha$ y $\beta$, contenga también $\alpha-\beta$, y, junto con $\alpha$, también $\lambda\alpha$, donde $\lambda$ es un elemento arbitrario del anillo; $\frako$ designará siempre el ideal unidad consistente en todos los elementos.

\textbf{2. Ideales transformados.} Un ideal $\frakm$ en $\frakS'$ se llama transformado si surge de un ideal $\barfrakm$ en $\frakS$ mediante (1) después de adjuntar las $u_{\mu\nu}$, o las $t_{\mu\nu}$. Así, cuando $\bar f(y)$ o $\bar g(y)$ recorre todos los polinomios de $\barfrakm$, $\frakm$ consiste en todas las combinaciones lineales
\[
 f(x)=\sum U_i(u)\bar f_i(y)=\sum U_i(u)f_i(x)
\]
o también
\[
 g(x)=\sum T_i(t)\bar g_i(y)=\sum T_i(t)g_i(x).
\]
En particular están contenidos todos los $f_i(x)=\bar f_i(y)$, o, lo que equivale a lo mismo, todos los $g_i(x)=\bar g_i(y)$. Puesto que $f(x)$, respectivamente $g(x)$, sólo quedan fijados salvo cantidades de $P(u)=P(t)$, los $U_i,T_i$ pueden suponerse, sin pérdida de generalidad, productos de potencias de las $u$, respectivamente de las $t$. Los polinomios $f(x),g(x)$ pasan de nuevo a polinomios de $\frakm$ si $U,T$ se reemplazan por productos de potencias arbitrarios, en particular al permutar las columnas de (1), respectivamente las filas de (1'); por tanto $\frakm$ contiene, junto con cualquier polinomio, todos los que se obtienen de él por permutar las $x$, siempre que se realicen las permutaciones correspondientes en los coeficientes dependientes de $u$, respectivamente de $t$. Todos los ideales que aparezcan en lo sucesivo se supondrán transformados salvo indicación expresa en contrario. Los ideales no transformados pasan, al componer (1) con
\[
 x_i=v_{i1}z_1+\cdots+v_{in}z_n,
\]
a ideales transformados en el dominio polinómico de las $z$ con coeficientes de $P(u,v)$, mientras que para ideales transformados esta composición se reduce a reemplazar las $u_{\mu\nu}$ por combinaciones bilineales de las $u,v$.

\textbf{3. Notación.} Por $f^{(i)},g^{(i)},\ldots,a^{(i)},b^{(i)},\ldots$ entenderemos siempre polinomios de $\frakS'$ libres de $x_i,\ldots,x_{i-1}$.

\textbf{4. Ideales fundamentales.} A cada ideal $\frakm$ se asignan $n$ ideales fundamentales $\frakg_0,\frakg_1,\ldots,\frakg_{n-1}$ mediante la siguiente estipulación: el ideal fundamental de etapa $(i-1)$, $\frakg_{i-1}$, contiene todos y sólo aquellos polinomios $G(x)$ para los cuales existe un polinomio $b^{(i)}\ne0$ ---en general dependiente de $G(x)$--- tal que $b^{(i)}G(x)\equiv0\pmod{\frakm}$. De la existencia de la base ideal se sigue también la existencia de al menos un $B^{(i)}$, fijo para todos los $G(x)$, de modo que
\begin{equation}
 B^{(i)}\frakg_{i-1}\equiv0\pmod{\frakm},\qquad B^{(i)}\ne0 . \tag{2}
\end{equation}
Los ideales fundamentales de ideales transformados vuelven a ser ideales transformados (H.-N., Teorema VI). El ideal fundamental de etapa $i$, $\frakg_i$, se define también como el ideal al que pasa $\frakm$ cuando $x_{i+1},\ldots,x_n$ se adjuntan a $P(u)$, siempre que uno se restrinja ---lo cual entonces no supone pérdida de generalidad--- a polinomios enteros en $x_{i+1},\ldots,x_n$; $\frakg_0$ es siempre igual al ideal unidad $\frako$.

\textbf{5. Representación modular de ideales, divisores elementales y normas individuales.} Si todo polinomio $f(x)$ se considera como forma lineal en los productos de potencias $\xi$ de $x_1,\ldots,x_{i-1}$, con polinomios $a^{(i)}$ como coeficientes, entonces el ideal $\frakm$ pasa a un módulo $\frakM_{i-1}$ de formas lineales respecto del dominio de los $a^{(i)}$; es decir, $\frakM_{i-1}$ contiene, junto con dos formas lineales cualesquiera, también su diferencia, y, junto con una forma lineal $l(\xi)$, también $a^{(i)}l(\xi)$ para un $a^{(i)}$ arbitrario. Como módulo fundamental de tal módulo $\frakM$ se designa la totalidad de las formas lineales $g(\xi)$ para las que $b^{(i)}g(\xi)\equiv0\pmod{\frakM}$, $b^{(i)}\ne0$; por tanto el ideal fundamental $\frakg_{i-1}$ pasa, bajo la representación modular, al módulo fundamental $\frakG_{i-1}$ de $\frakM_{i-1}$. $\frakM_{i-1}$ y $\frakG_{i-1}$ dependen de infinitas indeterminadas $\xi$, de tal modo que en cada forma lineal individual sólo aparecen finitas de estas indeterminadas. $\frakM_{i-1}$ tiene la propiedad característica de que, después de adjuntar $x_{i+1},\ldots,x_n$ a $P(u)$, posee sólo finitos divisores elementales distintos de la unidad; es decir, después de esta adjunción $\frakG_{i-1}$ y $\frakM_{i-1}$ admiten una representación de base ---donde $\zeta$ denota nuevas indeterminadas, relacionadas con las $\xi$ mediante transformaciones lineales invertibles enteras en $x_i$, de manera que en $\zeta_i=r_i(\xi)$ sólo aparecen cada vez finitas de estas indeterminadas---
\begin{align}
\frakG_{i-1}&=(\zeta_0,\zeta_1,\ldots,\zeta_\sigma,\zeta_{\sigma+1},\ldots,\zeta_{\sigma+\nu},\ldots), \notag\\
\frakM_{i-1}&=(E_0^{(i)}\zeta_0,E_1^{(i)}\zeta_1,\ldots,E_\sigma^{(i)}\zeta_\sigma,
\zeta_{\sigma+1},\ldots,\zeta_{\sigma+\nu},\ldots), \tag{3}
\end{align}
donde cada $E_\mu^{(i)}$ divide al precedente (H.-N. (24) y (28)).

El divisor elemental más alto $E^{(i)}$ se define también como el máximo común divisor ---en el sentido polinómico--- de todos los $B^{(i)}$ que aparecen en (2); puede suponerse entero y primitivo en $x_{i+1},\ldots,x_n$, y entonces es regular en $x_i$, es decir, $E^{(i)}$ contiene el término $x_i^\lambda$ con coeficiente no nulo si es de grado $\lambda$ en las $x$.

El producto $R^{(i)}$ de todos los divisores elementales que aparecen en (3) se llama norma de $\frakG_{i-1}$ respecto de $\frakM_{i-1}$, o también norma individual del ideal $\frakm$; en símbolos, teniendo en cuenta que $\frakm$ pasa a $\frakg_i$ al adjuntar $x_{i+1},\ldots,x_n$:
\begin{equation}
 R^{(i)}=E_0^{(i)}E_1^{(i)}\cdots E_\sigma^{(i)}
 =N(\frakG_{i-1}\mid\frakM_{i-1})=N(\frakg_{i-1}\mid\frakg_i).
\end{equation}
Como muestra (3), después de adjuntar $x_{i+1},\ldots,x_n$, $R^{(i)}$ se hace igual al determinante de la sustitución de paso de $\frakG_{i-1}$ a $\frakM_{i-1}$, y el grado de $R^{(i)}$ da el número de clases residuales, linealmente independientes sobre $P(u,x_{i+1},\ldots,x_n)$, de $\frakG_{i-1}$ módulo $\frakM_{i-1}$, por tanto de $\frakg_{i-1}$ módulo $\frakg_i$; $R^{(i)}$ puede suponerse entero y primitivo en $x_{i+1},\ldots,x_n$, y entonces es regular en $x_i$. $R^{(i)}$ puede calcularse como el máximo común divisor ---en el sentido polinómico--- de todos los determinantes de $\varrho$ filas de un módulo $\frakM^*_{i-1}$, siempre existente, de rango $\varrho$, con la propiedad de que el sistema de clases residuales $\frakG^*_{i-1}/\frakM^*_{i-1}$ es isomorfo a $\frakG_{i-1}/\frakM_{i-1}$, donde $\frakG^*_{i-1}$ se entiende como el módulo fundamental de $\frakM^*_{i-1}$ (H.-N., Teorema VII y § 5).

\textbf{6. Forma de divisores elementales y norma (forma resultante) de los ideales.} El producto $E_\frakm=E^{(1)}E^{(2)}\cdots E^{(n)}$ de todos los divisores elementales más altos se llama forma de divisores elementales, y el producto $R_\frakm=R^{(1)}R^{(2)}\cdots R^{(n)}$ de todas las normas individuales se llama norma (forma resultante)\footnote{En H.-N. se usa sólo la denominación forma resultante; la denominación norma corresponde a las propiedades aritméticas.} de $\frakm$. La forma de divisores elementales y la norma son divisibles por el ideal; la norma es divisible por la forma de divisores elementales, y una potencia de esta última por la norma. Más generalmente se tiene todavía:
\begin{align}
 E^{(i)}\frakg_{i-1}&\equiv0\pmod{\frakg_i},
& E^{(n)}\cdots E^{(i)}\frakg_{i-1}&\equiv0\pmod{\frakm},\notag\\
 R^{(i)}\frakg_{i-1}&\equiv0\pmod{\frakg_i},
& R^{(n)}\cdots R^{(i)}\frakg_{i-1}&\equiv0\pmod{\frakm}, \tag{4}
\end{align}
(H.-N., Teorema VIII), y además: si $\frakm$ es divisible por $\frakn$, y las normas $R_\frakm$ y $R_\frakn$ coinciden, entonces también coinciden los ideales $\frakm$ y $\frakn$ (H.-N., Teorema IX). Teniendo en cuenta que, para un divisor de $\frakm$, la coincidencia de $R^{(1)},R^{(2)},\ldots$ implica la coincidencia de los ideales fundamentales $\frakg_1,\frakg_2,\ldots$, y que siempre $\frakg_0=\frako$, se obtiene correspondientemente: si $\frakn$ es un divisor propio de $\frakm$, entonces la primera norma individual de $\frakn$ que difiere de la correspondiente de $\frakm$ se convierte en un divisor propio.\footnote{En cambio, los factores posteriores de $R_\frakn$ pueden convertirse en múltiplos de los factores correspondientes de $R_\frakm$; por ejemplo siempre que $\frakm$ es un ideal primo y $\frakn$ no es el ideal unidad.} Si $\frakm$ contiene un polinomio $G^{(i)}\ne0$, entonces por definición $\frakg_0,\frakg_1,\ldots,\frakg_{i-1}$ son iguales al ideal unidad; en consecuencia, por la propiedad de las normas individuales indicada en 5, a saber, determinar el número de clases residuales linealmente independientes, $R^{(1)},\ldots,R^{(i-1)}$ se hacen iguales a la unidad; por tanto también $E^{(1)},\ldots,E^{(i-1)}$ se hacen iguales a la unidad.

\textbf{7. Descomposición de las normas individuales y de los divisores elementales; componentes de los ideales fundamentales.} Por función prima se entiende un polinomio irreducible con respecto a $P(u)$; por función primaria, una potencia de una función prima. Una función primaria propia es aquella que no es al mismo tiempo función prima, de modo que intervienen potencias superiores a la primera. Una descomposición de un polinomio en mayores factores primarios es una descomposición en la que cada factor es una función primaria, pero el producto de dos factores cualesquiera ya no es primario.\footnote{La definición de función prima y función primaria también podría formularse en exacta analogía con 11, reemplazando sólo la palabra ideal por polinomio. Los mayores factores primarios son el análogo de las mayores componentes primarias de la descomposición en 12.} A la descomposición de la norma individual $R^{(i)}=N(\frakg_{i-1}\mid\frakg_i)$ en mayores factores primarios $Q_\nu^{(i)}$ corresponde una representación de $\frakg_i$ como mínimo común múltiplo
\[
 \frakg_i=[\frakr_1,\frakr_2,\ldots,\frakr_s],
\]
tal que $Q_\nu^{(i)}\frakg_{i-1}\equiv0\pmod{\frakr_\nu}$; aquí $\frakr_\nu$ posee a $\frakg_{i-1}$ como su ideal fundamental de etapa $(i-1)$, y coincide con su ideal fundamental de etapa $i$; los $\frakr_\nu$ se llaman componentes de los ideales fundamentales. El divisor elemental más alto de $\frakr_\nu$ se hace igual al máximo común divisor ---en el sentido polinómico--- de $Q_\nu^{(i)}$ y $E^{(i)}$, por tanto igual a un mayor factor primario de $E^{(i)}$. Los $\frakr_\nu$ quedan determinados unívocamente por su comportamiento, indicado arriba, respecto de los ideales fundamentales y por la exigencia de que la norma individual o el divisor elemental más alto sea un mayor factor primario de $R^{(i)}$, respectivamente de $E^{(i)}$ (H.-N., Teorema X).

\textbf{8. Concepto de dimensión.} Al ideal $\frakm$ se le asigna la dimensión máxima $n-i$ si $R^{(i)}$ es la primera norma individual distinta de la unidad ---o, lo que por 6 es equivalente, si $\frakg_i$ es el primer ideal fundamental distinto del ideal unidad; al ideal unidad $\frako$ se le asigna la dimensión $-1$. Si sólo hay una norma individual $R^{(i)}$ distinta de la unidad, de modo que $\frakg_i=\frakg_{i+1}=\cdots=\frakm$, entonces la dimensión máxima se llama también simplemente dimensión de $\frakm$. Si $\frakm$ contiene un polinomio $G^{(i)}\ne0$, entonces tiene dimensión máxima a lo sumo $n-i$, como muestra la observación al final de 6. Si $\frakm$ tiene dimensión máxima $n-i$ y $\frakn$ es un divisor de $\frakm$, entonces $\frakn$ tiene también dimensión máxima a lo sumo $n-i$.

\textbf{9. Norma (forma resultante) y ceros.} Por ceros de un ideal $\frakm$ se entienden todos los sistemas de valores de $x_1,\ldots,x_i$ pertenecientes a un cuerpo algebraico de extensión adecuado de $P(u;x_{i+1},\ldots,x_n)$ ($i=1,2,\ldots,n$) para los cuales ---después de adjuntar $x_{i+1},\ldots,x_n$ al dominio de coeficientes--- se anulan todos los polinomios de $\frakm$. A cada factor $R^{(i)}$ de la norma (forma resultante), bajo esta adjunción, corresponden tantos ceros de $\frakm$ como indica el grado de $R^{(i)}$; y $R^{(i)}$, escrito en combinaciones bilineales $w_{\mu\nu}$ de las $u_{\mu\nu}$ y de otras indeterminadas $v_{\mu\nu}$, admite la descomposición explícita
\begin{equation}
 R^{(i)}=\prod_\nu (z_i-z_i^{(\nu)})
 =\prod_\nu (v_{i1}x_1+\cdots+v_{in}x_n - z_i^{(\nu)}), \tag{5}
\end{equation}
donde $x_1^{(\nu)},\ldots,x_i^{(\nu)}$ denota un sistema de ceros asociados (H.-N., Teorema XIII; en el § 3 se dará una nueva fundamentación). Así, entre los ceros de un ideal de dimensión máxima $n-i$ hay ceros que dependen de $(n-i)$ parámetros, pero no hay ninguno que dependa de más parámetros; queda probado con ello el acuerdo del concepto de dimensión dado en 8 con la formulación usual.

\textbf{10. El dominio de anillos de la teoría general de ideales.} El dominio subyacente será un anillo conmutativo en el que vale el teorema de la cadena finita: toda cadena de ideales $\fraka_1,\fraka_2,\ldots,\fraka_k,\ldots$, donde cada $\fraka_i$ es un divisor propio ---divisor en el sentido ideal--- del inmediatamente precedente, se interrumpe después de finitos términos. Esta exigencia es equivalente a la otra, que todo ideal posea una base ideal (\emph{Idealtheorie}, Teorema I), y por tanto se cumple en particular para el dominio polinómico. En los números siguientes 11, 12, 13 se supone este dominio general de anillos.

\textbf{11. Ideales primos e ideales primarios.} Un ideal $\frakp$ se llama ideal primo si de la divisibilidad de un producto por $\frakp$ se sigue la divisibilidad de al menos un factor; $\frakq$ se llama ideal primario ---o primario--- si de la divisibilidad de un producto por $\frakq$ se sigue la divisibilidad de un factor o de una potencia de cada factor. En símbolos:
\[
 a\not\equiv0\pmod{\frakp},\quad b\not\equiv0\pmod{\frakp}\quad\hbox{implica}\quad ab\not\equiv0\pmod{\frakp};
\]
\[
 a\not\equiv0\pmod{\frakq},\quad b^x\not\equiv0\pmod{\frakq}\ \hbox{para toda potencia }x\quad\hbox{implica}\quad ab\not\equiv0\pmod{\frakq}.
\]
A cada ideal primario $\frakq$ le corresponde uno y sólo un ideal primo asociado $\frakp$, que es divisor de $\frakq$ y una potencia del cual es divisible por $\frakq$, $\frakp\ne\frako$, $\frakp^\rho\equiv0\pmod{\frakq}$; aquí $\frakp$ consiste en todos los elementos del anillo de los que alguna potencia es divisible por $\frakq$. Como muestra la definición, todo ideal primo es a la vez un ideal primario; por ideales primarios propios se entienden los que no son al mismo tiempo ideales primos, de modo que $\rho>1$.

\textbf{12. El teorema de descomposición.} Una representación $\fraka=[\frakq_1,\ldots,\frakq_s]$ de un ideal como mínimo común múltiplo se llama mínima si ningún $\frakq$ está contenido en el mínimo común múltiplo de los otros ---en su complemento; se llama representación por mayores componentes primarias si cada $\frakq$ es primario, pero el mínimo común múltiplo de dos $\frakq$ cualesquiera no es primario. Todo ideal admite una representación mínima por finitas mayores componentes primarias; para dos representaciones distintas de este tipo coinciden el número de componentes y los ideales primos asociados, todos mutuamente distintos. Los ideales primos así determinados unívocamente se llamarán los ideales primos, o ideales primos asociados, del ideal (\emph{Idealtheorie}, Teorema IX).

\textbf{13. Componentes aisladas de la descomposición.} Un ideal $\frakr=[\frakq_1,\ldots,\frakq_i]$ se llama componente aislada de la descomposición de $\fraka$ si los $\frakq_i$ aparecen en al menos una representación mínima de $\fraka$ por mayores componentes primarias y satisfacen la condición de que ninguno de sus ideales primos asociados esté contenido en uno de los otros ideales primos de $\fraka$. Las componentes aisladas de la descomposición están determinadas unívocamente por sus ideales primos; en particular, las mayores componentes primarias aisladas están determinadas unívocamente (\emph{Idealtheorie}, Teorema XIII).

\textbf{14. Característica, cuerpos perfectos e imperfectos, extensión de primera especie.} Un cuerpo $\Omega$ se dice de característica cero si el cuerpo primo contenido en $\Omega$ y derivado de la unidad es del tipo del cuerpo de los números racionales; tiene característica $p$ si este cuerpo primo es del tipo del sistema de clases residuales módulo un número primo $p$. Un cuerpo $\Omega$ se llama perfecto si toda función prima con coeficientes de $\Omega$ se descompone en un cuerpo de extensión adecuado en factores lineales distintos; de lo contrario se llama imperfecto. Todo cuerpo de característica cero es perfecto; un cuerpo de característica $p$ es perfecto si y sólo si, junto con cada elemento, contiene también su raíz $p$-ésima. Una función prima en un cuerpo imperfecto es una función entera de grado $r$ en $x^{p^f}$, y se descompone en un cuerpo de extensión adecuado en potencias $p^f$-ésimas de $r$ factores lineales distintos; $f$ se llama el exponente. Una función prima se dice de primera especie si se descompone en un cuerpo de extensión adecuado en factores lineales distintos, de modo que el exponente $f$ es cero; una extensión se llama de primera especie si todo elemento es de primera especie, es decir, es cero de una función prima de primera especie. Toda extensión que surge al adjuntar un elemento de primera especie es de primera especie; sobre cuerpos perfectos sólo hay extensiones de primera especie (Steinitz, § 11 y § 13).

\textbf{15. Grado reducido y exponente de una extensión finita.} Si $\mathcal K$ es una extensión finita de un cuerpo imperfecto $\Omega$, entonces $\mathcal K$ contiene un subcuerpo $\mathcal K_0$ de grado $r$ que es de primera especie respecto de $\Omega$ y consiste en todos y sólo los elementos de primera especie de $\mathcal K$. La potencia $p^f$-ésima de todo elemento de $\mathcal K$ pertenece a $\mathcal K_0$, y hay elementos en $\mathcal K$ que son exactamente de grado $rp^f$. $r$ se llama el grado reducido, $f$ el exponente de la extensión finita (Steinitz, § 14, 1 y § 14, 5).

\subsection*{§ 2. Forma de divisores elementales y norma de los ideales primos y primarios. Divisores de los ideales primos}

En adelante trataremos siempre del dominio polinómico definido en § 1, 1. Primero hay que mostrar que también para ideales primos e ideales primarios se puede uno restringir a ideales transformados, según

\textbf{Lema I.} El ideal transformado $\frakp$ de un ideal primo $\bar{\frakp}$ en $\frakS$ se convierte en un ideal primo; el transformado $\frakq$ de un ideal primario $\bar{\frakq}$ en $\frakS$ se convierte en un ideal primario. En otras palabras, la propiedad de ser ideal primo o ideal primario se conserva cuando se adjuntan las $u_{\mu\nu}$ a $P$.

En efecto, supóngase que $\fraka\not\equiv0\pmod{\frakp}$, $\frakb\not\equiv0\pmod{\frakp}$, donde $\fraka$ y $\frakb$ no tienen que ser ideales transformados; digamos que $a(x)\not\equiv0\pmod{\frakp}$ y $b(x)\not\equiv0\pmod{\frakp}$, con $a(x)$ y $b(x)$ polinomios de $\fraka$ y $\frakb$. Por inversión de (1) según (1') se obtiene ---entendiendo $T$, sin pérdida de generalidad, como productos de potencias de las $t_{\mu\nu}$---
\[
 a(x)=\sum T_i a_i(y),\qquad b(x)=\sum T_i' b_i(y),
\]
y al menos un $a_i(y)$ y un $b_i(y)$ no deben ser divisibles por $\bar{\frakp}$. Sean, por ejemplo, $a_h(y)\not\equiv0\pmod{\bar{\frakp}}$ y $b_k(y)\not\equiv0\pmod{\bar{\frakp}}$ los respectivos términos más altos bajo algún orden lexicográfico de las $T$; entonces también $a_h(y)b_k(y)\not\equiv0\pmod{\bar{\frakp}}$, y por consiguiente $a(x)b(x)\not\equiv0\pmod{\frakp}$, luego $\fraka\frakb\not\equiv0\pmod{\frakp}$, lo que prueba que $\frakp$ es primo. La demostración descansa evidentemente en que el sistema de clases residuales módulo un ideal primo forma un anillo sin divisores de cero, propiedad que se conserva al adjuntar indeterminadas.

Correspondientemente supóngase que $\fraka\not\equiv0\pmod{\frakq}$ y $\frakb^x\not\equiv0\pmod{\frakq}$ para toda potencia $x$; entonces de la existencia de la base ideal se sigue que deben existir un $a(x)$ de $\fraka$ y un $b(x)$ de $\frakb$ tales que $a(x)\not\equiv0\pmod{\frakq}$ y $b(x)^x\not\equiv0\pmod{\frakq}$ para toda potencia $x$; pues, bajo el supuesto contrario, una potencia de $\frakb$ sería divisible por $\frakq$.

Poniendo $a(x)b(x)=c(x)$, sean $\fraka^*,\frakb^*,\frakc^*$ los ideales derivados respectivamente de los coeficientes $a_i(y),b_j(y),c_k(y)$ de $a(x),b(x),c(x)$. Entonces $\frakb^{*x}\not\equiv0\pmod{\bar{\frakq}}$ para toda potencia $x$, pues de lo contrario $b(x)^x$ sería divisible por $\frakq$; puesto que $\fraka^*\not\equiv0\pmod{\bar{\frakq}}$, se debe tener también $\fraka^*\frakb^{*\lambda}\not\equiv0\pmod{\bar{\frakq}}$ para toda potencia $\lambda$. Pero por el teorema de Dedekind--Mertens\footnote{Compárese H.-N., p. 63 y nota 10).} existe un exponente $\lambda$ tal que $\fraka^*\frakb^{*\lambda}\equiv \frakc^*\frakb^{*\lambda-1}$; por tanto $\frakc^*\not\equiv0\pmod{\bar{\frakq}}$. Esto da $c(x)\not\equiv0\pmod{\frakq}$ y por consiguiente $\fraka\frakb\not\equiv0\pmod{\frakq}$, probando que $\frakq$ es primario.

También al representar un ideal como mínimo común múltiplo se puede uno restringir a ideales transformados, según

\textbf{Lema II.} De una representación $\barfrakm=[\frakc_1,\ldots,\frakc_s]$ en $\frakS$ se sigue una representación $\frakm=[\frakc_1',\ldots,\frakc_s']$ en $\frakS'$, donde $\frakc_i'$ designa el ideal transformado de $\frakc_i$. En particular, por tanto, en toda representación mínima por mayores componentes primarias éstas pueden suponerse transformadas.

En efecto, $\frakm$ es divisible por $[\frakc_1',\ldots,\frakc_s']$, pues todo polinomio $f(x)$ de $\frakm$ ---después de multiplicar, si es necesario, por una potencia adecuada de $U(u)$--- es de la forma $\sum U_i(u)f_i(y)$, donde cada $f_i(y)$, como polinomio de $\barfrakm$, es por hipótesis divisible por todos los $\frakc_i$. Recíprocamente, si $f(x)$ es divisible por todos los $\frakc_i'$, entonces tiene la forma anterior y por tanto es divisible por $\frakm$. Así, si se parte de una descomposición en mayores componentes primarias en $\frakS$, se obtiene una descomposición $\frakm=[\frakq_1',\ldots,\frakq_s']$ en $\frakS'$; aquí los $\frakq_i'$, como ideales transformados de ideales primarios, son primarios por el Lema I, y son mayores componentes primarias, puesto que ideales primos asociados distintos $\bar{\frakp}$ en $\frakS$ corresponden a distintos $\frakp$ en $\frakS'$.

Necesitamos ahora una primera caracterización de ideales primos e ideales primarios por medio de la forma de divisores elementales y la norma, todavía sin separación completa entre primo y primario. En exacta analogía con la teoría de ideales en cuerpos de números algebraicos,\footnote{Como divisor elemental más alto de ideales en cuerpos de números algebraicos hay que considerar el menor número racional entero divisible por el ideal, por tanto en particular el número primo divisible por un ideal primo o la potencia prima divisible por un ideal primario (una potencia de un ideal primo). De hecho, todo ideal $\fraka^*$ es también un módulo $A^*$ de formas lineales en los elementos $\omega_1,\ldots,\omega_n$ de una base del cuerpo, con $(\omega_1,\ldots,\omega_n)$ como módulo fundamental de $A^*$; el menor número racional entero divisible por $\fraka^*$ se convierte por tanto en el divisor elemental más alto de este módulo $A^*$, mientras que la norma de $\fraka^*$ se convierte en el producto de todos los divisores elementales de $A^*$.} se tiene

\textbf{Teorema I.} La forma de divisores elementales de un ideal primo se hace igual a una función prima, y la norma a una potencia de esta función prima. Para ideales primarios, la forma de divisores elementales y la norma se hacen funciones primarias, a saber, potencias de la misma función prima, que es ella misma la forma de divisores elementales del ideal primo asociado. Para ideales primos e ideales primarios, la dimensión máxima coincide con la dimensión simplemente dicha; esta dimensión coincide para un ideal primario y su ideal primo asociado.

El Teorema I se cumple evidentemente para el ideal unidad, cuya forma de divisores elementales y norma se hacen iguales a la unidad. Sea ahora el ideal primo $\frakp$ de dimensión máxima $(n-i)$; por (4), § 1, 6, se obtiene
\[
 E^{(i)}\frakg_{i-1}\not\equiv0\pmod{\frakp},\quad\hbox{pero}\quad
 E^{(i)}\frakg_i\equiv0\pmod{\frakp},
\]
pues de lo contrario $\frakp$ tendría, por § 1, 8, dimensión máxima a lo sumo $n-i-1$. De aquí $\frakg_i\equiv0\pmod{\frakp}$, de modo que $\frakp$ es idéntico a su ideal fundamental de etapa $i$ y por consiguiente también a $\frakg_{i+1},\frakg_{i+2},\ldots$. Así $R^{(i+1)},\ldots,R^{(n)}$ se hacen iguales a la unidad, y por tanto también $E^{(i+1)},E^{(i+2)},\ldots$, como divisores de $R^{(i+1)},R^{(i+2)},\ldots$; por consiguiente, la forma de divisores elementales se hace igual a $E^{(i)}$, que debe ser igual a una función prima $P^{(i)}$. Pues por § 1, 5 y teniendo en cuenta que $\frakg_{i-1}$ se convierte en el ideal unidad, $E^{(i)}$ se define como el máximo común divisor ---en el sentido polinómico--- de todos los $B^{(i)}$ divisibles por $\frakp$; de $E^{(i)}=E_1^{(i)}E_2^{(i)}\equiv0\pmod{\frakp}$ se seguiría que $E^{(i)}$ tendría que estar contenido en uno de sus factores $E_1^{(i)}$ o $E_2^{(i)}$. Pero, puesto que una potencia de $E^{(i)}$ es divisible por $R^{(i)}$, $R^{(i)}$ se hace una potencia ---posiblemente la primera--- de esta función prima $P^{(i)}$; al mismo tiempo $R^{(i)}$ se convierte en la norma de $\frakp$. Como sólo aparece una norma individual distinta de la unidad, la dimensión máxima de $\frakp$ coincide finalmente con su dimensión simplemente dicha.

Correspondientemente, sea el ideal primario $\frakq$ de dimensión máxima $n-i$; entonces, como arriba,
\[
 (E^{(i)}\frakg_{i-1})^x\not\equiv0\pmod{\frakq},\quad\hbox{pero}\quad
 (E^{(i)}\frakg_i)^x\equiv0\pmod{\frakq}
\]
para toda potencia $x$; y en consecuencia, como arriba, $\frakq=\frakg_i$, su forma de divisores elementales es $E^{(i)}$, su norma es $R^{(i)}$, y su dimensión máxima es su dimensión simplemente dicha. Esta dimensión coincide con la del ideal primo asociado; pues de $\frakp^\rho\equiv0\pmod{\frakq}$, $\frakq\equiv0\pmod{\frakp}$ se sigue que la dimensión $(n-i)$ de $\frakp$ es a lo sumo igual a la de $\frakq$, y que la de $\frakq$ es a lo sumo igual a la de $\frakp^\rho$; de $(P^{(i)})^\rho\equiv0\pmod{\frakp^\rho}$, donde $P^{(i)}$ designa la forma de divisores elementales de $\frakp$, se halla que esta última dimensión máxima es a lo sumo $n-i$, y por tanto $n-i$ es la dimensión de $\frakp$ y $\frakq$. De $(P^{(i)})^\rho\equiv0\pmod{\frakq}$ se sigue además que $E^{(i)}$ ---como máximo común divisor de todos los $B^{(i)}$ de $\frakq$--- se convierte en una potencia de $P^{(i)}$, que también puede ser la primera potencia, y que lo mismo vale para $R^{(i)}$. Esto prueba el Teorema I en todas sus partes.

Una caracterización de los ideales primos mediante el concepto de dimensión viene dada por

\textbf{Teorema II.} Un ideal primo no tiene ningún divisor propio de la misma dimensión máxima; recíprocamente, todo ideal con esta propiedad es primo.

La demostración descansa en

\textbf{Lema III.} Si un ideal primo $\frakp$ de polinomios con coeficientes de un cuerpo $\Omega$ tiene sólo finitas clases residuales linealmente independientes sobre $\Omega$, entonces $\frakp$ no tiene ningún divisor propio distinto del ideal unidad; en otras palabras, el sistema de clases residuales módulo $\frakp$ forma un cuerpo.

La hipótesis dice que hay un número finito $k$ tal que entre cualesquiera $(k+1)$ clases residuales hay una dependencia lineal con coeficientes de $\Omega$ ---más precisamente, con coeficientes del cuerpo de clases residuales $(\Omega)$ representado por todos los elementos de $\Omega$---, pero que existe al menos un sistema de $k$ clases residuales linealmente independientes sobre $\Omega$. De ello se sigue inmediatamente que todas las clases residuales pueden expresarse linealmente mediante cualquier sistema elegido de $k$ clases linealmente independientes. Sea ahora $\fraka$ un divisor propio de $\frakp$, y sea $a(z)\not\equiv0\pmod{\frakp}$ un polinomio de $\fraka$; de $c_1f_1(z)+\cdots+c_kf_k(z)\equiv0\pmod{\frakp}$ se sigue que
\[
 c_1a(z)f_1(z)+\cdots+c_ka(z)f_k(z)\equiv0\pmod{\frakp}.
\]
Esto significa que, junto con las clases residuales de $f_1,\ldots,f_k$, también las clases residuales de $af_1,\ldots,af_k$ forman un sistema de $k$ clases linealmente independientes, mediante el cual en particular la clase unidad puede representarse linealmente. Por tanto $\fraka$ es el ideal unidad; dicho de otro modo, el sistema de clases residuales forma un cuerpo, puesto que para $a(z)\not\equiv0\pmod{\frakp}$ siempre existe un $f(z)$ tal que $1\equiv a(z)f(z)\pmod{\frakp}$.

Para la demostración de la primera parte del Teorema II sólo resta observar que todo ideal primo de dimensión $(n-i)$ ---más generalmente todo ideal de dimensión máxima $(n-i)$--- tiene sólo finitas clases residuales linealmente independientes sobre $P(u;x_{i+1},\ldots,x_n)$, a saber, por § 1, 5, tantas como indica el grado de $R^{(i)}$, puesto que aquí $\frakg_{i-1}$ se convierte en el ideal unidad. Todo divisor propio $\fraka$ de $\frakp$ se convierte por tanto en el ideal unidad cuando $x_{i+1},\ldots,x_n$ se adjuntan a $P(u)$; dicho de otro modo, el ideal fundamental de etapa $i$ de $\fraka$ se convierte en el ideal unidad, lo cual significa que la dimensión máxima de $\fraka$ es a lo sumo $n-i-1$. Para que la afirmación valga también para un ideal no transformado $\fraka$ como divisor, basta pasar, según § 1, 2, a un nuevo dominio transformado.

Recíprocamente, supóngase ahora que $\frakp$ no tiene ningún divisor propio de la misma dimensión máxima $n-i$, y sean $\fraka\not\equiv0\pmod{\frakp}$ y $\frakb\not\equiv0\pmod{\frakp}$, con $\fraka$ y $\frakb$ de nuevo supuestos ideales transformados. Entonces, por hipótesis, puesto que $(\fraka,\frakp)$ y $(\frakb,\frakp)$, como máximos comunes divisores ---en el sentido ideal---, se convierten en divisores propios de $\frakp$, se tiene
\[
 E^{(n)}\cdots E^{(i+1)}\frakg_i\equiv0\pmod{(\fraka,\frakp)},
\qquad
 E^{(n)}\cdots E^{(i+1)}\frakg_i\equiv0\pmod{(\frakb,\frakp)}.
\]
Esto da
\[
 E^{(n)}\cdots E^{(i+1)}\frakg_i\equiv0\pmod{(\fraka\frakb,\frakp)},
\]
y por consiguiente necesariamente $\fraka\frakb\not\equiv0\pmod{\frakp}$, pues de lo contrario $\frakp$ tendría dimensión máxima menor que $n-i$. Como $\fraka,\frakb$ son ideales transformados, $\frakp$, y por el Lema I también $\bar{\frakp}$, son ideales primos. Esto prueba el Teorema II.

\subsection*{§ 3. Ceros de los ideales primos y de los ideales primarios}

El paso a una fundamentación directa de la teoría de los ceros (§ 1, 9), primero para ideales primos, lo constituye el siguiente lema, que amplía el Lema III y es válido para ideales primos en anillos arbitrarios.

\textbf{Lema IV.} El sistema de clases residuales módulo todo ideal primo distinto del ideal unidad puede ampliarse, por formación de cocientes, a un cuerpo, el cuerpo de clases residuales del ideal primo.

El sistema de clases residuales módulo un ideal arbitrario forma un anillo, puesto que suma, diferencia y producto pertenecen de nuevo al sistema, y las leyes asociativa, conmutativa y distributiva se conservan al pasar a clases residuales. Si el ideal es específicamente primo, este anillo no tiene divisores de cero; pues la definición de ideales primos (§ 1, 11) dice que el producto de clases residuales no nulas tampoco se anula. Si el ideal primo es distinto del ideal unidad, entonces este anillo contiene al menos un elemento distinto del elemento cero, y por tanto puede ampliarse a un cuerpo por formación de cocientes ---adjunción de pares de elementos;\footnote{Steinitz, § 3. Basta la hipótesis de un elemento no nulo en el dominio íntegro original, pues entonces la existencia de la unidad queda asegurada por formación de cocientes, y así el dominio original se convierte en subdominio del cuerpo; Steinitz supone la existencia de la unidad en el dominio íntegro.} así queda definido el cuerpo de clases residuales.

Fijemos primero la notación que se mantendrá en todo lo que sigue.

\textbf{Notación.} Las clases residuales se denotarán en general poniendo entre paréntesis un representante cualquiera de la clase, $(x),\ldots,(a(x)),\ldots$; de modo similar, los cuerpos de clases residuales se denotarán mediante paréntesis. En particular, $(P)$, $(P(u))$, $(P(u,x_{i+1},\ldots,x_n))$ denotan los cuerpos de clases residuales consistentes en todos y sólo aquellos residuos que pueden representarse por elementos de $P$, $P(u)$, $P(u,x_{i+1},\ldots,x_n)$.

Recordemos también lo siguiente. Se dice que un cuerpo $\mathcal R$ tiene rango algebraico (grado de trascendencia) $k$ respecto de un cuerpo base $\Omega$ si $\mathcal R$ contiene al menos un sistema de $k$ cantidades algebraicamente independientes sobre $\Omega$ ---cantidades trascendentes---, pero cualesquiera $k+1$ cantidades de $\mathcal R$ son algebraicamente dependientes sobre $\Omega$. Si $\mathcal R$ puede representarse como cuerpo de extensión algebraica de un subcuerpo $\Omega(z_1,\ldots,z_s)$, donde las $z$ son algebraicamente independientes ---trascendentes--- sobre $\Omega$, entonces necesariamente $s=k$;\footnote{Para una demostración de este hecho familiar válida en característica arbitraria, compárese Steinitz, § 22, 9. --- Al adjuntar las $u$, el rango algebraico no puede aumentar, puesto que sólo intervienen combinaciones racionales de los elementos originales con coeficientes de $\Omega(u)$; tampoco puede disminuir, puesto que toda relación entre elementos de $\mathcal R$ con coeficientes de $\Omega(u)$ se descompone en finitas relaciones con coeficientes de $\Omega$.} asimismo el rango algebraico de $\mathcal R(u)$ sobre $\Omega(u)$ es igual a $k$ cuando $u$ denota indeterminadas no contenidas en $\mathcal R$.

La construcción del cuerpo de ceros procede ahora en completa analogía con el camino usual para funciones primas de una indeterminada,\footnote{Steinitz, § 6.} según

\textbf{Teorema III.} Al cuerpo de clases residuales $(\mathcal R)$ de un ideal primo $\frakp$ de dimensión $n-i$ puede asignársele, isomórficamente, un cuerpo de ceros $R$ de $\frakp$, que tiene rango algebraico $n-i$ ---depende de $n-i$ parámetros--- y en particular puede considerarse como cuerpo de extensión finita de $P(u,x_{i+1},\ldots,x_n)$. El isomorfismo entre $(\mathcal R)$ y $R$ incluye el existente entre $(P)$ y $P$, $(P(u))$ y $P(u)$, y, cuando $x_{i+1},\ldots,x_n$ se toman como parámetros, también el existente entre $(P(u,x_{i+1},\ldots,x_n))$ y $P(u,x_{i+1},\ldots,x_n)$. Recíprocamente, todo cuerpo de ceros de rango algebraico $n-i$ ---dependiente de $n-i$ parámetros--- es isomorfo al cuerpo de clases residuales $(\mathcal R)$.

El Teorema III da como corolario el teorema conocido: si un ideal $\fraka$ se anula en todos los ceros de un ideal primo $\frakp$, entonces $\fraka$ es divisible por $\frakp$.

Para la demostración, obsérvese primero que el cuerpo de clases residuales $(\mathcal R)$ tiene rango algebraico $n-i$ respecto de $(P(u))$. En efecto, las clases $(x_{i+1}),\ldots,(x_n)$ son algebraicamente independientes respecto de $(P(u))$, puesto que $\frakp$, por ser de dimensión $n-i$, no contiene ningún polinomio libre de $x_1,\ldots,x_i$; toda clase ulterior, sin embargo, depende de éstas. Pues de la pertenencia de la forma de divisores elementales $P^{(i)}$ a $\frakp$ se sigue, por § 1, 2, también para $\lambda=1,2,\ldots,i$, que pertenecen a $\frakp$ polinomios que dependen respectivamente de $x_\lambda,x_{i+1},\ldots,x_n$. Así $(\mathcal R)$ se convierte en un cuerpo de extensión finita de $(P(u,x_{i+1},\ldots,x_n))$; por § 1, 5 ---teniendo en cuenta que $\frakg_{i-1}$ es igual al ideal unidad y $\frakg_i$, por el Teorema I, a $\frakp$--- el grado sobre este subcuerpo es igual al grado de la norma.

Para pasar a un cuerpo de ceros isomorfo, obsérvese además que $P(u,x_{i+1},\ldots,x_n)$ y $(P(u,x_{i+1},\ldots,x_n))$ están relacionados isomórficamente asignando a cada elemento de $P(u,x_{i+1},\ldots,x_n)$ la clase residual representada por ese elemento. Primero, bajo esta asignación, los dominios íntegros consistentes en polinomios en $u,x_{i+1},\ldots,x_n$ y respectivamente en $(u),(x_{i+1}),\ldots,(x_n)$ se corresponden isomórficamente, puesto que $\frakp$ no contiene ningún polinomio libre de $x_1,\ldots,x_i$, y por tanto polinomios distintos de $P(u,x_{i+1},\ldots,x_n)$ corresponden a clases distintas; de dominios íntegros isomorfos, sin embargo, surgen cuerpos isomorfos por formación de cocientes. La relación isomorfa entre $P(u,x_{i+1},\ldots,x_n)$ y $(P(u,x_{i+1},\ldots,x_n))$ incluye la de $P$ con $(P)$, y la de $P(u)$ con $(P(u))$. Para extender este isomorfismo a uno entre $(\mathcal R)$ y un cuerpo de ceros $R$, asígnense isomórficamente ciertos nuevos elementos $\xi_1,\ldots,\xi_i$ a las clases $(x_1),\ldots,(x_i)$; es decir, a cada polinomio $(a(x))$ corresponde un polinomio $a(\xi_1,\ldots,\xi_i,x_{i+1},\ldots,x_n)$, y las sumas y productos se corresponden con sumas y productos. Por formación de cocientes ---donde uno puede restringirse a denominadores del subcuerpo $(P(u,x_{i+1},\ldots,x_n))$, respectivamente $P(u,x_{i+1},\ldots,x_n)$--- al cuerpo de clases residuales $(\mathcal R)$ corresponde entonces un cuerpo $R$, una extensión finita de $P(u,x_{i+1},\ldots,x_n)$ de grado igual al de la norma; puede llamarse el cuerpo de ceros. Pues la asignación dice que $f(\xi_1,\ldots,\xi_i,x_{i+1},\ldots,x_n)$ se anula tan pronto como $f(x)\equiv0\pmod{\frakp}$. En lugar de $(P(u,x_{i+1},\ldots,x_n))$ podría aparecer en todo el argumento cualquier otro subcuerpo de $(\mathcal R)$ obtenido al adjuntar $n-i$ cantidades algebraicamente independientes.

Recíprocamente es fácil mostrar que todo cuerpo de ceros $R$ de rango algebraico $n-i$ es isomorfo al cuerpo de clases residuales. Puesto que $R$ puede derivarse racionalmente, con coeficientes de $P(u)$, de las cantidades $\xi_1,\ldots,\xi_i$, entre estas cantidades debe haber $n-i$ elementos algebraicamente independientes, y por tanto trascendentes, respecto de $P(u)$. En particular esto debe valer para $\xi_{i+1},\ldots,\xi_n$, pues, como $P^{(i)}\equiv0\pmod{\frakp}$, se sigue como arriba que $R$ es un cuerpo de extensión algebraica de $P(u,\xi_{i+1},\ldots,\xi_n)$. Puesto que por hipótesis $f(\xi)$ se anula para todo polinomio $f(x)$ de $\frakp$, a cada clase del dominio polinómico de las $(x)$ contenida en $(\mathcal R)$ corresponde uno y sólo un elemento de $R$ que es un polinomio en las $\xi$; y las sumas y productos se corresponden con sumas y productos. Pero clases distintas corresponden a elementos distintos de $R$; de lo contrario una clase no nula, y por tanto después de una multiplicación adecuada también una clase representada por un polinomio de $(P(u,x_{i+1},\ldots,x_n))$, correspondería al elemento cero en $R$, y contra la hipótesis las cantidades $\xi_{i+1},\ldots,\xi_n$ satisfacerían una dependencia algebraica. Esta asignación agota todos los polinomios en las $\xi$ contenidos en $R$; por formación de cocientes ---donde de nuevo uno puede restringirse a denominadores de $(P(u,x_{i+1},\ldots,x_n))$, respectivamente $P(u,x_{i+1},\ldots,x_n)$--- surgen cuerpos isomorfos de estos dominios íntegros isomorfos.

Sea ahora el ideal $\fraka$ nulo en todos los ceros de $\frakp$, en particular también en todos los ceros que dependen de $n-i$ parámetros. Si $a(x)$ es un polinomio cualquiera de $\fraka$, entonces $a(\xi)$ se anula; por el isomorfismo entre $R$ y $(\mathcal R)$, la clase $(a(x))$ es por tanto la clase cero, y $a(x)$, luego $\fraka$, es divisible por $\frakp$. Esto prueba el Teorema III y su corolario.

Para pasar del Teorema III a la descomposición dada en § 1, 9, primero para ideales primos, y para poder precisar los exponentes, se requiere una investigación adicional de la forma de divisores elementales, según

\textbf{Lema V.} Si en característica $p$ la forma de divisores elementales $E^{(i)}$ de un ideal primo o primario ---más generalmente, de un ideal $\fraka$ para el que coinciden dimensión máxima y dimensión simplemente dicha--- tiene exponente $f$ respecto de $x_i$, y por tanto es una función entera de $x_i^{p^f}$, entonces tiene también exponente $f$ respecto de $x_{i+1},\ldots,x_n$ y respecto de las indeterminadas $u_{\mu\nu}$, respectivamente $t_{\mu\nu}$, que aparecen en $E^{(i)}$. En particular, si el dominio de coeficientes es un cuerpo perfecto, la forma de divisores elementales $P^{(i)}$ de un ideal primo tiene siempre exponente cero respecto de $x_i$, y por tanto es una función prima de primera especie.

Para ello debe observarse que $P(u,x_{i+1},\ldots,x_n)$ se vuelve imperfecto tan pronto como $P$ es un cuerpo perfecto de característica $p$, de modo que no se sigue directamente de § 1, 14 que $P^{(i)}$ sea de primera especie. Sea $E^{(i)}$ de exponente $f$ respecto de $x_i$, pero supóngase que contiene otra indeterminada $x_{i+1}$ con exponente $f'$. Entonces ---permutando las $u_{\mu\nu}$, respectivamente $t_{\mu\nu}$--- existe también, por § 1, 2, un polinomio no nulo $G^{(i)}$, divisible por el ideal, que es una función entera de $x_{i+1}^{p^{f'}}$, y que, por definición de la forma de divisores elementales, es divisible por $E^{(i)}$. Pero, como $G^{(i)}$ tiene el mismo grado que $E^{(i)}$, debe coincidir con $E^{(i)}$ salvo un factor de $P(u)$; necesariamente, por tanto, $f'=f$. Así $E^{(i)}$ ---que puede suponerse entero y primitivo en las $t_{\mu\nu}$--- es de la forma
\[
 E^{(i)}=\sum c_\rho(t)\,x_i^{\alpha_\rho p^f}x_{i+1}^{\beta_\rho p^f}\cdots x_n^{\gamma_\rho p^f}
       =\sum c_\rho(t)\,X_\rho^{p^f},
\]
donde los $X_\rho$ son productos de potencias de las $x$, y los $X_\rho$ son divisibles por $\fraka$. Por tanto se obtiene de nuevo un polinomio perteneciente a $\fraka$ si en los $X_\rho$ todo exponente de cada $t$ se reemplaza por el mayor múltiplo de $p^f$ contenido en él; como muestra la representación exhibida, esto equivale a reemplazar, en $c_\rho(t)$, todo exponente de las $t$ por el mayor múltiplo de $p^f$ contenido en él. Mediante esta sustitución $E^{(i)}$ pasa a un polinomio en $x_i^{p^f},\ldots,x_n^{p^f}$ que, por definición, es divisible por $E^{(i)}$, y, por la primitividad supuesta, también respecto de las $t$; por tanto debe ser idéntico a $E^{(i)}$, ya que el grado no ha aumentado en ningún argumento. En particular, si $P$ es perfecto, entonces $E^{(i)}$, como función entera de $x_i^{p^f}$, es una potencia $p^f$-ésima; por tanto, si se trata de un ideal primo, necesariamente $f=0$, y la forma de divisores elementales $P^{(i)}$ se convierte en una función prima de primera especie.

La descomposición viene ahora dada por

\textbf{Teorema IV.} La forma de divisores elementales $P^{(i)}$ de un ideal primo admite, en el cuerpo de Galois derivado del cuerpo de ceros, la descomposición explícita
\[
 P^{(i)}=\prod_\nu (x_i-t_{i+1}x_{i+1}^{(\nu)}-\cdots-t_nx_n^{(\nu)})^e
 =\prod_\nu (t_i(y_i-y_i^{(\nu)})+\cdots+t_n(y_n-y_n^{(\nu)}))^e,
\]
donde $y_1^{(\nu)},\ldots,y_n^{(\nu)}$ denota un sistema asociado de ceros de las indeterminadas originales $y$, independiente de $t_i,\ldots,t_n$; distintos $\nu$ corresponden a factores lineales distintos; y el exponente $e$ tiene el valor uno cuando $P$ es perfecto y el valor $p^f$ ($f>0$) cuando es imperfecto. Con ello, después de adjuntar nuevas indeterminadas $v$, queda también determinada la descomposición de la forma de divisores elementales escrita en combinaciones bilineales de las $u$ y $v$ en lugar de las $u_{\mu\nu}$:
\[
 P^{(i)}(w)=\prod_\nu (z_i-z_i^{(\nu)})^e
 =\prod_\nu (v_{i1}(y_1-y_1^{(\nu)})+\cdots+v_{in}(y_n-y_n^{(\nu)}))^e,
\]
donde $e$ tiene el mismo significado que arriba. Asimismo queda dada la descomposición de la norma de un ideal primo, o de un ideal primario con $\frakp$ como ideal primo asociado ---como potencia de $P^{(i)}$---.

Como corolario del Teorema IV se obtiene también: si dos ideales primos coinciden en su forma de divisores elementales $P^{(i)}$, entonces son idénticos.

La demostración del Teorema IV consistirá en probar que la forma de divisores elementales $P^{(i)}$ es idéntica a la ecuación fundamental del cuerpo de extensión $(P(u,x_{i+1},\ldots,x_n),(x_i))$ sobre $(P(u,x_{i+1},\ldots,x_n))$. Primero, para obtener una visión de la dependencia respecto de las indeterminadas $u_{\mu\nu}$, respectivamente $t_{\mu\nu}$, se pasa al cuerpo de clases residuales $(\mathcal S)$ de $\frakp$, que, por lo dicho en la definición de rango algebraico, debe tener rango algebraico $n-i$ respecto de $(P)$, donde $(P)$ denota el subcuerpo representado por elementos de $P$. Pues $(\mathcal S)$ surge al adjuntar las indeterminadas $u$, respectivamente $t$, a un subcuerpo isomorfo a $(\mathcal R)$, obtenido por formación de cocientes a partir de todas y sólo aquellas clases representables por polinomios de $\frakS$, correspondiéndose $(P)$ y $(P)$ entre sí. El rango algebraico de $(\mathcal S)$ respecto de $(P(u))$ es por tanto igual al de este subcuerpo respecto de $(P)$, y de ahí, por el isomorfismo, igual al de $(\mathcal R)$ respecto de $(P)$. Como $(\mathcal S)$ puede derivarse racionalmente, con coeficientes de $(P)$, de las clases $(y_1),\ldots,(y_n)$, entre estas clases debe haber $n-i$ algebraicamente independientes, y $(\mathcal S)$ surge como cuerpo de extensión algebraica finita del subcuerpo generado al adjuntar estas $n-i$ clases a $(P)$.\footnote{Esta observación muestra que el Teorema III con su corolario podría haberse enunciado y demostrado directamente para el cuerpo de ceros $R$ de $\frakp$. Sólo al pasar al ideal transformado se consigue que en las mayores componentes primarias de igual dimensión de un ideal arbitrario aparezcan los mismos parámetros entre los ceros; y sólo así resulta la conexión con la norma y la forma de divisores elementales de un ideal arbitrario.}

Si se adjuntan sólo las indeterminadas $t_{\mu\nu}$ que aparecen en $x_{i+1},\ldots,x_n$, surge de nuevo un cuerpo de clases residuales que contiene las clases $(y_1),\ldots,(y_n)$ y $(x_{i+1}),\ldots,(x_n)$ y se convierte en un cuerpo de extensión algebraica finita de $(P(t_{\mu\nu},x_{i+1},\ldots,x_n))$; se trata de cuerpos isomorfos a subcuerpos de $(\mathcal S)$, no idénticos a ellos, pues para indeterminadas distintas las clases residuales están representadas por los mismos elementos, pero no son idénticas porque no contienen los mismos elementos. Así, en la dependencia de las clases $(y)$ respecto de $(P(t_{\mu\nu},x_{i+1},\ldots,x_n))$ sólo intervienen las indeterminadas $t_{\mu\nu}$. Ahora fórmese, adjuntando $t_{i1},\ldots,t_{in}$, la clase $(x_i)=(t_{i1}y_1+\cdots+t_{in}y_n)$; sean $(x_{i\nu})=(t_{i1}y_{1\nu}+\cdots+t_{in}y_{n\nu})$ los $r$ valores conjugados distintos que yacen en el cuerpo de Galois con respecto a $(P(t_{\mu\nu},x_{i+1},\ldots,x_n))$. Entonces ---según se trate o no de extensiones de primera especie--- el producto sobre los factores $x_i-(x_{i\nu})$, o sobre sus potencias $p^f$-ésimas, es una función prima $(G)$ respecto de $(P(t_{\mu\nu},x_{i+1},\ldots,x_n))$; pues, puesto que existen elementos de grado $r\cdot p^f$ (§ 1, 15), $(x_i)$ debe ser necesariamente tal elemento, ya que todo elemento surge por especialización de las $t_{\mu\nu}$; pero $(x_i)$ es un cero de $(G)$. Al pasar al cuerpo de ceros isomorfo de $\frakp$ y al cuerpo de Galois derivado de él, $G$ admite por tanto una descomposición en factores lineales $x_i-t_{i1}y_{1\nu}-\cdots-t_{in}y_{n\nu}$, respectivamente en sus potencias $p^f$-ésimas. Ahora, como $(G)$ se anula para $x_i=(x_i)$, $G(x_i)$ es divisible por $\frakp$ y por tanto por $P^{(i)}$; y, como función prima respecto de $P(t_{\mu\nu},x_{i+1},\ldots,x_n)$ y primitiva en $x_i$, $G$ debe hacerse idéntica a $P^{(i)}$. Así queda probada la primera descomposición del Teorema IV.

Para pasar a la segunda descomposición, obsérvese que la propiedad de ser, o no ser, una extensión de primera especie se conserva al adjuntar indeterminadas. Por tanto, después de adjuntar $v_1,\ldots,v_n$ y todas las $t_{\mu\nu}$, fórmese la clase $(z)=(v_1x_1+\cdots+v_ix_i)$ y las clases conjugadas $(z_\nu)=(v_1x_{1\nu}+\cdots+v_ix_{i\nu})$. Entonces el producto sobre todos los $z-(z_\nu)$, o sobre sus potencias $p^f$-ésimas, es de nuevo una función prima $(H)$ que se anula para $z=(z)$. Por tanto $H$ se representa como producto de factores lineales, como antes, y es idéntica a la forma de divisores elementales del ideal primo derivado de $\frakp$ al componer la transformación (1) con $z=v_1x_1+\cdots+v_nx_n$; esta forma de divisores elementales, sin embargo ---por la especialización recíproca--- se obtiene de $P^{(i)}$ reemplazando las $u_{\mu\nu}$ por ciertas combinaciones bilineales de las $u$ y $v$.\footnote{H.-N., § 7.} Con estas descomposiciones queda dada también la descomposición de la norma como potencia de $P^{(i)}$,\footnote{Para extensiones de primera especie, y por tanto en particular sobre cuerpos perfectos, interviene la primera potencia; sobre cuerpos imperfectos interviene la potencia $p^f$-ésima, como se mostrará en § 6, Teoremas XIII y XIV.} e igualmente la de la norma y la forma de divisores elementales de un ideal primario con $\frakp$ como ideal primo asociado.

Finalmente, si dos ideales primos coinciden en su forma de divisores elementales $P^{(i)}$, entonces, como consecuencia de la descomposición anterior, coinciden en sus ceros; por tanto son mutuamente divisibles y por consiguiente idénticos.


\clearpage

\subsection*{§ 4. Formulación aritmética del concepto de dimensión}

Una primera formulación del concepto de dimensión que es independiente de la introducción de ideales transformados viene dada por

\textbf{Teorema V.} \emph{La dimensión de un ideal primo viene dada por el rango algebraico de su cuerpo de clases residuales. La dimensión máxima de un ideal es igual al mayor de los números de dimensión de sus ideales primos asociados.}

La primera parte del Teorema V fue probada en la primera observación al Teorema IV, donde se mostró que el rango algebraico del cuerpo de clases residuales $(\mathcal R)$ de $\frakp$ es igual al del cuerpo de clases residuales $(\mathcal R)$ de $\frakp$, y por tanto igual a la dimensión de $\frakp$, de acuerdo con § 1, 8. Para la demostración de la segunda parte, sea $\frakm=[\frakq_1,\ldots,\frakq_a]$ una representación más corta por mayores componentes primarias. Entonces, ante todo, ningún $\frakq$, como divisor de $\frakm$, puede poseer dimensión más alta que $\frakm$. Pero el producto de todos los $\frakq$ es divisible por $\frakm$, y por tanto también lo es el producto de todas las formas de divisores elementales de los distintos $\frakq$. Si, entonces, $n-i$ es el mayor de los números de dimensión de los $\frakq$, entonces $\frakm$ contiene un polinomio $G^{(i)}\ne0$ y por consiguiente no puede tener dimensión mayor que $n-i$. Los números de dimensión de los $\frakq$ coinciden, por el Teorema I, con los de sus ideales primos asociados. Así, si, sin introducir ideales transformados, se define la dimensión de un ideal primo por el rango algebraico de su cuerpo de clases residuales, la segunda parte del Teorema V da la definición de la dimensión máxima de un ideal arbitrario.

Para captar, por otra parte, el concepto de dimensión mediante cadenas de ideales primos ---de nuevo independientemente de la introducción de ideales transformados---, hay que mostrar, completando el Teorema II:

\textbf{Teorema VI.} \emph{Todo ideal primo distinto del ideal unidad posee ---después de una posible adjunción de indeterminadas--- al menos un ideal primo como divisor cuya dimensión ha disminuido exactamente en una unidad.}

En efecto, si $\frakp$ tiene dimensión $(n-i)>0$ y $v$ denota una nueva indeterminada, entonces
\[
 \frakc=(\frakp,x_i-v)
\]
es un ideal con la propiedad requerida. Esto se sigue de una correspondencia isomorfa. Se tiene $\frakc=(\frakp^*,x_i-v)$, donde $\frakp^*$ surge de $\frakp$ reemplazando $x_i$ por la indeterminada $v$ y adjuntando $v$ al dominio de coeficientes; las clases residuales se obtienen del mismo modo reemplazando la clase $(x_i)$ por $(v)$. Pero $\frakp$ pasa a sí mismo cuando se adjunta $x_i$ a $P(u)$; pues de
\[
 h(x_i)f(x)\equiv0\pmod{\frakp},\qquad h(x_i)\ne0
\]
se obtiene necesariamente $f(x)\equiv0\pmod{\frakp}$. De lo contrario $\frakp$ contendría un polinomio que depende sólo de $x_i$, y por consiguiente también uno que depende sólo de $x_n$, contra la suposición de que tendría dimensión a lo sumo cero. Bajo la asignación $v\sim x_i$, por tanto, la clase cero módulo $\frakc$ corresponde necesariamente a la clase cero módulo $\frakp$, y por supuesto también a la inversa. Así hay un isomorfismo entre el sistema de clases residuales módulo $\frakp$ y el sistema módulo $\frakc$; este último no tiene divisores de cero, de modo que $\frakc$ es un ideal primo. Bajo este isomorfismo, la dependencia algebraica entre $(x_i)$ y $(P(u,x_{i+1},\ldots,x_n))$ corresponde a una relación entre $(x_{i+1}),\ldots,(x_n)$ respecto de $(P(u,v))$, mientras que para $\lambda=1,\ldots,i-1$ las dependencias entre $(x_\lambda)$ y $(P(u,x_{i+1},\ldots,x_n))$ permanecen como tales y, por tanto, no rebajan ulteriormente el rango algebraico. El rango algebraico ha disminuido exactamente en una unidad. Si $\frakp$ era de dimensión cero, $\frakc$ se convierte en el ideal unidad; el resultado anterior sigue siendo válido. En el dominio original $\bar{\frakS}$, en consecuencia,
\[
 \bar{\fraka}=(\bar{\frakp},v+v_1y_1+\cdots+v_ny_n),
\]
donde $v_\lambda$ se pone por $-t_{i\lambda}$, es un ideal del tipo requerido. Si, en particular, $P$ contiene infinitos elementos, entonces $v,v_1,\ldots,v_n$ pueden siempre especializarse en elementos de $P$ de modo que la norma y la forma de divisores elementales, y por tanto la dimensión de $\bar{\fraka}$, se hagan iguales a las del ideal especializado $\frakb$;\footnote{Compárese H.-N., el último párrafo del § 7.} la forma de divisores elementales, sin embargo, no tiene por qué seguir siendo una función prima. Por el Teorema V, $\frakb$ posee no obstante al menos un ideal primo asociado de la dimensión deseada; pasando de nuevo a $\bar{\frakS}$, el Teorema VI vale allí por tanto sin adjuntar indeterminadas.

Los Teoremas II y VI dan inmediatamente la formulación por cadenas buscada:

\textbf{Teorema VII.} \emph{Un ideal primo $\frakp$ tiene dimensión $n-i$ si ---después de una posible adjunción de indeterminadas--- existe al menos una cadena de $1+(n-i)+1$ ideales primos}
\[
 \frakp_0=\frakp,\quad \frakp_1,\ldots,\frakp_{n-i},\quad \frakp_{n-i+1},
\]
\emph{cada uno de los cuales es divisor propio del inmediatamente precedente, mientras que no existe ninguna cadena semejante con mayor número de miembros. Esta definición vale también en el dominio original, y sin introducir indeterminadas, cuando $P$ contiene infinitos elementos.}\footnote{Esta definición de la dimensión por cadenas da, en el caso de un cuerpo de números algebraicos, dimensión cero para los ideales primos ordinarios, dimensión $-1$ para el ideal unidad y dimensión uno para el ideal cero. En efecto, este último satisface también la definición de ideal primo: en un cuerpo un producto de cantidades no nulas es siempre distinto de cero. El Teorema II sigue siendo válido, con esta definición de dimensión, en los cuerpos de números algebraicos.}

Primero está claro que el último miembro de la cadena debe ser igual al ideal unidad ---que satisface la definición de ideal primo---, pues de lo contrario la cadena podría prolongarse adjuntando $\frako$; para $\frako$ mismo el Teorema VII da dimensión $-1$, de acuerdo con § 1, 8. Sea ahora $\frakp$ distinto de $\frako$ y de dimensión $n-i$. Entonces el número de miembros de la cadena puede ser a lo sumo $1+(n-i)+1$, ya que por el Teorema II no pueden aparecer dos ideales primos de la misma dimensión. Recíprocamente, por el Teorema VI, después de una posible adjunción de nuevas indeterminadas, existe al menos una cadena de esta longitud, pues la dimensión ha disminuido exactamente en una unidad en cada miembro de la cadena. Si se toma el Teorema VII como definición de dimensión ---definición que evidentemente puede enunciarse exactamente del mismo modo en el dominio original y para la cual, por el Teorema VI, la introducción de indeterminadas se vuelve innecesaria cuando $P$ contiene infinitos elementos---, entonces la definición de la dimensión máxima de un ideal arbitrario viene dada de nuevo por la segunda parte del Teorema V.

\subsection*{§ 5. Clasificación de los ideales fundamentales y de sus componentes en la teoría general de ideales. Ceros de ideales}

La cuestión es caracterizar los ideales fundamentales y sus componentes (§ 1, 7) mediante componentes aisladas de la descomposición (§ 1, 13) en el sentido de la teoría general de ideales. Esta caracterización mostrará el paralelismo entre los teoremas de descomposición para las normas y los de la teoría general de ideales, y hará posible enunciar, para ideales arbitrarios, la teoría de los ceros dada en § 3 para ideales primos e ideales primarios. Para los ideales fundamentales se tiene

\textbf{Teorema VIII.} \emph{Los ideales fundamentales $\frakg_i$ de etapa $i$ son las componentes aisladas de la descomposición que están determinadas unívocamente por la totalidad de los ideales primos asociados de dimensiones $n-1,n-2,\ldots,n-i$. Si todos los ideales primos asociados tienen la misma dimensión, entonces el ideal tiene también esta dimensión; es decir, sólo aparece un factor $R^{(i)}$ de la norma. En general la norma tiene tantos factores $R^{(i)},R^{(i+\sigma)},R^{(i+\tau)},\ldots$ distintos de la unidad como números de dimensión distintos aparecen entre los ideales primos asociados.}

Anteponemos a la demostración el siguiente

\textbf{Lema VI.} \emph{Si un ideal se representa como mínimo común múltiplo, $\frakm=[\fraka_1,\ldots,\fraka_a]$, entonces su ideal fundamental $\frakg_i$ es el mínimo común múltiplo de los ideales fundamentales $\frakh_i$ de etapa $i$ de los $\fraka$.}

Pues de
\[
 b^{(i+1)}f\equiv0\pmod{\frakm},\qquad b^{(i+1)}\ne0
\]
one obtains
\[
 b^{(i+1)}f\equiv0\pmod{\fraka_\lambda},\qquad b^{(i+1)}\ne0;
\]
por tanto $\frakg_i$ es divisible por $[\frakh_{i1},\ldots,\frakh_{ia}]$. Recíprocamente, de
\[
 c_\lambda^{(i+1)}f\equiv0\pmod{\fraka_\lambda},\qquad c_\lambda^{(i+1)}\ne0
\]
one obtains
\[
 c_1^{(i+1)}\cdots c_a^{(i+1)}f\equiv0\pmod{\frakm},\qquad
 c_1^{(i+1)}\cdots c_a^{(i+1)}\ne0,
\]
y de ahí la divisibilidad recíproca, lo que prueba el lema.

Ordenemos ahora por dimensión los ideales primos asociados de $\frakm$ (algunas dimensiones pueden, desde luego, estar ausentes, $j_\lambda=0$):
\[
 \frakp_{1,1},\ldots,\frakp_{1,j_1};\quad
 \frakp_{2,1},\ldots,\frakp_{2,j_2};\quad \ldots\quad
 \frakp_{n,1},\ldots,\frakp_{n,j_n},
\]
donde en general $\frakp_{\lambda,\kappa}$ denota un ideal primo de dimensión $n-\lambda$; sea $\frakq_{\lambda,\kappa}$ el ideal primario correspondiente que aparece en una descomposición fija dada de $\frakm$. Entonces, por definición, el ideal fundamental de etapa $i$ es igual al ideal unidad para todos los $\frakq_{\lambda,\kappa}$ para los que $\lambda>i$, mientras que para $\lambda\le i$ es, por el Teorema I, igual a $\frakq_{\lambda,\kappa}$. Por tanto, por el Lema VI,
\[
 \tag{5}
 \frakg_i=[\frakq_{1,1},\ldots,\frakq_{1,j_1};\ldots;\frakq_{i,1},\ldots,\frakq_{i,j_i}],
\]
y esta representación reconoce a $\frakg_i$ como una componente aislada de la descomposición, ya que ninguno de los ideales primos pertenecientes a $\frakg_i$ puede estar contenido en uno de los restantes ideales primos, todos los cuales tienen dimensión menor.

Si, en particular, sólo aparecen ideales primos asociados de dimensión $n-i$, entonces por (5) se tiene que $\frakg_0,\ldots,\frakg_{i-1}$ son iguales al ideal unidad, mientras que $\frakg_i=\frakg_{i+1}=\cdots=\frakm$; por tanto $R_\frakm=R^{(i)}$. Si, en cambio, aparecen ideales primos de dimensiones distintas, digamos que $j_i,j_{i+\sigma},j_{i+\tau},\ldots$ son no nulos, entonces por (5)
\[
 \frakg_0\ne \frakg_i\ne \frakg_{i+\sigma}\ne \frakg_{i+\tau}\ne\cdots,
\]
whereas
\[
 \frakg_0=\frakg_1=\cdots=\frakg_{i-1}=\frako,
 \qquad
 \frakg_i=\frakg_{i+1}=\cdots=\frakg_{i+\sigma-1},\ldots;
\]
por tanto los factores $R^{(i)},R^{(i+\sigma)},R^{(i+\tau)}$, y sólo éstos, son distintos de la unidad.\footnote{Teniendo en cuenta el Lema II, el Teorema VIII implica de nuevo que los ideales fundamentales ---como mínimos comunes múltiplos de ideales transformados--- se convierten en ideales transformados. Como el Teorema VIII descansa sólo en teoremas de H.-N. en los que este hecho no se usa, esto da al mismo tiempo una nueva demostración que revela la razón interna. En H.-N. el teorema se utiliza sólo en la teoría de los ceros, que aquí se desarrolla de nuevo.}

Por medio del Teorema VIII, el Teorema I se afina en

\textbf{Teorema IX.} \emph{Un ideal es primario si y sólo si su forma de divisores elementales ---y por tanto su norma--- es una función primaria.}

Que la condición se cumple para ideales primarios fue mostrado en el Teorema I. Recíprocamente, sea dada la forma de divisores elementales de $\frakm$ por $Q^{(i)}=(P^{(i)})^\rho$, donde $P^{(i)}$ denota una función prima. Por el Teorema VIII, todos los ideales primos de $\frakm$ tienen entonces la misma dimensión $n-i$. De $\frakm=[\frakq_1,\ldots,\frakq_a]$ se obtiene
\[
 Q^{(i)}\equiv0\pmod{\frakq_\lambda};
 \quad\hbox{por tanto}\quad
 (P^{(i)})^\rho\equiv0\pmod{\frakp_\lambda},
 \quad\hbox{y por consiguiente}\quad P^{(i)}\equiv0\pmod{\frakp_\lambda}.
\]
Como todos los $\frakp_\lambda$ tienen dimensión $n-i$ y $P^{(i)}$ es una función prima, $P^{(i)}$ se convierte en la forma de divisores elementales de cada $\frakp_\lambda$. Así, por el corolario al Teorema IV, todos los $\frakp_\lambda$ coinciden. En consecuencia, puesto que se trata de mayores componentes primarias, sólo puede aparecer un $\frakp$; $\frakm$ es primario, sin excluir el caso primo especial $\rho=1$. Los Teoremas VIII y IX conducen a la clasificación de las componentes de los ideales fundamentales de la manera siguiente.

\textbf{Teorema X.} \emph{La componente $\frakr_\lambda$ del ideal fundamental $\frakg_i$ correspondiente al mayor factor primario $Q_\lambda^{(i)}=(P_\lambda^{(i)})^{\rho_\lambda}$ de $R^{(i)}$ viene dada por la componente aislada de la descomposición que está determinada unívocamente por los ideales primos asociados de dimensiones $n-1,\ldots,n-i+1$ y por aquel ideal primo de dimensión $n-i$ cuya forma de divisores elementales se hace igual a $P_\lambda^{(i)}$. Los ideales primos asociados y los mayores factores primarios de la norma se corresponden uno a uno.}

Por § 1, 7 los $\frakr_\lambda$ coinciden con su ideal fundamental de etapa $i$ y poseen $\frakg_{i-1}$ como ideal fundamental de etapa $(i-1)$. Por tanto, por el Teorema VIII, respectivamente por (5), admiten una representación más corta
\[
 \frakr_\lambda=[\frakg_{i-1},\frakq_{\lambda,1},\ldots,\frakq_{\lambda,t_\lambda}],
\]
donde todos los $\frakq$ tienen dimensión $n-i$ y pertenecen a ideales primos distintos. Por la definición de $Q_\lambda^{(i)}$ ---teniendo en cuenta que $\frakr_\lambda$ coincide con su ideal fundamental de etapa $i$--- se tiene
\[
 Q_\lambda^{(i)}\frakg_{i-1}\equiv0\pmod{\frakq_{\lambda,\kappa}},
 \quad(\kappa=1,\ldots,t_\lambda),
\]
and consequently
\[
 (Q_\lambda^{(i)})^{\rho_\lambda}\frakr_\lambda\equiv0\pmod{\frakq_{\lambda,\kappa}};
 \quad\hbox{por tanto}\quad
 (P_\lambda^{(i)})^{\rho_\lambda}\frakr_\lambda\equiv0\pmod{\frakp_{\lambda,\kappa}},
\]
y así $P_\lambda^{(i)}\equiv0\pmod{\frakp_{\lambda,\kappa}}$. Como en la demostración del Teorema IX se sigue que en la representación de $\frakr_\lambda$ sólo puede aparecer un $\frakq_\lambda$ y un $\frakp_\lambda$ asociado de dimensión $n-i$, y que la forma de divisores elementales de $\frakp_\lambda$ viene dada por $P_\lambda^{(i)}$. Resta probar que $\frakp_\lambda$ es un ideal primo asociado de $\frakm$ y que $\frakq_\lambda$ aparece efectivamente en al menos una representación de $\frakm$ por mayores componentes primarias. Entonces $\frakr_\lambda$ queda reconocido también como la componente aislada de la descomposición, puesto que ningún ideal primo puede ser absorbido por uno de igual o menor dimensión, a saber, como la componente determinada por $\frakp_\lambda$ y por todos los ideales primos asociados de dimensión mayor.

Para esta prueba basta mostrar que $\frakq_\lambda$ aparece en una representación más corta de $\frakg_i$, puesto que tal representación siempre puede completarse, por el Teorema VIII, hasta una representación más corta de $\frakm$. Así, basta mostrar que
\[
 \frakg_i=[\frakr_1,\ldots,\frakr_a]=[\frakg_{i-1},\frakq_1,\ldots,\frakq_a]
\]
se convierte en una representación más corta. Si, por ejemplo, $\frakq_\lambda$ fuera absorbido por su complemento, de modo que $\frakg_{i-1}\frakq_1\cdots\frakq_{\lambda-1}\frakq_{\lambda+1}\cdots\frakq_a$ fuera divisible por $\frakq_\lambda$, entonces, como $\frakg_{i-1}\not\equiv0\pmod{\frakq_\lambda}$, se seguiría que una potencia de algún $\frakq_\mu$ es divisible por $\frakq_\lambda$ para $\mu\ne\lambda$. Así, los ideales primos asociados $\frakp_\mu$ y $\frakp_\lambda$, ambos de la misma dimensión, tendrían que coincidir por el Teorema II; pero esto es imposible, pues sus formas de divisores elementales $P_\mu^{(i)}$ y $P_\lambda^{(i)}$ son distintas por definición. La representación de $\frakg_i$, a la que por el Teorema VIII corresponden todos los ideales primos asociados de dimensión $n-i$, muestra además que a cada tal ideal primo le corresponde realmente un mayor factor primario de la norma. Queda probada por tanto la correspondencia biunívoca.

El Teorema X da inmediatamente el traslado del Teorema IV:

\textbf{Teorema XI.} \emph{Los mayores factores primarios individuales $Q_\lambda^{(i)}=(P_\lambda^{(i)})^{\rho_\lambda}$ de la norma admiten una representación como producto}
\[
 Q_\lambda^{(i)}=\prod \bigl(x_i-t_{1i}\bar y_{1\nu}-\cdots-t_{ni}\bar y_{n\nu}\bigr)^{\delta_\lambda\rho_\lambda}
 =\prod \bigl(t_{1i}(y_1-\bar y_{1\nu})+\cdots+t_{ni}(y_n-\bar y_{n\nu})\bigr)^{\delta_\lambda\rho_\lambda},
\]
\emph{donde $\bar y_{1\nu},\ldots,\bar y_{n\nu}$ representa un sistema asociado de ceros del ideal primo asociado $\frakp_\lambda$, en las indeterminadas originales $y$, independiente de $t_{1i},\ldots,t_{ni}$, y $\delta_\lambda$ tiene el valor uno o $p^{f_\lambda}$. Así queda dada completamente la descomposición de la norma en factores lineales. El grado de $Q_\lambda^{(i)}$ da, sobre el subcuerpo de clases residuales $(P(u,x_{i+1},\ldots,x_n))$ obtenido por formación de cocientes, el grado del subanillo de clases residuales, representado por clases residuales de $\frakg_{i-1}$, módulo la componente aislada $\frakr_\lambda$; para ideales primarios éste es por tanto el anillo de todas las clases residuales. Al adjuntar nuevas indeterminadas $v$ se obtiene además la descomposición}
\[
 Q_\lambda^{(i)}(u,v)=\prod (z-v_1\bar x_{1\nu}-\cdots-v_i\bar x_{i\nu})^{\delta_\lambda\rho_\lambda}.
\]

Como $P_\lambda^{(i)}$ ha sido reconocido por el Teorema X como la forma de divisores elementales de $\frakp_\lambda$, esto no es de hecho sino la sustitución de la descomposición del Teorema IV. La observación sobre el grado es sólo otra formulación de lo dicho en § 1, 5 y 7.\footnote{Asimismo, la formulación de la multiplicidad mencionada en la nota 2 de H.-N. es una consecuencia inmediata del Teorema X. Pues, después de adjuntar $x_{i+1},\ldots,x_n$, el sistema de clases residuales $\frakg_{i-1}\mid\frakr_\lambda$ es isomorfo a $\frakt_\lambda\mid\frakr_\lambda$, donde se pone $\frakt_\lambda=[\frakg_{i-1},\frakr_1,\ldots,\frakr_{\lambda-1},\frakr_{\lambda+1},\ldots,\frakr_a]$, como se sigue del retorno a módulos de formas lineales (H.-N., Teorema V), en vista de la primalidad relativa de los factores $Q_\mu^{(i)}$. Directamente se muestra que $\frakt_\lambda\mid\frakr_\lambda$ es isomorfo a $\frakt_\lambda\mid\frakq_\lambda$; aquí $\frakt_\lambda$ surge del complemento de $\frakq_\lambda$ después de adjuntar $x_{i+1},\ldots,x_n$.}

\subsection*{§ 6. Caracterización de los ideales primos y de los ideales primarios propios mediante forma de divisores elementales y norma}

En el Teorema IX los ideales primarios fueron caracterizados por forma de divisores elementales y norma, pero sólo de tal modo que el caso primo debía considerarse como caso especial del primario. Ahora necesitamos separar los ideales primos de los ideales primarios propios; las consideraciones correrán paralelas a las del § 3. Las condiciones que son necesarias o suficientes se enuncian fácilmente:

\textbf{Teorema XII.} \emph{Una condición necesaria para un ideal primo es que su forma de divisores elementales sea una función prima; una condición suficiente es que su norma sea una función prima. En otras palabras, la forma de divisores elementales de un ideal primo es siempre una función prima, y la norma de un ideal primario propio es siempre una función primaria propia.}

Que la forma de divisores elementales de un ideal primo es una función prima ya se mostró en el Teorema I. Supóngase ahora que la norma $R_{\frakq}$ es igual a una función prima. Hay que mostrar que, bajo esta hipótesis, todo divisor propio $\fraka$ de $\frakq$ tiene dimensión máxima menor; entonces $\frakq$ queda reconocido como ideal primo por el Teorema II. En efecto, por § 1, 6 el primer factor de la norma $R_\fraka$ que difiere del correspondiente de $R_\frakq$ es un divisor propio. Como $R_\frakq=P^{(i)}$, el factor $R^{(i)}$ de $R_\fraka$ debe ser por tanto igual a un divisor propio de $P^{(i)}$, y por ello a la unidad; $\fraka$ tiene dimensión máxima menor.

Una segunda demostración sencilla descansa en el siguiente lema, en el que se apoya también el próximo teorema.

\textbf{Lema VII.} \emph{El sistema de clases residuales representado por polinomios $H^{(i)}$ módulo la forma de divisores elementales $E^{(i)}$ de un ideal primario ---más generalmente, de un ideal que tiene sólo ideales primos asociados de dimensión $n-i$--- es isomorfo a un subsistema de las clases residuales módulo este ideal. El grado de la forma de divisores elementales da el grado de este sistema de clases residuales respecto del cuerpo de cocientes $(P(u,x_{i+1},\ldots,x_n))$.}

En efecto, se define una correspondencia asignando entre sí las clases representadas por los mismos polinomios $H^{(i)}$; las sumas y productos corresponden a sumas y productos. La asignación es además biunívoca. Pues todos los polinomios $H^{(i)}$ que pertenecen a la clase cero módulo el ideal son, por definición, divisibles por $E^{(i)}$; por tanto la clase cero corresponde a la clase cero módulo $E^{(i)}$, y también a la inversa.

Sea ahora la norma de un ideal una función prima $P^{(i)}$, por tanto idéntica a la forma de divisores elementales. Los grados de los sistemas de clases residuales módulo la forma de divisores elementales $P^{(i)}$ y módulo $\frakq$ coinciden entonces respecto de $(P(u,x_{i+1},\ldots,x_n))$; el subsistema isomorfo al sistema de clases residuales módulo $P^{(i)}$ agota por tanto las clases residuales módulo $\frakq$, y como este sistema no puede tener divisores de cero, $\frakq$ es un ideal primo.

En general el Teorema XII no puede afinarse hasta obtener condiciones necesarias y suficientes, como se mostrará abajo mediante ejemplos. Esto es posible, sin embargo, cuando el dominio de coeficientes $P$ es un cuerpo perfecto (§ 1, 14):

\textbf{Teorema XIII.} \emph{Si el dominio de coeficientes $P$ es un cuerpo perfecto, entonces la norma y la forma de divisores elementales de un ideal primo se convierten en funciones primas y por consiguiente coinciden; la norma y la forma de divisores elementales de un ideal primario propio se convierten en funciones primarias propias. Así, sobre un cuerpo perfecto, la condición de que la forma de divisores elementales sea una función prima es necesaria y suficiente para que el ideal sea primo; en la condición puede ponerse también la norma en lugar de la forma de divisores elementales.}

Teniendo en cuenta el Teorema XII, el Teorema XIII quedará probado una vez se muestre que, sobre un cuerpo perfecto, el ideal se vuelve primo tan pronto como su forma de divisores elementales es una función prima, y que entonces la norma es igual a esta función prima. En el Lema V, § 3, se mostró que, sobre un cuerpo perfecto, la forma de divisores elementales se convierte en una función prima de primera especie tan pronto como es en absoluto una función prima. Sea por tanto $(\mathcal R)$ el cuerpo de clases residuales obtenido al adjuntar $(x_i)$ a $(P(u,x_{i+1},\ldots,x_n))$, es decir, $(P(u,x_i,x_{i+1},\ldots,x_n))$ módulo $P^{(i)}$. Entonces $(x_i)$ es un cero de la función prima de primera especie $(P^{(i)})$, y de ahí, como muestra por ejemplo la fórmula de Lagrange, $(y_1),\ldots,(y_n)$ están contenidos en $(\mathcal R)$. El subsistema de las clases residuales módulo el ideal $\frakq$ considerado que, por el Lema VII, es isomorfo a $(\mathcal R)$ ---donde sólo aparecen denominadores de $(P(u,x_{i+1},\ldots,x_n))$--- contiene por tanto las clases residuales $(y_1),\ldots,(y_n)$ y agota así el sistema de clases residuales módulo $\frakq$. Como, al ser isomorfo a $(\mathcal R)$, no puede tener divisores de cero, $\frakq$ es un ideal primo. El grado de este cuerpo de clases residuales $(\mathcal R)$ de $\frakq$ respecto de $(P(u,x_{i+1},\ldots,x_n))$ es igual al grado de la norma de $\frakq$; por otra parte, por el isomorfismo con $(\mathcal R)$, es igual al grado de la forma de divisores elementales $P^{(i)}$. Por tanto, norma y forma de divisores elementales deben coincidir.

Para dominios de coeficientes imperfectos, el Teorema XII permite todavía el siguiente afinamiento:

\textbf{Teorema XIV.} \emph{Si el dominio de coeficientes $P$ es un cuerpo imperfecto de característica $p$, entonces la norma de un ideal primo es la potencia $p^g$-ésima de la forma de divisores elementales; se tiene $0\le g\le(i-1)f$, donde $f$ denota el exponente de la forma de divisores elementales y $n-i$ la dimensión del ideal primo.}

El Teorema XIV equivale a afirmar que el cuerpo de clases residuales $(\mathcal R)$ de $\frakp$ tiene grado $p^g$ ($g\ge0$) respecto del subcuerpo $(\mathcal R^{\prime})=(P(u,x_i,x_{i+1},\ldots,x_n))$ isomorfo a $(\mathcal R)$. Esto se sigue directamente de los teoremas de Steinitz enunciados en § 1, 15, más precisamente de

\textbf{Lema VIII.} \emph{Sea $\Lambda$ una extensión finita del cuerpo imperfecto $\Omega$ de grado reducido $r$ y exponente $f$; sea $\Lambda_0$ el cuerpo de primera especie contenido en $\Lambda$, y sea $M$ un cuerpo intermedio entre $\Lambda_0$ y $\Lambda$. Entonces el grado de todo elemento de $\Lambda$ respecto de $M$ es una potencia de $p$, donde $p$ denota la característica de $\Omega$.}

Pues la potencia $p^{f^{\prime}}$-ésima ($f^{\prime}\le f$) de todo elemento $z$ de $\Lambda$ pertenece a $\Lambda_0$; por tanto $z$ es un cero de una función prima $G(t)=(t-z)^{p^{f^{\prime}}}$ con coeficientes en $\Lambda_0$. Sea $H(t)=(t-z)^\delta$ la función prima en $M$ con cero $z$. Entonces $H(t)$ tiene también coeficientes en $\Lambda_0(z)$, y por tanto los coeficientes de $H(t)$ pertenecen al cuerpo intersección de $M$ y $\Lambda_0(z)$, luego a un cuerpo intermedio entre $\Lambda_0$ y $\Lambda_0(z)$. El grado de $z$ respecto de $M$ es por consiguiente igual al grado respecto de este cuerpo intermedio y de ahí, como divisor de $p^{f^{\prime}}$, es una potencia de $p$.

Para la demostración del Teorema XIV sólo resta observar que el cuerpo $(\mathcal R)=(P(u,x_i,x_{i+1},\ldots,x_n))$ debe tener grado $r\cdot p^f$ respecto de $(P(u,x_{i+1},\ldots,x_n))$, si $r$ denota el grado reducido y $f$ el exponente de $(\mathcal R)$; en consecuencia, el cuerpo de primera especie $(\mathcal R_0)$ contenido en $(\mathcal R)$ debe ser un subcuerpo de $(\mathcal R^{\prime})$. Puesto que en $(\mathcal R)$ existen elementos de grado $r\cdot p^f$ (§ 1, 15), $(x_i)$ debe ser necesariamente de este grado, ya que todos los elementos de $(\mathcal R)$ surgen especializando los $t_{\mu\nu}$ en $(x_i)$. El exponente de $(\mathcal R)$ coincide con el exponente de la forma de divisores elementales, y $(x_i)^{p^f}$ tiene grado $r$ y es un elemento de $(\mathcal R_0)$, por tanto un elemento primitivo de $(\mathcal R_0)$. Como $(\mathcal R)$ surge al adjuntar finitamente muchos elementos ---digamos $(x_1),\ldots,(x_{i-1})$--- a $(\mathcal R^{\prime})$, la aplicación finita y repetida del Lema VIII da la demostración deseada y al mismo tiempo la estimación para $g$.\footnote{Compárense las consideraciones paralelas en A. Ostrowski, Zur arithmetischen Theorie der algebraischen Größen, Gött. Nachr. 1919, pp. 279--298, en especial p. 288, donde el teorema correspondiente se enuncia sin demostración.}

Añadimos ahora los ejemplos que muestran que el Teorema XII no puede en general afinarse más allá del Teorema XIV. Conciernen a un ideal primo cuya norma es igual al cuadrado de la forma de divisores elementales ($g>0$), y a ideales primarios propios cuyas formas de divisores elementales son funciones primas. Este último ejemplo muestra también que aquí no existe un análogo del Teorema XIV, sino que la norma puede convertirse en una potencia arbitraria de la forma de divisores elementales.

Sea el dominio de coeficientes $P$ obtenido adjuntando dos indeterminadas $\lambda,\mu$ al cuerpo de clases residuales módulo $2$, y sea por tanto $P$ un cuerpo imperfecto de característica dos. Consideremos el ideal
\[
 \bar{\frakp}=(y_1^2+\lambda,\,y_2^2+\mu),
\]
que debe ser un ideal primo, puesto que el sistema de clases residuales $(\mathcal R)$ es isomorfo al cuerpo $P(\sqrt\lambda,\sqrt\mu)$. El cuerpo $(\mathcal R)$ tiene grado cuatro respecto de $(P)$, grado reducido $r=1$ y exponente $f=2$. Por tanto la forma de divisores elementales $P^{(2)}$ es de segundo grado, y la norma es de cuarto grado y en consecuencia el cuadrado de $P^{(2)}$, como puede comprobarse también directamente a partir de la representación modular (3), § 1, 5. Se obtiene
\[
 P^{(2)}=(u_{11}u_{22}+u_{12}u_{21})^2x_2^2+(u_{11}^2\mu+u_{21}^2\lambda),
\]
or, after a suitable multiplication,
\[
 x_2^2+(t_{12}^2\lambda+t_{22}^2\mu).
\]

Sea ahora $P$ el resultado de adjuntar la única indeterminada $\lambda$ al cuerpo de clases residuales módulo $2$, y tomemos como base el ideal
\[
 \bar{\frakq}=(y_1^2+\lambda,\,y_2^2+\lambda)=(y_1^2+\lambda,\,(y_1+y_2)^2).
\]
Como forma de divisores elementales, especializando $\mu=\lambda$ en la acabada de dar, se obtiene
\[
 P^{(2)}=x_2^2+\lambda(t_{12}^2+t_{22}^2),
\]
por tanto una función prima, ya que $t_{12}=(0),\ t_{22}=(1)$ conduce a la función prima $x_2^2+\lambda$. Pero $\bar{\frakq}$ es primario propio, puesto que contiene $(y_1+y_2)^2$ pero no $(y_1+y_2)$. La norma debe ser igual al cuadrado, y por tanto a la potencia $p$-ésima, de $P^{(2)}$, ya que el número de clases residuales de $\bar{\frakq}$ linealmente independientes sobre $P$ es cuatro.

Que este último análogo del Teorema XIV no siempre se cumple lo muestra el ejemplo siguiente. Sea $P$ el resultado de adjuntar la indeterminada $\lambda$ al cuerpo de clases residuales módulo $3$, y pongamos
\[
 \frakq=(y_1^3+\lambda,\,(y_1-y_2)^2).
\]
Entonces, como arriba, $\frakq$ es un ideal primario propio. La forma de divisores elementales se convierte ---teniendo en cuenta que $y_2^3+\lambda$ es también divisible por $\frakq$--- en
\[
 P^{(2)}=x_2^3+\lambda(t_{12}+t_{22})^3,
\]
una función prima; pero la norma es igual al cuadrado de $P^{(2)}$, puesto que hay seis clases residuales linealmente independientes módulo $\frakq$. Así, la potencia $p^g$-ésima de la forma de divisores elementales no aparece aquí.

El primer ejemplo muestra que, sobre un cuerpo imperfecto como dominio de coeficientes, pueden aparecer, exactamente como en los cuerpos de números algebraicos, ideales primos de grado superior al primero, donde $p^g$ debe llamarse el grado del ideal primo. Los ejemplos restantes deben entenderse como análogos de ideales de ramificación: un número primo puede ser divisible por una potencia superior de un ideal primo, por un ideal primario. Pues, como se mostró en la nota 10, la forma de divisores elementales corresponde al menor número racional entero divisible por un ideal en un cuerpo de números algebraicos.

\subsection*{§ 7. Ideales primos absolutos}

Un ideal primo con coeficientes en $P$ se llama \emph{ideal primo absoluto} si sigue siendo un ideal primo en el cuerpo algebraicamente cerrado $A$ al que puede extenderse $P$.\footnote{Un cuerpo algebraicamente cerrado es, como es sabido, un cuerpo en el que todo polinomio en una indeterminada se descompone en factores lineales. Para la existencia y la unicidad esencial del cuerpo algebraicamente cerrado perteneciente a un cuerpo arbitrario, véase Steinitz.} Correspondientemente, una función prima $P$ con coeficientes en $P$ se llama \emph{función prima absoluta} ---polinomio absolutamente irreducible--- si $P$ sigue siendo una función prima en $A$. La caracterización de los ideales primos absolutos viene dada por

\textbf{Teorema XV.} \emph{Una condición necesaria y suficiente para un ideal primo absoluto es que su forma de divisores elementales ---que puede suponerse entera y primitiva en las $u$--- se convierta en una función prima absoluta como polinomio en las $x$ y las $u$. La forma de divisores elementales se vuelve idéntica a la norma, de modo que la condición puede enunciarse también para la norma.}

Para la demostración obsérvese primero que, en general, norma y forma de divisores elementales se conservan bajo una extensión algebraica del dominio de coeficientes. En efecto, los factores individuales $R^{(i)}$, respectivamente $E^{(i)}$, están definidos como máximos comunes divisores en sentido polinómico, propiedad que se conserva por extensión algebraica del dominio de coeficientes. Así $P^{(i)}$ sigue siendo la forma de divisores elementales de $\frakp$ al pasar de $P$ a $A$, y por tanto al pasar a un cuerpo perfecto como dominio de coeficientes, ya que todo cuerpo algebraicamente cerrado es perfecto, como muestra inmediatamente la definición. Por el Teorema XIII, una condición necesaria y suficiente para que $\frakp$ sea un ideal primo absoluto es que $P^{(i)}$ se convierta en una función prima respecto de $A(u)$; esto es, en una función prima absoluta como polinomio en las $x$ y las $u$, si $P^{(i)}$ se supone entera y primitiva en las $u$. Al mismo tiempo, por el Teorema XIII, $P^{(i)}$ se vuelve idéntica a la norma de $\frakp$ tan pronto como se toma $A$ como dominio de coeficientes. Por lo que acaba de observarse, lo mismo vale sobre $P$. Esto prueba el Teorema XV.

Para funciones primas absolutas $P$ con coeficientes en un cuerpo de números algebraico (finito) $\mathfrak K$, se tiene ahora el teorema\footnote{A. Ostrowski, loc. cit. (nota 21), lema p. 296. Una demostración más sencilla está en E. Noether, Ein algebraisches Kriterium für absolute Irreduzibilität, Math. Ann. 85 (1922), pp. 26--33, núm. 7.} de que $P$ sigue siendo una función prima, más precisamente una función prima absoluta, módulo todo ideal primo de $\mathfrak K$, salvo a lo sumo finitos ideales primos de $\mathfrak K$. El traslado de este teorema a ideales primos absolutos descansa en

\textbf{Teorema XVI.} \emph{Si, como dominio de coeficientes, se toma, en lugar de un cuerpo, el anillo $\mathfrak o^*$ de todos los enteros algebraicos de un cuerpo de números algebraico (finito) $\mathfrak K$, entonces la norma y la forma de divisores elementales de un ideal polinómico se conservan como norma y forma de divisores elementales módulo todo ideal primo de $\mathfrak K$, salvo a lo sumo finitos ideales primos de $\mathfrak K$.}

Para la demostración debe modificarse, respecto de § 1, 1, el dominio polinómico en que todo se basa. Sea $\bar{\frakS}$ el conjunto de todos los polinomios en $y_1,\ldots,y_n$ con coeficientes en $\mathfrak o^*$, por tanto enteros algebraicos de $\mathfrak K$; sea el dominio de coeficientes de $\frakS$ $\mathfrak o^*[u]$, es decir, todos los polinomios en $u$ con coeficientes en $\mathfrak o^*$. Si $\bar{\frakm}$ es un ideal en $\bar{\frakS}$, pasa por (1) a un ideal transformado $\frakm$ en $\frakS$. Hay que mostrar que en la formación de la norma sólo aparecen finitamente muchos polinomios de $\mathfrak o^*[u]$ en los denominadores; de ahí se sigue directamente el Teorema XVI.

Obsérvese primero que la representación modular (3) de § 1, 5, para la formación de la norma individual, equivale a reducir potencias de $x_{i-1}$ módulo un polinomio regular en $x_{i-1}$,
\[
 C^{(i-1)}=U_{i-1}(u)x_{i-1}^{e_{i-1}}+\hbox{términos inferiores},
\]
\footnote{Compárese H.-N., § 4.} donde todos los coeficientes de $C^{(i-1)}$, en particular $U_{i-1}$, pueden suponerse elementos de $\mathfrak o^*[u]$. Como se trata de la reducción sucesiva de $x_1,\ldots,x_{i-1}$, la representación modular entera en las $x$ es también entera en $\mathfrak o^*[u]$, salvo por productos de potencias
\[
 U_1^{\lambda_1}\cdots U_{i-1}^{\lambda_{i-1}}
\]
en el denominador. El producto $U_1U_2\cdots U_n$, considerado como polinomio en las $u$, define por sus coeficientes un ideal $\mathfrak r^*$ de $\mathfrak K$, que por tanto es divisible sólo por finitos ideales primos de $\mathfrak K$. Si $\frakp^*$ se elige distinto de estos finitos ideales y se usa el cuerpo de clases residuales módulo $\frakp^*$ como dominio de coeficientes, la representación modular original se conserva reemplazando cada número de $\mathfrak o^*$ por su clase residual módulo $\frakp^*$.

Además, puesto que la norma y la forma de divisores elementales están determinadas sólo salvo factores de $\mathfrak o^*[u]$, también pueden suponerse divisibles por $\frakm$ respecto de $\mathfrak o^*[u]$. Supóngase, bajo esta convención, que por ejemplo
\[
 R^{(i)}=V_i(u)x_i^{e_i}+\hbox{términos inferiores},
\]
de modo que, en la divisibilidad de todos los determinantes de $\rho$ filas del módulo $\mathfrak M_{i-1}^*$ de rango $\rho$ determinado por $C^{(i-1)}$, sólo aparecen potencias de $V_i$ en el denominador además de los $U_1^{\lambda_1}\cdots U_{i-1}^{\lambda_{i-1}}$. Elíjase $\frakp^*$ además distinto de los finitos ideales primos de $\mathfrak K$ que aparecen en el ideal determinado por $V_1V_2\cdots V_n$. Entonces $R^{(i)}$ sigue siendo un divisor común de estos determinantes cuando se toma como dominio de coeficientes el cuerpo de clases residuales módulo $\frakp^*$, y el rango de $\mathfrak M_{i-1}^*$ se conserva, ya que $C^{(i)}$ era un determinante de $\rho$ filas de $\mathfrak M_{i-1}^*$.

Finalmente debe elegirse $\frakp^*$ de modo que $R^{(i)}$ siga siendo en cada caso el máximo común divisor. Recorra $D^{(i)}$ todos los determinantes de $\rho$ filas de $\mathfrak M_{i-1}^*$; por hipótesis
\[
 U_1^{\lambda_1}\cdots U_{i-1}^{\lambda_{i-1}}V_i^{\rho}D^{(i)}=R^{(i)}T^{(i)}
\]
con los $T^{(i)}$ sin divisor común que contenga $x$ en sentido polinómico. En consecuencia, el ideal derivado de todos los $T^{(i)}$ contiene un polinomio
\[
 G^{(i+1)}=W_i(u)x_{i+1}^{m_i}+\hbox{términos inferiores},\qquad W_i\ne0.
\]
Que $G^{(i+1)}$ pueda suponerse regular en $x_{i+1}$ se sigue de componer la transformación (1) con una de $x_{i+1},\ldots,x_n$ cuyos coeficientes son nuevas indeterminadas. Si, entonces, $\frakp^*$ se elige también distinto de los finitos ideales primos contenidos en el ideal determinado por $W_1W_2\cdots W_n$, cuando se usa como dominio de coeficientes el cuerpo de clases residuales módulo $\frakp^*$, $R_\frakm$ se conserva efectivamente como norma. De modo completamente correspondiente, $\frakp^*$ puede elegirse aún más restrictivamente para que $E_\frakm$ siga siendo también la forma de divisores elementales. Esto prueba el Teorema XVI.

De los Teoremas XV y XVI, como generalización del teorema sobre funciones primas absolutas y usando ese teorema, se obtiene

\textbf{Teorema XVII.} \emph{Si $\frakp$ es un ideal primo absoluto con enteros algebraicos de un cuerpo de números algebraico (finito) $\mathfrak K$ como coeficientes, entonces $\frakp$ sigue siendo un ideal primo, más precisamente un ideal primo absoluto, módulo todo ideal primo de $\mathfrak K$, salvo a lo sumo finitos ideales primos de $\mathfrak K$.}

Para la demostración, elíjase la norma $R_\frakp$ de $\frakp$, como en el Teorema XVI, de modo que sea divisible por $\frakp$ también respecto de $\mathfrak o^*[u]$. Entonces
\[
 R_\frakp=T(u)P^{(i)}(x),
\]
donde $P^{(i)}(x)$ puede suponerse entero y primitivo en las $u$. Por el Teorema XV, $P^{(i)}$, considerado como polinomio en $u$ y $x$ con coeficientes en $\mathfrak K$, es una función prima absoluta. Por el teorema sobre funciones primas absolutas conserva esta propiedad módulo todo ideal primo de $\mathfrak K$, salvo a lo sumo finitos ideales primos de $\mathfrak K$. Elíjase por tanto $\frakp^*$ distinto de esos finitos ideales primos y, por el Teorema XVI, también distinto de otros finitos ideales primos de $\mathfrak K$. Entonces $P^{(i)}$ ---con sus coeficientes de $\mathfrak o^*$ reemplazados por las clases residuales correspondientes módulo $\frakp^*$--- se convierte en la norma de $\frakp$ cuando se toma como dominio de coeficientes $P$ el cuerpo de clases residuales módulo $\frakp^*$, y esta norma es una función prima absoluta. Tomando $P$ como dominio de coeficientes, $\frakp$ sigue siendo por tanto un ideal primo absoluto por el Teorema XV. Esto prueba el Teorema XVII.

\begin{center}
Gotinga, 8 de mayo de 1923.\\[1em]
(Recibido el 9 de marzo de 1923.)
\end{center}



\clearpage
\section*{25. Teoría de la eliminación y teoría de ideales}

\begin{center}
J. Ber. d. DMV 33 (1924), pp. 116--120
\end{center}

Me propongo mostrar en lo que sigue cómo puede construirse la teoría de la eliminación en completo paralelismo con la teoría de los ceros de polinomios de una variable. Esta construcción se basa en enlazar la teoría de la eliminación dada por \emph{Hentzelt} con los métodos de la teoría de cuerpos, tal como fueron desarrollados por \emph{Steinitz}, y con la teoría de ideales, en la forma que yo misma he desarrollado.\footnote{Exposición detallada y referencias bibliográficas en Math. Ann. 90 y en un trabajo de próxima aparición.}

Primero quiero esbozar la cuestión en el caso de los polinomios de una variable. Como dominio de coeficientes se fija de una vez por todas un cuerpo $P$, al cual se refieren los conceptos de función prima, etc. Aquí $P$ puede suponerse definido abstractamente en el sentido de Steinitz, pero en el caso especial puede naturalmente coincidir con el cuerpo de los números racionales o con el cuerpo de todos los números complejos, según la hipótesis antes usual en la teoría de la eliminación. La teoría de los ceros de polinomios de una variable se divide entonces en cuatro pasos:

1. \emph{El teorema de descomposición}, que permite reducir la cuestión a una cuestión sobre funciones primas, es decir, sobre polinomios irreducibles en $P$. Pues de
\[
 a(x)=p_1(x)^{\varrho_1}\cdots p_r(x)^{\varrho_r}=q_1(x)\cdots q_r(x)
\]
--donde los $p(x)$ designan funciones primas distintas, y los $q(x)$ por tanto ``máximas funciones primarias''-- se sigue que $a(x)$ se anula en todos y sólo en los ceros de las funciones primas $p(x)$.

2. \emph{La construcción del cuerpo de ceros de una función prima $p(x)$.} Tal cuerpo de ceros $\mathfrak R=P(\alpha)$, donde $\alpha$ designa un cero de $p(x)$, puede construirse siempre, según \emph{Steinitz}, como un cuerpo de extensión de $P$ isomorfo al cuerpo de clases residuales $(\mathfrak R)$ de $p(x)$. Recíprocamente, todo cuerpo de ceros situado en un cuerpo de extensión de $P$ es isomorfo al cuerpo de clases residuales $(\mathfrak R)$; por tanto, cuerpo de ceros y cuerpo de clases residuales son abstractamente idénticos.

3. \emph{La descomposición de $p(x)$ en factores lineales en el cuerpo de Galois determinado por $p(x)$.} Repitiendo la construcción indicada en 2. se llega al cuerpo de Galois $P(\alpha_1,\ldots,\alpha_m)$, esencialmente determinado de manera única, que es precisamente suficiente para que $p(x)$ se descomponga en factores lineales.

4. \emph{El significado de la multiplicidad para el polinomio inicial $a(x)$.} Por 1., 2. y 3., para un polinomio arbitrario $a(x)$ quedan resueltas la cuestión de los ceros y la correspondiente descomposición en factores lineales. Como multiplicidad puede considerarse el grado de la función primaria individual $q(x)$, esto es, el número de clases residuales módulo $q(x)$ linealmente independientes respecto de $P$; pues, en el caso del dominio de coeficientes algebraicamente cerrado $P$, donde los $p(x)$ se vuelven lineales, este número coincide precisamente con la multiplicidad de los ceros.

Paso ahora a la teoría general de la eliminación, y tomo por base el dominio $\mathfrak S$ de todos los polinomios en $x_1,\ldots,x_n$ con coeficientes en $P$; la teoría de la eliminación puede definirse entonces como la \emph{teoría de los ceros de un ideal} de este dominio polinómico. Por ideal $\mathfrak a$ entiendo, como de costumbre, un sistema de polinomios tal que, junto con $f$ y $g$, también $f-g$, y junto con $f$ también $Af$, pertenece a $\mathfrak a$, donde $A$ es un polinomio arbitrario de $\mathfrak S$. Por cero de $\mathfrak a$ entiendo todo sistema de valores $\alpha_1,\ldots,\alpha_n$, perteneciente a un cuerpo de extensión de $P$, tal que $f(\alpha)$ se anula para todo polinomio $f$ perteneciente al ideal. Si $f_1,\ldots,f_t$ designa una base ideal, siempre existente, de $\mathfrak a$ --de modo que $f=A_1f_1+\cdots+A_tf_t$ para todo $f$ de $\mathfrak a$--, entonces los ceros de $\mathfrak a$ son evidentemente idénticos a los ceros de $f_1,\ldots,f_t$. Como, recíprocamente, todo sistema finito de polinomios puede considerarse como base del ideal derivado de él, la definición de ceros dada arriba para la teoría de la eliminación es idéntica a la usual, pero evita la dificultad que reside en distinguir arbitrariamente una base. Para polinomios de una variable, donde sólo intervienen ideales principales, esa dificultad no aparece, porque el polinomio básico está determinado de manera única salvo una unidad, es decir, salvo un factor tomado de $P$; aun allí, sin embargo, la concepción ideal dará una comprensión más exacta.

La teoría de la eliminación se divide ahora, correspondiendo al caso de una variable, en cuatro pasos:

I. \emph{El teorema de descomposición} --la representación de un ideal por máximas componentes primarias--, que permite reducir la cuestión a una cuestión sobre ideales primos. Un ideal se llama primo cuando divide a un producto sólo si divide al menos uno de los factores; se llama primario cuando divide al menos un factor o una potencia de cada factor. A todo ideal primario $\mathfrak q$ le pertenece uno y sólo un ideal primo asociado $\mathfrak p$, que es divisor de $\mathfrak q$ y una potencia del cual es divisible por $\mathfrak q$. Se tiene la representación como mínimo común múltiplo:
\[
 \mathfrak a=[\mathfrak q_1,\ldots,\mathfrak q_r];\qquad
 \mathfrak p_1^{\sigma_1}\cdots \mathfrak p_r^{\sigma_r}\equiv 0\pmod{\mathfrak a}.
\]
Aquí los $\mathfrak q$ son máximas componentes primarias; es decir, cada $\mathfrak q$ es primario, pero el mínimo común múltiplo de cualesquiera dos $\mathfrak q$ no es primario. Además, puede suponerse sin restricción de generalidad que ningún $\mathfrak q$ divide al mínimo común múltiplo de los restantes, de modo que ningún $\mathfrak q$ puede simplemente omitirse; los $\mathfrak p$ son los ideales primos asociados de los $\mathfrak q$. Ahora bien, no son los $\mathfrak q$, sino --y en esto reside el análogo del caso de una variable-- sus ideales primos asociados $\mathfrak p$ los que quedan determinados unívocamente por $\mathfrak a$; y como cada $\mathfrak p$ es divisor de $\mathfrak a$, y recíprocamente un producto de potencias de los $\mathfrak p$ resulta divisible por $\mathfrak a$, los ceros del ideal coinciden por tanto con los de sus ideales primos asociados.

II. \emph{La construcción del cuerpo de ceros de un ideal primo $\mathfrak p$.} Por definición, las clases residuales módulo un ideal primo forman un anillo sin divisores de cero, que contiene al menos un elemento distinto del elemento cero tan pronto como $\mathfrak p$ es distinto del ideal unidad; formando cocientes, este anillo puede por tanto ampliarse al cuerpo de clases residuales $(\mathfrak R)$ de $\mathfrak p$. Siempre puede construirse un cuerpo de extensión $\mathfrak R$ de $P$, isomorfo a $(\mathfrak R)$, que se convierte en el cuerpo de ceros de $\mathfrak p$; es decir, $\mathfrak R=P(\alpha_1,\ldots,\alpha_n)$, donde $\alpha_1,\ldots,\alpha_n$ designa un cero de $\mathfrak p$. Aquí $\mathfrak R$ tiene grado de trascendencia $k$ $(0\leq k<n)$ respecto de $P$; contiene por tanto $k$ elementos algebraicamente independientes respecto de $P$, mientras que cualesquiera $(k+1)$ elementos se vuelven algebraicamente dependientes respecto de $P$. Recíprocamente, todo cuerpo de ceros de grado de trascendencia $k$ situado en un cuerpo de extensión de $P$ es isomorfo al cuerpo de clases residuales $(\mathfrak R)$; cuerpo de ceros de grado de trascendencia $k$ y cuerpo de clases residuales son, por tanto, abstractamente idénticos.

III. \emph{Descomposición de la norma de $\mathfrak p$ en factores lineales en el cuerpo de Galois determinado por $\mathfrak p$.} Si las indeterminadas $x_1,\ldots,x_n$ se consideran generadas a partir de indeterminadas originarias $y_1,\ldots,y_n$ mediante una transformación lineal cuyos coeficientes son indeterminadas, entonces el cuerpo de ceros $\mathfrak R$ de grado de trascendencia $k$ puede considerarse como un cuerpo de extensión algebraico finito de $P(x_{i+1},\ldots,x_n)$ $(k=n-i)$. Se puede pasar así al cuerpo de Galois $\mathfrak R$ respecto de $P(x_{i+1},\ldots,x_n)$, en el cual la función prima $P^{(i)}(x_i,x_{i+1},\ldots,x_n)$ con el cero $\alpha_i$ se descompone en factores lineales. Esta función prima $P^{(i)}$, cuando se vuelve a las indeterminadas originarias $y$, desempeña el papel de la ecuación fundamental; pues $\alpha_i$ se vuelve igual a la forma fundamental $u_1\beta_1+\cdots+u_n\beta_n$, donde $\beta$ designa un sistema de ceros en las $y$ que depende algebraicamente de los parámetros $x_{i+1},\ldots,x_n$; la descomposición en factores lineales da, por tanto, el producto sobre todos los ceros. Pero la función prima $P^{(i)}$, o una potencia de ella, también puede considerarse como la norma $N(\mathfrak p)$ de $\mathfrak p$, considerando $\mathfrak p$ como módulo de formas lineales en los productos de potencias de $x_1,\ldots,x_{i-1}$, tomando $P(x_{i+1},\ldots,x_n)$ en lugar de $P$ como dominio de coeficientes y definiendo la norma como producto sobre todos los divisores elementales, de los cuales sólo finitamente muchos son distintos de la unidad. Si el dominio de coeficientes $P$ es un cuerpo perfecto, se obtiene siempre $N(\mathfrak p)=P^{(i)}$; si el cuerpo es imperfecto, puede aparecer como norma una potencia de $P^{(i)}$.

IV. \emph{La norma del ideal inicial y el significado de la multiplicidad.} Si, correspondientemente, se considera el ideal $\mathfrak a$ sucesivamente como módulo de formas lineales en los productos de potencias de $x_1,\ldots,x_{i-1}$ $(i=1,2,\ldots,n)$, entonces puede demostrarse que, al tomar $P(x_{i+1},\ldots,x_n)$ como base, este módulo posee en cada caso sólo finitamente muchos divisores elementales distintos de la unidad; ésta es una condición de finitud que se sigue de la existencia de la base ideal. El producto sobre estos finitamente muchos divisores elementales da las normas individuales; la norma $N(\mathfrak a)$ se define como el producto de todas las normas individuales. Los máximos factores primarios $Q^{(i)}(x_i,x_{i+1},\ldots,x_n)$ de $N(\mathfrak a)$ corresponden biunívocamente a los ideales primos asociados $\mathfrak p$ de $\mathfrak a$, de tal manera que $Q^{(i)}$ se convierte en una potencia de $P^{(i)}$, donde $P^{(i)}$ (o, si es necesario, una potencia de $P^{(i)}$) designa la norma $N(\mathfrak p)$ de $\mathfrak p$. Así se obtiene una descomposición
\[
 N(\mathfrak a)=N(\mathfrak p_1)^{\varrho_1}\cdots N(\mathfrak p_r)^{\varrho_r}=Q_1\cdots Q_r,
\]
--donde, si es necesario, $N(\mathfrak p)$ debe sustituirse por el divisor elemental más alto $E(\mathfrak p)=P^{(i)}$-- y con ello, por I, II y III, queda efectuada la descomposición de la norma en factores lineales para todos los ceros de $\mathfrak a$. Como multiplicidad de los ceros correspondientes al ideal primo individual $\mathfrak p$ de $\mathfrak a$ puede considerarse el grado de $Q^{(i)}$, que se interpreta como el número de clases residuales linealmente independientes sobre $P(x_{i+1},\ldots,x_n)$. Más precisamente, se trata de las clases residuales de ciertas componentes aisladas de la descomposición, determinadas unívocamente por $\mathfrak p$ y por los ideales primos de dimensión superior, a saber: las clases residuales del ideal fundamental de etapa $(i-1)$ módulo la componente del ideal fundamental de etapa $i$ determinada por $\mathfrak p$; esta última componente posee a $\mathfrak p$ y a todos los ideales primos de dimensión superior como ideales primos asociados, mientras que el primero posee sólo los ideales primos de dimensión superior.

Una \emph{inserción completa del caso de una variable} se obtiene al pasar también aquí al ideal principal. La descomposición dada en 1. se escinde entonces --según que $a(x)$ se considere como polinomio básico o como norma-- en dos descomposiciones separadas: en la representación del ideal principal como mínimo común múltiplo de máximas componentes primarias según I, que aquí equivale a una representación como producto de potencias de ideales primos; o en la descomposición de la norma como producto de potencias de las normas de los ideales primos individuales según IV. En la construcción del cuerpo de ceros según 2., lo que en realidad está en cuestión es el cuerpo de ceros del ideal primo definido por $p(x)$, mientras que en 3. se descompone la norma en factores lineales. La definición de la multiplicidad según 4. se subordina a la dada en IV, pues el ideal fundamental de etapa cero es igual al ideal unidad, mientras que para una variable el ideal fundamental de primera etapa ya es igual al propio ideal.

La \emph{inserción en la teoría usual de la eliminación} está dada por el hecho de que la norma del ideal tiene significado de resultante. En esta interpretación, el viejo problema de la teoría de la eliminación se convierte, por tanto, en una parte de la teoría de ideales del dominio polinómico, concretamente en aquella parte que, junto a los teoremas de descomposición de los ideales, trata también la descomposición de sus normas y la conexión entre ambas descomposiciones.

\begin{center}
(Recibido el 19 de octubre de 1923.)
\end{center}

\section*{26. Construcción abstracta de la teoría de ideales en el cuerpo algebraico de números}

\begin{center}
J. Ber. d. DMV 33 (1924), p. 102
\end{center}

4. E. Noether, Gotinga: Construcción abstracta de la teoría de ideales en el cuerpo algebraico de números.

Se indicaron las propiedades abstractas que son completamente equivalentes a la teoría de ideales en el cuerpo algebraico de números y que, por tanto, caracterizan todos los anillos que poseen una teoría de ideales de ese tipo. Son los anillos con elemento unidad y sin divisores de cero en los que vale el doble teorema de cadenas: toda cadena de divisores formada por ideales termina en lo finito, e igualmente toda cadena de múltiplos módulo cualquier ideal distinto del ideal cero; y que son algebraicamente íntegramente cerrados respecto de su cuerpo de cocientes. Si se suprime esta última hipótesis, se obtienen teoremas sobre la teoría de ideales en un orden finito, que en parte ya se encuentran sin demostración en Dedekind (3.ª ed.).

Una exposición detallada aparecerá en los \emph{Mathematische Annalen}.

\section*{27. Números hilbertianos en la teoría de ideales}

\begin{center}
J. Ber. d. DMV 34 (1925), p. 101
\end{center}

E. Noether: Números hilbertianos en la teoría de ideales. Por números hilbertianos se entienden en general los números que se refieren al recuento de clases residuales de los ideales. Hilbert mismo dio, en su ``función característica'', una función que, para los ideales del dominio polinómico, da en cada caso, a partir de cierta cota, el número exacto de clases residuales linealmente independientes. Para los grados menores, Macaulay complementó esta función mediante una función generatriz que Ostrowski pudo evaluar explícitamente. Pero Macaulay indicó también números de otro tipo, que resultaron ser los más esenciales para la teoría general de ideales, porque pueden formularse abstractamente. De acuerdo con los teoremas generales de descomposición, basta la definición para ideales primarios $\mathfrak q$. Se trata del número de clases residuales linealmente independientes respecto del cuerpo de clases residuales del ideal primo asociado $\mathfrak p$. Estos números se dividen en subsistemas dados respectivamente por el número de componentes irreducibles de $\mathfrak q,\mathfrak q/\mathfrak p,\mathfrak q/\mathfrak p^2,\ldots$. El número del $i$-ésimo subsistema se caracteriza también por el número de términos de una serie de composición que va de $\mathfrak q/\mathfrak p^i$ a $\mathfrak q/\mathfrak p^{i-1}$. La serie de composición total y sus números hilbertianos forman el equivalente de la representación de los ideales primarios como potencias de ideales primos, representación que en el caso general --por ejemplo, en los órdenes finitos de un cuerpo algebraico de números-- ya no es válida.

\section*{28. Caracteres de grupos y teoría de ideales}

\begin{center}
J. Ber. d. DMV 34 (1925), p. 144
\end{center}

3. E. Noether, Gotinga: Caracteres de grupos y teoría de ideales.

La teoría frobeniana de los caracteres de grupos --esto es, de la representación de grupos finitos-- se concibe como teoría de ideales de un anillo completamente reducible, el anillo de grupo. Por ello se desarrollan en general los teoremas de isomorfía para la descomposición aditiva de un anillo completamente reducible: la representación como suma directa de ideales simples directamente indescomponibles es única hasta isomorfía, y los ideales unilaterales indescomponibles que pertenecen al mismo ideal bilateral son isomorfos entre sí. Si, por tanto, se reúnen los ideales isomorfos en una clase, hay tantas clases distintas de ideales indescomponibles como ideales bilateralmente indescomponibles.

En la especialización a sistemas completamente reducibles de cantidades hipercomplejas --en particular al anillo de grupo--, las clases de ideales indescomponibles y las clases de representaciones irreducibles se corresponden mutuamente de manera unívoca. Además se corresponden los ideales bilateralmente indescomponibles y los caracteres; y el número de componentes aditivas indescomponibles de un ideal bilateral es igual al rango de los ideales indescomponibles. Con ello queda clasificada la teoría frobeniana.

Una exposición detallada aparecerá en los Math. Ann.

\section*{29. El teorema de finitud de los invariantes de grupos lineales finitos de característica $p$}

\begin{center}
Nachr. v. d. Ges. d. Wiss. zu Göttingen 1926, pp. 28--35
\end{center}

Presentado por R. Courant en la sesión del 30 de julio de 1926.

Muestro en lo que sigue que el conocido teorema de finitud de los invariantes de grupos finitos de sustituciones lineales\footnote{Para una demostración elemental de finitud, véase E. Noether, Der Endlichkeitssatz der Invarianten endlicher Gruppen. Math. Ann. 77 (1916), pp. 89--92.} sigue siendo válido cuando, como dominio de coeficientes, en lugar de un cuerpo numérico ordinario se toma un cuerpo arbitrario, en particular uno de característica $p$.\footnote{Para los conceptos de teoría de cuerpos, véase E. Steinitz, Algebraische Theorie der Körper. J. f. M. 137 (1910), pp. 167--309. Característica en § 4, cuerpo de raíces en § 12, cuerpo de cocientes en § 3.} Pero las demostraciones conocidas de finitud fallan tan pronto como el orden del grupo es divisible por $p$; hay que recurrir a hechos aritméticos más profundos, a saber, al siguiente criterio, conocido hasta ahora sólo para característica cero:

\textbf{Criterio de finitud}\footnote{Compárese E. Noether, Algebraische und Differentialinvarianten. Jahresber. d. d. Math.-Ver. 32 (1923), pp. 177--184, n.º 3, teorema 4, y la bibliografía allí indicada. En la última frase del criterio allí dice erróneamente ``base de módulo respecto de $\mathfrak R$'' (sistema fundamental de cantidades enteras) en lugar de ``respecto de cierto subanillo de $\mathfrak R$''.}: Un dominio íntegro $\mathfrak S$ de polinomios --más generalmente, de funciones racionales-- con coeficientes en un cuerpo es finito si y sólo si en $\mathfrak S$ está contenido un subanillo finito $\mathfrak R$ tal que $\mathfrak S$ depende algebraicamente y enteramente de $\mathfrak R$. Toda base íntegra de $\mathfrak R$ se completa entonces, mediante una base de módulo de $\mathfrak S$ respecto de cierto subanillo de $\mathfrak R$, hasta una base íntegra de $\mathfrak S$.

Este criterio, que tiene interés independiente por sí mismo, se basa en la existencia de la base racional y en el hecho de que el teorema de cadena de divisores se conserva al pasar a las cantidades enteras de una extensión finita. Este segundo teorema fue demostrado en la nota precedente de Artin y van der Waerden para característica arbitraria, sólo bajo cierta hipótesis de finitud sobre el dominio inicial, hipótesis que se cumple en el caso de los grupos finitos (el anillo de raíces representa una extensión finita). Para la existencia de la base racional doy una demostración válida para cuerpos arbitrarios como dominio de coeficientes; con ello puedo demostrar el criterio de finitud bajo la restricción anterior.

Que, para invariantes de grupos finitos, se cumpla el criterio de finitud descansa en el principio de las funciones simétricas. Mientras que en característica cero los coeficientes de la resolvente de Galois ya proporcionan la base íntegra, en el caso general sólo generan el subanillo finito exigido en el criterio de finitud. De la existencia de este subanillo finito se deduce, por los mismos razonamientos, la finitud de los invariantes ``relativos'' y de los invariantes ``modulares''.

Entre los invariantes considerados aquí caen como casos especiales los ``invariantes proyectivos mod $p$'' tratados por Dickson y su escuela.\footnote{Compárese Dickson, On Invariants and the theory of Numbers. The Madison Colloquium 1913. La bibliografía más reciente se encuentra en los volúmenes desde entonces publicados de las Transactions of the American Mathematical Society.} En ese caso se conocía el teorema de finitud para los invariantes modulares, pero no para los invariantes ``formales'' generales.

\subsection*{§ 1. El criterio de finitud}

1. \textbf{Teorema sobre la existencia de la base racional.} Primera formulación: todo sistema $S$ de funciones racionales de $n$ indeterminadas, con coeficientes en un cuerpo $P$, posee una base racional finita; esto es, existen finitamente muchas funciones del sistema tales que toda función del sistema puede expresarse racionalmente, con coeficientes en $P$, por medio de ellas.

Segunda formulación: sea $\mathfrak K$ un cuerpo intermedio entre $P$ y $P(x_1,\ldots,x_n)$ --donde $x_1,\ldots,x_n$ designan indeterminadas-- de grado de trascendencia $t\leq n$; sea $y_1,\ldots,y_t$ un sistema irreducible\footnote{Un sistema se llama irreducible, según Steinitz, si consta de funciones algebraicamente independientes. Un sistema se dice de grado de trascendencia $t$ si contiene al menos un sistema irreducible de $t$ elementos, pero ninguno de más elementos. Para los teoremas simples sobre sistemas irreducibles usados en lo que sigue, véase Steinitz, § 22.} de $\mathfrak K$. Entonces $\mathfrak K$ es una extensión algebraica finita de $P(y_1,\ldots,y_t)$.

Las dos formulaciones son de hecho equivalentes. Basta demostrar la existencia de la base racional para todos los cuerpos intermedios; pues el cuerpo derivado de un sistema $S$ se convierte en tal cuerpo intermedio $\mathfrak K$, cuya base racional da inmediatamente una base racional de $S$, ya que todo elemento de $\mathfrak K$ es combinación racional de finitamente muchas funciones de $S$. De la segunda formulación se sigue, por otra parte, que un sistema irreducible arbitrario $y_1,\ldots,y_t$ de $\mathfrak K$, aumentado por las finitamente muchas funciones que caracterizan a $\mathfrak K$ como extensión finita de $P(y_1,\ldots,y_t)$, da tal base racional. Recíprocamente, si según la primera formulación se da una base racional de $\mathfrak K$, entonces, por definición del grado de trascendencia, cada elemento de la base debe ser algebraicamente dependiente de $P(y_1,\ldots,y_t)$; de este modo $\mathfrak K$ queda caracterizado como extensión algebraica finita de $P(y_1,\ldots,y_t)$.

El teorema se demostrará en la segunda formulación. Supóngase primero que $\mathfrak K$ tiene el mismo grado de trascendencia $n$ que $P(x_1,\ldots,x_n)$; sea $y_1,\ldots,y_n$ un sistema irreducible de $\mathfrak K$ y, por consiguiente, de $P(x_1,\ldots,x_n)$. Por el grado de trascendencia $n$ de $y_1,\ldots,y_n$, cada $x_i$ es algebraico sobre $P(y_1,\ldots,y_n)$; por tanto $P(x_1,\ldots,x_n)$, y con él $\mathfrak K$ como cuerpo intermedio, es una extensión algebraica finita de $P(y_1,\ldots,y_n)$.

Sea ahora $t<n$; elíjase la numeración de las $x$ de tal modo que $y_1,\ldots,y_t,x_{t+1},\ldots,x_n$ formen un sistema irreducible de $P(x_1,\ldots,x_n)$. Entonces $x_{t+1},\ldots,x_n$ son irreducibles respecto de $P(y_1,\ldots,y_t)$ y, por la ley transitiva de la dependencia algebraica, también irreducibles respecto de $\mathfrak K$ y de todo cuerpo intermedio $\mathfrak M$ entre $P(y_1,\ldots,y_t)$ y $\mathfrak K$. Esto significa que $\overline{\mathfrak K}=\mathfrak K(x_{t+1},\ldots,x_n)$, respectivamente $\overline{\mathfrak M}=\mathfrak M(x_{t+1},\ldots,x_n)$, es una extensión puramente trascendente de $\mathfrak K$, respectivamente de $\mathfrak M$; por tanto todo elemento de $\overline{\mathfrak K}$ que no pertenezca a $\mathfrak K$ es trascendente respecto de $\mathfrak K$, y de igual modo todo elemento de $\overline{\mathfrak M}$ que no pertenezca a $\mathfrak M$ es trascendente respecto de $\mathfrak M$. Póngase asimismo $\overline P$ igual a la extensión puramente trascendente $P(x_{t+1},\ldots,x_n)$; entonces $y_1,\ldots,y_t$ son irreducibles respecto de $\overline P$.

La demostración puede reducirse ahora al caso ya resuelto tomando $\overline P$ como dominio de coeficientes. De $P<\mathfrak K<P(x_1,\ldots,x_t,x_{t+1},\ldots,x_n)$\footnote{$P<\mathfrak K$ significa: $P$ está contenido en $\mathfrak K$, es decir, es subcuerpo de $\mathfrak K$.} se sigue que $\overline P<\overline{\mathfrak K}<\overline P(x_1,\ldots,x_t)$, donde $\overline{\mathfrak K}$ tiene respecto de $\overline P$ el mismo grado de trascendencia $t$ que $\overline P(x_1,\ldots,x_t)$. Por tanto $\overline{\mathfrak K}$ posee, respecto de $\overline P$ como dominio de coeficientes, una base racional finita, digamos $y_1,\ldots,y_t; z_1,\ldots,z_s$. Los $z$ pueden escogerse de $\mathfrak K$, ya que $\overline{\mathfrak K}$ es el cuerpo compuesto de $\mathfrak K$ y $\overline P$. Resta mostrar que $\mathfrak K$ es igual a $\mathfrak M=P(y_1,\ldots,y_t;z_1,\ldots,z_s)$. Ahora bien, $\mathfrak M$ es un cuerpo intermedio entre $P(y_1,\ldots,y_t)$ y $\mathfrak K$; por tanto, todo elemento de $\mathfrak K$ es algebraico respecto de $\mathfrak M$. Por otra parte, $\overline{\mathfrak M}$ es igual a $\overline{\mathfrak K}$ y, por consiguiente, $\mathfrak K$ está contenido en $\overline{\mathfrak M}$. Así $\mathfrak K$ es un cuerpo intermedio entre $\mathfrak M$ y $\overline{\mathfrak M}$; como todo elemento de $\overline{\mathfrak M}$ que no pertenece a $\mathfrak M$ es trascendente respecto de $\mathfrak M$, se sigue necesariamente que $\mathfrak K$ es idéntico a $\mathfrak M$.

Corolario: si $P<\mathfrak K<\mathfrak L<P(x_1,\ldots,x_n)$ y $\mathfrak L$ es algebraica respecto de $\mathfrak K$, entonces $\mathfrak L$ es una extensión algebraica finita de $\mathfrak K$. Pues todo elemento de una base racional de $\mathfrak L$ es entonces algebraico respecto de $\mathfrak K$; por tanto $\mathfrak L$ queda caracterizado como extensión finita de $\mathfrak K$.

2. \textbf{Demostración del criterio de finitud.} La condición es evidentemente necesaria; pues si $\mathfrak S$ es un dominio íntegro finito, entonces $\mathfrak S$ mismo puede considerarse como el subanillo finito $\mathfrak R$ que aparece en el criterio.

La condición es suficiente bajo la hipótesis de que el cuerpo de raíces $P^{1/p}$ sea una extensión finita del cuerpo de coeficientes\footnote{Véase p. 28, nota 2.} cuando éste tiene característica $p$; para característica cero no se necesita ninguna hipótesis restrictiva. Para demostrarlo hay que probar que $\mathfrak S$ posee una base de módulo finita respecto de un subanillo de $\mathfrak R$.

Sea $\mathfrak K$ el cuerpo de cocientes\footnote{Véase p. 28, nota 2.} de $\mathfrak R$, y sea $\mathfrak L$ el cuerpo de cocientes de $\mathfrak S$. Estos cuerpos de cocientes existen porque se trata de anillos sin divisores de cero, y se tiene $P<\mathfrak K<\mathfrak L<P(x_1,\ldots,x_n)$. Así, por el corolario de 1., $\mathfrak L$ es una extensión algebraica finita de $\mathfrak K$. Además, por hipótesis, todo elemento de $\mathfrak S$ es entero respecto de $\mathfrak R$, y por tanto $\mathfrak S$ es subanillo del anillo $\mathfrak S'$ formado por todos los elementos de $\mathfrak L$ enteros sobre $\mathfrak R$. Pero, como no se supone nada sobre que $\mathfrak R$ esté íntegramente cerrado en $\mathfrak K$, no puede aplicarse directamente el hecho de que el teorema de cadena de divisores se conserva bajo extensión finita. Más bien, primero hay que pasar a un subanillo $\mathfrak T$ de $\mathfrak R$, igualmente finito, para el cual se cumpla esta clausura íntegra.

Supóngase primero que el dominio de coeficientes $P$ contiene infinitamente muchos elementos. Entonces puede elegirse de $\mathfrak R$ un sistema irreducible $g_1(x),\ldots,g_t(x)$ de tal manera que $\mathfrak R$, y por tanto $\mathfrak S$, sea entero sobre $\mathfrak T=P[g_1(x),\ldots,g_t(x)]$. En efecto, sean $f_1(x),\ldots,f_k(x)$ la base íntegra de $\mathfrak R$ existente por hipótesis; si ésta no forma un sistema irreducible, sea --tras una eventual renumeración de los índices-- $F(f_k;f_1,\ldots,f_{k-1})=0$ una relación satisfecha por $f_k$. Si se reemplazan $f_1,\ldots,f_{k-1}$ por
\[
 f_1+t_1f_k,\ldots,f_{k-1}+t_{k-1}f_k,
\]
entonces $f_k$, y por tanto también $f_1,\ldots,f_{k-1}$, depende enteramente de estas nuevas funciones cuando los $t_i$ se eligen como indeterminadas; como $P$ posee infinitamente muchos elementos, los $t_i$ pueden especializarse a elementos de $P$ de modo que esta dependencia entera se conserve. Tras finitamente muchos pasos, usando la ley transitiva de la dependencia entera, se llega así al anillo buscado $\mathfrak T=P[g_1(x),\ldots,g_t(x)]$. Entonces $\mathfrak L$ es un cuerpo de extensión finito del cuerpo de cocientes $P(g_1(x),\ldots,g_t(x))$ de $\mathfrak T$, y el anillo $\mathfrak S'$ de todos los elementos de $\mathfrak L$ enteros sobre $\mathfrak R$ es idéntico al anillo de todos los elementos de $\mathfrak L$ enteros sobre $\mathfrak T$.

En $\mathfrak T$ vale el teorema de cadena de divisores para el sistema de todos los ideales, y $\mathfrak T$ está íntegramente cerrado en su cuerpo de cocientes, ambas cosas porque $\mathfrak T$ es isomorfo al dominio polinómico en $t$ indeterminadas. De la hipótesis de que $P^{1/p}$ es finito respecto de $P$ se sigue además que $\mathfrak T^{1/p}$ es finito respecto de $\mathfrak T$.\footnote{Compárese Hilbert, Über die vollen Invariantensysteme, § 1. Math. Ann. 42 (1893), pp. 313--373.}\footnote{Compárese la nota de Artin y van der Waerden citada en la introducción.} Por tanto, en $\mathfrak S'$ vale el teorema de cadena de divisores para el sistema de todos los $\mathfrak T$-módulos; en particular, $\mathfrak S$, como subanillo de $\mathfrak S'$ que contiene a $\mathfrak T$, posee una base de $\mathfrak T$-módulo finita. La existencia de la base finita de $\mathfrak T$-módulo para $\mathfrak S$ sigue siendo válida sin ninguna hipótesis restrictiva sobre el dominio de coeficientes $P$ si $\mathfrak L$ es una extensión de primera especie de $P(g_1(x),\ldots,g_t(x))$; por tanto, en particular, siempre que $P$ tenga característica cero.\footnote{Compárese, por ejemplo, E. Noether, Abstrakter Aufbau der Idealtheorie in algebraischen Zahl- und Funktionenkörpern, § 3. Math. Ann. 96 (1926), pp. 26--61.}

Sean $h_1(x),\ldots,h_s(x)$ una base de $\mathfrak T$-módulo de $\mathfrak S$. Entonces la representación
\[
 f(x)=A_1(g)h_1(x)+\cdots+A_s(g)h_s(x)
\]
para un $f(x)$ arbitrario de $\mathfrak S$ --donde las $A$, como elementos de $\mathfrak T$, designan polinomios en las $g$-- muestra que $g_1(x),\ldots,g_t(x),h_1(x),\ldots,h_s(x)$ forman la base íntegra buscada.

Si el cuerpo $P$ contiene sólo finitamente muchos elementos, entonces el cuerpo $P(u)$ obtenido por adjunción de una indeterminada $u$ --es decir, de un elemento trascendente respecto de $\mathfrak S$-- contiene infinitamente muchos elementos. Así, el anillo $\mathfrak S(u)$ posee una base íntegra finita respecto de $P(u)$ como dominio de coeficientes; y como $\mathfrak S(u)$ es el anillo compuesto de $\mathfrak S$ y $P(u)$, esta base puede escogerse de $\mathfrak S$ separando según las potencias de $u$, aunque ahora los elementos de base $g_i$ de $\mathfrak S$ ya no formarán un sistema irreducible. En consecuencia, todo polinomio $f(x)$ de $\mathfrak S$ admite una representación $C(u)f(x)=G(u,g_i,h_i)$, donde $C(u)$ designa un polinomio no nulo de $P[u]$ y $G$ es entero en $u,g_i,h_i$. Comparando los coeficientes de una potencia adecuada de $u$ se ve que $g_i(x)$ y $h_i(x)$ dan una base íntegra de $\mathfrak S$. Con ello queda \textbf{demostrado el criterio de finitud en plena generalidad}.

\subsection*{§ 2. Finitud de los invariantes}

1. Por \textbf{grupo lineal finito de característica $p$} se entiende un grupo $\mathfrak H$ de $h$ sustituciones lineales $A_1,\ldots,A_h$ (de determinante no nulo) con coeficientes en un cuerpo $P$ de característica $p$. Aquí $A_k$ es la sustitución lineal
\[
 x_1^{(k)}=a_{11}^{(k)}x_1+\cdots+a_{1n}^{(k)}x_n,
 \ldots,
 x_n^{(k)}=a_{n1}^{(k)}x_1+\cdots+a_{nn}^{(k)}x_n,
\]
o, abreviadamente, $(x^{(k)})=A_k(x)$. Así, el grupo $\mathfrak H$ lleva la fila $(x)$ con elementos $x_1,\ldots,x_n$ a las filas $(x^{(k)})$ con elementos $x_1^{(k)},\ldots,x_n^{(k)}$. Como la identidad debe estar contenida entre $A_1,\ldots,A_h$, la fila $(x)$ está también contenida entre las filas $(x^{(k)})$.

Por \textbf{invariante racional entero (absoluto) del grupo} se entiende una función racional entera de $x_1,\ldots,x_n$, con coeficientes en $P$\footnote{Esto no constituye restricción de generalidad, en cuanto todo invariante con coeficientes en un cuerpo de característica $p$ --que debe poseer un cuerpo común que lo contenga junto con $P$ para que los invariantes estén definidos-- puede descomponerse linealmente en finitamente muchos invariantes con coeficientes en $P$. En el caso general basta expresar los coeficientes mediante otros linealmente independientes respecto de $P$; entonces la propiedad de invariancia se satisface idénticamente en estos coeficientes independientes.}, que permanece idénticamente inalterada al aplicar $A_1,\ldots,A_h$. Pertenecen así a estos invariantes todas las funciones simétricas, con coeficientes en $P$, de las filas $(x^{(k)})$. Recíprocamente, para todo invariante de este tipo se tiene
\[
 f(x)=f(x^{(1)})=\cdots=f(x^{(h)})
 \quad\text{o}\quad
 h\cdot f(x)=f(x^{(1)})+\cdots+f(x^{(h)}).
\]
De aquí se sigue que, en el caso en que el orden $h$ del grupo no sea divisible por la característica --en particular en característica cero--, las funciones simétricas de las filas $(x^{(k)})$, con coeficientes en $P$, agotan también la totalidad de los invariantes. Como estas funciones simétricas poseen una base íntegra finita, en este caso queda realizada la demostración de finitud; al mismo tiempo se obtiene la base íntegra como el sistema $G_{\alpha\alpha_1\ldots\alpha_n}(x)$ de los coeficientes de la resolvente de Galois:\footnote{Véase p. 28, nota 1.}
\[
 \Phi(z,u)=\prod_{k=1}^{h}\bigl(z-u_1x_1^{(k)}-\cdots-u_nx_n^{(k)}\bigr)
\]
\[
 =z^h+\sum G_{\alpha\alpha_1\ldots\alpha_n}(x)
 z^{\alpha}u_1^{\alpha_1}\cdots u_n^{\alpha_n}
 \quad(\alpha\ne h;\ \alpha+\alpha_1+\cdots+\alpha_n=h).
\]

2. En el caso en que $h$ sea divisible por $p$, no es necesario que todo invariante del grupo sea una función simétrica de las filas $(x^{(k)})$, pues ya no se puede concluir de $h\cdot f(x)$ a $f(x)$. Pero el anillo $\mathfrak R$ derivado de los finitamente muchos invariantes $G_{\alpha\alpha_1\ldots\alpha_n}(x)$ --los coeficientes de la resolvente de Galois-- forma un subanillo finito del sistema $\mathfrak S$ de todos los invariantes, y queda por mostrar que $\mathfrak S$, respecto de $\mathfrak R$, cumple las hipótesis del criterio de finitud.

Ante todo, sin restricción de generalidad,\footnote{Véase p. 33, nota.} puede tomarse como dominio de coeficientes el subcuerpo $\overline P$ de $P$ derivado de todos los coeficientes de sustitución $a_{\mu\nu}^{(k)}$. Este cuerpo, derivado de finitamente muchos elementos, se obtiene por adjunción de finitamente muchos elementos algebraicos o trascendentes al cuerpo primo; y como el cuerpo primo es perfecto, se sigue que $\overline P^{1/p}$ es finito respecto de $\overline P$.\footnote{Véase p. 32, nota 2.}

Queda por mostrar que $\mathfrak S$ depende algebraicamente y enteramente de $\mathfrak R$. Por 1., entre las $h$ filas $x^{(k)}$ aparece también la fila originaria $x_1,\ldots,x_n$; por tanto, la resolvente de Galois $\Phi(z,u)$ se anula idénticamente en $u$ para $z=u_1x_1+\cdots+u_nx_n$. Si se especializan en cada caso todos los $u$ a cero salvo $u_i$, que se pone igual a uno, se obtiene una relación
\[
 x_i^h+\sum x_i^\alpha G_{\alpha0\ldots\alpha_i\ldots0}(x)=0
 \quad (\alpha\ne h;\ \alpha+\alpha_i=h),
\]
que muestra que cada $x_i$ depende algebraicamente y enteramente de $\mathfrak R$. Pero entonces lo mismo vale para todo polinomio en $x$; en particular, $\mathfrak S$ es entero sobre $\mathfrak R$. Así se cumplen las hipótesis del criterio de finitud, y queda \textbf{demostrada la finitud de los invariantes}.

3. Debe observarse que los mismos razonamientos dan la \textbf{finitud de los invariantes relativos y modulares}. Un polinomio $f(x)$ se llama \textbf{invariante relativo} del grupo si todos los $f(x^{(k)})$ difieren de $f(x)$ sólo por un factor no nulo de $P$. Por tanto, el sistema de invariantes absolutos --en particular el anillo $\mathfrak R$ derivado de los $G_{\alpha\alpha_1\ldots\alpha_n}(x)$-- forma un subanillo finito del anillo derivado de todos los invariantes relativos; y, como arriba, se cumplen las hipótesis del criterio de finitud respecto de $\mathfrak R$.

Un polinomio $f(x)$ se llama \textbf{invariante modular} del grupo si $f(x)$ se vuelve igual a $f(x^{(k)})$ para todo sistema especial de valores $x_1=\alpha_1,\ldots,x_n=\alpha_n$ de $P$. Todo invariante absoluto es al mismo tiempo un invariante modular; sin embargo, si $P$ contiene sólo finitamente muchos elementos, no tiene por qué valer la recíproca. También aquí se cumplen las hipótesis del criterio de finitud respecto del subanillo $\mathfrak R$ derivado de los $G_{\alpha\alpha_1\ldots\alpha_n}(x)$.

\section*{30. Construcción abstracta de la teoría de ideales en cuerpos algebraicos de números y de funciones}
\noindent\emph{Math. Ann. 96 (1927), pp. 26--61.}

En lo que sigue se da una caracterización abstracta de todos aquellos anillos cuya teoría de ideales coincide con la teoría de ideales de todos los enteros de un cuerpo algebraico de números: es decir, aquellos anillos cuyos ideales pueden representarse de manera única como productos de potencias de ideales primos. Entre los ejemplos más conocidos de tales anillos se cuentan también el anillo de todos los enteros de un cuerpo algebraico de funciones de una indeterminada y, más generalmente, de varias indeterminadas tan pronto como, siguiendo a Kronecker, se restringe uno a ideales de dimensión máxima, es decir, se pasa a cierto anillo de cocientes, el dominio funcional.

Para el anillo conmutativo subyacente, los axiomas completamente equivalentes a esta teoría de ideales son los siguientes:\footnote{Para la definición de los conceptos básicos de la teoría de ideales compárese, por ejemplo, E. Noether, \emph{Idealtheorie in Ringbereichen}, Math. Ann. 83 (1921), pp. 24--66 (citado abajo como \emph{Teoría de ideales}). Por lo demás, los conceptos más esenciales se formulan también brevemente en el presente trabajo, especialmente en §§1, 4 y 5. Para completitud se incluye algún material que no se necesita para los fines inmediatos de este trabajo.}

\begin{description}
\item[I. Condición de cadena de divisores.] Toda cadena de ideales en la que cada ideal es un divisor propio del precedente termina después de finitamente muchos pasos; en otras palabras, en toda cadena de divisores de ideales hay un índice a partir del cual todos los ideales son iguales.
\item[II. Condición de cadena de múltiplos módulo todo ideal no nulo.] Toda cadena de ideales ---todos ellos divisores de un ideal fijo distinto del ideal cero--- en la que cada ideal es un múltiplo propio del precedente termina después de finitamente muchos pasos.
\item[III.] Existencia del elemento identidad para la multiplicación.
\item[IV.] Anillo sin divisores de cero.
\item[V. Clausura íntegra en el cuerpo de cocientes.] Todo elemento del cuerpo de cocientes que sea íntegro respecto del anillo pertenece al anillo.
\end{description}

A partir de estos axiomas desarrollo, paso a paso, una teoría de ideales cada vez más restringida hasta llegar a la deseada; recíprocamente, muestro que todos estos axiomas se satisfacen efectivamente allí.

De la condición de cadena de divisores I se sigue, como mostré en \emph{Teoría de ideales} §4, que un ideal puede representarse como mínimo común múltiplo de finitamente muchos ideales primarios pertenecientes a ideales primos distintos. Recuerdo aquí brevemente esta demostración (§6), tanto porque puede simplificarse utilizando el concepto de cociente ideal, como porque quiero corregir una omisión: la demostración original usa en un punto la existencia, no supuesta, del elemento identidad de la multiplicación.\footnote{P. Urysohn llamó mi atención sobre esta omisión; mi propia modificación de la demostración era algo más pesada que la que se da aquí, debida a E. Artin. Él ya había llenado antes la laguna para sí mismo y, de manera independiente de mí, había observado también la simplificación con el cociente ideal. Una simplificación ulterior, que evita introducir la ``representación reducida'', la tomo de una cuestión afín en W. Krull, \emph{Algebraische Theorie der Ringe III}, Math. Ann. 92 (1924), pp. 181--213, Teorema I.}

La hipótesis de la condición de cadena de múltiplos tiene como efecto (§7) que, en el anillo de clases residuales módulo todo ideal distinto del ideal cero, un ideal primo no puede tener un divisor propio, salvo el ideal unitario que contiene todos los elementos. En consecuencia, los componentes primarios de estos ideales quedan determinados de manera única, salvo los que pertenecen al ideal unitario. En una forma algo menos precisa este resultado ya se encuentra en Masazo Sono;\footnote{Masazo Sono, \emph{On the Reduction of Ideals}, Memoirs of the College of Science, Kyoto University, Series A, 7 (1924), pp. 191--204.} su hipótesis de existencia de una serie de composición en el anillo de clases residuales es equivalente a la ``doble condición de cadena'' en este anillo, como muestro brevemente al final (§10). Por doble condición de cadena entiendo la validez de la condición de cadena de divisores y de la condición de cadena de múltiplos sin restricción a ideales distintos del ideal cero.

Si se satisface además el axioma III ---existencia del elemento identidad---, el ideal unitario no aparece como ideal primo asociado. Por tanto los componentes primarios están determinados de manera única; son coprimos dos a dos, y su mínimo común múltiplo es igual a su producto. Los axiomas I, II, III valen para todos los órdenes finitos de un cuerpo algebraico de números; ya Dedekind (Dirichlet--Dedekind, \emph{Zahlentheorie}, 3.ª ed., suplemento 11, §172), sin dar la demostración con todo detalle, enunció aquí la representación única de un ideal como producto de ideales de una sola especie (componentes primarios) coprimos dos a dos.

De la hipótesis IV ---ausencia de divisores de cero--- se sigue que las distintas potencias de un ideal distinto tanto del ideal cero como del ideal unitario son todas distintas, y que existe el cuerpo de cocientes. Si finalmente se satisface también el axioma V de clausura íntegra en el cuerpo de cocientes, los ideales primarios ---determinados unívocamente en virtud de IV--- se vuelven potencias de ideales primos, y se obtiene el teorema usual de factorización (§8). Esta última demostración descansa en una consecuencia de la clausura íntegra que Dedekind ya extrajo (3.ª ed., §172) para los cuerpos de números: un elemento genuinamente fraccionario puede representarse como cociente de enteros de tal manera que ni siquiera el cuadrado del numerador sea divisible por el denominador. En exposiciones posteriores esta consecuencia es sustituida ---también como consecuencia de la clausura íntegra--- por el teorema gaussiano generalizado, mucho menos transparente, o por los correspondientes teoremas de módulos que dicen algo más; todos ellos son teoremas cuya aplicación presupone elementos de base, mientras que aquí se trabaja directamente con los ideales mismos.

En la dirección recíproca (§9), los axiomas distintos del axioma V se siguen casi directamente; V se sigue de la demostración de que un elemento genuinamente fraccionario siempre puede representarse de modo que ninguna potencia del numerador sea divisible por el denominador, lo cual es por tanto un reforzamiento de la consecuencia originalmente derivada de V. Esta demostración sólo tiene éxito después de extraer de la representación ideal las conclusiones habituales: la representación por ideales principales módulo un ideal fijo y la teoría de los ideales fraccionarios, es decir, de los cocientes ideales en el cuerpo. El hecho de que la clausura íntegra se sigue de la representación por potencias de ideales primos era ya conocido por Dedekind, como muestra una observación de la 3.ª edición, §172, p. 522 abajo.

Los axiomas I a V implican que las cuatro descomposiciones que en general son distintas ---en ideales mínimos coprimos dos a dos, en ideales mutuamente primos, en máximos componentes primarios y en ideales irreducibles--- coinciden; recíprocamente, de esta coincidencia y de los axiomas I--IV se sigue la clausura íntegra. Para anillos con divisores de cero, el cumplimiento de los axiomas restantes ---con V definido como clausura íntegra en el anillo de cocientes--- no implica necesariamente que las descomposiciones coincidan, como muestra el ejemplo del anillo de clases residuales módulo ciertos ideales de un orden finito (§9).

En §1 esbozo la teoría de las cantidades íntegras respecto de un anillo, hasta la consecuencia de Dedekind derivada de la clausura íntegra. En §2 muestro cómo las condiciones de cadena se trasladan del anillo a módulos formados por formas lineales, con este anillo como dominio de coeficientes y de multiplicadores.\footnote{Esta forma de los teoremas de módulos se debe a Prüfer (\emph{Neue Begründung der algebraischen Zahlentheorie}, §2, Math. Ann. 94 (1924), pp. 198--243) y es más transparente que el traslado usual de los teoremas de base.} De ello extraigo (§3) la conclusión de que, al pasar a los elementos íntegros de un cuerpo de extensión algebraica de primera especie, los axiomas I a IV siguen siendo válidos para todo orden, mientras que para el orden formado por todos los elementos íntegros vale también V. Este hecho ---por supuesto familiar para el cuerpo algebraico de números--- permite subordinar los cuerpos de funciones mencionados al comienzo; en particular, mediante la demostración sencilla de que el dominio funcional derivado de polinomios íntegros en varias indeterminadas satisface los axiomas, la teoría de ideales de Kronecker queda plenamente subsumida.\footnote{La diferencia aparece sólo en los teoremas de discriminantes, porque el anillo de clases residuales del dominio funcional módulo un número primo se vuelve un cuerpo imperfecto, y por ello pueden aparecer extensiones de segunda especie.}

La teoría de ideales desarrollada después no presupone estos §§2 y 3, que sirven sólo para la clasificación. En §4 y §5 se dan los fundamentos independientes de los axiomas I a V; así se pone de relieve con más nitidez que en la fundamentación original qué partes de la teoría son independientes de las hipótesis de finitud.

\subsection*{§1. Teoría de las cantidades íntegras}

Sea dado un anillo conmutativo\footnote{Como es sabido, un anillo queda definido una vez dada una relación de igualdad en un conjunto y, además, dos operaciones, suma y multiplicación, que satisfacen las leyes usuales: la suma es asociativa, conmutativa y reversiblemente unívoca; la multiplicación es asociativa ---conmutativa cuando el anillo es conmutativo--- y distributiva respecto de la suma. Si la definición subyacente de igualdad no es la identidad conjuntista, entonces, para que la misma definición de igualdad se conserve en todo subsistema, tal subsistema (subanillo, módulo, ideal) debe contener, junto con cualquier elemento, todos los elementos iguales a él. Del mismo modo, en toda extensión debe suponerse que la nueva definición de igualdad incluye la antigua. Todos los conceptos como ``finitamente muchos elementos'', ``asignación única'', etc. se entienden en el sentido de la definición de igualdad; por ejemplo, ``hay uno y sólo un elemento'' significa que hay un solo elemento salvo igualdad. Por lo demás, siempre se puede pasar de la igualdad a la identidad conjuntista reuniendo los elementos iguales en una clase e introduciendo estas clases como nuevos elementos. Las observaciones dadas aquí sobre la definición de igualdad se deben a R. Hölzer.} $T$ sin divisores de cero, con elemento identidad $e$ para la multiplicación (axiomas III y IV); en $T$ se distingue un subanillo fijo $R$ que contiene el elemento identidad. Los conceptos módulo, orden, íntegro se referirán siempre a este anillo fijo $R$, que más adelante se omitirá de la notación por brevedad.\footnote{La teoría de las cantidades íntegras sigue siendo válida si se admiten divisores de cero, siempre que en su lugar se suponga la condición de cadena de divisores en $R$, lo que por §2 también hace posible la demostración del lema. Como esta forma de la teoría de cantidades íntegras no se usa después, sólo la indico en notas. Todas las definiciones consideradas siguen siendo válidas en presencia de divisores de cero; la mayoría de ellas requiere, como es sabido y fácil de ver, hipótesis aún más débiles. Lo que no se conserva es la formulación de Dedekind: ``Un número es íntegro cuando tiene una envoltura''. Pues, con divisores de cero presentes, el cociente de dos módulos finitos no tiene por qué seguir siendo finito; por ello el orden se define aquí directamente y no como cociente de módulos.} Los elementos de $R$ se denotarán por letras latinas, los elementos de $T$ en general por letras griegas.

1. Un sistema $M$ de elementos de $T$ se llama, como de costumbre, un $R$-módulo ---brevemente, un módulo--- si, junto con dos elementos $\alpha$ y $\beta$, contiene también su diferencia, y junto con $\alpha$ contiene también $r\alpha$, donde $r$ es un elemento arbitrario de $R$. Se dice que $M$ es divisible por $N$,
\[
        M\equiv 0 \pmod N,
\]
y $M$ se llama múltiplo, $N$ divisor, si todo elemento de $M$ está también contenido en $N$. Los $R$-módulos todos cuyos elementos pertenecen a $R$ se llaman ideales de $R$.

Claramente, la intersección ---el mínimo común múltiplo $[\ldots,M_i,\ldots]$--- de cualquier número de módulos de $T$ vuelve a ser un módulo; asimismo $T$ mismo es un módulo. Por tanto, el módulo derivado de un sistema arbitrario $\Sigma$ de elementos de $T$ existe de manera única como la intersección de todos los módulos que contienen $\Sigma$. La suma ---el máximo común divisor $(\ldots,M_i,\ldots)$--- de un sistema arbitrario de módulos es el módulo derivado de su unión. El producto $MN$ de dos módulos se define como el módulo derivado de los elementos $\mu\nu$, donde $\mu$ recorre todos los elementos de $M$ y $\nu$ todos los elementos de $N$. El producto satisface por tanto las leyes asociativa y conmutativa, ya que éstas valen para los elementos del anillo. La ley distributiva en $T$ da además, para los módulos, la ley distributiva
\[
        (\ldots,M_i,\ldots)N=(\ldots,M_iN,\ldots),
\]
pues precisamente por esta ley en $T$ ambos lados son el módulo derivado de todos los elementos $\mu_i\nu$.

Como $R$ mismo es un módulo, la condición de que $R$ sea el dominio de multiplicadores puede expresarse por $RM\equiv0\pmod M$. Como por hipótesis $R$ contiene el elemento identidad, se tiene también $M\equiv0\pmod{RM}$, y por consiguiente $M=RM$.

Un módulo $M$ se llama finito si existe al menos un sistema $\Sigma$ de finitamente muchos elementos del cual se deriva $M$; los elementos $\mu_1,\ldots,\mu_n$ de $\Sigma$ se llaman una base de módulo de
\[
        M=(\mu_1,\ldots,\mu_n).
\]
La suma y el producto de dos módulos finitos, y por tanto de finitamente muchos, son de nuevo finitos, puesto que se derivan respectivamente de la unión y de los productos de los elementos de base. La última afirmación se sigue de la representación $r_1\mu_1+\cdots+r_n\mu_n$ de todos los elementos $\mu$ de $M$ por la base y de la ley conmutativa de la multiplicación en $T$.

2. Un anillo de extensión de $R$ situado en $T$, es decir, uno que contiene todos los elementos de $R$, se llama un $R$-orden, brevemente un orden. La intersección de arbitrariamente muchos órdenes en $T$ vuelve a ser un orden, y también $T$ es un orden; por tanto existe de manera única el orden derivado de cualquier sistema $\Sigma$ de elementos de $T$. Sumas y productos de órdenes se definen correspondientemente como para módulos; en particular $O^2=O$ para todo orden.

Todo orden es al mismo tiempo un $R$-módulo; un orden se llama finito cuando es un módulo finito. Puesto que sumas y productos surgen formando sumas y productos de módulos, las sumas y productos de órdenes finitos son de nuevo órdenes finitos por 1.; además vale la siguiente ley transitiva: si $A$ es un $R$-orden finito y $B$ es un $A$-orden finito, entonces $B$ es también un $R$-orden finito. En efecto, $B$ es al mismo tiempo un $R$-orden, luego un $R$-módulo. Si $\alpha_1,\ldots,\alpha_m$ forman una base de módulo de $A$ respecto de $R$, y $\beta_1,\ldots,\beta_n$ una de $B$ respecto de $A$, entonces los productos $\alpha_i\beta_j$ forman manifiestamente una base de módulo de $B$ respecto de $R$.

3. Un elemento $\alpha$ de $T$ se llama íntegro respecto de $R$, brevemente íntegro, si el orden $R_\alpha$ derivado de $\alpha$ es finito; equivalentemente, si la cadena de módulos
\[
        U_n=(\alpha^0,\alpha,\ldots,\alpha^{n-1})
\]
termina después de finitamente muchos pasos, $U_n=U_{n+1}=\cdots$.

La definición admite también una segunda formulación, que se usará en 4.: un elemento $\alpha$ de $T$ se llama íntegro si existe al menos un módulo principal $C$ distinto del módulo cero ---es decir, $C$ se deriva de un elemento $\gamma\ne0$--- tal que el producto del orden $R_\alpha$ con $C$ sea un módulo finito; en otras palabras, tal que la cadena de módulos $CU_n$ termine después de finitamente muchos pasos.\footnote{Si se admiten divisores de cero en $R$, entonces debe exigirse la existencia de un módulo principal regular; es decir, debe suponerse que $\gamma$ sea un no-divisor de cero, un elemento regular, que en anillos sin divisores de cero significa simplemente $\gamma\ne0$.}

Finalmente, la definición es también idéntica a la formulación usual: un elemento $\alpha$ de $T$ se llama íntegro si satisface al menos una ecuación
\[
        \alpha^n+r_1\alpha^{n-1}+\cdots+r_n=0,
\]
donde los $r$ son elementos de $R$.

Estas definiciones son en efecto idénticas. Puesto que $R_\alpha$ coincide con el módulo derivado de todas las potencias de $\alpha$, cuando $R_\alpha$ es finito siempre se puede escoger una base de módulo entre estas potencias. Si tal base viene dada, por ejemplo, por $e=\alpha^0,\alpha,\ldots,\alpha^{n-1}$, esto significa que la cadena de los $U_n$ termina en $U_n$; recíprocamente, $U_n=U_{n+1}=\cdots$ da esta base. Pero la terminación de la cadena de los $U_n$ es también idéntica a la existencia de la ecuación
\[
        \alpha^n+r_1\alpha^{n-1}+\cdots+r_n=0.
\]
Si $R_\alpha$ es finito, entonces $CR_\alpha$ es finito, de modo que la cadena $CU_n$ termina; recíprocamente, de la finitud de $CR_\alpha$, y por tanto de la terminación de los $CU_n$, se sigue que también los $U_n$ terminan y por consiguiente que $R_\alpha$ es finito. En efecto, como $C$ se supone un módulo principal distinto del módulo cero, $CU_n=CU_{n+1}$ implica también $U_n=U_{n+1}$.

La propiedad de anillo de las cantidades íntegras y la ley transitiva de la dependencia íntegra descansan en el siguiente

\medskip
\noindent\textbf{Lema.} Sea $O$ un orden finito en $T$, y sea $\alpha$ un elemento cualquiera de $O$. Entonces el orden $R_\alpha$ derivado de $\alpha$ es finito; en otras palabras, todos los elementos de un orden finito son íntegros.

La demostración sigue los argumentos usuales. Sea $\omega_1,\ldots,\omega_n$ una base de módulo de $O$. Como $O$ es un anillo, $\alpha\omega_i$ pertenece a $O$. Por tanto
\[
        \alpha\omega_i=r_{i1}\omega_1+\cdots+r_{in}\omega_n,
\]
y de aquí
\[
        \det(r_{ij}-e_{ij}\alpha)=0,
\]
donde $e_{ij}=0$ para $i\ne j$ y $e_{ii}=e$; esto prueba que $\alpha$ es íntegro. Para la anulación del determinante se usa la hipótesis de que $T$, y por tanto $O$, no tiene divisores de cero.\footnote{Si se supone la condición de cadena de divisores en $R$ pero se admiten divisores de cero, entonces el lema es una consecuencia inmediata del teorema de módulos de §2, según el cual la condición de cadena se transfiere a los módulos de $O$. Esta demostración se debe también, para el cuerpo de números, a Dedekind (4.ª ed., §173, II) y es la primera aplicación de condiciones de cadena en la literatura. En las consecuencias que se darán bajo 4., Dedekind usa en cambio el hecho más especial de que el número de clases residuales módulo un ideal en un cuerpo algebraico de números es finito.}

El sistema $G$ de todas las cantidades íntegras en $T$ forma un orden. Contiene a $R$, pues los elementos de $R$ son íntegros, ya que $R$ es un $R$-módulo finito con el elemento identidad como base. Pero $G$ tiene también la propiedad de anillo: el orden $R_{\alpha,\beta}$ derivado de dos cantidades íntegras arbitrarias $\alpha,\beta$ es igual al producto $R_\alpha R_\beta$ y por ello es finito. Por el lema, $\alpha+\beta$ y $\alpha\beta$ son por tanto íntegros como elementos de un orden finito.

\medskip
\noindent\textbf{Ley transitiva de la dependencia íntegra.} Si $O$ es un orden formado por cantidades íntegras, entonces todo elemento de $T$ íntegro respecto de $O$ es también íntegro respecto de $R$. Por definición, $\alpha$ satisface una ecuación
\[
        \alpha^m+\omega_1\alpha^{m-1}+\cdots+\omega_m=0,
\]
y por tanto es también íntegro respecto del orden $O'=R_{\omega_1,\ldots,\omega_m}$ derivado de $\omega_1,\ldots,\omega_m$, que es finito como producto $R_{\omega_1}\cdots R_{\omega_m}$. Por la ley transitiva dada en 2., el $O'$-orden finito $O'_\alpha$ derivado de $\alpha$ es entonces también un $R$-orden finito, y $\alpha$ es íntegro por el lema como elemento de un orden finito.

4. El anillo $R$ se llama íntegramente cerrado en $T$ si todo elemento de $T$ íntegro respecto de $R$ pertenece a $R$. Según la segunda formulación de la definición de cantidades íntegras en 3., un elemento $\alpha$ de $T$ pertenece a $R$ si y sólo si existe al menos un módulo principal $C$ distinto del módulo cero tal que $CR_\alpha$ sea un módulo finito.\footnote{Si se admiten divisores de cero en $R$, $C$ debe suponerse de nuevo un módulo principal regular; asimismo el elemento $c$ en la Consecuencia I debe suponerse regular.}

Del concepto de anillo íntegramente cerrado Dedekind (3.ª ed., §172) extrajo dos consecuencias importantes, que hacen posible usar este concepto para la teoría de ideales:

\medskip
\noindent\textbf{Consecuencia de Dedekind I.} Sea $R$ íntegramente cerrado en $T$, y supóngase además que en $R$ vale la condición de cadena de divisores para ideales (axioma I). Si $\beta$ es un elemento no íntegro de $T$, y $c\ne0$ un elemento de $R$, entonces existe un exponente $\rho>0$ tal que en la sucesión de elementos
\[
        c,\;c\beta,\ldots,c\beta^\nu,\ldots
\]
los elementos $c,\ldots,c\beta^\rho$ son íntegros, y todos los restantes son no íntegros.

Si todos los elementos $c\beta^\nu$ fueran íntegros, entonces por la clausura íntegra de $R$ pertenecerían a $R$. Entonces la cadena de módulos $C,CB_1,\ldots,CB_\nu,\ldots$ ---donde $C$ denota el módulo derivado de $c$, y $B_\nu$ el derivado de $1,\beta,\ldots,\beta^\nu$--- se convertiría en una cadena de ideales $C_\nu=(c,c\beta,\ldots,c\beta^\nu)$, que termina por hipótesis. Pero entonces, contra la suposición, $R_\beta$ sería finito y $\beta$ íntegro; por tanto, algunos de los elementos $c\beta^\nu$ son no íntegros. Además, junto con cualquier elemento no íntegro, todos los siguientes son no íntegros: pues si $c\beta^r$ es íntegro, y por tanto pertenece a $R$, entonces para $\rho<r$ el elemento $c\beta^\rho$ satisface
\[
        (c\beta^\rho)^r-c^{\,r-\rho}(c\beta^r)^\rho=0
\]
y es por tanto íntegro. Así, si $c\beta^{\rho+1}$ es el primer elemento no íntegro, entonces ---como $c$ es íntegro--- $\rho>0$, y éste es el exponente requerido.

Al especializar el anillo de extensión $T$ al cuerpo de cocientes de $R$ ---es decir, al cuerpo obtenido por adjuntar todos los pares de elementos\footnote{En la formulación familiar de Steinitz, J. f. M. 137. Si se admiten divisores de cero en $R$, entonces el anillo de cocientes debe sustituir al cuerpo de cocientes, adjuntando todos los pares cocientes cuyo denominador sea un elemento regular de $R$. Aquí la Consecuencia de Dedekind II ya no es válida, pues $n$ no será en general un elemento regular.}--- se obtiene lo siguiente:

\medskip
\noindent\textbf{Consecuencia de Dedekind II.} Sea $R$ íntegramente cerrado en su cuerpo de cocientes $K$, y supóngase que en $R$ vale la condición de cadena de divisores para ideales. Si $\beta$ es un elemento no íntegro de $K$, entonces $\beta$ admite una representación como cociente de elementos de $R$,
\[
        \beta=m/n,
\]
tal que $m^2/n$ también es no íntegro.

Póngase $\beta=b/c$ con $c\ne0$. Entonces en la sucesión $c,c\beta=b,c\beta^2,\ldots$ los dos primeros elementos son ciertamente íntegros, de modo que $\rho>1$, donde $\rho$ es el exponente de la Consecuencia I. Si se pone
\[
        m=c\beta^\rho,\qquad n=c\beta^{\rho-1},
\]
entonces $m/n=\beta$, y $m^2/n=c\beta^{\rho+1}$ es no íntegro.

La ley transitiva de la clausura íntegra tiene la siguiente forma: si $R$ es un anillo cualquiera en $T$ y $G$ denota el sistema de todas las cantidades de $T$ íntegras respecto de $R$, entonces $G$ consta también de todas las cantidades de $T$ íntegras respecto de $G$. Pues toda cantidad íntegra respecto de $G$ es también íntegra respecto de $R$, por la ley transitiva en 3., y por tanto pertenece a $G$.

\subsection*{§2. Condiciones de cadena en dominios de módulos finitos}

El teorema de módulos por el cual se transfieren las condiciones de cadena aparece, en la aplicación dada aquí, sólo para módulos de un anillo de extensión. Como la demostración es la misma en el caso de un dominio de módulos general, se tomará tal dominio como punto de partida.

Un sistema $M$ de elementos $\alpha,\beta,\ldots$ se llama dominio de módulos respecto de un anillo $R$ si se dan en $M$ dos operaciones, cada una de las cuales conduce unívocamente a elementos de $M$: una combinación aditiva de los elementos y una multiplicación de los elementos por elementos de $R$; si $M$ es un grupo abeliano respecto de la suma, mientras que la multiplicación satisface la ley asociativa; y si vale la ley distributiva en ambas formas.\footnote{Compárese \emph{Teoría de ideales}, §9.}

Para un dominio de módulos, todas las definiciones de módulo dadas en §1, 1. que no se refieren a la multiplicación siguen siendo manifiestamente válidas. En particular, $M$ se llama un dominio de módulos finito si existen finitamente muchos elementos $\xi_1,\ldots,\xi_k$ de $M$ tales que $M$ sea igual al módulo derivado de $\xi_1,\ldots,\xi_k$, y por tanto igual al sistema de todas las formas lineales $r_1\xi_1+\cdots+r_k\xi_k$ cuando $R$ tiene elemento identidad; en caso contrario aparecen términos adicionales $n_i\xi_i$.

\medskip
\noindent\textbf{Teorema de módulos.} Sea $M$ un dominio de módulos finito respecto de un anillo conmutativo $R$ con elemento identidad, y supóngase que en $R$ vale la condición de cadena de divisores, respectivamente de múltiplos, para ideales. Entonces en $M$ vale la condición de cadena de divisores, respectivamente de múltiplos, para módulos.

La demostración descansa en asignar unívocamente a cada módulo de $M$ un sistema de finitamente muchos ideales de $R$, de tal manera que a un divisor propio, respectivamente a un múltiplo propio, del módulo corresponda cada vez al menos un ideal divisor propio, respectivamente múltiplo propio.

Se dice que un elemento de $M$ tiene longitud $i$ si puede escribirse al menos de una manera como forma lineal en $\xi_1,\ldots,\xi_i$, mientras que no tiene representación como forma lineal en $\xi_1,\ldots,\xi_{i-1}$. A un módulo arbitrario $W$ de $M$ se le asignan ahora $k$ ideales $\mathfrak a_1,\ldots,\mathfrak a_k$ de $R$: sea $W_i$ el módulo de todos los elementos de $W$ de longitud $\le i$, y sea $\mathfrak a_i$ el sistema de todos los coeficientes de $\xi_i$ en $W_i$. Primero se obtiene:

Si $B$ es un divisor propio de $W$, y si $\mathfrak b_1,\ldots,\mathfrak b_k$ son los ideales asignados a $B$, entonces entre los $\mathfrak b_i$ hay al menos un divisor propio del correspondiente $\mathfrak a_i$. De $B\equiv0\pmod W$ se sigue que $B_i\equiv0\pmod {W_i}$, y por tanto $\mathfrak b_i\equiv0\pmod{\mathfrak a_i}$ para $i=1,2,\ldots,k$. Por hipótesis debe haber elementos en $B$ que no estén contenidos en $W$, y por tanto tales elementos de longitud mínima $\ell$. Entonces $\mathfrak b_\ell$ es un divisor propio de $\mathfrak a_\ell$; en efecto $\mathfrak b_\ell\equiv0\pmod{\mathfrak a_\ell}$ cuando
\[
        \beta=b_1\xi_1+\cdots+b_\ell\xi_\ell
\]
es tal elemento de longitud mínima en $B$. Si $\mathfrak b_\ell=\mathfrak a_\ell$, existe un elemento
\[
        \alpha=a_\ell\xi_\ell+\cdots+a_1\xi_1\equiv0\pmod W,
\]
y por tanto un elemento $\beta-\alpha\in B$, $\beta-\alpha\notin W$, de longitud menor que $\ell$.

Sea dada ahora una cadena de divisores o de múltiplos
\[
        W^{(1)},W^{(2)},\ldots,W^{(\nu)},\ldots;
\]
supóngase que la correspondiente condición de cadena vale en $R$. Si $\mathfrak a_{1\nu},\ldots,\mathfrak a_{k\nu}$ son los ideales asignados a $W^{(\nu)}$, entonces las sucesiones
\[
        \mathfrak a_{i1},\mathfrak a_{i2},\ldots
\]
son cadenas de divisores, respectivamente de múltiplos, para $i=1,\ldots,k$. Por hipótesis existe por tanto un índice $\nu_0$ tal que el ideal $\mathfrak a_{i\nu_0}$ es igual a todos los ideales siguientes para cada $i$. Entonces, por lo que se acaba de demostrar, el módulo $W^{(\nu_0)}$ es igual a todos los módulos siguientes; así se transfieren las condiciones de cadena.

De aquí se obtiene inmediatamente lo siguiente:

\medskip
\noindent\textbf{Consecuencia del teorema de módulos.} Si en $R$ vale la condición de cadena de múltiplos módulo todo ideal distinto del ideal cero, entonces en $M$ vale la condición de cadena de múltiplos módulo todo módulo $C$ cuyos ideales asignados $\mathfrak c_1,\ldots,\mathfrak c_k$ son todos distintos del ideal cero. Bajo estas hipótesis las cadenas de múltiplos $\mathfrak a_{i1},\ldots,\mathfrak a_{i\nu},\ldots$ terminan después de finitamente muchos pasos.

\medskip
\noindent\textbf{Observación adicional.} Si $R$ es un anillo sin elemento identidad, entonces la condición de cadena de divisores y la condición de cadena de múltiplos módulo ideales distintos del ideal cero siguen transfiriéndose. En efecto, en lugar de $M$ se considera el sistema $M'$ de todos los elementos de la forma
\[
        a_1\xi_1+\cdots+a_k\xi_k+n_1\xi_1+\cdots+n_k\xi_k,
\]
que vuelve a pasar a $M$ cuando los múltiplos enteros adicionales $n_i\xi_i$ se reemplazan por elementos de $M$. A este $M'$ se le pueden asignar, correspondientemente a lo anterior, $2k$ ideales, donde los $\mathfrak a_i$ son ideales de $R$ y los $\mathfrak n_i$ son ideales de los enteros racionales. Como para los $\mathfrak n_i$ se satisfacen tanto la condición de cadena de divisores como la condición de cadena de múltiplos módulo todo ideal no nulo, la demostración de la condición de cadena de divisores sigue siendo válida para $M'$ y por tanto para $M$; para la condición de cadena de múltiplos, sin embargo, debe exigirse que los $2k$ ideales asignados a $C$, respectivamente a $C'$, sean todos distintos del ideal cero.

\subsection*{§3. Paso a cuerpos de extensión finita}

Para la clasificación de los cuerpos de números y de funciones queda por mostrar cómo los axiomas I a V formulados en la introducción se transfieren al pasar a cuerpos de extensión finita.

1. \emph{Supóngase que en $\mR$ se satisfacen todos los axiomas I a V salvo el axioma II, la condición de cadena de múltiplos. Sea $\mK$ el cuerpo de cocientes de $\mR$, y sea $\mL$ un cuerpo de extensión finita de primera especie sobre $\mK$.\footnote{Siguiendo a Steinitz, un cuerpo de extensión algebraica se llama ``de primera especie'' si todo elemento es cero de una función prima que, en un cuerpo de extensión adecuado, se descompone en factores lineales distintos dos a dos.} Entonces en el sistema $\mS$ de todos los elementos de $\mL$ íntegros respecto de $\mR$ se satisfacen los mismos axiomas, y en todo orden de $\mS$ siguen satisfaciéndose todos estos axiomas con excepción de la clausura íntegra.}

Por §1, 3, $\mS$ forma un anillo para el cual se satisfacen los axiomas III y IV: existencia del elemento identidad y ausencia de divisores de cero. Por la conclusión de §1, 4, $\mS$ es íntegramente cerrado en $\mL$; al mismo tiempo, $\mL$ se convierte en el cuerpo de cocientes de $\mS$, pues todo elemento de $\mL$ se lleva a un elemento de $\mS$ multiplicándolo por un elemento adecuado de $\mR$. Así se satisface el axioma V.

La demostración de I sigue los argumentos conocidos. Como extensión de primera especie, $\mL$ se obtiene adjuntando a $\mK$ un solo elemento $\alpha$, que por lo anterior puede suponerse íntegro respecto de $\mR$. Junto con $\alpha$, también las cantidades conjugadas $\alpha',\alpha'',\ldots$ contenidas en el cuerpo de Galois de $\mL$ sobre $\mK$ son íntegras; por tanto también lo es su producto de diferencias, y también lo es el cuadrado $D$ de este producto de diferencias. Como $\alpha,\alpha',\ldots$ son todos distintos, $D$ es una cantidad no nula de $\mK$, y por la clausura íntegra de $\mR$ en $\mK$ es un elemento de $\mR$. Por el argumento usual, como $\mR$ es íntegramente cerrado, $\mS$ está contenido en el dominio de $\mR$-módulos finito $M$ generado por los elementos $\xi_i=\alpha^i/D$; para ello sólo hay que pasar, en la representación de un elemento de $\mS$ mediante los $\xi_i$, a los elementos conjugados y resolver el sistema resultante de ecuaciones lineales para los coeficientes de la representación. Por el teorema de módulos de §2, la condición de cadena de divisores vale para todos los $\mR$-módulos de $M$, y por tanto en particular para todos los ideales de $\mS$. La condición de cadena de divisores vale asimismo para los ideales de cualquier orden en $\mS$, pues también aquí se trata de $\mR$-módulos en $M$; además, para todo orden se satisfacen los axiomas III y IV.

2. \emph{Si en $\mR$ se satisfacen todos los axiomas I a V, lo mismo vale en $\mS$. En todo orden de $\mS$ siguen valiendo todos los axiomas salvo la clausura íntegra.}

En vista de 1, sólo queda probar el axioma II. Por el corolario del teorema de módulos de §2 basta probar lo siguiente. Si un ideal no nulo de $\mS$ se considera como un $\mR$-módulo $\mC$ en $M$, entonces los ideales $\mc_i$ de $\mR$ asociados a $\mC$ son todos distintos del ideal cero. Ahora bien, $\mC$ tiene el mismo rango lineal finito sobre $\mK$ que $M$; pues junto con cualquier elemento $\gamma$, $\mC$ contiene también $\gamma\alpha^i$. Así, para cada $\xi_i$ existe un elemento $c_i\ne0$ de $\mR$ tal que $c_i\xi_i$ pertenece a $\mC$; por la unicidad de la representación por los $\xi_i$ ---donde el índice debe recorrer sólo los valores menores que el grado del cuerpo--- $c_i$ pertenece a $\mc_i$, y por tanto $\mc_i$ no es el ideal cero. Este argumento sigue siendo válido para todos los órdenes del mismo rango que $M$; para un orden de rango menor se pasa a su cuerpo de cocientes, que como cuerpo intermedio entre $\mK$ y $\mL$ es también de primera especie, y cuyo grado sobre $\mK$ coincide con el rango lineal del orden. Por consiguiente se conserva el mismo argumento. Finalmente, debe observarse que un orden de $\mS$ distinto de $\mS$ nunca es íntegramente cerrado: como todo orden contiene también $\mR$, la clausura íntegra lo forzaría a contener $\mS$.

Para la subordinación de los cuerpos de números y de funciones basta, por tanto, verificar los axiomas I a V para los dominios básicos: los enteros racionales, los polinomios de una indeterminada y el dominio funcional de polinomios de varias indeterminadas.

3. \emph{Para el anillo $\mR$ de los enteros racionales, respectivamente de los polinomios de una indeterminada con coeficientes en un cuerpo, se satisfacen los axiomas I a V.} Aquí I y II se siguen ya sea del hecho de que, módulo cualquier número no nulo o cualquier tal polinomio no nulo de una variable, hay sólo finitamente muchas clases residuales, respectivamente finitamente muchas clases residuales linealmente independientes; ya sea del hecho de que todo ideal es principal. De esto se sigue la descomposición única de los ideales como producto de potencias de ideales primos, y por tanto un ideal no nulo es divisible sólo por finitamente muchos ideales. Los axiomas III y IV se satisfacen manifiestamente, el último como consecuencia de la definición de igualdad. La clausura íntegra V se sigue del hecho de que, por la factorización única de los elementos de $\mR$ en irreducibles salvo unidades ---divisores de uno---, todo elemento del cuerpo de cocientes no contenido en $\mR$ puede escribirse como cociente de dos elementos coprimos de $\mR$; por tanto ninguna potencia del numerador es divisible por el denominador. Tal elemento no puede entonces satisfacer ninguna ecuación que lo caracterice como íntegro respecto de $\mR$.

4. \emph{Sea ahora $\mR$ el anillo de todos los polinomios de varias indeterminadas con coeficientes en un cuerpo, o con coeficientes enteros racionales.} Entonces los axiomas III, IV, V se satisfacen como en 3, pues también aquí la factorización de los elementos en irreducibles es única salvo unidades. La condición de cadena de divisores I es una consecuencia inmediata del teorema de Hilbert sobre la existencia de una base ideal, mientras que la condición de cadena de múltiplos II no vale aquí.

Sin embargo, pasando al \emph{dominio funcional}, lo que equivale a restringirse a ideales de dimensión máxima, se puede obtener también la validez del axioma II; la validez de I se puede entonces demostrar como en 3, sin usar el teorema de Hilbert. En efecto, adjúntese a $\mR$ una indeterminada $u$ y considérese el anillo $\mR^*$ de todos los polinomios en $u$ con coeficientes en $\mR$, definiendo la igualdad coeficiente a coeficiente. Como de costumbre, llámese primitivo (respecto de $\mR$) a un polinomio de $\mR^*$ si el máximo común divisor de sus coeficientes ---divisor en el sentido de un polinomio de $\mR$, no de un divisor ideal--- es una unidad de $\mR$. Entonces todo polinomio de $\mR^*$ es producto de un elemento de $\mR$ y de un polinomio primitivo de $\mR^*$, y el producto de polinomios primitivos de $\mR^*$ vuelve a ser primitivo.

\emph{La totalidad de los elementos del cuerpo de cocientes de $\mR^*$ cuyos denominadores son polinomios primitivos de $\mR^*$ forma, por tanto, un anillo, el dominio funcional $\mF$ de $\mR$.} En este dominio funcional $\mF$, las unidades son exactamente los elementos cuyo numerador y denominador son polinomios primitivos de $\mR^*$; todo elemento de $\mF$ es por tanto representable de manera única como producto de un elemento de $\mR$ y de una unidad de $\mF$. De aquí que el teorema de factorización para los elementos de $\mR$ se transfiera a los elementos de $\mF$; para $\mF$, por tanto, además de los axiomas III y IV, también se satisface el axioma V como en 3.

Además, en $\mF$ todo ideal es principal. Pues junto con cualquier elemento de $\mF$ un ideal contiene también el elemento de $\mR$ obtenido de él por multiplicación por una unidad; y junto con dos elementos de $\mR$ contiene su máximo común divisor. En efecto, con $f(x)=t(x)\overline f(x)$ y $g(x)=t(x)\overline g(x)$ contiene también $t(x)(\overline f(x)+u\overline g(x))$, y por tanto $t(x)$, si $\overline f$ y $\overline g$ son coprimos y su combinación lineal es, en consecuencia, una unidad. Así, partiendo de un elemento arbitrario $f(x)$ del ideal, si en el ideal hay un elemento $g(x)$ que no es divisible por él, entonces el máximo común divisor $t(x)$ pertenece también al ideal y es un divisor propio de $f(x)$. La repetición finita de este proceso conduce a un elemento base del ideal, y así se reconoce que el ideal es principal. Por tanto en $\mF$ vale la descomposición única de los ideales como productos de potencias de ideales primos, correspondiente a la factorización de los elementos de $\mR$; como en 3, se sigue que se satisfacen los axiomas I y II. Los cuerpos de extensión del cuerpo de cocientes de $\mF$ que aquí entran en consideración son sólo los que pueden generarse por adjunción de un elemento de $\mS$.\footnote{El sistema de cantidades íntegras respecto de $\mF$ en estos cuerpos de extensión se obtiene también pasando de $\mS$ a un dominio funcional correspondiente, exactamente como se pasa de $\mR$ a $\mF$. Compárese J. König, \emph{Algebraische Größen}, Leipzig 1903, pp. 468/69.} Teniendo en cuenta 1 y 2, queda establecida la validez de los axiomas I a V para cuerpos de números y de funciones.

\subsection*{§4. Teoremas de isomorfía. Sumas directas}

En lo que sigue se supone solamente la propiedad de anillo, respectivamente la propiedad de módulo; no se presupone ningún axioma ulterior.

1. Si $M$ y $\overline M$ son dominios de módulos respecto de $\mR$ (§2), se dice que $M$ es homomorfo a $\overline M$ (más precisamente, homomorfo como módulo), $M\sim\overline M$, si a cada elemento de $M$ corresponde uno y sólo un elemento de $\overline M$, de modo que se agote $\overline M$;\footnote{En el sentido de la definición de igualdad vigente en $\overline M$; compárese la nota 6.} y si, bajo esta correspondencia, corresponden las diferencias y la multiplicación por el mismo elemento de $\mR$. Así, de $\beta\sim\overline\beta$ y $\gamma\sim\overline\gamma$ se sigue siempre que $(\beta-\gamma)\sim(\overline\beta-\overline\gamma)$ y $r\beta\sim r\overline\beta$.

A un $\mR$-módulo $\mB$ en $M$ le corresponde por tanto homomórficamente un $\mR$-módulo $\overline{\mB}$ en $\overline M$; y la totalidad $\mC^*$ de los elementos de $M$ que corresponden a elementos de un $\mR$-módulo $\overline{\mC}$ en $\overline M$ forma un $\mR$-módulo en $M$ unívocamente determinado por $\overline{\mC}$, el módulo $\mC^*$ asociado con $\overline{\mC}$, con $\mC^*\sim\overline{\mC}$. Si en particular $\mA$ es el módulo en $M$ asociado al elemento cero de $\overline M$, entonces $\mC^*$ es un divisor de $\mA$ y se descompone en clases de elementos de $M$ congruentes módulo $\mA$, de tal manera que estas clases corresponden biunívocamente a los elementos de $\overline{\mC}$. Si se pasa de $\mB$ en $M$ a $\overline{\mB}$ en $\overline M$ y luego de nuevo a $\mB^*$, se obtiene $\mB^*=(\mB,\mA)$, el máximo común divisor de $\mB$ y $\mA$; así, $\mB^*=\mB$ si $\mB$ es un divisor de $\mA$.

Si la correspondencia entre los elementos de $M$ y $\overline M$ es reversiblemente biunívoca, los dominios se llaman isomorfos (más precisamente, isomorfos como módulos), $M\cong\overline M$.

Si $\mA$ es un $\mR$-módulo arbitrario en $M$, surge un dominio de módulos $\overline M$ homomorfo a $M$ ---el módulo de clases residuales $M\mid\mA$--- tomando la congruencia módulo $\mA$ como nueva relación de igualdad. A cada elemento de $M$ se le asocian todos y sólo los elementos iguales a él en $\overline M$. Pasar en $\overline M$ de la definición de igualdad a la identidad significa reunir todos los elementos iguales de $\overline M$ en una clase ---una clase residual--- y tomar estas clases residuales como los nuevos elementos de $\overline M$.

\emph{Todo homomorfismo se produce por paso a un módulo de clases residuales}; pues si $M\sim\overline M$ y si $\mA$ es el módulo asociado al elemento cero de $\overline M$, entonces, como se mostró arriba, $\overline M$ es isomorfo al módulo de clases residuales $M\mid\mA$.

2. \textbf{Primer teorema de isomorfía.} \emph{Sea $\overline M$ el módulo de clases residuales $M\mid\mA$, y sea $\mC$ un divisor de $\mA$. Entonces hay un isomorfismo $\overline M\mid\overline{\mC}\cong M\mid\mC$.} En efecto, la congruencia módulo $\mA$ es al mismo tiempo congruencia módulo $\mC$; los elementos iguales módulo $\mA$ siguen siendo iguales módulo $\mC$. Por tanto se puede formar el módulo de clases residuales $M\mid\mC$ ---es decir, la igualdad módulo $\mC$--- identificando primero los elementos módulo $\mA$, pasando a $\overline M$, y reuniendo luego en $\overline M$ los elementos que son iguales módulo $\mC$; esto es lo mismo que identificar módulo $\overline{\mC}$ en $\overline M$, es decir, formar $\overline M\mid\overline{\mC}$.

\textbf{Segundo teorema de isomorfía.} \emph{Si $\mB$ y $\mA$ son módulos en $M$, entonces $(\mB,\mA)\mid\mA\cong\mB\mid[\mB,\mA]$.} Pues por 1, $\mB$ se vuelve homomorfo a $\overline{\mB}$ cuando $(\mB,\mA)\mid\mA$ se pone igual a $\overline{\mB}$; y como precisamente los elementos de $[\mB,\mA]$ corresponden al elemento cero en $\overline{\mB}$, el isomorfismo afirmado se sigue de nuevo de 1.

3. Si $\mR$ y $\overline{\mR}$ son anillos (conmutativos), se dice que $\mR$ es homomorfo a $\overline{\mR}$ (más precisamente, homomorfo como anillo), $\mR\sim\overline{\mR}$, si a cada elemento de $\mR$ corresponde uno y sólo un elemento de $\overline{\mR}$ de tal modo que se agote $\overline{\mR}$, y si bajo esta asignación corresponden las diferencias y los productos. Si la correspondencia de elementos es reversiblemente biunívoca, los anillos se llaman isomorfos (más precisamente, isomorfos como anillos), $\mR\cong\overline{\mR}$.

Todos los argumentos de 1 y 2 siguen siendo válidos cuando los módulos se sustituyen por ideales de $\mR$,\footnote{Para anillos no conmutativos deben usarse ideales bilaterales.} y cuando homomorfismo e isomorfismo de módulos se sustituyen por homomorfismo e isomorfismo de anillos. En particular, si $\mc$ es un divisor de $\ma$, el anillo de clases residuales $\mc\mid\ma$ se obtiene introduciendo la congruencia módulo $\ma$ como nueva relación de igualdad; aquí hay que observar que ahora también están definidos los productos de clases residuales.

Así, $\mc$ es homomorfo a $\mc\mid\ma$; todo homomorfismo se produce de este modo, y los teoremas de isomorfía valen como en 2:

\textbf{Primer teorema de isomorfía.} Si $\overline{\mR}$ denota el anillo de clases residuales $\mR\mid\ma$, y $\mc$ es un divisor de $\ma$, entonces $\overline{\mR}\mid\overline{\mc}\cong\mR\mid\mc$ en el sentido de isomorfismo de anillos.

\textbf{Segundo teorema de isomorfía.} Si $\mb$ y $\ma$ son ideales de $\mR$, entonces vale el isomorfismo de anillos $(\mb,\ma)\mid\ma\cong\mb\mid[\mb,\ma]$.\footnote{Estos teoremas de isomorfía, conocidos por la teoría de grupos, aparecen para módulos en una forma algo más especial por primera vez en Dedekind (compárese 3.ª ed., p. 484). La forma anterior para ideales aparece en Masazo Sono, \emph{On Congruences. I, II, III, IV}, Memoirs of the College of Science, Kyoto Imperial University, 2 (1917), 3 (1918, 1919), compárese 2, p. 215. En cada caso sólo se enuncia explícitamente el segundo teorema de isomorfía.}

4. Reúno ahora las reglas de cálculo para ideales coprimos\footnote{En la forma que se remonta a Dedekind (4.ª ed., §178, III a VIII). El cálculo con el elemento identidad que aparece en todas las exposiciones más recientes es mucho más engorroso.} de las que se siguen los teoremas sobre sumas directas en 5. En el anillo conmutativo $\mR$, supóngase la existencia de un elemento identidad multiplicativo $e$. Así $\ma=\mo\ma$ para todo ideal de $\mR$, donde $\mo$ denota el ideal unitario. Dos ideales $\ma,\mb$ se llaman coprimos si su máximo común divisor $(\ma,\mb)$ es $\mo$.

\emph{4$\alpha$. Si $(\ma,\mb)=\mo$, entonces $(\ma,\mb\mc)=(\ma,\mc)$.} Pues $(\ma,\mc)=(\ma,\mb)(\ma,\mc)=(\ma^2,\ma\mb,\ma\mc,\mb\mc)$ es divisible por $(\ma,\mb\mc)$ y recíprocamente. Por tanto:

\emph{4$\beta$. De $\mc\mb\equiv0\pmod{\ma}$ y $(\ma,\mb)=\mo$ se sigue $\mc\equiv0\pmod{\ma}$.} En efecto, $\mc\mb\equiv0\pmod{\ma}$ equivale a $(\ma,\mb\mc)=\ma$; como $(\ma,\mb)=\mo$, esto da también $(\ma,\mc)=\ma$, o sea $\mc\equiv0\pmod{\ma}$. Además:

\emph{4$\gamma$. Si cada uno de los ideales $\ma_1,\ldots,\ma_r$ es coprimo con cada uno de los ideales $\mb_1,\ldots,\mb_s$, entonces el producto de los $\ma_i$ es coprimo con el producto de los $\mb_i$.} Pues de $(\ma,\mb)=\mo$ y $(\ma,\mc)=\mo$ se obtiene $(\ma,\mb\mc)=\mo$, y la repetición finita da la afirmación.

De 4$\alpha$ y 4$\gamma$ se sigue:

\emph{4$\delta$. Si los ideales $\mb_1,\ldots,\mb_r$ son coprimos dos a dos, es decir, $(\mb_i,\mb_k)=\mo$ para $i\ne k$, entonces su mínimo común múltiplo es igual a su producto.} Sean $(\ma,\mb)=\mo$ y $\mv=[\ma,\mb]$. Entonces $\mv=\mo\mv=(\ma\mv,\mb\mv)=0(\ma,\mb)$; como $\ma\mb\equiv0\pmod{\mv}$, se sigue que $\mv=\ma\mb$. La repetición finita, usando 4$\gamma$, da la afirmación.

\emph{4$\varepsilon$. Si los ideales $\mb_1,\ldots,\mb_r$ son coprimos dos a dos y $\ma_i=[\mb_1,\ldots,\mb_{i-1},\mb_{i+1},\ldots,\mb_r]=\mb_1\cdots\mb_{i-1}\mb_{i+1}\cdots\mb_r$, entonces $(\ma_1,\ldots,\ma_{i-1},\ma_{i+1},\ldots,\ma_r)=\mb_i$, y por tanto $(\ma_1,\ldots,\ma_r)=\mo$.} En efecto, por la ley distributiva $(\ma_1,\ma_2)=\mb_3\cdots\mb_r(\mb_2,\mb_1)=\mb_3\cdots\mb_r$; luego $(\ma_1,\ma_2,\ma_3)=\mb_4\cdots\mb_r(\mb_3,\mb_1\mb_2)=\mb_4\cdots\mb_r$, y la repetición finita, tras una numeración adecuada, prueba la afirmación.

El ideal $\ma_i$ se llama el complemento de $\mb_i$ en la representación $\mm=[\mb_1,\ldots,\mb_r]$.

5. Sea de nuevo $\mR$ un anillo conmutativo con elemento identidad multiplicativo. Se dice que $\mR$ es suma directa de los ideales $\ma_1,\ldots,\ma_r$, en símbolos $\mR=\ma_1+\ma_2+\cdots+\ma_r$, si cada elemento $c$ de $\mR$ puede representarse de una y sólo una manera en la forma $c=a_1+a_2+\cdots+a_r$, con $a_i$ elemento de $\ma_i$.

\emph{Si el ideal cero es el mínimo común múltiplo de los ideales $\mb_1,\ldots,\mb_r$ coprimos dos a dos, y si $\ma_i$ es el complemento de $\mb_i$, entonces $\mR$ es la suma directa de los ideales $\ma_1,\ldots,\ma_r$.} Como $\mo=(\ma_1,\ldots,\ma_r)$, todo elemento de $\mR$ admite al menos una representación aditiva por los $\ma_i$. Como $[\ma_i,\mb_i]=(0)$, esta representación es única: de $0=a_1+a_2+\cdots+a_r=a_i+b_i$ se obtiene $a_i=0$. Por el segundo teorema de isomorfía, los $\ma_i$ son isomorfos a los anillos de clases residuales $\mR\mid\mb_i$. Valen las relaciones de ortogonalidad $\ma_i\ma_k=(0)$ para $i\ne k$, y $\ma_i^2=\ma_i$; esta última porque $\ma_i=\mo\ma_i$.

\emph{Recíprocamente, si existe una representación de $\mR$ como suma directa, $\mR=\ma_1+\ma_2+\cdots+\ma_r$, y si se pone $\mb_i=(\ma_1,\ldots,\ma_{i-1},\ma_{i+1},\ldots,\ma_r)=\ma_1+\cdots+\ma_{i-1}+\ma_{i+1}+\cdots+\ma_r$, entonces los $\mb_i$ son coprimos dos a dos y su mínimo común múltiplo es el ideal cero.} Pues de $\mo=(\ma_i,\mb_i)$ se sigue $\mo=(\mb_k,\mb_i)$ para $k\ne i$, ya que $\mb_k$ es divisor de $\ma_i$. Si además $c$ es divisible por todos los $\mb_i$, entonces
\[
 c=a_2^{(1)}+\cdots+a_r^{(1)}=a_1^{(2)}+a_3^{(2)}+\cdots+a_r^{(2)}=\cdots=a_1^{(r)}+\cdots+a_{r-1}^{(r)},
\]
donde en cada caso $a_i^{(\lambda)}\equiv0\pmod{\ma_i}$. Por la unicidad de la representación, como en cada representación una componente es cero, se sigue que $0=a_1^{(\lambda)},\ldots,0=a_r^{(\lambda)}$ para todo $\lambda$, y por tanto $c=0$.\footnote{Los teoremas dados en 5 son casos especiales de un teorema general sobre la conexión entre suma directa e intersección directa. Compárese H. Prüfer, \emph{Theorie der Abelschen Gruppen I}, Math. Zeitschr. 20 (1924), pp. 165--187, §6. Los argumentos dados allí siguen siendo válidos también para una formulación tan general del concepto de grupo que módulos e ideales caigan bajo él como casos especiales.}

\subsection*{§5. Ideales primos e ideales primarios}

Sea de nuevo $\mR$ un anillo conmutativo para el que no se supone ningún axioma ulterior, ni siquiera la existencia de un elemento identidad. Reunimos los hechos básicos sobre ideales primos e ideales primarios.\footnote{En \emph{Teoría de ideales}, §4, se supone el axioma I, la condición de cadena de divisores; por ello allí no se precisa qué propiedades de los ideales primos y primarios son independientes de este axioma.}

1. Un ideal $\mpideal$ de $\mR$ se llama \emph{ideal primo débil} si el anillo de clases residuales $\mR\mid\mpideal$ es un anillo sin divisores de cero; es decir, si de $\ma\not\equiv0\pmod{\mpideal}$ y $\mb\not\equiv0\pmod{\mpideal}$ se sigue siempre que $\ma\mb\not\equiv0\pmod{\mpideal}$. Un ideal $\mpideal^*$ de $\mR$ se llama \emph{ideal primo fuerte} si el anillo de clases residuales $\mR\mid\mpideal^*$ es un anillo sin ideales divisores de cero; es decir, si de $\ma\not\equiv0\pmod{\mpideal^*}$ y $\mb\not\equiv0\pmod{\mpideal^*}$ se sigue siempre que $\ma\mb\not\equiv0\pmod{\mpideal^*}$.

\emph{Todo ideal primo fuerte es al mismo tiempo un ideal primo débil}, como muestra la especialización de $\ma,\mb$ a ideales principales. Recíprocamente, \emph{todo ideal primo débil es también un ideal primo fuerte}; por tanto se puede hablar simplemente de ideales primos. Sea en efecto $\mpideal$ un ideal primo débil; sean $\ma\not\equiv0\pmod{\mpideal}$ y $\mb\not\equiv0\pmod{\mpideal}$; escójanse elementos $a$ y $b$ en $\ma$ y $\mb$ tales que $a\not\equiv0\pmod{\mpideal}$ y $b\not\equiv0\pmod{\mpideal}$. Entonces $ab\not\equiv0\pmod{\mpideal}$, de donde $\ma\mb\not\equiv0\pmod{\mpideal}$, y así $\mpideal$ es también un ideal primo fuerte.

2. Un ideal $\mq$ de $\mR$ se llama \emph{ideal primario débil} si en el anillo de clases residuales $\mR\mid\mq$ alguna potencia de cada divisor de cero se anula; equivalentemente, si de $a\not\equiv0\pmod{\mq}$ y $b^x\not\equiv0\pmod{\mq}$ para todo $x$ se sigue siempre que $ab\not\equiv0\pmod{\mq}$. El sistema $\mpideal$ de todos los elementos de $\mR$ que se vuelven divisores de cero en $\mR\mid\mq$ forma un ideal, y de hecho un ideal primo; es divisor de $\mq$ y se llama el ideal primo asociado. Pues junto con $a$, también $ra$ es divisor de cero; junto con $a$ y $b$ lo es también $a-b$, ya que con $a^x$ y $b^\lambda$ el elemento $(a-b)^{x+\lambda}$ es siempre divisible por $\mq$. Si además $a$ no es divisor de cero, entonces por definición ninguna potencia de $a$ es divisor de cero; de $a\not\equiv0\pmod{\mpideal}$ y $b\not\equiv0\pmod{\mpideal}$ se sigue $a^x b^\lambda=(ab)^x\not\equiv0\pmod{\mq}$, y por tanto $ab\not\equiv0\pmod{\mpideal}$.

Un ideal $\mq^*$ de $\mR$ se llama \emph{ideal primario fuerte} si en el anillo de clases residuales $\mR\mid\mq^*$ alguna potencia de cada ideal divisor de cero se anula; es decir, si de $\ma\not\equiv0\pmod{\mq^*}$ y $\mb^x\not\equiv0\pmod{\mq^*}$ para todo $x$ se sigue siempre que $\ma\mb\not\equiv0\pmod{\mq^*}$. Como para los ideales primarios débiles se muestra que el máximo común divisor de todos los ideales de $\mR$ que se vuelven ideales divisores de cero en $\mR\mid\mq^*$ forma un ideal primo, divisor de $\mq^*$, el ideal primo asociado. Recíprocamente, todos los ideales primarios con el mismo ideal primo asociado $\mpideal$ se dicen pertenecientes a $\mpideal$.

\emph{Todo ideal primario fuerte es al mismo tiempo un ideal primario débil}, como muestra la especialización de $\ma,\mb$ a ideales principales. En general, sin embargo, la recíproca es falsa.\footnote{Lo muestra el siguiente ejemplo comunicado por R. Hölzer. Sea $\mR$ el dominio de polinomios en numerablemente muchas indeterminadas $x_i$ con coeficientes en un cuerpo, y sea $\mq=(x_1^2,x_2^3,\ldots,x_\nu^{\nu+1},\ldots,x_i x_k\,[i\ne k]\ldots)$. Los divisores de cero en el anillo de clases residuales son exactamente los polinomios divisibles por $\mpideal=(x_1,x_2,\ldots,x_\nu,\ldots)$, y en cada caso una potencia de tal divisor de cero es divisible por $\mq$; por tanto $\mq$ es un ideal primario débil con ideal primo asociado $\mpideal$. Pero $\mq$ no es un ideal primario fuerte: poniendo $\ma=(x_1,x_3,x_5,\ldots,x_{2\nu+1},\ldots)$ y $\mb=(x_2,x_4,\ldots,x_{2\nu},\ldots)$, se tiene $\ma\mb\equiv0\pmod{\mq}$, pero ninguna potencia de $\ma$ ni de $\mb$ es divisible por $\mq$.} \emph{Si, no obstante, se supone en $\mR$ el axioma I, la condición de cadena de divisores, entonces todo ideal primario débil es también un ideal primario fuerte}; por tanto se puede hablar simplemente de ideales primarios. Sea $\mq$ un ideal primario débil; supóngase $\ma\not\equiv0\pmod{\mq}$ y $\mb^x\not\equiv0\pmod{\mq}$ para todo $x$. Entonces existen elementos $a$ de $\ma$ y $b$ de $\mb$ tales que $a\not\equiv0\pmod{\mq}$ y $b^x\not\equiv0\pmod{\mq}$ para todo $x$; luego $ab\not\equiv0\pmod{\mq}$ y por tanto $\ma\mb\not\equiv0\pmod{\mq}$. Pues si para todo $b$ de $\mb$ hubiera un exponente $\lambda$ tal que $b^\lambda\equiv0\pmod{\mq}$, entonces $\mb^{\lambda_1+\cdots+\lambda_r}\equiv0\pmod{\mq}$, con $\lambda_1,\ldots,\lambda_r$ los exponentes de los finitamente muchos elementos de base.\footnote{Para inferir la base ideal finita del axioma I debe usarse un buen orden fijo en $\mR$; compárese §6.} El mismo argumento muestra también que hay un exponente mínimo $\varrho$ tal que $\mpideal^\varrho\equiv0\pmod{\mq}$, donde $\mpideal$ es el ideal primo asociado; $\varrho$ se llama el exponente de $\mq$.

3. En lo que sigue la condición de cadena de divisores de nuevo \emph{no} se supone. El mínimo común múltiplo $[\ma_1,\ldots,\ma_r]$ de finitamente muchos ideales se llama una \emph{representación más corta} si ningún $\ma_i$ está contenido en el mínimo común múltiplo de los ideales restantes, es decir, si ningún $\ma_i$ puede omitirse.

\emph{El mínimo común múltiplo de finitamente muchos ideales primarios débiles, respectivamente fuertes, pertenecientes al mismo ideal primo $\mpideal$, vuelve a ser un ideal primario débil con $\mpideal$ como ideal primo asociado. Una representación más corta por finitamente muchos ideales primarios débiles, respectivamente fuertes, pertenecientes a ideales primos distintos no es un ideal primario débil, respectivamente fuerte.} Sea $\mf=[\mq_1,\ldots,\mq_r]$, y sea $\mpideal$ el ideal primo asociado de los ideales primarios débiles $\mq_i$. Entonces $\mpideal$ consta también exactamente de los elementos alguna de cuyas potencias es divisible por $\mf$. De $a\not\equiv0\pmod{\mf}$ y $b^x\not\equiv0\pmod{\mf}$ para todo $x$ se sigue que $b\not\equiv0\pmod{\mpideal}$ y $a\not\equiv0\pmod{\mq_i}$ para al menos un índice $i$; por tanto $ab\not\equiv0\pmod{\mq_i}$, y de ahí $ab\not\equiv0\pmod{\mf}$. Si en toda la argumentación se sustituyen elementos por ideales, ya no se puede inferir $\mb\not\equiv0\pmod{\mpideal}$.

Sea ahora $\mm=[\mq_1,\ldots,\mq_r]$ una representación más corta por ideales primarios débiles con $\mpideal_i\ne\mpideal_k$ para $i\ne k$, y sea $r\ge2$. Entonces hay al menos un ideal primo asociado, digamos $\mpideal_1$, que no está contenido en ninguno de los otros; pues toda cadena propia de múltiplos de los $\mpideal$ debe terminar después de a lo sumo $r$ pasos. Existen por tanto elementos $a_i$ de $\mpideal_i$ tales que $a_i^{\varrho_i}\equiv0\pmod{\mq_i}$, mientras que $a_2\not\equiv0\pmod{\mpideal_1},\ldots,a_r\not\equiv0\pmod{\mpideal_1}$. Si además $q_1\not\equiv0\pmod{\mm}$ es un elemento de $\mq_1$, entonces $q_1a_2^{\varrho_2}\cdots a_r^{\varrho_r}\equiv0\pmod{\mm}$, pero ninguna potencia de $a_2^{\varrho_2}\cdots a_r^{\varrho_r}$ es divisible por $\mm$, pues de lo contrario se seguiría la divisibilidad por $\mpideal_1$. Así $\mm$ no es un ideal primario débil. Si se sustituyen en toda la argumentación los elementos por ideales, es decir, $q_1$ por $\mq_1$ y $a_i$ por $\ma_i[\not\equiv0\pmod{\mpideal_1}$, pero $\ma_i^{\varrho_i}\equiv0\pmod{\mq_i}]$, se obtiene la demostración para ideales primarios fuertes.\footnote{Debo a B. L. van der Waerden esta modificación de mi demostración original, que hace innecesaria la condición de cadena de divisores.}

4. \emph{Cociente de módulos y cociente de ideales.} Sea $\mT$ un anillo de extensión de $\mR$, y sean $\mA,\mB,\mC,\ldots$ $\mR$-módulos en $\mT$. El cociente $\mC=\mA:\mB$ se define como un módulo que satisface las condiciones $\mC\mB\equiv0\pmod{\mA}$ y, si $\mD\mB\equiv0\pmod{\mA}$, entonces $\mD\equiv0\pmod{\mC}$; así $\mC$ es el módulo más pequeño para el cual vale $\mC\mB\equiv0\pmod{\mA}$. Esto determina el cociente unívocamente como el máximo común divisor de todos los módulos $\mD$ para los que $\mD\mB\equiv0\pmod{\mA}$; tales $\mD$ existen, pues el módulo cero es ciertamente uno de ellos.

De la definición se sigue: si $\mB$ es múltiplo de $\overline{\mB}$, entonces $\mA:\overline{\mB}$ es múltiplo de $\mA:\mB$. Si $\mA$ es múltiplo de $\overline{\mA}$, entonces $\mA:\mB$ es múltiplo de $\overline{\mA}:\mB$. Además
\[
        \mA:\mB\mC=(\mA:\mB):\mC=(\mA:\mC):\mB .
\]
Si el anillo de extensión $\mT$ coincide con $\mR$, el cociente de módulos se convierte en el cociente de ideales $\ma:\mb$; por tanto las reglas anteriores valen también aquí. Como $\ma\mb\equiv0\pmod{\ma}$, se tiene también $\ma\equiv0\pmod{\ma:\mb}$.

5. \emph{Si $\ma:\mb=\ma$, entonces $\mb$ se llama primo respecto de $\ma$.} Así, un ideal $\mb$ es primo respecto de $\ma$ si y sólo si $\md\mb\equiv0\pmod{\ma}$ implica siempre $\md\equiv0\pmod{\ma}$. De las reglas de cálculo de 4 se sigue: si $\mb$ y $\mc$ son primos respecto de $\ma$, también lo son su producto y su mínimo común múltiplo. Si $\mb$ es primo respecto de $\ma$ y $\ma$ es primo respecto de $\mb$, entonces $\ma$ y $\mb$ se llaman mutuamente primos. De §4, 4$\beta$, se sigue que los ideales coprimos son también mutuamente primos.

\setcounter{footnote}{25}

\subsection*{§6. Teoría de ideales bajo la hipótesis de la condición de cadena de divisores}

Sea como base un anillo conmutativo $\mR$ en el que se satisfaga el axioma I, la condición de cadena de divisores. En lo que sigue se supone además fijado un buen orden de todos los elementos de $\mR$. Esto determina también un buen orden de todos los ideales de $\mR$. En efecto, ambas hipótesis implican inmediatamente que todo ideal de $\mR$ posee una base ideal formada por finitamente muchos elementos. Pasando del buen orden de los elementos de $\mR$ a un buen orden cuasi-lexicográfico de todos los subconjuntos finitos de $\mR$,
\footnote{Esto debe entenderse así. Si los elementos $a_1,\ldots,a_n$ de un subconjunto finito $\mathfrak W$ se denotan de modo que el orden de los índices concuerde con el buen orden prescrito en $\mR$, entonces todo subconjunto con $n$ elementos precede a uno con $m>n$ elementos. Si dos subconjuntos tienen el mismo número de elementos, el primero precede al segundo cuando el primer elemento en que difieren sus listas ordenadas precede al elemento correspondiente de la otra lista. Así todo sistema de subconjuntos finitos tiene un primer elemento. Dado un buen orden de $\mR$, no se necesita ningún postulado adicional de elección; en particular, si $\mR$ es numerable, como en el caso de un cuerpo algebraico de números, no interviene ningún postulado de elección. El papel del postulado de elección en la teoría de ideales se mencionó, sin más detalle, en \emph{Teoría de ideales}, nota 5.}
y asignando a cada ideal, como base distinguida, la primera entre sus posibles bases en este buen orden de subconjuntos finitos, se obtiene un buen orden de los ideales junto con el de los subconjuntos finitos.

\textbf{Teorema I.} \emph{Bajo la hipótesis de la condición de cadena de divisores, todo ideal de $\mR$ admite una representación como mínimo común múltiplo de finitamente muchos ideales irreducibles, es decir, de ideales que no pueden representarse como mínimo común múltiplo de dos divisores propios.}

Para la demostración basta mostrar lo siguiente: si el Teorema I es falso para un ideal $\mm$, entonces $\mm$ posee un divisor propio para el cual el Teorema I es igualmente falso; a partir de esto puede construirse, contra la condición de cadena de divisores, una cadena de divisores que no termina en finitamente muchos pasos. En efecto, $\mm$ debe ser reducible, pues de lo contrario $\mm=[\mm]$ ya daría la representación deseada. Si $\mm=[\ma,\mb]$, el Teorema I no puede valer simultáneamente para $\ma$ y $\mb$, porque entonces se seguiría la representación correspondiente para $\mm$. Por tanto hay divisores propios de $\mm$ para los cuales el Teorema I es falso; sea $\ma_1$ el primero de ellos en el buen orden de los ideales. Construyendo del mismo modo, a partir de $\ma_1$, un divisor propio $\ma_2$, y en general a partir de $\ma_i$ un divisor propio $\ma_{i+1}$, se obtiene una cadena de divisores bien determinada que no termina, en contradicción con la hipótesis.
\footnote{El Teorema I vale manifiestamente del mismo modo para dominios de módulos, si se supone la condición de cadena de divisores para el sistema de todos los módulos. El teorema tiene carácter puramente conjuntista; es al mismo tiempo el único punto de la demostración en que se usa el buen orden de los ideales. En abstracto: sea $M$ un conjunto, y sea distinguida y bien ordenada una familia $\Sigma$ de subconjuntos. Supóngase la condición de cadena para los conjuntos $\Sigma$: toda cadena $X_1,X_2,\ldots$ en la que $X_i$ es un superconjunto propio de $X_{i-1}$ termina. Un conjunto $\Sigma$ es reducible si es la intersección de dos conjuntos $\Sigma$ que son ambos superconjuntos propios; de lo contrario es irreducible. Entonces todo conjunto $\Sigma$ es intersección de finitamente muchos conjuntos $\Sigma$ irreducibles.}

Como se supone la condición de cadena de divisores, §5, 2 permite hablar simplemente de ideales primarios. La relación entre primario e irreducible está dada por lo siguiente.

\textbf{Teorema II.} \emph{Bajo la hipótesis de la condición de cadena de divisores, todo ideal irreducible es primario; equivalentemente, todo ideal no primario es reducible.}

Al pasar al anillo de clases residuales $\mR/\mm$, la condición de cadena de divisores se conserva por homomorfía. La representación de $\mm$ como mínimo común múltiplo corresponde a la representación del ideal cero; y por el primer teorema de isomorfía (§4, 3), los divisores primarios de $\mm$ corresponden de nuevo a ideales primarios en $\mR/\mm$, y recíprocamente. Por tanto basta considerar la descomposición del ideal cero en el anillo de clases residuales. Como este anillo es de nuevo un anillo de la misma generalidad, desde el principio se puede suponer que el ideal que hay que descomponer es el ideal cero de $\mR$.

Supóngase entonces que el ideal cero de $\mR$ no es primario. Hay al menos un par de ideales $\ma,\mb$ --y $\mb$ puede suponerse principal-- tales que
\[
\ma\ne(0),\qquad \mb^x\ne(0)\quad\text{para todo }x,
\qquad\text{pero}\qquad \ma\mb=(0).
\]
Formemos la cadena de cocientes de ideales
\[
\ma,\quad \ma:\mb,\quad \ldots,\quad \ma:\mb^\nu,\quad\ldots .
\]
Esta cadena termina; sea, por ejemplo,
\[
\mt=\ma:\mb^{m-1}=\ma:\mb^m=\cdots .
\]
Así $\mt=\mt:\mb$, o sea, $\mb$ es primo respecto de $\mt$. Además, $\mt$ es un divisor de $\ma$ distinto del ideal cero, y $\mb^x$ es distinto de cero para todo exponente $x$. Por tanto, para demostrar el Teorema II basta establecer
\[
        (0)=[\mt,\mb^{m+1}].
\]
Por definición,
\[
        \mb^{m-1}\mt\equiv0\pmod{\ma},\qquad\text{y por tanto}\qquad \mb^m\mt=(0),
\]
de modo que sólo resta mostrar
\[
        \mc=[\mt,\mb^{m+1}]\equiv0\pmod{\mb^m\mt}.
\]
Esto se sigue de las hipótesis de que $\mb$ es principal y primo respecto de $\mt$. Todo elemento $c$ de $\mc$, por ser divisible por $\mb^{m+1}$, tiene una representación --donde $n$ denota un símbolo entero--
\[
        c=k b^{m+1}+n b^{m+1}=r b^m,
\]
donde $r$ vuelve a ser un elemento de $\mR$. De $c\equiv0\pmod{\mt}$ se sigue $r b^m\equiv0\pmod{\mt}$, y de ahí $r\equiv0\pmod{\mt}$, puesto que $\mt:\mb=\mt$. Por tanto $c\equiv0\pmod{\mt\mb^m}$, y queda demostrado el Teorema II.

\textbf{Teorema III.} \emph{Bajo la hipótesis de la condición de cadena de divisores, todo ideal admite una representación más corta como mínimo común múltiplo de finitamente muchos máximos componentes primarios pertenecientes a ideales primos distintos.}

Para obtener tal representación, basta sustituir, siempre que sea posible, la representación por ideales irreducibles dada por el Teorema I --por tanto, por el Teorema II, por ideales primarios-- por una representación más corta. Si se agrupan los ideales primarios pertenecientes al mismo ideal primo, §5, 3 da la representación afirmada. Los componentes primarios pueden llamarse máximos porque el mínimo común múltiplo de cualesquiera de ellos ya no es primario.
\footnote{Los ideales primos correspondientes y los componentes aislados están determinados de manera única. Compárese \emph{Teoría de ideales}, o W. Krull, ``Ein neuer Beweis für die Hauptsätze der allgemeinen Idealtheorie'', Math. Ann. 90 (1923), pp. 55--64. En §7 se da una demostración directa de unicidad para el caso especial simple que allí aparece. Los ideales primarios allí reconocidos como únicos son componentes aislados.}

\subsection*{§7. Teoría de ideales bajo la hipótesis de la doble condición de cadena}

Las simplificaciones, frente a la teoría desarrollada hasta ahora, descansan en los resultados auxiliares dados en 1 y 2.

1. \emph{Si en un anillo conmutativo sin divisores de cero se satisface la condición de cadena de múltiplos, entonces el anillo ya es un cuerpo.} Hay que mostrar que, para $a\ne0$, la ecuación $ax=b$ siempre tiene solución en el anillo. Que no puede tener más de una solución se sigue de la ausencia de divisores de cero.

Sea $\ma$ el ideal principal generado por $a$. Por hipótesis, la sucesión de potencias termina; sea $\ma^m$ igual a todas las potencias siguientes. Si $\mb$ denota el ideal principal generado por $b$, entonces
\[
        \ma^m\mb=\ma^{m+1}\mb,
        \qquad \ma^{m+1}\mb=(a^{m+1}b).
\]
En particular, el elemento $a^m b$ tiene una representación --de nuevo con $n$ como símbolo entero--
\[
        a^m b=k a^{m+1}b+n a^{m+1}b=a^{m+1}c,
\]
donde $c$ es un elemento de $\mR$. Como no hay divisores de cero, esto da $b=ac$, lo que prueba que el anillo es un cuerpo.

2. De 1 se sigue inmediatamente el resultado auxiliar: \emph{si la condición de cadena de múltiplos se satisface en un anillo conmutativo, entonces un ideal primo no tiene ningún divisor propio distinto de $\mo$.}

Del mismo modo se obtiene el complemento: \emph{si sólo se satisface el axioma II, es decir, la condición de cadena de múltiplos módulo todo ideal distinto del ideal cero, entonces todo ideal primo distinto del ideal cero no tiene ningún divisor propio distinto de $\mo$.}

En efecto, el anillo de clases residuales módulo un ideal primo $\mpideal$ es, por definición, un anillo sin divisores de cero; como la condición de cadena de múltiplos se conserva allí también por homomorfía, por 1 es un cuerpo. Por la correspondencia biunívoca entre los divisores de $\mpideal$ y los ideales en el anillo de clases residuales, $\mpideal$ no tiene divisor propio salvo $\mo$.
\footnote{No es necesario que exista un elemento identidad en $\mR$, ni por tanto en $\mo$, el ideal que consta de todos los elementos de $\mR$, aunque por 1 exista una clase identidad en el anillo de clases residuales. El ejemplo más simple de esta clase, en el caso más débil del complemento, lo proporciona el sistema de todos los enteros pares.}

\textbf{Teorema IV.} \emph{Si en un anillo conmutativo se satisface la doble condición de cadena, entonces en toda representación más corta de un ideal están determinados unívocamente los componentes primarios que no pertenecen al ideal unitario $\mo$.}

\emph{Complemento.} Bajo las hipótesis de los axiomas I y II, el Teorema IV vale para todo ideal distinto del ideal cero.

Sean
\[
        \mm=[\mq,\mq_1,\ldots,\mq_r]=[\mbar\mq,\mbar\mq_1,\ldots,\mbar\mq_s]
\]
representaciones más cortas de $\mm$ por máximos componentes primarios, con ideales primos asociados
\[
        \mo,\mpideal_1,\ldots,\mpideal_r,
        \qquad \mo,\mbar\mpideal_1,\ldots,\mbar\mpideal_s.
\]
El componente $\mq$, respectivamente $\mbar\mq$, debe omitirse cuando no aparece ningún ideal primario perteneciente a $\mo$. Como
\[
        \mo\,\mpideal_1\cdots\mpideal_r\equiv0\pmod{\mm},
\]
cada $\mpideal_i$ debe estar contenido en algún $\mbar\mpideal_j$, y por el resultado auxiliar es idéntico a él. Recíprocamente, cada $\mbar\mpideal_j$ debe coincidir con algún $\mpideal_i$; por tanto los ideales primos distintos de $\mo$ coinciden. Además,
\[
        \mq\,\mq_1\cdots\mq_r\equiv0\pmod{\mq_i}
\]
implica $\mbar\mq_i\equiv0\pmod{\mq_i}$, porque los componentes restantes no son divisibles por $\mpideal_i$ y por tanto son primos respecto de $\mq_i$. La congruencia recíproca se obtiene del mismo modo; así los componentes primarios que no pertenecen a $\mo$ son únicos. Si $\mq$ aparece efectivamente, entonces $\mbar\mq$ debe aparecer también efectivamente, porque las representaciones son más cortas; esto da también la unicidad de los ideales primos asociados. La demostración muestra asimismo que toda representación sin componente primario perteneciente a $\mo$ ya es más corta.

3. Si a la doble condición de cadena se añade la existencia del elemento identidad, el Teorema IV se agudiza en el siguiente.

\textbf{Teorema V.} \emph{En un anillo conmutativo con identidad en el que se satisface la doble condición de cadena, todo ideal puede representarse de manera única como producto de finitamente muchos ideales primarios coprimos dos a dos.}

\emph{Complemento.} Bajo los axiomas I, II, III, el Teorema V vale para todo ideal distinto del ideal cero.

Por la existencia del elemento identidad se tiene $\mo^2=\mo$; por tanto todo ideal primario perteneciente a $\mo$ es igual a $\mo$ y no puede aparecer en una representación más corta de un ideal distinto de $\mo$. Así, por el Teorema IV, los componentes primarios están determinados de manera única, y al mismo tiempo toda representación se vuelve más corta después de omitir $\mo$. Por el resultado auxiliar de 2, dos ideales primos distintos son siempre coprimos; por las reglas de cálculo de §4, 4, lo mismo vale para sus componentes primarios. El mínimo común múltiplo se convierte por tanto en su producto.

De las mismas hipótesis se sigue el resultado auxiliar: \emph{en un anillo conmutativo con identidad y la doble condición de cadena, los ideales primos que son primos respecto de un ideal $\mm$ son exactamente, aparte del ideal unitario, los que no están contenidos en $\mm$. Las nociones de primo respecto de y coprimo coinciden.}

\emph{Complemento.} Bajo los axiomas I, II, III, debe suponerse que el ideal $\mm$ es distinto del ideal cero.

Si $\mpideal$ no está contenido en $\mm$, entonces es distinto de todos los ideales primos pertenecientes a $\mm$, y por el resultado auxiliar de 2 no es divisible por ninguno de ellos. Por tanto es primo respecto de cada componente primario y, por consiguiente, respecto de $\mm$; al mismo tiempo es coprimo con cada componente y por tanto con $\mm$.

Si, en cambio, $\mpideal\ne\mo$ está contenido en $\mm$, entonces debe coincidir con uno de los ideales primos asociados. Supóngase que pertenece al componente $\mq=\mq_1$. Entonces $\mq:\mpideal$ es un divisor propio de $\mq$, pues
\[
        \mpideal^{\varrho-1}\equiv0\pmod{\mq:\mpideal},\qquad
        \mpideal^{\varrho-1}\not\equiv0\pmod{\mq}.
\]
Al mismo tiempo, $\mq:\mpideal$, como divisor de $\mpideal^{\varrho-1}$, sólo es divisible por el ideal primo $\mpideal\ne\mo$ y por tanto es primario. Por la unicidad de la descomposición, el producto
\[
        (\mq:\mpideal)\mq_2\cdots\mq_r,
\]
que es divisible por $\mm:\mpideal$, es un divisor propio de $\mm$; por consiguiente $\mpideal$ no es primo respecto de $\mm$.

Así, un ideal $\mt$ es primo respecto de $\mm$ si y sólo si ninguno de los ideales primos pertenecientes a $\mt$ está contenido en $\mm$; pero en tal caso $\mt$ también es coprimo con $\mm$. Recíprocamente, los ideales coprimos son siempre mutuamente primos (§5, 5), de modo que bajo estas hipótesis ambos conceptos coinciden.

\subsection*{§8. Teoría de ideales bajo clausura íntegra en el cuerpo de cocientes}

La teoría de ideales desarrollada hasta aquí debe ahora agudizarse hasta la usual añadiendo los dos últimos axiomas.

1. \emph{Resultado auxiliar.} Sea $\mR$ un anillo conmutativo sin divisores de cero y con elemento identidad. Si se supone en $\mR$ la condición de cadena de divisores (axiomas I, III, IV), entonces de $\ma=\ma\mb$ y $\ma\ne(0)$ se sigue $\mb=\mo$. Por tanto
\[
        \mc\ne\mc\mt
\]
para todos los ideales distintos del ideal cero y del ideal unitario.
\footnote{Por otra parte, bajo las mismas hipótesis no se puede inferir $\mb=\mc$ de $\ma\mb=\ma\mc$. Por ejemplo, sea $\mR$ el dominio de polinomios en $x,y$ sobre un cuerpo, con $\ma=(xy)$, $\mb=(x^3,xy,y^2)$, $\mc=(x^2,y^2)$. Entonces $\ma\mb=\ma\mc=\ma^2$, pero $\mb\ne\mc$.}
La demostración es la misma que para el resultado auxiliar §1, 3, pues $\ma\mb$ puede considerarse como módulo finito respecto de $\mb$. Sea $\alpha_1,\ldots,\alpha_n$ una base ideal de $\ma$, existente por I. De $\ma=\ma\mb$ se obtiene el sistema
\[
        \alpha_i=b_{i1}\alpha_1+\cdots+b_{in}\alpha_n,
        \qquad b_{ij}\equiv0\pmod{\mb},\quad i=1,\ldots,n .
\]
Como $\mR$ no tiene divisores de cero, el determinante $|b_{ij}-e_{ij}|$ se anula; aquí $e_{ij}=0$ para $i\ne j$ y $e_{ii}=e$. De
\[
        e-b_{11}-\cdots\equiv0\pmod{\mb}
\]
se sigue $e\equiv0\pmod{\mb}$ y por tanto $\mb=\mo$.

2. Sólo resta mostrar que, al añadir la clausura íntegra en el cuerpo de cocientes (axioma V) a los axiomas I a IV, todo ideal primario distinto del ideal cero es una potencia de su ideal primo asociado. La unicidad de la representación resultante
\[
        \mm=\mpideal_1^{\varrho_1}\cdots\mpideal_r^{\varrho_r}
\]
para todo ideal distinto del ideal cero ya fue demostrada por el Teorema V y el resultado auxiliar de 1; al mismo tiempo se mostró que estos ideales primos no poseen ningún divisor propio distinto del ideal unitario. La demostración se dará primero para el exponente dos, recurriendo a la consecuencia II de Dedekind (§1, 4), y luego en general por inducción completa.

\emph{Resultado auxiliar.} \emph{Bajo los axiomas I a V no hay ideales primarios de exponente dos distintos de $\mq=\mpideal^2$.}

Hay que mostrar que de
\[
        \mpideal^2\equiv0\pmod{\mq},
        \qquad \mq\equiv0\pmod{\mpideal},
        \qquad \mq\not\equiv0\pmod{\mpideal^2}
\]
se sigue necesariamente $\mq=\mpideal^2$. Escójase
\[
        \mc\equiv0\pmod{\mq},
        \qquad \mc\not\equiv0\pmod{\mpideal},
        \qquad \mc\equiv0\pmod{\mpideal^2}.
\]
Del resultado auxiliar §7, 3 se sigue que $\mc:\mpideal$ es un divisor propio de $\mc$; por tanto existe un elemento $b$ tal que
\[
        b\mpideal\equiv0\pmod{\mc}\equiv0\pmod{\mq},
        \qquad b\not\equiv0\pmod{\mc}.
\]
Así $y=b/c$ pertenece al cuerpo de cocientes pero no a $\mR$; no es íntegro. Debemos probar $b\not\equiv0\pmod{\mpideal}$; entonces, como $\mq$ es primario y pertenece a $\mpideal$, se sigue que $\mpideal\equiv0\pmod{\mq}$.

Por la consecuencia II de Dedekind existen elementos $m,n$ de $\mR$ tales que
\[
        y=b/c=m/n,
        \qquad m^2/n \text{ no es íntegro};
\]
es decir,
\[
        bn=mc,
        \qquad m\not\equiv0\pmod{\mn},
        \qquad m^2\not\equiv0\pmod{\mn}.
\]
Multiplicando $b\mpideal$ por $m$ se obtiene
\[
        mb\mpideal\equiv0\pmod{\mn c},
        \qquad\text{y por tanto}\qquad m\mpideal\equiv0\pmod{\mn}.
\]
Por consiguiente $m\not\equiv0\pmod{\mpideal}$, porque $m^2\not\equiv0\pmod{\mn}$ y $\mn\not\equiv0\pmod{\mpideal}$; esto último se sigue nuevamente de §7, 3, puesto que $\mn:\mpideal$ es un divisor propio de $\mn$. Las cuatro relaciones
\[
        m\not\equiv0\pmod{\mpideal},\qquad
        n\not\equiv0\pmod{\mpideal},\qquad
        c\equiv0\pmod{\mpideal},\qquad
        c\not\equiv0\pmod{\mpideal^2},\qquad
        bn=mc
\]
implican ahora $b\not\equiv0\pmod{\mpideal}$. En efecto, como $\mpideal^2$ es primario, se obtiene primero $mc\not\equiv0\pmod{\mpideal^2}$, y luego de $bn=mc$ se sigue $b\not\equiv0\pmod{\mpideal}$. Esto prueba el resultado auxiliar.

3. \emph{Resultado auxiliar.} \emph{Bajo los axiomas I a V no hay ideales primarios de exponente $\varrho$ distintos de $\mq=\mpideal^\varrho$.}

Esto se sigue del resultado auxiliar de 2 sin nuevo uso de los axiomas.
\footnote{La demostración del resultado auxiliar 3 bajo la hipótesis de validez del resultado auxiliar 2 se encuentra en Masazo Sono, \emph{On Congruences II}, §§9 y 10. Compárese la nota 17.}
Primero se prueba:
\[
        \text{si }\mc\equiv0\pmod{\mpideal},\ \mc\not\equiv0\pmod{\mpideal^2},
        \quad\text{entonces}\quad
        \mpideal^\varrho=(\mc^\varrho,\mpideal^{\varrho+1})
\]
para todo $\varrho$. Por el resultado de 2, $\mpideal=(\mc,\mpideal^2)$. Suponiendo $\mpideal^{\varrho-1}=(\mc^{\varrho-1},\mpideal^\varrho)$, la multiplicación por $\mpideal$ y la ley distributiva dan
\[
        \mpideal^\varrho=(\mc^\varrho,\mpideal^{\varrho+1}).
\]
El resultado deseado puede escribirse en la siguiente segunda forma: si
\[
        \mq\equiv0\pmod{\mpideal^\varrho},
        \qquad \mq\not\equiv0\pmod{\mpideal^{\varrho+1}},
        \qquad \mpideal^{\varrho+1}\equiv0\pmod{\mq},
        \quad \varrho\ge1,
\]
entonces ya $\mpideal^\varrho\equiv0\pmod{\mq}$. En efecto, para todo ideal primario hay tal exponente $\varrho>1$; la condición $\mq\not\equiv0\pmod{\mpideal^{\varrho+1}}$ se sigue de $\mpideal^\varrho\equiv0\pmod{\mq}$ y de $\mpideal^\varrho\ne\mpideal^{\varrho+1}$ por el resultado de 1. La aplicación repetida de la segunda forma da $\mpideal^\varrho\equiv0\pmod{\mq}$ y por tanto $\mq=\mpideal^\varrho$.

A partir de la hipótesis de la segunda forma y de la representación de $\mpideal^\varrho$, cada elemento $q$ de $\mq$ puede escribirse
\[
        q=b c^\varrho+p^{\varrho+1},
        \qquad b\equiv0\pmod{\mpideal},
        \quad\text{o}\quad b c^\varrho\equiv0\pmod{\mq,\mpideal^{\varrho+1}}.
\]
Como $(\mq,\mpideal^{\varrho+1})$ es primario, se sigue
\[
        c^\varrho\equiv0\pmod{\mq,\mpideal^{\varrho+1}},
        \qquad\text{o}\qquad
        c^{\varrho+1}\equiv0\pmod{\mq,\mpideal^{\varrho+2}},
\]
y por tanto $\mpideal^\varrho\equiv0\pmod{\mq}$.

4. Resumiendo, se obtiene:

\textbf{Teorema VI.} \emph{Si en un anillo conmutativo se satisfacen los axiomas I a V, entonces todo ideal distinto del ideal cero y del ideal unitario puede representarse de manera única como producto de potencias de finitamente muchos ideales primos distintos del ideal cero y del ideal unitario. Estos ideales primos no poseen ningún divisor propio distinto del ideal unitario.}

\subsection*{§9. Los axiomas como consecuencia de la descomposición supuesta}

Recíprocamente, para poder inferir los axiomas a partir de la existencia de la descomposición ideal usual --que por el Teorema VI se sigue de los axiomas I a V--, la descomposición debe suponerse en la siguiente forma.

\emph{Hipótesis.} En el anillo conmutativo $\mR$, todo ideal no generado por el elemento cero ni por el elemento unitario puede representarse de manera única como producto de potencias de ideales primos. Estos ideales primos son ideales simples no generados por el elemento unitario, es decir, no tienen ningún divisor propio distinto de $\mo$; recíprocamente, todos los ideales simples son ideales primos. La unicidad debe valer en la forma aguda de que
\[
        \ma^\varrho\ne\ma^{\varrho+1}
\]
para todo ideal no generado por el elemento cero ni por el elemento unitario.

Aquí no se supone la existencia de un elemento identidad. Si no existe elemento identidad, la condición ``no generado por el elemento unitario'' no es una restricción.

1. \emph{Demostración del axioma III, existencia del elemento identidad.} Supóngase que no se satisface el axioma III. Hay que mostrar que $\mo^2$ se vuelve un ideal simple sin ser primo, contra la hipótesis. Por la hipótesis aguda de unicidad, $\mo\ne\mo^2$; de ahí $\mo\equiv0\pmod{\mo^2}$, pero $\mo^2\not\equiv0\pmod{\mo^2}$, y por tanto $\mo^2$ no es primo. Es, sin embargo, simple: no puede ser divisible por ningún ideal primo distinto de $\mo$, y por la representación producto supuesta no tiene divisores salvo $\mo$ y $\mo^2$.

2. \emph{Demostración del axioma IV, ausencia de divisores de cero.} Si $a$ es un divisor de cero, digamos $ab=0$ con $a\ne0$ y $b\ne0$, entonces la multiplicación de las representaciones producto de los ideales principales generados por $a$ y por $b$ da
\[
        (0)=\mpideal_1^{\varrho_1}\cdots\mpideal_r^{\varrho_r},
\]
donde los $\mpideal_i$ son distintos del ideal cero y del ideal unitario; esto último porque ya se ha probado la existencia del elemento identidad. Tómese la representación como más corta, en el sentido de que suprimir cualquier factor da un divisor propio del ideal cero; esto no tiene por qué ser automático, pues no se ha supuesto unicidad alguna para el ideal cero. Si aparece sólo un ideal primo, entonces $(0)=\mpideal^\varrho=\mpideal^{\varrho+1}$, contra la exigencia aguda de unicidad. Si aparece más de un ideal primo, entonces
\[
\begin{aligned}
\mpideal_1^{\varrho_1}
 &= (\mpideal_1^{\varrho_1},\mpideal_2^{\varrho_2}\cdots\mpideal_r^{\varrho_r})\mpideal_1^{\varrho_1}  \\
 &= \mpideal_1^{\varrho_1}\mpideal_2^{\varrho_2}\cdots\mpideal_r^{\varrho_r},
\end{aligned}
\]
porque, por la existencia del elemento identidad, $\mpideal_1$ es coprimo con todos los demás ideales primos. Esto contradice de nuevo la condición aguda de unicidad.

3. \emph{Demostración de los axiomas I y II, las condiciones de cadena.} Las hipótesis dan, para todo ideal distinto del ideal cero: la divisibilidad se refleja en la representación producto. Así, de $\mm\equiv0\pmod{\ma}$ y $\ma\ne\mo$, con
\[
        \mm=\mpideal_1^{\varrho_1}\cdots\mpideal_r^{\varrho_r},
        \qquad
        \ma=\mbar\mpideal_1^{\bar\varrho_1}\cdots\mbar\mpideal_s^{\bar\varrho_s},
\]
se sigue que cada $\mbar\mpideal_j$ es idéntico a algún $\mpideal_i$ y que $0\le\bar\varrho_i\le\varrho_i$. Por tanto sólo hay finitamente muchos divisores de $\mm$, correspondientes a las finitamente muchas combinaciones de exponentes. En consecuencia, la doble condición de cadena vale en el anillo de clases residuales módulo $\mm$. Quedan probados los axiomas I y II: la condición de cadena de divisores vale también en $\mR$ mismo, ya que todo divisor propio del ideal cero es distinto del ideal cero.

La demostración del axioma V, la clausura íntegra en el cuerpo de cocientes, usa las consecuencias estándar de la descomposición de ideales, que por tanto deben derivarse primero.

4. \emph{El anillo de clases residuales módulo todo ideal distinto del ideal cero es un anillo de ideales principales.} Puede suponerse que el ideal es distinto del ideal unitario, pues para el anillo de clases residuales que consta sólo del elemento cero la afirmación es inmediata.

Si el ideal es primero primario, $\mq=\mpideal^\varrho$, y si $\mc\equiv0\pmod{\mpideal}$ pero $\mc\not\equiv0\pmod{\mpideal^2}$, entonces por 3 se tiene $\mc=\mpideal\ma$ con $(\ma,\mpideal)=\mo$. En consecuencia
\[
        \mpideal^\sigma=(\mc^\sigma,\mpideal^{\sigma+1})
        \qquad(\sigma<\varrho),
\]
de modo que todo ideal en el anillo de clases residuales módulo un ideal primario es principal. En general, si
\[
        \mm=\mq_1\cdots\mq_r
\]
es la representación y
\[
        \mR/\mm=\ma_1+\cdots+\ma_r
\]
la representación correspondiente del anillo de clases residuales como suma directa (§4, 5), entonces $\ma_i$ es isomorfo a $\mR/\mq_i$; por tanto todo ideal de $\ma_i$ es principal. Si $\mc$ es un ideal cualquiera del anillo de clases residuales, entonces $\ma_i\mc$ es un ideal principal $\ma_i c_i$, y
\[
        \mc=\ma_1c_1+\cdots+\ma_r c_r=\mo c,
        \qquad c=c_1+\cdots+c_r.
\]
En efecto, $\ma_i\mo c_i$ y $\ma_i\mc$ coinciden en $\ma_i$, porque $\ma_i\ma_k=(0)$ para $i\ne k$ y $c_i\equiv0\pmod{\ma_i}$.

El teorema sobre anillos residuales de ideales principales admite también la siguiente forma. Si $\mc$ es un ideal cualquiera, puede transformarse en un ideal principal multiplicándolo por un ideal primo respecto de un ideal prescrito $\mb$. Pues, poniendo $\mm=\mb\mc$, el ideal $\mc$ se vuelve principal módulo $\mm$; es decir,
\[
        \mc=(\mo c,\mm)=(\mc\ma,\mc\mb)=\mc(\ma,\mb),
\]
de modo que $\mo c=\mc\ma$ y $(\ma,\mb)=\mo$.

5. \emph{Teoría de los ideales fraccionarios.} Un ideal fraccionario significa un $\mR$-módulo finito en el cuerpo de cocientes $\mK$.

\emph{Los $\mR$-módulos finitos no nulos en $\mK$ forman un grupo abeliano respecto de la multiplicación.}

El producto de dos $\mR$-módulos finitos es de nuevo finito; la multiplicación es asociativa y conmutativa; como $\mK=\mo\mK$, el ideal unitario es el elemento identidad. Sólo resta mostrar que la ecuación
\[
        \mathfrak U X=\mo
\]
siempre tiene solución en el sistema de los $\mR$-módulos finitos.

Observación preliminar: si la ecuación $\mathfrak A T=\mathfrak B$ tiene una y sólo una solución, entonces $T$ es el cociente de módulos $\mathfrak B:\mathfrak A$ (§5, 4). Pues
\[
        T\equiv0\pmod{\mathfrak B:\mathfrak A},
        \qquad
        \mathfrak B=\mathfrak A T\equiv0\pmod{\mathfrak A(\mathfrak B:\mathfrak A)}\equiv0\pmod{\mathfrak B}.
\]

Si primero $\mathfrak A$ es principal, $\mathfrak A=(\alpha)$, entonces $X=(e/\alpha)$ es una solución de $\mathfrak A X=\mo$, y $X$ es de nuevo un $\mR$-módulo finito. Se sigue que todo $\mR$-módulo finito es el cociente de módulos de dos ideales de $\mR$, lo que justifica el término ideal fraccionario. En efecto, si $t_1,\ldots,t_n$ es una base de módulo de $T$, pongamos $t_i=\tau_i/a$ y sea $\mt$ el ideal generado por los $\tau_i$ en $\mR$. Entonces $\mo a\,T=\mt$, y por la observación preliminar $T=(\mt:\mo a)$, tomado el cociente en $\mK$.

La solución de $\mathfrak A X=\mo$ puede darse ahora en general. Póngase $\mathfrak A=\ma:\mo c$, y escójase $\mb$ de modo que $\ma\mb$ sea el ideal principal $\mo a$. Entonces
\[
        \mathfrak A\,\mo c=\ma,
        \qquad \mathfrak A\,\mb\mo c=\mo a,
        \qquad \mathfrak A(\mb\mo c:\mo a)=\mo.
\]
Aquí $X$, como cociente de un ideal por un ideal principal, vuelve a ser un $\mR$-módulo finito. La propiedad de grupo queda demostrada. También se sigue que el cociente de módulos de dos ideales cualesquiera, $\ma:\mb$, es un $\mR$-módulo finito, como solución de $\mb X=\ma$.

6. De la propiedad de grupo se siguen otras consecuencias.

6\,$\alpha$. \emph{La cancelación es posible: de $\mathfrak T=\mathfrak A\mathfrak M:\mathfrak B\mathfrak M$ se sigue $\mathfrak T=\mathfrak A:\mathfrak B$, y recíprocamente.} Pues $\mathfrak T\mathfrak B\mathfrak M=\mathfrak A\mathfrak M$ da $\mathfrak T\mathfrak B=\mathfrak A$, y recíprocamente.

6\,$\beta$. \emph{Todo ideal fraccionario admite una representación $\mathfrak C=\ma:\mb$, donde $\ma$ y $\mb$ son coprimos; esta condición determina $\ma$ y $\mb$ de manera única.} La existencia se sigue de 5 y de 6\,$\alpha$. De $\mathfrak C=\ma:\mb=\mbar\ma:\mbar\mb$ se sigue $\ma\mbar\mb=\mb\mbar\ma$, y de ahí la unicidad cuando se supone que ambos pares $\ma,\mb$ y $\mbar\ma,\mbar\mb$ son coprimos.

6\,$\gamma$. \emph{Todo módulo principal generado por un elemento no íntegro, es decir, por un elemento de $\mK$ que no pertenece a $\mR$, admite una representación como cociente de ideales principales tal que ninguna potencia del numerador sea divisible por el denominador.} Sea $\mo n=\mb:\mc$ la representación reducida, de modo que $(\mb,\mc)=\mo$; elíjase $\ma$ según 4 de modo que $\mo b=\mb\ma$ y $(\ma,\mc)=\mo$. Entonces
\[
        \mo n=\ma\mb:\ma\mc=\mo b:\ma\mc,
\]
y como $\ma\mc=\mo b:\mo n$, el ideal $\ma\mc$ es también principal; escribamos
\[
        \mo n=\mo b:\mo c.
\]
Ninguna potencia de $b$ es divisible por $c$, porque $(\ma,\mc)=\mo$ y $(\mb,\mc)=\mo$.

7. \emph{Demostración del axioma V, clausura íntegra en el cuerpo de cocientes.} Por 6\,$\gamma$, todo elemento no íntegro de $\mK$ tiene una representación como cociente
\[
        \eta=a/c
\]
tal que, en $\mR$, ninguna potencia del numerador es divisible por el denominador. Tal elemento no puede satisfacer una ecuación que lo caracterice como íntegro sobre $\mR$; pues de
\[
        \eta^n+r_1\eta^{n-1}+\cdots+r_n=0
\]
se obtendría $a^n\equiv0\pmod{\mo c}$ en $\mR$.

8. \emph{Consecuencia.} Si en un anillo conmutativo $\mR$ se satisfacen los axiomas I a IV y todo ideal primario es irreducible, entonces también se satisface el axioma V, la clausura íntegra en el cuerpo de cocientes. Debido a los axiomas I a III, toda potencia de un ideal primo $\mpideal^\varrho$ es primaria; en particular $\mpideal^2$ es primaria y por hipótesis irreducible. A partir de esto debe probarse
\[
        \mpideal=(\mo c,\mpideal^2).
\]
Entonces, por el resultado auxiliar §8, 3, todo ideal primario es una potencia de ideal primo; junto con las demás hipótesis, 7 da la clausura íntegra.

Por la condición de cadena de divisores,
\[
        \mpideal=(\mo c_1,\ldots,\mo c_n),
\]
donde sólo importan las clases residuales módulo $\mpideal$ como multiplicadores de los $c_i$, y los $c_i$ pueden suponerse linealmente independientes sobre el cuerpo de clases residuales módulo $\mpideal$. Si $n>1$, esta independencia lineal da una representación
\[
        \mpideal^2=[\mq_1,\mq_2],
        \qquad
        \mq_1=(\mo c_1,\mpideal^2),
        \qquad
        \mq_2=(\mo c_2,\ldots,\mo c_n,\mpideal^2),
\]
con $\mq_1$ y $\mq_2$ divisores propios de $\mpideal^2$; por tanto $\mpideal^2$ sería reducible, contra la hipótesis. Así $\mpideal=(\mo c,\mpideal^2)$, y la consecuencia afirmada queda probada.

Recíprocamente, como consecuencia de los axiomas I a V, todo ideal primario es inmediatamente irreducible: una potencia de ideal primo $\mpideal^\varrho$ no puede ser mínimo común múltiplo de divisores propios de la forma $\mpideal^\sigma$ y $\mpideal^\tau$.

9. Si se permiten divisores de cero en el anillo, entonces los axiomas I a III y la clausura íntegra en el anillo de cocientes no implican que todo ideal primario sea irreducible.
\footnote{Compárese la nota 18.}
Sea $\mT$ un anillo en el que valen todos los axiomas I a IV salvo la clausura íntegra; tales anillos existen, por ejemplo, los órdenes finitos distintos del sistema completo de cantidades íntegras en un cuerpo de extensión finita del cuerpo de cocientes de $\mR$, cuando en $\mR$ valen los axiomas I a V (§3, 2). Por la consecuencia de 8 hay en $\mT$ un ideal primario reducible $\mq$. Sea $\mpideal^\varrho$ un múltiplo propio de $\mq$, y sea $\mR$ el anillo de clases residuales $\mT/\mpideal^\varrho$, en el cual el ideal correspondiente a $\mq$ es un ideal primario reducible no nulo. En $\mR$ valen los axiomas I a III, y también la clausura íntegra en el anillo de cocientes. En efecto, todo elemento de $\mT$ que no es divisible por $\mpideal$ se vuelve coprimo con $\mpideal^\varrho$; por tanto todos los elementos regulares de $\mR$ son unidades, y $\mR$ coincide con su anillo de cocientes.

\subsection*{§10. Doble condición de cadena y serie de composición}

Ahora se muestra que, para dominios de módulos arbitrarios (§2) supuestos bien ordenados, la validez de la doble condición de cadena --toda cadena de divisores y toda cadena de múltiplos de módulos termina después de finitamente muchos pasos-- es equivalente a la existencia de una serie de composición.
\footnote{Los módulos del dominio son grupos abelianos respecto de la adición, de hecho grupos abelianos generalizados, pues el dominio de multiplicadores consta de los elementos de $\mR$. Como son grupos abelianos, el concepto de serie de composición coincide aquí con el de serie principal. Los teoremas de este parágrafo siguen siendo válidos si, por ``módulos'', se entiende la totalidad de los divisores normales de un grupo no conmutativo arbitrario, incluso generalizado. La notación, ligeramente distinta de la anterior, pretende recordar la teoría de grupos.}
Al especializar el dominio de módulos a un anillo conmutativo se obtienen los hechos correspondientes para el sistema de todos los ideales del anillo; en 2, en todas partes, el isomorfismo de módulos debe sustituirse entonces por isomorfismo de anillos.

Un dominio de módulos $\mG$ se llama \emph{simple} si en $\mG$ no hay $\mR$-módulos aparte de $\mG$ mismo y del módulo cero $\mE$. Un módulo $\mA$ se llama \emph{simple en $\mG$} si $\mG/\mA$ es simple.

Una cadena de múltiplos
\[
        \mG,\mU_1,\ldots,\mU_r,\mE
\]
se llama una serie de composición de $\mG$ de longitud $r$ si todos los módulos de la cadena son distintos y si cada módulo es simple en el precedente; equivalentemente, si los módulos de clases residuales, los grupos cociente, $\mU_{i-1}/\mU_i$ son dominios de módulos simples, al igual que $\mG/\mU_1$ y $\mU_r/\mE$. Formando el módulo de clases residuales $\mG/\mF$, el caso en que existe una serie de composición desde $\mG$ hasta $\mF\ne\mG$ se reduce al caso absoluto.

1. \emph{Si se supone la doble condición de cadena en $\mG$, entonces existe una serie de composición en $\mG$}, siempre que $\mG\ne\mE$. Del buen orden supuesto de $\mG$ y de la condición de cadena de divisores se obtiene, como en §6, mediante ordenación cuasi-lexicográfica, un buen orden del sistema de todos los módulos. Este buen orden, junto con la condición de cadena de divisores, da la existencia de al menos un módulo simple en $\mG$. En efecto, sea $\mG_1$ el primer divisor propio de $\mG$ en el buen orden; en general, sea $\mG_i$ el primer divisor propio de $\mG_{i-1}$. La cadena bien determinada
\[
        \mG,\mG_1,\ldots,\mG_i,\ldots
\]
termina después de finitamente muchos pasos y por tanto conduce necesariamente a un módulo $\mA$ simple en $\mG$. Si $\mU_1\ne\mE$, el mismo procedimiento da un módulo $\mU_2$ simple en $\mU_1$; en general, si $\mU_{i-1}\ne\mE$, hay un módulo $\mU_i$ simple en $\mU_{i-1}$. Esto da una cadena de múltiplos bien determinada
\[
        \mG,\mU_1,\ldots,\mU_i,
\]
que termina por la condición de cadena de múltiplos y es por tanto una serie de composición de $\mG$.

2. \emph{Si existe una serie de composición en $\mG$, entonces la doble condición de cadena vale en $\mG$.} La demostración es por inducción completa y no requiere un buen orden de $\mG$. En el curso de la inducción se obtiene al mismo tiempo una prueba del teorema de Jordan--Hölder bajo la única hipótesis de que existe una serie de composición.
\footnote{Compárese Masazo Sono (nota 21), \emph{On Congruences I}, §§11--13, para el caso de los anillos. Allí se trata también el análogo de una serie de composición en grupos no conmutativos: series $\mR,\mU_1,\ldots,\mU_r,\mE$, donde $\mU_i$ es un ideal y simple en $\mU_{i-1}$, sin que necesariamente sea un ideal en $\mR$. Los argumentos presentes siguen siendo válidos para grupos no conmutativos si por cadena de múltiplos se entiende una cadena en la que cada $\mU_i$ es normal en $\mU_{i-1}$ y las cadenas de divisores se definen correspondientemente. Compárese también Dedekind, ``Über die von drei Moduln erzeugte Dualgruppe'', Math. Ann. 53 (1900), pp. 371--403. Allí, para dominios mucho más generales, se prueba igualmente el teorema de Jordan--Hölder sólo bajo la hipótesis de existencia; en lugar del segundo teorema de isomorfía aparece una relación de correspondencia algo más débil, y por consiguiente el enunciado del teorema también es algo más débil. El principio de la demostración es idéntico al dado aquí.}

Si $\mG$ es simple, de modo que $r=0$ y la única serie de composición es $\mG,\mE$, entonces valen las siguientes afirmaciones.

$\alpha$) Por todo módulo simple en $\mG$ puede trazarse una serie de composición.

$\beta$) Teorema de Jordan--Hölder: toda serie de composición tiene la misma longitud, y el sistema de grupos cociente (módulos de clases residuales) coincide salvo el orden; los grupos cociente correspondientes son isomorfos.

$\gamma$) Por todo módulo distinto de $\mG$ puede trazarse una serie de composición.

$\delta$) Vale la condición de cadena de múltiplos.

$\varepsilon$) Vale la condición de cadena de divisores.

Supóngase por tanto que $\alpha$)--$\varepsilon$) se han demostrado para todo dominio de módulos que posee una serie de composición de longitud $r<r+1$. Las demostraremos bajo la hipótesis de que en $\mG$ hay una serie de composición
\[
        \mG,\mA,\mA_1,\ldots,\mA_r,\mE
\]
de longitud $r+1$.

2$\alpha$. Si no hay ningún módulo simple en $\mG$ aparte de $\mA$, no hay nada que demostrar. Sea $\mB$ simple en $\mG$ y $\mB\ne\mA$. Como $\mA$ y $\mB$ son ambos simples en $\mG$ y son distintos, necesariamente $(\mA,\mB)=\mG$. Pongamos $\mM=[\mA,\mB]$. Por el segundo teorema de isomorfía (§4, 2),
\[
        \mG/\mB\cong \mA/\mM,
        \qquad
        \mG/\mA\cong \mB/\mM.
\]
Así $\mM$ es simple en $\mA$ y en $\mB$. Como $\mA$ tiene una serie de composición de longitud $r<r+1$, las hipótesis de inducción $\alpha$ y $\delta$ dan en $\mA$ una serie de composición que pasa por $\mM$, digamos
\[
        \mA,\mM,\mM_1,\ldots,\mM_{r-1},\mE.
\]
Por consiguiente
\[
        \mG,\mB,\mM,\mM_1,\ldots,\mM_{r-1},\mE
\]
es una serie de composición de $\mG$ que pasa por $\mB$.

2$\beta$. Para probar el teorema de Jordan--Hölder, sean
\[
\begin{aligned}
&(1)\quad \mG,\mA,\mA_1,\ldots,\mA_r,\mE,\qquad\text{y}\qquad
(2)\quad \mG,\mB,\mB_1,\ldots,\mB_s,\mE
\end{aligned}
\]
dos series de composición de $\mG$. Si $\mA=\mB$, el resultado se sigue de la hipótesis de inducción para la longitud $r<r+1$. En caso contrario se compara con las series construidas en 2$\alpha$,
\[
        (3)\quad \mG,\mB,\mM,\mM_1,\ldots,\mM_{r-1},\mE,
        \qquad
        (4)\quad \mG,\mA,\mM,\mM_1,\ldots,\mM_{r-1},\mE.
\]
Por la hipótesis de inducción, el sistema de grupos cociente de (1) y (3) coincide; lo mismo ocurre con el de (2) y (4), puesto que (4), tras pasar al cociente adecuado, es una serie de composición de longitud $r$ que pasa por $\mB$. Así $s=r$. El isomorfismo de 2$\alpha$ da la coincidencia entre (3) y (4), y por tanto entre (1) y (2). Esto prueba el teorema de Jordan--Hölder.

2$\gamma$. Observación preliminar. Sean $\mA,\mB,\mH$ módulos de $\mG$, y supóngase que $\mB$ es simple en $\mA$. Entonces los módulos $(\mB,\mH)$ y $(\mA,\mH)$ coinciden, o bien $(\mB,\mH)$ es simple en $(\mA,\mH)$. Por el segundo teorema de isomorfía,
\[
        (\mA,\mB,\mH)/(\mB,\mH)
        \cong
        \mA/[\mA,(\mB,\mH)].
\]
Como $\mB\equiv0\pmod{\mA}$, esto es
\[
        (\mA,\mH)/(\mB,\mH)\cong \mA/\mC,
        \qquad \mC=[\mA,(\mB,\mH)].
\]
Aquí $\mC\equiv0\pmod{\mA}$; por otra parte $\mB\equiv0\pmod{\mC}$, porque $\mB\equiv0\pmod{\mA}$ y $\mB\equiv0\pmod{(\mB,\mH)}$. Por tanto $\mC$ se sitúa entre $\mA$ y $\mB$ y, por hipótesis, es idéntico a $\mA$ o a $\mB$. Esto demuestra la observación preliminar.

Sea ahora
\[
        \mG,\mU,\mU_1,\ldots,\mU_r,\mE
\]
una serie de composición de $\mG$, y sea $\mH$ un módulo arbitrario de $\mG$ distinto de $\mG$. Para mostrar que una serie de composición pasa por $\mH$, fórmese
\[
        \mG,(\mU,\mH),(\mU_1,\mH),\ldots,(\mU_r,\mH),\mH.
\]
Por la observación preliminar, los módulos distintos entre éstos forman una serie de composición desde $\mG$ hasta $\mH$. Así el primer término propio es simple en $\mG$ (éste es el caso especial 2$\alpha$ cuando coincide con $\mB$). Por 2$\alpha$ hay en $\mG$ una serie de composición de longitud $r<r+1$ que pasa por ese término, y por consiguiente, por la hipótesis de inducción, una serie de composición que pasa por $\mH$; añadiendo $\mG$ se obtiene una serie de composición de $\mG$ que pasa por $\mH$. Repitiendo el argumento se muestra que tal serie puede elegirse, en particular, como prolongación de la serie recién construida desde $\mG$ hasta $\mH$.

2$\delta$. Demostración de la condición de cadena de múltiplos. Supóngase de nuevo que existe una serie de composición de longitud $r+1$ en $\mG$, y sea
\[
        \mG,\mB,\mB_1,\ldots,\mB_i,\ldots
\]
una cadena de múltiplos, cada término un múltiplo propio del precedente. Que la cadena comience con $\mG$ no es una restricción de generalidad, pues $\mG$ siempre puede adjuntarse como primer término. Por 2$\gamma$ existe una serie de composición que pasa por $\mB$, y por tanto $\mB$ tiene una serie de composición de menor longitud. Por la hipótesis de inducción, la cadena
\[
        \mB,\mB_1,\ldots,\mB_i,
\]
termina después de finitamente muchos pasos. Por consiguiente lo mismo es cierto después de adjuntar $\mG$.

2$\varepsilon$. Demostración de la condición de cadena de divisores. Supóngase otra vez una serie de composición de longitud $r+1$ en $\mG$, y sea
\[
        \mG,\mC,\mC_1,\ldots,\mC_i,
\]
una cadena de divisores, cada término un divisor propio del precedente. Tampoco hay aquí restricción en suponer que la cadena comienza con $\mG$. Por 2$\gamma$ existe una serie de composición que pasa por $\mC$, y por tanto en $\mG/\mC$ hay una serie de composición de menor longitud. Por la hipótesis de inducción, la cadena de divisores
\[
        \mG/\mC,\ \mC_1/\mC,\ldots,\mC_i/\mC,
\]
termina después de finitamente muchos pasos. Por la correspondencia biunívoca dada por el primer teorema de isomorfía (§4, 2), lo mismo vale para la cadena original que comienza con $\mG$, y por tanto también para la cadena con $\mE$ adjuntado. Queda así probada la equivalencia de las dos hipótesis: existencia de una serie de composición y doble condición de cadena.

\emph{Recibido el 13 de agosto de 1925.}

\end{document}
